% bverfge_layout.tex
% 共用排版基礎設施，供 Schwangerschaftsabbruch I 與 II 共享。
% 由各文件以 \documentclass 後 % bverfge_layout.tex
% 共用排版基礎設施，供 Schwangerschaftsabbruch I 與 II 共享。
% 由各文件以 \documentclass 後 % bverfge_layout.tex
% 共用排版基礎設施，供 Schwangerschaftsabbruch I 與 II 共享。
% 由各文件以 \documentclass 後 % bverfge_layout.tex
% 共用排版基礎設施，供 Schwangerschaftsabbruch I 與 II 共享。
% 由各文件以 \documentclass 後 \input{../../共用LaTeX/bverfge_layout} 載入。
% 不含：\documentclass、\documentversion、\ggversionde/zh、\tocde/\toczh。
% 日期與首頁版本列由本檔自動產生；各文件只需手動指定 \documentversion。

\usepackage{xeCJK}
\usepackage{fontspec}
\usepackage{geometry}
\usepackage{setspace}
\usepackage{hyperref}
\usepackage{xurl}
\usepackage{bookmark}
\usepackage{enumitem}
\usepackage{xcolor}
\usepackage{titlesec}
\usepackage{array}
\usepackage{longtable}
\usepackage{tcolorbox}
\usepackage{environ}
\usepackage{pdflscape}
\usepackage{etoolbox}
\usepackage{ragged2e}
\usepackage{paracol}
\usepackage{multicol}
\usepackage{needspace}
\usepackage{fancyhdr}
\usepackage{tikz}
\usepackage{zref-abspage}
\tcbuselibrary{breakable,skins}
\usetikzlibrary{calc}
\usepackage{pgfcalendar}

\hypersetup{
  unicode=true,
  bookmarksopen=true,
  bookmarksnumbered=false,
  hidelinks
}

% ===== 自動推導，不要手動改下面 =====
% （\docmode 與 \bilingualpaper 由各文件在 \input{bverfge_layout} 之前設定；
%   預設 chinese / a4）
\ifdefined\docmode\else\def\docmode{chinese}\fi
\ifdefined\bilingualpaper\else\def\bilingualpaper{a4}\fi
\newif\ifshoworiginal
\newif\ifshowtranslation
\newif\ifbilinguallandscape
\newif\ifbilingualcolumns
\newif\ifintranslation
\newif\iffirstbilingualsection

\showoriginalfalse
\showtranslationfalse
\bilinguallandscapefalse
\bilingualcolumnsfalse
\intranslationfalse
\firstbilingualsectionfalse

\def\modechinese{chinese}
\def\modegerman{german}
\def\modebilingual{bilingual}
\def\bilingualpaperaThree{a3}
\def\bilingualpaperaFour{a4}

\ifx\docmode\modechinese
  \showtranslationtrue
\fi

\ifx\docmode\modegerman
  \showoriginaltrue
\fi

\ifx\docmode\modebilingual
  \showoriginaltrue
  \showtranslationtrue
  \bilinguallandscapetrue
  \bilingualcolumnstrue
\fi

% 編排模式說明：
% chinese  → \showoriginalfalse  \showtranslationtrue  （A4 直式，純中文）
% german   → \showoriginaltrue   \showtranslationfalse （A4 直式，純德文）
% bilingual→ \showoriginaltrue   \showtranslationtrue  \bilingualcolumnstrue（A3/A4 橫式對照）

\ifbilingualcolumns
  \ifx\bilingualpaper\bilingualpaperaFour
    \geometry{
      a4paper,
      landscape,
      left=1.25cm,
      right=2.25cm,
      top=1.25cm,
      bottom=1.75cm,
      footskip=8.5mm,
      marginparwidth=0.75cm,
      marginparsep=0.25cm
    }
  \else
    \geometry{
      a3paper,
      landscape,
      left=1.55cm,
      right=2.75cm,
      top=1.55cm,
      bottom=2.05cm,
      footskip=9.5mm,
      marginparwidth=0.9cm,
      marginparsep=0.35cm
    }
  \fi
\else
\geometry{
  a4paper,
  portrait,
  left=2.2cm,
  right=2.6cm,
  top=2.0cm,
  bottom=2.0cm,
  marginparwidth=0.8cm,
  marginparsep=0.3cm
}
\fi
% ===== 時間戳基礎設施 =====
\newcount\stamptimeval
\newcount\stampminuteval
\newcommand{\stamphh}{%
  \stamptimeval=\time
  \divide\stamptimeval by 60
  \ifnum\stamptimeval<10 0\fi
  \the\stamptimeval}
\newcommand{\stampmm}{%
  \stamptimeval=\time
  \stampminuteval=\time
  \divide\stampminuteval by 60
  \multiply\stampminuteval by 60
  \advance\stamptimeval by -\stampminuteval
  \ifnum\stamptimeval<10 0\fi
  \the\stamptimeval}
\newcommand{\stampwkde}[1]{%
  \ifcase#1Montag\or Dienstag\or Mittwoch\or Donnerstag\or Freitag\or Samstag\or Sonntag\fi}
\newcommand{\stampwkzh}[1]{%
  \ifcase#1 一\or 二\or 三\or 四\or 五\or 六\or 日\fi}
\newcommand{\stampmonthde}[1]{%
  \ifcase#1\or Januar\or Februar\or März\or April\or Mai\or Juni\or
  Juli\or August\or September\or Oktober\or November\or Dezember\fi}
\newcount\stampjdval
\newcount\stampwdval
\pgfcalendardatetojulian{\the\year-\the\month-\the\day}{\stampjdval}
\pgfcalendarjuliantoweekday{\the\stampjdval}{\stampwdval}
\newcommand{\timestampzh}{%
  \songfont 最後修改：\the\year 年\the\month 月\the\day 日%
  % （星期\stampwkzh{\the\stampwdval}）%
  % \space\stamphh:\stampmm
  }

\newcommand{\documentdatezh}{\the\year 年 \the\month 月 \the\day 日}
\newcommand{\documentdatede}{\the\day. \stampmonthde{\the\month} \the\year}
\newcommand{\bverfgetranslationnote}{%
  {\songfont 翻譯說明：初譯使用 Claude Code 與 GPT 協助迭代；仍請以德文原文為準。}%
}
\newcommand{\bverfgecourseinfo}{%
  {\songfont 課程：114 學年度第 2 學期；憲法專題研究（四）；授課老師：湯德宗教授；導讀日期：115 年 6 月 1 日。}%
}
\newcommand{\bverfgeeditorinfo}{%
  {\songfont 編者：王逸帆（R10A21126），國立台灣大學法律系研究所碩士班。}%
}
\newcommand{\bverfgetitleversionline}{%
  \ifbilingualcolumns
    Fassung \documentversion\enspace\textperiodcentered\enspace Stand: \documentdatede
    \quad
    {\songfont 版本 \documentversion\enspace\textperiodcentered\enspace 修訂日期：\documentdatezh\par}%
  \else
    \ifshowtranslation
      {\songfont 版本 \documentversion\enspace\textperiodcentered\enspace 修訂日期：\documentdatezh\par}%
    \else
      Fassung \documentversion\enspace\textperiodcentered\enspace Stand: \documentdatede\par
    \fi
  \fi
}
\newcommand{\bverfgetitlecornerinfo}{%
  \ifshowtranslation
    \vspace{0.72em}%
    \noindent\hfill
    \begin{minipage}{0.66\textwidth}
      \raggedleft\tiny\color{pendingink}\songfont
      \bverfgecourseinfo\par
      \bverfgeeditorinfo
    \end{minipage}%
  \fi
}
\newcommand{\bverfgetitlemeta}{%
  \vfill
  \begin{center}%
    \footnotesize\color{pendingink}%
    \bverfgetitleversionline
    \ifshowtranslation
      \vspace{0.28em}{\scriptsize\bverfgetranslationnote\par}%
    \fi
  \end{center}%
  \bverfgetitlecornerinfo
}

\pagestyle{fancy}
\fancyhf{}
\renewcommand{\headrulewidth}{0pt}
\renewcommand{\footrulewidth}{0pt}
\fancyfoot[C]{\thepage}
\ifbilingualcolumns
  \fancyfoot[R]{\scriptsize\color{pendingink}\timestampzh}%
\else
  \ifshoworiginal
  \else
    \fancyfoot[R]{\scriptsize\color{pendingink}\timestampzh\hspace{0.45cm}}%
  \fi
\fi
\setmainfont{Times New Roman}
\setsansfont{Helvetica Neue}
\setmonofont{Menlo}
\IfFileExists{/Users/iw/Library/Fonts/kaiu.ttf}{
  \setCJKmainfont[Path=/Users/iw/Library/Fonts/,AutoFakeBold=4]{kaiu.ttf}
}{
  \IfFontExistsTF{標楷體}{
    \setCJKmainfont[AutoFakeBold=4]{標楷體}
  }{
    \setCJKmainfont{Songti TC}
  }
}
\IfFontExistsTF{Songti TC}{
  \setCJKfamilyfont{song}{Songti TC}
  \setCJKmonofont{Songti TC}
}{
  \setCJKfamilyfont{song}[Path=/Users/iw/Library/Fonts/,AutoFakeBold=4]{kaiu.ttf}
  \setCJKmonofont[Path=/Users/iw/Library/Fonts/,AutoFakeBold=4]{kaiu.ttf}
}
\newcommand{\songfont}{\CJKfamily{song}}

% ===== 封面／目錄專用打字機字體（僅限封面、目錄頁使用，不影響正文）=====
\IfFontExistsTF{Radio Newsman}{%
  \newfontfamily\bverfgetitlede{Radio Newsman}%
}{%
  \newfontfamily\bverfgetitlede{Courier New}%
}
\IfFontExistsTF{Erikas Farbband}{%
  \newfontfamily\bverfgebodyde{Erikas Farbband}%
}{%
  \let\bverfgebodyde\bverfgetitlede
}
\IfFontExistsTF{Maoken Heavy Labourer}{%
  \setCJKfamilyfont{bverfgetitlecjk}{Maoken Heavy Labourer}%
}{%
  \setCJKfamilyfont{bverfgetitlecjk}{Songti TC}%
}
\IfFontExistsTF{Huiwen-mincho}{%
  \setCJKfamilyfont{bverfgebodycjk}{Huiwen-mincho}%
}{%
  \setCJKfamilyfont{bverfgebodycjk}{Songti TC}%
}
\newcommand{\bverfgetitlezh}{\bverfgetitlede\CJKfamily{bverfgetitlecjk}}
\newcommand{\bverfgebodyzh}{\bverfgebodyde\CJKfamily{bverfgebodycjk}}

\setstretch{1.35}
\setlist{itemsep=0.2em, topsep=0.4em}
\setlength{\parindent}{0pt}
\setlength{\parskip}{0.45em}
\raggedbottom
\setcounter{secnumdepth}{0}
\setcounter{tocdepth}{2}

\newcommand{\lawboxsourcefont}{\footnotesize}
\newcommand{\lawboxtransfont}{\footnotesize}
\newcommand{\lawboxsourcestretch}{1.02}
\newcommand{\lawboxtransstretch}{0.92}

\titleformat{\section}{\centering\Large\bfseries\songfont}{\thesection}{1em}{}
\titleformat{\subsection}{\large\bfseries\songfont}{\thesubsection}{1em}{}
\titleformat{\subsubsection}{\normalsize\bfseries\songfont}{\thesubsubsection}{1em}{}
\titleformat{\paragraph}{\normalsize\bfseries\songfont}{\theparagraph}{1em}{}
\titlespacing*{\section}{0pt}{1.8em}{1.0em}
\titlespacing*{\subsection}{0pt}{1.2em}{0.6em}
\titlespacing*{\subsubsection}{0pt}{1.0em}{0.45em}
\titlespacing*{\paragraph}{0pt}{0.6em}{0.4em}

\definecolor{noteblue}{HTML}{A9C9F5}
\definecolor{noteteal}{HTML}{AEE1D6}
\definecolor{notegray}{HTML}{D7DADF}
\definecolor{caseink}{HTML}{000000}
\definecolor{casepaper}{HTML}{FCFEFD}
\definecolor{sourcepaper}{HTML}{FAFBF8}
\definecolor{sourceline}{HTML}{D7DED8}
\definecolor{sourceink}{HTML}{000000}
\definecolor{pendingink}{HTML}{000000}

\tcbset{
  caseboxbase/.style={
    enhanced,
    breakable,
    width=0.985\linewidth,
    colback=casepaper,
    coltitle=caseink,
    fonttitle=\bfseries\songfont,
    boxrule=0.28pt,
    arc=1.2mm,
    left=2.4mm,
    right=2.4mm,
    top=1.5mm,
    bottom=1.5mm,
    boxsep=1.5mm,
    before skip=0.65em,
    after skip=0.65em,
    before upper={\setlength{\parindent}{0pt}\setlength{\parskip}{0.35em}},
  }
}

\newcommand{\bilingualstart}{}
\newcommand{\bilingualstop}{}
\ifbilingualcolumns
  \renewcommand{\bilingualstart}{\small\setlength{\columnsep}{1.25cm}\setlength{\columnseprule}{0.12pt}\columnratio{0.5,0.5}\multicolumnfootnotes\flushbottom\sloppy\interlinepenalty=80\clubpenalty=800\widowpenalty=800\displaywidowpenalty=800\begin{paracol}{2}\global\intranslationfalse\global\firstbilingualsectiontrue{\footnotesize\color{sourceink}Deutsch\par}\switchcolumn{\footnotesize\songfont\color{sourceink}中文譯文\par}\switchcolumn*}
  \renewcommand{\bilingualstop}{\ifintranslation\switchcolumn*\global\intranslationfalse\fi\end{paracol}\clearpage\raggedbottom}
\fi
\newif\ifrandnumbers
\randnumbersfalse
\newcounter{randnumber}
\newcounter{randnumberlocal}
\newcounter{lastsourcerandstart}
\newif\iflastsourcehasrandnumbers
\lastsourcehasrandnumbersfalse
\newif\ifinlawbox
\inlawboxfalse
\newif\iflawboxhaspendingrn
\lawboxhaspendingrnfalse
\newcommand{\lawboxpendingrn}{}
\newcommand{\startrandnumbers}{\global\randnumberstrue\setcounter{randnumber}{1}\global\lastsourcehasrandnumbersfalse}
\newcommand{\stoprandnumbers}{\global\randnumbersfalse\global\lastsourcehasrandnumbersfalse}
\newlength{\randnumberwidth}
\setlength{\randnumberwidth}{0.95cm}
\newlength{\randnumberrightoffset}
\setlength{\randnumberrightoffset}{0.85cm}
\newlength{\randnumberlawboxrightoffset}
\setlength{\randnumberlawboxrightoffset}{1.40cm}
\newcommand{\randnumberbox}[1]{%
  \leavevmode
  \makebox[0pt][l]{%
    \smash{%
      \tikz[remember picture,overlay,baseline=(rnmark.base)]{%
        \coordinate (rnmark) at (0,0);%
        \ifinlawbox
          \path let \p1=(current page.east), \p2=(rnmark.base) in
            node[
              anchor=base east,
              inner sep=0pt,
              outer sep=0pt,
              text width=\randnumberwidth,
              align=right,
              text=pendingink,
              font=\normalfont\mdseries\scriptsize\ttfamily
            ] at (\x1-\randnumberlawboxrightoffset,\y2) {{\normalfont\mdseries\scriptsize\ttfamily #1}};%
        \else
          \path let \p1=(current page.east), \p2=(rnmark.base) in
            node[
              anchor=base east,
              inner sep=0pt,
              outer sep=0pt,
              text width=\randnumberwidth,
              align=right,
              text=pendingink,
              font=\normalfont\mdseries\scriptsize\ttfamily
            ] at (\x1-\randnumberrightoffset,\y2) {{\normalfont\mdseries\scriptsize\ttfamily #1}};%
        \fi
      }%
    }%
  }%
  \ignorespaces
}
\newcommand{\lawboxstorerandnumber}[1]{%
  \gdef\lawboxpendingrn{#1}%
  \global\lawboxhaspendingrntrue
  \ignorespaces
}
\newcommand{\lawboxuserandnumber}{%
  \iflawboxhaspendingrn
    \randnumberbox{\lawboxpendingrn}%
    \global\lawboxhaspendingrnfalse
  \fi
}
\newcommand{\sourcern}[1]{%
  \ifrandnumbers
    \ifshoworiginal
      \ifshowtranslation\else
        \ifinlawbox\lawboxstorerandnumber{#1}\else\randnumberbox{#1}\fi
      \fi
    \fi
  \fi
}
\newcommand{\transrn}[1]{%
  \ifrandnumbers
    \ifshowtranslation
      \ifinlawbox\lawboxstorerandnumber{#1}\else\randnumberbox{#1}\fi
    \fi
  \fi
}
% ===== 課堂模式：德文原文縮寫註解 =====
% 日常切換只改各判決檔的一行：
%   \classroommodeon        開啟課堂版；註解掉這行即回到一般版。
% 模板預設：
%   1. 課堂模式關閉。
%   2. 課堂模式開啟時，縮寫註解包含「德文全稱＋中文譯名」。
%   3. 課堂模式開啟時，GG/條文編者註仍會自動顯示。
% 進階微調才需要在單案檔覆寫：
%   \classroomabbrzhoff     縮寫註解僅顯示德文全稱。
%   \showeditornotestrue    一般版也顯示編者註。
% 註解策略：同一實際 PDF 頁面中，同一縮寫只在第一次出現處加註。
% 這裡用 zref-abspage 的絕對頁序，不用 \thepage，避免文件內頁碼重置時誤判。
\newif\ifclassroommode
\newif\ifclassroomabbrzh
\classroommodefalse
\classroomabbrzhtrue
\newcommand{\classroommodeon}{\classroommodetrue}
\newcommand{\classroommodeoff}{\classroommodefalse}
\newcommand{\classroomabbrzhon}{\classroomabbrzhtrue}
\newcommand{\classroomabbrzhoff}{\classroomabbrzhfalse}
\newcommand{\abbrnotelabel}{\emph{Abk.}: }
\newcommand{\abbrnote}[3]{%
  \ifclassroommode
    \ifshoworiginal
      \ifintranslation\else
        \footnote{\normalfont\footnotesize\abbrnotelabel #2\ifclassroomabbrzh；{\songfont #3}\fi}%
      \fi
    \fi
  \fi
}
\newcommand{\abbrnoteonce}[4]{%
  #2%
  \ifclassroommode
    \ifshoworiginal
      \ifintranslation\else
        \begingroup
          \edef\abbrcurrentabspage{\number\numexpr\value{abspage}+1\relax}%
          \ifcsname abbrnoteseen@#1@\abbrcurrentabspage\endcsname\else
            \global\expandafter\let\csname abbrnoteseen@#1@\abbrcurrentabspage\endcsname\empty
            \abbrnote{#1}{#3}{#4}%
          \fi
        \endgroup
      \fi
    \fi
  \fi
}
\newcommand{\abbrAbs}{\abbrnoteonce{abs}{Abs.}{Absatz}{項}}
\newcommand{\abbrAAO}{\abbrnoteonce{aao}{a.a.O.}{am angegebenen Ort}{前引處}}
\newcommand{\abbrAE}{\abbrnoteonce{ae}{AE}{Alternativ-Entwurf}{替代草案}}
\newcommand{\abbrArt}{\abbrnoteonce{art}{Art.}{Artikel}{條}}
\newcommand{\abbrAufl}{\abbrnoteonce{aufl}{Aufl.}{Auflage}{版}}
\newcommand{\abbrBd}{\abbrnoteonce{bd}{Bd.}{Band}{卷}}
\newcommand{\abbrBGBl}{\abbrnoteonce{bgbl}{BGBl.}{Bundesgesetzblatt}{聯邦法律公報}}
\newcommand{\abbrBGHSt}{\abbrnoteonce{bghst}{BGHSt}{Entscheidungen des Bundesgerichtshofes in Strafsachen}{聯邦普通法院刑事判決彙編}}
\newcommand{\abbrBRDrucks}{\abbrnoteonce{brdrucks}{BRDrucks.}{Bundesrats-Drucksache}{聯邦參議院文件}}
\newcommand{\abbrBTDrucks}{\abbrnoteonce{btdrucks}{BTDrucks.}{Bundestags-Drucksache}{聯邦議院文件}}
\newcommand{\abbrBundesgesetzbl}{\abbrnoteonce{bundesgesetzbl}{Bundesgesetzbl.}{Bundesgesetzblatt}{聯邦法律公報}}
\newcommand{\abbrBvF}{\abbrnoteonce{bvf}{BvF}{Bundesverfassungsgerichtsverfahren}{聯邦憲法法院程序案號}}
\newcommand{\abbrBvL}{\abbrnoteonce{bvl}{BvL}{Bundesverfassungsgerichtsverfahren}{聯邦憲法法院程序案號}}
\newcommand{\abbrBvR}{\abbrnoteonce{bvr}{BvR}{Bundesverfassungsgerichtsverfahren}{聯邦憲法法院程序案號}}
\newcommand{\abbrBVerfG}{\abbrnoteonce{bverfg}{BVerfG}{Bundesverfassungsgericht}{聯邦憲法法院}}
\newcommand{\abbrBVerfGE}{\abbrnoteonce{bverfge}{BVerfGE}{Entscheidungen des Bundesverfassungsgerichts}{聯邦憲法法院判決集}}
\newcommand{\abbrBVerfGG}{\abbrnoteonce{bverfgg}{BVerfGG}{Bundesverfassungsgerichtsgesetz}{聯邦憲法法院法}}
\newcommand{\abbrDDR}{\abbrnoteonce{ddr}{DDR}{Deutsche Demokratische Republik}{德意志民主共和國；東德}}
\newcommand{\abbrDJT}{\abbrnoteonce{djt}{DJT}{Deutscher Juristentag}{德國法學家大會}}
\newcommand{\abbrDr}{\abbrnoteonce{dr}{Dr.}{Doktor}{Dr.}}
\newcommand{\abbrEuGRZ}{\abbrnoteonce{eugrz}{EuGRZ}{Europäische Grundrechte-Zeitschrift}{歐洲基本權期刊}}
\newcommand{\abbrEV}{\abbrnoteonce{ev}{EV}{Einigungsvertrag}{統一條約}}
\newcommand{\abbrF}{\abbrnoteonce{f}{f.}{folgende}{以下一頁；下一段}}
\newcommand{\abbrFF}{\abbrnoteonce{ff}{ff.}{fortfolgende}{以下數頁；以下各段}}
\newcommand{\abbrGBlDDR}{\abbrnoteonce{gblddr}{GBl. DDR}{Gesetzblatt der Deutschen Demokratischen Republik}{德意志民主共和國法律公報}}
\newcommand{\abbrGG}{\abbrnoteonce{gg}{GG}{Grundgesetz}{基本法}}
\newcommand{\abbrGVBl}{\abbrnoteonce{gvbl}{GVBl.}{Gesetz- und Verordnungsblatt}{法律及命令公報}}
\newcommand{\abbrJR}{\abbrnoteonce{jr}{JR}{Juristische Rundschau}{法學評論}}
\newcommand{\abbrJZ}{\abbrnoteonce{jz}{JZ}{JuristenZeitung}{法學家報}}
\newcommand{\abbrIVm}{\abbrnoteonce{ivm}{i.V.m.}{in Verbindung mit}{結合；連同}}
\newcommand{\abbrMdB}{\abbrnoteonce{mdb}{MdB}{Mitglied des Bundestages}{聯邦議院議員}}
\newcommand{\abbrMwN}{\abbrnoteonce{mwn}{m.w.N.}{mit weiteren Nachweisen}{附進一步文獻／判例出處}}
\newcommand{\abbrNF}{\abbrnoteonce{nf}{n.F.}{neue Fassung}{新版本}}
\newcommand{\abbrAF}{\abbrnoteonce{af}{a.F.}{alte Fassung}{舊版本}}
\newcommand{\abbrNr}{\abbrnoteonce{nr}{Nr.}{Nummer}{款、目或號，依語境}}
\newcommand{\abbrProf}{\abbrnoteonce{prof}{Prof.}{Professor}{教授}}
\newcommand{\abbrRdnr}{\abbrnoteonce{rdnr}{Rdnr.}{Randnummer}{邊碼；段碼}}
\newcommand{\abbrRdnrn}{\abbrnoteonce{rdnrn}{Rdnrn.}{Randnummern}{邊碼；段碼}}
\newcommand{\abbrRGBl}{\abbrnoteonce{rgbl}{RGBl.}{Reichsgesetzblatt}{帝國法律公報}}
\newcommand{\abbrRGSt}{\abbrnoteonce{rgst}{RGSt}{Entscheidungen des Reichsgerichts in Strafsachen}{帝國法院刑事判決彙編}}
\newcommand{\abbrRspr}{\abbrnoteonce{rspr}{Rspr.}{Rechtsprechung}{判例；司法實務}}
\newcommand{\abbrRVO}{\abbrnoteonce{rvo}{RVO}{Reichsversicherungsordnung}{帝國保險法}}
\newcommand{\abbrS}{\abbrnoteonce{s}{S.}{Seite}{頁}}
\newcommand{\abbrSFHG}{\abbrnoteonce{sfhg}{SFHG}{Schwangeren- und Familienhilfegesetz}{孕婦及家庭扶助法}}
\newcommand{\abbrSGBV}{\abbrnoteonce{sgbv}{SGB V}{Fünftes Buch Sozialgesetzbuch}{社會法典第五編}}
\newcommand{\abbrStAG}{\abbrnoteonce{staeg}{StÄG}{Strafrechtsänderungsgesetz}{刑法修正法}}
\newcommand{\abbrStenBer}{\abbrnoteonce{stenber}{StenBer.}{Stenographischer Bericht}{速記紀錄}}
\newcommand{\abbrStGB}{\abbrnoteonce{stgb}{StGB}{Strafgesetzbuch}{刑法}}
\newcommand{\abbrStREG}{\abbrnoteonce{streg}{StREG}{Strafrechtsreform-Ergänzungsgesetz}{刑法改革補充法}}
\newcommand{\abbrStrRG}{\abbrnoteonce{strrg}{StrRG}{Strafrechtsreformgesetz}{刑法改革法}}
\newcommand{\abbrUS}{\abbrnoteonce{us}{U.S.}{United States Reports}{美國最高法院判決彙編}}
\newcommand{\abbrVgl}{\abbrnoteonce{vgl}{vgl.}{vergleiche}{參見}}
\newcommand{\abbrVVDStRL}{\abbrnoteonce{vvdstrl}{VVDStRL}{Veröffentlichungen der Vereinigung der Deutschen Staatsrechtslehrer}{德國公法教師協會論文集}}
\newcommand{\abbrWP}{\abbrnoteonce{wp}{WP.}{Wahlperiode}{屆期}}
\newcommand{\abbrWp}{\abbrnoteonce{wp}{Wp.}{Wahlperiode}{屆期}}
\newcommand{\abbrZStrW}{\abbrnoteonce{zstrw}{ZStrW}{Zeitschrift für die gesamte Strafrechtswissenschaft}{整體刑法學期刊}}
\newcommand{\recordsourcerandstart}{%
  \ifrandnumbers
    \setcounter{lastsourcerandstart}{\value{randnumber}}%
    \global\lastsourcehasrandnumberstrue
  \else
    \global\lastsourcehasrandnumbersfalse
  \fi
}
\newlength{\biparagraphskip}
\setlength{\biparagraphskip}{0.52em}
\newlength{\lawboxblockbeforeskip}
\setlength{\lawboxblockbeforeskip}{0.72em}
\newlength{\lawboxblockafterskip}
\setlength{\lawboxblockafterskip}{0.92em}
\newif\iflawboxpartendedblock
\lawboxpartendedblockfalse
\newcommand{\sourceparstyle}{\footnotesize\color{sourceink}\setstretch{1.12}\RaggedRight\emergencystretch=3em\setlength{\parskip}{0.42em}\obeylines\selectfont}
\newcommand{\transparstyle}{\songfont\color{caseink}\setstretch{1.12}\setlength{\parindent}{0pt}\setlength{\parskip}{0.42em}}
\newcommand{\bilingualpairskip}{%
  \par\addvspace{\biparagraphskip}%
  \switchcolumn
  \par\addvspace{\biparagraphskip}%
  \switchcolumn*%
  \global\intranslationfalse
}
\newcommand{\bilingualboxpairskip}{%
  \par
  \switchcolumn
  \par
  \switchcolumn*%
  \global\intranslationfalse
}
\newcommand{\bilinguallawboxblockskip}{%
  \switchcolumn*%
  \par\addvspace{\lawboxblockafterskip}%
  \switchcolumn
  \par\addvspace{\lawboxblockafterskip}%
  \switchcolumn*%
  \global\intranslationfalse
}
\newcommand{\bikeepwithnext}[1][4]{%
  \ifbilingualcolumns
    \syncleft
    \Needspace{#1\baselineskip}%
    \switchcolumn
    \Needspace{#1\baselineskip}%
    \switchcolumn*%
    \global\intranslationfalse
  \else
    \Needspace{#1\baselineskip}%
  \fi
}
\NewEnviron{sourcepara}[1][]{
  \recordsourcerandstart
  \ifshoworiginal
    \ifbilingualcolumns
      \ifintranslation\switchcolumn*\global\intranslationfalse\fi
    \else
      \Needspace{5\baselineskip}
    \fi
    \ifbilingualcolumns\else\par\addvspace{0.35em}\fi
    {\sourceparstyle
    \noindent\BODY\par}
    \ifbilingualcolumns\else\addvspace{0.25em}\fi
    \ifbilingualcolumns
      \switchcolumn
      \global\intranslationtrue
    \fi
  \else
    \ifrandnumbers
      \begingroup
        \setbox0=\vbox{\hsize=\linewidth\sourceparstyle\noindent\BODY\par}%
      \endgroup
    \fi
  \fi
}
\NewEnviron{transpara}{
  \ifshowtranslation
    \setcounter{randnumberlocal}{\value{lastsourcerandstart}}%
    \ifbilingualcolumns
      \ifintranslation\else\switchcolumn\global\intranslationtrue\fi
      {\transparstyle\setlength{\leftskip}{0.55em}\setlength{\rightskip}{0.15em plus 1em}\addtolength{\linewidth}{-0.55em}\BODY\par}
      \switchcolumn*
      \bilingualpairskip
    \else
      {\transparstyle\BODY\par}
    \fi
  \else
    \ifbilingualcolumns
      \ifintranslation\switchcolumn*\global\intranslationfalse\fi
    \fi
  \fi
}
\NewEnviron{sourceboxpara}[1][]{
  \global\lawboxpartendedblockfalse
  \recordsourcerandstart
  \ifshoworiginal
    \ifbilingualcolumns
      \ifintranslation\switchcolumn*\global\intranslationfalse\fi
    \else
      \Needspace{3\baselineskip}
    \fi
    {\sourceparstyle
    \BODY\par}
    \ifbilingualcolumns\else
      \iflawboxpartendedblock\addvspace{\lawboxblockafterskip}\fi
    \fi
    \ifbilingualcolumns
      \switchcolumn
      \global\intranslationtrue
    \fi
  \fi
}
\NewEnviron{transboxpara}{
  \global\lawboxpartendedblockfalse
  \ifshowtranslation
    \setcounter{randnumberlocal}{\value{lastsourcerandstart}}%
    \ifbilingualcolumns
      \ifintranslation\else\switchcolumn\global\intranslationtrue\fi
      {\transparstyle\BODY\par}
      \iflawboxpartendedblock
        \bilinguallawboxblockskip
      \else
        \switchcolumn*
        \bilingualboxpairskip
      \fi
    \else
      {\transparstyle\BODY\par}
      \iflawboxpartendedblock\addvspace{\lawboxblockafterskip}\fi
    \fi
  \else
    \ifbilingualcolumns
      \ifintranslation\switchcolumn*\global\intranslationfalse\fi
    \fi
  \fi
}
\newcommand{\leitsatznum}[1]{\noindent\hangindent=3.0em\hangafter=1\makebox[2.25em][r]{\bfseries #1.}\hspace{0.55em}\ignorespaces}
\newcommand{\syncleft}{\ifbilingualcolumns\ifintranslation\switchcolumn*\global\intranslationfalse\fi\fi}
\pretocmd{\section}{\syncleft\clearpage}{}{}
\pretocmd{\subsection}{\syncleft}{}{}
\pretocmd{\subsubsection}{\syncleft}{}{}
\newcommand{\biheadingfont}{\songfont}
\newcommand{\bitocentry}[2]{%
  \ifshoworiginal
    \ifshowtranslation
      \ifstrequal{#1}{#2}{#1}{#1\space #2}%
    \else
      #1%
    \fi
  \else
    #2%
  \fi
}
\pdfstringdefDisableCommands{%
  \def\bitocentry#1#2{#1}%
}
\newcommand{\biheading}[7]{%
  \syncleft#1\bikeepwithnext[#3]\phantomsection\addcontentsline{toc}{#2}{\bitocentry{#6}{#7}}%
  \ifbilingualcolumns
    {#4 #6\par}%
    \switchcolumn
    {#4\biheadingfont #7\par}%
    \switchcolumn*%
    \global\intranslationfalse
    \par\addvspace{#5}%
    \switchcolumn
    \par\addvspace{#5}%
    \switchcolumn*%
  \else
    \ifshoworiginal{#4 #6\par}\fi
    \ifshowtranslation{#4\biheadingfont #7\par}\fi
    \addvspace{#5}%
  \fi
}
\newcommand{\termsectionheading}[1]{%
  \par\Needspace{9\baselineskip}%
  \vspace{0.65em}%
  {\songfont\bfseries #1\par}%
  \vspace{0.2em}%
}
% 章節換頁：
%   雙語 paracol 欄內不要用 \clearpage；它會在同步兩欄時吐出只有欄線/頁碼的空頁。
%   \newpage 足以讓下一個 section 起新頁，又不會強制清空所有待排材料。
%   單語模式不在 paracol 內，仍可用 \clearpage。
\newcommand{\bisectionpagebreak}{\ifbilingualcolumns\newpage\else\clearpage\fi}
\newcommand{\bisection}[2]{%
  \biheading{\iffirstbilingualsection\global\firstbilingualsectionfalse\else\bisectionpagebreak\fi}{section}{6}{\centering\Large\bfseries}{0.8em}{#1}{#2}%
}
\newcommand{\bisubsection}[2]{%
  \biheading{}{subsection}{5}{\large\bfseries}{0.55em}{#1}{#2}%
}
\newcommand{\bisubsubsection}[2]{%
  \biheading{}{subsubsection}{4}{\normalsize\bfseries}{0.45em}{#1}{#2}%
}
\newcommand{\bicompositeheading}[4]{%
  \biheading{}{subsection}{5}{\large\bfseries}{0.16em}{#1}{#2}%
  \biheading{}{subsubsection}{3}{\normalsize\bfseries}{0.45em}{#3}{#4}%
}
\newif\iflawboxfirstitem
\lawboxfirstitemfalse
\newcommand{\lawboxprespace}[1]{%
  \par
  \iflawboxfirstitem
    \global\lawboxfirstitemfalse
  \else
    \addvspace{#1}%
  \fi
}
\newtcolorbox{lawboxinner}{
  caseboxbase,
  colframe=sourceline!85!black,
  colback=white,
  left=1.05em,
  right=1.05em,
  top=0.62em,
  bottom=0.62em,
  before upper={\global\inlawboxtrue\global\lawboxfirstitemtrue\global\lawboxhaspendingrnfalse\ifintranslation\lawboxtransfont\else\lawboxsourcefont\fi\RaggedRight\setlength{\parindent}{0pt}\setlength{\parskip}{0pt}\ifintranslation\setstretch{\lawboxtransstretch}\else\setstretch{\lawboxsourcestretch}\fi},
  after upper={\global\lawboxhaspendingrnfalse\global\inlawboxfalse}
}
\tcbset{
  lawboxpartbase/.style={
    enhanced,
    breakable=false,
    width=0.985\linewidth,
    colback=white,
    frame hidden,
    boxrule=0pt,
    arc=0pt,
    left=1.05em,
    right=1.05em,
    boxsep=1.5mm,
    before skip=0pt,
    after skip=0pt,
    before upper={\global\inlawboxtrue\global\lawboxfirstitemtrue\global\lawboxhaspendingrnfalse\ifintranslation\lawboxtransfont\else\lawboxsourcefont\fi\RaggedRight\setlength{\parindent}{0pt}\setlength{\parskip}{0pt}\ifintranslation\setstretch{\lawboxtransstretch}\else\setstretch{\lawboxsourcestretch}\fi},
    after upper={\global\lawboxhaspendingrnfalse\global\inlawboxfalse},
    borderline west={0.28pt}{0pt}{sourceline!85!black},
    borderline east={0.28pt}{0pt}{sourceline!85!black},
  },
  lawboxpartfirst/.style={
    lawboxpartbase,
    top=0.62em,
    bottom=0.04em,
    before skip=\lawboxblockbeforeskip,
    borderline north={0.28pt}{0pt}{sourceline!85!black},
  },
  lawboxpartmiddle/.style={
    lawboxpartbase,
    top=0.04em,
    bottom=0.04em,
  },
  lawboxpartlast/.style={
    lawboxpartbase,
    top=0.04em,
    bottom=0.62em,
    after skip=0pt,
    borderline south={0.28pt}{0pt}{sourceline!85!black},
  }
}
\newtcolorbox{lawboxpartinner}[1]{#1}
\newtcolorbox{termboxinner}{
  caseboxbase,
  colframe=noteblue!70!black,
  colback=casepaper,
  before upper={\songfont\setlength{\parindent}{0pt}\setlength{\parskip}{0.35em}}
}
\newtcolorbox{deboxinner}{
  caseboxbase,
  colframe=notegray!72!black,
  colback=white,
  left=1.05em,
  right=1.05em,
  top=0.62em,
  bottom=0.62em,
  before upper={\global\inlawboxtrue\global\lawboxfirstitemtrue\global\lawboxhaspendingrnfalse\small\setlength{\parindent}{0pt}\setlength{\parskip}{0.35em}},
  after upper={\global\lawboxhaspendingrnfalse\global\inlawboxfalse}
}
\NewEnviron{lawbox}{
  \begin{lawboxinner}\BODY\end{lawboxinner}
}
\NewEnviron{lawboxpart}[2]{
  \begin{lawboxpartinner}{lawboxpart#1,equal height group=#2}\BODY\end{lawboxpartinner}%
  \ifstrequal{#1}{last}{\global\lawboxpartendedblocktrue}{}%
}
\NewEnviron{termbox}{
  \par\addvspace{0.55em}
  {\color{noteblue!70!black}\rule{\linewidth}{0.28pt}}\par
  \begingroup
    \songfont\small\color{caseink}
    \setlength{\parindent}{0pt}
    \setlength{\parskip}{0.24em}
    \BODY
    \par
  \endgroup
  {\color{noteblue!70!black}\rule{\linewidth}{0.28pt}}\par
  \addvspace{0.55em}
}
\NewEnviron{debox}{
  \begin{deboxinner}\BODY\end{deboxinner}
}
\providecommand{\paragraphmark}[1]{}
\newcommand{\lawboxtitleafter}{\addvspace{0.06em}}
\newcommand{\lawtitle}[1]{\lawboxprespace{0.22em}{\lawboxtransfont\bfseries\lawboxuserandnumber #1}\par\lawboxtitleafter}
\newcommand{\absno}[2]{\lawboxprespace{0.04em}\noindent\lawboxuserandnumber\hangindent=2.6em\hangafter=1\makebox[2.6em][l]{(#1)}#2\par}
\newcommand{\nrno}[2]{\lawboxprespace{0pt}\noindent\lawboxuserandnumber\hspace*{2.6em}\hangindent=4.4em\hangafter=1\makebox[1.8em][l]{#1.}#2\par}
\newcommand{\letterno}[2]{\lawboxprespace{0pt}\noindent\lawboxuserandnumber\hspace*{4.4em}\hangindent=6.0em\hangafter=1\makebox[1.6em][l]{#1)}#2\par}
\newcommand{\sourcelawtitle}[1]{\lawboxprespace{0.22em}{\lawboxsourcefont\bfseries\lawboxuserandnumber #1}\par\lawboxtitleafter}
\newcommand{\sourcelawsubtitle}[1]{\lawboxprespace{0.04em}{\lawboxsourcefont\bfseries\lawboxuserandnumber #1}\par\lawboxtitleafter}
\newcommand{\sourceabsno}[2]{\lawboxprespace{0.04em}\noindent\lawboxuserandnumber\hangindent=2.45em\hangafter=1\makebox[2.45em][l]{(#1)}#2\par}
\newcommand{\sourcenrno}[2]{\lawboxprespace{0pt}\noindent\lawboxuserandnumber\hspace*{2.45em}\hangindent=4.25em\hangafter=1\makebox[1.8em][l]{#1.}#2\par}
\newcommand{\sourceletterno}[2]{\lawboxprespace{0pt}\noindent\lawboxuserandnumber\hspace*{4.25em}\hangindent=5.85em\hangafter=1\makebox[1.6em][l]{#1)}#2\par}
\newcommand{\sourcequotepara}[1]{\lawboxprespace{0.04em}\noindent\lawboxuserandnumber\hangindent=1.2em\hangafter=1 #1\par}
\newcommand{\sourcelawline}[1]{\lawboxprespace{0.04em}\noindent\lawboxuserandnumber #1\par}
\newcommand{\lawline}[1]{\lawboxprespace{0.04em}\noindent\lawboxuserandnumber #1\par}
\newcommand{\pendingtranslation}{\par{\songfont\footnotesize\color{pendingink}（中文譯文待補。）}\par}
\newcommand{\tocpart}[1]{\par\addvspace{0.52em}{\footnotesize\bfseries\color{pendingink}#1\par}\addvspace{0.16em}}
\newcommand{\tocdashline}{\par\addvspace{0.48em}\noindent{\color{notegray}\xleaders\hbox to 0.82em{\hfil\rule{0.42em}{0.28pt}\hfil}\hskip\linewidth}\par\addvspace{0.38em}}
% \tocpanel{font-cmd}{content}：將目錄置中於所在欄寬（雙語欄適用）。
% A3 橫式使用 0.58\linewidth，A4 橫式使用 0.84\linewidth，
% 使兩種紙張下目錄欄內容實際寬度相近（均約 100–105 mm）。
\newcommand{\tocpanel}[2]{%
  \makebox[\linewidth][c]{%
    \ifx\bilingualpaper\bilingualpaperaThree
      \begin{minipage}{0.58\linewidth}%
    \else
      \begin{minipage}{0.84\linewidth}%
    \fi
    \toccolstyle
    #1#2%
    \end{minipage}%
  }%
}
% ===== 首頁簡目錄的兩層 description list 縮排 =====
% enumitem 的 description list 可把每一列拆成：
%   [label]  labelsep  body text
% 這裡的「起點」都以所在 minipage/欄位的左緣為 0。
%
% 核心規則：
%   labelindent = label 欄位從左緣開始的位置。
%   labelwidth  = label 欄位本身寬度；標籤太長（如 Abw. Meinung）就加大它。
%   labelsep    = label 欄位右緣到正文的距離。
%   leftmargin  = 正文 body text 的起點。判斷層級視覺時主要看這個。
%
% 所以：
%   想讓「Begründung / 判決理由」整列正文往右或往左：改 tocitems 的 leftmargin。
%   想讓「A. Verfahren und Gegenstand / A. 程序與審查標的」正文往右或往左：
%     改 tocsubitems 的 leftmargin。
%   想讓 A./B./C. 這些標籤本身往右或往左，但正文不動：改 tocsubitems 的 labelindent。
%   想讓 A./B./C. 與正文間距變大或變小：改 tocsubitems 的 labelsep。
%
% 層級原則：
%   tocsubitems.leftmargin 必須大於 tocitems.leftmargin，
%   否則子層級正文會看起來比「Gründe / 理由」還靠左。
\newlist{tocitems}{description}{1}
\setlist[tocitems]{%
  labelindent=0pt,    % 第一層 label 欄位起點；0pt 表示貼齊目錄欄左緣。
  leftmargin=2.54cm,  % 第一層正文起點；例如 Begründung / 判決理由 的左緣。
  labelwidth=2.16cm,  % 第一層 label 欄位寬度；需容納 Leitsätze、Abw. Meinung、不同意見等。
  labelsep=0.32cm,    % 第一層 label 與正文的水平間距；只改間距，不改正文起點。
  itemsep=0.28em,     % 第一層各項之間的垂直距離。
  topsep=0.20em,      % 第一層 list 上下與前後內容的垂直距離。
  parsep=0pt,         % 同一 item 內段落之間的距離；目錄項通常不需要。
  partopsep=0pt,      % list 若緊接段落開始時的額外距離；關閉以免目錄高度飄動。
  align=left,         % label 欄內左對齊；長短標籤不靠右貼齊。
  font=\normalfont\bfseries % label 的字型；正文不受這行影響。
}
\newlist{tocsubitems}{description}{1}
\setlist[tocsubitems]{%
  labelindent=1cm,    % 第二層 label 起點；只控制 A./B./C. 本身的位置。
  leftmargin=3.00cm,  % 第二層正文起點；必須大於 2.54cm，才會比 Begründung / 判決理由更右。
  labelwidth=0.5cm,   % 第二層 label 欄寬；A./B./C. 很短，可小於第一層。
  labelsep=1.5cm,    % 第二層 label 與正文間距；A. 與 Verfahren / 程序 的距離看這裡。
  itemsep=0.12em,     % 第二層各項較密，避免 A.–E. 佔太高。
  topsep=0.04em,      % 第二層 list 與 Gründe / 理由之間的垂直距離。
  parsep=0pt,         % 同一第二層 item 內段落距離；目錄項通常不需要。
  partopsep=0pt,      % 關閉額外段落起始距離，避免版面不穩。
  align=left,         % A./B./C. label 欄內左對齊。
  font=\normalfont\bfseries, % A./B./C. label 加粗；正文不受這行影響。
  before={}           % 不在第二層 list 前自動插入額外內容。
}
\newlist{termitems}{description}{1}
\ifbilingualcolumns
  \setlist[termitems]{%
    leftmargin=5.2cm,
    labelwidth=4.8cm,
    labelsep=0.35cm,
    itemsep=0.16em,
    topsep=0.2em,
    parsep=0pt,
    partopsep=0pt,
    font=\normalfont\bfseries,
    style=nextline
  }
\else
  \setlist[termitems]{%
    leftmargin=0pt,
    labelwidth=0pt,
    labelsep=0pt,
    itemsep=0.24em,
    topsep=0.2em,
    parsep=0pt,
    partopsep=0pt,
    font=\normalfont\bfseries,
    style=nextline
  }
\fi
\newlist{abbritems}{description}{1}
\setlist[abbritems]{%
  leftmargin=2.38cm,
  labelwidth=2.02cm,
  labelsep=0.28cm,
  itemsep=0.16em,
  topsep=0.2em,
  parsep=0pt,
  partopsep=0pt,
  align=left,
  font=\normalfont\bfseries
}
\ifbilingualcolumns
  \newenvironment{termcolumns}{\begin{multicols}{2}}{\end{multicols}}
\else
  \newenvironment{termcolumns}{}{}
\fi

% ===== 編者註腳開關 =====
% 一般版預設不顯示編者註，避免正文腳註過密。
% 若開啟 \classroommodeon，GG/條文編者註仍會自動顯示，無需另開本開關。
% 若一般版也要顯示編者註，可在單案檔 \input 後加 \showeditornotestrue。
\newif\ifshoweditornotes
\showeditornotesfalse

% ===== 目錄排版輔助 =====
\newcommand{\toccolstyle}{\raggedright\arraybackslash}
\newcommand{\tocsingleindent}{\leftskip=1.45cm\rightskip=1.2cm}

% ===== 行內引文環境（德文原文與中文譯文各一，帶左側灰色邊線）=====
\newcommand{\quoteparbreak}{\par\addvspace{0.48em}}
\newcommand{\sourcequoteinner}[1]{%
  \begin{tcolorbox}[
    enhanced, breakable, width=\dimexpr\linewidth-0.8em\relax, frame hidden,
    colback=white, coltext=caseink, boxrule=0pt,
    borderline west={0.55pt}{0.88em}{notegray},
    left=1.78em, right=0.72em, top=0.18em, bottom=0.18em,
    boxsep=0pt, before skip=0.24em, after skip=0.24em,
    before upper={\normalfont\footnotesize\color{caseink}\setlength{\parindent}{0pt}\setlength{\parskip}{0.34em}\raggedright\setlength{\rightskip}{0pt plus 1em}}
  ]%
  #1\par
  \end{tcolorbox}%
}
\newcommand{\transquoteinner}[1]{%
  \begin{tcolorbox}[
    enhanced, breakable, width=\dimexpr\linewidth-0.8em\relax, frame hidden,
    colback=white, coltext=caseink, boxrule=0pt,
    borderline west={0.55pt}{0.88em}{notegray},
    left=1.78em, right=0.72em, top=0.18em, bottom=0.18em,
    boxsep=0pt, before skip=0.24em, after skip=0.24em,
    before upper={\songfont\small\color{caseink}\setlength{\parindent}{0pt}\setlength{\parskip}{0.34em}\raggedright\setlength{\rightskip}{0pt plus 1em}}
  ]%
  #1\par
  \end{tcolorbox}%
}

% ===== 大綱（Gliederung）Rn 輔助命令 =====
\newcommand{\outrn}[2]{\enspace{\normalfont\footnotesize\color{pendingink}(Rn\,#1–#2)}}
\newcommand{\outrnsingle}[1]{\enspace{\normalfont\footnotesize\color{pendingink}(Rn\,#1)}}
\newcommand{\outrnzh}[2]{\enspace{\normalfont\footnotesize\color{pendingink}（第\,#1–#2\,段）}}
\newcommand{\outrnzhs}[1]{\enspace{\normalfont\footnotesize\color{pendingink}（第\,#1\,段）}}
% ===== 大綱雙語對照表格輔助命令（三欄：段碼 │ 德文 │ 中文）=====
% 欄 1：段碼（由欄定義自動格式化：小字、等寬、右對齊、pendingink）
% 欄 2：德文標籤＋正文
% 欄 3：中文標籤＋正文
%
% 無段碼的列（\omainrow、\osecrow），欄 1 以 & 空置。
% 有段碼的列（\osecrowrn、\osubrow、\ointrow），#1 為預格式化字串
%   例：\osubrow{2–49}{I.}{...}{I.}{...}
%       \ointrow{107}{...}{...}
%
% \opartrow：段組標題（橫跨三欄）
\newcommand{\opartrow}[2]{%
  \noalign{\vspace{0.30em}}%
  \multicolumn{3}{@{}l@{}}{%
    \footnotesize\bfseries\color{pendingink}#1\hfill\songfont#2%
  }\\[0.06em]%
}
% \odashrow：虛線分隔（橫跨三欄）
\newcommand{\odashrow}{%
  \noalign{\vspace{0.28em}}%
  \multicolumn{3}{@{}l@{}}{\color{notegray}%
    \xleaders\hbox to 0.82em{\hfil\rule{0.42em}{0.28pt}\hfil}\hskip 0.88\linewidth}%
  \\[0.24em]%
}
% \odissentpart：不同意見書區段標頭。比一般 \opartrow 更像附錄性轉場。
\newcommand{\odissentpart}[2]{%
  \noalign{\vspace{0.42em}}%
  \multicolumn{3}{@{}p{0.88\textwidth}@{}}{%
    \color{notegray}\rule{\linewidth}{0.18pt}%
  }\\[-0.18em]%
  \multicolumn{3}{@{}p{0.88\textwidth}@{}}{%
    \makebox[\linewidth][s]{%
      \footnotesize\bfseries\color{pendingink}#1%
      \hfill{\songfont #2}%
    }%
  }\\[-0.02em]%
}
% \odissentopinion：不同意見書的署名與一行摘要，右欄保留段碼範圍。
\newcommand{\odissentopinion}[5]{%
  \footnotesize\hangindent=0.52cm\hangafter=1%
    {\bfseries\color{caseink}#2}\enspace{\color{notegray}\textperiodcentered}\enspace\color{caseink}#3%
  &\footnotesize\hangindent=0.52cm\hangafter=1%
    {\songfont\bfseries\color{caseink}#4}\enspace{\color{notegray}\textperiodcentered}\enspace\songfont\color{caseink}#5%
  &#1%
  \\[0.12em]%
}
% \omainrow：頂層項目，無段碼（Leitsätze、Urteil、Gründe …）
\newcommand{\omainrow}[4]{%
  \footnotesize\hangindent=1.92cm\hangafter=1%
    \makebox[1.92cm][l]{\bfseries\color{caseink}#1}\color{caseink}#2%
  &\footnotesize\hangindent=1.92cm\hangafter=1%
    \makebox[1.92cm][l]{\bfseries\color{caseink}#3}\songfont\color{caseink}#4%
  &%
  \\[0.15em]%
}
% \omainrowrn：頂層項目，有段碼（不同意見等標籤寬於 \osecrow 框者）
\newcommand{\omainrowrn}[5]{%
  \footnotesize\hangindent=1.92cm\hangafter=1%
    \makebox[1.92cm][l]{\bfseries\color{caseink}#2}\color{caseink}#3%
  &\footnotesize\hangindent=1.92cm\hangafter=1%
    \makebox[1.92cm][l]{\bfseries\color{caseink}#4}\songfont\color{caseink}#5%
  &#1%
  \\[0.15em]%
}
% \osecrow：一級子項，無段碼（A.、C.、D. 等有子項者）
\newcommand{\osecrow}[4]{%
  \footnotesize\hspace{1.10cm}\hangindent=1.98cm\hangafter=1%
    \makebox[0.86cm][l]{\bfseries\color{caseink}#1}\color{caseink}#2%
  &\footnotesize\hspace{1.10cm}\hangindent=1.98cm\hangafter=1%
    \makebox[0.86cm][l]{\bfseries\color{caseink}#3}\songfont\color{caseink}#4%
  &%
  \\[0.11em]%
}
% \osecrowrn：一級子項，有段碼（B.、E. 等完整段落範圍者）
% #1=段碼字串（如 101–106）, #2=DE標籤, #3=DE正文, #4=ZH標籤, #5=ZH正文
\newcommand{\osecrowrn}[5]{%
  \footnotesize\hspace{1.10cm}\hangindent=1.98cm\hangafter=1%
    \makebox[0.86cm][l]{\bfseries\color{caseink}#2}\color{caseink}#3%
  &\footnotesize\hspace{1.10cm}\hangindent=1.98cm\hangafter=1%
    \makebox[0.86cm][l]{\bfseries\color{caseink}#4}\songfont\color{caseink}#5%
  &#1%
  \\[0.11em]%
}
% \osubrow：二級子項，有段碼範圍（I.、II.、A.I. …）
% #1=段碼字串, #2=DE標籤, #3=DE正文, #4=ZH標籤, #5=ZH正文
\newcommand{\osubrow}[5]{%
  \footnotesize\hspace{1.40cm}\hangindent=2.30cm\hangafter=1%
    \makebox[0.90cm][l]{\bfseries\color{caseink}#2}\color{caseink}#3%
  &\footnotesize\hspace{1.40cm}\hangindent=2.30cm\hangafter=1%
    \makebox[0.90cm][l]{\bfseries\color{caseink}#4}\songfont\color{caseink}#5%
  &#1%
  \\[0.09em]%
}
% \ointrow：引入段落，有單一段碼，無標籤
% #1=段碼字串, #2=DE正文, #3=ZH正文
\newcommand{\ointrow}[3]{%
  \footnotesize\hspace{0.52cm}\hangindent=0.52cm\hangafter=1\color{caseink}#2%
  &\footnotesize\hspace{0.52cm}\hangindent=0.52cm\hangafter=1\songfont\color{caseink}#3%
  &#1%
  \\[0.09em]%
}
% \osubsubrow：三級子項（1.、2.、3.）
% #1=段碼字串, #2=DE標籤, #3=DE正文, #4=ZH標籤, #5=ZH正文
\newcommand{\osubsubrow}[5]{%
  \footnotesize\hspace{1.80cm}\hangindent=2.48cm\hangafter=1%
    \makebox[0.68cm][l]{\bfseries\color{caseink}#2}\color{caseink}#3%
  &\footnotesize\hspace{1.80cm}\hangindent=2.48cm\hangafter=1%
    \makebox[0.68cm][l]{\bfseries\color{caseink}#4}\songfont\color{caseink}#5%
  &#1%
  \\[0.08em]%
}
% \osubsubsubrow：四級子項（a）、b））
% #1=段碼字串, #2=DE標籤, #3=DE正文, #4=ZH標籤, #5=ZH正文
\newcommand{\osubsubsubrow}[5]{%
  \footnotesize\hspace{2.10cm}\hangindent=2.68cm\hangafter=1%
    \makebox[0.58cm][l]{\bfseries\color{caseink}#2}\color{caseink}#3%
  &\footnotesize\hspace{2.10cm}\hangindent=2.68cm\hangafter=1%
    \makebox[0.58cm][l]{\bfseries\color{caseink}#4}\songfont\color{caseink}#5%
  &#1%
  \\[0.07em]%
}
% \osubsubsubsubrow：五級子項（aa）、bb））
% 標籤框 0.62cm：足以容納「aa)」在 footnotesize bold 下（約 0.50cm）並留有間距
% #1=段碼字串, #2=DE標籤, #3=DE正文, #4=ZH標籤, #5=ZH正文
\newcommand{\osubsubsubsubrow}[5]{%
  \footnotesize\hspace{2.38cm}\hangindent=3.06cm\hangafter=1%
    \makebox[0.68cm][l]{\bfseries\color{caseink}#2}\color{caseink}#3%
  &\footnotesize\hspace{2.38cm}\hangindent=3.06cm\hangafter=1%
    \makebox[0.68cm][l]{\bfseries\color{caseink}#4}\songfont\color{caseink}#5%
  &#1%
  \\[0.06em]%
}
% \outlinepage：雙語對照大綱頁包裝環境（longtable 自動跨頁）
% 三欄：德文欄 + 中文欄 + 段碼欄（最右，不折行）
% 段碼欄寬度：1.80cm，足以容納最長段碼如「309–348」
% DE/ZH 欄寬：(0.87\textwidth - 2.08cm) / 2
%   驗算：2×欄 + 0.05\textwidth + 0.28cm gap + 1.80cm = 0.87tw - 2.08cm + 0.05tw + 2.08cm = 0.92tw ✓
\newenvironment{outlinepage}{%
  \vspace*{0.40cm}%
  \setlength{\LTleft}{0.04\textwidth}%
  \setlength{\LTright}{0.04\textwidth}%
  \setlength{\tabcolsep}{0pt}%
  % 大綱列距由 longtable 的 \arraystretch 與各 row macro 結尾的 \\[...em] 共同決定。
  % 這裡控制「每一列本身的表格 strut 高度」；數值越大，A./I./1. 各項之間越像有空行。
  % A3 原本用 0.62，視覺上沒有空行；A4 也沿用同值，避免切到 A4 後項目被撐開。
  % 若日後覺得大綱太擠，只微調這個值（例如 0.68），不要改各 \omainrow/\osubrow 的縮排。
  \renewcommand{\arraystretch}{0.62}%
  \begin{longtable}{%
    >{\vspace{0pt}\raggedright\arraybackslash}p{\dimexpr(0.87\textwidth-2.08cm)/2\relax}%
    @{\hspace{0.025\textwidth}}%
    !{\color{notegray}\vrule width 0.12pt}%
    @{\hspace{0.025\textwidth}}%
    >{\vspace{0pt}\raggedright\arraybackslash}p{\dimexpr(0.87\textwidth-2.08cm)/2\relax}%
    @{\hspace{0.28cm}}%
    >{\vspace{0pt}\footnotesize\normalfont\color{pendingink}\raggedright\arraybackslash}p{1.80cm}%
    @{}%
  }%
  % 首頁欄標籤：不要橫跨置中；分別放在德文欄與中文欄的實際欄位起點。
  \footnotesize\bfseries\color{sourceink}Gliederung%
  &\footnotesize\bfseries\songfont\color{sourceink}大綱%
  &%
  \\[0.36em]%
  \endfirsthead%
  % 續頁欄標籤：同樣分欄定位，以灰色細體區分；無需標注「續」。
  \footnotesize\color{pendingink}Gliederung%
  &\footnotesize\songfont\color{pendingink}大綱%
  &%
  \\[0.24em]%
  \endhead%
}{%
  \end{longtable}%
}
 載入。
% 不含：\documentclass、\documentversion、\ggversionde/zh、\tocde/\toczh。
% 日期與首頁版本列由本檔自動產生；各文件只需手動指定 \documentversion。

\usepackage{xeCJK}
\usepackage{fontspec}
\usepackage{geometry}
\usepackage{setspace}
\usepackage{hyperref}
\usepackage{xurl}
\usepackage{bookmark}
\usepackage{enumitem}
\usepackage{xcolor}
\usepackage{titlesec}
\usepackage{array}
\usepackage{longtable}
\usepackage{tcolorbox}
\usepackage{environ}
\usepackage{pdflscape}
\usepackage{etoolbox}
\usepackage{ragged2e}
\usepackage{paracol}
\usepackage{multicol}
\usepackage{needspace}
\usepackage{fancyhdr}
\usepackage{tikz}
\usepackage{zref-abspage}
\tcbuselibrary{breakable,skins}
\usetikzlibrary{calc}
\usepackage{pgfcalendar}

\hypersetup{
  unicode=true,
  bookmarksopen=true,
  bookmarksnumbered=false,
  hidelinks
}

% ===== 自動推導，不要手動改下面 =====
% （\docmode 與 \bilingualpaper 由各文件在 % bverfge_layout.tex
% 共用排版基礎設施，供 Schwangerschaftsabbruch I 與 II 共享。
% 由各文件以 \documentclass 後 \input{../../共用LaTeX/bverfge_layout} 載入。
% 不含：\documentclass、\documentversion、\ggversionde/zh、\tocde/\toczh。
% 日期與首頁版本列由本檔自動產生；各文件只需手動指定 \documentversion。

\usepackage{xeCJK}
\usepackage{fontspec}
\usepackage{geometry}
\usepackage{setspace}
\usepackage{hyperref}
\usepackage{xurl}
\usepackage{bookmark}
\usepackage{enumitem}
\usepackage{xcolor}
\usepackage{titlesec}
\usepackage{array}
\usepackage{longtable}
\usepackage{tcolorbox}
\usepackage{environ}
\usepackage{pdflscape}
\usepackage{etoolbox}
\usepackage{ragged2e}
\usepackage{paracol}
\usepackage{multicol}
\usepackage{needspace}
\usepackage{fancyhdr}
\usepackage{tikz}
\usepackage{zref-abspage}
\tcbuselibrary{breakable,skins}
\usetikzlibrary{calc}
\usepackage{pgfcalendar}

\hypersetup{
  unicode=true,
  bookmarksopen=true,
  bookmarksnumbered=false,
  hidelinks
}

% ===== 自動推導，不要手動改下面 =====
% （\docmode 與 \bilingualpaper 由各文件在 \input{bverfge_layout} 之前設定；
%   預設 chinese / a4）
\ifdefined\docmode\else\def\docmode{chinese}\fi
\ifdefined\bilingualpaper\else\def\bilingualpaper{a4}\fi
\newif\ifshoworiginal
\newif\ifshowtranslation
\newif\ifbilinguallandscape
\newif\ifbilingualcolumns
\newif\ifintranslation
\newif\iffirstbilingualsection

\showoriginalfalse
\showtranslationfalse
\bilinguallandscapefalse
\bilingualcolumnsfalse
\intranslationfalse
\firstbilingualsectionfalse

\def\modechinese{chinese}
\def\modegerman{german}
\def\modebilingual{bilingual}
\def\bilingualpaperaThree{a3}
\def\bilingualpaperaFour{a4}

\ifx\docmode\modechinese
  \showtranslationtrue
\fi

\ifx\docmode\modegerman
  \showoriginaltrue
\fi

\ifx\docmode\modebilingual
  \showoriginaltrue
  \showtranslationtrue
  \bilinguallandscapetrue
  \bilingualcolumnstrue
\fi

% 編排模式說明：
% chinese  → \showoriginalfalse  \showtranslationtrue  （A4 直式，純中文）
% german   → \showoriginaltrue   \showtranslationfalse （A4 直式，純德文）
% bilingual→ \showoriginaltrue   \showtranslationtrue  \bilingualcolumnstrue（A3/A4 橫式對照）

\ifbilingualcolumns
  \ifx\bilingualpaper\bilingualpaperaFour
    \geometry{
      a4paper,
      landscape,
      left=1.25cm,
      right=2.25cm,
      top=1.25cm,
      bottom=1.75cm,
      footskip=8.5mm,
      marginparwidth=0.75cm,
      marginparsep=0.25cm
    }
  \else
    \geometry{
      a3paper,
      landscape,
      left=1.55cm,
      right=2.75cm,
      top=1.55cm,
      bottom=2.05cm,
      footskip=9.5mm,
      marginparwidth=0.9cm,
      marginparsep=0.35cm
    }
  \fi
\else
\geometry{
  a4paper,
  portrait,
  left=2.2cm,
  right=2.6cm,
  top=2.0cm,
  bottom=2.0cm,
  marginparwidth=0.8cm,
  marginparsep=0.3cm
}
\fi
% ===== 時間戳基礎設施 =====
\newcount\stamptimeval
\newcount\stampminuteval
\newcommand{\stamphh}{%
  \stamptimeval=\time
  \divide\stamptimeval by 60
  \ifnum\stamptimeval<10 0\fi
  \the\stamptimeval}
\newcommand{\stampmm}{%
  \stamptimeval=\time
  \stampminuteval=\time
  \divide\stampminuteval by 60
  \multiply\stampminuteval by 60
  \advance\stamptimeval by -\stampminuteval
  \ifnum\stamptimeval<10 0\fi
  \the\stamptimeval}
\newcommand{\stampwkde}[1]{%
  \ifcase#1Montag\or Dienstag\or Mittwoch\or Donnerstag\or Freitag\or Samstag\or Sonntag\fi}
\newcommand{\stampwkzh}[1]{%
  \ifcase#1 一\or 二\or 三\or 四\or 五\or 六\or 日\fi}
\newcommand{\stampmonthde}[1]{%
  \ifcase#1\or Januar\or Februar\or März\or April\or Mai\or Juni\or
  Juli\or August\or September\or Oktober\or November\or Dezember\fi}
\newcount\stampjdval
\newcount\stampwdval
\pgfcalendardatetojulian{\the\year-\the\month-\the\day}{\stampjdval}
\pgfcalendarjuliantoweekday{\the\stampjdval}{\stampwdval}
\newcommand{\timestampzh}{%
  \songfont 最後修改：\the\year 年\the\month 月\the\day 日%
  % （星期\stampwkzh{\the\stampwdval}）%
  % \space\stamphh:\stampmm
  }

\newcommand{\documentdatezh}{\the\year 年 \the\month 月 \the\day 日}
\newcommand{\documentdatede}{\the\day. \stampmonthde{\the\month} \the\year}
\newcommand{\bverfgetranslationnote}{%
  {\songfont 翻譯說明：初譯使用 Claude Code 與 GPT 協助迭代；仍請以德文原文為準。}%
}
\newcommand{\bverfgecourseinfo}{%
  {\songfont 課程：114 學年度第 2 學期；憲法專題研究（四）；授課老師：湯德宗教授；導讀日期：115 年 6 月 1 日。}%
}
\newcommand{\bverfgeeditorinfo}{%
  {\songfont 編者：王逸帆（R10A21126），國立台灣大學法律系研究所碩士班。}%
}
\newcommand{\bverfgetitleversionline}{%
  \ifbilingualcolumns
    Fassung \documentversion\enspace\textperiodcentered\enspace Stand: \documentdatede
    \quad
    {\songfont 版本 \documentversion\enspace\textperiodcentered\enspace 修訂日期：\documentdatezh\par}%
  \else
    \ifshowtranslation
      {\songfont 版本 \documentversion\enspace\textperiodcentered\enspace 修訂日期：\documentdatezh\par}%
    \else
      Fassung \documentversion\enspace\textperiodcentered\enspace Stand: \documentdatede\par
    \fi
  \fi
}
\newcommand{\bverfgetitlecornerinfo}{%
  \ifshowtranslation
    \vspace{0.72em}%
    \noindent\hfill
    \begin{minipage}{0.66\textwidth}
      \raggedleft\tiny\color{pendingink}\songfont
      \bverfgecourseinfo\par
      \bverfgeeditorinfo
    \end{minipage}%
  \fi
}
\newcommand{\bverfgetitlemeta}{%
  \vfill
  \begin{center}%
    \footnotesize\color{pendingink}%
    \bverfgetitleversionline
    \ifshowtranslation
      \vspace{0.28em}{\scriptsize\bverfgetranslationnote\par}%
    \fi
  \end{center}%
  \bverfgetitlecornerinfo
}

\pagestyle{fancy}
\fancyhf{}
\renewcommand{\headrulewidth}{0pt}
\renewcommand{\footrulewidth}{0pt}
\fancyfoot[C]{\thepage}
\ifbilingualcolumns
  \fancyfoot[R]{\scriptsize\color{pendingink}\timestampzh}%
\else
  \ifshoworiginal
  \else
    \fancyfoot[R]{\scriptsize\color{pendingink}\timestampzh\hspace{0.45cm}}%
  \fi
\fi
\setmainfont{Times New Roman}
\setsansfont{Helvetica Neue}
\setmonofont{Menlo}
\IfFileExists{/Users/iw/Library/Fonts/kaiu.ttf}{
  \setCJKmainfont[Path=/Users/iw/Library/Fonts/,AutoFakeBold=4]{kaiu.ttf}
}{
  \IfFontExistsTF{標楷體}{
    \setCJKmainfont[AutoFakeBold=4]{標楷體}
  }{
    \setCJKmainfont{Songti TC}
  }
}
\IfFontExistsTF{Songti TC}{
  \setCJKfamilyfont{song}{Songti TC}
  \setCJKmonofont{Songti TC}
}{
  \setCJKfamilyfont{song}[Path=/Users/iw/Library/Fonts/,AutoFakeBold=4]{kaiu.ttf}
  \setCJKmonofont[Path=/Users/iw/Library/Fonts/,AutoFakeBold=4]{kaiu.ttf}
}
\newcommand{\songfont}{\CJKfamily{song}}

% ===== 封面／目錄專用打字機字體（僅限封面、目錄頁使用，不影響正文）=====
\IfFontExistsTF{Radio Newsman}{%
  \newfontfamily\bverfgetitlede{Radio Newsman}%
}{%
  \newfontfamily\bverfgetitlede{Courier New}%
}
\IfFontExistsTF{Erikas Farbband}{%
  \newfontfamily\bverfgebodyde{Erikas Farbband}%
}{%
  \let\bverfgebodyde\bverfgetitlede
}
\IfFontExistsTF{Maoken Heavy Labourer}{%
  \setCJKfamilyfont{bverfgetitlecjk}{Maoken Heavy Labourer}%
}{%
  \setCJKfamilyfont{bverfgetitlecjk}{Songti TC}%
}
\IfFontExistsTF{Huiwen-mincho}{%
  \setCJKfamilyfont{bverfgebodycjk}{Huiwen-mincho}%
}{%
  \setCJKfamilyfont{bverfgebodycjk}{Songti TC}%
}
\newcommand{\bverfgetitlezh}{\bverfgetitlede\CJKfamily{bverfgetitlecjk}}
\newcommand{\bverfgebodyzh}{\bverfgebodyde\CJKfamily{bverfgebodycjk}}

\setstretch{1.35}
\setlist{itemsep=0.2em, topsep=0.4em}
\setlength{\parindent}{0pt}
\setlength{\parskip}{0.45em}
\raggedbottom
\setcounter{secnumdepth}{0}
\setcounter{tocdepth}{2}

\newcommand{\lawboxsourcefont}{\footnotesize}
\newcommand{\lawboxtransfont}{\footnotesize}
\newcommand{\lawboxsourcestretch}{1.02}
\newcommand{\lawboxtransstretch}{0.92}

\titleformat{\section}{\centering\Large\bfseries\songfont}{\thesection}{1em}{}
\titleformat{\subsection}{\large\bfseries\songfont}{\thesubsection}{1em}{}
\titleformat{\subsubsection}{\normalsize\bfseries\songfont}{\thesubsubsection}{1em}{}
\titleformat{\paragraph}{\normalsize\bfseries\songfont}{\theparagraph}{1em}{}
\titlespacing*{\section}{0pt}{1.8em}{1.0em}
\titlespacing*{\subsection}{0pt}{1.2em}{0.6em}
\titlespacing*{\subsubsection}{0pt}{1.0em}{0.45em}
\titlespacing*{\paragraph}{0pt}{0.6em}{0.4em}

\definecolor{noteblue}{HTML}{A9C9F5}
\definecolor{noteteal}{HTML}{AEE1D6}
\definecolor{notegray}{HTML}{D7DADF}
\definecolor{caseink}{HTML}{000000}
\definecolor{casepaper}{HTML}{FCFEFD}
\definecolor{sourcepaper}{HTML}{FAFBF8}
\definecolor{sourceline}{HTML}{D7DED8}
\definecolor{sourceink}{HTML}{000000}
\definecolor{pendingink}{HTML}{000000}

\tcbset{
  caseboxbase/.style={
    enhanced,
    breakable,
    width=0.985\linewidth,
    colback=casepaper,
    coltitle=caseink,
    fonttitle=\bfseries\songfont,
    boxrule=0.28pt,
    arc=1.2mm,
    left=2.4mm,
    right=2.4mm,
    top=1.5mm,
    bottom=1.5mm,
    boxsep=1.5mm,
    before skip=0.65em,
    after skip=0.65em,
    before upper={\setlength{\parindent}{0pt}\setlength{\parskip}{0.35em}},
  }
}

\newcommand{\bilingualstart}{}
\newcommand{\bilingualstop}{}
\ifbilingualcolumns
  \renewcommand{\bilingualstart}{\small\setlength{\columnsep}{1.25cm}\setlength{\columnseprule}{0.12pt}\columnratio{0.5,0.5}\multicolumnfootnotes\flushbottom\sloppy\interlinepenalty=80\clubpenalty=800\widowpenalty=800\displaywidowpenalty=800\begin{paracol}{2}\global\intranslationfalse\global\firstbilingualsectiontrue{\footnotesize\color{sourceink}Deutsch\par}\switchcolumn{\footnotesize\songfont\color{sourceink}中文譯文\par}\switchcolumn*}
  \renewcommand{\bilingualstop}{\ifintranslation\switchcolumn*\global\intranslationfalse\fi\end{paracol}\clearpage\raggedbottom}
\fi
\newif\ifrandnumbers
\randnumbersfalse
\newcounter{randnumber}
\newcounter{randnumberlocal}
\newcounter{lastsourcerandstart}
\newif\iflastsourcehasrandnumbers
\lastsourcehasrandnumbersfalse
\newif\ifinlawbox
\inlawboxfalse
\newif\iflawboxhaspendingrn
\lawboxhaspendingrnfalse
\newcommand{\lawboxpendingrn}{}
\newcommand{\startrandnumbers}{\global\randnumberstrue\setcounter{randnumber}{1}\global\lastsourcehasrandnumbersfalse}
\newcommand{\stoprandnumbers}{\global\randnumbersfalse\global\lastsourcehasrandnumbersfalse}
\newlength{\randnumberwidth}
\setlength{\randnumberwidth}{0.95cm}
\newlength{\randnumberrightoffset}
\setlength{\randnumberrightoffset}{0.85cm}
\newlength{\randnumberlawboxrightoffset}
\setlength{\randnumberlawboxrightoffset}{1.40cm}
\newcommand{\randnumberbox}[1]{%
  \leavevmode
  \makebox[0pt][l]{%
    \smash{%
      \tikz[remember picture,overlay,baseline=(rnmark.base)]{%
        \coordinate (rnmark) at (0,0);%
        \ifinlawbox
          \path let \p1=(current page.east), \p2=(rnmark.base) in
            node[
              anchor=base east,
              inner sep=0pt,
              outer sep=0pt,
              text width=\randnumberwidth,
              align=right,
              text=pendingink,
              font=\normalfont\mdseries\scriptsize\ttfamily
            ] at (\x1-\randnumberlawboxrightoffset,\y2) {{\normalfont\mdseries\scriptsize\ttfamily #1}};%
        \else
          \path let \p1=(current page.east), \p2=(rnmark.base) in
            node[
              anchor=base east,
              inner sep=0pt,
              outer sep=0pt,
              text width=\randnumberwidth,
              align=right,
              text=pendingink,
              font=\normalfont\mdseries\scriptsize\ttfamily
            ] at (\x1-\randnumberrightoffset,\y2) {{\normalfont\mdseries\scriptsize\ttfamily #1}};%
        \fi
      }%
    }%
  }%
  \ignorespaces
}
\newcommand{\lawboxstorerandnumber}[1]{%
  \gdef\lawboxpendingrn{#1}%
  \global\lawboxhaspendingrntrue
  \ignorespaces
}
\newcommand{\lawboxuserandnumber}{%
  \iflawboxhaspendingrn
    \randnumberbox{\lawboxpendingrn}%
    \global\lawboxhaspendingrnfalse
  \fi
}
\newcommand{\sourcern}[1]{%
  \ifrandnumbers
    \ifshoworiginal
      \ifshowtranslation\else
        \ifinlawbox\lawboxstorerandnumber{#1}\else\randnumberbox{#1}\fi
      \fi
    \fi
  \fi
}
\newcommand{\transrn}[1]{%
  \ifrandnumbers
    \ifshowtranslation
      \ifinlawbox\lawboxstorerandnumber{#1}\else\randnumberbox{#1}\fi
    \fi
  \fi
}
% ===== 課堂模式：德文原文縮寫註解 =====
% 日常切換只改各判決檔的一行：
%   \classroommodeon        開啟課堂版；註解掉這行即回到一般版。
% 模板預設：
%   1. 課堂模式關閉。
%   2. 課堂模式開啟時，縮寫註解包含「德文全稱＋中文譯名」。
%   3. 課堂模式開啟時，GG/條文編者註仍會自動顯示。
% 進階微調才需要在單案檔覆寫：
%   \classroomabbrzhoff     縮寫註解僅顯示德文全稱。
%   \showeditornotestrue    一般版也顯示編者註。
% 註解策略：同一實際 PDF 頁面中，同一縮寫只在第一次出現處加註。
% 這裡用 zref-abspage 的絕對頁序，不用 \thepage，避免文件內頁碼重置時誤判。
\newif\ifclassroommode
\newif\ifclassroomabbrzh
\classroommodefalse
\classroomabbrzhtrue
\newcommand{\classroommodeon}{\classroommodetrue}
\newcommand{\classroommodeoff}{\classroommodefalse}
\newcommand{\classroomabbrzhon}{\classroomabbrzhtrue}
\newcommand{\classroomabbrzhoff}{\classroomabbrzhfalse}
\newcommand{\abbrnotelabel}{\emph{Abk.}: }
\newcommand{\abbrnote}[3]{%
  \ifclassroommode
    \ifshoworiginal
      \ifintranslation\else
        \footnote{\normalfont\footnotesize\abbrnotelabel #2\ifclassroomabbrzh；{\songfont #3}\fi}%
      \fi
    \fi
  \fi
}
\newcommand{\abbrnoteonce}[4]{%
  #2%
  \ifclassroommode
    \ifshoworiginal
      \ifintranslation\else
        \begingroup
          \edef\abbrcurrentabspage{\number\numexpr\value{abspage}+1\relax}%
          \ifcsname abbrnoteseen@#1@\abbrcurrentabspage\endcsname\else
            \global\expandafter\let\csname abbrnoteseen@#1@\abbrcurrentabspage\endcsname\empty
            \abbrnote{#1}{#3}{#4}%
          \fi
        \endgroup
      \fi
    \fi
  \fi
}
\newcommand{\abbrAbs}{\abbrnoteonce{abs}{Abs.}{Absatz}{項}}
\newcommand{\abbrAAO}{\abbrnoteonce{aao}{a.a.O.}{am angegebenen Ort}{前引處}}
\newcommand{\abbrAE}{\abbrnoteonce{ae}{AE}{Alternativ-Entwurf}{替代草案}}
\newcommand{\abbrArt}{\abbrnoteonce{art}{Art.}{Artikel}{條}}
\newcommand{\abbrAufl}{\abbrnoteonce{aufl}{Aufl.}{Auflage}{版}}
\newcommand{\abbrBd}{\abbrnoteonce{bd}{Bd.}{Band}{卷}}
\newcommand{\abbrBGBl}{\abbrnoteonce{bgbl}{BGBl.}{Bundesgesetzblatt}{聯邦法律公報}}
\newcommand{\abbrBGHSt}{\abbrnoteonce{bghst}{BGHSt}{Entscheidungen des Bundesgerichtshofes in Strafsachen}{聯邦普通法院刑事判決彙編}}
\newcommand{\abbrBRDrucks}{\abbrnoteonce{brdrucks}{BRDrucks.}{Bundesrats-Drucksache}{聯邦參議院文件}}
\newcommand{\abbrBTDrucks}{\abbrnoteonce{btdrucks}{BTDrucks.}{Bundestags-Drucksache}{聯邦議院文件}}
\newcommand{\abbrBundesgesetzbl}{\abbrnoteonce{bundesgesetzbl}{Bundesgesetzbl.}{Bundesgesetzblatt}{聯邦法律公報}}
\newcommand{\abbrBvF}{\abbrnoteonce{bvf}{BvF}{Bundesverfassungsgerichtsverfahren}{聯邦憲法法院程序案號}}
\newcommand{\abbrBvL}{\abbrnoteonce{bvl}{BvL}{Bundesverfassungsgerichtsverfahren}{聯邦憲法法院程序案號}}
\newcommand{\abbrBvR}{\abbrnoteonce{bvr}{BvR}{Bundesverfassungsgerichtsverfahren}{聯邦憲法法院程序案號}}
\newcommand{\abbrBVerfG}{\abbrnoteonce{bverfg}{BVerfG}{Bundesverfassungsgericht}{聯邦憲法法院}}
\newcommand{\abbrBVerfGE}{\abbrnoteonce{bverfge}{BVerfGE}{Entscheidungen des Bundesverfassungsgerichts}{聯邦憲法法院判決集}}
\newcommand{\abbrBVerfGG}{\abbrnoteonce{bverfgg}{BVerfGG}{Bundesverfassungsgerichtsgesetz}{聯邦憲法法院法}}
\newcommand{\abbrDDR}{\abbrnoteonce{ddr}{DDR}{Deutsche Demokratische Republik}{德意志民主共和國；東德}}
\newcommand{\abbrDJT}{\abbrnoteonce{djt}{DJT}{Deutscher Juristentag}{德國法學家大會}}
\newcommand{\abbrDr}{\abbrnoteonce{dr}{Dr.}{Doktor}{Dr.}}
\newcommand{\abbrEuGRZ}{\abbrnoteonce{eugrz}{EuGRZ}{Europäische Grundrechte-Zeitschrift}{歐洲基本權期刊}}
\newcommand{\abbrEV}{\abbrnoteonce{ev}{EV}{Einigungsvertrag}{統一條約}}
\newcommand{\abbrF}{\abbrnoteonce{f}{f.}{folgende}{以下一頁；下一段}}
\newcommand{\abbrFF}{\abbrnoteonce{ff}{ff.}{fortfolgende}{以下數頁；以下各段}}
\newcommand{\abbrGBlDDR}{\abbrnoteonce{gblddr}{GBl. DDR}{Gesetzblatt der Deutschen Demokratischen Republik}{德意志民主共和國法律公報}}
\newcommand{\abbrGG}{\abbrnoteonce{gg}{GG}{Grundgesetz}{基本法}}
\newcommand{\abbrGVBl}{\abbrnoteonce{gvbl}{GVBl.}{Gesetz- und Verordnungsblatt}{法律及命令公報}}
\newcommand{\abbrJR}{\abbrnoteonce{jr}{JR}{Juristische Rundschau}{法學評論}}
\newcommand{\abbrJZ}{\abbrnoteonce{jz}{JZ}{JuristenZeitung}{法學家報}}
\newcommand{\abbrIVm}{\abbrnoteonce{ivm}{i.V.m.}{in Verbindung mit}{結合；連同}}
\newcommand{\abbrMdB}{\abbrnoteonce{mdb}{MdB}{Mitglied des Bundestages}{聯邦議院議員}}
\newcommand{\abbrMwN}{\abbrnoteonce{mwn}{m.w.N.}{mit weiteren Nachweisen}{附進一步文獻／判例出處}}
\newcommand{\abbrNF}{\abbrnoteonce{nf}{n.F.}{neue Fassung}{新版本}}
\newcommand{\abbrAF}{\abbrnoteonce{af}{a.F.}{alte Fassung}{舊版本}}
\newcommand{\abbrNr}{\abbrnoteonce{nr}{Nr.}{Nummer}{款、目或號，依語境}}
\newcommand{\abbrProf}{\abbrnoteonce{prof}{Prof.}{Professor}{教授}}
\newcommand{\abbrRdnr}{\abbrnoteonce{rdnr}{Rdnr.}{Randnummer}{邊碼；段碼}}
\newcommand{\abbrRdnrn}{\abbrnoteonce{rdnrn}{Rdnrn.}{Randnummern}{邊碼；段碼}}
\newcommand{\abbrRGBl}{\abbrnoteonce{rgbl}{RGBl.}{Reichsgesetzblatt}{帝國法律公報}}
\newcommand{\abbrRGSt}{\abbrnoteonce{rgst}{RGSt}{Entscheidungen des Reichsgerichts in Strafsachen}{帝國法院刑事判決彙編}}
\newcommand{\abbrRspr}{\abbrnoteonce{rspr}{Rspr.}{Rechtsprechung}{判例；司法實務}}
\newcommand{\abbrRVO}{\abbrnoteonce{rvo}{RVO}{Reichsversicherungsordnung}{帝國保險法}}
\newcommand{\abbrS}{\abbrnoteonce{s}{S.}{Seite}{頁}}
\newcommand{\abbrSFHG}{\abbrnoteonce{sfhg}{SFHG}{Schwangeren- und Familienhilfegesetz}{孕婦及家庭扶助法}}
\newcommand{\abbrSGBV}{\abbrnoteonce{sgbv}{SGB V}{Fünftes Buch Sozialgesetzbuch}{社會法典第五編}}
\newcommand{\abbrStAG}{\abbrnoteonce{staeg}{StÄG}{Strafrechtsänderungsgesetz}{刑法修正法}}
\newcommand{\abbrStenBer}{\abbrnoteonce{stenber}{StenBer.}{Stenographischer Bericht}{速記紀錄}}
\newcommand{\abbrStGB}{\abbrnoteonce{stgb}{StGB}{Strafgesetzbuch}{刑法}}
\newcommand{\abbrStREG}{\abbrnoteonce{streg}{StREG}{Strafrechtsreform-Ergänzungsgesetz}{刑法改革補充法}}
\newcommand{\abbrStrRG}{\abbrnoteonce{strrg}{StrRG}{Strafrechtsreformgesetz}{刑法改革法}}
\newcommand{\abbrUS}{\abbrnoteonce{us}{U.S.}{United States Reports}{美國最高法院判決彙編}}
\newcommand{\abbrVgl}{\abbrnoteonce{vgl}{vgl.}{vergleiche}{參見}}
\newcommand{\abbrVVDStRL}{\abbrnoteonce{vvdstrl}{VVDStRL}{Veröffentlichungen der Vereinigung der Deutschen Staatsrechtslehrer}{德國公法教師協會論文集}}
\newcommand{\abbrWP}{\abbrnoteonce{wp}{WP.}{Wahlperiode}{屆期}}
\newcommand{\abbrWp}{\abbrnoteonce{wp}{Wp.}{Wahlperiode}{屆期}}
\newcommand{\abbrZStrW}{\abbrnoteonce{zstrw}{ZStrW}{Zeitschrift für die gesamte Strafrechtswissenschaft}{整體刑法學期刊}}
\newcommand{\recordsourcerandstart}{%
  \ifrandnumbers
    \setcounter{lastsourcerandstart}{\value{randnumber}}%
    \global\lastsourcehasrandnumberstrue
  \else
    \global\lastsourcehasrandnumbersfalse
  \fi
}
\newlength{\biparagraphskip}
\setlength{\biparagraphskip}{0.52em}
\newlength{\lawboxblockbeforeskip}
\setlength{\lawboxblockbeforeskip}{0.72em}
\newlength{\lawboxblockafterskip}
\setlength{\lawboxblockafterskip}{0.92em}
\newif\iflawboxpartendedblock
\lawboxpartendedblockfalse
\newcommand{\sourceparstyle}{\footnotesize\color{sourceink}\setstretch{1.12}\RaggedRight\emergencystretch=3em\setlength{\parskip}{0.42em}\obeylines\selectfont}
\newcommand{\transparstyle}{\songfont\color{caseink}\setstretch{1.12}\setlength{\parindent}{0pt}\setlength{\parskip}{0.42em}}
\newcommand{\bilingualpairskip}{%
  \par\addvspace{\biparagraphskip}%
  \switchcolumn
  \par\addvspace{\biparagraphskip}%
  \switchcolumn*%
  \global\intranslationfalse
}
\newcommand{\bilingualboxpairskip}{%
  \par
  \switchcolumn
  \par
  \switchcolumn*%
  \global\intranslationfalse
}
\newcommand{\bilinguallawboxblockskip}{%
  \switchcolumn*%
  \par\addvspace{\lawboxblockafterskip}%
  \switchcolumn
  \par\addvspace{\lawboxblockafterskip}%
  \switchcolumn*%
  \global\intranslationfalse
}
\newcommand{\bikeepwithnext}[1][4]{%
  \ifbilingualcolumns
    \syncleft
    \Needspace{#1\baselineskip}%
    \switchcolumn
    \Needspace{#1\baselineskip}%
    \switchcolumn*%
    \global\intranslationfalse
  \else
    \Needspace{#1\baselineskip}%
  \fi
}
\NewEnviron{sourcepara}[1][]{
  \recordsourcerandstart
  \ifshoworiginal
    \ifbilingualcolumns
      \ifintranslation\switchcolumn*\global\intranslationfalse\fi
    \else
      \Needspace{5\baselineskip}
    \fi
    \ifbilingualcolumns\else\par\addvspace{0.35em}\fi
    {\sourceparstyle
    \noindent\BODY\par}
    \ifbilingualcolumns\else\addvspace{0.25em}\fi
    \ifbilingualcolumns
      \switchcolumn
      \global\intranslationtrue
    \fi
  \else
    \ifrandnumbers
      \begingroup
        \setbox0=\vbox{\hsize=\linewidth\sourceparstyle\noindent\BODY\par}%
      \endgroup
    \fi
  \fi
}
\NewEnviron{transpara}{
  \ifshowtranslation
    \setcounter{randnumberlocal}{\value{lastsourcerandstart}}%
    \ifbilingualcolumns
      \ifintranslation\else\switchcolumn\global\intranslationtrue\fi
      {\transparstyle\setlength{\leftskip}{0.55em}\setlength{\rightskip}{0.15em plus 1em}\addtolength{\linewidth}{-0.55em}\BODY\par}
      \switchcolumn*
      \bilingualpairskip
    \else
      {\transparstyle\BODY\par}
    \fi
  \else
    \ifbilingualcolumns
      \ifintranslation\switchcolumn*\global\intranslationfalse\fi
    \fi
  \fi
}
\NewEnviron{sourceboxpara}[1][]{
  \global\lawboxpartendedblockfalse
  \recordsourcerandstart
  \ifshoworiginal
    \ifbilingualcolumns
      \ifintranslation\switchcolumn*\global\intranslationfalse\fi
    \else
      \Needspace{3\baselineskip}
    \fi
    {\sourceparstyle
    \BODY\par}
    \ifbilingualcolumns\else
      \iflawboxpartendedblock\addvspace{\lawboxblockafterskip}\fi
    \fi
    \ifbilingualcolumns
      \switchcolumn
      \global\intranslationtrue
    \fi
  \fi
}
\NewEnviron{transboxpara}{
  \global\lawboxpartendedblockfalse
  \ifshowtranslation
    \setcounter{randnumberlocal}{\value{lastsourcerandstart}}%
    \ifbilingualcolumns
      \ifintranslation\else\switchcolumn\global\intranslationtrue\fi
      {\transparstyle\BODY\par}
      \iflawboxpartendedblock
        \bilinguallawboxblockskip
      \else
        \switchcolumn*
        \bilingualboxpairskip
      \fi
    \else
      {\transparstyle\BODY\par}
      \iflawboxpartendedblock\addvspace{\lawboxblockafterskip}\fi
    \fi
  \else
    \ifbilingualcolumns
      \ifintranslation\switchcolumn*\global\intranslationfalse\fi
    \fi
  \fi
}
\newcommand{\leitsatznum}[1]{\noindent\hangindent=3.0em\hangafter=1\makebox[2.25em][r]{\bfseries #1.}\hspace{0.55em}\ignorespaces}
\newcommand{\syncleft}{\ifbilingualcolumns\ifintranslation\switchcolumn*\global\intranslationfalse\fi\fi}
\pretocmd{\section}{\syncleft\clearpage}{}{}
\pretocmd{\subsection}{\syncleft}{}{}
\pretocmd{\subsubsection}{\syncleft}{}{}
\newcommand{\biheadingfont}{\songfont}
\newcommand{\bitocentry}[2]{%
  \ifshoworiginal
    \ifshowtranslation
      \ifstrequal{#1}{#2}{#1}{#1\space #2}%
    \else
      #1%
    \fi
  \else
    #2%
  \fi
}
\pdfstringdefDisableCommands{%
  \def\bitocentry#1#2{#1}%
}
\newcommand{\biheading}[7]{%
  \syncleft#1\bikeepwithnext[#3]\phantomsection\addcontentsline{toc}{#2}{\bitocentry{#6}{#7}}%
  \ifbilingualcolumns
    {#4 #6\par}%
    \switchcolumn
    {#4\biheadingfont #7\par}%
    \switchcolumn*%
    \global\intranslationfalse
    \par\addvspace{#5}%
    \switchcolumn
    \par\addvspace{#5}%
    \switchcolumn*%
  \else
    \ifshoworiginal{#4 #6\par}\fi
    \ifshowtranslation{#4\biheadingfont #7\par}\fi
    \addvspace{#5}%
  \fi
}
\newcommand{\termsectionheading}[1]{%
  \par\Needspace{9\baselineskip}%
  \vspace{0.65em}%
  {\songfont\bfseries #1\par}%
  \vspace{0.2em}%
}
% 章節換頁：
%   雙語 paracol 欄內不要用 \clearpage；它會在同步兩欄時吐出只有欄線/頁碼的空頁。
%   \newpage 足以讓下一個 section 起新頁，又不會強制清空所有待排材料。
%   單語模式不在 paracol 內，仍可用 \clearpage。
\newcommand{\bisectionpagebreak}{\ifbilingualcolumns\newpage\else\clearpage\fi}
\newcommand{\bisection}[2]{%
  \biheading{\iffirstbilingualsection\global\firstbilingualsectionfalse\else\bisectionpagebreak\fi}{section}{6}{\centering\Large\bfseries}{0.8em}{#1}{#2}%
}
\newcommand{\bisubsection}[2]{%
  \biheading{}{subsection}{5}{\large\bfseries}{0.55em}{#1}{#2}%
}
\newcommand{\bisubsubsection}[2]{%
  \biheading{}{subsubsection}{4}{\normalsize\bfseries}{0.45em}{#1}{#2}%
}
\newcommand{\bicompositeheading}[4]{%
  \biheading{}{subsection}{5}{\large\bfseries}{0.16em}{#1}{#2}%
  \biheading{}{subsubsection}{3}{\normalsize\bfseries}{0.45em}{#3}{#4}%
}
\newif\iflawboxfirstitem
\lawboxfirstitemfalse
\newcommand{\lawboxprespace}[1]{%
  \par
  \iflawboxfirstitem
    \global\lawboxfirstitemfalse
  \else
    \addvspace{#1}%
  \fi
}
\newtcolorbox{lawboxinner}{
  caseboxbase,
  colframe=sourceline!85!black,
  colback=white,
  left=1.05em,
  right=1.05em,
  top=0.62em,
  bottom=0.62em,
  before upper={\global\inlawboxtrue\global\lawboxfirstitemtrue\global\lawboxhaspendingrnfalse\ifintranslation\lawboxtransfont\else\lawboxsourcefont\fi\RaggedRight\setlength{\parindent}{0pt}\setlength{\parskip}{0pt}\ifintranslation\setstretch{\lawboxtransstretch}\else\setstretch{\lawboxsourcestretch}\fi},
  after upper={\global\lawboxhaspendingrnfalse\global\inlawboxfalse}
}
\tcbset{
  lawboxpartbase/.style={
    enhanced,
    breakable=false,
    width=0.985\linewidth,
    colback=white,
    frame hidden,
    boxrule=0pt,
    arc=0pt,
    left=1.05em,
    right=1.05em,
    boxsep=1.5mm,
    before skip=0pt,
    after skip=0pt,
    before upper={\global\inlawboxtrue\global\lawboxfirstitemtrue\global\lawboxhaspendingrnfalse\ifintranslation\lawboxtransfont\else\lawboxsourcefont\fi\RaggedRight\setlength{\parindent}{0pt}\setlength{\parskip}{0pt}\ifintranslation\setstretch{\lawboxtransstretch}\else\setstretch{\lawboxsourcestretch}\fi},
    after upper={\global\lawboxhaspendingrnfalse\global\inlawboxfalse},
    borderline west={0.28pt}{0pt}{sourceline!85!black},
    borderline east={0.28pt}{0pt}{sourceline!85!black},
  },
  lawboxpartfirst/.style={
    lawboxpartbase,
    top=0.62em,
    bottom=0.04em,
    before skip=\lawboxblockbeforeskip,
    borderline north={0.28pt}{0pt}{sourceline!85!black},
  },
  lawboxpartmiddle/.style={
    lawboxpartbase,
    top=0.04em,
    bottom=0.04em,
  },
  lawboxpartlast/.style={
    lawboxpartbase,
    top=0.04em,
    bottom=0.62em,
    after skip=0pt,
    borderline south={0.28pt}{0pt}{sourceline!85!black},
  }
}
\newtcolorbox{lawboxpartinner}[1]{#1}
\newtcolorbox{termboxinner}{
  caseboxbase,
  colframe=noteblue!70!black,
  colback=casepaper,
  before upper={\songfont\setlength{\parindent}{0pt}\setlength{\parskip}{0.35em}}
}
\newtcolorbox{deboxinner}{
  caseboxbase,
  colframe=notegray!72!black,
  colback=white,
  left=1.05em,
  right=1.05em,
  top=0.62em,
  bottom=0.62em,
  before upper={\global\inlawboxtrue\global\lawboxfirstitemtrue\global\lawboxhaspendingrnfalse\small\setlength{\parindent}{0pt}\setlength{\parskip}{0.35em}},
  after upper={\global\lawboxhaspendingrnfalse\global\inlawboxfalse}
}
\NewEnviron{lawbox}{
  \begin{lawboxinner}\BODY\end{lawboxinner}
}
\NewEnviron{lawboxpart}[2]{
  \begin{lawboxpartinner}{lawboxpart#1,equal height group=#2}\BODY\end{lawboxpartinner}%
  \ifstrequal{#1}{last}{\global\lawboxpartendedblocktrue}{}%
}
\NewEnviron{termbox}{
  \par\addvspace{0.55em}
  {\color{noteblue!70!black}\rule{\linewidth}{0.28pt}}\par
  \begingroup
    \songfont\small\color{caseink}
    \setlength{\parindent}{0pt}
    \setlength{\parskip}{0.24em}
    \BODY
    \par
  \endgroup
  {\color{noteblue!70!black}\rule{\linewidth}{0.28pt}}\par
  \addvspace{0.55em}
}
\NewEnviron{debox}{
  \begin{deboxinner}\BODY\end{deboxinner}
}
\providecommand{\paragraphmark}[1]{}
\newcommand{\lawboxtitleafter}{\addvspace{0.06em}}
\newcommand{\lawtitle}[1]{\lawboxprespace{0.22em}{\lawboxtransfont\bfseries\lawboxuserandnumber #1}\par\lawboxtitleafter}
\newcommand{\absno}[2]{\lawboxprespace{0.04em}\noindent\lawboxuserandnumber\hangindent=2.6em\hangafter=1\makebox[2.6em][l]{(#1)}#2\par}
\newcommand{\nrno}[2]{\lawboxprespace{0pt}\noindent\lawboxuserandnumber\hspace*{2.6em}\hangindent=4.4em\hangafter=1\makebox[1.8em][l]{#1.}#2\par}
\newcommand{\letterno}[2]{\lawboxprespace{0pt}\noindent\lawboxuserandnumber\hspace*{4.4em}\hangindent=6.0em\hangafter=1\makebox[1.6em][l]{#1)}#2\par}
\newcommand{\sourcelawtitle}[1]{\lawboxprespace{0.22em}{\lawboxsourcefont\bfseries\lawboxuserandnumber #1}\par\lawboxtitleafter}
\newcommand{\sourcelawsubtitle}[1]{\lawboxprespace{0.04em}{\lawboxsourcefont\bfseries\lawboxuserandnumber #1}\par\lawboxtitleafter}
\newcommand{\sourceabsno}[2]{\lawboxprespace{0.04em}\noindent\lawboxuserandnumber\hangindent=2.45em\hangafter=1\makebox[2.45em][l]{(#1)}#2\par}
\newcommand{\sourcenrno}[2]{\lawboxprespace{0pt}\noindent\lawboxuserandnumber\hspace*{2.45em}\hangindent=4.25em\hangafter=1\makebox[1.8em][l]{#1.}#2\par}
\newcommand{\sourceletterno}[2]{\lawboxprespace{0pt}\noindent\lawboxuserandnumber\hspace*{4.25em}\hangindent=5.85em\hangafter=1\makebox[1.6em][l]{#1)}#2\par}
\newcommand{\sourcequotepara}[1]{\lawboxprespace{0.04em}\noindent\lawboxuserandnumber\hangindent=1.2em\hangafter=1 #1\par}
\newcommand{\sourcelawline}[1]{\lawboxprespace{0.04em}\noindent\lawboxuserandnumber #1\par}
\newcommand{\lawline}[1]{\lawboxprespace{0.04em}\noindent\lawboxuserandnumber #1\par}
\newcommand{\pendingtranslation}{\par{\songfont\footnotesize\color{pendingink}（中文譯文待補。）}\par}
\newcommand{\tocpart}[1]{\par\addvspace{0.52em}{\footnotesize\bfseries\color{pendingink}#1\par}\addvspace{0.16em}}
\newcommand{\tocdashline}{\par\addvspace{0.48em}\noindent{\color{notegray}\xleaders\hbox to 0.82em{\hfil\rule{0.42em}{0.28pt}\hfil}\hskip\linewidth}\par\addvspace{0.38em}}
% \tocpanel{font-cmd}{content}：將目錄置中於所在欄寬（雙語欄適用）。
% A3 橫式使用 0.58\linewidth，A4 橫式使用 0.84\linewidth，
% 使兩種紙張下目錄欄內容實際寬度相近（均約 100–105 mm）。
\newcommand{\tocpanel}[2]{%
  \makebox[\linewidth][c]{%
    \ifx\bilingualpaper\bilingualpaperaThree
      \begin{minipage}{0.58\linewidth}%
    \else
      \begin{minipage}{0.84\linewidth}%
    \fi
    \toccolstyle
    #1#2%
    \end{minipage}%
  }%
}
% ===== 首頁簡目錄的兩層 description list 縮排 =====
% enumitem 的 description list 可把每一列拆成：
%   [label]  labelsep  body text
% 這裡的「起點」都以所在 minipage/欄位的左緣為 0。
%
% 核心規則：
%   labelindent = label 欄位從左緣開始的位置。
%   labelwidth  = label 欄位本身寬度；標籤太長（如 Abw. Meinung）就加大它。
%   labelsep    = label 欄位右緣到正文的距離。
%   leftmargin  = 正文 body text 的起點。判斷層級視覺時主要看這個。
%
% 所以：
%   想讓「Begründung / 判決理由」整列正文往右或往左：改 tocitems 的 leftmargin。
%   想讓「A. Verfahren und Gegenstand / A. 程序與審查標的」正文往右或往左：
%     改 tocsubitems 的 leftmargin。
%   想讓 A./B./C. 這些標籤本身往右或往左，但正文不動：改 tocsubitems 的 labelindent。
%   想讓 A./B./C. 與正文間距變大或變小：改 tocsubitems 的 labelsep。
%
% 層級原則：
%   tocsubitems.leftmargin 必須大於 tocitems.leftmargin，
%   否則子層級正文會看起來比「Gründe / 理由」還靠左。
\newlist{tocitems}{description}{1}
\setlist[tocitems]{%
  labelindent=0pt,    % 第一層 label 欄位起點；0pt 表示貼齊目錄欄左緣。
  leftmargin=2.54cm,  % 第一層正文起點；例如 Begründung / 判決理由 的左緣。
  labelwidth=2.16cm,  % 第一層 label 欄位寬度；需容納 Leitsätze、Abw. Meinung、不同意見等。
  labelsep=0.32cm,    % 第一層 label 與正文的水平間距；只改間距，不改正文起點。
  itemsep=0.28em,     % 第一層各項之間的垂直距離。
  topsep=0.20em,      % 第一層 list 上下與前後內容的垂直距離。
  parsep=0pt,         % 同一 item 內段落之間的距離；目錄項通常不需要。
  partopsep=0pt,      % list 若緊接段落開始時的額外距離；關閉以免目錄高度飄動。
  align=left,         % label 欄內左對齊；長短標籤不靠右貼齊。
  font=\normalfont\bfseries % label 的字型；正文不受這行影響。
}
\newlist{tocsubitems}{description}{1}
\setlist[tocsubitems]{%
  labelindent=1cm,    % 第二層 label 起點；只控制 A./B./C. 本身的位置。
  leftmargin=3.00cm,  % 第二層正文起點；必須大於 2.54cm，才會比 Begründung / 判決理由更右。
  labelwidth=0.5cm,   % 第二層 label 欄寬；A./B./C. 很短，可小於第一層。
  labelsep=1.5cm,    % 第二層 label 與正文間距；A. 與 Verfahren / 程序 的距離看這裡。
  itemsep=0.12em,     % 第二層各項較密，避免 A.–E. 佔太高。
  topsep=0.04em,      % 第二層 list 與 Gründe / 理由之間的垂直距離。
  parsep=0pt,         % 同一第二層 item 內段落距離；目錄項通常不需要。
  partopsep=0pt,      % 關閉額外段落起始距離，避免版面不穩。
  align=left,         % A./B./C. label 欄內左對齊。
  font=\normalfont\bfseries, % A./B./C. label 加粗；正文不受這行影響。
  before={}           % 不在第二層 list 前自動插入額外內容。
}
\newlist{termitems}{description}{1}
\ifbilingualcolumns
  \setlist[termitems]{%
    leftmargin=5.2cm,
    labelwidth=4.8cm,
    labelsep=0.35cm,
    itemsep=0.16em,
    topsep=0.2em,
    parsep=0pt,
    partopsep=0pt,
    font=\normalfont\bfseries,
    style=nextline
  }
\else
  \setlist[termitems]{%
    leftmargin=0pt,
    labelwidth=0pt,
    labelsep=0pt,
    itemsep=0.24em,
    topsep=0.2em,
    parsep=0pt,
    partopsep=0pt,
    font=\normalfont\bfseries,
    style=nextline
  }
\fi
\newlist{abbritems}{description}{1}
\setlist[abbritems]{%
  leftmargin=2.38cm,
  labelwidth=2.02cm,
  labelsep=0.28cm,
  itemsep=0.16em,
  topsep=0.2em,
  parsep=0pt,
  partopsep=0pt,
  align=left,
  font=\normalfont\bfseries
}
\ifbilingualcolumns
  \newenvironment{termcolumns}{\begin{multicols}{2}}{\end{multicols}}
\else
  \newenvironment{termcolumns}{}{}
\fi

% ===== 編者註腳開關 =====
% 一般版預設不顯示編者註，避免正文腳註過密。
% 若開啟 \classroommodeon，GG/條文編者註仍會自動顯示，無需另開本開關。
% 若一般版也要顯示編者註，可在單案檔 \input 後加 \showeditornotestrue。
\newif\ifshoweditornotes
\showeditornotesfalse

% ===== 目錄排版輔助 =====
\newcommand{\toccolstyle}{\raggedright\arraybackslash}
\newcommand{\tocsingleindent}{\leftskip=1.45cm\rightskip=1.2cm}

% ===== 行內引文環境（德文原文與中文譯文各一，帶左側灰色邊線）=====
\newcommand{\quoteparbreak}{\par\addvspace{0.48em}}
\newcommand{\sourcequoteinner}[1]{%
  \begin{tcolorbox}[
    enhanced, breakable, width=\dimexpr\linewidth-0.8em\relax, frame hidden,
    colback=white, coltext=caseink, boxrule=0pt,
    borderline west={0.55pt}{0.88em}{notegray},
    left=1.78em, right=0.72em, top=0.18em, bottom=0.18em,
    boxsep=0pt, before skip=0.24em, after skip=0.24em,
    before upper={\normalfont\footnotesize\color{caseink}\setlength{\parindent}{0pt}\setlength{\parskip}{0.34em}\raggedright\setlength{\rightskip}{0pt plus 1em}}
  ]%
  #1\par
  \end{tcolorbox}%
}
\newcommand{\transquoteinner}[1]{%
  \begin{tcolorbox}[
    enhanced, breakable, width=\dimexpr\linewidth-0.8em\relax, frame hidden,
    colback=white, coltext=caseink, boxrule=0pt,
    borderline west={0.55pt}{0.88em}{notegray},
    left=1.78em, right=0.72em, top=0.18em, bottom=0.18em,
    boxsep=0pt, before skip=0.24em, after skip=0.24em,
    before upper={\songfont\small\color{caseink}\setlength{\parindent}{0pt}\setlength{\parskip}{0.34em}\raggedright\setlength{\rightskip}{0pt plus 1em}}
  ]%
  #1\par
  \end{tcolorbox}%
}

% ===== 大綱（Gliederung）Rn 輔助命令 =====
\newcommand{\outrn}[2]{\enspace{\normalfont\footnotesize\color{pendingink}(Rn\,#1–#2)}}
\newcommand{\outrnsingle}[1]{\enspace{\normalfont\footnotesize\color{pendingink}(Rn\,#1)}}
\newcommand{\outrnzh}[2]{\enspace{\normalfont\footnotesize\color{pendingink}（第\,#1–#2\,段）}}
\newcommand{\outrnzhs}[1]{\enspace{\normalfont\footnotesize\color{pendingink}（第\,#1\,段）}}
% ===== 大綱雙語對照表格輔助命令（三欄：段碼 │ 德文 │ 中文）=====
% 欄 1：段碼（由欄定義自動格式化：小字、等寬、右對齊、pendingink）
% 欄 2：德文標籤＋正文
% 欄 3：中文標籤＋正文
%
% 無段碼的列（\omainrow、\osecrow），欄 1 以 & 空置。
% 有段碼的列（\osecrowrn、\osubrow、\ointrow），#1 為預格式化字串
%   例：\osubrow{2–49}{I.}{...}{I.}{...}
%       \ointrow{107}{...}{...}
%
% \opartrow：段組標題（橫跨三欄）
\newcommand{\opartrow}[2]{%
  \noalign{\vspace{0.30em}}%
  \multicolumn{3}{@{}l@{}}{%
    \footnotesize\bfseries\color{pendingink}#1\hfill\songfont#2%
  }\\[0.06em]%
}
% \odashrow：虛線分隔（橫跨三欄）
\newcommand{\odashrow}{%
  \noalign{\vspace{0.28em}}%
  \multicolumn{3}{@{}l@{}}{\color{notegray}%
    \xleaders\hbox to 0.82em{\hfil\rule{0.42em}{0.28pt}\hfil}\hskip 0.88\linewidth}%
  \\[0.24em]%
}
% \odissentpart：不同意見書區段標頭。比一般 \opartrow 更像附錄性轉場。
\newcommand{\odissentpart}[2]{%
  \noalign{\vspace{0.42em}}%
  \multicolumn{3}{@{}p{0.88\textwidth}@{}}{%
    \color{notegray}\rule{\linewidth}{0.18pt}%
  }\\[-0.18em]%
  \multicolumn{3}{@{}p{0.88\textwidth}@{}}{%
    \makebox[\linewidth][s]{%
      \footnotesize\bfseries\color{pendingink}#1%
      \hfill{\songfont #2}%
    }%
  }\\[-0.02em]%
}
% \odissentopinion：不同意見書的署名與一行摘要，右欄保留段碼範圍。
\newcommand{\odissentopinion}[5]{%
  \footnotesize\hangindent=0.52cm\hangafter=1%
    {\bfseries\color{caseink}#2}\enspace{\color{notegray}\textperiodcentered}\enspace\color{caseink}#3%
  &\footnotesize\hangindent=0.52cm\hangafter=1%
    {\songfont\bfseries\color{caseink}#4}\enspace{\color{notegray}\textperiodcentered}\enspace\songfont\color{caseink}#5%
  &#1%
  \\[0.12em]%
}
% \omainrow：頂層項目，無段碼（Leitsätze、Urteil、Gründe …）
\newcommand{\omainrow}[4]{%
  \footnotesize\hangindent=1.92cm\hangafter=1%
    \makebox[1.92cm][l]{\bfseries\color{caseink}#1}\color{caseink}#2%
  &\footnotesize\hangindent=1.92cm\hangafter=1%
    \makebox[1.92cm][l]{\bfseries\color{caseink}#3}\songfont\color{caseink}#4%
  &%
  \\[0.15em]%
}
% \omainrowrn：頂層項目，有段碼（不同意見等標籤寬於 \osecrow 框者）
\newcommand{\omainrowrn}[5]{%
  \footnotesize\hangindent=1.92cm\hangafter=1%
    \makebox[1.92cm][l]{\bfseries\color{caseink}#2}\color{caseink}#3%
  &\footnotesize\hangindent=1.92cm\hangafter=1%
    \makebox[1.92cm][l]{\bfseries\color{caseink}#4}\songfont\color{caseink}#5%
  &#1%
  \\[0.15em]%
}
% \osecrow：一級子項，無段碼（A.、C.、D. 等有子項者）
\newcommand{\osecrow}[4]{%
  \footnotesize\hspace{1.10cm}\hangindent=1.98cm\hangafter=1%
    \makebox[0.86cm][l]{\bfseries\color{caseink}#1}\color{caseink}#2%
  &\footnotesize\hspace{1.10cm}\hangindent=1.98cm\hangafter=1%
    \makebox[0.86cm][l]{\bfseries\color{caseink}#3}\songfont\color{caseink}#4%
  &%
  \\[0.11em]%
}
% \osecrowrn：一級子項，有段碼（B.、E. 等完整段落範圍者）
% #1=段碼字串（如 101–106）, #2=DE標籤, #3=DE正文, #4=ZH標籤, #5=ZH正文
\newcommand{\osecrowrn}[5]{%
  \footnotesize\hspace{1.10cm}\hangindent=1.98cm\hangafter=1%
    \makebox[0.86cm][l]{\bfseries\color{caseink}#2}\color{caseink}#3%
  &\footnotesize\hspace{1.10cm}\hangindent=1.98cm\hangafter=1%
    \makebox[0.86cm][l]{\bfseries\color{caseink}#4}\songfont\color{caseink}#5%
  &#1%
  \\[0.11em]%
}
% \osubrow：二級子項，有段碼範圍（I.、II.、A.I. …）
% #1=段碼字串, #2=DE標籤, #3=DE正文, #4=ZH標籤, #5=ZH正文
\newcommand{\osubrow}[5]{%
  \footnotesize\hspace{1.40cm}\hangindent=2.30cm\hangafter=1%
    \makebox[0.90cm][l]{\bfseries\color{caseink}#2}\color{caseink}#3%
  &\footnotesize\hspace{1.40cm}\hangindent=2.30cm\hangafter=1%
    \makebox[0.90cm][l]{\bfseries\color{caseink}#4}\songfont\color{caseink}#5%
  &#1%
  \\[0.09em]%
}
% \ointrow：引入段落，有單一段碼，無標籤
% #1=段碼字串, #2=DE正文, #3=ZH正文
\newcommand{\ointrow}[3]{%
  \footnotesize\hspace{0.52cm}\hangindent=0.52cm\hangafter=1\color{caseink}#2%
  &\footnotesize\hspace{0.52cm}\hangindent=0.52cm\hangafter=1\songfont\color{caseink}#3%
  &#1%
  \\[0.09em]%
}
% \osubsubrow：三級子項（1.、2.、3.）
% #1=段碼字串, #2=DE標籤, #3=DE正文, #4=ZH標籤, #5=ZH正文
\newcommand{\osubsubrow}[5]{%
  \footnotesize\hspace{1.80cm}\hangindent=2.48cm\hangafter=1%
    \makebox[0.68cm][l]{\bfseries\color{caseink}#2}\color{caseink}#3%
  &\footnotesize\hspace{1.80cm}\hangindent=2.48cm\hangafter=1%
    \makebox[0.68cm][l]{\bfseries\color{caseink}#4}\songfont\color{caseink}#5%
  &#1%
  \\[0.08em]%
}
% \osubsubsubrow：四級子項（a）、b））
% #1=段碼字串, #2=DE標籤, #3=DE正文, #4=ZH標籤, #5=ZH正文
\newcommand{\osubsubsubrow}[5]{%
  \footnotesize\hspace{2.10cm}\hangindent=2.68cm\hangafter=1%
    \makebox[0.58cm][l]{\bfseries\color{caseink}#2}\color{caseink}#3%
  &\footnotesize\hspace{2.10cm}\hangindent=2.68cm\hangafter=1%
    \makebox[0.58cm][l]{\bfseries\color{caseink}#4}\songfont\color{caseink}#5%
  &#1%
  \\[0.07em]%
}
% \osubsubsubsubrow：五級子項（aa）、bb））
% 標籤框 0.62cm：足以容納「aa)」在 footnotesize bold 下（約 0.50cm）並留有間距
% #1=段碼字串, #2=DE標籤, #3=DE正文, #4=ZH標籤, #5=ZH正文
\newcommand{\osubsubsubsubrow}[5]{%
  \footnotesize\hspace{2.38cm}\hangindent=3.06cm\hangafter=1%
    \makebox[0.68cm][l]{\bfseries\color{caseink}#2}\color{caseink}#3%
  &\footnotesize\hspace{2.38cm}\hangindent=3.06cm\hangafter=1%
    \makebox[0.68cm][l]{\bfseries\color{caseink}#4}\songfont\color{caseink}#5%
  &#1%
  \\[0.06em]%
}
% \outlinepage：雙語對照大綱頁包裝環境（longtable 自動跨頁）
% 三欄：德文欄 + 中文欄 + 段碼欄（最右，不折行）
% 段碼欄寬度：1.80cm，足以容納最長段碼如「309–348」
% DE/ZH 欄寬：(0.87\textwidth - 2.08cm) / 2
%   驗算：2×欄 + 0.05\textwidth + 0.28cm gap + 1.80cm = 0.87tw - 2.08cm + 0.05tw + 2.08cm = 0.92tw ✓
\newenvironment{outlinepage}{%
  \vspace*{0.40cm}%
  \setlength{\LTleft}{0.04\textwidth}%
  \setlength{\LTright}{0.04\textwidth}%
  \setlength{\tabcolsep}{0pt}%
  % 大綱列距由 longtable 的 \arraystretch 與各 row macro 結尾的 \\[...em] 共同決定。
  % 這裡控制「每一列本身的表格 strut 高度」；數值越大，A./I./1. 各項之間越像有空行。
  % A3 原本用 0.62，視覺上沒有空行；A4 也沿用同值，避免切到 A4 後項目被撐開。
  % 若日後覺得大綱太擠，只微調這個值（例如 0.68），不要改各 \omainrow/\osubrow 的縮排。
  \renewcommand{\arraystretch}{0.62}%
  \begin{longtable}{%
    >{\vspace{0pt}\raggedright\arraybackslash}p{\dimexpr(0.87\textwidth-2.08cm)/2\relax}%
    @{\hspace{0.025\textwidth}}%
    !{\color{notegray}\vrule width 0.12pt}%
    @{\hspace{0.025\textwidth}}%
    >{\vspace{0pt}\raggedright\arraybackslash}p{\dimexpr(0.87\textwidth-2.08cm)/2\relax}%
    @{\hspace{0.28cm}}%
    >{\vspace{0pt}\footnotesize\normalfont\color{pendingink}\raggedright\arraybackslash}p{1.80cm}%
    @{}%
  }%
  % 首頁欄標籤：不要橫跨置中；分別放在德文欄與中文欄的實際欄位起點。
  \footnotesize\bfseries\color{sourceink}Gliederung%
  &\footnotesize\bfseries\songfont\color{sourceink}大綱%
  &%
  \\[0.36em]%
  \endfirsthead%
  % 續頁欄標籤：同樣分欄定位，以灰色細體區分；無需標注「續」。
  \footnotesize\color{pendingink}Gliederung%
  &\footnotesize\songfont\color{pendingink}大綱%
  &%
  \\[0.24em]%
  \endhead%
}{%
  \end{longtable}%
}
 之前設定；
%   預設 chinese / a4）
\ifdefined\docmode\else\def\docmode{chinese}\fi
\ifdefined\bilingualpaper\else\def\bilingualpaper{a4}\fi
\newif\ifshoworiginal
\newif\ifshowtranslation
\newif\ifbilinguallandscape
\newif\ifbilingualcolumns
\newif\ifintranslation
\newif\iffirstbilingualsection

\showoriginalfalse
\showtranslationfalse
\bilinguallandscapefalse
\bilingualcolumnsfalse
\intranslationfalse
\firstbilingualsectionfalse

\def\modechinese{chinese}
\def\modegerman{german}
\def\modebilingual{bilingual}
\def\bilingualpaperaThree{a3}
\def\bilingualpaperaFour{a4}

\ifx\docmode\modechinese
  \showtranslationtrue
\fi

\ifx\docmode\modegerman
  \showoriginaltrue
\fi

\ifx\docmode\modebilingual
  \showoriginaltrue
  \showtranslationtrue
  \bilinguallandscapetrue
  \bilingualcolumnstrue
\fi

% 編排模式說明：
% chinese  → \showoriginalfalse  \showtranslationtrue  （A4 直式，純中文）
% german   → \showoriginaltrue   \showtranslationfalse （A4 直式，純德文）
% bilingual→ \showoriginaltrue   \showtranslationtrue  \bilingualcolumnstrue（A3/A4 橫式對照）

\ifbilingualcolumns
  \ifx\bilingualpaper\bilingualpaperaFour
    \geometry{
      a4paper,
      landscape,
      left=1.25cm,
      right=2.25cm,
      top=1.25cm,
      bottom=1.75cm,
      footskip=8.5mm,
      marginparwidth=0.75cm,
      marginparsep=0.25cm
    }
  \else
    \geometry{
      a3paper,
      landscape,
      left=1.55cm,
      right=2.75cm,
      top=1.55cm,
      bottom=2.05cm,
      footskip=9.5mm,
      marginparwidth=0.9cm,
      marginparsep=0.35cm
    }
  \fi
\else
\geometry{
  a4paper,
  portrait,
  left=2.2cm,
  right=2.6cm,
  top=2.0cm,
  bottom=2.0cm,
  marginparwidth=0.8cm,
  marginparsep=0.3cm
}
\fi
% ===== 時間戳基礎設施 =====
\newcount\stamptimeval
\newcount\stampminuteval
\newcommand{\stamphh}{%
  \stamptimeval=\time
  \divide\stamptimeval by 60
  \ifnum\stamptimeval<10 0\fi
  \the\stamptimeval}
\newcommand{\stampmm}{%
  \stamptimeval=\time
  \stampminuteval=\time
  \divide\stampminuteval by 60
  \multiply\stampminuteval by 60
  \advance\stamptimeval by -\stampminuteval
  \ifnum\stamptimeval<10 0\fi
  \the\stamptimeval}
\newcommand{\stampwkde}[1]{%
  \ifcase#1Montag\or Dienstag\or Mittwoch\or Donnerstag\or Freitag\or Samstag\or Sonntag\fi}
\newcommand{\stampwkzh}[1]{%
  \ifcase#1 一\or 二\or 三\or 四\or 五\or 六\or 日\fi}
\newcommand{\stampmonthde}[1]{%
  \ifcase#1\or Januar\or Februar\or März\or April\or Mai\or Juni\or
  Juli\or August\or September\or Oktober\or November\or Dezember\fi}
\newcount\stampjdval
\newcount\stampwdval
\pgfcalendardatetojulian{\the\year-\the\month-\the\day}{\stampjdval}
\pgfcalendarjuliantoweekday{\the\stampjdval}{\stampwdval}
\newcommand{\timestampzh}{%
  \songfont 最後修改：\the\year 年\the\month 月\the\day 日%
  % （星期\stampwkzh{\the\stampwdval}）%
  % \space\stamphh:\stampmm
  }

\newcommand{\documentdatezh}{\the\year 年 \the\month 月 \the\day 日}
\newcommand{\documentdatede}{\the\day. \stampmonthde{\the\month} \the\year}
\newcommand{\bverfgetranslationnote}{%
  {\songfont 翻譯說明：初譯使用 Claude Code 與 GPT 協助迭代；仍請以德文原文為準。}%
}
\newcommand{\bverfgecourseinfo}{%
  {\songfont 課程：114 學年度第 2 學期；憲法專題研究（四）；授課老師：湯德宗教授；導讀日期：115 年 6 月 1 日。}%
}
\newcommand{\bverfgeeditorinfo}{%
  {\songfont 編者：王逸帆（R10A21126），國立台灣大學法律系研究所碩士班。}%
}
\newcommand{\bverfgetitleversionline}{%
  \ifbilingualcolumns
    Fassung \documentversion\enspace\textperiodcentered\enspace Stand: \documentdatede
    \quad
    {\songfont 版本 \documentversion\enspace\textperiodcentered\enspace 修訂日期：\documentdatezh\par}%
  \else
    \ifshowtranslation
      {\songfont 版本 \documentversion\enspace\textperiodcentered\enspace 修訂日期：\documentdatezh\par}%
    \else
      Fassung \documentversion\enspace\textperiodcentered\enspace Stand: \documentdatede\par
    \fi
  \fi
}
\newcommand{\bverfgetitlecornerinfo}{%
  \ifshowtranslation
    \vspace{0.72em}%
    \noindent\hfill
    \begin{minipage}{0.66\textwidth}
      \raggedleft\tiny\color{pendingink}\songfont
      \bverfgecourseinfo\par
      \bverfgeeditorinfo
    \end{minipage}%
  \fi
}
\newcommand{\bverfgetitlemeta}{%
  \vfill
  \begin{center}%
    \footnotesize\color{pendingink}%
    \bverfgetitleversionline
    \ifshowtranslation
      \vspace{0.28em}{\scriptsize\bverfgetranslationnote\par}%
    \fi
  \end{center}%
  \bverfgetitlecornerinfo
}

\pagestyle{fancy}
\fancyhf{}
\renewcommand{\headrulewidth}{0pt}
\renewcommand{\footrulewidth}{0pt}
\fancyfoot[C]{\thepage}
\ifbilingualcolumns
  \fancyfoot[R]{\scriptsize\color{pendingink}\timestampzh}%
\else
  \ifshoworiginal
  \else
    \fancyfoot[R]{\scriptsize\color{pendingink}\timestampzh\hspace{0.45cm}}%
  \fi
\fi
\setmainfont{Times New Roman}
\setsansfont{Helvetica Neue}
\setmonofont{Menlo}
\IfFileExists{/Users/iw/Library/Fonts/kaiu.ttf}{
  \setCJKmainfont[Path=/Users/iw/Library/Fonts/,AutoFakeBold=4]{kaiu.ttf}
}{
  \IfFontExistsTF{標楷體}{
    \setCJKmainfont[AutoFakeBold=4]{標楷體}
  }{
    \setCJKmainfont{Songti TC}
  }
}
\IfFontExistsTF{Songti TC}{
  \setCJKfamilyfont{song}{Songti TC}
  \setCJKmonofont{Songti TC}
}{
  \setCJKfamilyfont{song}[Path=/Users/iw/Library/Fonts/,AutoFakeBold=4]{kaiu.ttf}
  \setCJKmonofont[Path=/Users/iw/Library/Fonts/,AutoFakeBold=4]{kaiu.ttf}
}
\newcommand{\songfont}{\CJKfamily{song}}

% ===== 封面／目錄專用打字機字體（僅限封面、目錄頁使用，不影響正文）=====
\IfFontExistsTF{Radio Newsman}{%
  \newfontfamily\bverfgetitlede{Radio Newsman}%
}{%
  \newfontfamily\bverfgetitlede{Courier New}%
}
\IfFontExistsTF{Erikas Farbband}{%
  \newfontfamily\bverfgebodyde{Erikas Farbband}%
}{%
  \let\bverfgebodyde\bverfgetitlede
}
\IfFontExistsTF{Maoken Heavy Labourer}{%
  \setCJKfamilyfont{bverfgetitlecjk}{Maoken Heavy Labourer}%
}{%
  \setCJKfamilyfont{bverfgetitlecjk}{Songti TC}%
}
\IfFontExistsTF{Huiwen-mincho}{%
  \setCJKfamilyfont{bverfgebodycjk}{Huiwen-mincho}%
}{%
  \setCJKfamilyfont{bverfgebodycjk}{Songti TC}%
}
\newcommand{\bverfgetitlezh}{\bverfgetitlede\CJKfamily{bverfgetitlecjk}}
\newcommand{\bverfgebodyzh}{\bverfgebodyde\CJKfamily{bverfgebodycjk}}

\setstretch{1.35}
\setlist{itemsep=0.2em, topsep=0.4em}
\setlength{\parindent}{0pt}
\setlength{\parskip}{0.45em}
\raggedbottom
\setcounter{secnumdepth}{0}
\setcounter{tocdepth}{2}

\newcommand{\lawboxsourcefont}{\footnotesize}
\newcommand{\lawboxtransfont}{\footnotesize}
\newcommand{\lawboxsourcestretch}{1.02}
\newcommand{\lawboxtransstretch}{0.92}

\titleformat{\section}{\centering\Large\bfseries\songfont}{\thesection}{1em}{}
\titleformat{\subsection}{\large\bfseries\songfont}{\thesubsection}{1em}{}
\titleformat{\subsubsection}{\normalsize\bfseries\songfont}{\thesubsubsection}{1em}{}
\titleformat{\paragraph}{\normalsize\bfseries\songfont}{\theparagraph}{1em}{}
\titlespacing*{\section}{0pt}{1.8em}{1.0em}
\titlespacing*{\subsection}{0pt}{1.2em}{0.6em}
\titlespacing*{\subsubsection}{0pt}{1.0em}{0.45em}
\titlespacing*{\paragraph}{0pt}{0.6em}{0.4em}

\definecolor{noteblue}{HTML}{A9C9F5}
\definecolor{noteteal}{HTML}{AEE1D6}
\definecolor{notegray}{HTML}{D7DADF}
\definecolor{caseink}{HTML}{000000}
\definecolor{casepaper}{HTML}{FCFEFD}
\definecolor{sourcepaper}{HTML}{FAFBF8}
\definecolor{sourceline}{HTML}{D7DED8}
\definecolor{sourceink}{HTML}{000000}
\definecolor{pendingink}{HTML}{000000}

\tcbset{
  caseboxbase/.style={
    enhanced,
    breakable,
    width=0.985\linewidth,
    colback=casepaper,
    coltitle=caseink,
    fonttitle=\bfseries\songfont,
    boxrule=0.28pt,
    arc=1.2mm,
    left=2.4mm,
    right=2.4mm,
    top=1.5mm,
    bottom=1.5mm,
    boxsep=1.5mm,
    before skip=0.65em,
    after skip=0.65em,
    before upper={\setlength{\parindent}{0pt}\setlength{\parskip}{0.35em}},
  }
}

\newcommand{\bilingualstart}{}
\newcommand{\bilingualstop}{}
\ifbilingualcolumns
  \renewcommand{\bilingualstart}{\small\setlength{\columnsep}{1.25cm}\setlength{\columnseprule}{0.12pt}\columnratio{0.5,0.5}\multicolumnfootnotes\flushbottom\sloppy\interlinepenalty=80\clubpenalty=800\widowpenalty=800\displaywidowpenalty=800\begin{paracol}{2}\global\intranslationfalse\global\firstbilingualsectiontrue{\footnotesize\color{sourceink}Deutsch\par}\switchcolumn{\footnotesize\songfont\color{sourceink}中文譯文\par}\switchcolumn*}
  \renewcommand{\bilingualstop}{\ifintranslation\switchcolumn*\global\intranslationfalse\fi\end{paracol}\clearpage\raggedbottom}
\fi
\newif\ifrandnumbers
\randnumbersfalse
\newcounter{randnumber}
\newcounter{randnumberlocal}
\newcounter{lastsourcerandstart}
\newif\iflastsourcehasrandnumbers
\lastsourcehasrandnumbersfalse
\newif\ifinlawbox
\inlawboxfalse
\newif\iflawboxhaspendingrn
\lawboxhaspendingrnfalse
\newcommand{\lawboxpendingrn}{}
\newcommand{\startrandnumbers}{\global\randnumberstrue\setcounter{randnumber}{1}\global\lastsourcehasrandnumbersfalse}
\newcommand{\stoprandnumbers}{\global\randnumbersfalse\global\lastsourcehasrandnumbersfalse}
\newlength{\randnumberwidth}
\setlength{\randnumberwidth}{0.95cm}
\newlength{\randnumberrightoffset}
\setlength{\randnumberrightoffset}{0.85cm}
\newlength{\randnumberlawboxrightoffset}
\setlength{\randnumberlawboxrightoffset}{1.40cm}
\newcommand{\randnumberbox}[1]{%
  \leavevmode
  \makebox[0pt][l]{%
    \smash{%
      \tikz[remember picture,overlay,baseline=(rnmark.base)]{%
        \coordinate (rnmark) at (0,0);%
        \ifinlawbox
          \path let \p1=(current page.east), \p2=(rnmark.base) in
            node[
              anchor=base east,
              inner sep=0pt,
              outer sep=0pt,
              text width=\randnumberwidth,
              align=right,
              text=pendingink,
              font=\normalfont\mdseries\scriptsize\ttfamily
            ] at (\x1-\randnumberlawboxrightoffset,\y2) {{\normalfont\mdseries\scriptsize\ttfamily #1}};%
        \else
          \path let \p1=(current page.east), \p2=(rnmark.base) in
            node[
              anchor=base east,
              inner sep=0pt,
              outer sep=0pt,
              text width=\randnumberwidth,
              align=right,
              text=pendingink,
              font=\normalfont\mdseries\scriptsize\ttfamily
            ] at (\x1-\randnumberrightoffset,\y2) {{\normalfont\mdseries\scriptsize\ttfamily #1}};%
        \fi
      }%
    }%
  }%
  \ignorespaces
}
\newcommand{\lawboxstorerandnumber}[1]{%
  \gdef\lawboxpendingrn{#1}%
  \global\lawboxhaspendingrntrue
  \ignorespaces
}
\newcommand{\lawboxuserandnumber}{%
  \iflawboxhaspendingrn
    \randnumberbox{\lawboxpendingrn}%
    \global\lawboxhaspendingrnfalse
  \fi
}
\newcommand{\sourcern}[1]{%
  \ifrandnumbers
    \ifshoworiginal
      \ifshowtranslation\else
        \ifinlawbox\lawboxstorerandnumber{#1}\else\randnumberbox{#1}\fi
      \fi
    \fi
  \fi
}
\newcommand{\transrn}[1]{%
  \ifrandnumbers
    \ifshowtranslation
      \ifinlawbox\lawboxstorerandnumber{#1}\else\randnumberbox{#1}\fi
    \fi
  \fi
}
% ===== 課堂模式：德文原文縮寫註解 =====
% 日常切換只改各判決檔的一行：
%   \classroommodeon        開啟課堂版；註解掉這行即回到一般版。
% 模板預設：
%   1. 課堂模式關閉。
%   2. 課堂模式開啟時，縮寫註解包含「德文全稱＋中文譯名」。
%   3. 課堂模式開啟時，GG/條文編者註仍會自動顯示。
% 進階微調才需要在單案檔覆寫：
%   \classroomabbrzhoff     縮寫註解僅顯示德文全稱。
%   \showeditornotestrue    一般版也顯示編者註。
% 註解策略：同一實際 PDF 頁面中，同一縮寫只在第一次出現處加註。
% 這裡用 zref-abspage 的絕對頁序，不用 \thepage，避免文件內頁碼重置時誤判。
\newif\ifclassroommode
\newif\ifclassroomabbrzh
\classroommodefalse
\classroomabbrzhtrue
\newcommand{\classroommodeon}{\classroommodetrue}
\newcommand{\classroommodeoff}{\classroommodefalse}
\newcommand{\classroomabbrzhon}{\classroomabbrzhtrue}
\newcommand{\classroomabbrzhoff}{\classroomabbrzhfalse}
\newcommand{\abbrnotelabel}{\emph{Abk.}: }
\newcommand{\abbrnote}[3]{%
  \ifclassroommode
    \ifshoworiginal
      \ifintranslation\else
        \footnote{\normalfont\footnotesize\abbrnotelabel #2\ifclassroomabbrzh；{\songfont #3}\fi}%
      \fi
    \fi
  \fi
}
\newcommand{\abbrnoteonce}[4]{%
  #2%
  \ifclassroommode
    \ifshoworiginal
      \ifintranslation\else
        \begingroup
          \edef\abbrcurrentabspage{\number\numexpr\value{abspage}+1\relax}%
          \ifcsname abbrnoteseen@#1@\abbrcurrentabspage\endcsname\else
            \global\expandafter\let\csname abbrnoteseen@#1@\abbrcurrentabspage\endcsname\empty
            \abbrnote{#1}{#3}{#4}%
          \fi
        \endgroup
      \fi
    \fi
  \fi
}
\newcommand{\abbrAbs}{\abbrnoteonce{abs}{Abs.}{Absatz}{項}}
\newcommand{\abbrAAO}{\abbrnoteonce{aao}{a.a.O.}{am angegebenen Ort}{前引處}}
\newcommand{\abbrAE}{\abbrnoteonce{ae}{AE}{Alternativ-Entwurf}{替代草案}}
\newcommand{\abbrArt}{\abbrnoteonce{art}{Art.}{Artikel}{條}}
\newcommand{\abbrAufl}{\abbrnoteonce{aufl}{Aufl.}{Auflage}{版}}
\newcommand{\abbrBd}{\abbrnoteonce{bd}{Bd.}{Band}{卷}}
\newcommand{\abbrBGBl}{\abbrnoteonce{bgbl}{BGBl.}{Bundesgesetzblatt}{聯邦法律公報}}
\newcommand{\abbrBGHSt}{\abbrnoteonce{bghst}{BGHSt}{Entscheidungen des Bundesgerichtshofes in Strafsachen}{聯邦普通法院刑事判決彙編}}
\newcommand{\abbrBRDrucks}{\abbrnoteonce{brdrucks}{BRDrucks.}{Bundesrats-Drucksache}{聯邦參議院文件}}
\newcommand{\abbrBTDrucks}{\abbrnoteonce{btdrucks}{BTDrucks.}{Bundestags-Drucksache}{聯邦議院文件}}
\newcommand{\abbrBundesgesetzbl}{\abbrnoteonce{bundesgesetzbl}{Bundesgesetzbl.}{Bundesgesetzblatt}{聯邦法律公報}}
\newcommand{\abbrBvF}{\abbrnoteonce{bvf}{BvF}{Bundesverfassungsgerichtsverfahren}{聯邦憲法法院程序案號}}
\newcommand{\abbrBvL}{\abbrnoteonce{bvl}{BvL}{Bundesverfassungsgerichtsverfahren}{聯邦憲法法院程序案號}}
\newcommand{\abbrBvR}{\abbrnoteonce{bvr}{BvR}{Bundesverfassungsgerichtsverfahren}{聯邦憲法法院程序案號}}
\newcommand{\abbrBVerfG}{\abbrnoteonce{bverfg}{BVerfG}{Bundesverfassungsgericht}{聯邦憲法法院}}
\newcommand{\abbrBVerfGE}{\abbrnoteonce{bverfge}{BVerfGE}{Entscheidungen des Bundesverfassungsgerichts}{聯邦憲法法院判決集}}
\newcommand{\abbrBVerfGG}{\abbrnoteonce{bverfgg}{BVerfGG}{Bundesverfassungsgerichtsgesetz}{聯邦憲法法院法}}
\newcommand{\abbrDDR}{\abbrnoteonce{ddr}{DDR}{Deutsche Demokratische Republik}{德意志民主共和國；東德}}
\newcommand{\abbrDJT}{\abbrnoteonce{djt}{DJT}{Deutscher Juristentag}{德國法學家大會}}
\newcommand{\abbrDr}{\abbrnoteonce{dr}{Dr.}{Doktor}{Dr.}}
\newcommand{\abbrEuGRZ}{\abbrnoteonce{eugrz}{EuGRZ}{Europäische Grundrechte-Zeitschrift}{歐洲基本權期刊}}
\newcommand{\abbrEV}{\abbrnoteonce{ev}{EV}{Einigungsvertrag}{統一條約}}
\newcommand{\abbrF}{\abbrnoteonce{f}{f.}{folgende}{以下一頁；下一段}}
\newcommand{\abbrFF}{\abbrnoteonce{ff}{ff.}{fortfolgende}{以下數頁；以下各段}}
\newcommand{\abbrGBlDDR}{\abbrnoteonce{gblddr}{GBl. DDR}{Gesetzblatt der Deutschen Demokratischen Republik}{德意志民主共和國法律公報}}
\newcommand{\abbrGG}{\abbrnoteonce{gg}{GG}{Grundgesetz}{基本法}}
\newcommand{\abbrGVBl}{\abbrnoteonce{gvbl}{GVBl.}{Gesetz- und Verordnungsblatt}{法律及命令公報}}
\newcommand{\abbrJR}{\abbrnoteonce{jr}{JR}{Juristische Rundschau}{法學評論}}
\newcommand{\abbrJZ}{\abbrnoteonce{jz}{JZ}{JuristenZeitung}{法學家報}}
\newcommand{\abbrIVm}{\abbrnoteonce{ivm}{i.V.m.}{in Verbindung mit}{結合；連同}}
\newcommand{\abbrMdB}{\abbrnoteonce{mdb}{MdB}{Mitglied des Bundestages}{聯邦議院議員}}
\newcommand{\abbrMwN}{\abbrnoteonce{mwn}{m.w.N.}{mit weiteren Nachweisen}{附進一步文獻／判例出處}}
\newcommand{\abbrNF}{\abbrnoteonce{nf}{n.F.}{neue Fassung}{新版本}}
\newcommand{\abbrAF}{\abbrnoteonce{af}{a.F.}{alte Fassung}{舊版本}}
\newcommand{\abbrNr}{\abbrnoteonce{nr}{Nr.}{Nummer}{款、目或號，依語境}}
\newcommand{\abbrProf}{\abbrnoteonce{prof}{Prof.}{Professor}{教授}}
\newcommand{\abbrRdnr}{\abbrnoteonce{rdnr}{Rdnr.}{Randnummer}{邊碼；段碼}}
\newcommand{\abbrRdnrn}{\abbrnoteonce{rdnrn}{Rdnrn.}{Randnummern}{邊碼；段碼}}
\newcommand{\abbrRGBl}{\abbrnoteonce{rgbl}{RGBl.}{Reichsgesetzblatt}{帝國法律公報}}
\newcommand{\abbrRGSt}{\abbrnoteonce{rgst}{RGSt}{Entscheidungen des Reichsgerichts in Strafsachen}{帝國法院刑事判決彙編}}
\newcommand{\abbrRspr}{\abbrnoteonce{rspr}{Rspr.}{Rechtsprechung}{判例；司法實務}}
\newcommand{\abbrRVO}{\abbrnoteonce{rvo}{RVO}{Reichsversicherungsordnung}{帝國保險法}}
\newcommand{\abbrS}{\abbrnoteonce{s}{S.}{Seite}{頁}}
\newcommand{\abbrSFHG}{\abbrnoteonce{sfhg}{SFHG}{Schwangeren- und Familienhilfegesetz}{孕婦及家庭扶助法}}
\newcommand{\abbrSGBV}{\abbrnoteonce{sgbv}{SGB V}{Fünftes Buch Sozialgesetzbuch}{社會法典第五編}}
\newcommand{\abbrStAG}{\abbrnoteonce{staeg}{StÄG}{Strafrechtsänderungsgesetz}{刑法修正法}}
\newcommand{\abbrStenBer}{\abbrnoteonce{stenber}{StenBer.}{Stenographischer Bericht}{速記紀錄}}
\newcommand{\abbrStGB}{\abbrnoteonce{stgb}{StGB}{Strafgesetzbuch}{刑法}}
\newcommand{\abbrStREG}{\abbrnoteonce{streg}{StREG}{Strafrechtsreform-Ergänzungsgesetz}{刑法改革補充法}}
\newcommand{\abbrStrRG}{\abbrnoteonce{strrg}{StrRG}{Strafrechtsreformgesetz}{刑法改革法}}
\newcommand{\abbrUS}{\abbrnoteonce{us}{U.S.}{United States Reports}{美國最高法院判決彙編}}
\newcommand{\abbrVgl}{\abbrnoteonce{vgl}{vgl.}{vergleiche}{參見}}
\newcommand{\abbrVVDStRL}{\abbrnoteonce{vvdstrl}{VVDStRL}{Veröffentlichungen der Vereinigung der Deutschen Staatsrechtslehrer}{德國公法教師協會論文集}}
\newcommand{\abbrWP}{\abbrnoteonce{wp}{WP.}{Wahlperiode}{屆期}}
\newcommand{\abbrWp}{\abbrnoteonce{wp}{Wp.}{Wahlperiode}{屆期}}
\newcommand{\abbrZStrW}{\abbrnoteonce{zstrw}{ZStrW}{Zeitschrift für die gesamte Strafrechtswissenschaft}{整體刑法學期刊}}
\newcommand{\recordsourcerandstart}{%
  \ifrandnumbers
    \setcounter{lastsourcerandstart}{\value{randnumber}}%
    \global\lastsourcehasrandnumberstrue
  \else
    \global\lastsourcehasrandnumbersfalse
  \fi
}
\newlength{\biparagraphskip}
\setlength{\biparagraphskip}{0.52em}
\newlength{\lawboxblockbeforeskip}
\setlength{\lawboxblockbeforeskip}{0.72em}
\newlength{\lawboxblockafterskip}
\setlength{\lawboxblockafterskip}{0.92em}
\newif\iflawboxpartendedblock
\lawboxpartendedblockfalse
\newcommand{\sourceparstyle}{\footnotesize\color{sourceink}\setstretch{1.12}\RaggedRight\emergencystretch=3em\setlength{\parskip}{0.42em}\obeylines\selectfont}
\newcommand{\transparstyle}{\songfont\color{caseink}\setstretch{1.12}\setlength{\parindent}{0pt}\setlength{\parskip}{0.42em}}
\newcommand{\bilingualpairskip}{%
  \par\addvspace{\biparagraphskip}%
  \switchcolumn
  \par\addvspace{\biparagraphskip}%
  \switchcolumn*%
  \global\intranslationfalse
}
\newcommand{\bilingualboxpairskip}{%
  \par
  \switchcolumn
  \par
  \switchcolumn*%
  \global\intranslationfalse
}
\newcommand{\bilinguallawboxblockskip}{%
  \switchcolumn*%
  \par\addvspace{\lawboxblockafterskip}%
  \switchcolumn
  \par\addvspace{\lawboxblockafterskip}%
  \switchcolumn*%
  \global\intranslationfalse
}
\newcommand{\bikeepwithnext}[1][4]{%
  \ifbilingualcolumns
    \syncleft
    \Needspace{#1\baselineskip}%
    \switchcolumn
    \Needspace{#1\baselineskip}%
    \switchcolumn*%
    \global\intranslationfalse
  \else
    \Needspace{#1\baselineskip}%
  \fi
}
\NewEnviron{sourcepara}[1][]{
  \recordsourcerandstart
  \ifshoworiginal
    \ifbilingualcolumns
      \ifintranslation\switchcolumn*\global\intranslationfalse\fi
    \else
      \Needspace{5\baselineskip}
    \fi
    \ifbilingualcolumns\else\par\addvspace{0.35em}\fi
    {\sourceparstyle
    \noindent\BODY\par}
    \ifbilingualcolumns\else\addvspace{0.25em}\fi
    \ifbilingualcolumns
      \switchcolumn
      \global\intranslationtrue
    \fi
  \else
    \ifrandnumbers
      \begingroup
        \setbox0=\vbox{\hsize=\linewidth\sourceparstyle\noindent\BODY\par}%
      \endgroup
    \fi
  \fi
}
\NewEnviron{transpara}{
  \ifshowtranslation
    \setcounter{randnumberlocal}{\value{lastsourcerandstart}}%
    \ifbilingualcolumns
      \ifintranslation\else\switchcolumn\global\intranslationtrue\fi
      {\transparstyle\setlength{\leftskip}{0.55em}\setlength{\rightskip}{0.15em plus 1em}\addtolength{\linewidth}{-0.55em}\BODY\par}
      \switchcolumn*
      \bilingualpairskip
    \else
      {\transparstyle\BODY\par}
    \fi
  \else
    \ifbilingualcolumns
      \ifintranslation\switchcolumn*\global\intranslationfalse\fi
    \fi
  \fi
}
\NewEnviron{sourceboxpara}[1][]{
  \global\lawboxpartendedblockfalse
  \recordsourcerandstart
  \ifshoworiginal
    \ifbilingualcolumns
      \ifintranslation\switchcolumn*\global\intranslationfalse\fi
    \else
      \Needspace{3\baselineskip}
    \fi
    {\sourceparstyle
    \BODY\par}
    \ifbilingualcolumns\else
      \iflawboxpartendedblock\addvspace{\lawboxblockafterskip}\fi
    \fi
    \ifbilingualcolumns
      \switchcolumn
      \global\intranslationtrue
    \fi
  \fi
}
\NewEnviron{transboxpara}{
  \global\lawboxpartendedblockfalse
  \ifshowtranslation
    \setcounter{randnumberlocal}{\value{lastsourcerandstart}}%
    \ifbilingualcolumns
      \ifintranslation\else\switchcolumn\global\intranslationtrue\fi
      {\transparstyle\BODY\par}
      \iflawboxpartendedblock
        \bilinguallawboxblockskip
      \else
        \switchcolumn*
        \bilingualboxpairskip
      \fi
    \else
      {\transparstyle\BODY\par}
      \iflawboxpartendedblock\addvspace{\lawboxblockafterskip}\fi
    \fi
  \else
    \ifbilingualcolumns
      \ifintranslation\switchcolumn*\global\intranslationfalse\fi
    \fi
  \fi
}
\newcommand{\leitsatznum}[1]{\noindent\hangindent=3.0em\hangafter=1\makebox[2.25em][r]{\bfseries #1.}\hspace{0.55em}\ignorespaces}
\newcommand{\syncleft}{\ifbilingualcolumns\ifintranslation\switchcolumn*\global\intranslationfalse\fi\fi}
\pretocmd{\section}{\syncleft\clearpage}{}{}
\pretocmd{\subsection}{\syncleft}{}{}
\pretocmd{\subsubsection}{\syncleft}{}{}
\newcommand{\biheadingfont}{\songfont}
\newcommand{\bitocentry}[2]{%
  \ifshoworiginal
    \ifshowtranslation
      \ifstrequal{#1}{#2}{#1}{#1\space #2}%
    \else
      #1%
    \fi
  \else
    #2%
  \fi
}
\pdfstringdefDisableCommands{%
  \def\bitocentry#1#2{#1}%
}
\newcommand{\biheading}[7]{%
  \syncleft#1\bikeepwithnext[#3]\phantomsection\addcontentsline{toc}{#2}{\bitocentry{#6}{#7}}%
  \ifbilingualcolumns
    {#4 #6\par}%
    \switchcolumn
    {#4\biheadingfont #7\par}%
    \switchcolumn*%
    \global\intranslationfalse
    \par\addvspace{#5}%
    \switchcolumn
    \par\addvspace{#5}%
    \switchcolumn*%
  \else
    \ifshoworiginal{#4 #6\par}\fi
    \ifshowtranslation{#4\biheadingfont #7\par}\fi
    \addvspace{#5}%
  \fi
}
\newcommand{\termsectionheading}[1]{%
  \par\Needspace{9\baselineskip}%
  \vspace{0.65em}%
  {\songfont\bfseries #1\par}%
  \vspace{0.2em}%
}
% 章節換頁：
%   雙語 paracol 欄內不要用 \clearpage；它會在同步兩欄時吐出只有欄線/頁碼的空頁。
%   \newpage 足以讓下一個 section 起新頁，又不會強制清空所有待排材料。
%   單語模式不在 paracol 內，仍可用 \clearpage。
\newcommand{\bisectionpagebreak}{\ifbilingualcolumns\newpage\else\clearpage\fi}
\newcommand{\bisection}[2]{%
  \biheading{\iffirstbilingualsection\global\firstbilingualsectionfalse\else\bisectionpagebreak\fi}{section}{6}{\centering\Large\bfseries}{0.8em}{#1}{#2}%
}
\newcommand{\bisubsection}[2]{%
  \biheading{}{subsection}{5}{\large\bfseries}{0.55em}{#1}{#2}%
}
\newcommand{\bisubsubsection}[2]{%
  \biheading{}{subsubsection}{4}{\normalsize\bfseries}{0.45em}{#1}{#2}%
}
\newcommand{\bicompositeheading}[4]{%
  \biheading{}{subsection}{5}{\large\bfseries}{0.16em}{#1}{#2}%
  \biheading{}{subsubsection}{3}{\normalsize\bfseries}{0.45em}{#3}{#4}%
}
\newif\iflawboxfirstitem
\lawboxfirstitemfalse
\newcommand{\lawboxprespace}[1]{%
  \par
  \iflawboxfirstitem
    \global\lawboxfirstitemfalse
  \else
    \addvspace{#1}%
  \fi
}
\newtcolorbox{lawboxinner}{
  caseboxbase,
  colframe=sourceline!85!black,
  colback=white,
  left=1.05em,
  right=1.05em,
  top=0.62em,
  bottom=0.62em,
  before upper={\global\inlawboxtrue\global\lawboxfirstitemtrue\global\lawboxhaspendingrnfalse\ifintranslation\lawboxtransfont\else\lawboxsourcefont\fi\RaggedRight\setlength{\parindent}{0pt}\setlength{\parskip}{0pt}\ifintranslation\setstretch{\lawboxtransstretch}\else\setstretch{\lawboxsourcestretch}\fi},
  after upper={\global\lawboxhaspendingrnfalse\global\inlawboxfalse}
}
\tcbset{
  lawboxpartbase/.style={
    enhanced,
    breakable=false,
    width=0.985\linewidth,
    colback=white,
    frame hidden,
    boxrule=0pt,
    arc=0pt,
    left=1.05em,
    right=1.05em,
    boxsep=1.5mm,
    before skip=0pt,
    after skip=0pt,
    before upper={\global\inlawboxtrue\global\lawboxfirstitemtrue\global\lawboxhaspendingrnfalse\ifintranslation\lawboxtransfont\else\lawboxsourcefont\fi\RaggedRight\setlength{\parindent}{0pt}\setlength{\parskip}{0pt}\ifintranslation\setstretch{\lawboxtransstretch}\else\setstretch{\lawboxsourcestretch}\fi},
    after upper={\global\lawboxhaspendingrnfalse\global\inlawboxfalse},
    borderline west={0.28pt}{0pt}{sourceline!85!black},
    borderline east={0.28pt}{0pt}{sourceline!85!black},
  },
  lawboxpartfirst/.style={
    lawboxpartbase,
    top=0.62em,
    bottom=0.04em,
    before skip=\lawboxblockbeforeskip,
    borderline north={0.28pt}{0pt}{sourceline!85!black},
  },
  lawboxpartmiddle/.style={
    lawboxpartbase,
    top=0.04em,
    bottom=0.04em,
  },
  lawboxpartlast/.style={
    lawboxpartbase,
    top=0.04em,
    bottom=0.62em,
    after skip=0pt,
    borderline south={0.28pt}{0pt}{sourceline!85!black},
  }
}
\newtcolorbox{lawboxpartinner}[1]{#1}
\newtcolorbox{termboxinner}{
  caseboxbase,
  colframe=noteblue!70!black,
  colback=casepaper,
  before upper={\songfont\setlength{\parindent}{0pt}\setlength{\parskip}{0.35em}}
}
\newtcolorbox{deboxinner}{
  caseboxbase,
  colframe=notegray!72!black,
  colback=white,
  left=1.05em,
  right=1.05em,
  top=0.62em,
  bottom=0.62em,
  before upper={\global\inlawboxtrue\global\lawboxfirstitemtrue\global\lawboxhaspendingrnfalse\small\setlength{\parindent}{0pt}\setlength{\parskip}{0.35em}},
  after upper={\global\lawboxhaspendingrnfalse\global\inlawboxfalse}
}
\NewEnviron{lawbox}{
  \begin{lawboxinner}\BODY\end{lawboxinner}
}
\NewEnviron{lawboxpart}[2]{
  \begin{lawboxpartinner}{lawboxpart#1,equal height group=#2}\BODY\end{lawboxpartinner}%
  \ifstrequal{#1}{last}{\global\lawboxpartendedblocktrue}{}%
}
\NewEnviron{termbox}{
  \par\addvspace{0.55em}
  {\color{noteblue!70!black}\rule{\linewidth}{0.28pt}}\par
  \begingroup
    \songfont\small\color{caseink}
    \setlength{\parindent}{0pt}
    \setlength{\parskip}{0.24em}
    \BODY
    \par
  \endgroup
  {\color{noteblue!70!black}\rule{\linewidth}{0.28pt}}\par
  \addvspace{0.55em}
}
\NewEnviron{debox}{
  \begin{deboxinner}\BODY\end{deboxinner}
}
\providecommand{\paragraphmark}[1]{}
\newcommand{\lawboxtitleafter}{\addvspace{0.06em}}
\newcommand{\lawtitle}[1]{\lawboxprespace{0.22em}{\lawboxtransfont\bfseries\lawboxuserandnumber #1}\par\lawboxtitleafter}
\newcommand{\absno}[2]{\lawboxprespace{0.04em}\noindent\lawboxuserandnumber\hangindent=2.6em\hangafter=1\makebox[2.6em][l]{(#1)}#2\par}
\newcommand{\nrno}[2]{\lawboxprespace{0pt}\noindent\lawboxuserandnumber\hspace*{2.6em}\hangindent=4.4em\hangafter=1\makebox[1.8em][l]{#1.}#2\par}
\newcommand{\letterno}[2]{\lawboxprespace{0pt}\noindent\lawboxuserandnumber\hspace*{4.4em}\hangindent=6.0em\hangafter=1\makebox[1.6em][l]{#1)}#2\par}
\newcommand{\sourcelawtitle}[1]{\lawboxprespace{0.22em}{\lawboxsourcefont\bfseries\lawboxuserandnumber #1}\par\lawboxtitleafter}
\newcommand{\sourcelawsubtitle}[1]{\lawboxprespace{0.04em}{\lawboxsourcefont\bfseries\lawboxuserandnumber #1}\par\lawboxtitleafter}
\newcommand{\sourceabsno}[2]{\lawboxprespace{0.04em}\noindent\lawboxuserandnumber\hangindent=2.45em\hangafter=1\makebox[2.45em][l]{(#1)}#2\par}
\newcommand{\sourcenrno}[2]{\lawboxprespace{0pt}\noindent\lawboxuserandnumber\hspace*{2.45em}\hangindent=4.25em\hangafter=1\makebox[1.8em][l]{#1.}#2\par}
\newcommand{\sourceletterno}[2]{\lawboxprespace{0pt}\noindent\lawboxuserandnumber\hspace*{4.25em}\hangindent=5.85em\hangafter=1\makebox[1.6em][l]{#1)}#2\par}
\newcommand{\sourcequotepara}[1]{\lawboxprespace{0.04em}\noindent\lawboxuserandnumber\hangindent=1.2em\hangafter=1 #1\par}
\newcommand{\sourcelawline}[1]{\lawboxprespace{0.04em}\noindent\lawboxuserandnumber #1\par}
\newcommand{\lawline}[1]{\lawboxprespace{0.04em}\noindent\lawboxuserandnumber #1\par}
\newcommand{\pendingtranslation}{\par{\songfont\footnotesize\color{pendingink}（中文譯文待補。）}\par}
\newcommand{\tocpart}[1]{\par\addvspace{0.52em}{\footnotesize\bfseries\color{pendingink}#1\par}\addvspace{0.16em}}
\newcommand{\tocdashline}{\par\addvspace{0.48em}\noindent{\color{notegray}\xleaders\hbox to 0.82em{\hfil\rule{0.42em}{0.28pt}\hfil}\hskip\linewidth}\par\addvspace{0.38em}}
% \tocpanel{font-cmd}{content}：將目錄置中於所在欄寬（雙語欄適用）。
% A3 橫式使用 0.58\linewidth，A4 橫式使用 0.84\linewidth，
% 使兩種紙張下目錄欄內容實際寬度相近（均約 100–105 mm）。
\newcommand{\tocpanel}[2]{%
  \makebox[\linewidth][c]{%
    \ifx\bilingualpaper\bilingualpaperaThree
      \begin{minipage}{0.58\linewidth}%
    \else
      \begin{minipage}{0.84\linewidth}%
    \fi
    \toccolstyle
    #1#2%
    \end{minipage}%
  }%
}
% ===== 首頁簡目錄的兩層 description list 縮排 =====
% enumitem 的 description list 可把每一列拆成：
%   [label]  labelsep  body text
% 這裡的「起點」都以所在 minipage/欄位的左緣為 0。
%
% 核心規則：
%   labelindent = label 欄位從左緣開始的位置。
%   labelwidth  = label 欄位本身寬度；標籤太長（如 Abw. Meinung）就加大它。
%   labelsep    = label 欄位右緣到正文的距離。
%   leftmargin  = 正文 body text 的起點。判斷層級視覺時主要看這個。
%
% 所以：
%   想讓「Begründung / 判決理由」整列正文往右或往左：改 tocitems 的 leftmargin。
%   想讓「A. Verfahren und Gegenstand / A. 程序與審查標的」正文往右或往左：
%     改 tocsubitems 的 leftmargin。
%   想讓 A./B./C. 這些標籤本身往右或往左，但正文不動：改 tocsubitems 的 labelindent。
%   想讓 A./B./C. 與正文間距變大或變小：改 tocsubitems 的 labelsep。
%
% 層級原則：
%   tocsubitems.leftmargin 必須大於 tocitems.leftmargin，
%   否則子層級正文會看起來比「Gründe / 理由」還靠左。
\newlist{tocitems}{description}{1}
\setlist[tocitems]{%
  labelindent=0pt,    % 第一層 label 欄位起點；0pt 表示貼齊目錄欄左緣。
  leftmargin=2.54cm,  % 第一層正文起點；例如 Begründung / 判決理由 的左緣。
  labelwidth=2.16cm,  % 第一層 label 欄位寬度；需容納 Leitsätze、Abw. Meinung、不同意見等。
  labelsep=0.32cm,    % 第一層 label 與正文的水平間距；只改間距，不改正文起點。
  itemsep=0.28em,     % 第一層各項之間的垂直距離。
  topsep=0.20em,      % 第一層 list 上下與前後內容的垂直距離。
  parsep=0pt,         % 同一 item 內段落之間的距離；目錄項通常不需要。
  partopsep=0pt,      % list 若緊接段落開始時的額外距離；關閉以免目錄高度飄動。
  align=left,         % label 欄內左對齊；長短標籤不靠右貼齊。
  font=\normalfont\bfseries % label 的字型；正文不受這行影響。
}
\newlist{tocsubitems}{description}{1}
\setlist[tocsubitems]{%
  labelindent=1cm,    % 第二層 label 起點；只控制 A./B./C. 本身的位置。
  leftmargin=3.00cm,  % 第二層正文起點；必須大於 2.54cm，才會比 Begründung / 判決理由更右。
  labelwidth=0.5cm,   % 第二層 label 欄寬；A./B./C. 很短，可小於第一層。
  labelsep=1.5cm,    % 第二層 label 與正文間距；A. 與 Verfahren / 程序 的距離看這裡。
  itemsep=0.12em,     % 第二層各項較密，避免 A.–E. 佔太高。
  topsep=0.04em,      % 第二層 list 與 Gründe / 理由之間的垂直距離。
  parsep=0pt,         % 同一第二層 item 內段落距離；目錄項通常不需要。
  partopsep=0pt,      % 關閉額外段落起始距離，避免版面不穩。
  align=left,         % A./B./C. label 欄內左對齊。
  font=\normalfont\bfseries, % A./B./C. label 加粗；正文不受這行影響。
  before={}           % 不在第二層 list 前自動插入額外內容。
}
\newlist{termitems}{description}{1}
\ifbilingualcolumns
  \setlist[termitems]{%
    leftmargin=5.2cm,
    labelwidth=4.8cm,
    labelsep=0.35cm,
    itemsep=0.16em,
    topsep=0.2em,
    parsep=0pt,
    partopsep=0pt,
    font=\normalfont\bfseries,
    style=nextline
  }
\else
  \setlist[termitems]{%
    leftmargin=0pt,
    labelwidth=0pt,
    labelsep=0pt,
    itemsep=0.24em,
    topsep=0.2em,
    parsep=0pt,
    partopsep=0pt,
    font=\normalfont\bfseries,
    style=nextline
  }
\fi
\newlist{abbritems}{description}{1}
\setlist[abbritems]{%
  leftmargin=2.38cm,
  labelwidth=2.02cm,
  labelsep=0.28cm,
  itemsep=0.16em,
  topsep=0.2em,
  parsep=0pt,
  partopsep=0pt,
  align=left,
  font=\normalfont\bfseries
}
\ifbilingualcolumns
  \newenvironment{termcolumns}{\begin{multicols}{2}}{\end{multicols}}
\else
  \newenvironment{termcolumns}{}{}
\fi

% ===== 編者註腳開關 =====
% 一般版預設不顯示編者註，避免正文腳註過密。
% 若開啟 \classroommodeon，GG/條文編者註仍會自動顯示，無需另開本開關。
% 若一般版也要顯示編者註，可在單案檔 \input 後加 \showeditornotestrue。
\newif\ifshoweditornotes
\showeditornotesfalse

% ===== 目錄排版輔助 =====
\newcommand{\toccolstyle}{\raggedright\arraybackslash}
\newcommand{\tocsingleindent}{\leftskip=1.45cm\rightskip=1.2cm}

% ===== 行內引文環境（德文原文與中文譯文各一，帶左側灰色邊線）=====
\newcommand{\quoteparbreak}{\par\addvspace{0.48em}}
\newcommand{\sourcequoteinner}[1]{%
  \begin{tcolorbox}[
    enhanced, breakable, width=\dimexpr\linewidth-0.8em\relax, frame hidden,
    colback=white, coltext=caseink, boxrule=0pt,
    borderline west={0.55pt}{0.88em}{notegray},
    left=1.78em, right=0.72em, top=0.18em, bottom=0.18em,
    boxsep=0pt, before skip=0.24em, after skip=0.24em,
    before upper={\normalfont\footnotesize\color{caseink}\setlength{\parindent}{0pt}\setlength{\parskip}{0.34em}\raggedright\setlength{\rightskip}{0pt plus 1em}}
  ]%
  #1\par
  \end{tcolorbox}%
}
\newcommand{\transquoteinner}[1]{%
  \begin{tcolorbox}[
    enhanced, breakable, width=\dimexpr\linewidth-0.8em\relax, frame hidden,
    colback=white, coltext=caseink, boxrule=0pt,
    borderline west={0.55pt}{0.88em}{notegray},
    left=1.78em, right=0.72em, top=0.18em, bottom=0.18em,
    boxsep=0pt, before skip=0.24em, after skip=0.24em,
    before upper={\songfont\small\color{caseink}\setlength{\parindent}{0pt}\setlength{\parskip}{0.34em}\raggedright\setlength{\rightskip}{0pt plus 1em}}
  ]%
  #1\par
  \end{tcolorbox}%
}

% ===== 大綱（Gliederung）Rn 輔助命令 =====
\newcommand{\outrn}[2]{\enspace{\normalfont\footnotesize\color{pendingink}(Rn\,#1–#2)}}
\newcommand{\outrnsingle}[1]{\enspace{\normalfont\footnotesize\color{pendingink}(Rn\,#1)}}
\newcommand{\outrnzh}[2]{\enspace{\normalfont\footnotesize\color{pendingink}（第\,#1–#2\,段）}}
\newcommand{\outrnzhs}[1]{\enspace{\normalfont\footnotesize\color{pendingink}（第\,#1\,段）}}
% ===== 大綱雙語對照表格輔助命令（三欄：段碼 │ 德文 │ 中文）=====
% 欄 1：段碼（由欄定義自動格式化：小字、等寬、右對齊、pendingink）
% 欄 2：德文標籤＋正文
% 欄 3：中文標籤＋正文
%
% 無段碼的列（\omainrow、\osecrow），欄 1 以 & 空置。
% 有段碼的列（\osecrowrn、\osubrow、\ointrow），#1 為預格式化字串
%   例：\osubrow{2–49}{I.}{...}{I.}{...}
%       \ointrow{107}{...}{...}
%
% \opartrow：段組標題（橫跨三欄）
\newcommand{\opartrow}[2]{%
  \noalign{\vspace{0.30em}}%
  \multicolumn{3}{@{}l@{}}{%
    \footnotesize\bfseries\color{pendingink}#1\hfill\songfont#2%
  }\\[0.06em]%
}
% \odashrow：虛線分隔（橫跨三欄）
\newcommand{\odashrow}{%
  \noalign{\vspace{0.28em}}%
  \multicolumn{3}{@{}l@{}}{\color{notegray}%
    \xleaders\hbox to 0.82em{\hfil\rule{0.42em}{0.28pt}\hfil}\hskip 0.88\linewidth}%
  \\[0.24em]%
}
% \odissentpart：不同意見書區段標頭。比一般 \opartrow 更像附錄性轉場。
\newcommand{\odissentpart}[2]{%
  \noalign{\vspace{0.42em}}%
  \multicolumn{3}{@{}p{0.88\textwidth}@{}}{%
    \color{notegray}\rule{\linewidth}{0.18pt}%
  }\\[-0.18em]%
  \multicolumn{3}{@{}p{0.88\textwidth}@{}}{%
    \makebox[\linewidth][s]{%
      \footnotesize\bfseries\color{pendingink}#1%
      \hfill{\songfont #2}%
    }%
  }\\[-0.02em]%
}
% \odissentopinion：不同意見書的署名與一行摘要，右欄保留段碼範圍。
\newcommand{\odissentopinion}[5]{%
  \footnotesize\hangindent=0.52cm\hangafter=1%
    {\bfseries\color{caseink}#2}\enspace{\color{notegray}\textperiodcentered}\enspace\color{caseink}#3%
  &\footnotesize\hangindent=0.52cm\hangafter=1%
    {\songfont\bfseries\color{caseink}#4}\enspace{\color{notegray}\textperiodcentered}\enspace\songfont\color{caseink}#5%
  &#1%
  \\[0.12em]%
}
% \omainrow：頂層項目，無段碼（Leitsätze、Urteil、Gründe …）
\newcommand{\omainrow}[4]{%
  \footnotesize\hangindent=1.92cm\hangafter=1%
    \makebox[1.92cm][l]{\bfseries\color{caseink}#1}\color{caseink}#2%
  &\footnotesize\hangindent=1.92cm\hangafter=1%
    \makebox[1.92cm][l]{\bfseries\color{caseink}#3}\songfont\color{caseink}#4%
  &%
  \\[0.15em]%
}
% \omainrowrn：頂層項目，有段碼（不同意見等標籤寬於 \osecrow 框者）
\newcommand{\omainrowrn}[5]{%
  \footnotesize\hangindent=1.92cm\hangafter=1%
    \makebox[1.92cm][l]{\bfseries\color{caseink}#2}\color{caseink}#3%
  &\footnotesize\hangindent=1.92cm\hangafter=1%
    \makebox[1.92cm][l]{\bfseries\color{caseink}#4}\songfont\color{caseink}#5%
  &#1%
  \\[0.15em]%
}
% \osecrow：一級子項，無段碼（A.、C.、D. 等有子項者）
\newcommand{\osecrow}[4]{%
  \footnotesize\hspace{1.10cm}\hangindent=1.98cm\hangafter=1%
    \makebox[0.86cm][l]{\bfseries\color{caseink}#1}\color{caseink}#2%
  &\footnotesize\hspace{1.10cm}\hangindent=1.98cm\hangafter=1%
    \makebox[0.86cm][l]{\bfseries\color{caseink}#3}\songfont\color{caseink}#4%
  &%
  \\[0.11em]%
}
% \osecrowrn：一級子項，有段碼（B.、E. 等完整段落範圍者）
% #1=段碼字串（如 101–106）, #2=DE標籤, #3=DE正文, #4=ZH標籤, #5=ZH正文
\newcommand{\osecrowrn}[5]{%
  \footnotesize\hspace{1.10cm}\hangindent=1.98cm\hangafter=1%
    \makebox[0.86cm][l]{\bfseries\color{caseink}#2}\color{caseink}#3%
  &\footnotesize\hspace{1.10cm}\hangindent=1.98cm\hangafter=1%
    \makebox[0.86cm][l]{\bfseries\color{caseink}#4}\songfont\color{caseink}#5%
  &#1%
  \\[0.11em]%
}
% \osubrow：二級子項，有段碼範圍（I.、II.、A.I. …）
% #1=段碼字串, #2=DE標籤, #3=DE正文, #4=ZH標籤, #5=ZH正文
\newcommand{\osubrow}[5]{%
  \footnotesize\hspace{1.40cm}\hangindent=2.30cm\hangafter=1%
    \makebox[0.90cm][l]{\bfseries\color{caseink}#2}\color{caseink}#3%
  &\footnotesize\hspace{1.40cm}\hangindent=2.30cm\hangafter=1%
    \makebox[0.90cm][l]{\bfseries\color{caseink}#4}\songfont\color{caseink}#5%
  &#1%
  \\[0.09em]%
}
% \ointrow：引入段落，有單一段碼，無標籤
% #1=段碼字串, #2=DE正文, #3=ZH正文
\newcommand{\ointrow}[3]{%
  \footnotesize\hspace{0.52cm}\hangindent=0.52cm\hangafter=1\color{caseink}#2%
  &\footnotesize\hspace{0.52cm}\hangindent=0.52cm\hangafter=1\songfont\color{caseink}#3%
  &#1%
  \\[0.09em]%
}
% \osubsubrow：三級子項（1.、2.、3.）
% #1=段碼字串, #2=DE標籤, #3=DE正文, #4=ZH標籤, #5=ZH正文
\newcommand{\osubsubrow}[5]{%
  \footnotesize\hspace{1.80cm}\hangindent=2.48cm\hangafter=1%
    \makebox[0.68cm][l]{\bfseries\color{caseink}#2}\color{caseink}#3%
  &\footnotesize\hspace{1.80cm}\hangindent=2.48cm\hangafter=1%
    \makebox[0.68cm][l]{\bfseries\color{caseink}#4}\songfont\color{caseink}#5%
  &#1%
  \\[0.08em]%
}
% \osubsubsubrow：四級子項（a）、b））
% #1=段碼字串, #2=DE標籤, #3=DE正文, #4=ZH標籤, #5=ZH正文
\newcommand{\osubsubsubrow}[5]{%
  \footnotesize\hspace{2.10cm}\hangindent=2.68cm\hangafter=1%
    \makebox[0.58cm][l]{\bfseries\color{caseink}#2}\color{caseink}#3%
  &\footnotesize\hspace{2.10cm}\hangindent=2.68cm\hangafter=1%
    \makebox[0.58cm][l]{\bfseries\color{caseink}#4}\songfont\color{caseink}#5%
  &#1%
  \\[0.07em]%
}
% \osubsubsubsubrow：五級子項（aa）、bb））
% 標籤框 0.62cm：足以容納「aa)」在 footnotesize bold 下（約 0.50cm）並留有間距
% #1=段碼字串, #2=DE標籤, #3=DE正文, #4=ZH標籤, #5=ZH正文
\newcommand{\osubsubsubsubrow}[5]{%
  \footnotesize\hspace{2.38cm}\hangindent=3.06cm\hangafter=1%
    \makebox[0.68cm][l]{\bfseries\color{caseink}#2}\color{caseink}#3%
  &\footnotesize\hspace{2.38cm}\hangindent=3.06cm\hangafter=1%
    \makebox[0.68cm][l]{\bfseries\color{caseink}#4}\songfont\color{caseink}#5%
  &#1%
  \\[0.06em]%
}
% \outlinepage：雙語對照大綱頁包裝環境（longtable 自動跨頁）
% 三欄：德文欄 + 中文欄 + 段碼欄（最右，不折行）
% 段碼欄寬度：1.80cm，足以容納最長段碼如「309–348」
% DE/ZH 欄寬：(0.87\textwidth - 2.08cm) / 2
%   驗算：2×欄 + 0.05\textwidth + 0.28cm gap + 1.80cm = 0.87tw - 2.08cm + 0.05tw + 2.08cm = 0.92tw ✓
\newenvironment{outlinepage}{%
  \vspace*{0.40cm}%
  \setlength{\LTleft}{0.04\textwidth}%
  \setlength{\LTright}{0.04\textwidth}%
  \setlength{\tabcolsep}{0pt}%
  % 大綱列距由 longtable 的 \arraystretch 與各 row macro 結尾的 \\[...em] 共同決定。
  % 這裡控制「每一列本身的表格 strut 高度」；數值越大，A./I./1. 各項之間越像有空行。
  % A3 原本用 0.62，視覺上沒有空行；A4 也沿用同值，避免切到 A4 後項目被撐開。
  % 若日後覺得大綱太擠，只微調這個值（例如 0.68），不要改各 \omainrow/\osubrow 的縮排。
  \renewcommand{\arraystretch}{0.62}%
  \begin{longtable}{%
    >{\vspace{0pt}\raggedright\arraybackslash}p{\dimexpr(0.87\textwidth-2.08cm)/2\relax}%
    @{\hspace{0.025\textwidth}}%
    !{\color{notegray}\vrule width 0.12pt}%
    @{\hspace{0.025\textwidth}}%
    >{\vspace{0pt}\raggedright\arraybackslash}p{\dimexpr(0.87\textwidth-2.08cm)/2\relax}%
    @{\hspace{0.28cm}}%
    >{\vspace{0pt}\footnotesize\normalfont\color{pendingink}\raggedright\arraybackslash}p{1.80cm}%
    @{}%
  }%
  % 首頁欄標籤：不要橫跨置中；分別放在德文欄與中文欄的實際欄位起點。
  \footnotesize\bfseries\color{sourceink}Gliederung%
  &\footnotesize\bfseries\songfont\color{sourceink}大綱%
  &%
  \\[0.36em]%
  \endfirsthead%
  % 續頁欄標籤：同樣分欄定位，以灰色細體區分；無需標注「續」。
  \footnotesize\color{pendingink}Gliederung%
  &\footnotesize\songfont\color{pendingink}大綱%
  &%
  \\[0.24em]%
  \endhead%
}{%
  \end{longtable}%
}
 載入。
% 不含：\documentclass、\documentversion、\ggversionde/zh、\tocde/\toczh。
% 日期與首頁版本列由本檔自動產生；各文件只需手動指定 \documentversion。

\usepackage{xeCJK}
\usepackage{fontspec}
\usepackage{geometry}
\usepackage{setspace}
\usepackage{hyperref}
\usepackage{xurl}
\usepackage{bookmark}
\usepackage{enumitem}
\usepackage{xcolor}
\usepackage{titlesec}
\usepackage{array}
\usepackage{longtable}
\usepackage{tcolorbox}
\usepackage{environ}
\usepackage{pdflscape}
\usepackage{etoolbox}
\usepackage{ragged2e}
\usepackage{paracol}
\usepackage{multicol}
\usepackage{needspace}
\usepackage{fancyhdr}
\usepackage{tikz}
\usepackage{zref-abspage}
\tcbuselibrary{breakable,skins}
\usetikzlibrary{calc}
\usepackage{pgfcalendar}

\hypersetup{
  unicode=true,
  bookmarksopen=true,
  bookmarksnumbered=false,
  hidelinks
}

% ===== 自動推導，不要手動改下面 =====
% （\docmode 與 \bilingualpaper 由各文件在 % bverfge_layout.tex
% 共用排版基礎設施，供 Schwangerschaftsabbruch I 與 II 共享。
% 由各文件以 \documentclass 後 % bverfge_layout.tex
% 共用排版基礎設施，供 Schwangerschaftsabbruch I 與 II 共享。
% 由各文件以 \documentclass 後 \input{../../共用LaTeX/bverfge_layout} 載入。
% 不含：\documentclass、\documentversion、\ggversionde/zh、\tocde/\toczh。
% 日期與首頁版本列由本檔自動產生；各文件只需手動指定 \documentversion。

\usepackage{xeCJK}
\usepackage{fontspec}
\usepackage{geometry}
\usepackage{setspace}
\usepackage{hyperref}
\usepackage{xurl}
\usepackage{bookmark}
\usepackage{enumitem}
\usepackage{xcolor}
\usepackage{titlesec}
\usepackage{array}
\usepackage{longtable}
\usepackage{tcolorbox}
\usepackage{environ}
\usepackage{pdflscape}
\usepackage{etoolbox}
\usepackage{ragged2e}
\usepackage{paracol}
\usepackage{multicol}
\usepackage{needspace}
\usepackage{fancyhdr}
\usepackage{tikz}
\usepackage{zref-abspage}
\tcbuselibrary{breakable,skins}
\usetikzlibrary{calc}
\usepackage{pgfcalendar}

\hypersetup{
  unicode=true,
  bookmarksopen=true,
  bookmarksnumbered=false,
  hidelinks
}

% ===== 自動推導，不要手動改下面 =====
% （\docmode 與 \bilingualpaper 由各文件在 \input{bverfge_layout} 之前設定；
%   預設 chinese / a4）
\ifdefined\docmode\else\def\docmode{chinese}\fi
\ifdefined\bilingualpaper\else\def\bilingualpaper{a4}\fi
\newif\ifshoworiginal
\newif\ifshowtranslation
\newif\ifbilinguallandscape
\newif\ifbilingualcolumns
\newif\ifintranslation
\newif\iffirstbilingualsection

\showoriginalfalse
\showtranslationfalse
\bilinguallandscapefalse
\bilingualcolumnsfalse
\intranslationfalse
\firstbilingualsectionfalse

\def\modechinese{chinese}
\def\modegerman{german}
\def\modebilingual{bilingual}
\def\bilingualpaperaThree{a3}
\def\bilingualpaperaFour{a4}

\ifx\docmode\modechinese
  \showtranslationtrue
\fi

\ifx\docmode\modegerman
  \showoriginaltrue
\fi

\ifx\docmode\modebilingual
  \showoriginaltrue
  \showtranslationtrue
  \bilinguallandscapetrue
  \bilingualcolumnstrue
\fi

% 編排模式說明：
% chinese  → \showoriginalfalse  \showtranslationtrue  （A4 直式，純中文）
% german   → \showoriginaltrue   \showtranslationfalse （A4 直式，純德文）
% bilingual→ \showoriginaltrue   \showtranslationtrue  \bilingualcolumnstrue（A3/A4 橫式對照）

\ifbilingualcolumns
  \ifx\bilingualpaper\bilingualpaperaFour
    \geometry{
      a4paper,
      landscape,
      left=1.25cm,
      right=2.25cm,
      top=1.25cm,
      bottom=1.75cm,
      footskip=8.5mm,
      marginparwidth=0.75cm,
      marginparsep=0.25cm
    }
  \else
    \geometry{
      a3paper,
      landscape,
      left=1.55cm,
      right=2.75cm,
      top=1.55cm,
      bottom=2.05cm,
      footskip=9.5mm,
      marginparwidth=0.9cm,
      marginparsep=0.35cm
    }
  \fi
\else
\geometry{
  a4paper,
  portrait,
  left=2.2cm,
  right=2.6cm,
  top=2.0cm,
  bottom=2.0cm,
  marginparwidth=0.8cm,
  marginparsep=0.3cm
}
\fi
% ===== 時間戳基礎設施 =====
\newcount\stamptimeval
\newcount\stampminuteval
\newcommand{\stamphh}{%
  \stamptimeval=\time
  \divide\stamptimeval by 60
  \ifnum\stamptimeval<10 0\fi
  \the\stamptimeval}
\newcommand{\stampmm}{%
  \stamptimeval=\time
  \stampminuteval=\time
  \divide\stampminuteval by 60
  \multiply\stampminuteval by 60
  \advance\stamptimeval by -\stampminuteval
  \ifnum\stamptimeval<10 0\fi
  \the\stamptimeval}
\newcommand{\stampwkde}[1]{%
  \ifcase#1Montag\or Dienstag\or Mittwoch\or Donnerstag\or Freitag\or Samstag\or Sonntag\fi}
\newcommand{\stampwkzh}[1]{%
  \ifcase#1 一\or 二\or 三\or 四\or 五\or 六\or 日\fi}
\newcommand{\stampmonthde}[1]{%
  \ifcase#1\or Januar\or Februar\or März\or April\or Mai\or Juni\or
  Juli\or August\or September\or Oktober\or November\or Dezember\fi}
\newcount\stampjdval
\newcount\stampwdval
\pgfcalendardatetojulian{\the\year-\the\month-\the\day}{\stampjdval}
\pgfcalendarjuliantoweekday{\the\stampjdval}{\stampwdval}
\newcommand{\timestampzh}{%
  \songfont 最後修改：\the\year 年\the\month 月\the\day 日%
  % （星期\stampwkzh{\the\stampwdval}）%
  % \space\stamphh:\stampmm
  }

\newcommand{\documentdatezh}{\the\year 年 \the\month 月 \the\day 日}
\newcommand{\documentdatede}{\the\day. \stampmonthde{\the\month} \the\year}
\newcommand{\bverfgetranslationnote}{%
  {\songfont 翻譯說明：初譯使用 Claude Code 與 GPT 協助迭代；仍請以德文原文為準。}%
}
\newcommand{\bverfgecourseinfo}{%
  {\songfont 課程：114 學年度第 2 學期；憲法專題研究（四）；授課老師：湯德宗教授；導讀日期：115 年 6 月 1 日。}%
}
\newcommand{\bverfgeeditorinfo}{%
  {\songfont 編者：王逸帆（R10A21126），國立台灣大學法律系研究所碩士班。}%
}
\newcommand{\bverfgetitleversionline}{%
  \ifbilingualcolumns
    Fassung \documentversion\enspace\textperiodcentered\enspace Stand: \documentdatede
    \quad
    {\songfont 版本 \documentversion\enspace\textperiodcentered\enspace 修訂日期：\documentdatezh\par}%
  \else
    \ifshowtranslation
      {\songfont 版本 \documentversion\enspace\textperiodcentered\enspace 修訂日期：\documentdatezh\par}%
    \else
      Fassung \documentversion\enspace\textperiodcentered\enspace Stand: \documentdatede\par
    \fi
  \fi
}
\newcommand{\bverfgetitlecornerinfo}{%
  \ifshowtranslation
    \vspace{0.72em}%
    \noindent\hfill
    \begin{minipage}{0.66\textwidth}
      \raggedleft\tiny\color{pendingink}\songfont
      \bverfgecourseinfo\par
      \bverfgeeditorinfo
    \end{minipage}%
  \fi
}
\newcommand{\bverfgetitlemeta}{%
  \vfill
  \begin{center}%
    \footnotesize\color{pendingink}%
    \bverfgetitleversionline
    \ifshowtranslation
      \vspace{0.28em}{\scriptsize\bverfgetranslationnote\par}%
    \fi
  \end{center}%
  \bverfgetitlecornerinfo
}

\pagestyle{fancy}
\fancyhf{}
\renewcommand{\headrulewidth}{0pt}
\renewcommand{\footrulewidth}{0pt}
\fancyfoot[C]{\thepage}
\ifbilingualcolumns
  \fancyfoot[R]{\scriptsize\color{pendingink}\timestampzh}%
\else
  \ifshoworiginal
  \else
    \fancyfoot[R]{\scriptsize\color{pendingink}\timestampzh\hspace{0.45cm}}%
  \fi
\fi
\setmainfont{Times New Roman}
\setsansfont{Helvetica Neue}
\setmonofont{Menlo}
\IfFileExists{/Users/iw/Library/Fonts/kaiu.ttf}{
  \setCJKmainfont[Path=/Users/iw/Library/Fonts/,AutoFakeBold=4]{kaiu.ttf}
}{
  \IfFontExistsTF{標楷體}{
    \setCJKmainfont[AutoFakeBold=4]{標楷體}
  }{
    \setCJKmainfont{Songti TC}
  }
}
\IfFontExistsTF{Songti TC}{
  \setCJKfamilyfont{song}{Songti TC}
  \setCJKmonofont{Songti TC}
}{
  \setCJKfamilyfont{song}[Path=/Users/iw/Library/Fonts/,AutoFakeBold=4]{kaiu.ttf}
  \setCJKmonofont[Path=/Users/iw/Library/Fonts/,AutoFakeBold=4]{kaiu.ttf}
}
\newcommand{\songfont}{\CJKfamily{song}}

% ===== 封面／目錄專用打字機字體（僅限封面、目錄頁使用，不影響正文）=====
\IfFontExistsTF{Radio Newsman}{%
  \newfontfamily\bverfgetitlede{Radio Newsman}%
}{%
  \newfontfamily\bverfgetitlede{Courier New}%
}
\IfFontExistsTF{Erikas Farbband}{%
  \newfontfamily\bverfgebodyde{Erikas Farbband}%
}{%
  \let\bverfgebodyde\bverfgetitlede
}
\IfFontExistsTF{Maoken Heavy Labourer}{%
  \setCJKfamilyfont{bverfgetitlecjk}{Maoken Heavy Labourer}%
}{%
  \setCJKfamilyfont{bverfgetitlecjk}{Songti TC}%
}
\IfFontExistsTF{Huiwen-mincho}{%
  \setCJKfamilyfont{bverfgebodycjk}{Huiwen-mincho}%
}{%
  \setCJKfamilyfont{bverfgebodycjk}{Songti TC}%
}
\newcommand{\bverfgetitlezh}{\bverfgetitlede\CJKfamily{bverfgetitlecjk}}
\newcommand{\bverfgebodyzh}{\bverfgebodyde\CJKfamily{bverfgebodycjk}}

\setstretch{1.35}
\setlist{itemsep=0.2em, topsep=0.4em}
\setlength{\parindent}{0pt}
\setlength{\parskip}{0.45em}
\raggedbottom
\setcounter{secnumdepth}{0}
\setcounter{tocdepth}{2}

\newcommand{\lawboxsourcefont}{\footnotesize}
\newcommand{\lawboxtransfont}{\footnotesize}
\newcommand{\lawboxsourcestretch}{1.02}
\newcommand{\lawboxtransstretch}{0.92}

\titleformat{\section}{\centering\Large\bfseries\songfont}{\thesection}{1em}{}
\titleformat{\subsection}{\large\bfseries\songfont}{\thesubsection}{1em}{}
\titleformat{\subsubsection}{\normalsize\bfseries\songfont}{\thesubsubsection}{1em}{}
\titleformat{\paragraph}{\normalsize\bfseries\songfont}{\theparagraph}{1em}{}
\titlespacing*{\section}{0pt}{1.8em}{1.0em}
\titlespacing*{\subsection}{0pt}{1.2em}{0.6em}
\titlespacing*{\subsubsection}{0pt}{1.0em}{0.45em}
\titlespacing*{\paragraph}{0pt}{0.6em}{0.4em}

\definecolor{noteblue}{HTML}{A9C9F5}
\definecolor{noteteal}{HTML}{AEE1D6}
\definecolor{notegray}{HTML}{D7DADF}
\definecolor{caseink}{HTML}{000000}
\definecolor{casepaper}{HTML}{FCFEFD}
\definecolor{sourcepaper}{HTML}{FAFBF8}
\definecolor{sourceline}{HTML}{D7DED8}
\definecolor{sourceink}{HTML}{000000}
\definecolor{pendingink}{HTML}{000000}

\tcbset{
  caseboxbase/.style={
    enhanced,
    breakable,
    width=0.985\linewidth,
    colback=casepaper,
    coltitle=caseink,
    fonttitle=\bfseries\songfont,
    boxrule=0.28pt,
    arc=1.2mm,
    left=2.4mm,
    right=2.4mm,
    top=1.5mm,
    bottom=1.5mm,
    boxsep=1.5mm,
    before skip=0.65em,
    after skip=0.65em,
    before upper={\setlength{\parindent}{0pt}\setlength{\parskip}{0.35em}},
  }
}

\newcommand{\bilingualstart}{}
\newcommand{\bilingualstop}{}
\ifbilingualcolumns
  \renewcommand{\bilingualstart}{\small\setlength{\columnsep}{1.25cm}\setlength{\columnseprule}{0.12pt}\columnratio{0.5,0.5}\multicolumnfootnotes\flushbottom\sloppy\interlinepenalty=80\clubpenalty=800\widowpenalty=800\displaywidowpenalty=800\begin{paracol}{2}\global\intranslationfalse\global\firstbilingualsectiontrue{\footnotesize\color{sourceink}Deutsch\par}\switchcolumn{\footnotesize\songfont\color{sourceink}中文譯文\par}\switchcolumn*}
  \renewcommand{\bilingualstop}{\ifintranslation\switchcolumn*\global\intranslationfalse\fi\end{paracol}\clearpage\raggedbottom}
\fi
\newif\ifrandnumbers
\randnumbersfalse
\newcounter{randnumber}
\newcounter{randnumberlocal}
\newcounter{lastsourcerandstart}
\newif\iflastsourcehasrandnumbers
\lastsourcehasrandnumbersfalse
\newif\ifinlawbox
\inlawboxfalse
\newif\iflawboxhaspendingrn
\lawboxhaspendingrnfalse
\newcommand{\lawboxpendingrn}{}
\newcommand{\startrandnumbers}{\global\randnumberstrue\setcounter{randnumber}{1}\global\lastsourcehasrandnumbersfalse}
\newcommand{\stoprandnumbers}{\global\randnumbersfalse\global\lastsourcehasrandnumbersfalse}
\newlength{\randnumberwidth}
\setlength{\randnumberwidth}{0.95cm}
\newlength{\randnumberrightoffset}
\setlength{\randnumberrightoffset}{0.85cm}
\newlength{\randnumberlawboxrightoffset}
\setlength{\randnumberlawboxrightoffset}{1.40cm}
\newcommand{\randnumberbox}[1]{%
  \leavevmode
  \makebox[0pt][l]{%
    \smash{%
      \tikz[remember picture,overlay,baseline=(rnmark.base)]{%
        \coordinate (rnmark) at (0,0);%
        \ifinlawbox
          \path let \p1=(current page.east), \p2=(rnmark.base) in
            node[
              anchor=base east,
              inner sep=0pt,
              outer sep=0pt,
              text width=\randnumberwidth,
              align=right,
              text=pendingink,
              font=\normalfont\mdseries\scriptsize\ttfamily
            ] at (\x1-\randnumberlawboxrightoffset,\y2) {{\normalfont\mdseries\scriptsize\ttfamily #1}};%
        \else
          \path let \p1=(current page.east), \p2=(rnmark.base) in
            node[
              anchor=base east,
              inner sep=0pt,
              outer sep=0pt,
              text width=\randnumberwidth,
              align=right,
              text=pendingink,
              font=\normalfont\mdseries\scriptsize\ttfamily
            ] at (\x1-\randnumberrightoffset,\y2) {{\normalfont\mdseries\scriptsize\ttfamily #1}};%
        \fi
      }%
    }%
  }%
  \ignorespaces
}
\newcommand{\lawboxstorerandnumber}[1]{%
  \gdef\lawboxpendingrn{#1}%
  \global\lawboxhaspendingrntrue
  \ignorespaces
}
\newcommand{\lawboxuserandnumber}{%
  \iflawboxhaspendingrn
    \randnumberbox{\lawboxpendingrn}%
    \global\lawboxhaspendingrnfalse
  \fi
}
\newcommand{\sourcern}[1]{%
  \ifrandnumbers
    \ifshoworiginal
      \ifshowtranslation\else
        \ifinlawbox\lawboxstorerandnumber{#1}\else\randnumberbox{#1}\fi
      \fi
    \fi
  \fi
}
\newcommand{\transrn}[1]{%
  \ifrandnumbers
    \ifshowtranslation
      \ifinlawbox\lawboxstorerandnumber{#1}\else\randnumberbox{#1}\fi
    \fi
  \fi
}
% ===== 課堂模式：德文原文縮寫註解 =====
% 日常切換只改各判決檔的一行：
%   \classroommodeon        開啟課堂版；註解掉這行即回到一般版。
% 模板預設：
%   1. 課堂模式關閉。
%   2. 課堂模式開啟時，縮寫註解包含「德文全稱＋中文譯名」。
%   3. 課堂模式開啟時，GG/條文編者註仍會自動顯示。
% 進階微調才需要在單案檔覆寫：
%   \classroomabbrzhoff     縮寫註解僅顯示德文全稱。
%   \showeditornotestrue    一般版也顯示編者註。
% 註解策略：同一實際 PDF 頁面中，同一縮寫只在第一次出現處加註。
% 這裡用 zref-abspage 的絕對頁序，不用 \thepage，避免文件內頁碼重置時誤判。
\newif\ifclassroommode
\newif\ifclassroomabbrzh
\classroommodefalse
\classroomabbrzhtrue
\newcommand{\classroommodeon}{\classroommodetrue}
\newcommand{\classroommodeoff}{\classroommodefalse}
\newcommand{\classroomabbrzhon}{\classroomabbrzhtrue}
\newcommand{\classroomabbrzhoff}{\classroomabbrzhfalse}
\newcommand{\abbrnotelabel}{\emph{Abk.}: }
\newcommand{\abbrnote}[3]{%
  \ifclassroommode
    \ifshoworiginal
      \ifintranslation\else
        \footnote{\normalfont\footnotesize\abbrnotelabel #2\ifclassroomabbrzh；{\songfont #3}\fi}%
      \fi
    \fi
  \fi
}
\newcommand{\abbrnoteonce}[4]{%
  #2%
  \ifclassroommode
    \ifshoworiginal
      \ifintranslation\else
        \begingroup
          \edef\abbrcurrentabspage{\number\numexpr\value{abspage}+1\relax}%
          \ifcsname abbrnoteseen@#1@\abbrcurrentabspage\endcsname\else
            \global\expandafter\let\csname abbrnoteseen@#1@\abbrcurrentabspage\endcsname\empty
            \abbrnote{#1}{#3}{#4}%
          \fi
        \endgroup
      \fi
    \fi
  \fi
}
\newcommand{\abbrAbs}{\abbrnoteonce{abs}{Abs.}{Absatz}{項}}
\newcommand{\abbrAAO}{\abbrnoteonce{aao}{a.a.O.}{am angegebenen Ort}{前引處}}
\newcommand{\abbrAE}{\abbrnoteonce{ae}{AE}{Alternativ-Entwurf}{替代草案}}
\newcommand{\abbrArt}{\abbrnoteonce{art}{Art.}{Artikel}{條}}
\newcommand{\abbrAufl}{\abbrnoteonce{aufl}{Aufl.}{Auflage}{版}}
\newcommand{\abbrBd}{\abbrnoteonce{bd}{Bd.}{Band}{卷}}
\newcommand{\abbrBGBl}{\abbrnoteonce{bgbl}{BGBl.}{Bundesgesetzblatt}{聯邦法律公報}}
\newcommand{\abbrBGHSt}{\abbrnoteonce{bghst}{BGHSt}{Entscheidungen des Bundesgerichtshofes in Strafsachen}{聯邦普通法院刑事判決彙編}}
\newcommand{\abbrBRDrucks}{\abbrnoteonce{brdrucks}{BRDrucks.}{Bundesrats-Drucksache}{聯邦參議院文件}}
\newcommand{\abbrBTDrucks}{\abbrnoteonce{btdrucks}{BTDrucks.}{Bundestags-Drucksache}{聯邦議院文件}}
\newcommand{\abbrBundesgesetzbl}{\abbrnoteonce{bundesgesetzbl}{Bundesgesetzbl.}{Bundesgesetzblatt}{聯邦法律公報}}
\newcommand{\abbrBvF}{\abbrnoteonce{bvf}{BvF}{Bundesverfassungsgerichtsverfahren}{聯邦憲法法院程序案號}}
\newcommand{\abbrBvL}{\abbrnoteonce{bvl}{BvL}{Bundesverfassungsgerichtsverfahren}{聯邦憲法法院程序案號}}
\newcommand{\abbrBvR}{\abbrnoteonce{bvr}{BvR}{Bundesverfassungsgerichtsverfahren}{聯邦憲法法院程序案號}}
\newcommand{\abbrBVerfG}{\abbrnoteonce{bverfg}{BVerfG}{Bundesverfassungsgericht}{聯邦憲法法院}}
\newcommand{\abbrBVerfGE}{\abbrnoteonce{bverfge}{BVerfGE}{Entscheidungen des Bundesverfassungsgerichts}{聯邦憲法法院判決集}}
\newcommand{\abbrBVerfGG}{\abbrnoteonce{bverfgg}{BVerfGG}{Bundesverfassungsgerichtsgesetz}{聯邦憲法法院法}}
\newcommand{\abbrDDR}{\abbrnoteonce{ddr}{DDR}{Deutsche Demokratische Republik}{德意志民主共和國；東德}}
\newcommand{\abbrDJT}{\abbrnoteonce{djt}{DJT}{Deutscher Juristentag}{德國法學家大會}}
\newcommand{\abbrDr}{\abbrnoteonce{dr}{Dr.}{Doktor}{Dr.}}
\newcommand{\abbrEuGRZ}{\abbrnoteonce{eugrz}{EuGRZ}{Europäische Grundrechte-Zeitschrift}{歐洲基本權期刊}}
\newcommand{\abbrEV}{\abbrnoteonce{ev}{EV}{Einigungsvertrag}{統一條約}}
\newcommand{\abbrF}{\abbrnoteonce{f}{f.}{folgende}{以下一頁；下一段}}
\newcommand{\abbrFF}{\abbrnoteonce{ff}{ff.}{fortfolgende}{以下數頁；以下各段}}
\newcommand{\abbrGBlDDR}{\abbrnoteonce{gblddr}{GBl. DDR}{Gesetzblatt der Deutschen Demokratischen Republik}{德意志民主共和國法律公報}}
\newcommand{\abbrGG}{\abbrnoteonce{gg}{GG}{Grundgesetz}{基本法}}
\newcommand{\abbrGVBl}{\abbrnoteonce{gvbl}{GVBl.}{Gesetz- und Verordnungsblatt}{法律及命令公報}}
\newcommand{\abbrJR}{\abbrnoteonce{jr}{JR}{Juristische Rundschau}{法學評論}}
\newcommand{\abbrJZ}{\abbrnoteonce{jz}{JZ}{JuristenZeitung}{法學家報}}
\newcommand{\abbrIVm}{\abbrnoteonce{ivm}{i.V.m.}{in Verbindung mit}{結合；連同}}
\newcommand{\abbrMdB}{\abbrnoteonce{mdb}{MdB}{Mitglied des Bundestages}{聯邦議院議員}}
\newcommand{\abbrMwN}{\abbrnoteonce{mwn}{m.w.N.}{mit weiteren Nachweisen}{附進一步文獻／判例出處}}
\newcommand{\abbrNF}{\abbrnoteonce{nf}{n.F.}{neue Fassung}{新版本}}
\newcommand{\abbrAF}{\abbrnoteonce{af}{a.F.}{alte Fassung}{舊版本}}
\newcommand{\abbrNr}{\abbrnoteonce{nr}{Nr.}{Nummer}{款、目或號，依語境}}
\newcommand{\abbrProf}{\abbrnoteonce{prof}{Prof.}{Professor}{教授}}
\newcommand{\abbrRdnr}{\abbrnoteonce{rdnr}{Rdnr.}{Randnummer}{邊碼；段碼}}
\newcommand{\abbrRdnrn}{\abbrnoteonce{rdnrn}{Rdnrn.}{Randnummern}{邊碼；段碼}}
\newcommand{\abbrRGBl}{\abbrnoteonce{rgbl}{RGBl.}{Reichsgesetzblatt}{帝國法律公報}}
\newcommand{\abbrRGSt}{\abbrnoteonce{rgst}{RGSt}{Entscheidungen des Reichsgerichts in Strafsachen}{帝國法院刑事判決彙編}}
\newcommand{\abbrRspr}{\abbrnoteonce{rspr}{Rspr.}{Rechtsprechung}{判例；司法實務}}
\newcommand{\abbrRVO}{\abbrnoteonce{rvo}{RVO}{Reichsversicherungsordnung}{帝國保險法}}
\newcommand{\abbrS}{\abbrnoteonce{s}{S.}{Seite}{頁}}
\newcommand{\abbrSFHG}{\abbrnoteonce{sfhg}{SFHG}{Schwangeren- und Familienhilfegesetz}{孕婦及家庭扶助法}}
\newcommand{\abbrSGBV}{\abbrnoteonce{sgbv}{SGB V}{Fünftes Buch Sozialgesetzbuch}{社會法典第五編}}
\newcommand{\abbrStAG}{\abbrnoteonce{staeg}{StÄG}{Strafrechtsänderungsgesetz}{刑法修正法}}
\newcommand{\abbrStenBer}{\abbrnoteonce{stenber}{StenBer.}{Stenographischer Bericht}{速記紀錄}}
\newcommand{\abbrStGB}{\abbrnoteonce{stgb}{StGB}{Strafgesetzbuch}{刑法}}
\newcommand{\abbrStREG}{\abbrnoteonce{streg}{StREG}{Strafrechtsreform-Ergänzungsgesetz}{刑法改革補充法}}
\newcommand{\abbrStrRG}{\abbrnoteonce{strrg}{StrRG}{Strafrechtsreformgesetz}{刑法改革法}}
\newcommand{\abbrUS}{\abbrnoteonce{us}{U.S.}{United States Reports}{美國最高法院判決彙編}}
\newcommand{\abbrVgl}{\abbrnoteonce{vgl}{vgl.}{vergleiche}{參見}}
\newcommand{\abbrVVDStRL}{\abbrnoteonce{vvdstrl}{VVDStRL}{Veröffentlichungen der Vereinigung der Deutschen Staatsrechtslehrer}{德國公法教師協會論文集}}
\newcommand{\abbrWP}{\abbrnoteonce{wp}{WP.}{Wahlperiode}{屆期}}
\newcommand{\abbrWp}{\abbrnoteonce{wp}{Wp.}{Wahlperiode}{屆期}}
\newcommand{\abbrZStrW}{\abbrnoteonce{zstrw}{ZStrW}{Zeitschrift für die gesamte Strafrechtswissenschaft}{整體刑法學期刊}}
\newcommand{\recordsourcerandstart}{%
  \ifrandnumbers
    \setcounter{lastsourcerandstart}{\value{randnumber}}%
    \global\lastsourcehasrandnumberstrue
  \else
    \global\lastsourcehasrandnumbersfalse
  \fi
}
\newlength{\biparagraphskip}
\setlength{\biparagraphskip}{0.52em}
\newlength{\lawboxblockbeforeskip}
\setlength{\lawboxblockbeforeskip}{0.72em}
\newlength{\lawboxblockafterskip}
\setlength{\lawboxblockafterskip}{0.92em}
\newif\iflawboxpartendedblock
\lawboxpartendedblockfalse
\newcommand{\sourceparstyle}{\footnotesize\color{sourceink}\setstretch{1.12}\RaggedRight\emergencystretch=3em\setlength{\parskip}{0.42em}\obeylines\selectfont}
\newcommand{\transparstyle}{\songfont\color{caseink}\setstretch{1.12}\setlength{\parindent}{0pt}\setlength{\parskip}{0.42em}}
\newcommand{\bilingualpairskip}{%
  \par\addvspace{\biparagraphskip}%
  \switchcolumn
  \par\addvspace{\biparagraphskip}%
  \switchcolumn*%
  \global\intranslationfalse
}
\newcommand{\bilingualboxpairskip}{%
  \par
  \switchcolumn
  \par
  \switchcolumn*%
  \global\intranslationfalse
}
\newcommand{\bilinguallawboxblockskip}{%
  \switchcolumn*%
  \par\addvspace{\lawboxblockafterskip}%
  \switchcolumn
  \par\addvspace{\lawboxblockafterskip}%
  \switchcolumn*%
  \global\intranslationfalse
}
\newcommand{\bikeepwithnext}[1][4]{%
  \ifbilingualcolumns
    \syncleft
    \Needspace{#1\baselineskip}%
    \switchcolumn
    \Needspace{#1\baselineskip}%
    \switchcolumn*%
    \global\intranslationfalse
  \else
    \Needspace{#1\baselineskip}%
  \fi
}
\NewEnviron{sourcepara}[1][]{
  \recordsourcerandstart
  \ifshoworiginal
    \ifbilingualcolumns
      \ifintranslation\switchcolumn*\global\intranslationfalse\fi
    \else
      \Needspace{5\baselineskip}
    \fi
    \ifbilingualcolumns\else\par\addvspace{0.35em}\fi
    {\sourceparstyle
    \noindent\BODY\par}
    \ifbilingualcolumns\else\addvspace{0.25em}\fi
    \ifbilingualcolumns
      \switchcolumn
      \global\intranslationtrue
    \fi
  \else
    \ifrandnumbers
      \begingroup
        \setbox0=\vbox{\hsize=\linewidth\sourceparstyle\noindent\BODY\par}%
      \endgroup
    \fi
  \fi
}
\NewEnviron{transpara}{
  \ifshowtranslation
    \setcounter{randnumberlocal}{\value{lastsourcerandstart}}%
    \ifbilingualcolumns
      \ifintranslation\else\switchcolumn\global\intranslationtrue\fi
      {\transparstyle\setlength{\leftskip}{0.55em}\setlength{\rightskip}{0.15em plus 1em}\addtolength{\linewidth}{-0.55em}\BODY\par}
      \switchcolumn*
      \bilingualpairskip
    \else
      {\transparstyle\BODY\par}
    \fi
  \else
    \ifbilingualcolumns
      \ifintranslation\switchcolumn*\global\intranslationfalse\fi
    \fi
  \fi
}
\NewEnviron{sourceboxpara}[1][]{
  \global\lawboxpartendedblockfalse
  \recordsourcerandstart
  \ifshoworiginal
    \ifbilingualcolumns
      \ifintranslation\switchcolumn*\global\intranslationfalse\fi
    \else
      \Needspace{3\baselineskip}
    \fi
    {\sourceparstyle
    \BODY\par}
    \ifbilingualcolumns\else
      \iflawboxpartendedblock\addvspace{\lawboxblockafterskip}\fi
    \fi
    \ifbilingualcolumns
      \switchcolumn
      \global\intranslationtrue
    \fi
  \fi
}
\NewEnviron{transboxpara}{
  \global\lawboxpartendedblockfalse
  \ifshowtranslation
    \setcounter{randnumberlocal}{\value{lastsourcerandstart}}%
    \ifbilingualcolumns
      \ifintranslation\else\switchcolumn\global\intranslationtrue\fi
      {\transparstyle\BODY\par}
      \iflawboxpartendedblock
        \bilinguallawboxblockskip
      \else
        \switchcolumn*
        \bilingualboxpairskip
      \fi
    \else
      {\transparstyle\BODY\par}
      \iflawboxpartendedblock\addvspace{\lawboxblockafterskip}\fi
    \fi
  \else
    \ifbilingualcolumns
      \ifintranslation\switchcolumn*\global\intranslationfalse\fi
    \fi
  \fi
}
\newcommand{\leitsatznum}[1]{\noindent\hangindent=3.0em\hangafter=1\makebox[2.25em][r]{\bfseries #1.}\hspace{0.55em}\ignorespaces}
\newcommand{\syncleft}{\ifbilingualcolumns\ifintranslation\switchcolumn*\global\intranslationfalse\fi\fi}
\pretocmd{\section}{\syncleft\clearpage}{}{}
\pretocmd{\subsection}{\syncleft}{}{}
\pretocmd{\subsubsection}{\syncleft}{}{}
\newcommand{\biheadingfont}{\songfont}
\newcommand{\bitocentry}[2]{%
  \ifshoworiginal
    \ifshowtranslation
      \ifstrequal{#1}{#2}{#1}{#1\space #2}%
    \else
      #1%
    \fi
  \else
    #2%
  \fi
}
\pdfstringdefDisableCommands{%
  \def\bitocentry#1#2{#1}%
}
\newcommand{\biheading}[7]{%
  \syncleft#1\bikeepwithnext[#3]\phantomsection\addcontentsline{toc}{#2}{\bitocentry{#6}{#7}}%
  \ifbilingualcolumns
    {#4 #6\par}%
    \switchcolumn
    {#4\biheadingfont #7\par}%
    \switchcolumn*%
    \global\intranslationfalse
    \par\addvspace{#5}%
    \switchcolumn
    \par\addvspace{#5}%
    \switchcolumn*%
  \else
    \ifshoworiginal{#4 #6\par}\fi
    \ifshowtranslation{#4\biheadingfont #7\par}\fi
    \addvspace{#5}%
  \fi
}
\newcommand{\termsectionheading}[1]{%
  \par\Needspace{9\baselineskip}%
  \vspace{0.65em}%
  {\songfont\bfseries #1\par}%
  \vspace{0.2em}%
}
% 章節換頁：
%   雙語 paracol 欄內不要用 \clearpage；它會在同步兩欄時吐出只有欄線/頁碼的空頁。
%   \newpage 足以讓下一個 section 起新頁，又不會強制清空所有待排材料。
%   單語模式不在 paracol 內，仍可用 \clearpage。
\newcommand{\bisectionpagebreak}{\ifbilingualcolumns\newpage\else\clearpage\fi}
\newcommand{\bisection}[2]{%
  \biheading{\iffirstbilingualsection\global\firstbilingualsectionfalse\else\bisectionpagebreak\fi}{section}{6}{\centering\Large\bfseries}{0.8em}{#1}{#2}%
}
\newcommand{\bisubsection}[2]{%
  \biheading{}{subsection}{5}{\large\bfseries}{0.55em}{#1}{#2}%
}
\newcommand{\bisubsubsection}[2]{%
  \biheading{}{subsubsection}{4}{\normalsize\bfseries}{0.45em}{#1}{#2}%
}
\newcommand{\bicompositeheading}[4]{%
  \biheading{}{subsection}{5}{\large\bfseries}{0.16em}{#1}{#2}%
  \biheading{}{subsubsection}{3}{\normalsize\bfseries}{0.45em}{#3}{#4}%
}
\newif\iflawboxfirstitem
\lawboxfirstitemfalse
\newcommand{\lawboxprespace}[1]{%
  \par
  \iflawboxfirstitem
    \global\lawboxfirstitemfalse
  \else
    \addvspace{#1}%
  \fi
}
\newtcolorbox{lawboxinner}{
  caseboxbase,
  colframe=sourceline!85!black,
  colback=white,
  left=1.05em,
  right=1.05em,
  top=0.62em,
  bottom=0.62em,
  before upper={\global\inlawboxtrue\global\lawboxfirstitemtrue\global\lawboxhaspendingrnfalse\ifintranslation\lawboxtransfont\else\lawboxsourcefont\fi\RaggedRight\setlength{\parindent}{0pt}\setlength{\parskip}{0pt}\ifintranslation\setstretch{\lawboxtransstretch}\else\setstretch{\lawboxsourcestretch}\fi},
  after upper={\global\lawboxhaspendingrnfalse\global\inlawboxfalse}
}
\tcbset{
  lawboxpartbase/.style={
    enhanced,
    breakable=false,
    width=0.985\linewidth,
    colback=white,
    frame hidden,
    boxrule=0pt,
    arc=0pt,
    left=1.05em,
    right=1.05em,
    boxsep=1.5mm,
    before skip=0pt,
    after skip=0pt,
    before upper={\global\inlawboxtrue\global\lawboxfirstitemtrue\global\lawboxhaspendingrnfalse\ifintranslation\lawboxtransfont\else\lawboxsourcefont\fi\RaggedRight\setlength{\parindent}{0pt}\setlength{\parskip}{0pt}\ifintranslation\setstretch{\lawboxtransstretch}\else\setstretch{\lawboxsourcestretch}\fi},
    after upper={\global\lawboxhaspendingrnfalse\global\inlawboxfalse},
    borderline west={0.28pt}{0pt}{sourceline!85!black},
    borderline east={0.28pt}{0pt}{sourceline!85!black},
  },
  lawboxpartfirst/.style={
    lawboxpartbase,
    top=0.62em,
    bottom=0.04em,
    before skip=\lawboxblockbeforeskip,
    borderline north={0.28pt}{0pt}{sourceline!85!black},
  },
  lawboxpartmiddle/.style={
    lawboxpartbase,
    top=0.04em,
    bottom=0.04em,
  },
  lawboxpartlast/.style={
    lawboxpartbase,
    top=0.04em,
    bottom=0.62em,
    after skip=0pt,
    borderline south={0.28pt}{0pt}{sourceline!85!black},
  }
}
\newtcolorbox{lawboxpartinner}[1]{#1}
\newtcolorbox{termboxinner}{
  caseboxbase,
  colframe=noteblue!70!black,
  colback=casepaper,
  before upper={\songfont\setlength{\parindent}{0pt}\setlength{\parskip}{0.35em}}
}
\newtcolorbox{deboxinner}{
  caseboxbase,
  colframe=notegray!72!black,
  colback=white,
  left=1.05em,
  right=1.05em,
  top=0.62em,
  bottom=0.62em,
  before upper={\global\inlawboxtrue\global\lawboxfirstitemtrue\global\lawboxhaspendingrnfalse\small\setlength{\parindent}{0pt}\setlength{\parskip}{0.35em}},
  after upper={\global\lawboxhaspendingrnfalse\global\inlawboxfalse}
}
\NewEnviron{lawbox}{
  \begin{lawboxinner}\BODY\end{lawboxinner}
}
\NewEnviron{lawboxpart}[2]{
  \begin{lawboxpartinner}{lawboxpart#1,equal height group=#2}\BODY\end{lawboxpartinner}%
  \ifstrequal{#1}{last}{\global\lawboxpartendedblocktrue}{}%
}
\NewEnviron{termbox}{
  \par\addvspace{0.55em}
  {\color{noteblue!70!black}\rule{\linewidth}{0.28pt}}\par
  \begingroup
    \songfont\small\color{caseink}
    \setlength{\parindent}{0pt}
    \setlength{\parskip}{0.24em}
    \BODY
    \par
  \endgroup
  {\color{noteblue!70!black}\rule{\linewidth}{0.28pt}}\par
  \addvspace{0.55em}
}
\NewEnviron{debox}{
  \begin{deboxinner}\BODY\end{deboxinner}
}
\providecommand{\paragraphmark}[1]{}
\newcommand{\lawboxtitleafter}{\addvspace{0.06em}}
\newcommand{\lawtitle}[1]{\lawboxprespace{0.22em}{\lawboxtransfont\bfseries\lawboxuserandnumber #1}\par\lawboxtitleafter}
\newcommand{\absno}[2]{\lawboxprespace{0.04em}\noindent\lawboxuserandnumber\hangindent=2.6em\hangafter=1\makebox[2.6em][l]{(#1)}#2\par}
\newcommand{\nrno}[2]{\lawboxprespace{0pt}\noindent\lawboxuserandnumber\hspace*{2.6em}\hangindent=4.4em\hangafter=1\makebox[1.8em][l]{#1.}#2\par}
\newcommand{\letterno}[2]{\lawboxprespace{0pt}\noindent\lawboxuserandnumber\hspace*{4.4em}\hangindent=6.0em\hangafter=1\makebox[1.6em][l]{#1)}#2\par}
\newcommand{\sourcelawtitle}[1]{\lawboxprespace{0.22em}{\lawboxsourcefont\bfseries\lawboxuserandnumber #1}\par\lawboxtitleafter}
\newcommand{\sourcelawsubtitle}[1]{\lawboxprespace{0.04em}{\lawboxsourcefont\bfseries\lawboxuserandnumber #1}\par\lawboxtitleafter}
\newcommand{\sourceabsno}[2]{\lawboxprespace{0.04em}\noindent\lawboxuserandnumber\hangindent=2.45em\hangafter=1\makebox[2.45em][l]{(#1)}#2\par}
\newcommand{\sourcenrno}[2]{\lawboxprespace{0pt}\noindent\lawboxuserandnumber\hspace*{2.45em}\hangindent=4.25em\hangafter=1\makebox[1.8em][l]{#1.}#2\par}
\newcommand{\sourceletterno}[2]{\lawboxprespace{0pt}\noindent\lawboxuserandnumber\hspace*{4.25em}\hangindent=5.85em\hangafter=1\makebox[1.6em][l]{#1)}#2\par}
\newcommand{\sourcequotepara}[1]{\lawboxprespace{0.04em}\noindent\lawboxuserandnumber\hangindent=1.2em\hangafter=1 #1\par}
\newcommand{\sourcelawline}[1]{\lawboxprespace{0.04em}\noindent\lawboxuserandnumber #1\par}
\newcommand{\lawline}[1]{\lawboxprespace{0.04em}\noindent\lawboxuserandnumber #1\par}
\newcommand{\pendingtranslation}{\par{\songfont\footnotesize\color{pendingink}（中文譯文待補。）}\par}
\newcommand{\tocpart}[1]{\par\addvspace{0.52em}{\footnotesize\bfseries\color{pendingink}#1\par}\addvspace{0.16em}}
\newcommand{\tocdashline}{\par\addvspace{0.48em}\noindent{\color{notegray}\xleaders\hbox to 0.82em{\hfil\rule{0.42em}{0.28pt}\hfil}\hskip\linewidth}\par\addvspace{0.38em}}
% \tocpanel{font-cmd}{content}：將目錄置中於所在欄寬（雙語欄適用）。
% A3 橫式使用 0.58\linewidth，A4 橫式使用 0.84\linewidth，
% 使兩種紙張下目錄欄內容實際寬度相近（均約 100–105 mm）。
\newcommand{\tocpanel}[2]{%
  \makebox[\linewidth][c]{%
    \ifx\bilingualpaper\bilingualpaperaThree
      \begin{minipage}{0.58\linewidth}%
    \else
      \begin{minipage}{0.84\linewidth}%
    \fi
    \toccolstyle
    #1#2%
    \end{minipage}%
  }%
}
% ===== 首頁簡目錄的兩層 description list 縮排 =====
% enumitem 的 description list 可把每一列拆成：
%   [label]  labelsep  body text
% 這裡的「起點」都以所在 minipage/欄位的左緣為 0。
%
% 核心規則：
%   labelindent = label 欄位從左緣開始的位置。
%   labelwidth  = label 欄位本身寬度；標籤太長（如 Abw. Meinung）就加大它。
%   labelsep    = label 欄位右緣到正文的距離。
%   leftmargin  = 正文 body text 的起點。判斷層級視覺時主要看這個。
%
% 所以：
%   想讓「Begründung / 判決理由」整列正文往右或往左：改 tocitems 的 leftmargin。
%   想讓「A. Verfahren und Gegenstand / A. 程序與審查標的」正文往右或往左：
%     改 tocsubitems 的 leftmargin。
%   想讓 A./B./C. 這些標籤本身往右或往左，但正文不動：改 tocsubitems 的 labelindent。
%   想讓 A./B./C. 與正文間距變大或變小：改 tocsubitems 的 labelsep。
%
% 層級原則：
%   tocsubitems.leftmargin 必須大於 tocitems.leftmargin，
%   否則子層級正文會看起來比「Gründe / 理由」還靠左。
\newlist{tocitems}{description}{1}
\setlist[tocitems]{%
  labelindent=0pt,    % 第一層 label 欄位起點；0pt 表示貼齊目錄欄左緣。
  leftmargin=2.54cm,  % 第一層正文起點；例如 Begründung / 判決理由 的左緣。
  labelwidth=2.16cm,  % 第一層 label 欄位寬度；需容納 Leitsätze、Abw. Meinung、不同意見等。
  labelsep=0.32cm,    % 第一層 label 與正文的水平間距；只改間距，不改正文起點。
  itemsep=0.28em,     % 第一層各項之間的垂直距離。
  topsep=0.20em,      % 第一層 list 上下與前後內容的垂直距離。
  parsep=0pt,         % 同一 item 內段落之間的距離；目錄項通常不需要。
  partopsep=0pt,      % list 若緊接段落開始時的額外距離；關閉以免目錄高度飄動。
  align=left,         % label 欄內左對齊；長短標籤不靠右貼齊。
  font=\normalfont\bfseries % label 的字型；正文不受這行影響。
}
\newlist{tocsubitems}{description}{1}
\setlist[tocsubitems]{%
  labelindent=1cm,    % 第二層 label 起點；只控制 A./B./C. 本身的位置。
  leftmargin=3.00cm,  % 第二層正文起點；必須大於 2.54cm，才會比 Begründung / 判決理由更右。
  labelwidth=0.5cm,   % 第二層 label 欄寬；A./B./C. 很短，可小於第一層。
  labelsep=1.5cm,    % 第二層 label 與正文間距；A. 與 Verfahren / 程序 的距離看這裡。
  itemsep=0.12em,     % 第二層各項較密，避免 A.–E. 佔太高。
  topsep=0.04em,      % 第二層 list 與 Gründe / 理由之間的垂直距離。
  parsep=0pt,         % 同一第二層 item 內段落距離；目錄項通常不需要。
  partopsep=0pt,      % 關閉額外段落起始距離，避免版面不穩。
  align=left,         % A./B./C. label 欄內左對齊。
  font=\normalfont\bfseries, % A./B./C. label 加粗；正文不受這行影響。
  before={}           % 不在第二層 list 前自動插入額外內容。
}
\newlist{termitems}{description}{1}
\ifbilingualcolumns
  \setlist[termitems]{%
    leftmargin=5.2cm,
    labelwidth=4.8cm,
    labelsep=0.35cm,
    itemsep=0.16em,
    topsep=0.2em,
    parsep=0pt,
    partopsep=0pt,
    font=\normalfont\bfseries,
    style=nextline
  }
\else
  \setlist[termitems]{%
    leftmargin=0pt,
    labelwidth=0pt,
    labelsep=0pt,
    itemsep=0.24em,
    topsep=0.2em,
    parsep=0pt,
    partopsep=0pt,
    font=\normalfont\bfseries,
    style=nextline
  }
\fi
\newlist{abbritems}{description}{1}
\setlist[abbritems]{%
  leftmargin=2.38cm,
  labelwidth=2.02cm,
  labelsep=0.28cm,
  itemsep=0.16em,
  topsep=0.2em,
  parsep=0pt,
  partopsep=0pt,
  align=left,
  font=\normalfont\bfseries
}
\ifbilingualcolumns
  \newenvironment{termcolumns}{\begin{multicols}{2}}{\end{multicols}}
\else
  \newenvironment{termcolumns}{}{}
\fi

% ===== 編者註腳開關 =====
% 一般版預設不顯示編者註，避免正文腳註過密。
% 若開啟 \classroommodeon，GG/條文編者註仍會自動顯示，無需另開本開關。
% 若一般版也要顯示編者註，可在單案檔 \input 後加 \showeditornotestrue。
\newif\ifshoweditornotes
\showeditornotesfalse

% ===== 目錄排版輔助 =====
\newcommand{\toccolstyle}{\raggedright\arraybackslash}
\newcommand{\tocsingleindent}{\leftskip=1.45cm\rightskip=1.2cm}

% ===== 行內引文環境（德文原文與中文譯文各一，帶左側灰色邊線）=====
\newcommand{\quoteparbreak}{\par\addvspace{0.48em}}
\newcommand{\sourcequoteinner}[1]{%
  \begin{tcolorbox}[
    enhanced, breakable, width=\dimexpr\linewidth-0.8em\relax, frame hidden,
    colback=white, coltext=caseink, boxrule=0pt,
    borderline west={0.55pt}{0.88em}{notegray},
    left=1.78em, right=0.72em, top=0.18em, bottom=0.18em,
    boxsep=0pt, before skip=0.24em, after skip=0.24em,
    before upper={\normalfont\footnotesize\color{caseink}\setlength{\parindent}{0pt}\setlength{\parskip}{0.34em}\raggedright\setlength{\rightskip}{0pt plus 1em}}
  ]%
  #1\par
  \end{tcolorbox}%
}
\newcommand{\transquoteinner}[1]{%
  \begin{tcolorbox}[
    enhanced, breakable, width=\dimexpr\linewidth-0.8em\relax, frame hidden,
    colback=white, coltext=caseink, boxrule=0pt,
    borderline west={0.55pt}{0.88em}{notegray},
    left=1.78em, right=0.72em, top=0.18em, bottom=0.18em,
    boxsep=0pt, before skip=0.24em, after skip=0.24em,
    before upper={\songfont\small\color{caseink}\setlength{\parindent}{0pt}\setlength{\parskip}{0.34em}\raggedright\setlength{\rightskip}{0pt plus 1em}}
  ]%
  #1\par
  \end{tcolorbox}%
}

% ===== 大綱（Gliederung）Rn 輔助命令 =====
\newcommand{\outrn}[2]{\enspace{\normalfont\footnotesize\color{pendingink}(Rn\,#1–#2)}}
\newcommand{\outrnsingle}[1]{\enspace{\normalfont\footnotesize\color{pendingink}(Rn\,#1)}}
\newcommand{\outrnzh}[2]{\enspace{\normalfont\footnotesize\color{pendingink}（第\,#1–#2\,段）}}
\newcommand{\outrnzhs}[1]{\enspace{\normalfont\footnotesize\color{pendingink}（第\,#1\,段）}}
% ===== 大綱雙語對照表格輔助命令（三欄：段碼 │ 德文 │ 中文）=====
% 欄 1：段碼（由欄定義自動格式化：小字、等寬、右對齊、pendingink）
% 欄 2：德文標籤＋正文
% 欄 3：中文標籤＋正文
%
% 無段碼的列（\omainrow、\osecrow），欄 1 以 & 空置。
% 有段碼的列（\osecrowrn、\osubrow、\ointrow），#1 為預格式化字串
%   例：\osubrow{2–49}{I.}{...}{I.}{...}
%       \ointrow{107}{...}{...}
%
% \opartrow：段組標題（橫跨三欄）
\newcommand{\opartrow}[2]{%
  \noalign{\vspace{0.30em}}%
  \multicolumn{3}{@{}l@{}}{%
    \footnotesize\bfseries\color{pendingink}#1\hfill\songfont#2%
  }\\[0.06em]%
}
% \odashrow：虛線分隔（橫跨三欄）
\newcommand{\odashrow}{%
  \noalign{\vspace{0.28em}}%
  \multicolumn{3}{@{}l@{}}{\color{notegray}%
    \xleaders\hbox to 0.82em{\hfil\rule{0.42em}{0.28pt}\hfil}\hskip 0.88\linewidth}%
  \\[0.24em]%
}
% \odissentpart：不同意見書區段標頭。比一般 \opartrow 更像附錄性轉場。
\newcommand{\odissentpart}[2]{%
  \noalign{\vspace{0.42em}}%
  \multicolumn{3}{@{}p{0.88\textwidth}@{}}{%
    \color{notegray}\rule{\linewidth}{0.18pt}%
  }\\[-0.18em]%
  \multicolumn{3}{@{}p{0.88\textwidth}@{}}{%
    \makebox[\linewidth][s]{%
      \footnotesize\bfseries\color{pendingink}#1%
      \hfill{\songfont #2}%
    }%
  }\\[-0.02em]%
}
% \odissentopinion：不同意見書的署名與一行摘要，右欄保留段碼範圍。
\newcommand{\odissentopinion}[5]{%
  \footnotesize\hangindent=0.52cm\hangafter=1%
    {\bfseries\color{caseink}#2}\enspace{\color{notegray}\textperiodcentered}\enspace\color{caseink}#3%
  &\footnotesize\hangindent=0.52cm\hangafter=1%
    {\songfont\bfseries\color{caseink}#4}\enspace{\color{notegray}\textperiodcentered}\enspace\songfont\color{caseink}#5%
  &#1%
  \\[0.12em]%
}
% \omainrow：頂層項目，無段碼（Leitsätze、Urteil、Gründe …）
\newcommand{\omainrow}[4]{%
  \footnotesize\hangindent=1.92cm\hangafter=1%
    \makebox[1.92cm][l]{\bfseries\color{caseink}#1}\color{caseink}#2%
  &\footnotesize\hangindent=1.92cm\hangafter=1%
    \makebox[1.92cm][l]{\bfseries\color{caseink}#3}\songfont\color{caseink}#4%
  &%
  \\[0.15em]%
}
% \omainrowrn：頂層項目，有段碼（不同意見等標籤寬於 \osecrow 框者）
\newcommand{\omainrowrn}[5]{%
  \footnotesize\hangindent=1.92cm\hangafter=1%
    \makebox[1.92cm][l]{\bfseries\color{caseink}#2}\color{caseink}#3%
  &\footnotesize\hangindent=1.92cm\hangafter=1%
    \makebox[1.92cm][l]{\bfseries\color{caseink}#4}\songfont\color{caseink}#5%
  &#1%
  \\[0.15em]%
}
% \osecrow：一級子項，無段碼（A.、C.、D. 等有子項者）
\newcommand{\osecrow}[4]{%
  \footnotesize\hspace{1.10cm}\hangindent=1.98cm\hangafter=1%
    \makebox[0.86cm][l]{\bfseries\color{caseink}#1}\color{caseink}#2%
  &\footnotesize\hspace{1.10cm}\hangindent=1.98cm\hangafter=1%
    \makebox[0.86cm][l]{\bfseries\color{caseink}#3}\songfont\color{caseink}#4%
  &%
  \\[0.11em]%
}
% \osecrowrn：一級子項，有段碼（B.、E. 等完整段落範圍者）
% #1=段碼字串（如 101–106）, #2=DE標籤, #3=DE正文, #4=ZH標籤, #5=ZH正文
\newcommand{\osecrowrn}[5]{%
  \footnotesize\hspace{1.10cm}\hangindent=1.98cm\hangafter=1%
    \makebox[0.86cm][l]{\bfseries\color{caseink}#2}\color{caseink}#3%
  &\footnotesize\hspace{1.10cm}\hangindent=1.98cm\hangafter=1%
    \makebox[0.86cm][l]{\bfseries\color{caseink}#4}\songfont\color{caseink}#5%
  &#1%
  \\[0.11em]%
}
% \osubrow：二級子項，有段碼範圍（I.、II.、A.I. …）
% #1=段碼字串, #2=DE標籤, #3=DE正文, #4=ZH標籤, #5=ZH正文
\newcommand{\osubrow}[5]{%
  \footnotesize\hspace{1.40cm}\hangindent=2.30cm\hangafter=1%
    \makebox[0.90cm][l]{\bfseries\color{caseink}#2}\color{caseink}#3%
  &\footnotesize\hspace{1.40cm}\hangindent=2.30cm\hangafter=1%
    \makebox[0.90cm][l]{\bfseries\color{caseink}#4}\songfont\color{caseink}#5%
  &#1%
  \\[0.09em]%
}
% \ointrow：引入段落，有單一段碼，無標籤
% #1=段碼字串, #2=DE正文, #3=ZH正文
\newcommand{\ointrow}[3]{%
  \footnotesize\hspace{0.52cm}\hangindent=0.52cm\hangafter=1\color{caseink}#2%
  &\footnotesize\hspace{0.52cm}\hangindent=0.52cm\hangafter=1\songfont\color{caseink}#3%
  &#1%
  \\[0.09em]%
}
% \osubsubrow：三級子項（1.、2.、3.）
% #1=段碼字串, #2=DE標籤, #3=DE正文, #4=ZH標籤, #5=ZH正文
\newcommand{\osubsubrow}[5]{%
  \footnotesize\hspace{1.80cm}\hangindent=2.48cm\hangafter=1%
    \makebox[0.68cm][l]{\bfseries\color{caseink}#2}\color{caseink}#3%
  &\footnotesize\hspace{1.80cm}\hangindent=2.48cm\hangafter=1%
    \makebox[0.68cm][l]{\bfseries\color{caseink}#4}\songfont\color{caseink}#5%
  &#1%
  \\[0.08em]%
}
% \osubsubsubrow：四級子項（a）、b））
% #1=段碼字串, #2=DE標籤, #3=DE正文, #4=ZH標籤, #5=ZH正文
\newcommand{\osubsubsubrow}[5]{%
  \footnotesize\hspace{2.10cm}\hangindent=2.68cm\hangafter=1%
    \makebox[0.58cm][l]{\bfseries\color{caseink}#2}\color{caseink}#3%
  &\footnotesize\hspace{2.10cm}\hangindent=2.68cm\hangafter=1%
    \makebox[0.58cm][l]{\bfseries\color{caseink}#4}\songfont\color{caseink}#5%
  &#1%
  \\[0.07em]%
}
% \osubsubsubsubrow：五級子項（aa）、bb））
% 標籤框 0.62cm：足以容納「aa)」在 footnotesize bold 下（約 0.50cm）並留有間距
% #1=段碼字串, #2=DE標籤, #3=DE正文, #4=ZH標籤, #5=ZH正文
\newcommand{\osubsubsubsubrow}[5]{%
  \footnotesize\hspace{2.38cm}\hangindent=3.06cm\hangafter=1%
    \makebox[0.68cm][l]{\bfseries\color{caseink}#2}\color{caseink}#3%
  &\footnotesize\hspace{2.38cm}\hangindent=3.06cm\hangafter=1%
    \makebox[0.68cm][l]{\bfseries\color{caseink}#4}\songfont\color{caseink}#5%
  &#1%
  \\[0.06em]%
}
% \outlinepage：雙語對照大綱頁包裝環境（longtable 自動跨頁）
% 三欄：德文欄 + 中文欄 + 段碼欄（最右，不折行）
% 段碼欄寬度：1.80cm，足以容納最長段碼如「309–348」
% DE/ZH 欄寬：(0.87\textwidth - 2.08cm) / 2
%   驗算：2×欄 + 0.05\textwidth + 0.28cm gap + 1.80cm = 0.87tw - 2.08cm + 0.05tw + 2.08cm = 0.92tw ✓
\newenvironment{outlinepage}{%
  \vspace*{0.40cm}%
  \setlength{\LTleft}{0.04\textwidth}%
  \setlength{\LTright}{0.04\textwidth}%
  \setlength{\tabcolsep}{0pt}%
  % 大綱列距由 longtable 的 \arraystretch 與各 row macro 結尾的 \\[...em] 共同決定。
  % 這裡控制「每一列本身的表格 strut 高度」；數值越大，A./I./1. 各項之間越像有空行。
  % A3 原本用 0.62，視覺上沒有空行；A4 也沿用同值，避免切到 A4 後項目被撐開。
  % 若日後覺得大綱太擠，只微調這個值（例如 0.68），不要改各 \omainrow/\osubrow 的縮排。
  \renewcommand{\arraystretch}{0.62}%
  \begin{longtable}{%
    >{\vspace{0pt}\raggedright\arraybackslash}p{\dimexpr(0.87\textwidth-2.08cm)/2\relax}%
    @{\hspace{0.025\textwidth}}%
    !{\color{notegray}\vrule width 0.12pt}%
    @{\hspace{0.025\textwidth}}%
    >{\vspace{0pt}\raggedright\arraybackslash}p{\dimexpr(0.87\textwidth-2.08cm)/2\relax}%
    @{\hspace{0.28cm}}%
    >{\vspace{0pt}\footnotesize\normalfont\color{pendingink}\raggedright\arraybackslash}p{1.80cm}%
    @{}%
  }%
  % 首頁欄標籤：不要橫跨置中；分別放在德文欄與中文欄的實際欄位起點。
  \footnotesize\bfseries\color{sourceink}Gliederung%
  &\footnotesize\bfseries\songfont\color{sourceink}大綱%
  &%
  \\[0.36em]%
  \endfirsthead%
  % 續頁欄標籤：同樣分欄定位，以灰色細體區分；無需標注「續」。
  \footnotesize\color{pendingink}Gliederung%
  &\footnotesize\songfont\color{pendingink}大綱%
  &%
  \\[0.24em]%
  \endhead%
}{%
  \end{longtable}%
}
 載入。
% 不含：\documentclass、\documentversion、\ggversionde/zh、\tocde/\toczh。
% 日期與首頁版本列由本檔自動產生；各文件只需手動指定 \documentversion。

\usepackage{xeCJK}
\usepackage{fontspec}
\usepackage{geometry}
\usepackage{setspace}
\usepackage{hyperref}
\usepackage{xurl}
\usepackage{bookmark}
\usepackage{enumitem}
\usepackage{xcolor}
\usepackage{titlesec}
\usepackage{array}
\usepackage{longtable}
\usepackage{tcolorbox}
\usepackage{environ}
\usepackage{pdflscape}
\usepackage{etoolbox}
\usepackage{ragged2e}
\usepackage{paracol}
\usepackage{multicol}
\usepackage{needspace}
\usepackage{fancyhdr}
\usepackage{tikz}
\usepackage{zref-abspage}
\tcbuselibrary{breakable,skins}
\usetikzlibrary{calc}
\usepackage{pgfcalendar}

\hypersetup{
  unicode=true,
  bookmarksopen=true,
  bookmarksnumbered=false,
  hidelinks
}

% ===== 自動推導，不要手動改下面 =====
% （\docmode 與 \bilingualpaper 由各文件在 % bverfge_layout.tex
% 共用排版基礎設施，供 Schwangerschaftsabbruch I 與 II 共享。
% 由各文件以 \documentclass 後 \input{../../共用LaTeX/bverfge_layout} 載入。
% 不含：\documentclass、\documentversion、\ggversionde/zh、\tocde/\toczh。
% 日期與首頁版本列由本檔自動產生；各文件只需手動指定 \documentversion。

\usepackage{xeCJK}
\usepackage{fontspec}
\usepackage{geometry}
\usepackage{setspace}
\usepackage{hyperref}
\usepackage{xurl}
\usepackage{bookmark}
\usepackage{enumitem}
\usepackage{xcolor}
\usepackage{titlesec}
\usepackage{array}
\usepackage{longtable}
\usepackage{tcolorbox}
\usepackage{environ}
\usepackage{pdflscape}
\usepackage{etoolbox}
\usepackage{ragged2e}
\usepackage{paracol}
\usepackage{multicol}
\usepackage{needspace}
\usepackage{fancyhdr}
\usepackage{tikz}
\usepackage{zref-abspage}
\tcbuselibrary{breakable,skins}
\usetikzlibrary{calc}
\usepackage{pgfcalendar}

\hypersetup{
  unicode=true,
  bookmarksopen=true,
  bookmarksnumbered=false,
  hidelinks
}

% ===== 自動推導，不要手動改下面 =====
% （\docmode 與 \bilingualpaper 由各文件在 \input{bverfge_layout} 之前設定；
%   預設 chinese / a4）
\ifdefined\docmode\else\def\docmode{chinese}\fi
\ifdefined\bilingualpaper\else\def\bilingualpaper{a4}\fi
\newif\ifshoworiginal
\newif\ifshowtranslation
\newif\ifbilinguallandscape
\newif\ifbilingualcolumns
\newif\ifintranslation
\newif\iffirstbilingualsection

\showoriginalfalse
\showtranslationfalse
\bilinguallandscapefalse
\bilingualcolumnsfalse
\intranslationfalse
\firstbilingualsectionfalse

\def\modechinese{chinese}
\def\modegerman{german}
\def\modebilingual{bilingual}
\def\bilingualpaperaThree{a3}
\def\bilingualpaperaFour{a4}

\ifx\docmode\modechinese
  \showtranslationtrue
\fi

\ifx\docmode\modegerman
  \showoriginaltrue
\fi

\ifx\docmode\modebilingual
  \showoriginaltrue
  \showtranslationtrue
  \bilinguallandscapetrue
  \bilingualcolumnstrue
\fi

% 編排模式說明：
% chinese  → \showoriginalfalse  \showtranslationtrue  （A4 直式，純中文）
% german   → \showoriginaltrue   \showtranslationfalse （A4 直式，純德文）
% bilingual→ \showoriginaltrue   \showtranslationtrue  \bilingualcolumnstrue（A3/A4 橫式對照）

\ifbilingualcolumns
  \ifx\bilingualpaper\bilingualpaperaFour
    \geometry{
      a4paper,
      landscape,
      left=1.25cm,
      right=2.25cm,
      top=1.25cm,
      bottom=1.75cm,
      footskip=8.5mm,
      marginparwidth=0.75cm,
      marginparsep=0.25cm
    }
  \else
    \geometry{
      a3paper,
      landscape,
      left=1.55cm,
      right=2.75cm,
      top=1.55cm,
      bottom=2.05cm,
      footskip=9.5mm,
      marginparwidth=0.9cm,
      marginparsep=0.35cm
    }
  \fi
\else
\geometry{
  a4paper,
  portrait,
  left=2.2cm,
  right=2.6cm,
  top=2.0cm,
  bottom=2.0cm,
  marginparwidth=0.8cm,
  marginparsep=0.3cm
}
\fi
% ===== 時間戳基礎設施 =====
\newcount\stamptimeval
\newcount\stampminuteval
\newcommand{\stamphh}{%
  \stamptimeval=\time
  \divide\stamptimeval by 60
  \ifnum\stamptimeval<10 0\fi
  \the\stamptimeval}
\newcommand{\stampmm}{%
  \stamptimeval=\time
  \stampminuteval=\time
  \divide\stampminuteval by 60
  \multiply\stampminuteval by 60
  \advance\stamptimeval by -\stampminuteval
  \ifnum\stamptimeval<10 0\fi
  \the\stamptimeval}
\newcommand{\stampwkde}[1]{%
  \ifcase#1Montag\or Dienstag\or Mittwoch\or Donnerstag\or Freitag\or Samstag\or Sonntag\fi}
\newcommand{\stampwkzh}[1]{%
  \ifcase#1 一\or 二\or 三\or 四\or 五\or 六\or 日\fi}
\newcommand{\stampmonthde}[1]{%
  \ifcase#1\or Januar\or Februar\or März\or April\or Mai\or Juni\or
  Juli\or August\or September\or Oktober\or November\or Dezember\fi}
\newcount\stampjdval
\newcount\stampwdval
\pgfcalendardatetojulian{\the\year-\the\month-\the\day}{\stampjdval}
\pgfcalendarjuliantoweekday{\the\stampjdval}{\stampwdval}
\newcommand{\timestampzh}{%
  \songfont 最後修改：\the\year 年\the\month 月\the\day 日%
  % （星期\stampwkzh{\the\stampwdval}）%
  % \space\stamphh:\stampmm
  }

\newcommand{\documentdatezh}{\the\year 年 \the\month 月 \the\day 日}
\newcommand{\documentdatede}{\the\day. \stampmonthde{\the\month} \the\year}
\newcommand{\bverfgetranslationnote}{%
  {\songfont 翻譯說明：初譯使用 Claude Code 與 GPT 協助迭代；仍請以德文原文為準。}%
}
\newcommand{\bverfgecourseinfo}{%
  {\songfont 課程：114 學年度第 2 學期；憲法專題研究（四）；授課老師：湯德宗教授；導讀日期：115 年 6 月 1 日。}%
}
\newcommand{\bverfgeeditorinfo}{%
  {\songfont 編者：王逸帆（R10A21126），國立台灣大學法律系研究所碩士班。}%
}
\newcommand{\bverfgetitleversionline}{%
  \ifbilingualcolumns
    Fassung \documentversion\enspace\textperiodcentered\enspace Stand: \documentdatede
    \quad
    {\songfont 版本 \documentversion\enspace\textperiodcentered\enspace 修訂日期：\documentdatezh\par}%
  \else
    \ifshowtranslation
      {\songfont 版本 \documentversion\enspace\textperiodcentered\enspace 修訂日期：\documentdatezh\par}%
    \else
      Fassung \documentversion\enspace\textperiodcentered\enspace Stand: \documentdatede\par
    \fi
  \fi
}
\newcommand{\bverfgetitlecornerinfo}{%
  \ifshowtranslation
    \vspace{0.72em}%
    \noindent\hfill
    \begin{minipage}{0.66\textwidth}
      \raggedleft\tiny\color{pendingink}\songfont
      \bverfgecourseinfo\par
      \bverfgeeditorinfo
    \end{minipage}%
  \fi
}
\newcommand{\bverfgetitlemeta}{%
  \vfill
  \begin{center}%
    \footnotesize\color{pendingink}%
    \bverfgetitleversionline
    \ifshowtranslation
      \vspace{0.28em}{\scriptsize\bverfgetranslationnote\par}%
    \fi
  \end{center}%
  \bverfgetitlecornerinfo
}

\pagestyle{fancy}
\fancyhf{}
\renewcommand{\headrulewidth}{0pt}
\renewcommand{\footrulewidth}{0pt}
\fancyfoot[C]{\thepage}
\ifbilingualcolumns
  \fancyfoot[R]{\scriptsize\color{pendingink}\timestampzh}%
\else
  \ifshoworiginal
  \else
    \fancyfoot[R]{\scriptsize\color{pendingink}\timestampzh\hspace{0.45cm}}%
  \fi
\fi
\setmainfont{Times New Roman}
\setsansfont{Helvetica Neue}
\setmonofont{Menlo}
\IfFileExists{/Users/iw/Library/Fonts/kaiu.ttf}{
  \setCJKmainfont[Path=/Users/iw/Library/Fonts/,AutoFakeBold=4]{kaiu.ttf}
}{
  \IfFontExistsTF{標楷體}{
    \setCJKmainfont[AutoFakeBold=4]{標楷體}
  }{
    \setCJKmainfont{Songti TC}
  }
}
\IfFontExistsTF{Songti TC}{
  \setCJKfamilyfont{song}{Songti TC}
  \setCJKmonofont{Songti TC}
}{
  \setCJKfamilyfont{song}[Path=/Users/iw/Library/Fonts/,AutoFakeBold=4]{kaiu.ttf}
  \setCJKmonofont[Path=/Users/iw/Library/Fonts/,AutoFakeBold=4]{kaiu.ttf}
}
\newcommand{\songfont}{\CJKfamily{song}}

% ===== 封面／目錄專用打字機字體（僅限封面、目錄頁使用，不影響正文）=====
\IfFontExistsTF{Radio Newsman}{%
  \newfontfamily\bverfgetitlede{Radio Newsman}%
}{%
  \newfontfamily\bverfgetitlede{Courier New}%
}
\IfFontExistsTF{Erikas Farbband}{%
  \newfontfamily\bverfgebodyde{Erikas Farbband}%
}{%
  \let\bverfgebodyde\bverfgetitlede
}
\IfFontExistsTF{Maoken Heavy Labourer}{%
  \setCJKfamilyfont{bverfgetitlecjk}{Maoken Heavy Labourer}%
}{%
  \setCJKfamilyfont{bverfgetitlecjk}{Songti TC}%
}
\IfFontExistsTF{Huiwen-mincho}{%
  \setCJKfamilyfont{bverfgebodycjk}{Huiwen-mincho}%
}{%
  \setCJKfamilyfont{bverfgebodycjk}{Songti TC}%
}
\newcommand{\bverfgetitlezh}{\bverfgetitlede\CJKfamily{bverfgetitlecjk}}
\newcommand{\bverfgebodyzh}{\bverfgebodyde\CJKfamily{bverfgebodycjk}}

\setstretch{1.35}
\setlist{itemsep=0.2em, topsep=0.4em}
\setlength{\parindent}{0pt}
\setlength{\parskip}{0.45em}
\raggedbottom
\setcounter{secnumdepth}{0}
\setcounter{tocdepth}{2}

\newcommand{\lawboxsourcefont}{\footnotesize}
\newcommand{\lawboxtransfont}{\footnotesize}
\newcommand{\lawboxsourcestretch}{1.02}
\newcommand{\lawboxtransstretch}{0.92}

\titleformat{\section}{\centering\Large\bfseries\songfont}{\thesection}{1em}{}
\titleformat{\subsection}{\large\bfseries\songfont}{\thesubsection}{1em}{}
\titleformat{\subsubsection}{\normalsize\bfseries\songfont}{\thesubsubsection}{1em}{}
\titleformat{\paragraph}{\normalsize\bfseries\songfont}{\theparagraph}{1em}{}
\titlespacing*{\section}{0pt}{1.8em}{1.0em}
\titlespacing*{\subsection}{0pt}{1.2em}{0.6em}
\titlespacing*{\subsubsection}{0pt}{1.0em}{0.45em}
\titlespacing*{\paragraph}{0pt}{0.6em}{0.4em}

\definecolor{noteblue}{HTML}{A9C9F5}
\definecolor{noteteal}{HTML}{AEE1D6}
\definecolor{notegray}{HTML}{D7DADF}
\definecolor{caseink}{HTML}{000000}
\definecolor{casepaper}{HTML}{FCFEFD}
\definecolor{sourcepaper}{HTML}{FAFBF8}
\definecolor{sourceline}{HTML}{D7DED8}
\definecolor{sourceink}{HTML}{000000}
\definecolor{pendingink}{HTML}{000000}

\tcbset{
  caseboxbase/.style={
    enhanced,
    breakable,
    width=0.985\linewidth,
    colback=casepaper,
    coltitle=caseink,
    fonttitle=\bfseries\songfont,
    boxrule=0.28pt,
    arc=1.2mm,
    left=2.4mm,
    right=2.4mm,
    top=1.5mm,
    bottom=1.5mm,
    boxsep=1.5mm,
    before skip=0.65em,
    after skip=0.65em,
    before upper={\setlength{\parindent}{0pt}\setlength{\parskip}{0.35em}},
  }
}

\newcommand{\bilingualstart}{}
\newcommand{\bilingualstop}{}
\ifbilingualcolumns
  \renewcommand{\bilingualstart}{\small\setlength{\columnsep}{1.25cm}\setlength{\columnseprule}{0.12pt}\columnratio{0.5,0.5}\multicolumnfootnotes\flushbottom\sloppy\interlinepenalty=80\clubpenalty=800\widowpenalty=800\displaywidowpenalty=800\begin{paracol}{2}\global\intranslationfalse\global\firstbilingualsectiontrue{\footnotesize\color{sourceink}Deutsch\par}\switchcolumn{\footnotesize\songfont\color{sourceink}中文譯文\par}\switchcolumn*}
  \renewcommand{\bilingualstop}{\ifintranslation\switchcolumn*\global\intranslationfalse\fi\end{paracol}\clearpage\raggedbottom}
\fi
\newif\ifrandnumbers
\randnumbersfalse
\newcounter{randnumber}
\newcounter{randnumberlocal}
\newcounter{lastsourcerandstart}
\newif\iflastsourcehasrandnumbers
\lastsourcehasrandnumbersfalse
\newif\ifinlawbox
\inlawboxfalse
\newif\iflawboxhaspendingrn
\lawboxhaspendingrnfalse
\newcommand{\lawboxpendingrn}{}
\newcommand{\startrandnumbers}{\global\randnumberstrue\setcounter{randnumber}{1}\global\lastsourcehasrandnumbersfalse}
\newcommand{\stoprandnumbers}{\global\randnumbersfalse\global\lastsourcehasrandnumbersfalse}
\newlength{\randnumberwidth}
\setlength{\randnumberwidth}{0.95cm}
\newlength{\randnumberrightoffset}
\setlength{\randnumberrightoffset}{0.85cm}
\newlength{\randnumberlawboxrightoffset}
\setlength{\randnumberlawboxrightoffset}{1.40cm}
\newcommand{\randnumberbox}[1]{%
  \leavevmode
  \makebox[0pt][l]{%
    \smash{%
      \tikz[remember picture,overlay,baseline=(rnmark.base)]{%
        \coordinate (rnmark) at (0,0);%
        \ifinlawbox
          \path let \p1=(current page.east), \p2=(rnmark.base) in
            node[
              anchor=base east,
              inner sep=0pt,
              outer sep=0pt,
              text width=\randnumberwidth,
              align=right,
              text=pendingink,
              font=\normalfont\mdseries\scriptsize\ttfamily
            ] at (\x1-\randnumberlawboxrightoffset,\y2) {{\normalfont\mdseries\scriptsize\ttfamily #1}};%
        \else
          \path let \p1=(current page.east), \p2=(rnmark.base) in
            node[
              anchor=base east,
              inner sep=0pt,
              outer sep=0pt,
              text width=\randnumberwidth,
              align=right,
              text=pendingink,
              font=\normalfont\mdseries\scriptsize\ttfamily
            ] at (\x1-\randnumberrightoffset,\y2) {{\normalfont\mdseries\scriptsize\ttfamily #1}};%
        \fi
      }%
    }%
  }%
  \ignorespaces
}
\newcommand{\lawboxstorerandnumber}[1]{%
  \gdef\lawboxpendingrn{#1}%
  \global\lawboxhaspendingrntrue
  \ignorespaces
}
\newcommand{\lawboxuserandnumber}{%
  \iflawboxhaspendingrn
    \randnumberbox{\lawboxpendingrn}%
    \global\lawboxhaspendingrnfalse
  \fi
}
\newcommand{\sourcern}[1]{%
  \ifrandnumbers
    \ifshoworiginal
      \ifshowtranslation\else
        \ifinlawbox\lawboxstorerandnumber{#1}\else\randnumberbox{#1}\fi
      \fi
    \fi
  \fi
}
\newcommand{\transrn}[1]{%
  \ifrandnumbers
    \ifshowtranslation
      \ifinlawbox\lawboxstorerandnumber{#1}\else\randnumberbox{#1}\fi
    \fi
  \fi
}
% ===== 課堂模式：德文原文縮寫註解 =====
% 日常切換只改各判決檔的一行：
%   \classroommodeon        開啟課堂版；註解掉這行即回到一般版。
% 模板預設：
%   1. 課堂模式關閉。
%   2. 課堂模式開啟時，縮寫註解包含「德文全稱＋中文譯名」。
%   3. 課堂模式開啟時，GG/條文編者註仍會自動顯示。
% 進階微調才需要在單案檔覆寫：
%   \classroomabbrzhoff     縮寫註解僅顯示德文全稱。
%   \showeditornotestrue    一般版也顯示編者註。
% 註解策略：同一實際 PDF 頁面中，同一縮寫只在第一次出現處加註。
% 這裡用 zref-abspage 的絕對頁序，不用 \thepage，避免文件內頁碼重置時誤判。
\newif\ifclassroommode
\newif\ifclassroomabbrzh
\classroommodefalse
\classroomabbrzhtrue
\newcommand{\classroommodeon}{\classroommodetrue}
\newcommand{\classroommodeoff}{\classroommodefalse}
\newcommand{\classroomabbrzhon}{\classroomabbrzhtrue}
\newcommand{\classroomabbrzhoff}{\classroomabbrzhfalse}
\newcommand{\abbrnotelabel}{\emph{Abk.}: }
\newcommand{\abbrnote}[3]{%
  \ifclassroommode
    \ifshoworiginal
      \ifintranslation\else
        \footnote{\normalfont\footnotesize\abbrnotelabel #2\ifclassroomabbrzh；{\songfont #3}\fi}%
      \fi
    \fi
  \fi
}
\newcommand{\abbrnoteonce}[4]{%
  #2%
  \ifclassroommode
    \ifshoworiginal
      \ifintranslation\else
        \begingroup
          \edef\abbrcurrentabspage{\number\numexpr\value{abspage}+1\relax}%
          \ifcsname abbrnoteseen@#1@\abbrcurrentabspage\endcsname\else
            \global\expandafter\let\csname abbrnoteseen@#1@\abbrcurrentabspage\endcsname\empty
            \abbrnote{#1}{#3}{#4}%
          \fi
        \endgroup
      \fi
    \fi
  \fi
}
\newcommand{\abbrAbs}{\abbrnoteonce{abs}{Abs.}{Absatz}{項}}
\newcommand{\abbrAAO}{\abbrnoteonce{aao}{a.a.O.}{am angegebenen Ort}{前引處}}
\newcommand{\abbrAE}{\abbrnoteonce{ae}{AE}{Alternativ-Entwurf}{替代草案}}
\newcommand{\abbrArt}{\abbrnoteonce{art}{Art.}{Artikel}{條}}
\newcommand{\abbrAufl}{\abbrnoteonce{aufl}{Aufl.}{Auflage}{版}}
\newcommand{\abbrBd}{\abbrnoteonce{bd}{Bd.}{Band}{卷}}
\newcommand{\abbrBGBl}{\abbrnoteonce{bgbl}{BGBl.}{Bundesgesetzblatt}{聯邦法律公報}}
\newcommand{\abbrBGHSt}{\abbrnoteonce{bghst}{BGHSt}{Entscheidungen des Bundesgerichtshofes in Strafsachen}{聯邦普通法院刑事判決彙編}}
\newcommand{\abbrBRDrucks}{\abbrnoteonce{brdrucks}{BRDrucks.}{Bundesrats-Drucksache}{聯邦參議院文件}}
\newcommand{\abbrBTDrucks}{\abbrnoteonce{btdrucks}{BTDrucks.}{Bundestags-Drucksache}{聯邦議院文件}}
\newcommand{\abbrBundesgesetzbl}{\abbrnoteonce{bundesgesetzbl}{Bundesgesetzbl.}{Bundesgesetzblatt}{聯邦法律公報}}
\newcommand{\abbrBvF}{\abbrnoteonce{bvf}{BvF}{Bundesverfassungsgerichtsverfahren}{聯邦憲法法院程序案號}}
\newcommand{\abbrBvL}{\abbrnoteonce{bvl}{BvL}{Bundesverfassungsgerichtsverfahren}{聯邦憲法法院程序案號}}
\newcommand{\abbrBvR}{\abbrnoteonce{bvr}{BvR}{Bundesverfassungsgerichtsverfahren}{聯邦憲法法院程序案號}}
\newcommand{\abbrBVerfG}{\abbrnoteonce{bverfg}{BVerfG}{Bundesverfassungsgericht}{聯邦憲法法院}}
\newcommand{\abbrBVerfGE}{\abbrnoteonce{bverfge}{BVerfGE}{Entscheidungen des Bundesverfassungsgerichts}{聯邦憲法法院判決集}}
\newcommand{\abbrBVerfGG}{\abbrnoteonce{bverfgg}{BVerfGG}{Bundesverfassungsgerichtsgesetz}{聯邦憲法法院法}}
\newcommand{\abbrDDR}{\abbrnoteonce{ddr}{DDR}{Deutsche Demokratische Republik}{德意志民主共和國；東德}}
\newcommand{\abbrDJT}{\abbrnoteonce{djt}{DJT}{Deutscher Juristentag}{德國法學家大會}}
\newcommand{\abbrDr}{\abbrnoteonce{dr}{Dr.}{Doktor}{Dr.}}
\newcommand{\abbrEuGRZ}{\abbrnoteonce{eugrz}{EuGRZ}{Europäische Grundrechte-Zeitschrift}{歐洲基本權期刊}}
\newcommand{\abbrEV}{\abbrnoteonce{ev}{EV}{Einigungsvertrag}{統一條約}}
\newcommand{\abbrF}{\abbrnoteonce{f}{f.}{folgende}{以下一頁；下一段}}
\newcommand{\abbrFF}{\abbrnoteonce{ff}{ff.}{fortfolgende}{以下數頁；以下各段}}
\newcommand{\abbrGBlDDR}{\abbrnoteonce{gblddr}{GBl. DDR}{Gesetzblatt der Deutschen Demokratischen Republik}{德意志民主共和國法律公報}}
\newcommand{\abbrGG}{\abbrnoteonce{gg}{GG}{Grundgesetz}{基本法}}
\newcommand{\abbrGVBl}{\abbrnoteonce{gvbl}{GVBl.}{Gesetz- und Verordnungsblatt}{法律及命令公報}}
\newcommand{\abbrJR}{\abbrnoteonce{jr}{JR}{Juristische Rundschau}{法學評論}}
\newcommand{\abbrJZ}{\abbrnoteonce{jz}{JZ}{JuristenZeitung}{法學家報}}
\newcommand{\abbrIVm}{\abbrnoteonce{ivm}{i.V.m.}{in Verbindung mit}{結合；連同}}
\newcommand{\abbrMdB}{\abbrnoteonce{mdb}{MdB}{Mitglied des Bundestages}{聯邦議院議員}}
\newcommand{\abbrMwN}{\abbrnoteonce{mwn}{m.w.N.}{mit weiteren Nachweisen}{附進一步文獻／判例出處}}
\newcommand{\abbrNF}{\abbrnoteonce{nf}{n.F.}{neue Fassung}{新版本}}
\newcommand{\abbrAF}{\abbrnoteonce{af}{a.F.}{alte Fassung}{舊版本}}
\newcommand{\abbrNr}{\abbrnoteonce{nr}{Nr.}{Nummer}{款、目或號，依語境}}
\newcommand{\abbrProf}{\abbrnoteonce{prof}{Prof.}{Professor}{教授}}
\newcommand{\abbrRdnr}{\abbrnoteonce{rdnr}{Rdnr.}{Randnummer}{邊碼；段碼}}
\newcommand{\abbrRdnrn}{\abbrnoteonce{rdnrn}{Rdnrn.}{Randnummern}{邊碼；段碼}}
\newcommand{\abbrRGBl}{\abbrnoteonce{rgbl}{RGBl.}{Reichsgesetzblatt}{帝國法律公報}}
\newcommand{\abbrRGSt}{\abbrnoteonce{rgst}{RGSt}{Entscheidungen des Reichsgerichts in Strafsachen}{帝國法院刑事判決彙編}}
\newcommand{\abbrRspr}{\abbrnoteonce{rspr}{Rspr.}{Rechtsprechung}{判例；司法實務}}
\newcommand{\abbrRVO}{\abbrnoteonce{rvo}{RVO}{Reichsversicherungsordnung}{帝國保險法}}
\newcommand{\abbrS}{\abbrnoteonce{s}{S.}{Seite}{頁}}
\newcommand{\abbrSFHG}{\abbrnoteonce{sfhg}{SFHG}{Schwangeren- und Familienhilfegesetz}{孕婦及家庭扶助法}}
\newcommand{\abbrSGBV}{\abbrnoteonce{sgbv}{SGB V}{Fünftes Buch Sozialgesetzbuch}{社會法典第五編}}
\newcommand{\abbrStAG}{\abbrnoteonce{staeg}{StÄG}{Strafrechtsänderungsgesetz}{刑法修正法}}
\newcommand{\abbrStenBer}{\abbrnoteonce{stenber}{StenBer.}{Stenographischer Bericht}{速記紀錄}}
\newcommand{\abbrStGB}{\abbrnoteonce{stgb}{StGB}{Strafgesetzbuch}{刑法}}
\newcommand{\abbrStREG}{\abbrnoteonce{streg}{StREG}{Strafrechtsreform-Ergänzungsgesetz}{刑法改革補充法}}
\newcommand{\abbrStrRG}{\abbrnoteonce{strrg}{StrRG}{Strafrechtsreformgesetz}{刑法改革法}}
\newcommand{\abbrUS}{\abbrnoteonce{us}{U.S.}{United States Reports}{美國最高法院判決彙編}}
\newcommand{\abbrVgl}{\abbrnoteonce{vgl}{vgl.}{vergleiche}{參見}}
\newcommand{\abbrVVDStRL}{\abbrnoteonce{vvdstrl}{VVDStRL}{Veröffentlichungen der Vereinigung der Deutschen Staatsrechtslehrer}{德國公法教師協會論文集}}
\newcommand{\abbrWP}{\abbrnoteonce{wp}{WP.}{Wahlperiode}{屆期}}
\newcommand{\abbrWp}{\abbrnoteonce{wp}{Wp.}{Wahlperiode}{屆期}}
\newcommand{\abbrZStrW}{\abbrnoteonce{zstrw}{ZStrW}{Zeitschrift für die gesamte Strafrechtswissenschaft}{整體刑法學期刊}}
\newcommand{\recordsourcerandstart}{%
  \ifrandnumbers
    \setcounter{lastsourcerandstart}{\value{randnumber}}%
    \global\lastsourcehasrandnumberstrue
  \else
    \global\lastsourcehasrandnumbersfalse
  \fi
}
\newlength{\biparagraphskip}
\setlength{\biparagraphskip}{0.52em}
\newlength{\lawboxblockbeforeskip}
\setlength{\lawboxblockbeforeskip}{0.72em}
\newlength{\lawboxblockafterskip}
\setlength{\lawboxblockafterskip}{0.92em}
\newif\iflawboxpartendedblock
\lawboxpartendedblockfalse
\newcommand{\sourceparstyle}{\footnotesize\color{sourceink}\setstretch{1.12}\RaggedRight\emergencystretch=3em\setlength{\parskip}{0.42em}\obeylines\selectfont}
\newcommand{\transparstyle}{\songfont\color{caseink}\setstretch{1.12}\setlength{\parindent}{0pt}\setlength{\parskip}{0.42em}}
\newcommand{\bilingualpairskip}{%
  \par\addvspace{\biparagraphskip}%
  \switchcolumn
  \par\addvspace{\biparagraphskip}%
  \switchcolumn*%
  \global\intranslationfalse
}
\newcommand{\bilingualboxpairskip}{%
  \par
  \switchcolumn
  \par
  \switchcolumn*%
  \global\intranslationfalse
}
\newcommand{\bilinguallawboxblockskip}{%
  \switchcolumn*%
  \par\addvspace{\lawboxblockafterskip}%
  \switchcolumn
  \par\addvspace{\lawboxblockafterskip}%
  \switchcolumn*%
  \global\intranslationfalse
}
\newcommand{\bikeepwithnext}[1][4]{%
  \ifbilingualcolumns
    \syncleft
    \Needspace{#1\baselineskip}%
    \switchcolumn
    \Needspace{#1\baselineskip}%
    \switchcolumn*%
    \global\intranslationfalse
  \else
    \Needspace{#1\baselineskip}%
  \fi
}
\NewEnviron{sourcepara}[1][]{
  \recordsourcerandstart
  \ifshoworiginal
    \ifbilingualcolumns
      \ifintranslation\switchcolumn*\global\intranslationfalse\fi
    \else
      \Needspace{5\baselineskip}
    \fi
    \ifbilingualcolumns\else\par\addvspace{0.35em}\fi
    {\sourceparstyle
    \noindent\BODY\par}
    \ifbilingualcolumns\else\addvspace{0.25em}\fi
    \ifbilingualcolumns
      \switchcolumn
      \global\intranslationtrue
    \fi
  \else
    \ifrandnumbers
      \begingroup
        \setbox0=\vbox{\hsize=\linewidth\sourceparstyle\noindent\BODY\par}%
      \endgroup
    \fi
  \fi
}
\NewEnviron{transpara}{
  \ifshowtranslation
    \setcounter{randnumberlocal}{\value{lastsourcerandstart}}%
    \ifbilingualcolumns
      \ifintranslation\else\switchcolumn\global\intranslationtrue\fi
      {\transparstyle\setlength{\leftskip}{0.55em}\setlength{\rightskip}{0.15em plus 1em}\addtolength{\linewidth}{-0.55em}\BODY\par}
      \switchcolumn*
      \bilingualpairskip
    \else
      {\transparstyle\BODY\par}
    \fi
  \else
    \ifbilingualcolumns
      \ifintranslation\switchcolumn*\global\intranslationfalse\fi
    \fi
  \fi
}
\NewEnviron{sourceboxpara}[1][]{
  \global\lawboxpartendedblockfalse
  \recordsourcerandstart
  \ifshoworiginal
    \ifbilingualcolumns
      \ifintranslation\switchcolumn*\global\intranslationfalse\fi
    \else
      \Needspace{3\baselineskip}
    \fi
    {\sourceparstyle
    \BODY\par}
    \ifbilingualcolumns\else
      \iflawboxpartendedblock\addvspace{\lawboxblockafterskip}\fi
    \fi
    \ifbilingualcolumns
      \switchcolumn
      \global\intranslationtrue
    \fi
  \fi
}
\NewEnviron{transboxpara}{
  \global\lawboxpartendedblockfalse
  \ifshowtranslation
    \setcounter{randnumberlocal}{\value{lastsourcerandstart}}%
    \ifbilingualcolumns
      \ifintranslation\else\switchcolumn\global\intranslationtrue\fi
      {\transparstyle\BODY\par}
      \iflawboxpartendedblock
        \bilinguallawboxblockskip
      \else
        \switchcolumn*
        \bilingualboxpairskip
      \fi
    \else
      {\transparstyle\BODY\par}
      \iflawboxpartendedblock\addvspace{\lawboxblockafterskip}\fi
    \fi
  \else
    \ifbilingualcolumns
      \ifintranslation\switchcolumn*\global\intranslationfalse\fi
    \fi
  \fi
}
\newcommand{\leitsatznum}[1]{\noindent\hangindent=3.0em\hangafter=1\makebox[2.25em][r]{\bfseries #1.}\hspace{0.55em}\ignorespaces}
\newcommand{\syncleft}{\ifbilingualcolumns\ifintranslation\switchcolumn*\global\intranslationfalse\fi\fi}
\pretocmd{\section}{\syncleft\clearpage}{}{}
\pretocmd{\subsection}{\syncleft}{}{}
\pretocmd{\subsubsection}{\syncleft}{}{}
\newcommand{\biheadingfont}{\songfont}
\newcommand{\bitocentry}[2]{%
  \ifshoworiginal
    \ifshowtranslation
      \ifstrequal{#1}{#2}{#1}{#1\space #2}%
    \else
      #1%
    \fi
  \else
    #2%
  \fi
}
\pdfstringdefDisableCommands{%
  \def\bitocentry#1#2{#1}%
}
\newcommand{\biheading}[7]{%
  \syncleft#1\bikeepwithnext[#3]\phantomsection\addcontentsline{toc}{#2}{\bitocentry{#6}{#7}}%
  \ifbilingualcolumns
    {#4 #6\par}%
    \switchcolumn
    {#4\biheadingfont #7\par}%
    \switchcolumn*%
    \global\intranslationfalse
    \par\addvspace{#5}%
    \switchcolumn
    \par\addvspace{#5}%
    \switchcolumn*%
  \else
    \ifshoworiginal{#4 #6\par}\fi
    \ifshowtranslation{#4\biheadingfont #7\par}\fi
    \addvspace{#5}%
  \fi
}
\newcommand{\termsectionheading}[1]{%
  \par\Needspace{9\baselineskip}%
  \vspace{0.65em}%
  {\songfont\bfseries #1\par}%
  \vspace{0.2em}%
}
% 章節換頁：
%   雙語 paracol 欄內不要用 \clearpage；它會在同步兩欄時吐出只有欄線/頁碼的空頁。
%   \newpage 足以讓下一個 section 起新頁，又不會強制清空所有待排材料。
%   單語模式不在 paracol 內，仍可用 \clearpage。
\newcommand{\bisectionpagebreak}{\ifbilingualcolumns\newpage\else\clearpage\fi}
\newcommand{\bisection}[2]{%
  \biheading{\iffirstbilingualsection\global\firstbilingualsectionfalse\else\bisectionpagebreak\fi}{section}{6}{\centering\Large\bfseries}{0.8em}{#1}{#2}%
}
\newcommand{\bisubsection}[2]{%
  \biheading{}{subsection}{5}{\large\bfseries}{0.55em}{#1}{#2}%
}
\newcommand{\bisubsubsection}[2]{%
  \biheading{}{subsubsection}{4}{\normalsize\bfseries}{0.45em}{#1}{#2}%
}
\newcommand{\bicompositeheading}[4]{%
  \biheading{}{subsection}{5}{\large\bfseries}{0.16em}{#1}{#2}%
  \biheading{}{subsubsection}{3}{\normalsize\bfseries}{0.45em}{#3}{#4}%
}
\newif\iflawboxfirstitem
\lawboxfirstitemfalse
\newcommand{\lawboxprespace}[1]{%
  \par
  \iflawboxfirstitem
    \global\lawboxfirstitemfalse
  \else
    \addvspace{#1}%
  \fi
}
\newtcolorbox{lawboxinner}{
  caseboxbase,
  colframe=sourceline!85!black,
  colback=white,
  left=1.05em,
  right=1.05em,
  top=0.62em,
  bottom=0.62em,
  before upper={\global\inlawboxtrue\global\lawboxfirstitemtrue\global\lawboxhaspendingrnfalse\ifintranslation\lawboxtransfont\else\lawboxsourcefont\fi\RaggedRight\setlength{\parindent}{0pt}\setlength{\parskip}{0pt}\ifintranslation\setstretch{\lawboxtransstretch}\else\setstretch{\lawboxsourcestretch}\fi},
  after upper={\global\lawboxhaspendingrnfalse\global\inlawboxfalse}
}
\tcbset{
  lawboxpartbase/.style={
    enhanced,
    breakable=false,
    width=0.985\linewidth,
    colback=white,
    frame hidden,
    boxrule=0pt,
    arc=0pt,
    left=1.05em,
    right=1.05em,
    boxsep=1.5mm,
    before skip=0pt,
    after skip=0pt,
    before upper={\global\inlawboxtrue\global\lawboxfirstitemtrue\global\lawboxhaspendingrnfalse\ifintranslation\lawboxtransfont\else\lawboxsourcefont\fi\RaggedRight\setlength{\parindent}{0pt}\setlength{\parskip}{0pt}\ifintranslation\setstretch{\lawboxtransstretch}\else\setstretch{\lawboxsourcestretch}\fi},
    after upper={\global\lawboxhaspendingrnfalse\global\inlawboxfalse},
    borderline west={0.28pt}{0pt}{sourceline!85!black},
    borderline east={0.28pt}{0pt}{sourceline!85!black},
  },
  lawboxpartfirst/.style={
    lawboxpartbase,
    top=0.62em,
    bottom=0.04em,
    before skip=\lawboxblockbeforeskip,
    borderline north={0.28pt}{0pt}{sourceline!85!black},
  },
  lawboxpartmiddle/.style={
    lawboxpartbase,
    top=0.04em,
    bottom=0.04em,
  },
  lawboxpartlast/.style={
    lawboxpartbase,
    top=0.04em,
    bottom=0.62em,
    after skip=0pt,
    borderline south={0.28pt}{0pt}{sourceline!85!black},
  }
}
\newtcolorbox{lawboxpartinner}[1]{#1}
\newtcolorbox{termboxinner}{
  caseboxbase,
  colframe=noteblue!70!black,
  colback=casepaper,
  before upper={\songfont\setlength{\parindent}{0pt}\setlength{\parskip}{0.35em}}
}
\newtcolorbox{deboxinner}{
  caseboxbase,
  colframe=notegray!72!black,
  colback=white,
  left=1.05em,
  right=1.05em,
  top=0.62em,
  bottom=0.62em,
  before upper={\global\inlawboxtrue\global\lawboxfirstitemtrue\global\lawboxhaspendingrnfalse\small\setlength{\parindent}{0pt}\setlength{\parskip}{0.35em}},
  after upper={\global\lawboxhaspendingrnfalse\global\inlawboxfalse}
}
\NewEnviron{lawbox}{
  \begin{lawboxinner}\BODY\end{lawboxinner}
}
\NewEnviron{lawboxpart}[2]{
  \begin{lawboxpartinner}{lawboxpart#1,equal height group=#2}\BODY\end{lawboxpartinner}%
  \ifstrequal{#1}{last}{\global\lawboxpartendedblocktrue}{}%
}
\NewEnviron{termbox}{
  \par\addvspace{0.55em}
  {\color{noteblue!70!black}\rule{\linewidth}{0.28pt}}\par
  \begingroup
    \songfont\small\color{caseink}
    \setlength{\parindent}{0pt}
    \setlength{\parskip}{0.24em}
    \BODY
    \par
  \endgroup
  {\color{noteblue!70!black}\rule{\linewidth}{0.28pt}}\par
  \addvspace{0.55em}
}
\NewEnviron{debox}{
  \begin{deboxinner}\BODY\end{deboxinner}
}
\providecommand{\paragraphmark}[1]{}
\newcommand{\lawboxtitleafter}{\addvspace{0.06em}}
\newcommand{\lawtitle}[1]{\lawboxprespace{0.22em}{\lawboxtransfont\bfseries\lawboxuserandnumber #1}\par\lawboxtitleafter}
\newcommand{\absno}[2]{\lawboxprespace{0.04em}\noindent\lawboxuserandnumber\hangindent=2.6em\hangafter=1\makebox[2.6em][l]{(#1)}#2\par}
\newcommand{\nrno}[2]{\lawboxprespace{0pt}\noindent\lawboxuserandnumber\hspace*{2.6em}\hangindent=4.4em\hangafter=1\makebox[1.8em][l]{#1.}#2\par}
\newcommand{\letterno}[2]{\lawboxprespace{0pt}\noindent\lawboxuserandnumber\hspace*{4.4em}\hangindent=6.0em\hangafter=1\makebox[1.6em][l]{#1)}#2\par}
\newcommand{\sourcelawtitle}[1]{\lawboxprespace{0.22em}{\lawboxsourcefont\bfseries\lawboxuserandnumber #1}\par\lawboxtitleafter}
\newcommand{\sourcelawsubtitle}[1]{\lawboxprespace{0.04em}{\lawboxsourcefont\bfseries\lawboxuserandnumber #1}\par\lawboxtitleafter}
\newcommand{\sourceabsno}[2]{\lawboxprespace{0.04em}\noindent\lawboxuserandnumber\hangindent=2.45em\hangafter=1\makebox[2.45em][l]{(#1)}#2\par}
\newcommand{\sourcenrno}[2]{\lawboxprespace{0pt}\noindent\lawboxuserandnumber\hspace*{2.45em}\hangindent=4.25em\hangafter=1\makebox[1.8em][l]{#1.}#2\par}
\newcommand{\sourceletterno}[2]{\lawboxprespace{0pt}\noindent\lawboxuserandnumber\hspace*{4.25em}\hangindent=5.85em\hangafter=1\makebox[1.6em][l]{#1)}#2\par}
\newcommand{\sourcequotepara}[1]{\lawboxprespace{0.04em}\noindent\lawboxuserandnumber\hangindent=1.2em\hangafter=1 #1\par}
\newcommand{\sourcelawline}[1]{\lawboxprespace{0.04em}\noindent\lawboxuserandnumber #1\par}
\newcommand{\lawline}[1]{\lawboxprespace{0.04em}\noindent\lawboxuserandnumber #1\par}
\newcommand{\pendingtranslation}{\par{\songfont\footnotesize\color{pendingink}（中文譯文待補。）}\par}
\newcommand{\tocpart}[1]{\par\addvspace{0.52em}{\footnotesize\bfseries\color{pendingink}#1\par}\addvspace{0.16em}}
\newcommand{\tocdashline}{\par\addvspace{0.48em}\noindent{\color{notegray}\xleaders\hbox to 0.82em{\hfil\rule{0.42em}{0.28pt}\hfil}\hskip\linewidth}\par\addvspace{0.38em}}
% \tocpanel{font-cmd}{content}：將目錄置中於所在欄寬（雙語欄適用）。
% A3 橫式使用 0.58\linewidth，A4 橫式使用 0.84\linewidth，
% 使兩種紙張下目錄欄內容實際寬度相近（均約 100–105 mm）。
\newcommand{\tocpanel}[2]{%
  \makebox[\linewidth][c]{%
    \ifx\bilingualpaper\bilingualpaperaThree
      \begin{minipage}{0.58\linewidth}%
    \else
      \begin{minipage}{0.84\linewidth}%
    \fi
    \toccolstyle
    #1#2%
    \end{minipage}%
  }%
}
% ===== 首頁簡目錄的兩層 description list 縮排 =====
% enumitem 的 description list 可把每一列拆成：
%   [label]  labelsep  body text
% 這裡的「起點」都以所在 minipage/欄位的左緣為 0。
%
% 核心規則：
%   labelindent = label 欄位從左緣開始的位置。
%   labelwidth  = label 欄位本身寬度；標籤太長（如 Abw. Meinung）就加大它。
%   labelsep    = label 欄位右緣到正文的距離。
%   leftmargin  = 正文 body text 的起點。判斷層級視覺時主要看這個。
%
% 所以：
%   想讓「Begründung / 判決理由」整列正文往右或往左：改 tocitems 的 leftmargin。
%   想讓「A. Verfahren und Gegenstand / A. 程序與審查標的」正文往右或往左：
%     改 tocsubitems 的 leftmargin。
%   想讓 A./B./C. 這些標籤本身往右或往左，但正文不動：改 tocsubitems 的 labelindent。
%   想讓 A./B./C. 與正文間距變大或變小：改 tocsubitems 的 labelsep。
%
% 層級原則：
%   tocsubitems.leftmargin 必須大於 tocitems.leftmargin，
%   否則子層級正文會看起來比「Gründe / 理由」還靠左。
\newlist{tocitems}{description}{1}
\setlist[tocitems]{%
  labelindent=0pt,    % 第一層 label 欄位起點；0pt 表示貼齊目錄欄左緣。
  leftmargin=2.54cm,  % 第一層正文起點；例如 Begründung / 判決理由 的左緣。
  labelwidth=2.16cm,  % 第一層 label 欄位寬度；需容納 Leitsätze、Abw. Meinung、不同意見等。
  labelsep=0.32cm,    % 第一層 label 與正文的水平間距；只改間距，不改正文起點。
  itemsep=0.28em,     % 第一層各項之間的垂直距離。
  topsep=0.20em,      % 第一層 list 上下與前後內容的垂直距離。
  parsep=0pt,         % 同一 item 內段落之間的距離；目錄項通常不需要。
  partopsep=0pt,      % list 若緊接段落開始時的額外距離；關閉以免目錄高度飄動。
  align=left,         % label 欄內左對齊；長短標籤不靠右貼齊。
  font=\normalfont\bfseries % label 的字型；正文不受這行影響。
}
\newlist{tocsubitems}{description}{1}
\setlist[tocsubitems]{%
  labelindent=1cm,    % 第二層 label 起點；只控制 A./B./C. 本身的位置。
  leftmargin=3.00cm,  % 第二層正文起點；必須大於 2.54cm，才會比 Begründung / 判決理由更右。
  labelwidth=0.5cm,   % 第二層 label 欄寬；A./B./C. 很短，可小於第一層。
  labelsep=1.5cm,    % 第二層 label 與正文間距；A. 與 Verfahren / 程序 的距離看這裡。
  itemsep=0.12em,     % 第二層各項較密，避免 A.–E. 佔太高。
  topsep=0.04em,      % 第二層 list 與 Gründe / 理由之間的垂直距離。
  parsep=0pt,         % 同一第二層 item 內段落距離；目錄項通常不需要。
  partopsep=0pt,      % 關閉額外段落起始距離，避免版面不穩。
  align=left,         % A./B./C. label 欄內左對齊。
  font=\normalfont\bfseries, % A./B./C. label 加粗；正文不受這行影響。
  before={}           % 不在第二層 list 前自動插入額外內容。
}
\newlist{termitems}{description}{1}
\ifbilingualcolumns
  \setlist[termitems]{%
    leftmargin=5.2cm,
    labelwidth=4.8cm,
    labelsep=0.35cm,
    itemsep=0.16em,
    topsep=0.2em,
    parsep=0pt,
    partopsep=0pt,
    font=\normalfont\bfseries,
    style=nextline
  }
\else
  \setlist[termitems]{%
    leftmargin=0pt,
    labelwidth=0pt,
    labelsep=0pt,
    itemsep=0.24em,
    topsep=0.2em,
    parsep=0pt,
    partopsep=0pt,
    font=\normalfont\bfseries,
    style=nextline
  }
\fi
\newlist{abbritems}{description}{1}
\setlist[abbritems]{%
  leftmargin=2.38cm,
  labelwidth=2.02cm,
  labelsep=0.28cm,
  itemsep=0.16em,
  topsep=0.2em,
  parsep=0pt,
  partopsep=0pt,
  align=left,
  font=\normalfont\bfseries
}
\ifbilingualcolumns
  \newenvironment{termcolumns}{\begin{multicols}{2}}{\end{multicols}}
\else
  \newenvironment{termcolumns}{}{}
\fi

% ===== 編者註腳開關 =====
% 一般版預設不顯示編者註，避免正文腳註過密。
% 若開啟 \classroommodeon，GG/條文編者註仍會自動顯示，無需另開本開關。
% 若一般版也要顯示編者註，可在單案檔 \input 後加 \showeditornotestrue。
\newif\ifshoweditornotes
\showeditornotesfalse

% ===== 目錄排版輔助 =====
\newcommand{\toccolstyle}{\raggedright\arraybackslash}
\newcommand{\tocsingleindent}{\leftskip=1.45cm\rightskip=1.2cm}

% ===== 行內引文環境（德文原文與中文譯文各一，帶左側灰色邊線）=====
\newcommand{\quoteparbreak}{\par\addvspace{0.48em}}
\newcommand{\sourcequoteinner}[1]{%
  \begin{tcolorbox}[
    enhanced, breakable, width=\dimexpr\linewidth-0.8em\relax, frame hidden,
    colback=white, coltext=caseink, boxrule=0pt,
    borderline west={0.55pt}{0.88em}{notegray},
    left=1.78em, right=0.72em, top=0.18em, bottom=0.18em,
    boxsep=0pt, before skip=0.24em, after skip=0.24em,
    before upper={\normalfont\footnotesize\color{caseink}\setlength{\parindent}{0pt}\setlength{\parskip}{0.34em}\raggedright\setlength{\rightskip}{0pt plus 1em}}
  ]%
  #1\par
  \end{tcolorbox}%
}
\newcommand{\transquoteinner}[1]{%
  \begin{tcolorbox}[
    enhanced, breakable, width=\dimexpr\linewidth-0.8em\relax, frame hidden,
    colback=white, coltext=caseink, boxrule=0pt,
    borderline west={0.55pt}{0.88em}{notegray},
    left=1.78em, right=0.72em, top=0.18em, bottom=0.18em,
    boxsep=0pt, before skip=0.24em, after skip=0.24em,
    before upper={\songfont\small\color{caseink}\setlength{\parindent}{0pt}\setlength{\parskip}{0.34em}\raggedright\setlength{\rightskip}{0pt plus 1em}}
  ]%
  #1\par
  \end{tcolorbox}%
}

% ===== 大綱（Gliederung）Rn 輔助命令 =====
\newcommand{\outrn}[2]{\enspace{\normalfont\footnotesize\color{pendingink}(Rn\,#1–#2)}}
\newcommand{\outrnsingle}[1]{\enspace{\normalfont\footnotesize\color{pendingink}(Rn\,#1)}}
\newcommand{\outrnzh}[2]{\enspace{\normalfont\footnotesize\color{pendingink}（第\,#1–#2\,段）}}
\newcommand{\outrnzhs}[1]{\enspace{\normalfont\footnotesize\color{pendingink}（第\,#1\,段）}}
% ===== 大綱雙語對照表格輔助命令（三欄：段碼 │ 德文 │ 中文）=====
% 欄 1：段碼（由欄定義自動格式化：小字、等寬、右對齊、pendingink）
% 欄 2：德文標籤＋正文
% 欄 3：中文標籤＋正文
%
% 無段碼的列（\omainrow、\osecrow），欄 1 以 & 空置。
% 有段碼的列（\osecrowrn、\osubrow、\ointrow），#1 為預格式化字串
%   例：\osubrow{2–49}{I.}{...}{I.}{...}
%       \ointrow{107}{...}{...}
%
% \opartrow：段組標題（橫跨三欄）
\newcommand{\opartrow}[2]{%
  \noalign{\vspace{0.30em}}%
  \multicolumn{3}{@{}l@{}}{%
    \footnotesize\bfseries\color{pendingink}#1\hfill\songfont#2%
  }\\[0.06em]%
}
% \odashrow：虛線分隔（橫跨三欄）
\newcommand{\odashrow}{%
  \noalign{\vspace{0.28em}}%
  \multicolumn{3}{@{}l@{}}{\color{notegray}%
    \xleaders\hbox to 0.82em{\hfil\rule{0.42em}{0.28pt}\hfil}\hskip 0.88\linewidth}%
  \\[0.24em]%
}
% \odissentpart：不同意見書區段標頭。比一般 \opartrow 更像附錄性轉場。
\newcommand{\odissentpart}[2]{%
  \noalign{\vspace{0.42em}}%
  \multicolumn{3}{@{}p{0.88\textwidth}@{}}{%
    \color{notegray}\rule{\linewidth}{0.18pt}%
  }\\[-0.18em]%
  \multicolumn{3}{@{}p{0.88\textwidth}@{}}{%
    \makebox[\linewidth][s]{%
      \footnotesize\bfseries\color{pendingink}#1%
      \hfill{\songfont #2}%
    }%
  }\\[-0.02em]%
}
% \odissentopinion：不同意見書的署名與一行摘要，右欄保留段碼範圍。
\newcommand{\odissentopinion}[5]{%
  \footnotesize\hangindent=0.52cm\hangafter=1%
    {\bfseries\color{caseink}#2}\enspace{\color{notegray}\textperiodcentered}\enspace\color{caseink}#3%
  &\footnotesize\hangindent=0.52cm\hangafter=1%
    {\songfont\bfseries\color{caseink}#4}\enspace{\color{notegray}\textperiodcentered}\enspace\songfont\color{caseink}#5%
  &#1%
  \\[0.12em]%
}
% \omainrow：頂層項目，無段碼（Leitsätze、Urteil、Gründe …）
\newcommand{\omainrow}[4]{%
  \footnotesize\hangindent=1.92cm\hangafter=1%
    \makebox[1.92cm][l]{\bfseries\color{caseink}#1}\color{caseink}#2%
  &\footnotesize\hangindent=1.92cm\hangafter=1%
    \makebox[1.92cm][l]{\bfseries\color{caseink}#3}\songfont\color{caseink}#4%
  &%
  \\[0.15em]%
}
% \omainrowrn：頂層項目，有段碼（不同意見等標籤寬於 \osecrow 框者）
\newcommand{\omainrowrn}[5]{%
  \footnotesize\hangindent=1.92cm\hangafter=1%
    \makebox[1.92cm][l]{\bfseries\color{caseink}#2}\color{caseink}#3%
  &\footnotesize\hangindent=1.92cm\hangafter=1%
    \makebox[1.92cm][l]{\bfseries\color{caseink}#4}\songfont\color{caseink}#5%
  &#1%
  \\[0.15em]%
}
% \osecrow：一級子項，無段碼（A.、C.、D. 等有子項者）
\newcommand{\osecrow}[4]{%
  \footnotesize\hspace{1.10cm}\hangindent=1.98cm\hangafter=1%
    \makebox[0.86cm][l]{\bfseries\color{caseink}#1}\color{caseink}#2%
  &\footnotesize\hspace{1.10cm}\hangindent=1.98cm\hangafter=1%
    \makebox[0.86cm][l]{\bfseries\color{caseink}#3}\songfont\color{caseink}#4%
  &%
  \\[0.11em]%
}
% \osecrowrn：一級子項，有段碼（B.、E. 等完整段落範圍者）
% #1=段碼字串（如 101–106）, #2=DE標籤, #3=DE正文, #4=ZH標籤, #5=ZH正文
\newcommand{\osecrowrn}[5]{%
  \footnotesize\hspace{1.10cm}\hangindent=1.98cm\hangafter=1%
    \makebox[0.86cm][l]{\bfseries\color{caseink}#2}\color{caseink}#3%
  &\footnotesize\hspace{1.10cm}\hangindent=1.98cm\hangafter=1%
    \makebox[0.86cm][l]{\bfseries\color{caseink}#4}\songfont\color{caseink}#5%
  &#1%
  \\[0.11em]%
}
% \osubrow：二級子項，有段碼範圍（I.、II.、A.I. …）
% #1=段碼字串, #2=DE標籤, #3=DE正文, #4=ZH標籤, #5=ZH正文
\newcommand{\osubrow}[5]{%
  \footnotesize\hspace{1.40cm}\hangindent=2.30cm\hangafter=1%
    \makebox[0.90cm][l]{\bfseries\color{caseink}#2}\color{caseink}#3%
  &\footnotesize\hspace{1.40cm}\hangindent=2.30cm\hangafter=1%
    \makebox[0.90cm][l]{\bfseries\color{caseink}#4}\songfont\color{caseink}#5%
  &#1%
  \\[0.09em]%
}
% \ointrow：引入段落，有單一段碼，無標籤
% #1=段碼字串, #2=DE正文, #3=ZH正文
\newcommand{\ointrow}[3]{%
  \footnotesize\hspace{0.52cm}\hangindent=0.52cm\hangafter=1\color{caseink}#2%
  &\footnotesize\hspace{0.52cm}\hangindent=0.52cm\hangafter=1\songfont\color{caseink}#3%
  &#1%
  \\[0.09em]%
}
% \osubsubrow：三級子項（1.、2.、3.）
% #1=段碼字串, #2=DE標籤, #3=DE正文, #4=ZH標籤, #5=ZH正文
\newcommand{\osubsubrow}[5]{%
  \footnotesize\hspace{1.80cm}\hangindent=2.48cm\hangafter=1%
    \makebox[0.68cm][l]{\bfseries\color{caseink}#2}\color{caseink}#3%
  &\footnotesize\hspace{1.80cm}\hangindent=2.48cm\hangafter=1%
    \makebox[0.68cm][l]{\bfseries\color{caseink}#4}\songfont\color{caseink}#5%
  &#1%
  \\[0.08em]%
}
% \osubsubsubrow：四級子項（a）、b））
% #1=段碼字串, #2=DE標籤, #3=DE正文, #4=ZH標籤, #5=ZH正文
\newcommand{\osubsubsubrow}[5]{%
  \footnotesize\hspace{2.10cm}\hangindent=2.68cm\hangafter=1%
    \makebox[0.58cm][l]{\bfseries\color{caseink}#2}\color{caseink}#3%
  &\footnotesize\hspace{2.10cm}\hangindent=2.68cm\hangafter=1%
    \makebox[0.58cm][l]{\bfseries\color{caseink}#4}\songfont\color{caseink}#5%
  &#1%
  \\[0.07em]%
}
% \osubsubsubsubrow：五級子項（aa）、bb））
% 標籤框 0.62cm：足以容納「aa)」在 footnotesize bold 下（約 0.50cm）並留有間距
% #1=段碼字串, #2=DE標籤, #3=DE正文, #4=ZH標籤, #5=ZH正文
\newcommand{\osubsubsubsubrow}[5]{%
  \footnotesize\hspace{2.38cm}\hangindent=3.06cm\hangafter=1%
    \makebox[0.68cm][l]{\bfseries\color{caseink}#2}\color{caseink}#3%
  &\footnotesize\hspace{2.38cm}\hangindent=3.06cm\hangafter=1%
    \makebox[0.68cm][l]{\bfseries\color{caseink}#4}\songfont\color{caseink}#5%
  &#1%
  \\[0.06em]%
}
% \outlinepage：雙語對照大綱頁包裝環境（longtable 自動跨頁）
% 三欄：德文欄 + 中文欄 + 段碼欄（最右，不折行）
% 段碼欄寬度：1.80cm，足以容納最長段碼如「309–348」
% DE/ZH 欄寬：(0.87\textwidth - 2.08cm) / 2
%   驗算：2×欄 + 0.05\textwidth + 0.28cm gap + 1.80cm = 0.87tw - 2.08cm + 0.05tw + 2.08cm = 0.92tw ✓
\newenvironment{outlinepage}{%
  \vspace*{0.40cm}%
  \setlength{\LTleft}{0.04\textwidth}%
  \setlength{\LTright}{0.04\textwidth}%
  \setlength{\tabcolsep}{0pt}%
  % 大綱列距由 longtable 的 \arraystretch 與各 row macro 結尾的 \\[...em] 共同決定。
  % 這裡控制「每一列本身的表格 strut 高度」；數值越大，A./I./1. 各項之間越像有空行。
  % A3 原本用 0.62，視覺上沒有空行；A4 也沿用同值，避免切到 A4 後項目被撐開。
  % 若日後覺得大綱太擠，只微調這個值（例如 0.68），不要改各 \omainrow/\osubrow 的縮排。
  \renewcommand{\arraystretch}{0.62}%
  \begin{longtable}{%
    >{\vspace{0pt}\raggedright\arraybackslash}p{\dimexpr(0.87\textwidth-2.08cm)/2\relax}%
    @{\hspace{0.025\textwidth}}%
    !{\color{notegray}\vrule width 0.12pt}%
    @{\hspace{0.025\textwidth}}%
    >{\vspace{0pt}\raggedright\arraybackslash}p{\dimexpr(0.87\textwidth-2.08cm)/2\relax}%
    @{\hspace{0.28cm}}%
    >{\vspace{0pt}\footnotesize\normalfont\color{pendingink}\raggedright\arraybackslash}p{1.80cm}%
    @{}%
  }%
  % 首頁欄標籤：不要橫跨置中；分別放在德文欄與中文欄的實際欄位起點。
  \footnotesize\bfseries\color{sourceink}Gliederung%
  &\footnotesize\bfseries\songfont\color{sourceink}大綱%
  &%
  \\[0.36em]%
  \endfirsthead%
  % 續頁欄標籤：同樣分欄定位，以灰色細體區分；無需標注「續」。
  \footnotesize\color{pendingink}Gliederung%
  &\footnotesize\songfont\color{pendingink}大綱%
  &%
  \\[0.24em]%
  \endhead%
}{%
  \end{longtable}%
}
 之前設定；
%   預設 chinese / a4）
\ifdefined\docmode\else\def\docmode{chinese}\fi
\ifdefined\bilingualpaper\else\def\bilingualpaper{a4}\fi
\newif\ifshoworiginal
\newif\ifshowtranslation
\newif\ifbilinguallandscape
\newif\ifbilingualcolumns
\newif\ifintranslation
\newif\iffirstbilingualsection

\showoriginalfalse
\showtranslationfalse
\bilinguallandscapefalse
\bilingualcolumnsfalse
\intranslationfalse
\firstbilingualsectionfalse

\def\modechinese{chinese}
\def\modegerman{german}
\def\modebilingual{bilingual}
\def\bilingualpaperaThree{a3}
\def\bilingualpaperaFour{a4}

\ifx\docmode\modechinese
  \showtranslationtrue
\fi

\ifx\docmode\modegerman
  \showoriginaltrue
\fi

\ifx\docmode\modebilingual
  \showoriginaltrue
  \showtranslationtrue
  \bilinguallandscapetrue
  \bilingualcolumnstrue
\fi

% 編排模式說明：
% chinese  → \showoriginalfalse  \showtranslationtrue  （A4 直式，純中文）
% german   → \showoriginaltrue   \showtranslationfalse （A4 直式，純德文）
% bilingual→ \showoriginaltrue   \showtranslationtrue  \bilingualcolumnstrue（A3/A4 橫式對照）

\ifbilingualcolumns
  \ifx\bilingualpaper\bilingualpaperaFour
    \geometry{
      a4paper,
      landscape,
      left=1.25cm,
      right=2.25cm,
      top=1.25cm,
      bottom=1.75cm,
      footskip=8.5mm,
      marginparwidth=0.75cm,
      marginparsep=0.25cm
    }
  \else
    \geometry{
      a3paper,
      landscape,
      left=1.55cm,
      right=2.75cm,
      top=1.55cm,
      bottom=2.05cm,
      footskip=9.5mm,
      marginparwidth=0.9cm,
      marginparsep=0.35cm
    }
  \fi
\else
\geometry{
  a4paper,
  portrait,
  left=2.2cm,
  right=2.6cm,
  top=2.0cm,
  bottom=2.0cm,
  marginparwidth=0.8cm,
  marginparsep=0.3cm
}
\fi
% ===== 時間戳基礎設施 =====
\newcount\stamptimeval
\newcount\stampminuteval
\newcommand{\stamphh}{%
  \stamptimeval=\time
  \divide\stamptimeval by 60
  \ifnum\stamptimeval<10 0\fi
  \the\stamptimeval}
\newcommand{\stampmm}{%
  \stamptimeval=\time
  \stampminuteval=\time
  \divide\stampminuteval by 60
  \multiply\stampminuteval by 60
  \advance\stamptimeval by -\stampminuteval
  \ifnum\stamptimeval<10 0\fi
  \the\stamptimeval}
\newcommand{\stampwkde}[1]{%
  \ifcase#1Montag\or Dienstag\or Mittwoch\or Donnerstag\or Freitag\or Samstag\or Sonntag\fi}
\newcommand{\stampwkzh}[1]{%
  \ifcase#1 一\or 二\or 三\or 四\or 五\or 六\or 日\fi}
\newcommand{\stampmonthde}[1]{%
  \ifcase#1\or Januar\or Februar\or März\or April\or Mai\or Juni\or
  Juli\or August\or September\or Oktober\or November\or Dezember\fi}
\newcount\stampjdval
\newcount\stampwdval
\pgfcalendardatetojulian{\the\year-\the\month-\the\day}{\stampjdval}
\pgfcalendarjuliantoweekday{\the\stampjdval}{\stampwdval}
\newcommand{\timestampzh}{%
  \songfont 最後修改：\the\year 年\the\month 月\the\day 日%
  % （星期\stampwkzh{\the\stampwdval}）%
  % \space\stamphh:\stampmm
  }

\newcommand{\documentdatezh}{\the\year 年 \the\month 月 \the\day 日}
\newcommand{\documentdatede}{\the\day. \stampmonthde{\the\month} \the\year}
\newcommand{\bverfgetranslationnote}{%
  {\songfont 翻譯說明：初譯使用 Claude Code 與 GPT 協助迭代；仍請以德文原文為準。}%
}
\newcommand{\bverfgecourseinfo}{%
  {\songfont 課程：114 學年度第 2 學期；憲法專題研究（四）；授課老師：湯德宗教授；導讀日期：115 年 6 月 1 日。}%
}
\newcommand{\bverfgeeditorinfo}{%
  {\songfont 編者：王逸帆（R10A21126），國立台灣大學法律系研究所碩士班。}%
}
\newcommand{\bverfgetitleversionline}{%
  \ifbilingualcolumns
    Fassung \documentversion\enspace\textperiodcentered\enspace Stand: \documentdatede
    \quad
    {\songfont 版本 \documentversion\enspace\textperiodcentered\enspace 修訂日期：\documentdatezh\par}%
  \else
    \ifshowtranslation
      {\songfont 版本 \documentversion\enspace\textperiodcentered\enspace 修訂日期：\documentdatezh\par}%
    \else
      Fassung \documentversion\enspace\textperiodcentered\enspace Stand: \documentdatede\par
    \fi
  \fi
}
\newcommand{\bverfgetitlecornerinfo}{%
  \ifshowtranslation
    \vspace{0.72em}%
    \noindent\hfill
    \begin{minipage}{0.66\textwidth}
      \raggedleft\tiny\color{pendingink}\songfont
      \bverfgecourseinfo\par
      \bverfgeeditorinfo
    \end{minipage}%
  \fi
}
\newcommand{\bverfgetitlemeta}{%
  \vfill
  \begin{center}%
    \footnotesize\color{pendingink}%
    \bverfgetitleversionline
    \ifshowtranslation
      \vspace{0.28em}{\scriptsize\bverfgetranslationnote\par}%
    \fi
  \end{center}%
  \bverfgetitlecornerinfo
}

\pagestyle{fancy}
\fancyhf{}
\renewcommand{\headrulewidth}{0pt}
\renewcommand{\footrulewidth}{0pt}
\fancyfoot[C]{\thepage}
\ifbilingualcolumns
  \fancyfoot[R]{\scriptsize\color{pendingink}\timestampzh}%
\else
  \ifshoworiginal
  \else
    \fancyfoot[R]{\scriptsize\color{pendingink}\timestampzh\hspace{0.45cm}}%
  \fi
\fi
\setmainfont{Times New Roman}
\setsansfont{Helvetica Neue}
\setmonofont{Menlo}
\IfFileExists{/Users/iw/Library/Fonts/kaiu.ttf}{
  \setCJKmainfont[Path=/Users/iw/Library/Fonts/,AutoFakeBold=4]{kaiu.ttf}
}{
  \IfFontExistsTF{標楷體}{
    \setCJKmainfont[AutoFakeBold=4]{標楷體}
  }{
    \setCJKmainfont{Songti TC}
  }
}
\IfFontExistsTF{Songti TC}{
  \setCJKfamilyfont{song}{Songti TC}
  \setCJKmonofont{Songti TC}
}{
  \setCJKfamilyfont{song}[Path=/Users/iw/Library/Fonts/,AutoFakeBold=4]{kaiu.ttf}
  \setCJKmonofont[Path=/Users/iw/Library/Fonts/,AutoFakeBold=4]{kaiu.ttf}
}
\newcommand{\songfont}{\CJKfamily{song}}

% ===== 封面／目錄專用打字機字體（僅限封面、目錄頁使用，不影響正文）=====
\IfFontExistsTF{Radio Newsman}{%
  \newfontfamily\bverfgetitlede{Radio Newsman}%
}{%
  \newfontfamily\bverfgetitlede{Courier New}%
}
\IfFontExistsTF{Erikas Farbband}{%
  \newfontfamily\bverfgebodyde{Erikas Farbband}%
}{%
  \let\bverfgebodyde\bverfgetitlede
}
\IfFontExistsTF{Maoken Heavy Labourer}{%
  \setCJKfamilyfont{bverfgetitlecjk}{Maoken Heavy Labourer}%
}{%
  \setCJKfamilyfont{bverfgetitlecjk}{Songti TC}%
}
\IfFontExistsTF{Huiwen-mincho}{%
  \setCJKfamilyfont{bverfgebodycjk}{Huiwen-mincho}%
}{%
  \setCJKfamilyfont{bverfgebodycjk}{Songti TC}%
}
\newcommand{\bverfgetitlezh}{\bverfgetitlede\CJKfamily{bverfgetitlecjk}}
\newcommand{\bverfgebodyzh}{\bverfgebodyde\CJKfamily{bverfgebodycjk}}

\setstretch{1.35}
\setlist{itemsep=0.2em, topsep=0.4em}
\setlength{\parindent}{0pt}
\setlength{\parskip}{0.45em}
\raggedbottom
\setcounter{secnumdepth}{0}
\setcounter{tocdepth}{2}

\newcommand{\lawboxsourcefont}{\footnotesize}
\newcommand{\lawboxtransfont}{\footnotesize}
\newcommand{\lawboxsourcestretch}{1.02}
\newcommand{\lawboxtransstretch}{0.92}

\titleformat{\section}{\centering\Large\bfseries\songfont}{\thesection}{1em}{}
\titleformat{\subsection}{\large\bfseries\songfont}{\thesubsection}{1em}{}
\titleformat{\subsubsection}{\normalsize\bfseries\songfont}{\thesubsubsection}{1em}{}
\titleformat{\paragraph}{\normalsize\bfseries\songfont}{\theparagraph}{1em}{}
\titlespacing*{\section}{0pt}{1.8em}{1.0em}
\titlespacing*{\subsection}{0pt}{1.2em}{0.6em}
\titlespacing*{\subsubsection}{0pt}{1.0em}{0.45em}
\titlespacing*{\paragraph}{0pt}{0.6em}{0.4em}

\definecolor{noteblue}{HTML}{A9C9F5}
\definecolor{noteteal}{HTML}{AEE1D6}
\definecolor{notegray}{HTML}{D7DADF}
\definecolor{caseink}{HTML}{000000}
\definecolor{casepaper}{HTML}{FCFEFD}
\definecolor{sourcepaper}{HTML}{FAFBF8}
\definecolor{sourceline}{HTML}{D7DED8}
\definecolor{sourceink}{HTML}{000000}
\definecolor{pendingink}{HTML}{000000}

\tcbset{
  caseboxbase/.style={
    enhanced,
    breakable,
    width=0.985\linewidth,
    colback=casepaper,
    coltitle=caseink,
    fonttitle=\bfseries\songfont,
    boxrule=0.28pt,
    arc=1.2mm,
    left=2.4mm,
    right=2.4mm,
    top=1.5mm,
    bottom=1.5mm,
    boxsep=1.5mm,
    before skip=0.65em,
    after skip=0.65em,
    before upper={\setlength{\parindent}{0pt}\setlength{\parskip}{0.35em}},
  }
}

\newcommand{\bilingualstart}{}
\newcommand{\bilingualstop}{}
\ifbilingualcolumns
  \renewcommand{\bilingualstart}{\small\setlength{\columnsep}{1.25cm}\setlength{\columnseprule}{0.12pt}\columnratio{0.5,0.5}\multicolumnfootnotes\flushbottom\sloppy\interlinepenalty=80\clubpenalty=800\widowpenalty=800\displaywidowpenalty=800\begin{paracol}{2}\global\intranslationfalse\global\firstbilingualsectiontrue{\footnotesize\color{sourceink}Deutsch\par}\switchcolumn{\footnotesize\songfont\color{sourceink}中文譯文\par}\switchcolumn*}
  \renewcommand{\bilingualstop}{\ifintranslation\switchcolumn*\global\intranslationfalse\fi\end{paracol}\clearpage\raggedbottom}
\fi
\newif\ifrandnumbers
\randnumbersfalse
\newcounter{randnumber}
\newcounter{randnumberlocal}
\newcounter{lastsourcerandstart}
\newif\iflastsourcehasrandnumbers
\lastsourcehasrandnumbersfalse
\newif\ifinlawbox
\inlawboxfalse
\newif\iflawboxhaspendingrn
\lawboxhaspendingrnfalse
\newcommand{\lawboxpendingrn}{}
\newcommand{\startrandnumbers}{\global\randnumberstrue\setcounter{randnumber}{1}\global\lastsourcehasrandnumbersfalse}
\newcommand{\stoprandnumbers}{\global\randnumbersfalse\global\lastsourcehasrandnumbersfalse}
\newlength{\randnumberwidth}
\setlength{\randnumberwidth}{0.95cm}
\newlength{\randnumberrightoffset}
\setlength{\randnumberrightoffset}{0.85cm}
\newlength{\randnumberlawboxrightoffset}
\setlength{\randnumberlawboxrightoffset}{1.40cm}
\newcommand{\randnumberbox}[1]{%
  \leavevmode
  \makebox[0pt][l]{%
    \smash{%
      \tikz[remember picture,overlay,baseline=(rnmark.base)]{%
        \coordinate (rnmark) at (0,0);%
        \ifinlawbox
          \path let \p1=(current page.east), \p2=(rnmark.base) in
            node[
              anchor=base east,
              inner sep=0pt,
              outer sep=0pt,
              text width=\randnumberwidth,
              align=right,
              text=pendingink,
              font=\normalfont\mdseries\scriptsize\ttfamily
            ] at (\x1-\randnumberlawboxrightoffset,\y2) {{\normalfont\mdseries\scriptsize\ttfamily #1}};%
        \else
          \path let \p1=(current page.east), \p2=(rnmark.base) in
            node[
              anchor=base east,
              inner sep=0pt,
              outer sep=0pt,
              text width=\randnumberwidth,
              align=right,
              text=pendingink,
              font=\normalfont\mdseries\scriptsize\ttfamily
            ] at (\x1-\randnumberrightoffset,\y2) {{\normalfont\mdseries\scriptsize\ttfamily #1}};%
        \fi
      }%
    }%
  }%
  \ignorespaces
}
\newcommand{\lawboxstorerandnumber}[1]{%
  \gdef\lawboxpendingrn{#1}%
  \global\lawboxhaspendingrntrue
  \ignorespaces
}
\newcommand{\lawboxuserandnumber}{%
  \iflawboxhaspendingrn
    \randnumberbox{\lawboxpendingrn}%
    \global\lawboxhaspendingrnfalse
  \fi
}
\newcommand{\sourcern}[1]{%
  \ifrandnumbers
    \ifshoworiginal
      \ifshowtranslation\else
        \ifinlawbox\lawboxstorerandnumber{#1}\else\randnumberbox{#1}\fi
      \fi
    \fi
  \fi
}
\newcommand{\transrn}[1]{%
  \ifrandnumbers
    \ifshowtranslation
      \ifinlawbox\lawboxstorerandnumber{#1}\else\randnumberbox{#1}\fi
    \fi
  \fi
}
% ===== 課堂模式：德文原文縮寫註解 =====
% 日常切換只改各判決檔的一行：
%   \classroommodeon        開啟課堂版；註解掉這行即回到一般版。
% 模板預設：
%   1. 課堂模式關閉。
%   2. 課堂模式開啟時，縮寫註解包含「德文全稱＋中文譯名」。
%   3. 課堂模式開啟時，GG/條文編者註仍會自動顯示。
% 進階微調才需要在單案檔覆寫：
%   \classroomabbrzhoff     縮寫註解僅顯示德文全稱。
%   \showeditornotestrue    一般版也顯示編者註。
% 註解策略：同一實際 PDF 頁面中，同一縮寫只在第一次出現處加註。
% 這裡用 zref-abspage 的絕對頁序，不用 \thepage，避免文件內頁碼重置時誤判。
\newif\ifclassroommode
\newif\ifclassroomabbrzh
\classroommodefalse
\classroomabbrzhtrue
\newcommand{\classroommodeon}{\classroommodetrue}
\newcommand{\classroommodeoff}{\classroommodefalse}
\newcommand{\classroomabbrzhon}{\classroomabbrzhtrue}
\newcommand{\classroomabbrzhoff}{\classroomabbrzhfalse}
\newcommand{\abbrnotelabel}{\emph{Abk.}: }
\newcommand{\abbrnote}[3]{%
  \ifclassroommode
    \ifshoworiginal
      \ifintranslation\else
        \footnote{\normalfont\footnotesize\abbrnotelabel #2\ifclassroomabbrzh；{\songfont #3}\fi}%
      \fi
    \fi
  \fi
}
\newcommand{\abbrnoteonce}[4]{%
  #2%
  \ifclassroommode
    \ifshoworiginal
      \ifintranslation\else
        \begingroup
          \edef\abbrcurrentabspage{\number\numexpr\value{abspage}+1\relax}%
          \ifcsname abbrnoteseen@#1@\abbrcurrentabspage\endcsname\else
            \global\expandafter\let\csname abbrnoteseen@#1@\abbrcurrentabspage\endcsname\empty
            \abbrnote{#1}{#3}{#4}%
          \fi
        \endgroup
      \fi
    \fi
  \fi
}
\newcommand{\abbrAbs}{\abbrnoteonce{abs}{Abs.}{Absatz}{項}}
\newcommand{\abbrAAO}{\abbrnoteonce{aao}{a.a.O.}{am angegebenen Ort}{前引處}}
\newcommand{\abbrAE}{\abbrnoteonce{ae}{AE}{Alternativ-Entwurf}{替代草案}}
\newcommand{\abbrArt}{\abbrnoteonce{art}{Art.}{Artikel}{條}}
\newcommand{\abbrAufl}{\abbrnoteonce{aufl}{Aufl.}{Auflage}{版}}
\newcommand{\abbrBd}{\abbrnoteonce{bd}{Bd.}{Band}{卷}}
\newcommand{\abbrBGBl}{\abbrnoteonce{bgbl}{BGBl.}{Bundesgesetzblatt}{聯邦法律公報}}
\newcommand{\abbrBGHSt}{\abbrnoteonce{bghst}{BGHSt}{Entscheidungen des Bundesgerichtshofes in Strafsachen}{聯邦普通法院刑事判決彙編}}
\newcommand{\abbrBRDrucks}{\abbrnoteonce{brdrucks}{BRDrucks.}{Bundesrats-Drucksache}{聯邦參議院文件}}
\newcommand{\abbrBTDrucks}{\abbrnoteonce{btdrucks}{BTDrucks.}{Bundestags-Drucksache}{聯邦議院文件}}
\newcommand{\abbrBundesgesetzbl}{\abbrnoteonce{bundesgesetzbl}{Bundesgesetzbl.}{Bundesgesetzblatt}{聯邦法律公報}}
\newcommand{\abbrBvF}{\abbrnoteonce{bvf}{BvF}{Bundesverfassungsgerichtsverfahren}{聯邦憲法法院程序案號}}
\newcommand{\abbrBvL}{\abbrnoteonce{bvl}{BvL}{Bundesverfassungsgerichtsverfahren}{聯邦憲法法院程序案號}}
\newcommand{\abbrBvR}{\abbrnoteonce{bvr}{BvR}{Bundesverfassungsgerichtsverfahren}{聯邦憲法法院程序案號}}
\newcommand{\abbrBVerfG}{\abbrnoteonce{bverfg}{BVerfG}{Bundesverfassungsgericht}{聯邦憲法法院}}
\newcommand{\abbrBVerfGE}{\abbrnoteonce{bverfge}{BVerfGE}{Entscheidungen des Bundesverfassungsgerichts}{聯邦憲法法院判決集}}
\newcommand{\abbrBVerfGG}{\abbrnoteonce{bverfgg}{BVerfGG}{Bundesverfassungsgerichtsgesetz}{聯邦憲法法院法}}
\newcommand{\abbrDDR}{\abbrnoteonce{ddr}{DDR}{Deutsche Demokratische Republik}{德意志民主共和國；東德}}
\newcommand{\abbrDJT}{\abbrnoteonce{djt}{DJT}{Deutscher Juristentag}{德國法學家大會}}
\newcommand{\abbrDr}{\abbrnoteonce{dr}{Dr.}{Doktor}{Dr.}}
\newcommand{\abbrEuGRZ}{\abbrnoteonce{eugrz}{EuGRZ}{Europäische Grundrechte-Zeitschrift}{歐洲基本權期刊}}
\newcommand{\abbrEV}{\abbrnoteonce{ev}{EV}{Einigungsvertrag}{統一條約}}
\newcommand{\abbrF}{\abbrnoteonce{f}{f.}{folgende}{以下一頁；下一段}}
\newcommand{\abbrFF}{\abbrnoteonce{ff}{ff.}{fortfolgende}{以下數頁；以下各段}}
\newcommand{\abbrGBlDDR}{\abbrnoteonce{gblddr}{GBl. DDR}{Gesetzblatt der Deutschen Demokratischen Republik}{德意志民主共和國法律公報}}
\newcommand{\abbrGG}{\abbrnoteonce{gg}{GG}{Grundgesetz}{基本法}}
\newcommand{\abbrGVBl}{\abbrnoteonce{gvbl}{GVBl.}{Gesetz- und Verordnungsblatt}{法律及命令公報}}
\newcommand{\abbrJR}{\abbrnoteonce{jr}{JR}{Juristische Rundschau}{法學評論}}
\newcommand{\abbrJZ}{\abbrnoteonce{jz}{JZ}{JuristenZeitung}{法學家報}}
\newcommand{\abbrIVm}{\abbrnoteonce{ivm}{i.V.m.}{in Verbindung mit}{結合；連同}}
\newcommand{\abbrMdB}{\abbrnoteonce{mdb}{MdB}{Mitglied des Bundestages}{聯邦議院議員}}
\newcommand{\abbrMwN}{\abbrnoteonce{mwn}{m.w.N.}{mit weiteren Nachweisen}{附進一步文獻／判例出處}}
\newcommand{\abbrNF}{\abbrnoteonce{nf}{n.F.}{neue Fassung}{新版本}}
\newcommand{\abbrAF}{\abbrnoteonce{af}{a.F.}{alte Fassung}{舊版本}}
\newcommand{\abbrNr}{\abbrnoteonce{nr}{Nr.}{Nummer}{款、目或號，依語境}}
\newcommand{\abbrProf}{\abbrnoteonce{prof}{Prof.}{Professor}{教授}}
\newcommand{\abbrRdnr}{\abbrnoteonce{rdnr}{Rdnr.}{Randnummer}{邊碼；段碼}}
\newcommand{\abbrRdnrn}{\abbrnoteonce{rdnrn}{Rdnrn.}{Randnummern}{邊碼；段碼}}
\newcommand{\abbrRGBl}{\abbrnoteonce{rgbl}{RGBl.}{Reichsgesetzblatt}{帝國法律公報}}
\newcommand{\abbrRGSt}{\abbrnoteonce{rgst}{RGSt}{Entscheidungen des Reichsgerichts in Strafsachen}{帝國法院刑事判決彙編}}
\newcommand{\abbrRspr}{\abbrnoteonce{rspr}{Rspr.}{Rechtsprechung}{判例；司法實務}}
\newcommand{\abbrRVO}{\abbrnoteonce{rvo}{RVO}{Reichsversicherungsordnung}{帝國保險法}}
\newcommand{\abbrS}{\abbrnoteonce{s}{S.}{Seite}{頁}}
\newcommand{\abbrSFHG}{\abbrnoteonce{sfhg}{SFHG}{Schwangeren- und Familienhilfegesetz}{孕婦及家庭扶助法}}
\newcommand{\abbrSGBV}{\abbrnoteonce{sgbv}{SGB V}{Fünftes Buch Sozialgesetzbuch}{社會法典第五編}}
\newcommand{\abbrStAG}{\abbrnoteonce{staeg}{StÄG}{Strafrechtsänderungsgesetz}{刑法修正法}}
\newcommand{\abbrStenBer}{\abbrnoteonce{stenber}{StenBer.}{Stenographischer Bericht}{速記紀錄}}
\newcommand{\abbrStGB}{\abbrnoteonce{stgb}{StGB}{Strafgesetzbuch}{刑法}}
\newcommand{\abbrStREG}{\abbrnoteonce{streg}{StREG}{Strafrechtsreform-Ergänzungsgesetz}{刑法改革補充法}}
\newcommand{\abbrStrRG}{\abbrnoteonce{strrg}{StrRG}{Strafrechtsreformgesetz}{刑法改革法}}
\newcommand{\abbrUS}{\abbrnoteonce{us}{U.S.}{United States Reports}{美國最高法院判決彙編}}
\newcommand{\abbrVgl}{\abbrnoteonce{vgl}{vgl.}{vergleiche}{參見}}
\newcommand{\abbrVVDStRL}{\abbrnoteonce{vvdstrl}{VVDStRL}{Veröffentlichungen der Vereinigung der Deutschen Staatsrechtslehrer}{德國公法教師協會論文集}}
\newcommand{\abbrWP}{\abbrnoteonce{wp}{WP.}{Wahlperiode}{屆期}}
\newcommand{\abbrWp}{\abbrnoteonce{wp}{Wp.}{Wahlperiode}{屆期}}
\newcommand{\abbrZStrW}{\abbrnoteonce{zstrw}{ZStrW}{Zeitschrift für die gesamte Strafrechtswissenschaft}{整體刑法學期刊}}
\newcommand{\recordsourcerandstart}{%
  \ifrandnumbers
    \setcounter{lastsourcerandstart}{\value{randnumber}}%
    \global\lastsourcehasrandnumberstrue
  \else
    \global\lastsourcehasrandnumbersfalse
  \fi
}
\newlength{\biparagraphskip}
\setlength{\biparagraphskip}{0.52em}
\newlength{\lawboxblockbeforeskip}
\setlength{\lawboxblockbeforeskip}{0.72em}
\newlength{\lawboxblockafterskip}
\setlength{\lawboxblockafterskip}{0.92em}
\newif\iflawboxpartendedblock
\lawboxpartendedblockfalse
\newcommand{\sourceparstyle}{\footnotesize\color{sourceink}\setstretch{1.12}\RaggedRight\emergencystretch=3em\setlength{\parskip}{0.42em}\obeylines\selectfont}
\newcommand{\transparstyle}{\songfont\color{caseink}\setstretch{1.12}\setlength{\parindent}{0pt}\setlength{\parskip}{0.42em}}
\newcommand{\bilingualpairskip}{%
  \par\addvspace{\biparagraphskip}%
  \switchcolumn
  \par\addvspace{\biparagraphskip}%
  \switchcolumn*%
  \global\intranslationfalse
}
\newcommand{\bilingualboxpairskip}{%
  \par
  \switchcolumn
  \par
  \switchcolumn*%
  \global\intranslationfalse
}
\newcommand{\bilinguallawboxblockskip}{%
  \switchcolumn*%
  \par\addvspace{\lawboxblockafterskip}%
  \switchcolumn
  \par\addvspace{\lawboxblockafterskip}%
  \switchcolumn*%
  \global\intranslationfalse
}
\newcommand{\bikeepwithnext}[1][4]{%
  \ifbilingualcolumns
    \syncleft
    \Needspace{#1\baselineskip}%
    \switchcolumn
    \Needspace{#1\baselineskip}%
    \switchcolumn*%
    \global\intranslationfalse
  \else
    \Needspace{#1\baselineskip}%
  \fi
}
\NewEnviron{sourcepara}[1][]{
  \recordsourcerandstart
  \ifshoworiginal
    \ifbilingualcolumns
      \ifintranslation\switchcolumn*\global\intranslationfalse\fi
    \else
      \Needspace{5\baselineskip}
    \fi
    \ifbilingualcolumns\else\par\addvspace{0.35em}\fi
    {\sourceparstyle
    \noindent\BODY\par}
    \ifbilingualcolumns\else\addvspace{0.25em}\fi
    \ifbilingualcolumns
      \switchcolumn
      \global\intranslationtrue
    \fi
  \else
    \ifrandnumbers
      \begingroup
        \setbox0=\vbox{\hsize=\linewidth\sourceparstyle\noindent\BODY\par}%
      \endgroup
    \fi
  \fi
}
\NewEnviron{transpara}{
  \ifshowtranslation
    \setcounter{randnumberlocal}{\value{lastsourcerandstart}}%
    \ifbilingualcolumns
      \ifintranslation\else\switchcolumn\global\intranslationtrue\fi
      {\transparstyle\setlength{\leftskip}{0.55em}\setlength{\rightskip}{0.15em plus 1em}\addtolength{\linewidth}{-0.55em}\BODY\par}
      \switchcolumn*
      \bilingualpairskip
    \else
      {\transparstyle\BODY\par}
    \fi
  \else
    \ifbilingualcolumns
      \ifintranslation\switchcolumn*\global\intranslationfalse\fi
    \fi
  \fi
}
\NewEnviron{sourceboxpara}[1][]{
  \global\lawboxpartendedblockfalse
  \recordsourcerandstart
  \ifshoworiginal
    \ifbilingualcolumns
      \ifintranslation\switchcolumn*\global\intranslationfalse\fi
    \else
      \Needspace{3\baselineskip}
    \fi
    {\sourceparstyle
    \BODY\par}
    \ifbilingualcolumns\else
      \iflawboxpartendedblock\addvspace{\lawboxblockafterskip}\fi
    \fi
    \ifbilingualcolumns
      \switchcolumn
      \global\intranslationtrue
    \fi
  \fi
}
\NewEnviron{transboxpara}{
  \global\lawboxpartendedblockfalse
  \ifshowtranslation
    \setcounter{randnumberlocal}{\value{lastsourcerandstart}}%
    \ifbilingualcolumns
      \ifintranslation\else\switchcolumn\global\intranslationtrue\fi
      {\transparstyle\BODY\par}
      \iflawboxpartendedblock
        \bilinguallawboxblockskip
      \else
        \switchcolumn*
        \bilingualboxpairskip
      \fi
    \else
      {\transparstyle\BODY\par}
      \iflawboxpartendedblock\addvspace{\lawboxblockafterskip}\fi
    \fi
  \else
    \ifbilingualcolumns
      \ifintranslation\switchcolumn*\global\intranslationfalse\fi
    \fi
  \fi
}
\newcommand{\leitsatznum}[1]{\noindent\hangindent=3.0em\hangafter=1\makebox[2.25em][r]{\bfseries #1.}\hspace{0.55em}\ignorespaces}
\newcommand{\syncleft}{\ifbilingualcolumns\ifintranslation\switchcolumn*\global\intranslationfalse\fi\fi}
\pretocmd{\section}{\syncleft\clearpage}{}{}
\pretocmd{\subsection}{\syncleft}{}{}
\pretocmd{\subsubsection}{\syncleft}{}{}
\newcommand{\biheadingfont}{\songfont}
\newcommand{\bitocentry}[2]{%
  \ifshoworiginal
    \ifshowtranslation
      \ifstrequal{#1}{#2}{#1}{#1\space #2}%
    \else
      #1%
    \fi
  \else
    #2%
  \fi
}
\pdfstringdefDisableCommands{%
  \def\bitocentry#1#2{#1}%
}
\newcommand{\biheading}[7]{%
  \syncleft#1\bikeepwithnext[#3]\phantomsection\addcontentsline{toc}{#2}{\bitocentry{#6}{#7}}%
  \ifbilingualcolumns
    {#4 #6\par}%
    \switchcolumn
    {#4\biheadingfont #7\par}%
    \switchcolumn*%
    \global\intranslationfalse
    \par\addvspace{#5}%
    \switchcolumn
    \par\addvspace{#5}%
    \switchcolumn*%
  \else
    \ifshoworiginal{#4 #6\par}\fi
    \ifshowtranslation{#4\biheadingfont #7\par}\fi
    \addvspace{#5}%
  \fi
}
\newcommand{\termsectionheading}[1]{%
  \par\Needspace{9\baselineskip}%
  \vspace{0.65em}%
  {\songfont\bfseries #1\par}%
  \vspace{0.2em}%
}
% 章節換頁：
%   雙語 paracol 欄內不要用 \clearpage；它會在同步兩欄時吐出只有欄線/頁碼的空頁。
%   \newpage 足以讓下一個 section 起新頁，又不會強制清空所有待排材料。
%   單語模式不在 paracol 內，仍可用 \clearpage。
\newcommand{\bisectionpagebreak}{\ifbilingualcolumns\newpage\else\clearpage\fi}
\newcommand{\bisection}[2]{%
  \biheading{\iffirstbilingualsection\global\firstbilingualsectionfalse\else\bisectionpagebreak\fi}{section}{6}{\centering\Large\bfseries}{0.8em}{#1}{#2}%
}
\newcommand{\bisubsection}[2]{%
  \biheading{}{subsection}{5}{\large\bfseries}{0.55em}{#1}{#2}%
}
\newcommand{\bisubsubsection}[2]{%
  \biheading{}{subsubsection}{4}{\normalsize\bfseries}{0.45em}{#1}{#2}%
}
\newcommand{\bicompositeheading}[4]{%
  \biheading{}{subsection}{5}{\large\bfseries}{0.16em}{#1}{#2}%
  \biheading{}{subsubsection}{3}{\normalsize\bfseries}{0.45em}{#3}{#4}%
}
\newif\iflawboxfirstitem
\lawboxfirstitemfalse
\newcommand{\lawboxprespace}[1]{%
  \par
  \iflawboxfirstitem
    \global\lawboxfirstitemfalse
  \else
    \addvspace{#1}%
  \fi
}
\newtcolorbox{lawboxinner}{
  caseboxbase,
  colframe=sourceline!85!black,
  colback=white,
  left=1.05em,
  right=1.05em,
  top=0.62em,
  bottom=0.62em,
  before upper={\global\inlawboxtrue\global\lawboxfirstitemtrue\global\lawboxhaspendingrnfalse\ifintranslation\lawboxtransfont\else\lawboxsourcefont\fi\RaggedRight\setlength{\parindent}{0pt}\setlength{\parskip}{0pt}\ifintranslation\setstretch{\lawboxtransstretch}\else\setstretch{\lawboxsourcestretch}\fi},
  after upper={\global\lawboxhaspendingrnfalse\global\inlawboxfalse}
}
\tcbset{
  lawboxpartbase/.style={
    enhanced,
    breakable=false,
    width=0.985\linewidth,
    colback=white,
    frame hidden,
    boxrule=0pt,
    arc=0pt,
    left=1.05em,
    right=1.05em,
    boxsep=1.5mm,
    before skip=0pt,
    after skip=0pt,
    before upper={\global\inlawboxtrue\global\lawboxfirstitemtrue\global\lawboxhaspendingrnfalse\ifintranslation\lawboxtransfont\else\lawboxsourcefont\fi\RaggedRight\setlength{\parindent}{0pt}\setlength{\parskip}{0pt}\ifintranslation\setstretch{\lawboxtransstretch}\else\setstretch{\lawboxsourcestretch}\fi},
    after upper={\global\lawboxhaspendingrnfalse\global\inlawboxfalse},
    borderline west={0.28pt}{0pt}{sourceline!85!black},
    borderline east={0.28pt}{0pt}{sourceline!85!black},
  },
  lawboxpartfirst/.style={
    lawboxpartbase,
    top=0.62em,
    bottom=0.04em,
    before skip=\lawboxblockbeforeskip,
    borderline north={0.28pt}{0pt}{sourceline!85!black},
  },
  lawboxpartmiddle/.style={
    lawboxpartbase,
    top=0.04em,
    bottom=0.04em,
  },
  lawboxpartlast/.style={
    lawboxpartbase,
    top=0.04em,
    bottom=0.62em,
    after skip=0pt,
    borderline south={0.28pt}{0pt}{sourceline!85!black},
  }
}
\newtcolorbox{lawboxpartinner}[1]{#1}
\newtcolorbox{termboxinner}{
  caseboxbase,
  colframe=noteblue!70!black,
  colback=casepaper,
  before upper={\songfont\setlength{\parindent}{0pt}\setlength{\parskip}{0.35em}}
}
\newtcolorbox{deboxinner}{
  caseboxbase,
  colframe=notegray!72!black,
  colback=white,
  left=1.05em,
  right=1.05em,
  top=0.62em,
  bottom=0.62em,
  before upper={\global\inlawboxtrue\global\lawboxfirstitemtrue\global\lawboxhaspendingrnfalse\small\setlength{\parindent}{0pt}\setlength{\parskip}{0.35em}},
  after upper={\global\lawboxhaspendingrnfalse\global\inlawboxfalse}
}
\NewEnviron{lawbox}{
  \begin{lawboxinner}\BODY\end{lawboxinner}
}
\NewEnviron{lawboxpart}[2]{
  \begin{lawboxpartinner}{lawboxpart#1,equal height group=#2}\BODY\end{lawboxpartinner}%
  \ifstrequal{#1}{last}{\global\lawboxpartendedblocktrue}{}%
}
\NewEnviron{termbox}{
  \par\addvspace{0.55em}
  {\color{noteblue!70!black}\rule{\linewidth}{0.28pt}}\par
  \begingroup
    \songfont\small\color{caseink}
    \setlength{\parindent}{0pt}
    \setlength{\parskip}{0.24em}
    \BODY
    \par
  \endgroup
  {\color{noteblue!70!black}\rule{\linewidth}{0.28pt}}\par
  \addvspace{0.55em}
}
\NewEnviron{debox}{
  \begin{deboxinner}\BODY\end{deboxinner}
}
\providecommand{\paragraphmark}[1]{}
\newcommand{\lawboxtitleafter}{\addvspace{0.06em}}
\newcommand{\lawtitle}[1]{\lawboxprespace{0.22em}{\lawboxtransfont\bfseries\lawboxuserandnumber #1}\par\lawboxtitleafter}
\newcommand{\absno}[2]{\lawboxprespace{0.04em}\noindent\lawboxuserandnumber\hangindent=2.6em\hangafter=1\makebox[2.6em][l]{(#1)}#2\par}
\newcommand{\nrno}[2]{\lawboxprespace{0pt}\noindent\lawboxuserandnumber\hspace*{2.6em}\hangindent=4.4em\hangafter=1\makebox[1.8em][l]{#1.}#2\par}
\newcommand{\letterno}[2]{\lawboxprespace{0pt}\noindent\lawboxuserandnumber\hspace*{4.4em}\hangindent=6.0em\hangafter=1\makebox[1.6em][l]{#1)}#2\par}
\newcommand{\sourcelawtitle}[1]{\lawboxprespace{0.22em}{\lawboxsourcefont\bfseries\lawboxuserandnumber #1}\par\lawboxtitleafter}
\newcommand{\sourcelawsubtitle}[1]{\lawboxprespace{0.04em}{\lawboxsourcefont\bfseries\lawboxuserandnumber #1}\par\lawboxtitleafter}
\newcommand{\sourceabsno}[2]{\lawboxprespace{0.04em}\noindent\lawboxuserandnumber\hangindent=2.45em\hangafter=1\makebox[2.45em][l]{(#1)}#2\par}
\newcommand{\sourcenrno}[2]{\lawboxprespace{0pt}\noindent\lawboxuserandnumber\hspace*{2.45em}\hangindent=4.25em\hangafter=1\makebox[1.8em][l]{#1.}#2\par}
\newcommand{\sourceletterno}[2]{\lawboxprespace{0pt}\noindent\lawboxuserandnumber\hspace*{4.25em}\hangindent=5.85em\hangafter=1\makebox[1.6em][l]{#1)}#2\par}
\newcommand{\sourcequotepara}[1]{\lawboxprespace{0.04em}\noindent\lawboxuserandnumber\hangindent=1.2em\hangafter=1 #1\par}
\newcommand{\sourcelawline}[1]{\lawboxprespace{0.04em}\noindent\lawboxuserandnumber #1\par}
\newcommand{\lawline}[1]{\lawboxprespace{0.04em}\noindent\lawboxuserandnumber #1\par}
\newcommand{\pendingtranslation}{\par{\songfont\footnotesize\color{pendingink}（中文譯文待補。）}\par}
\newcommand{\tocpart}[1]{\par\addvspace{0.52em}{\footnotesize\bfseries\color{pendingink}#1\par}\addvspace{0.16em}}
\newcommand{\tocdashline}{\par\addvspace{0.48em}\noindent{\color{notegray}\xleaders\hbox to 0.82em{\hfil\rule{0.42em}{0.28pt}\hfil}\hskip\linewidth}\par\addvspace{0.38em}}
% \tocpanel{font-cmd}{content}：將目錄置中於所在欄寬（雙語欄適用）。
% A3 橫式使用 0.58\linewidth，A4 橫式使用 0.84\linewidth，
% 使兩種紙張下目錄欄內容實際寬度相近（均約 100–105 mm）。
\newcommand{\tocpanel}[2]{%
  \makebox[\linewidth][c]{%
    \ifx\bilingualpaper\bilingualpaperaThree
      \begin{minipage}{0.58\linewidth}%
    \else
      \begin{minipage}{0.84\linewidth}%
    \fi
    \toccolstyle
    #1#2%
    \end{minipage}%
  }%
}
% ===== 首頁簡目錄的兩層 description list 縮排 =====
% enumitem 的 description list 可把每一列拆成：
%   [label]  labelsep  body text
% 這裡的「起點」都以所在 minipage/欄位的左緣為 0。
%
% 核心規則：
%   labelindent = label 欄位從左緣開始的位置。
%   labelwidth  = label 欄位本身寬度；標籤太長（如 Abw. Meinung）就加大它。
%   labelsep    = label 欄位右緣到正文的距離。
%   leftmargin  = 正文 body text 的起點。判斷層級視覺時主要看這個。
%
% 所以：
%   想讓「Begründung / 判決理由」整列正文往右或往左：改 tocitems 的 leftmargin。
%   想讓「A. Verfahren und Gegenstand / A. 程序與審查標的」正文往右或往左：
%     改 tocsubitems 的 leftmargin。
%   想讓 A./B./C. 這些標籤本身往右或往左，但正文不動：改 tocsubitems 的 labelindent。
%   想讓 A./B./C. 與正文間距變大或變小：改 tocsubitems 的 labelsep。
%
% 層級原則：
%   tocsubitems.leftmargin 必須大於 tocitems.leftmargin，
%   否則子層級正文會看起來比「Gründe / 理由」還靠左。
\newlist{tocitems}{description}{1}
\setlist[tocitems]{%
  labelindent=0pt,    % 第一層 label 欄位起點；0pt 表示貼齊目錄欄左緣。
  leftmargin=2.54cm,  % 第一層正文起點；例如 Begründung / 判決理由 的左緣。
  labelwidth=2.16cm,  % 第一層 label 欄位寬度；需容納 Leitsätze、Abw. Meinung、不同意見等。
  labelsep=0.32cm,    % 第一層 label 與正文的水平間距；只改間距，不改正文起點。
  itemsep=0.28em,     % 第一層各項之間的垂直距離。
  topsep=0.20em,      % 第一層 list 上下與前後內容的垂直距離。
  parsep=0pt,         % 同一 item 內段落之間的距離；目錄項通常不需要。
  partopsep=0pt,      % list 若緊接段落開始時的額外距離；關閉以免目錄高度飄動。
  align=left,         % label 欄內左對齊；長短標籤不靠右貼齊。
  font=\normalfont\bfseries % label 的字型；正文不受這行影響。
}
\newlist{tocsubitems}{description}{1}
\setlist[tocsubitems]{%
  labelindent=1cm,    % 第二層 label 起點；只控制 A./B./C. 本身的位置。
  leftmargin=3.00cm,  % 第二層正文起點；必須大於 2.54cm，才會比 Begründung / 判決理由更右。
  labelwidth=0.5cm,   % 第二層 label 欄寬；A./B./C. 很短，可小於第一層。
  labelsep=1.5cm,    % 第二層 label 與正文間距；A. 與 Verfahren / 程序 的距離看這裡。
  itemsep=0.12em,     % 第二層各項較密，避免 A.–E. 佔太高。
  topsep=0.04em,      % 第二層 list 與 Gründe / 理由之間的垂直距離。
  parsep=0pt,         % 同一第二層 item 內段落距離；目錄項通常不需要。
  partopsep=0pt,      % 關閉額外段落起始距離，避免版面不穩。
  align=left,         % A./B./C. label 欄內左對齊。
  font=\normalfont\bfseries, % A./B./C. label 加粗；正文不受這行影響。
  before={}           % 不在第二層 list 前自動插入額外內容。
}
\newlist{termitems}{description}{1}
\ifbilingualcolumns
  \setlist[termitems]{%
    leftmargin=5.2cm,
    labelwidth=4.8cm,
    labelsep=0.35cm,
    itemsep=0.16em,
    topsep=0.2em,
    parsep=0pt,
    partopsep=0pt,
    font=\normalfont\bfseries,
    style=nextline
  }
\else
  \setlist[termitems]{%
    leftmargin=0pt,
    labelwidth=0pt,
    labelsep=0pt,
    itemsep=0.24em,
    topsep=0.2em,
    parsep=0pt,
    partopsep=0pt,
    font=\normalfont\bfseries,
    style=nextline
  }
\fi
\newlist{abbritems}{description}{1}
\setlist[abbritems]{%
  leftmargin=2.38cm,
  labelwidth=2.02cm,
  labelsep=0.28cm,
  itemsep=0.16em,
  topsep=0.2em,
  parsep=0pt,
  partopsep=0pt,
  align=left,
  font=\normalfont\bfseries
}
\ifbilingualcolumns
  \newenvironment{termcolumns}{\begin{multicols}{2}}{\end{multicols}}
\else
  \newenvironment{termcolumns}{}{}
\fi

% ===== 編者註腳開關 =====
% 一般版預設不顯示編者註，避免正文腳註過密。
% 若開啟 \classroommodeon，GG/條文編者註仍會自動顯示，無需另開本開關。
% 若一般版也要顯示編者註，可在單案檔 \input 後加 \showeditornotestrue。
\newif\ifshoweditornotes
\showeditornotesfalse

% ===== 目錄排版輔助 =====
\newcommand{\toccolstyle}{\raggedright\arraybackslash}
\newcommand{\tocsingleindent}{\leftskip=1.45cm\rightskip=1.2cm}

% ===== 行內引文環境（德文原文與中文譯文各一，帶左側灰色邊線）=====
\newcommand{\quoteparbreak}{\par\addvspace{0.48em}}
\newcommand{\sourcequoteinner}[1]{%
  \begin{tcolorbox}[
    enhanced, breakable, width=\dimexpr\linewidth-0.8em\relax, frame hidden,
    colback=white, coltext=caseink, boxrule=0pt,
    borderline west={0.55pt}{0.88em}{notegray},
    left=1.78em, right=0.72em, top=0.18em, bottom=0.18em,
    boxsep=0pt, before skip=0.24em, after skip=0.24em,
    before upper={\normalfont\footnotesize\color{caseink}\setlength{\parindent}{0pt}\setlength{\parskip}{0.34em}\raggedright\setlength{\rightskip}{0pt plus 1em}}
  ]%
  #1\par
  \end{tcolorbox}%
}
\newcommand{\transquoteinner}[1]{%
  \begin{tcolorbox}[
    enhanced, breakable, width=\dimexpr\linewidth-0.8em\relax, frame hidden,
    colback=white, coltext=caseink, boxrule=0pt,
    borderline west={0.55pt}{0.88em}{notegray},
    left=1.78em, right=0.72em, top=0.18em, bottom=0.18em,
    boxsep=0pt, before skip=0.24em, after skip=0.24em,
    before upper={\songfont\small\color{caseink}\setlength{\parindent}{0pt}\setlength{\parskip}{0.34em}\raggedright\setlength{\rightskip}{0pt plus 1em}}
  ]%
  #1\par
  \end{tcolorbox}%
}

% ===== 大綱（Gliederung）Rn 輔助命令 =====
\newcommand{\outrn}[2]{\enspace{\normalfont\footnotesize\color{pendingink}(Rn\,#1–#2)}}
\newcommand{\outrnsingle}[1]{\enspace{\normalfont\footnotesize\color{pendingink}(Rn\,#1)}}
\newcommand{\outrnzh}[2]{\enspace{\normalfont\footnotesize\color{pendingink}（第\,#1–#2\,段）}}
\newcommand{\outrnzhs}[1]{\enspace{\normalfont\footnotesize\color{pendingink}（第\,#1\,段）}}
% ===== 大綱雙語對照表格輔助命令（三欄：段碼 │ 德文 │ 中文）=====
% 欄 1：段碼（由欄定義自動格式化：小字、等寬、右對齊、pendingink）
% 欄 2：德文標籤＋正文
% 欄 3：中文標籤＋正文
%
% 無段碼的列（\omainrow、\osecrow），欄 1 以 & 空置。
% 有段碼的列（\osecrowrn、\osubrow、\ointrow），#1 為預格式化字串
%   例：\osubrow{2–49}{I.}{...}{I.}{...}
%       \ointrow{107}{...}{...}
%
% \opartrow：段組標題（橫跨三欄）
\newcommand{\opartrow}[2]{%
  \noalign{\vspace{0.30em}}%
  \multicolumn{3}{@{}l@{}}{%
    \footnotesize\bfseries\color{pendingink}#1\hfill\songfont#2%
  }\\[0.06em]%
}
% \odashrow：虛線分隔（橫跨三欄）
\newcommand{\odashrow}{%
  \noalign{\vspace{0.28em}}%
  \multicolumn{3}{@{}l@{}}{\color{notegray}%
    \xleaders\hbox to 0.82em{\hfil\rule{0.42em}{0.28pt}\hfil}\hskip 0.88\linewidth}%
  \\[0.24em]%
}
% \odissentpart：不同意見書區段標頭。比一般 \opartrow 更像附錄性轉場。
\newcommand{\odissentpart}[2]{%
  \noalign{\vspace{0.42em}}%
  \multicolumn{3}{@{}p{0.88\textwidth}@{}}{%
    \color{notegray}\rule{\linewidth}{0.18pt}%
  }\\[-0.18em]%
  \multicolumn{3}{@{}p{0.88\textwidth}@{}}{%
    \makebox[\linewidth][s]{%
      \footnotesize\bfseries\color{pendingink}#1%
      \hfill{\songfont #2}%
    }%
  }\\[-0.02em]%
}
% \odissentopinion：不同意見書的署名與一行摘要，右欄保留段碼範圍。
\newcommand{\odissentopinion}[5]{%
  \footnotesize\hangindent=0.52cm\hangafter=1%
    {\bfseries\color{caseink}#2}\enspace{\color{notegray}\textperiodcentered}\enspace\color{caseink}#3%
  &\footnotesize\hangindent=0.52cm\hangafter=1%
    {\songfont\bfseries\color{caseink}#4}\enspace{\color{notegray}\textperiodcentered}\enspace\songfont\color{caseink}#5%
  &#1%
  \\[0.12em]%
}
% \omainrow：頂層項目，無段碼（Leitsätze、Urteil、Gründe …）
\newcommand{\omainrow}[4]{%
  \footnotesize\hangindent=1.92cm\hangafter=1%
    \makebox[1.92cm][l]{\bfseries\color{caseink}#1}\color{caseink}#2%
  &\footnotesize\hangindent=1.92cm\hangafter=1%
    \makebox[1.92cm][l]{\bfseries\color{caseink}#3}\songfont\color{caseink}#4%
  &%
  \\[0.15em]%
}
% \omainrowrn：頂層項目，有段碼（不同意見等標籤寬於 \osecrow 框者）
\newcommand{\omainrowrn}[5]{%
  \footnotesize\hangindent=1.92cm\hangafter=1%
    \makebox[1.92cm][l]{\bfseries\color{caseink}#2}\color{caseink}#3%
  &\footnotesize\hangindent=1.92cm\hangafter=1%
    \makebox[1.92cm][l]{\bfseries\color{caseink}#4}\songfont\color{caseink}#5%
  &#1%
  \\[0.15em]%
}
% \osecrow：一級子項，無段碼（A.、C.、D. 等有子項者）
\newcommand{\osecrow}[4]{%
  \footnotesize\hspace{1.10cm}\hangindent=1.98cm\hangafter=1%
    \makebox[0.86cm][l]{\bfseries\color{caseink}#1}\color{caseink}#2%
  &\footnotesize\hspace{1.10cm}\hangindent=1.98cm\hangafter=1%
    \makebox[0.86cm][l]{\bfseries\color{caseink}#3}\songfont\color{caseink}#4%
  &%
  \\[0.11em]%
}
% \osecrowrn：一級子項，有段碼（B.、E. 等完整段落範圍者）
% #1=段碼字串（如 101–106）, #2=DE標籤, #3=DE正文, #4=ZH標籤, #5=ZH正文
\newcommand{\osecrowrn}[5]{%
  \footnotesize\hspace{1.10cm}\hangindent=1.98cm\hangafter=1%
    \makebox[0.86cm][l]{\bfseries\color{caseink}#2}\color{caseink}#3%
  &\footnotesize\hspace{1.10cm}\hangindent=1.98cm\hangafter=1%
    \makebox[0.86cm][l]{\bfseries\color{caseink}#4}\songfont\color{caseink}#5%
  &#1%
  \\[0.11em]%
}
% \osubrow：二級子項，有段碼範圍（I.、II.、A.I. …）
% #1=段碼字串, #2=DE標籤, #3=DE正文, #4=ZH標籤, #5=ZH正文
\newcommand{\osubrow}[5]{%
  \footnotesize\hspace{1.40cm}\hangindent=2.30cm\hangafter=1%
    \makebox[0.90cm][l]{\bfseries\color{caseink}#2}\color{caseink}#3%
  &\footnotesize\hspace{1.40cm}\hangindent=2.30cm\hangafter=1%
    \makebox[0.90cm][l]{\bfseries\color{caseink}#4}\songfont\color{caseink}#5%
  &#1%
  \\[0.09em]%
}
% \ointrow：引入段落，有單一段碼，無標籤
% #1=段碼字串, #2=DE正文, #3=ZH正文
\newcommand{\ointrow}[3]{%
  \footnotesize\hspace{0.52cm}\hangindent=0.52cm\hangafter=1\color{caseink}#2%
  &\footnotesize\hspace{0.52cm}\hangindent=0.52cm\hangafter=1\songfont\color{caseink}#3%
  &#1%
  \\[0.09em]%
}
% \osubsubrow：三級子項（1.、2.、3.）
% #1=段碼字串, #2=DE標籤, #3=DE正文, #4=ZH標籤, #5=ZH正文
\newcommand{\osubsubrow}[5]{%
  \footnotesize\hspace{1.80cm}\hangindent=2.48cm\hangafter=1%
    \makebox[0.68cm][l]{\bfseries\color{caseink}#2}\color{caseink}#3%
  &\footnotesize\hspace{1.80cm}\hangindent=2.48cm\hangafter=1%
    \makebox[0.68cm][l]{\bfseries\color{caseink}#4}\songfont\color{caseink}#5%
  &#1%
  \\[0.08em]%
}
% \osubsubsubrow：四級子項（a）、b））
% #1=段碼字串, #2=DE標籤, #3=DE正文, #4=ZH標籤, #5=ZH正文
\newcommand{\osubsubsubrow}[5]{%
  \footnotesize\hspace{2.10cm}\hangindent=2.68cm\hangafter=1%
    \makebox[0.58cm][l]{\bfseries\color{caseink}#2}\color{caseink}#3%
  &\footnotesize\hspace{2.10cm}\hangindent=2.68cm\hangafter=1%
    \makebox[0.58cm][l]{\bfseries\color{caseink}#4}\songfont\color{caseink}#5%
  &#1%
  \\[0.07em]%
}
% \osubsubsubsubrow：五級子項（aa）、bb））
% 標籤框 0.62cm：足以容納「aa)」在 footnotesize bold 下（約 0.50cm）並留有間距
% #1=段碼字串, #2=DE標籤, #3=DE正文, #4=ZH標籤, #5=ZH正文
\newcommand{\osubsubsubsubrow}[5]{%
  \footnotesize\hspace{2.38cm}\hangindent=3.06cm\hangafter=1%
    \makebox[0.68cm][l]{\bfseries\color{caseink}#2}\color{caseink}#3%
  &\footnotesize\hspace{2.38cm}\hangindent=3.06cm\hangafter=1%
    \makebox[0.68cm][l]{\bfseries\color{caseink}#4}\songfont\color{caseink}#5%
  &#1%
  \\[0.06em]%
}
% \outlinepage：雙語對照大綱頁包裝環境（longtable 自動跨頁）
% 三欄：德文欄 + 中文欄 + 段碼欄（最右，不折行）
% 段碼欄寬度：1.80cm，足以容納最長段碼如「309–348」
% DE/ZH 欄寬：(0.87\textwidth - 2.08cm) / 2
%   驗算：2×欄 + 0.05\textwidth + 0.28cm gap + 1.80cm = 0.87tw - 2.08cm + 0.05tw + 2.08cm = 0.92tw ✓
\newenvironment{outlinepage}{%
  \vspace*{0.40cm}%
  \setlength{\LTleft}{0.04\textwidth}%
  \setlength{\LTright}{0.04\textwidth}%
  \setlength{\tabcolsep}{0pt}%
  % 大綱列距由 longtable 的 \arraystretch 與各 row macro 結尾的 \\[...em] 共同決定。
  % 這裡控制「每一列本身的表格 strut 高度」；數值越大，A./I./1. 各項之間越像有空行。
  % A3 原本用 0.62，視覺上沒有空行；A4 也沿用同值，避免切到 A4 後項目被撐開。
  % 若日後覺得大綱太擠，只微調這個值（例如 0.68），不要改各 \omainrow/\osubrow 的縮排。
  \renewcommand{\arraystretch}{0.62}%
  \begin{longtable}{%
    >{\vspace{0pt}\raggedright\arraybackslash}p{\dimexpr(0.87\textwidth-2.08cm)/2\relax}%
    @{\hspace{0.025\textwidth}}%
    !{\color{notegray}\vrule width 0.12pt}%
    @{\hspace{0.025\textwidth}}%
    >{\vspace{0pt}\raggedright\arraybackslash}p{\dimexpr(0.87\textwidth-2.08cm)/2\relax}%
    @{\hspace{0.28cm}}%
    >{\vspace{0pt}\footnotesize\normalfont\color{pendingink}\raggedright\arraybackslash}p{1.80cm}%
    @{}%
  }%
  % 首頁欄標籤：不要橫跨置中；分別放在德文欄與中文欄的實際欄位起點。
  \footnotesize\bfseries\color{sourceink}Gliederung%
  &\footnotesize\bfseries\songfont\color{sourceink}大綱%
  &%
  \\[0.36em]%
  \endfirsthead%
  % 續頁欄標籤：同樣分欄定位，以灰色細體區分；無需標注「續」。
  \footnotesize\color{pendingink}Gliederung%
  &\footnotesize\songfont\color{pendingink}大綱%
  &%
  \\[0.24em]%
  \endhead%
}{%
  \end{longtable}%
}
 之前設定；
%   預設 chinese / a4）
\ifdefined\docmode\else\def\docmode{chinese}\fi
\ifdefined\bilingualpaper\else\def\bilingualpaper{a4}\fi
\newif\ifshoworiginal
\newif\ifshowtranslation
\newif\ifbilinguallandscape
\newif\ifbilingualcolumns
\newif\ifintranslation
\newif\iffirstbilingualsection

\showoriginalfalse
\showtranslationfalse
\bilinguallandscapefalse
\bilingualcolumnsfalse
\intranslationfalse
\firstbilingualsectionfalse

\def\modechinese{chinese}
\def\modegerman{german}
\def\modebilingual{bilingual}
\def\bilingualpaperaThree{a3}
\def\bilingualpaperaFour{a4}

\ifx\docmode\modechinese
  \showtranslationtrue
\fi

\ifx\docmode\modegerman
  \showoriginaltrue
\fi

\ifx\docmode\modebilingual
  \showoriginaltrue
  \showtranslationtrue
  \bilinguallandscapetrue
  \bilingualcolumnstrue
\fi

% 編排模式說明：
% chinese  → \showoriginalfalse  \showtranslationtrue  （A4 直式，純中文）
% german   → \showoriginaltrue   \showtranslationfalse （A4 直式，純德文）
% bilingual→ \showoriginaltrue   \showtranslationtrue  \bilingualcolumnstrue（A3/A4 橫式對照）

\ifbilingualcolumns
  \ifx\bilingualpaper\bilingualpaperaFour
    \geometry{
      a4paper,
      landscape,
      left=1.25cm,
      right=2.25cm,
      top=1.25cm,
      bottom=1.75cm,
      footskip=8.5mm,
      marginparwidth=0.75cm,
      marginparsep=0.25cm
    }
  \else
    \geometry{
      a3paper,
      landscape,
      left=1.55cm,
      right=2.75cm,
      top=1.55cm,
      bottom=2.05cm,
      footskip=9.5mm,
      marginparwidth=0.9cm,
      marginparsep=0.35cm
    }
  \fi
\else
\geometry{
  a4paper,
  portrait,
  left=2.2cm,
  right=2.6cm,
  top=2.0cm,
  bottom=2.0cm,
  marginparwidth=0.8cm,
  marginparsep=0.3cm
}
\fi
% ===== 時間戳基礎設施 =====
\newcount\stamptimeval
\newcount\stampminuteval
\newcommand{\stamphh}{%
  \stamptimeval=\time
  \divide\stamptimeval by 60
  \ifnum\stamptimeval<10 0\fi
  \the\stamptimeval}
\newcommand{\stampmm}{%
  \stamptimeval=\time
  \stampminuteval=\time
  \divide\stampminuteval by 60
  \multiply\stampminuteval by 60
  \advance\stamptimeval by -\stampminuteval
  \ifnum\stamptimeval<10 0\fi
  \the\stamptimeval}
\newcommand{\stampwkde}[1]{%
  \ifcase#1Montag\or Dienstag\or Mittwoch\or Donnerstag\or Freitag\or Samstag\or Sonntag\fi}
\newcommand{\stampwkzh}[1]{%
  \ifcase#1 一\or 二\or 三\or 四\or 五\or 六\or 日\fi}
\newcommand{\stampmonthde}[1]{%
  \ifcase#1\or Januar\or Februar\or März\or April\or Mai\or Juni\or
  Juli\or August\or September\or Oktober\or November\or Dezember\fi}
\newcount\stampjdval
\newcount\stampwdval
\pgfcalendardatetojulian{\the\year-\the\month-\the\day}{\stampjdval}
\pgfcalendarjuliantoweekday{\the\stampjdval}{\stampwdval}
\newcommand{\timestampzh}{%
  \songfont 最後修改：\the\year 年\the\month 月\the\day 日%
  % （星期\stampwkzh{\the\stampwdval}）%
  % \space\stamphh:\stampmm
  }

\newcommand{\documentdatezh}{\the\year 年 \the\month 月 \the\day 日}
\newcommand{\documentdatede}{\the\day. \stampmonthde{\the\month} \the\year}
\newcommand{\bverfgetranslationnote}{%
  {\songfont 翻譯說明：初譯使用 Claude Code 與 GPT 協助迭代；仍請以德文原文為準。}%
}
\newcommand{\bverfgecourseinfo}{%
  {\songfont 課程：114 學年度第 2 學期；憲法專題研究（四）；授課老師：湯德宗教授；導讀日期：115 年 6 月 1 日。}%
}
\newcommand{\bverfgeeditorinfo}{%
  {\songfont 編者：王逸帆（R10A21126），國立台灣大學法律系研究所碩士班。}%
}
\newcommand{\bverfgetitleversionline}{%
  \ifbilingualcolumns
    Fassung \documentversion\enspace\textperiodcentered\enspace Stand: \documentdatede
    \quad
    {\songfont 版本 \documentversion\enspace\textperiodcentered\enspace 修訂日期：\documentdatezh\par}%
  \else
    \ifshowtranslation
      {\songfont 版本 \documentversion\enspace\textperiodcentered\enspace 修訂日期：\documentdatezh\par}%
    \else
      Fassung \documentversion\enspace\textperiodcentered\enspace Stand: \documentdatede\par
    \fi
  \fi
}
\newcommand{\bverfgetitlecornerinfo}{%
  \ifshowtranslation
    \vspace{0.72em}%
    \noindent\hfill
    \begin{minipage}{0.66\textwidth}
      \raggedleft\tiny\color{pendingink}\songfont
      \bverfgecourseinfo\par
      \bverfgeeditorinfo
    \end{minipage}%
  \fi
}
\newcommand{\bverfgetitlemeta}{%
  \vfill
  \begin{center}%
    \footnotesize\color{pendingink}%
    \bverfgetitleversionline
    \ifshowtranslation
      \vspace{0.28em}{\scriptsize\bverfgetranslationnote\par}%
    \fi
  \end{center}%
  \bverfgetitlecornerinfo
}

\pagestyle{fancy}
\fancyhf{}
\renewcommand{\headrulewidth}{0pt}
\renewcommand{\footrulewidth}{0pt}
\fancyfoot[C]{\thepage}
\ifbilingualcolumns
  \fancyfoot[R]{\scriptsize\color{pendingink}\timestampzh}%
\else
  \ifshoworiginal
  \else
    \fancyfoot[R]{\scriptsize\color{pendingink}\timestampzh\hspace{0.45cm}}%
  \fi
\fi
\setmainfont{Times New Roman}
\setsansfont{Helvetica Neue}
\setmonofont{Menlo}
\IfFileExists{/Users/iw/Library/Fonts/kaiu.ttf}{
  \setCJKmainfont[Path=/Users/iw/Library/Fonts/,AutoFakeBold=4]{kaiu.ttf}
}{
  \IfFontExistsTF{標楷體}{
    \setCJKmainfont[AutoFakeBold=4]{標楷體}
  }{
    \setCJKmainfont{Songti TC}
  }
}
\IfFontExistsTF{Songti TC}{
  \setCJKfamilyfont{song}{Songti TC}
  \setCJKmonofont{Songti TC}
}{
  \setCJKfamilyfont{song}[Path=/Users/iw/Library/Fonts/,AutoFakeBold=4]{kaiu.ttf}
  \setCJKmonofont[Path=/Users/iw/Library/Fonts/,AutoFakeBold=4]{kaiu.ttf}
}
\newcommand{\songfont}{\CJKfamily{song}}

% ===== 封面／目錄專用打字機字體（僅限封面、目錄頁使用，不影響正文）=====
\IfFontExistsTF{Radio Newsman}{%
  \newfontfamily\bverfgetitlede{Radio Newsman}%
}{%
  \newfontfamily\bverfgetitlede{Courier New}%
}
\IfFontExistsTF{Erikas Farbband}{%
  \newfontfamily\bverfgebodyde{Erikas Farbband}%
}{%
  \let\bverfgebodyde\bverfgetitlede
}
\IfFontExistsTF{Maoken Heavy Labourer}{%
  \setCJKfamilyfont{bverfgetitlecjk}{Maoken Heavy Labourer}%
}{%
  \setCJKfamilyfont{bverfgetitlecjk}{Songti TC}%
}
\IfFontExistsTF{Huiwen-mincho}{%
  \setCJKfamilyfont{bverfgebodycjk}{Huiwen-mincho}%
}{%
  \setCJKfamilyfont{bverfgebodycjk}{Songti TC}%
}
\newcommand{\bverfgetitlezh}{\bverfgetitlede\CJKfamily{bverfgetitlecjk}}
\newcommand{\bverfgebodyzh}{\bverfgebodyde\CJKfamily{bverfgebodycjk}}

\setstretch{1.35}
\setlist{itemsep=0.2em, topsep=0.4em}
\setlength{\parindent}{0pt}
\setlength{\parskip}{0.45em}
\raggedbottom
\setcounter{secnumdepth}{0}
\setcounter{tocdepth}{2}

\newcommand{\lawboxsourcefont}{\footnotesize}
\newcommand{\lawboxtransfont}{\footnotesize}
\newcommand{\lawboxsourcestretch}{1.02}
\newcommand{\lawboxtransstretch}{0.92}

\titleformat{\section}{\centering\Large\bfseries\songfont}{\thesection}{1em}{}
\titleformat{\subsection}{\large\bfseries\songfont}{\thesubsection}{1em}{}
\titleformat{\subsubsection}{\normalsize\bfseries\songfont}{\thesubsubsection}{1em}{}
\titleformat{\paragraph}{\normalsize\bfseries\songfont}{\theparagraph}{1em}{}
\titlespacing*{\section}{0pt}{1.8em}{1.0em}
\titlespacing*{\subsection}{0pt}{1.2em}{0.6em}
\titlespacing*{\subsubsection}{0pt}{1.0em}{0.45em}
\titlespacing*{\paragraph}{0pt}{0.6em}{0.4em}

\definecolor{noteblue}{HTML}{A9C9F5}
\definecolor{noteteal}{HTML}{AEE1D6}
\definecolor{notegray}{HTML}{D7DADF}
\definecolor{caseink}{HTML}{000000}
\definecolor{casepaper}{HTML}{FCFEFD}
\definecolor{sourcepaper}{HTML}{FAFBF8}
\definecolor{sourceline}{HTML}{D7DED8}
\definecolor{sourceink}{HTML}{000000}
\definecolor{pendingink}{HTML}{000000}

\tcbset{
  caseboxbase/.style={
    enhanced,
    breakable,
    width=0.985\linewidth,
    colback=casepaper,
    coltitle=caseink,
    fonttitle=\bfseries\songfont,
    boxrule=0.28pt,
    arc=1.2mm,
    left=2.4mm,
    right=2.4mm,
    top=1.5mm,
    bottom=1.5mm,
    boxsep=1.5mm,
    before skip=0.65em,
    after skip=0.65em,
    before upper={\setlength{\parindent}{0pt}\setlength{\parskip}{0.35em}},
  }
}

\newcommand{\bilingualstart}{}
\newcommand{\bilingualstop}{}
\ifbilingualcolumns
  \renewcommand{\bilingualstart}{\small\setlength{\columnsep}{1.25cm}\setlength{\columnseprule}{0.12pt}\columnratio{0.5,0.5}\multicolumnfootnotes\flushbottom\sloppy\interlinepenalty=80\clubpenalty=800\widowpenalty=800\displaywidowpenalty=800\begin{paracol}{2}\global\intranslationfalse\global\firstbilingualsectiontrue{\footnotesize\color{sourceink}Deutsch\par}\switchcolumn{\footnotesize\songfont\color{sourceink}中文譯文\par}\switchcolumn*}
  \renewcommand{\bilingualstop}{\ifintranslation\switchcolumn*\global\intranslationfalse\fi\end{paracol}\clearpage\raggedbottom}
\fi
\newif\ifrandnumbers
\randnumbersfalse
\newcounter{randnumber}
\newcounter{randnumberlocal}
\newcounter{lastsourcerandstart}
\newif\iflastsourcehasrandnumbers
\lastsourcehasrandnumbersfalse
\newif\ifinlawbox
\inlawboxfalse
\newif\iflawboxhaspendingrn
\lawboxhaspendingrnfalse
\newcommand{\lawboxpendingrn}{}
\newcommand{\startrandnumbers}{\global\randnumberstrue\setcounter{randnumber}{1}\global\lastsourcehasrandnumbersfalse}
\newcommand{\stoprandnumbers}{\global\randnumbersfalse\global\lastsourcehasrandnumbersfalse}
\newlength{\randnumberwidth}
\setlength{\randnumberwidth}{0.95cm}
\newlength{\randnumberrightoffset}
\setlength{\randnumberrightoffset}{0.85cm}
\newlength{\randnumberlawboxrightoffset}
\setlength{\randnumberlawboxrightoffset}{1.40cm}
\newcommand{\randnumberbox}[1]{%
  \leavevmode
  \makebox[0pt][l]{%
    \smash{%
      \tikz[remember picture,overlay,baseline=(rnmark.base)]{%
        \coordinate (rnmark) at (0,0);%
        \ifinlawbox
          \path let \p1=(current page.east), \p2=(rnmark.base) in
            node[
              anchor=base east,
              inner sep=0pt,
              outer sep=0pt,
              text width=\randnumberwidth,
              align=right,
              text=pendingink,
              font=\normalfont\mdseries\scriptsize\ttfamily
            ] at (\x1-\randnumberlawboxrightoffset,\y2) {{\normalfont\mdseries\scriptsize\ttfamily #1}};%
        \else
          \path let \p1=(current page.east), \p2=(rnmark.base) in
            node[
              anchor=base east,
              inner sep=0pt,
              outer sep=0pt,
              text width=\randnumberwidth,
              align=right,
              text=pendingink,
              font=\normalfont\mdseries\scriptsize\ttfamily
            ] at (\x1-\randnumberrightoffset,\y2) {{\normalfont\mdseries\scriptsize\ttfamily #1}};%
        \fi
      }%
    }%
  }%
  \ignorespaces
}
\newcommand{\lawboxstorerandnumber}[1]{%
  \gdef\lawboxpendingrn{#1}%
  \global\lawboxhaspendingrntrue
  \ignorespaces
}
\newcommand{\lawboxuserandnumber}{%
  \iflawboxhaspendingrn
    \randnumberbox{\lawboxpendingrn}%
    \global\lawboxhaspendingrnfalse
  \fi
}
\newcommand{\sourcern}[1]{%
  \ifrandnumbers
    \ifshoworiginal
      \ifshowtranslation\else
        \ifinlawbox\lawboxstorerandnumber{#1}\else\randnumberbox{#1}\fi
      \fi
    \fi
  \fi
}
\newcommand{\transrn}[1]{%
  \ifrandnumbers
    \ifshowtranslation
      \ifinlawbox\lawboxstorerandnumber{#1}\else\randnumberbox{#1}\fi
    \fi
  \fi
}
% ===== 課堂模式：德文原文縮寫註解 =====
% 日常切換只改各判決檔的一行：
%   \classroommodeon        開啟課堂版；註解掉這行即回到一般版。
% 模板預設：
%   1. 課堂模式關閉。
%   2. 課堂模式開啟時，縮寫註解包含「德文全稱＋中文譯名」。
%   3. 課堂模式開啟時，GG/條文編者註仍會自動顯示。
% 進階微調才需要在單案檔覆寫：
%   \classroomabbrzhoff     縮寫註解僅顯示德文全稱。
%   \showeditornotestrue    一般版也顯示編者註。
% 註解策略：同一實際 PDF 頁面中，同一縮寫只在第一次出現處加註。
% 這裡用 zref-abspage 的絕對頁序，不用 \thepage，避免文件內頁碼重置時誤判。
\newif\ifclassroommode
\newif\ifclassroomabbrzh
\classroommodefalse
\classroomabbrzhtrue
\newcommand{\classroommodeon}{\classroommodetrue}
\newcommand{\classroommodeoff}{\classroommodefalse}
\newcommand{\classroomabbrzhon}{\classroomabbrzhtrue}
\newcommand{\classroomabbrzhoff}{\classroomabbrzhfalse}
\newcommand{\abbrnotelabel}{\emph{Abk.}: }
\newcommand{\abbrnote}[3]{%
  \ifclassroommode
    \ifshoworiginal
      \ifintranslation\else
        \footnote{\normalfont\footnotesize\abbrnotelabel #2\ifclassroomabbrzh；{\songfont #3}\fi}%
      \fi
    \fi
  \fi
}
\newcommand{\abbrnoteonce}[4]{%
  #2%
  \ifclassroommode
    \ifshoworiginal
      \ifintranslation\else
        \begingroup
          \edef\abbrcurrentabspage{\number\numexpr\value{abspage}+1\relax}%
          \ifcsname abbrnoteseen@#1@\abbrcurrentabspage\endcsname\else
            \global\expandafter\let\csname abbrnoteseen@#1@\abbrcurrentabspage\endcsname\empty
            \abbrnote{#1}{#3}{#4}%
          \fi
        \endgroup
      \fi
    \fi
  \fi
}
\newcommand{\abbrAbs}{\abbrnoteonce{abs}{Abs.}{Absatz}{項}}
\newcommand{\abbrAAO}{\abbrnoteonce{aao}{a.a.O.}{am angegebenen Ort}{前引處}}
\newcommand{\abbrAE}{\abbrnoteonce{ae}{AE}{Alternativ-Entwurf}{替代草案}}
\newcommand{\abbrArt}{\abbrnoteonce{art}{Art.}{Artikel}{條}}
\newcommand{\abbrAufl}{\abbrnoteonce{aufl}{Aufl.}{Auflage}{版}}
\newcommand{\abbrBd}{\abbrnoteonce{bd}{Bd.}{Band}{卷}}
\newcommand{\abbrBGBl}{\abbrnoteonce{bgbl}{BGBl.}{Bundesgesetzblatt}{聯邦法律公報}}
\newcommand{\abbrBGHSt}{\abbrnoteonce{bghst}{BGHSt}{Entscheidungen des Bundesgerichtshofes in Strafsachen}{聯邦普通法院刑事判決彙編}}
\newcommand{\abbrBRDrucks}{\abbrnoteonce{brdrucks}{BRDrucks.}{Bundesrats-Drucksache}{聯邦參議院文件}}
\newcommand{\abbrBTDrucks}{\abbrnoteonce{btdrucks}{BTDrucks.}{Bundestags-Drucksache}{聯邦議院文件}}
\newcommand{\abbrBundesgesetzbl}{\abbrnoteonce{bundesgesetzbl}{Bundesgesetzbl.}{Bundesgesetzblatt}{聯邦法律公報}}
\newcommand{\abbrBvF}{\abbrnoteonce{bvf}{BvF}{Bundesverfassungsgerichtsverfahren}{聯邦憲法法院程序案號}}
\newcommand{\abbrBvL}{\abbrnoteonce{bvl}{BvL}{Bundesverfassungsgerichtsverfahren}{聯邦憲法法院程序案號}}
\newcommand{\abbrBvR}{\abbrnoteonce{bvr}{BvR}{Bundesverfassungsgerichtsverfahren}{聯邦憲法法院程序案號}}
\newcommand{\abbrBVerfG}{\abbrnoteonce{bverfg}{BVerfG}{Bundesverfassungsgericht}{聯邦憲法法院}}
\newcommand{\abbrBVerfGE}{\abbrnoteonce{bverfge}{BVerfGE}{Entscheidungen des Bundesverfassungsgerichts}{聯邦憲法法院判決集}}
\newcommand{\abbrBVerfGG}{\abbrnoteonce{bverfgg}{BVerfGG}{Bundesverfassungsgerichtsgesetz}{聯邦憲法法院法}}
\newcommand{\abbrDDR}{\abbrnoteonce{ddr}{DDR}{Deutsche Demokratische Republik}{德意志民主共和國；東德}}
\newcommand{\abbrDJT}{\abbrnoteonce{djt}{DJT}{Deutscher Juristentag}{德國法學家大會}}
\newcommand{\abbrDr}{\abbrnoteonce{dr}{Dr.}{Doktor}{Dr.}}
\newcommand{\abbrEuGRZ}{\abbrnoteonce{eugrz}{EuGRZ}{Europäische Grundrechte-Zeitschrift}{歐洲基本權期刊}}
\newcommand{\abbrEV}{\abbrnoteonce{ev}{EV}{Einigungsvertrag}{統一條約}}
\newcommand{\abbrF}{\abbrnoteonce{f}{f.}{folgende}{以下一頁；下一段}}
\newcommand{\abbrFF}{\abbrnoteonce{ff}{ff.}{fortfolgende}{以下數頁；以下各段}}
\newcommand{\abbrGBlDDR}{\abbrnoteonce{gblddr}{GBl. DDR}{Gesetzblatt der Deutschen Demokratischen Republik}{德意志民主共和國法律公報}}
\newcommand{\abbrGG}{\abbrnoteonce{gg}{GG}{Grundgesetz}{基本法}}
\newcommand{\abbrGVBl}{\abbrnoteonce{gvbl}{GVBl.}{Gesetz- und Verordnungsblatt}{法律及命令公報}}
\newcommand{\abbrJR}{\abbrnoteonce{jr}{JR}{Juristische Rundschau}{法學評論}}
\newcommand{\abbrJZ}{\abbrnoteonce{jz}{JZ}{JuristenZeitung}{法學家報}}
\newcommand{\abbrIVm}{\abbrnoteonce{ivm}{i.V.m.}{in Verbindung mit}{結合；連同}}
\newcommand{\abbrMdB}{\abbrnoteonce{mdb}{MdB}{Mitglied des Bundestages}{聯邦議院議員}}
\newcommand{\abbrMwN}{\abbrnoteonce{mwn}{m.w.N.}{mit weiteren Nachweisen}{附進一步文獻／判例出處}}
\newcommand{\abbrNF}{\abbrnoteonce{nf}{n.F.}{neue Fassung}{新版本}}
\newcommand{\abbrAF}{\abbrnoteonce{af}{a.F.}{alte Fassung}{舊版本}}
\newcommand{\abbrNr}{\abbrnoteonce{nr}{Nr.}{Nummer}{款、目或號，依語境}}
\newcommand{\abbrProf}{\abbrnoteonce{prof}{Prof.}{Professor}{教授}}
\newcommand{\abbrRdnr}{\abbrnoteonce{rdnr}{Rdnr.}{Randnummer}{邊碼；段碼}}
\newcommand{\abbrRdnrn}{\abbrnoteonce{rdnrn}{Rdnrn.}{Randnummern}{邊碼；段碼}}
\newcommand{\abbrRGBl}{\abbrnoteonce{rgbl}{RGBl.}{Reichsgesetzblatt}{帝國法律公報}}
\newcommand{\abbrRGSt}{\abbrnoteonce{rgst}{RGSt}{Entscheidungen des Reichsgerichts in Strafsachen}{帝國法院刑事判決彙編}}
\newcommand{\abbrRspr}{\abbrnoteonce{rspr}{Rspr.}{Rechtsprechung}{判例；司法實務}}
\newcommand{\abbrRVO}{\abbrnoteonce{rvo}{RVO}{Reichsversicherungsordnung}{帝國保險法}}
\newcommand{\abbrS}{\abbrnoteonce{s}{S.}{Seite}{頁}}
\newcommand{\abbrSFHG}{\abbrnoteonce{sfhg}{SFHG}{Schwangeren- und Familienhilfegesetz}{孕婦及家庭扶助法}}
\newcommand{\abbrSGBV}{\abbrnoteonce{sgbv}{SGB V}{Fünftes Buch Sozialgesetzbuch}{社會法典第五編}}
\newcommand{\abbrStAG}{\abbrnoteonce{staeg}{StÄG}{Strafrechtsänderungsgesetz}{刑法修正法}}
\newcommand{\abbrStenBer}{\abbrnoteonce{stenber}{StenBer.}{Stenographischer Bericht}{速記紀錄}}
\newcommand{\abbrStGB}{\abbrnoteonce{stgb}{StGB}{Strafgesetzbuch}{刑法}}
\newcommand{\abbrStREG}{\abbrnoteonce{streg}{StREG}{Strafrechtsreform-Ergänzungsgesetz}{刑法改革補充法}}
\newcommand{\abbrStrRG}{\abbrnoteonce{strrg}{StrRG}{Strafrechtsreformgesetz}{刑法改革法}}
\newcommand{\abbrUS}{\abbrnoteonce{us}{U.S.}{United States Reports}{美國最高法院判決彙編}}
\newcommand{\abbrVgl}{\abbrnoteonce{vgl}{vgl.}{vergleiche}{參見}}
\newcommand{\abbrVVDStRL}{\abbrnoteonce{vvdstrl}{VVDStRL}{Veröffentlichungen der Vereinigung der Deutschen Staatsrechtslehrer}{德國公法教師協會論文集}}
\newcommand{\abbrWP}{\abbrnoteonce{wp}{WP.}{Wahlperiode}{屆期}}
\newcommand{\abbrWp}{\abbrnoteonce{wp}{Wp.}{Wahlperiode}{屆期}}
\newcommand{\abbrZStrW}{\abbrnoteonce{zstrw}{ZStrW}{Zeitschrift für die gesamte Strafrechtswissenschaft}{整體刑法學期刊}}
\newcommand{\recordsourcerandstart}{%
  \ifrandnumbers
    \setcounter{lastsourcerandstart}{\value{randnumber}}%
    \global\lastsourcehasrandnumberstrue
  \else
    \global\lastsourcehasrandnumbersfalse
  \fi
}
\newlength{\biparagraphskip}
\setlength{\biparagraphskip}{0.52em}
\newlength{\lawboxblockbeforeskip}
\setlength{\lawboxblockbeforeskip}{0.72em}
\newlength{\lawboxblockafterskip}
\setlength{\lawboxblockafterskip}{0.92em}
\newif\iflawboxpartendedblock
\lawboxpartendedblockfalse
\newcommand{\sourceparstyle}{\footnotesize\color{sourceink}\setstretch{1.12}\RaggedRight\emergencystretch=3em\setlength{\parskip}{0.42em}\obeylines\selectfont}
\newcommand{\transparstyle}{\songfont\color{caseink}\setstretch{1.12}\setlength{\parindent}{0pt}\setlength{\parskip}{0.42em}}
\newcommand{\bilingualpairskip}{%
  \par\addvspace{\biparagraphskip}%
  \switchcolumn
  \par\addvspace{\biparagraphskip}%
  \switchcolumn*%
  \global\intranslationfalse
}
\newcommand{\bilingualboxpairskip}{%
  \par
  \switchcolumn
  \par
  \switchcolumn*%
  \global\intranslationfalse
}
\newcommand{\bilinguallawboxblockskip}{%
  \switchcolumn*%
  \par\addvspace{\lawboxblockafterskip}%
  \switchcolumn
  \par\addvspace{\lawboxblockafterskip}%
  \switchcolumn*%
  \global\intranslationfalse
}
\newcommand{\bikeepwithnext}[1][4]{%
  \ifbilingualcolumns
    \syncleft
    \Needspace{#1\baselineskip}%
    \switchcolumn
    \Needspace{#1\baselineskip}%
    \switchcolumn*%
    \global\intranslationfalse
  \else
    \Needspace{#1\baselineskip}%
  \fi
}
\NewEnviron{sourcepara}[1][]{
  \recordsourcerandstart
  \ifshoworiginal
    \ifbilingualcolumns
      \ifintranslation\switchcolumn*\global\intranslationfalse\fi
    \else
      \Needspace{5\baselineskip}
    \fi
    \ifbilingualcolumns\else\par\addvspace{0.35em}\fi
    {\sourceparstyle
    \noindent\BODY\par}
    \ifbilingualcolumns\else\addvspace{0.25em}\fi
    \ifbilingualcolumns
      \switchcolumn
      \global\intranslationtrue
    \fi
  \else
    \ifrandnumbers
      \begingroup
        \setbox0=\vbox{\hsize=\linewidth\sourceparstyle\noindent\BODY\par}%
      \endgroup
    \fi
  \fi
}
\NewEnviron{transpara}{
  \ifshowtranslation
    \setcounter{randnumberlocal}{\value{lastsourcerandstart}}%
    \ifbilingualcolumns
      \ifintranslation\else\switchcolumn\global\intranslationtrue\fi
      {\transparstyle\setlength{\leftskip}{0.55em}\setlength{\rightskip}{0.15em plus 1em}\addtolength{\linewidth}{-0.55em}\BODY\par}
      \switchcolumn*
      \bilingualpairskip
    \else
      {\transparstyle\BODY\par}
    \fi
  \else
    \ifbilingualcolumns
      \ifintranslation\switchcolumn*\global\intranslationfalse\fi
    \fi
  \fi
}
\NewEnviron{sourceboxpara}[1][]{
  \global\lawboxpartendedblockfalse
  \recordsourcerandstart
  \ifshoworiginal
    \ifbilingualcolumns
      \ifintranslation\switchcolumn*\global\intranslationfalse\fi
    \else
      \Needspace{3\baselineskip}
    \fi
    {\sourceparstyle
    \BODY\par}
    \ifbilingualcolumns\else
      \iflawboxpartendedblock\addvspace{\lawboxblockafterskip}\fi
    \fi
    \ifbilingualcolumns
      \switchcolumn
      \global\intranslationtrue
    \fi
  \fi
}
\NewEnviron{transboxpara}{
  \global\lawboxpartendedblockfalse
  \ifshowtranslation
    \setcounter{randnumberlocal}{\value{lastsourcerandstart}}%
    \ifbilingualcolumns
      \ifintranslation\else\switchcolumn\global\intranslationtrue\fi
      {\transparstyle\BODY\par}
      \iflawboxpartendedblock
        \bilinguallawboxblockskip
      \else
        \switchcolumn*
        \bilingualboxpairskip
      \fi
    \else
      {\transparstyle\BODY\par}
      \iflawboxpartendedblock\addvspace{\lawboxblockafterskip}\fi
    \fi
  \else
    \ifbilingualcolumns
      \ifintranslation\switchcolumn*\global\intranslationfalse\fi
    \fi
  \fi
}
\newcommand{\leitsatznum}[1]{\noindent\hangindent=3.0em\hangafter=1\makebox[2.25em][r]{\bfseries #1.}\hspace{0.55em}\ignorespaces}
\newcommand{\syncleft}{\ifbilingualcolumns\ifintranslation\switchcolumn*\global\intranslationfalse\fi\fi}
\pretocmd{\section}{\syncleft\clearpage}{}{}
\pretocmd{\subsection}{\syncleft}{}{}
\pretocmd{\subsubsection}{\syncleft}{}{}
\newcommand{\biheadingfont}{\songfont}
\newcommand{\bitocentry}[2]{%
  \ifshoworiginal
    \ifshowtranslation
      \ifstrequal{#1}{#2}{#1}{#1\space #2}%
    \else
      #1%
    \fi
  \else
    #2%
  \fi
}
\pdfstringdefDisableCommands{%
  \def\bitocentry#1#2{#1}%
}
\newcommand{\biheading}[7]{%
  \syncleft#1\bikeepwithnext[#3]\phantomsection\addcontentsline{toc}{#2}{\bitocentry{#6}{#7}}%
  \ifbilingualcolumns
    {#4 #6\par}%
    \switchcolumn
    {#4\biheadingfont #7\par}%
    \switchcolumn*%
    \global\intranslationfalse
    \par\addvspace{#5}%
    \switchcolumn
    \par\addvspace{#5}%
    \switchcolumn*%
  \else
    \ifshoworiginal{#4 #6\par}\fi
    \ifshowtranslation{#4\biheadingfont #7\par}\fi
    \addvspace{#5}%
  \fi
}
\newcommand{\termsectionheading}[1]{%
  \par\Needspace{9\baselineskip}%
  \vspace{0.65em}%
  {\songfont\bfseries #1\par}%
  \vspace{0.2em}%
}
% 章節換頁：
%   雙語 paracol 欄內不要用 \clearpage；它會在同步兩欄時吐出只有欄線/頁碼的空頁。
%   \newpage 足以讓下一個 section 起新頁，又不會強制清空所有待排材料。
%   單語模式不在 paracol 內，仍可用 \clearpage。
\newcommand{\bisectionpagebreak}{\ifbilingualcolumns\newpage\else\clearpage\fi}
\newcommand{\bisection}[2]{%
  \biheading{\iffirstbilingualsection\global\firstbilingualsectionfalse\else\bisectionpagebreak\fi}{section}{6}{\centering\Large\bfseries}{0.8em}{#1}{#2}%
}
\newcommand{\bisubsection}[2]{%
  \biheading{}{subsection}{5}{\large\bfseries}{0.55em}{#1}{#2}%
}
\newcommand{\bisubsubsection}[2]{%
  \biheading{}{subsubsection}{4}{\normalsize\bfseries}{0.45em}{#1}{#2}%
}
\newcommand{\bicompositeheading}[4]{%
  \biheading{}{subsection}{5}{\large\bfseries}{0.16em}{#1}{#2}%
  \biheading{}{subsubsection}{3}{\normalsize\bfseries}{0.45em}{#3}{#4}%
}
\newif\iflawboxfirstitem
\lawboxfirstitemfalse
\newcommand{\lawboxprespace}[1]{%
  \par
  \iflawboxfirstitem
    \global\lawboxfirstitemfalse
  \else
    \addvspace{#1}%
  \fi
}
\newtcolorbox{lawboxinner}{
  caseboxbase,
  colframe=sourceline!85!black,
  colback=white,
  left=1.05em,
  right=1.05em,
  top=0.62em,
  bottom=0.62em,
  before upper={\global\inlawboxtrue\global\lawboxfirstitemtrue\global\lawboxhaspendingrnfalse\ifintranslation\lawboxtransfont\else\lawboxsourcefont\fi\RaggedRight\setlength{\parindent}{0pt}\setlength{\parskip}{0pt}\ifintranslation\setstretch{\lawboxtransstretch}\else\setstretch{\lawboxsourcestretch}\fi},
  after upper={\global\lawboxhaspendingrnfalse\global\inlawboxfalse}
}
\tcbset{
  lawboxpartbase/.style={
    enhanced,
    breakable=false,
    width=0.985\linewidth,
    colback=white,
    frame hidden,
    boxrule=0pt,
    arc=0pt,
    left=1.05em,
    right=1.05em,
    boxsep=1.5mm,
    before skip=0pt,
    after skip=0pt,
    before upper={\global\inlawboxtrue\global\lawboxfirstitemtrue\global\lawboxhaspendingrnfalse\ifintranslation\lawboxtransfont\else\lawboxsourcefont\fi\RaggedRight\setlength{\parindent}{0pt}\setlength{\parskip}{0pt}\ifintranslation\setstretch{\lawboxtransstretch}\else\setstretch{\lawboxsourcestretch}\fi},
    after upper={\global\lawboxhaspendingrnfalse\global\inlawboxfalse},
    borderline west={0.28pt}{0pt}{sourceline!85!black},
    borderline east={0.28pt}{0pt}{sourceline!85!black},
  },
  lawboxpartfirst/.style={
    lawboxpartbase,
    top=0.62em,
    bottom=0.04em,
    before skip=\lawboxblockbeforeskip,
    borderline north={0.28pt}{0pt}{sourceline!85!black},
  },
  lawboxpartmiddle/.style={
    lawboxpartbase,
    top=0.04em,
    bottom=0.04em,
  },
  lawboxpartlast/.style={
    lawboxpartbase,
    top=0.04em,
    bottom=0.62em,
    after skip=0pt,
    borderline south={0.28pt}{0pt}{sourceline!85!black},
  }
}
\newtcolorbox{lawboxpartinner}[1]{#1}
\newtcolorbox{termboxinner}{
  caseboxbase,
  colframe=noteblue!70!black,
  colback=casepaper,
  before upper={\songfont\setlength{\parindent}{0pt}\setlength{\parskip}{0.35em}}
}
\newtcolorbox{deboxinner}{
  caseboxbase,
  colframe=notegray!72!black,
  colback=white,
  left=1.05em,
  right=1.05em,
  top=0.62em,
  bottom=0.62em,
  before upper={\global\inlawboxtrue\global\lawboxfirstitemtrue\global\lawboxhaspendingrnfalse\small\setlength{\parindent}{0pt}\setlength{\parskip}{0.35em}},
  after upper={\global\lawboxhaspendingrnfalse\global\inlawboxfalse}
}
\NewEnviron{lawbox}{
  \begin{lawboxinner}\BODY\end{lawboxinner}
}
\NewEnviron{lawboxpart}[2]{
  \begin{lawboxpartinner}{lawboxpart#1,equal height group=#2}\BODY\end{lawboxpartinner}%
  \ifstrequal{#1}{last}{\global\lawboxpartendedblocktrue}{}%
}
\NewEnviron{termbox}{
  \par\addvspace{0.55em}
  {\color{noteblue!70!black}\rule{\linewidth}{0.28pt}}\par
  \begingroup
    \songfont\small\color{caseink}
    \setlength{\parindent}{0pt}
    \setlength{\parskip}{0.24em}
    \BODY
    \par
  \endgroup
  {\color{noteblue!70!black}\rule{\linewidth}{0.28pt}}\par
  \addvspace{0.55em}
}
\NewEnviron{debox}{
  \begin{deboxinner}\BODY\end{deboxinner}
}
\providecommand{\paragraphmark}[1]{}
\newcommand{\lawboxtitleafter}{\addvspace{0.06em}}
\newcommand{\lawtitle}[1]{\lawboxprespace{0.22em}{\lawboxtransfont\bfseries\lawboxuserandnumber #1}\par\lawboxtitleafter}
\newcommand{\absno}[2]{\lawboxprespace{0.04em}\noindent\lawboxuserandnumber\hangindent=2.6em\hangafter=1\makebox[2.6em][l]{(#1)}#2\par}
\newcommand{\nrno}[2]{\lawboxprespace{0pt}\noindent\lawboxuserandnumber\hspace*{2.6em}\hangindent=4.4em\hangafter=1\makebox[1.8em][l]{#1.}#2\par}
\newcommand{\letterno}[2]{\lawboxprespace{0pt}\noindent\lawboxuserandnumber\hspace*{4.4em}\hangindent=6.0em\hangafter=1\makebox[1.6em][l]{#1)}#2\par}
\newcommand{\sourcelawtitle}[1]{\lawboxprespace{0.22em}{\lawboxsourcefont\bfseries\lawboxuserandnumber #1}\par\lawboxtitleafter}
\newcommand{\sourcelawsubtitle}[1]{\lawboxprespace{0.04em}{\lawboxsourcefont\bfseries\lawboxuserandnumber #1}\par\lawboxtitleafter}
\newcommand{\sourceabsno}[2]{\lawboxprespace{0.04em}\noindent\lawboxuserandnumber\hangindent=2.45em\hangafter=1\makebox[2.45em][l]{(#1)}#2\par}
\newcommand{\sourcenrno}[2]{\lawboxprespace{0pt}\noindent\lawboxuserandnumber\hspace*{2.45em}\hangindent=4.25em\hangafter=1\makebox[1.8em][l]{#1.}#2\par}
\newcommand{\sourceletterno}[2]{\lawboxprespace{0pt}\noindent\lawboxuserandnumber\hspace*{4.25em}\hangindent=5.85em\hangafter=1\makebox[1.6em][l]{#1)}#2\par}
\newcommand{\sourcequotepara}[1]{\lawboxprespace{0.04em}\noindent\lawboxuserandnumber\hangindent=1.2em\hangafter=1 #1\par}
\newcommand{\sourcelawline}[1]{\lawboxprespace{0.04em}\noindent\lawboxuserandnumber #1\par}
\newcommand{\lawline}[1]{\lawboxprespace{0.04em}\noindent\lawboxuserandnumber #1\par}
\newcommand{\pendingtranslation}{\par{\songfont\footnotesize\color{pendingink}（中文譯文待補。）}\par}
\newcommand{\tocpart}[1]{\par\addvspace{0.52em}{\footnotesize\bfseries\color{pendingink}#1\par}\addvspace{0.16em}}
\newcommand{\tocdashline}{\par\addvspace{0.48em}\noindent{\color{notegray}\xleaders\hbox to 0.82em{\hfil\rule{0.42em}{0.28pt}\hfil}\hskip\linewidth}\par\addvspace{0.38em}}
% \tocpanel{font-cmd}{content}：將目錄置中於所在欄寬（雙語欄適用）。
% A3 橫式使用 0.58\linewidth，A4 橫式使用 0.84\linewidth，
% 使兩種紙張下目錄欄內容實際寬度相近（均約 100–105 mm）。
\newcommand{\tocpanel}[2]{%
  \makebox[\linewidth][c]{%
    \ifx\bilingualpaper\bilingualpaperaThree
      \begin{minipage}{0.58\linewidth}%
    \else
      \begin{minipage}{0.84\linewidth}%
    \fi
    \toccolstyle
    #1#2%
    \end{minipage}%
  }%
}
% ===== 首頁簡目錄的兩層 description list 縮排 =====
% enumitem 的 description list 可把每一列拆成：
%   [label]  labelsep  body text
% 這裡的「起點」都以所在 minipage/欄位的左緣為 0。
%
% 核心規則：
%   labelindent = label 欄位從左緣開始的位置。
%   labelwidth  = label 欄位本身寬度；標籤太長（如 Abw. Meinung）就加大它。
%   labelsep    = label 欄位右緣到正文的距離。
%   leftmargin  = 正文 body text 的起點。判斷層級視覺時主要看這個。
%
% 所以：
%   想讓「Begründung / 判決理由」整列正文往右或往左：改 tocitems 的 leftmargin。
%   想讓「A. Verfahren und Gegenstand / A. 程序與審查標的」正文往右或往左：
%     改 tocsubitems 的 leftmargin。
%   想讓 A./B./C. 這些標籤本身往右或往左，但正文不動：改 tocsubitems 的 labelindent。
%   想讓 A./B./C. 與正文間距變大或變小：改 tocsubitems 的 labelsep。
%
% 層級原則：
%   tocsubitems.leftmargin 必須大於 tocitems.leftmargin，
%   否則子層級正文會看起來比「Gründe / 理由」還靠左。
\newlist{tocitems}{description}{1}
\setlist[tocitems]{%
  labelindent=0pt,    % 第一層 label 欄位起點；0pt 表示貼齊目錄欄左緣。
  leftmargin=2.54cm,  % 第一層正文起點；例如 Begründung / 判決理由 的左緣。
  labelwidth=2.16cm,  % 第一層 label 欄位寬度；需容納 Leitsätze、Abw. Meinung、不同意見等。
  labelsep=0.32cm,    % 第一層 label 與正文的水平間距；只改間距，不改正文起點。
  itemsep=0.28em,     % 第一層各項之間的垂直距離。
  topsep=0.20em,      % 第一層 list 上下與前後內容的垂直距離。
  parsep=0pt,         % 同一 item 內段落之間的距離；目錄項通常不需要。
  partopsep=0pt,      % list 若緊接段落開始時的額外距離；關閉以免目錄高度飄動。
  align=left,         % label 欄內左對齊；長短標籤不靠右貼齊。
  font=\normalfont\bfseries % label 的字型；正文不受這行影響。
}
\newlist{tocsubitems}{description}{1}
\setlist[tocsubitems]{%
  labelindent=1cm,    % 第二層 label 起點；只控制 A./B./C. 本身的位置。
  leftmargin=3.00cm,  % 第二層正文起點；必須大於 2.54cm，才會比 Begründung / 判決理由更右。
  labelwidth=0.5cm,   % 第二層 label 欄寬；A./B./C. 很短，可小於第一層。
  labelsep=1.5cm,    % 第二層 label 與正文間距；A. 與 Verfahren / 程序 的距離看這裡。
  itemsep=0.12em,     % 第二層各項較密，避免 A.–E. 佔太高。
  topsep=0.04em,      % 第二層 list 與 Gründe / 理由之間的垂直距離。
  parsep=0pt,         % 同一第二層 item 內段落距離；目錄項通常不需要。
  partopsep=0pt,      % 關閉額外段落起始距離，避免版面不穩。
  align=left,         % A./B./C. label 欄內左對齊。
  font=\normalfont\bfseries, % A./B./C. label 加粗；正文不受這行影響。
  before={}           % 不在第二層 list 前自動插入額外內容。
}
\newlist{termitems}{description}{1}
\ifbilingualcolumns
  \setlist[termitems]{%
    leftmargin=5.2cm,
    labelwidth=4.8cm,
    labelsep=0.35cm,
    itemsep=0.16em,
    topsep=0.2em,
    parsep=0pt,
    partopsep=0pt,
    font=\normalfont\bfseries,
    style=nextline
  }
\else
  \setlist[termitems]{%
    leftmargin=0pt,
    labelwidth=0pt,
    labelsep=0pt,
    itemsep=0.24em,
    topsep=0.2em,
    parsep=0pt,
    partopsep=0pt,
    font=\normalfont\bfseries,
    style=nextline
  }
\fi
\newlist{abbritems}{description}{1}
\setlist[abbritems]{%
  leftmargin=2.38cm,
  labelwidth=2.02cm,
  labelsep=0.28cm,
  itemsep=0.16em,
  topsep=0.2em,
  parsep=0pt,
  partopsep=0pt,
  align=left,
  font=\normalfont\bfseries
}
\ifbilingualcolumns
  \newenvironment{termcolumns}{\begin{multicols}{2}}{\end{multicols}}
\else
  \newenvironment{termcolumns}{}{}
\fi

% ===== 編者註腳開關 =====
% 一般版預設不顯示編者註，避免正文腳註過密。
% 若開啟 \classroommodeon，GG/條文編者註仍會自動顯示，無需另開本開關。
% 若一般版也要顯示編者註，可在單案檔 \input 後加 \showeditornotestrue。
\newif\ifshoweditornotes
\showeditornotesfalse

% ===== 目錄排版輔助 =====
\newcommand{\toccolstyle}{\raggedright\arraybackslash}
\newcommand{\tocsingleindent}{\leftskip=1.45cm\rightskip=1.2cm}

% ===== 行內引文環境（德文原文與中文譯文各一，帶左側灰色邊線）=====
\newcommand{\quoteparbreak}{\par\addvspace{0.48em}}
\newcommand{\sourcequoteinner}[1]{%
  \begin{tcolorbox}[
    enhanced, breakable, width=\dimexpr\linewidth-0.8em\relax, frame hidden,
    colback=white, coltext=caseink, boxrule=0pt,
    borderline west={0.55pt}{0.88em}{notegray},
    left=1.78em, right=0.72em, top=0.18em, bottom=0.18em,
    boxsep=0pt, before skip=0.24em, after skip=0.24em,
    before upper={\normalfont\footnotesize\color{caseink}\setlength{\parindent}{0pt}\setlength{\parskip}{0.34em}\raggedright\setlength{\rightskip}{0pt plus 1em}}
  ]%
  #1\par
  \end{tcolorbox}%
}
\newcommand{\transquoteinner}[1]{%
  \begin{tcolorbox}[
    enhanced, breakable, width=\dimexpr\linewidth-0.8em\relax, frame hidden,
    colback=white, coltext=caseink, boxrule=0pt,
    borderline west={0.55pt}{0.88em}{notegray},
    left=1.78em, right=0.72em, top=0.18em, bottom=0.18em,
    boxsep=0pt, before skip=0.24em, after skip=0.24em,
    before upper={\songfont\small\color{caseink}\setlength{\parindent}{0pt}\setlength{\parskip}{0.34em}\raggedright\setlength{\rightskip}{0pt plus 1em}}
  ]%
  #1\par
  \end{tcolorbox}%
}

% ===== 大綱（Gliederung）Rn 輔助命令 =====
\newcommand{\outrn}[2]{\enspace{\normalfont\footnotesize\color{pendingink}(Rn\,#1–#2)}}
\newcommand{\outrnsingle}[1]{\enspace{\normalfont\footnotesize\color{pendingink}(Rn\,#1)}}
\newcommand{\outrnzh}[2]{\enspace{\normalfont\footnotesize\color{pendingink}（第\,#1–#2\,段）}}
\newcommand{\outrnzhs}[1]{\enspace{\normalfont\footnotesize\color{pendingink}（第\,#1\,段）}}
% ===== 大綱雙語對照表格輔助命令（三欄：段碼 │ 德文 │ 中文）=====
% 欄 1：段碼（由欄定義自動格式化：小字、等寬、右對齊、pendingink）
% 欄 2：德文標籤＋正文
% 欄 3：中文標籤＋正文
%
% 無段碼的列（\omainrow、\osecrow），欄 1 以 & 空置。
% 有段碼的列（\osecrowrn、\osubrow、\ointrow），#1 為預格式化字串
%   例：\osubrow{2–49}{I.}{...}{I.}{...}
%       \ointrow{107}{...}{...}
%
% \opartrow：段組標題（橫跨三欄）
\newcommand{\opartrow}[2]{%
  \noalign{\vspace{0.30em}}%
  \multicolumn{3}{@{}l@{}}{%
    \footnotesize\bfseries\color{pendingink}#1\hfill\songfont#2%
  }\\[0.06em]%
}
% \odashrow：虛線分隔（橫跨三欄）
\newcommand{\odashrow}{%
  \noalign{\vspace{0.28em}}%
  \multicolumn{3}{@{}l@{}}{\color{notegray}%
    \xleaders\hbox to 0.82em{\hfil\rule{0.42em}{0.28pt}\hfil}\hskip 0.88\linewidth}%
  \\[0.24em]%
}
% \odissentpart：不同意見書區段標頭。比一般 \opartrow 更像附錄性轉場。
\newcommand{\odissentpart}[2]{%
  \noalign{\vspace{0.42em}}%
  \multicolumn{3}{@{}p{0.88\textwidth}@{}}{%
    \color{notegray}\rule{\linewidth}{0.18pt}%
  }\\[-0.18em]%
  \multicolumn{3}{@{}p{0.88\textwidth}@{}}{%
    \makebox[\linewidth][s]{%
      \footnotesize\bfseries\color{pendingink}#1%
      \hfill{\songfont #2}%
    }%
  }\\[-0.02em]%
}
% \odissentopinion：不同意見書的署名與一行摘要，右欄保留段碼範圍。
\newcommand{\odissentopinion}[5]{%
  \footnotesize\hangindent=0.52cm\hangafter=1%
    {\bfseries\color{caseink}#2}\enspace{\color{notegray}\textperiodcentered}\enspace\color{caseink}#3%
  &\footnotesize\hangindent=0.52cm\hangafter=1%
    {\songfont\bfseries\color{caseink}#4}\enspace{\color{notegray}\textperiodcentered}\enspace\songfont\color{caseink}#5%
  &#1%
  \\[0.12em]%
}
% \omainrow：頂層項目，無段碼（Leitsätze、Urteil、Gründe …）
\newcommand{\omainrow}[4]{%
  \footnotesize\hangindent=1.92cm\hangafter=1%
    \makebox[1.92cm][l]{\bfseries\color{caseink}#1}\color{caseink}#2%
  &\footnotesize\hangindent=1.92cm\hangafter=1%
    \makebox[1.92cm][l]{\bfseries\color{caseink}#3}\songfont\color{caseink}#4%
  &%
  \\[0.15em]%
}
% \omainrowrn：頂層項目，有段碼（不同意見等標籤寬於 \osecrow 框者）
\newcommand{\omainrowrn}[5]{%
  \footnotesize\hangindent=1.92cm\hangafter=1%
    \makebox[1.92cm][l]{\bfseries\color{caseink}#2}\color{caseink}#3%
  &\footnotesize\hangindent=1.92cm\hangafter=1%
    \makebox[1.92cm][l]{\bfseries\color{caseink}#4}\songfont\color{caseink}#5%
  &#1%
  \\[0.15em]%
}
% \osecrow：一級子項，無段碼（A.、C.、D. 等有子項者）
\newcommand{\osecrow}[4]{%
  \footnotesize\hspace{1.10cm}\hangindent=1.98cm\hangafter=1%
    \makebox[0.86cm][l]{\bfseries\color{caseink}#1}\color{caseink}#2%
  &\footnotesize\hspace{1.10cm}\hangindent=1.98cm\hangafter=1%
    \makebox[0.86cm][l]{\bfseries\color{caseink}#3}\songfont\color{caseink}#4%
  &%
  \\[0.11em]%
}
% \osecrowrn：一級子項，有段碼（B.、E. 等完整段落範圍者）
% #1=段碼字串（如 101–106）, #2=DE標籤, #3=DE正文, #4=ZH標籤, #5=ZH正文
\newcommand{\osecrowrn}[5]{%
  \footnotesize\hspace{1.10cm}\hangindent=1.98cm\hangafter=1%
    \makebox[0.86cm][l]{\bfseries\color{caseink}#2}\color{caseink}#3%
  &\footnotesize\hspace{1.10cm}\hangindent=1.98cm\hangafter=1%
    \makebox[0.86cm][l]{\bfseries\color{caseink}#4}\songfont\color{caseink}#5%
  &#1%
  \\[0.11em]%
}
% \osubrow：二級子項，有段碼範圍（I.、II.、A.I. …）
% #1=段碼字串, #2=DE標籤, #3=DE正文, #4=ZH標籤, #5=ZH正文
\newcommand{\osubrow}[5]{%
  \footnotesize\hspace{1.40cm}\hangindent=2.30cm\hangafter=1%
    \makebox[0.90cm][l]{\bfseries\color{caseink}#2}\color{caseink}#3%
  &\footnotesize\hspace{1.40cm}\hangindent=2.30cm\hangafter=1%
    \makebox[0.90cm][l]{\bfseries\color{caseink}#4}\songfont\color{caseink}#5%
  &#1%
  \\[0.09em]%
}
% \ointrow：引入段落，有單一段碼，無標籤
% #1=段碼字串, #2=DE正文, #3=ZH正文
\newcommand{\ointrow}[3]{%
  \footnotesize\hspace{0.52cm}\hangindent=0.52cm\hangafter=1\color{caseink}#2%
  &\footnotesize\hspace{0.52cm}\hangindent=0.52cm\hangafter=1\songfont\color{caseink}#3%
  &#1%
  \\[0.09em]%
}
% \osubsubrow：三級子項（1.、2.、3.）
% #1=段碼字串, #2=DE標籤, #3=DE正文, #4=ZH標籤, #5=ZH正文
\newcommand{\osubsubrow}[5]{%
  \footnotesize\hspace{1.80cm}\hangindent=2.48cm\hangafter=1%
    \makebox[0.68cm][l]{\bfseries\color{caseink}#2}\color{caseink}#3%
  &\footnotesize\hspace{1.80cm}\hangindent=2.48cm\hangafter=1%
    \makebox[0.68cm][l]{\bfseries\color{caseink}#4}\songfont\color{caseink}#5%
  &#1%
  \\[0.08em]%
}
% \osubsubsubrow：四級子項（a）、b））
% #1=段碼字串, #2=DE標籤, #3=DE正文, #4=ZH標籤, #5=ZH正文
\newcommand{\osubsubsubrow}[5]{%
  \footnotesize\hspace{2.10cm}\hangindent=2.68cm\hangafter=1%
    \makebox[0.58cm][l]{\bfseries\color{caseink}#2}\color{caseink}#3%
  &\footnotesize\hspace{2.10cm}\hangindent=2.68cm\hangafter=1%
    \makebox[0.58cm][l]{\bfseries\color{caseink}#4}\songfont\color{caseink}#5%
  &#1%
  \\[0.07em]%
}
% \osubsubsubsubrow：五級子項（aa）、bb））
% 標籤框 0.62cm：足以容納「aa)」在 footnotesize bold 下（約 0.50cm）並留有間距
% #1=段碼字串, #2=DE標籤, #3=DE正文, #4=ZH標籤, #5=ZH正文
\newcommand{\osubsubsubsubrow}[5]{%
  \footnotesize\hspace{2.38cm}\hangindent=3.06cm\hangafter=1%
    \makebox[0.68cm][l]{\bfseries\color{caseink}#2}\color{caseink}#3%
  &\footnotesize\hspace{2.38cm}\hangindent=3.06cm\hangafter=1%
    \makebox[0.68cm][l]{\bfseries\color{caseink}#4}\songfont\color{caseink}#5%
  &#1%
  \\[0.06em]%
}
% \outlinepage：雙語對照大綱頁包裝環境（longtable 自動跨頁）
% 三欄：德文欄 + 中文欄 + 段碼欄（最右，不折行）
% 段碼欄寬度：1.80cm，足以容納最長段碼如「309–348」
% DE/ZH 欄寬：(0.87\textwidth - 2.08cm) / 2
%   驗算：2×欄 + 0.05\textwidth + 0.28cm gap + 1.80cm = 0.87tw - 2.08cm + 0.05tw + 2.08cm = 0.92tw ✓
\newenvironment{outlinepage}{%
  \vspace*{0.40cm}%
  \setlength{\LTleft}{0.04\textwidth}%
  \setlength{\LTright}{0.04\textwidth}%
  \setlength{\tabcolsep}{0pt}%
  % 大綱列距由 longtable 的 \arraystretch 與各 row macro 結尾的 \\[...em] 共同決定。
  % 這裡控制「每一列本身的表格 strut 高度」；數值越大，A./I./1. 各項之間越像有空行。
  % A3 原本用 0.62，視覺上沒有空行；A4 也沿用同值，避免切到 A4 後項目被撐開。
  % 若日後覺得大綱太擠，只微調這個值（例如 0.68），不要改各 \omainrow/\osubrow 的縮排。
  \renewcommand{\arraystretch}{0.62}%
  \begin{longtable}{%
    >{\vspace{0pt}\raggedright\arraybackslash}p{\dimexpr(0.87\textwidth-2.08cm)/2\relax}%
    @{\hspace{0.025\textwidth}}%
    !{\color{notegray}\vrule width 0.12pt}%
    @{\hspace{0.025\textwidth}}%
    >{\vspace{0pt}\raggedright\arraybackslash}p{\dimexpr(0.87\textwidth-2.08cm)/2\relax}%
    @{\hspace{0.28cm}}%
    >{\vspace{0pt}\footnotesize\normalfont\color{pendingink}\raggedright\arraybackslash}p{1.80cm}%
    @{}%
  }%
  % 首頁欄標籤：不要橫跨置中；分別放在德文欄與中文欄的實際欄位起點。
  \footnotesize\bfseries\color{sourceink}Gliederung%
  &\footnotesize\bfseries\songfont\color{sourceink}大綱%
  &%
  \\[0.36em]%
  \endfirsthead%
  % 續頁欄標籤：同樣分欄定位，以灰色細體區分；無需標注「續」。
  \footnotesize\color{pendingink}Gliederung%
  &\footnotesize\songfont\color{pendingink}大綱%
  &%
  \\[0.24em]%
  \endhead%
}{%
  \end{longtable}%
}
 載入。
% 不含：\documentclass、\documentversion、\ggversionde/zh、\tocde/\toczh。
% 日期與首頁版本列由本檔自動產生；各文件只需手動指定 \documentversion。

\usepackage{xeCJK}
\usepackage{fontspec}
\usepackage{geometry}
\usepackage{setspace}
\usepackage{hyperref}
\usepackage{xurl}
\usepackage{bookmark}
\usepackage{enumitem}
\usepackage{xcolor}
\usepackage{titlesec}
\usepackage{array}
\usepackage{longtable}
\usepackage{tcolorbox}
\usepackage{environ}
\usepackage{pdflscape}
\usepackage{etoolbox}
\usepackage{ragged2e}
\usepackage{paracol}
\usepackage{multicol}
\usepackage{needspace}
\usepackage{fancyhdr}
\usepackage{tikz}
\usepackage{zref-abspage}
\tcbuselibrary{breakable,skins}
\usetikzlibrary{calc}
\usepackage{pgfcalendar}

\hypersetup{
  unicode=true,
  bookmarksopen=true,
  bookmarksnumbered=false,
  hidelinks
}

% ===== 自動推導，不要手動改下面 =====
% （\docmode 與 \bilingualpaper 由各文件在 % bverfge_layout.tex
% 共用排版基礎設施，供 Schwangerschaftsabbruch I 與 II 共享。
% 由各文件以 \documentclass 後 % bverfge_layout.tex
% 共用排版基礎設施，供 Schwangerschaftsabbruch I 與 II 共享。
% 由各文件以 \documentclass 後 % bverfge_layout.tex
% 共用排版基礎設施，供 Schwangerschaftsabbruch I 與 II 共享。
% 由各文件以 \documentclass 後 \input{../../共用LaTeX/bverfge_layout} 載入。
% 不含：\documentclass、\documentversion、\ggversionde/zh、\tocde/\toczh。
% 日期與首頁版本列由本檔自動產生；各文件只需手動指定 \documentversion。

\usepackage{xeCJK}
\usepackage{fontspec}
\usepackage{geometry}
\usepackage{setspace}
\usepackage{hyperref}
\usepackage{xurl}
\usepackage{bookmark}
\usepackage{enumitem}
\usepackage{xcolor}
\usepackage{titlesec}
\usepackage{array}
\usepackage{longtable}
\usepackage{tcolorbox}
\usepackage{environ}
\usepackage{pdflscape}
\usepackage{etoolbox}
\usepackage{ragged2e}
\usepackage{paracol}
\usepackage{multicol}
\usepackage{needspace}
\usepackage{fancyhdr}
\usepackage{tikz}
\usepackage{zref-abspage}
\tcbuselibrary{breakable,skins}
\usetikzlibrary{calc}
\usepackage{pgfcalendar}

\hypersetup{
  unicode=true,
  bookmarksopen=true,
  bookmarksnumbered=false,
  hidelinks
}

% ===== 自動推導，不要手動改下面 =====
% （\docmode 與 \bilingualpaper 由各文件在 \input{bverfge_layout} 之前設定；
%   預設 chinese / a4）
\ifdefined\docmode\else\def\docmode{chinese}\fi
\ifdefined\bilingualpaper\else\def\bilingualpaper{a4}\fi
\newif\ifshoworiginal
\newif\ifshowtranslation
\newif\ifbilinguallandscape
\newif\ifbilingualcolumns
\newif\ifintranslation
\newif\iffirstbilingualsection

\showoriginalfalse
\showtranslationfalse
\bilinguallandscapefalse
\bilingualcolumnsfalse
\intranslationfalse
\firstbilingualsectionfalse

\def\modechinese{chinese}
\def\modegerman{german}
\def\modebilingual{bilingual}
\def\bilingualpaperaThree{a3}
\def\bilingualpaperaFour{a4}

\ifx\docmode\modechinese
  \showtranslationtrue
\fi

\ifx\docmode\modegerman
  \showoriginaltrue
\fi

\ifx\docmode\modebilingual
  \showoriginaltrue
  \showtranslationtrue
  \bilinguallandscapetrue
  \bilingualcolumnstrue
\fi

% 編排模式說明：
% chinese  → \showoriginalfalse  \showtranslationtrue  （A4 直式，純中文）
% german   → \showoriginaltrue   \showtranslationfalse （A4 直式，純德文）
% bilingual→ \showoriginaltrue   \showtranslationtrue  \bilingualcolumnstrue（A3/A4 橫式對照）

\ifbilingualcolumns
  \ifx\bilingualpaper\bilingualpaperaFour
    \geometry{
      a4paper,
      landscape,
      left=1.25cm,
      right=2.25cm,
      top=1.25cm,
      bottom=1.75cm,
      footskip=8.5mm,
      marginparwidth=0.75cm,
      marginparsep=0.25cm
    }
  \else
    \geometry{
      a3paper,
      landscape,
      left=1.55cm,
      right=2.75cm,
      top=1.55cm,
      bottom=2.05cm,
      footskip=9.5mm,
      marginparwidth=0.9cm,
      marginparsep=0.35cm
    }
  \fi
\else
\geometry{
  a4paper,
  portrait,
  left=2.2cm,
  right=2.6cm,
  top=2.0cm,
  bottom=2.0cm,
  marginparwidth=0.8cm,
  marginparsep=0.3cm
}
\fi
% ===== 時間戳基礎設施 =====
\newcount\stamptimeval
\newcount\stampminuteval
\newcommand{\stamphh}{%
  \stamptimeval=\time
  \divide\stamptimeval by 60
  \ifnum\stamptimeval<10 0\fi
  \the\stamptimeval}
\newcommand{\stampmm}{%
  \stamptimeval=\time
  \stampminuteval=\time
  \divide\stampminuteval by 60
  \multiply\stampminuteval by 60
  \advance\stamptimeval by -\stampminuteval
  \ifnum\stamptimeval<10 0\fi
  \the\stamptimeval}
\newcommand{\stampwkde}[1]{%
  \ifcase#1Montag\or Dienstag\or Mittwoch\or Donnerstag\or Freitag\or Samstag\or Sonntag\fi}
\newcommand{\stampwkzh}[1]{%
  \ifcase#1 一\or 二\or 三\or 四\or 五\or 六\or 日\fi}
\newcommand{\stampmonthde}[1]{%
  \ifcase#1\or Januar\or Februar\or März\or April\or Mai\or Juni\or
  Juli\or August\or September\or Oktober\or November\or Dezember\fi}
\newcount\stampjdval
\newcount\stampwdval
\pgfcalendardatetojulian{\the\year-\the\month-\the\day}{\stampjdval}
\pgfcalendarjuliantoweekday{\the\stampjdval}{\stampwdval}
\newcommand{\timestampzh}{%
  \songfont 最後修改：\the\year 年\the\month 月\the\day 日%
  % （星期\stampwkzh{\the\stampwdval}）%
  % \space\stamphh:\stampmm
  }

\newcommand{\documentdatezh}{\the\year 年 \the\month 月 \the\day 日}
\newcommand{\documentdatede}{\the\day. \stampmonthde{\the\month} \the\year}
\newcommand{\bverfgetranslationnote}{%
  {\songfont 翻譯說明：初譯使用 Claude Code 與 GPT 協助迭代；仍請以德文原文為準。}%
}
\newcommand{\bverfgecourseinfo}{%
  {\songfont 課程：114 學年度第 2 學期；憲法專題研究（四）；授課老師：湯德宗教授；導讀日期：115 年 6 月 1 日。}%
}
\newcommand{\bverfgeeditorinfo}{%
  {\songfont 編者：王逸帆（R10A21126），國立台灣大學法律系研究所碩士班。}%
}
\newcommand{\bverfgetitleversionline}{%
  \ifbilingualcolumns
    Fassung \documentversion\enspace\textperiodcentered\enspace Stand: \documentdatede
    \quad
    {\songfont 版本 \documentversion\enspace\textperiodcentered\enspace 修訂日期：\documentdatezh\par}%
  \else
    \ifshowtranslation
      {\songfont 版本 \documentversion\enspace\textperiodcentered\enspace 修訂日期：\documentdatezh\par}%
    \else
      Fassung \documentversion\enspace\textperiodcentered\enspace Stand: \documentdatede\par
    \fi
  \fi
}
\newcommand{\bverfgetitlecornerinfo}{%
  \ifshowtranslation
    \vspace{0.72em}%
    \noindent\hfill
    \begin{minipage}{0.66\textwidth}
      \raggedleft\tiny\color{pendingink}\songfont
      \bverfgecourseinfo\par
      \bverfgeeditorinfo
    \end{minipage}%
  \fi
}
\newcommand{\bverfgetitlemeta}{%
  \vfill
  \begin{center}%
    \footnotesize\color{pendingink}%
    \bverfgetitleversionline
    \ifshowtranslation
      \vspace{0.28em}{\scriptsize\bverfgetranslationnote\par}%
    \fi
  \end{center}%
  \bverfgetitlecornerinfo
}

\pagestyle{fancy}
\fancyhf{}
\renewcommand{\headrulewidth}{0pt}
\renewcommand{\footrulewidth}{0pt}
\fancyfoot[C]{\thepage}
\ifbilingualcolumns
  \fancyfoot[R]{\scriptsize\color{pendingink}\timestampzh}%
\else
  \ifshoworiginal
  \else
    \fancyfoot[R]{\scriptsize\color{pendingink}\timestampzh\hspace{0.45cm}}%
  \fi
\fi
\setmainfont{Times New Roman}
\setsansfont{Helvetica Neue}
\setmonofont{Menlo}
\IfFileExists{/Users/iw/Library/Fonts/kaiu.ttf}{
  \setCJKmainfont[Path=/Users/iw/Library/Fonts/,AutoFakeBold=4]{kaiu.ttf}
}{
  \IfFontExistsTF{標楷體}{
    \setCJKmainfont[AutoFakeBold=4]{標楷體}
  }{
    \setCJKmainfont{Songti TC}
  }
}
\IfFontExistsTF{Songti TC}{
  \setCJKfamilyfont{song}{Songti TC}
  \setCJKmonofont{Songti TC}
}{
  \setCJKfamilyfont{song}[Path=/Users/iw/Library/Fonts/,AutoFakeBold=4]{kaiu.ttf}
  \setCJKmonofont[Path=/Users/iw/Library/Fonts/,AutoFakeBold=4]{kaiu.ttf}
}
\newcommand{\songfont}{\CJKfamily{song}}

% ===== 封面／目錄專用打字機字體（僅限封面、目錄頁使用，不影響正文）=====
\IfFontExistsTF{Radio Newsman}{%
  \newfontfamily\bverfgetitlede{Radio Newsman}%
}{%
  \newfontfamily\bverfgetitlede{Courier New}%
}
\IfFontExistsTF{Erikas Farbband}{%
  \newfontfamily\bverfgebodyde{Erikas Farbband}%
}{%
  \let\bverfgebodyde\bverfgetitlede
}
\IfFontExistsTF{Maoken Heavy Labourer}{%
  \setCJKfamilyfont{bverfgetitlecjk}{Maoken Heavy Labourer}%
}{%
  \setCJKfamilyfont{bverfgetitlecjk}{Songti TC}%
}
\IfFontExistsTF{Huiwen-mincho}{%
  \setCJKfamilyfont{bverfgebodycjk}{Huiwen-mincho}%
}{%
  \setCJKfamilyfont{bverfgebodycjk}{Songti TC}%
}
\newcommand{\bverfgetitlezh}{\bverfgetitlede\CJKfamily{bverfgetitlecjk}}
\newcommand{\bverfgebodyzh}{\bverfgebodyde\CJKfamily{bverfgebodycjk}}

\setstretch{1.35}
\setlist{itemsep=0.2em, topsep=0.4em}
\setlength{\parindent}{0pt}
\setlength{\parskip}{0.45em}
\raggedbottom
\setcounter{secnumdepth}{0}
\setcounter{tocdepth}{2}

\newcommand{\lawboxsourcefont}{\footnotesize}
\newcommand{\lawboxtransfont}{\footnotesize}
\newcommand{\lawboxsourcestretch}{1.02}
\newcommand{\lawboxtransstretch}{0.92}

\titleformat{\section}{\centering\Large\bfseries\songfont}{\thesection}{1em}{}
\titleformat{\subsection}{\large\bfseries\songfont}{\thesubsection}{1em}{}
\titleformat{\subsubsection}{\normalsize\bfseries\songfont}{\thesubsubsection}{1em}{}
\titleformat{\paragraph}{\normalsize\bfseries\songfont}{\theparagraph}{1em}{}
\titlespacing*{\section}{0pt}{1.8em}{1.0em}
\titlespacing*{\subsection}{0pt}{1.2em}{0.6em}
\titlespacing*{\subsubsection}{0pt}{1.0em}{0.45em}
\titlespacing*{\paragraph}{0pt}{0.6em}{0.4em}

\definecolor{noteblue}{HTML}{A9C9F5}
\definecolor{noteteal}{HTML}{AEE1D6}
\definecolor{notegray}{HTML}{D7DADF}
\definecolor{caseink}{HTML}{000000}
\definecolor{casepaper}{HTML}{FCFEFD}
\definecolor{sourcepaper}{HTML}{FAFBF8}
\definecolor{sourceline}{HTML}{D7DED8}
\definecolor{sourceink}{HTML}{000000}
\definecolor{pendingink}{HTML}{000000}

\tcbset{
  caseboxbase/.style={
    enhanced,
    breakable,
    width=0.985\linewidth,
    colback=casepaper,
    coltitle=caseink,
    fonttitle=\bfseries\songfont,
    boxrule=0.28pt,
    arc=1.2mm,
    left=2.4mm,
    right=2.4mm,
    top=1.5mm,
    bottom=1.5mm,
    boxsep=1.5mm,
    before skip=0.65em,
    after skip=0.65em,
    before upper={\setlength{\parindent}{0pt}\setlength{\parskip}{0.35em}},
  }
}

\newcommand{\bilingualstart}{}
\newcommand{\bilingualstop}{}
\ifbilingualcolumns
  \renewcommand{\bilingualstart}{\small\setlength{\columnsep}{1.25cm}\setlength{\columnseprule}{0.12pt}\columnratio{0.5,0.5}\multicolumnfootnotes\flushbottom\sloppy\interlinepenalty=80\clubpenalty=800\widowpenalty=800\displaywidowpenalty=800\begin{paracol}{2}\global\intranslationfalse\global\firstbilingualsectiontrue{\footnotesize\color{sourceink}Deutsch\par}\switchcolumn{\footnotesize\songfont\color{sourceink}中文譯文\par}\switchcolumn*}
  \renewcommand{\bilingualstop}{\ifintranslation\switchcolumn*\global\intranslationfalse\fi\end{paracol}\clearpage\raggedbottom}
\fi
\newif\ifrandnumbers
\randnumbersfalse
\newcounter{randnumber}
\newcounter{randnumberlocal}
\newcounter{lastsourcerandstart}
\newif\iflastsourcehasrandnumbers
\lastsourcehasrandnumbersfalse
\newif\ifinlawbox
\inlawboxfalse
\newif\iflawboxhaspendingrn
\lawboxhaspendingrnfalse
\newcommand{\lawboxpendingrn}{}
\newcommand{\startrandnumbers}{\global\randnumberstrue\setcounter{randnumber}{1}\global\lastsourcehasrandnumbersfalse}
\newcommand{\stoprandnumbers}{\global\randnumbersfalse\global\lastsourcehasrandnumbersfalse}
\newlength{\randnumberwidth}
\setlength{\randnumberwidth}{0.95cm}
\newlength{\randnumberrightoffset}
\setlength{\randnumberrightoffset}{0.85cm}
\newlength{\randnumberlawboxrightoffset}
\setlength{\randnumberlawboxrightoffset}{1.40cm}
\newcommand{\randnumberbox}[1]{%
  \leavevmode
  \makebox[0pt][l]{%
    \smash{%
      \tikz[remember picture,overlay,baseline=(rnmark.base)]{%
        \coordinate (rnmark) at (0,0);%
        \ifinlawbox
          \path let \p1=(current page.east), \p2=(rnmark.base) in
            node[
              anchor=base east,
              inner sep=0pt,
              outer sep=0pt,
              text width=\randnumberwidth,
              align=right,
              text=pendingink,
              font=\normalfont\mdseries\scriptsize\ttfamily
            ] at (\x1-\randnumberlawboxrightoffset,\y2) {{\normalfont\mdseries\scriptsize\ttfamily #1}};%
        \else
          \path let \p1=(current page.east), \p2=(rnmark.base) in
            node[
              anchor=base east,
              inner sep=0pt,
              outer sep=0pt,
              text width=\randnumberwidth,
              align=right,
              text=pendingink,
              font=\normalfont\mdseries\scriptsize\ttfamily
            ] at (\x1-\randnumberrightoffset,\y2) {{\normalfont\mdseries\scriptsize\ttfamily #1}};%
        \fi
      }%
    }%
  }%
  \ignorespaces
}
\newcommand{\lawboxstorerandnumber}[1]{%
  \gdef\lawboxpendingrn{#1}%
  \global\lawboxhaspendingrntrue
  \ignorespaces
}
\newcommand{\lawboxuserandnumber}{%
  \iflawboxhaspendingrn
    \randnumberbox{\lawboxpendingrn}%
    \global\lawboxhaspendingrnfalse
  \fi
}
\newcommand{\sourcern}[1]{%
  \ifrandnumbers
    \ifshoworiginal
      \ifshowtranslation\else
        \ifinlawbox\lawboxstorerandnumber{#1}\else\randnumberbox{#1}\fi
      \fi
    \fi
  \fi
}
\newcommand{\transrn}[1]{%
  \ifrandnumbers
    \ifshowtranslation
      \ifinlawbox\lawboxstorerandnumber{#1}\else\randnumberbox{#1}\fi
    \fi
  \fi
}
% ===== 課堂模式：德文原文縮寫註解 =====
% 日常切換只改各判決檔的一行：
%   \classroommodeon        開啟課堂版；註解掉這行即回到一般版。
% 模板預設：
%   1. 課堂模式關閉。
%   2. 課堂模式開啟時，縮寫註解包含「德文全稱＋中文譯名」。
%   3. 課堂模式開啟時，GG/條文編者註仍會自動顯示。
% 進階微調才需要在單案檔覆寫：
%   \classroomabbrzhoff     縮寫註解僅顯示德文全稱。
%   \showeditornotestrue    一般版也顯示編者註。
% 註解策略：同一實際 PDF 頁面中，同一縮寫只在第一次出現處加註。
% 這裡用 zref-abspage 的絕對頁序，不用 \thepage，避免文件內頁碼重置時誤判。
\newif\ifclassroommode
\newif\ifclassroomabbrzh
\classroommodefalse
\classroomabbrzhtrue
\newcommand{\classroommodeon}{\classroommodetrue}
\newcommand{\classroommodeoff}{\classroommodefalse}
\newcommand{\classroomabbrzhon}{\classroomabbrzhtrue}
\newcommand{\classroomabbrzhoff}{\classroomabbrzhfalse}
\newcommand{\abbrnotelabel}{\emph{Abk.}: }
\newcommand{\abbrnote}[3]{%
  \ifclassroommode
    \ifshoworiginal
      \ifintranslation\else
        \footnote{\normalfont\footnotesize\abbrnotelabel #2\ifclassroomabbrzh；{\songfont #3}\fi}%
      \fi
    \fi
  \fi
}
\newcommand{\abbrnoteonce}[4]{%
  #2%
  \ifclassroommode
    \ifshoworiginal
      \ifintranslation\else
        \begingroup
          \edef\abbrcurrentabspage{\number\numexpr\value{abspage}+1\relax}%
          \ifcsname abbrnoteseen@#1@\abbrcurrentabspage\endcsname\else
            \global\expandafter\let\csname abbrnoteseen@#1@\abbrcurrentabspage\endcsname\empty
            \abbrnote{#1}{#3}{#4}%
          \fi
        \endgroup
      \fi
    \fi
  \fi
}
\newcommand{\abbrAbs}{\abbrnoteonce{abs}{Abs.}{Absatz}{項}}
\newcommand{\abbrAAO}{\abbrnoteonce{aao}{a.a.O.}{am angegebenen Ort}{前引處}}
\newcommand{\abbrAE}{\abbrnoteonce{ae}{AE}{Alternativ-Entwurf}{替代草案}}
\newcommand{\abbrArt}{\abbrnoteonce{art}{Art.}{Artikel}{條}}
\newcommand{\abbrAufl}{\abbrnoteonce{aufl}{Aufl.}{Auflage}{版}}
\newcommand{\abbrBd}{\abbrnoteonce{bd}{Bd.}{Band}{卷}}
\newcommand{\abbrBGBl}{\abbrnoteonce{bgbl}{BGBl.}{Bundesgesetzblatt}{聯邦法律公報}}
\newcommand{\abbrBGHSt}{\abbrnoteonce{bghst}{BGHSt}{Entscheidungen des Bundesgerichtshofes in Strafsachen}{聯邦普通法院刑事判決彙編}}
\newcommand{\abbrBRDrucks}{\abbrnoteonce{brdrucks}{BRDrucks.}{Bundesrats-Drucksache}{聯邦參議院文件}}
\newcommand{\abbrBTDrucks}{\abbrnoteonce{btdrucks}{BTDrucks.}{Bundestags-Drucksache}{聯邦議院文件}}
\newcommand{\abbrBundesgesetzbl}{\abbrnoteonce{bundesgesetzbl}{Bundesgesetzbl.}{Bundesgesetzblatt}{聯邦法律公報}}
\newcommand{\abbrBvF}{\abbrnoteonce{bvf}{BvF}{Bundesverfassungsgerichtsverfahren}{聯邦憲法法院程序案號}}
\newcommand{\abbrBvL}{\abbrnoteonce{bvl}{BvL}{Bundesverfassungsgerichtsverfahren}{聯邦憲法法院程序案號}}
\newcommand{\abbrBvR}{\abbrnoteonce{bvr}{BvR}{Bundesverfassungsgerichtsverfahren}{聯邦憲法法院程序案號}}
\newcommand{\abbrBVerfG}{\abbrnoteonce{bverfg}{BVerfG}{Bundesverfassungsgericht}{聯邦憲法法院}}
\newcommand{\abbrBVerfGE}{\abbrnoteonce{bverfge}{BVerfGE}{Entscheidungen des Bundesverfassungsgerichts}{聯邦憲法法院判決集}}
\newcommand{\abbrBVerfGG}{\abbrnoteonce{bverfgg}{BVerfGG}{Bundesverfassungsgerichtsgesetz}{聯邦憲法法院法}}
\newcommand{\abbrDDR}{\abbrnoteonce{ddr}{DDR}{Deutsche Demokratische Republik}{德意志民主共和國；東德}}
\newcommand{\abbrDJT}{\abbrnoteonce{djt}{DJT}{Deutscher Juristentag}{德國法學家大會}}
\newcommand{\abbrDr}{\abbrnoteonce{dr}{Dr.}{Doktor}{Dr.}}
\newcommand{\abbrEuGRZ}{\abbrnoteonce{eugrz}{EuGRZ}{Europäische Grundrechte-Zeitschrift}{歐洲基本權期刊}}
\newcommand{\abbrEV}{\abbrnoteonce{ev}{EV}{Einigungsvertrag}{統一條約}}
\newcommand{\abbrF}{\abbrnoteonce{f}{f.}{folgende}{以下一頁；下一段}}
\newcommand{\abbrFF}{\abbrnoteonce{ff}{ff.}{fortfolgende}{以下數頁；以下各段}}
\newcommand{\abbrGBlDDR}{\abbrnoteonce{gblddr}{GBl. DDR}{Gesetzblatt der Deutschen Demokratischen Republik}{德意志民主共和國法律公報}}
\newcommand{\abbrGG}{\abbrnoteonce{gg}{GG}{Grundgesetz}{基本法}}
\newcommand{\abbrGVBl}{\abbrnoteonce{gvbl}{GVBl.}{Gesetz- und Verordnungsblatt}{法律及命令公報}}
\newcommand{\abbrJR}{\abbrnoteonce{jr}{JR}{Juristische Rundschau}{法學評論}}
\newcommand{\abbrJZ}{\abbrnoteonce{jz}{JZ}{JuristenZeitung}{法學家報}}
\newcommand{\abbrIVm}{\abbrnoteonce{ivm}{i.V.m.}{in Verbindung mit}{結合；連同}}
\newcommand{\abbrMdB}{\abbrnoteonce{mdb}{MdB}{Mitglied des Bundestages}{聯邦議院議員}}
\newcommand{\abbrMwN}{\abbrnoteonce{mwn}{m.w.N.}{mit weiteren Nachweisen}{附進一步文獻／判例出處}}
\newcommand{\abbrNF}{\abbrnoteonce{nf}{n.F.}{neue Fassung}{新版本}}
\newcommand{\abbrAF}{\abbrnoteonce{af}{a.F.}{alte Fassung}{舊版本}}
\newcommand{\abbrNr}{\abbrnoteonce{nr}{Nr.}{Nummer}{款、目或號，依語境}}
\newcommand{\abbrProf}{\abbrnoteonce{prof}{Prof.}{Professor}{教授}}
\newcommand{\abbrRdnr}{\abbrnoteonce{rdnr}{Rdnr.}{Randnummer}{邊碼；段碼}}
\newcommand{\abbrRdnrn}{\abbrnoteonce{rdnrn}{Rdnrn.}{Randnummern}{邊碼；段碼}}
\newcommand{\abbrRGBl}{\abbrnoteonce{rgbl}{RGBl.}{Reichsgesetzblatt}{帝國法律公報}}
\newcommand{\abbrRGSt}{\abbrnoteonce{rgst}{RGSt}{Entscheidungen des Reichsgerichts in Strafsachen}{帝國法院刑事判決彙編}}
\newcommand{\abbrRspr}{\abbrnoteonce{rspr}{Rspr.}{Rechtsprechung}{判例；司法實務}}
\newcommand{\abbrRVO}{\abbrnoteonce{rvo}{RVO}{Reichsversicherungsordnung}{帝國保險法}}
\newcommand{\abbrS}{\abbrnoteonce{s}{S.}{Seite}{頁}}
\newcommand{\abbrSFHG}{\abbrnoteonce{sfhg}{SFHG}{Schwangeren- und Familienhilfegesetz}{孕婦及家庭扶助法}}
\newcommand{\abbrSGBV}{\abbrnoteonce{sgbv}{SGB V}{Fünftes Buch Sozialgesetzbuch}{社會法典第五編}}
\newcommand{\abbrStAG}{\abbrnoteonce{staeg}{StÄG}{Strafrechtsänderungsgesetz}{刑法修正法}}
\newcommand{\abbrStenBer}{\abbrnoteonce{stenber}{StenBer.}{Stenographischer Bericht}{速記紀錄}}
\newcommand{\abbrStGB}{\abbrnoteonce{stgb}{StGB}{Strafgesetzbuch}{刑法}}
\newcommand{\abbrStREG}{\abbrnoteonce{streg}{StREG}{Strafrechtsreform-Ergänzungsgesetz}{刑法改革補充法}}
\newcommand{\abbrStrRG}{\abbrnoteonce{strrg}{StrRG}{Strafrechtsreformgesetz}{刑法改革法}}
\newcommand{\abbrUS}{\abbrnoteonce{us}{U.S.}{United States Reports}{美國最高法院判決彙編}}
\newcommand{\abbrVgl}{\abbrnoteonce{vgl}{vgl.}{vergleiche}{參見}}
\newcommand{\abbrVVDStRL}{\abbrnoteonce{vvdstrl}{VVDStRL}{Veröffentlichungen der Vereinigung der Deutschen Staatsrechtslehrer}{德國公法教師協會論文集}}
\newcommand{\abbrWP}{\abbrnoteonce{wp}{WP.}{Wahlperiode}{屆期}}
\newcommand{\abbrWp}{\abbrnoteonce{wp}{Wp.}{Wahlperiode}{屆期}}
\newcommand{\abbrZStrW}{\abbrnoteonce{zstrw}{ZStrW}{Zeitschrift für die gesamte Strafrechtswissenschaft}{整體刑法學期刊}}
\newcommand{\recordsourcerandstart}{%
  \ifrandnumbers
    \setcounter{lastsourcerandstart}{\value{randnumber}}%
    \global\lastsourcehasrandnumberstrue
  \else
    \global\lastsourcehasrandnumbersfalse
  \fi
}
\newlength{\biparagraphskip}
\setlength{\biparagraphskip}{0.52em}
\newlength{\lawboxblockbeforeskip}
\setlength{\lawboxblockbeforeskip}{0.72em}
\newlength{\lawboxblockafterskip}
\setlength{\lawboxblockafterskip}{0.92em}
\newif\iflawboxpartendedblock
\lawboxpartendedblockfalse
\newcommand{\sourceparstyle}{\footnotesize\color{sourceink}\setstretch{1.12}\RaggedRight\emergencystretch=3em\setlength{\parskip}{0.42em}\obeylines\selectfont}
\newcommand{\transparstyle}{\songfont\color{caseink}\setstretch{1.12}\setlength{\parindent}{0pt}\setlength{\parskip}{0.42em}}
\newcommand{\bilingualpairskip}{%
  \par\addvspace{\biparagraphskip}%
  \switchcolumn
  \par\addvspace{\biparagraphskip}%
  \switchcolumn*%
  \global\intranslationfalse
}
\newcommand{\bilingualboxpairskip}{%
  \par
  \switchcolumn
  \par
  \switchcolumn*%
  \global\intranslationfalse
}
\newcommand{\bilinguallawboxblockskip}{%
  \switchcolumn*%
  \par\addvspace{\lawboxblockafterskip}%
  \switchcolumn
  \par\addvspace{\lawboxblockafterskip}%
  \switchcolumn*%
  \global\intranslationfalse
}
\newcommand{\bikeepwithnext}[1][4]{%
  \ifbilingualcolumns
    \syncleft
    \Needspace{#1\baselineskip}%
    \switchcolumn
    \Needspace{#1\baselineskip}%
    \switchcolumn*%
    \global\intranslationfalse
  \else
    \Needspace{#1\baselineskip}%
  \fi
}
\NewEnviron{sourcepara}[1][]{
  \recordsourcerandstart
  \ifshoworiginal
    \ifbilingualcolumns
      \ifintranslation\switchcolumn*\global\intranslationfalse\fi
    \else
      \Needspace{5\baselineskip}
    \fi
    \ifbilingualcolumns\else\par\addvspace{0.35em}\fi
    {\sourceparstyle
    \noindent\BODY\par}
    \ifbilingualcolumns\else\addvspace{0.25em}\fi
    \ifbilingualcolumns
      \switchcolumn
      \global\intranslationtrue
    \fi
  \else
    \ifrandnumbers
      \begingroup
        \setbox0=\vbox{\hsize=\linewidth\sourceparstyle\noindent\BODY\par}%
      \endgroup
    \fi
  \fi
}
\NewEnviron{transpara}{
  \ifshowtranslation
    \setcounter{randnumberlocal}{\value{lastsourcerandstart}}%
    \ifbilingualcolumns
      \ifintranslation\else\switchcolumn\global\intranslationtrue\fi
      {\transparstyle\setlength{\leftskip}{0.55em}\setlength{\rightskip}{0.15em plus 1em}\addtolength{\linewidth}{-0.55em}\BODY\par}
      \switchcolumn*
      \bilingualpairskip
    \else
      {\transparstyle\BODY\par}
    \fi
  \else
    \ifbilingualcolumns
      \ifintranslation\switchcolumn*\global\intranslationfalse\fi
    \fi
  \fi
}
\NewEnviron{sourceboxpara}[1][]{
  \global\lawboxpartendedblockfalse
  \recordsourcerandstart
  \ifshoworiginal
    \ifbilingualcolumns
      \ifintranslation\switchcolumn*\global\intranslationfalse\fi
    \else
      \Needspace{3\baselineskip}
    \fi
    {\sourceparstyle
    \BODY\par}
    \ifbilingualcolumns\else
      \iflawboxpartendedblock\addvspace{\lawboxblockafterskip}\fi
    \fi
    \ifbilingualcolumns
      \switchcolumn
      \global\intranslationtrue
    \fi
  \fi
}
\NewEnviron{transboxpara}{
  \global\lawboxpartendedblockfalse
  \ifshowtranslation
    \setcounter{randnumberlocal}{\value{lastsourcerandstart}}%
    \ifbilingualcolumns
      \ifintranslation\else\switchcolumn\global\intranslationtrue\fi
      {\transparstyle\BODY\par}
      \iflawboxpartendedblock
        \bilinguallawboxblockskip
      \else
        \switchcolumn*
        \bilingualboxpairskip
      \fi
    \else
      {\transparstyle\BODY\par}
      \iflawboxpartendedblock\addvspace{\lawboxblockafterskip}\fi
    \fi
  \else
    \ifbilingualcolumns
      \ifintranslation\switchcolumn*\global\intranslationfalse\fi
    \fi
  \fi
}
\newcommand{\leitsatznum}[1]{\noindent\hangindent=3.0em\hangafter=1\makebox[2.25em][r]{\bfseries #1.}\hspace{0.55em}\ignorespaces}
\newcommand{\syncleft}{\ifbilingualcolumns\ifintranslation\switchcolumn*\global\intranslationfalse\fi\fi}
\pretocmd{\section}{\syncleft\clearpage}{}{}
\pretocmd{\subsection}{\syncleft}{}{}
\pretocmd{\subsubsection}{\syncleft}{}{}
\newcommand{\biheadingfont}{\songfont}
\newcommand{\bitocentry}[2]{%
  \ifshoworiginal
    \ifshowtranslation
      \ifstrequal{#1}{#2}{#1}{#1\space #2}%
    \else
      #1%
    \fi
  \else
    #2%
  \fi
}
\pdfstringdefDisableCommands{%
  \def\bitocentry#1#2{#1}%
}
\newcommand{\biheading}[7]{%
  \syncleft#1\bikeepwithnext[#3]\phantomsection\addcontentsline{toc}{#2}{\bitocentry{#6}{#7}}%
  \ifbilingualcolumns
    {#4 #6\par}%
    \switchcolumn
    {#4\biheadingfont #7\par}%
    \switchcolumn*%
    \global\intranslationfalse
    \par\addvspace{#5}%
    \switchcolumn
    \par\addvspace{#5}%
    \switchcolumn*%
  \else
    \ifshoworiginal{#4 #6\par}\fi
    \ifshowtranslation{#4\biheadingfont #7\par}\fi
    \addvspace{#5}%
  \fi
}
\newcommand{\termsectionheading}[1]{%
  \par\Needspace{9\baselineskip}%
  \vspace{0.65em}%
  {\songfont\bfseries #1\par}%
  \vspace{0.2em}%
}
% 章節換頁：
%   雙語 paracol 欄內不要用 \clearpage；它會在同步兩欄時吐出只有欄線/頁碼的空頁。
%   \newpage 足以讓下一個 section 起新頁，又不會強制清空所有待排材料。
%   單語模式不在 paracol 內，仍可用 \clearpage。
\newcommand{\bisectionpagebreak}{\ifbilingualcolumns\newpage\else\clearpage\fi}
\newcommand{\bisection}[2]{%
  \biheading{\iffirstbilingualsection\global\firstbilingualsectionfalse\else\bisectionpagebreak\fi}{section}{6}{\centering\Large\bfseries}{0.8em}{#1}{#2}%
}
\newcommand{\bisubsection}[2]{%
  \biheading{}{subsection}{5}{\large\bfseries}{0.55em}{#1}{#2}%
}
\newcommand{\bisubsubsection}[2]{%
  \biheading{}{subsubsection}{4}{\normalsize\bfseries}{0.45em}{#1}{#2}%
}
\newcommand{\bicompositeheading}[4]{%
  \biheading{}{subsection}{5}{\large\bfseries}{0.16em}{#1}{#2}%
  \biheading{}{subsubsection}{3}{\normalsize\bfseries}{0.45em}{#3}{#4}%
}
\newif\iflawboxfirstitem
\lawboxfirstitemfalse
\newcommand{\lawboxprespace}[1]{%
  \par
  \iflawboxfirstitem
    \global\lawboxfirstitemfalse
  \else
    \addvspace{#1}%
  \fi
}
\newtcolorbox{lawboxinner}{
  caseboxbase,
  colframe=sourceline!85!black,
  colback=white,
  left=1.05em,
  right=1.05em,
  top=0.62em,
  bottom=0.62em,
  before upper={\global\inlawboxtrue\global\lawboxfirstitemtrue\global\lawboxhaspendingrnfalse\ifintranslation\lawboxtransfont\else\lawboxsourcefont\fi\RaggedRight\setlength{\parindent}{0pt}\setlength{\parskip}{0pt}\ifintranslation\setstretch{\lawboxtransstretch}\else\setstretch{\lawboxsourcestretch}\fi},
  after upper={\global\lawboxhaspendingrnfalse\global\inlawboxfalse}
}
\tcbset{
  lawboxpartbase/.style={
    enhanced,
    breakable=false,
    width=0.985\linewidth,
    colback=white,
    frame hidden,
    boxrule=0pt,
    arc=0pt,
    left=1.05em,
    right=1.05em,
    boxsep=1.5mm,
    before skip=0pt,
    after skip=0pt,
    before upper={\global\inlawboxtrue\global\lawboxfirstitemtrue\global\lawboxhaspendingrnfalse\ifintranslation\lawboxtransfont\else\lawboxsourcefont\fi\RaggedRight\setlength{\parindent}{0pt}\setlength{\parskip}{0pt}\ifintranslation\setstretch{\lawboxtransstretch}\else\setstretch{\lawboxsourcestretch}\fi},
    after upper={\global\lawboxhaspendingrnfalse\global\inlawboxfalse},
    borderline west={0.28pt}{0pt}{sourceline!85!black},
    borderline east={0.28pt}{0pt}{sourceline!85!black},
  },
  lawboxpartfirst/.style={
    lawboxpartbase,
    top=0.62em,
    bottom=0.04em,
    before skip=\lawboxblockbeforeskip,
    borderline north={0.28pt}{0pt}{sourceline!85!black},
  },
  lawboxpartmiddle/.style={
    lawboxpartbase,
    top=0.04em,
    bottom=0.04em,
  },
  lawboxpartlast/.style={
    lawboxpartbase,
    top=0.04em,
    bottom=0.62em,
    after skip=0pt,
    borderline south={0.28pt}{0pt}{sourceline!85!black},
  }
}
\newtcolorbox{lawboxpartinner}[1]{#1}
\newtcolorbox{termboxinner}{
  caseboxbase,
  colframe=noteblue!70!black,
  colback=casepaper,
  before upper={\songfont\setlength{\parindent}{0pt}\setlength{\parskip}{0.35em}}
}
\newtcolorbox{deboxinner}{
  caseboxbase,
  colframe=notegray!72!black,
  colback=white,
  left=1.05em,
  right=1.05em,
  top=0.62em,
  bottom=0.62em,
  before upper={\global\inlawboxtrue\global\lawboxfirstitemtrue\global\lawboxhaspendingrnfalse\small\setlength{\parindent}{0pt}\setlength{\parskip}{0.35em}},
  after upper={\global\lawboxhaspendingrnfalse\global\inlawboxfalse}
}
\NewEnviron{lawbox}{
  \begin{lawboxinner}\BODY\end{lawboxinner}
}
\NewEnviron{lawboxpart}[2]{
  \begin{lawboxpartinner}{lawboxpart#1,equal height group=#2}\BODY\end{lawboxpartinner}%
  \ifstrequal{#1}{last}{\global\lawboxpartendedblocktrue}{}%
}
\NewEnviron{termbox}{
  \par\addvspace{0.55em}
  {\color{noteblue!70!black}\rule{\linewidth}{0.28pt}}\par
  \begingroup
    \songfont\small\color{caseink}
    \setlength{\parindent}{0pt}
    \setlength{\parskip}{0.24em}
    \BODY
    \par
  \endgroup
  {\color{noteblue!70!black}\rule{\linewidth}{0.28pt}}\par
  \addvspace{0.55em}
}
\NewEnviron{debox}{
  \begin{deboxinner}\BODY\end{deboxinner}
}
\providecommand{\paragraphmark}[1]{}
\newcommand{\lawboxtitleafter}{\addvspace{0.06em}}
\newcommand{\lawtitle}[1]{\lawboxprespace{0.22em}{\lawboxtransfont\bfseries\lawboxuserandnumber #1}\par\lawboxtitleafter}
\newcommand{\absno}[2]{\lawboxprespace{0.04em}\noindent\lawboxuserandnumber\hangindent=2.6em\hangafter=1\makebox[2.6em][l]{(#1)}#2\par}
\newcommand{\nrno}[2]{\lawboxprespace{0pt}\noindent\lawboxuserandnumber\hspace*{2.6em}\hangindent=4.4em\hangafter=1\makebox[1.8em][l]{#1.}#2\par}
\newcommand{\letterno}[2]{\lawboxprespace{0pt}\noindent\lawboxuserandnumber\hspace*{4.4em}\hangindent=6.0em\hangafter=1\makebox[1.6em][l]{#1)}#2\par}
\newcommand{\sourcelawtitle}[1]{\lawboxprespace{0.22em}{\lawboxsourcefont\bfseries\lawboxuserandnumber #1}\par\lawboxtitleafter}
\newcommand{\sourcelawsubtitle}[1]{\lawboxprespace{0.04em}{\lawboxsourcefont\bfseries\lawboxuserandnumber #1}\par\lawboxtitleafter}
\newcommand{\sourceabsno}[2]{\lawboxprespace{0.04em}\noindent\lawboxuserandnumber\hangindent=2.45em\hangafter=1\makebox[2.45em][l]{(#1)}#2\par}
\newcommand{\sourcenrno}[2]{\lawboxprespace{0pt}\noindent\lawboxuserandnumber\hspace*{2.45em}\hangindent=4.25em\hangafter=1\makebox[1.8em][l]{#1.}#2\par}
\newcommand{\sourceletterno}[2]{\lawboxprespace{0pt}\noindent\lawboxuserandnumber\hspace*{4.25em}\hangindent=5.85em\hangafter=1\makebox[1.6em][l]{#1)}#2\par}
\newcommand{\sourcequotepara}[1]{\lawboxprespace{0.04em}\noindent\lawboxuserandnumber\hangindent=1.2em\hangafter=1 #1\par}
\newcommand{\sourcelawline}[1]{\lawboxprespace{0.04em}\noindent\lawboxuserandnumber #1\par}
\newcommand{\lawline}[1]{\lawboxprespace{0.04em}\noindent\lawboxuserandnumber #1\par}
\newcommand{\pendingtranslation}{\par{\songfont\footnotesize\color{pendingink}（中文譯文待補。）}\par}
\newcommand{\tocpart}[1]{\par\addvspace{0.52em}{\footnotesize\bfseries\color{pendingink}#1\par}\addvspace{0.16em}}
\newcommand{\tocdashline}{\par\addvspace{0.48em}\noindent{\color{notegray}\xleaders\hbox to 0.82em{\hfil\rule{0.42em}{0.28pt}\hfil}\hskip\linewidth}\par\addvspace{0.38em}}
% \tocpanel{font-cmd}{content}：將目錄置中於所在欄寬（雙語欄適用）。
% A3 橫式使用 0.58\linewidth，A4 橫式使用 0.84\linewidth，
% 使兩種紙張下目錄欄內容實際寬度相近（均約 100–105 mm）。
\newcommand{\tocpanel}[2]{%
  \makebox[\linewidth][c]{%
    \ifx\bilingualpaper\bilingualpaperaThree
      \begin{minipage}{0.58\linewidth}%
    \else
      \begin{minipage}{0.84\linewidth}%
    \fi
    \toccolstyle
    #1#2%
    \end{minipage}%
  }%
}
% ===== 首頁簡目錄的兩層 description list 縮排 =====
% enumitem 的 description list 可把每一列拆成：
%   [label]  labelsep  body text
% 這裡的「起點」都以所在 minipage/欄位的左緣為 0。
%
% 核心規則：
%   labelindent = label 欄位從左緣開始的位置。
%   labelwidth  = label 欄位本身寬度；標籤太長（如 Abw. Meinung）就加大它。
%   labelsep    = label 欄位右緣到正文的距離。
%   leftmargin  = 正文 body text 的起點。判斷層級視覺時主要看這個。
%
% 所以：
%   想讓「Begründung / 判決理由」整列正文往右或往左：改 tocitems 的 leftmargin。
%   想讓「A. Verfahren und Gegenstand / A. 程序與審查標的」正文往右或往左：
%     改 tocsubitems 的 leftmargin。
%   想讓 A./B./C. 這些標籤本身往右或往左，但正文不動：改 tocsubitems 的 labelindent。
%   想讓 A./B./C. 與正文間距變大或變小：改 tocsubitems 的 labelsep。
%
% 層級原則：
%   tocsubitems.leftmargin 必須大於 tocitems.leftmargin，
%   否則子層級正文會看起來比「Gründe / 理由」還靠左。
\newlist{tocitems}{description}{1}
\setlist[tocitems]{%
  labelindent=0pt,    % 第一層 label 欄位起點；0pt 表示貼齊目錄欄左緣。
  leftmargin=2.54cm,  % 第一層正文起點；例如 Begründung / 判決理由 的左緣。
  labelwidth=2.16cm,  % 第一層 label 欄位寬度；需容納 Leitsätze、Abw. Meinung、不同意見等。
  labelsep=0.32cm,    % 第一層 label 與正文的水平間距；只改間距，不改正文起點。
  itemsep=0.28em,     % 第一層各項之間的垂直距離。
  topsep=0.20em,      % 第一層 list 上下與前後內容的垂直距離。
  parsep=0pt,         % 同一 item 內段落之間的距離；目錄項通常不需要。
  partopsep=0pt,      % list 若緊接段落開始時的額外距離；關閉以免目錄高度飄動。
  align=left,         % label 欄內左對齊；長短標籤不靠右貼齊。
  font=\normalfont\bfseries % label 的字型；正文不受這行影響。
}
\newlist{tocsubitems}{description}{1}
\setlist[tocsubitems]{%
  labelindent=1cm,    % 第二層 label 起點；只控制 A./B./C. 本身的位置。
  leftmargin=3.00cm,  % 第二層正文起點；必須大於 2.54cm，才會比 Begründung / 判決理由更右。
  labelwidth=0.5cm,   % 第二層 label 欄寬；A./B./C. 很短，可小於第一層。
  labelsep=1.5cm,    % 第二層 label 與正文間距；A. 與 Verfahren / 程序 的距離看這裡。
  itemsep=0.12em,     % 第二層各項較密，避免 A.–E. 佔太高。
  topsep=0.04em,      % 第二層 list 與 Gründe / 理由之間的垂直距離。
  parsep=0pt,         % 同一第二層 item 內段落距離；目錄項通常不需要。
  partopsep=0pt,      % 關閉額外段落起始距離，避免版面不穩。
  align=left,         % A./B./C. label 欄內左對齊。
  font=\normalfont\bfseries, % A./B./C. label 加粗；正文不受這行影響。
  before={}           % 不在第二層 list 前自動插入額外內容。
}
\newlist{termitems}{description}{1}
\ifbilingualcolumns
  \setlist[termitems]{%
    leftmargin=5.2cm,
    labelwidth=4.8cm,
    labelsep=0.35cm,
    itemsep=0.16em,
    topsep=0.2em,
    parsep=0pt,
    partopsep=0pt,
    font=\normalfont\bfseries,
    style=nextline
  }
\else
  \setlist[termitems]{%
    leftmargin=0pt,
    labelwidth=0pt,
    labelsep=0pt,
    itemsep=0.24em,
    topsep=0.2em,
    parsep=0pt,
    partopsep=0pt,
    font=\normalfont\bfseries,
    style=nextline
  }
\fi
\newlist{abbritems}{description}{1}
\setlist[abbritems]{%
  leftmargin=2.38cm,
  labelwidth=2.02cm,
  labelsep=0.28cm,
  itemsep=0.16em,
  topsep=0.2em,
  parsep=0pt,
  partopsep=0pt,
  align=left,
  font=\normalfont\bfseries
}
\ifbilingualcolumns
  \newenvironment{termcolumns}{\begin{multicols}{2}}{\end{multicols}}
\else
  \newenvironment{termcolumns}{}{}
\fi

% ===== 編者註腳開關 =====
% 一般版預設不顯示編者註，避免正文腳註過密。
% 若開啟 \classroommodeon，GG/條文編者註仍會自動顯示，無需另開本開關。
% 若一般版也要顯示編者註，可在單案檔 \input 後加 \showeditornotestrue。
\newif\ifshoweditornotes
\showeditornotesfalse

% ===== 目錄排版輔助 =====
\newcommand{\toccolstyle}{\raggedright\arraybackslash}
\newcommand{\tocsingleindent}{\leftskip=1.45cm\rightskip=1.2cm}

% ===== 行內引文環境（德文原文與中文譯文各一，帶左側灰色邊線）=====
\newcommand{\quoteparbreak}{\par\addvspace{0.48em}}
\newcommand{\sourcequoteinner}[1]{%
  \begin{tcolorbox}[
    enhanced, breakable, width=\dimexpr\linewidth-0.8em\relax, frame hidden,
    colback=white, coltext=caseink, boxrule=0pt,
    borderline west={0.55pt}{0.88em}{notegray},
    left=1.78em, right=0.72em, top=0.18em, bottom=0.18em,
    boxsep=0pt, before skip=0.24em, after skip=0.24em,
    before upper={\normalfont\footnotesize\color{caseink}\setlength{\parindent}{0pt}\setlength{\parskip}{0.34em}\raggedright\setlength{\rightskip}{0pt plus 1em}}
  ]%
  #1\par
  \end{tcolorbox}%
}
\newcommand{\transquoteinner}[1]{%
  \begin{tcolorbox}[
    enhanced, breakable, width=\dimexpr\linewidth-0.8em\relax, frame hidden,
    colback=white, coltext=caseink, boxrule=0pt,
    borderline west={0.55pt}{0.88em}{notegray},
    left=1.78em, right=0.72em, top=0.18em, bottom=0.18em,
    boxsep=0pt, before skip=0.24em, after skip=0.24em,
    before upper={\songfont\small\color{caseink}\setlength{\parindent}{0pt}\setlength{\parskip}{0.34em}\raggedright\setlength{\rightskip}{0pt plus 1em}}
  ]%
  #1\par
  \end{tcolorbox}%
}

% ===== 大綱（Gliederung）Rn 輔助命令 =====
\newcommand{\outrn}[2]{\enspace{\normalfont\footnotesize\color{pendingink}(Rn\,#1–#2)}}
\newcommand{\outrnsingle}[1]{\enspace{\normalfont\footnotesize\color{pendingink}(Rn\,#1)}}
\newcommand{\outrnzh}[2]{\enspace{\normalfont\footnotesize\color{pendingink}（第\,#1–#2\,段）}}
\newcommand{\outrnzhs}[1]{\enspace{\normalfont\footnotesize\color{pendingink}（第\,#1\,段）}}
% ===== 大綱雙語對照表格輔助命令（三欄：段碼 │ 德文 │ 中文）=====
% 欄 1：段碼（由欄定義自動格式化：小字、等寬、右對齊、pendingink）
% 欄 2：德文標籤＋正文
% 欄 3：中文標籤＋正文
%
% 無段碼的列（\omainrow、\osecrow），欄 1 以 & 空置。
% 有段碼的列（\osecrowrn、\osubrow、\ointrow），#1 為預格式化字串
%   例：\osubrow{2–49}{I.}{...}{I.}{...}
%       \ointrow{107}{...}{...}
%
% \opartrow：段組標題（橫跨三欄）
\newcommand{\opartrow}[2]{%
  \noalign{\vspace{0.30em}}%
  \multicolumn{3}{@{}l@{}}{%
    \footnotesize\bfseries\color{pendingink}#1\hfill\songfont#2%
  }\\[0.06em]%
}
% \odashrow：虛線分隔（橫跨三欄）
\newcommand{\odashrow}{%
  \noalign{\vspace{0.28em}}%
  \multicolumn{3}{@{}l@{}}{\color{notegray}%
    \xleaders\hbox to 0.82em{\hfil\rule{0.42em}{0.28pt}\hfil}\hskip 0.88\linewidth}%
  \\[0.24em]%
}
% \odissentpart：不同意見書區段標頭。比一般 \opartrow 更像附錄性轉場。
\newcommand{\odissentpart}[2]{%
  \noalign{\vspace{0.42em}}%
  \multicolumn{3}{@{}p{0.88\textwidth}@{}}{%
    \color{notegray}\rule{\linewidth}{0.18pt}%
  }\\[-0.18em]%
  \multicolumn{3}{@{}p{0.88\textwidth}@{}}{%
    \makebox[\linewidth][s]{%
      \footnotesize\bfseries\color{pendingink}#1%
      \hfill{\songfont #2}%
    }%
  }\\[-0.02em]%
}
% \odissentopinion：不同意見書的署名與一行摘要，右欄保留段碼範圍。
\newcommand{\odissentopinion}[5]{%
  \footnotesize\hangindent=0.52cm\hangafter=1%
    {\bfseries\color{caseink}#2}\enspace{\color{notegray}\textperiodcentered}\enspace\color{caseink}#3%
  &\footnotesize\hangindent=0.52cm\hangafter=1%
    {\songfont\bfseries\color{caseink}#4}\enspace{\color{notegray}\textperiodcentered}\enspace\songfont\color{caseink}#5%
  &#1%
  \\[0.12em]%
}
% \omainrow：頂層項目，無段碼（Leitsätze、Urteil、Gründe …）
\newcommand{\omainrow}[4]{%
  \footnotesize\hangindent=1.92cm\hangafter=1%
    \makebox[1.92cm][l]{\bfseries\color{caseink}#1}\color{caseink}#2%
  &\footnotesize\hangindent=1.92cm\hangafter=1%
    \makebox[1.92cm][l]{\bfseries\color{caseink}#3}\songfont\color{caseink}#4%
  &%
  \\[0.15em]%
}
% \omainrowrn：頂層項目，有段碼（不同意見等標籤寬於 \osecrow 框者）
\newcommand{\omainrowrn}[5]{%
  \footnotesize\hangindent=1.92cm\hangafter=1%
    \makebox[1.92cm][l]{\bfseries\color{caseink}#2}\color{caseink}#3%
  &\footnotesize\hangindent=1.92cm\hangafter=1%
    \makebox[1.92cm][l]{\bfseries\color{caseink}#4}\songfont\color{caseink}#5%
  &#1%
  \\[0.15em]%
}
% \osecrow：一級子項，無段碼（A.、C.、D. 等有子項者）
\newcommand{\osecrow}[4]{%
  \footnotesize\hspace{1.10cm}\hangindent=1.98cm\hangafter=1%
    \makebox[0.86cm][l]{\bfseries\color{caseink}#1}\color{caseink}#2%
  &\footnotesize\hspace{1.10cm}\hangindent=1.98cm\hangafter=1%
    \makebox[0.86cm][l]{\bfseries\color{caseink}#3}\songfont\color{caseink}#4%
  &%
  \\[0.11em]%
}
% \osecrowrn：一級子項，有段碼（B.、E. 等完整段落範圍者）
% #1=段碼字串（如 101–106）, #2=DE標籤, #3=DE正文, #4=ZH標籤, #5=ZH正文
\newcommand{\osecrowrn}[5]{%
  \footnotesize\hspace{1.10cm}\hangindent=1.98cm\hangafter=1%
    \makebox[0.86cm][l]{\bfseries\color{caseink}#2}\color{caseink}#3%
  &\footnotesize\hspace{1.10cm}\hangindent=1.98cm\hangafter=1%
    \makebox[0.86cm][l]{\bfseries\color{caseink}#4}\songfont\color{caseink}#5%
  &#1%
  \\[0.11em]%
}
% \osubrow：二級子項，有段碼範圍（I.、II.、A.I. …）
% #1=段碼字串, #2=DE標籤, #3=DE正文, #4=ZH標籤, #5=ZH正文
\newcommand{\osubrow}[5]{%
  \footnotesize\hspace{1.40cm}\hangindent=2.30cm\hangafter=1%
    \makebox[0.90cm][l]{\bfseries\color{caseink}#2}\color{caseink}#3%
  &\footnotesize\hspace{1.40cm}\hangindent=2.30cm\hangafter=1%
    \makebox[0.90cm][l]{\bfseries\color{caseink}#4}\songfont\color{caseink}#5%
  &#1%
  \\[0.09em]%
}
% \ointrow：引入段落，有單一段碼，無標籤
% #1=段碼字串, #2=DE正文, #3=ZH正文
\newcommand{\ointrow}[3]{%
  \footnotesize\hspace{0.52cm}\hangindent=0.52cm\hangafter=1\color{caseink}#2%
  &\footnotesize\hspace{0.52cm}\hangindent=0.52cm\hangafter=1\songfont\color{caseink}#3%
  &#1%
  \\[0.09em]%
}
% \osubsubrow：三級子項（1.、2.、3.）
% #1=段碼字串, #2=DE標籤, #3=DE正文, #4=ZH標籤, #5=ZH正文
\newcommand{\osubsubrow}[5]{%
  \footnotesize\hspace{1.80cm}\hangindent=2.48cm\hangafter=1%
    \makebox[0.68cm][l]{\bfseries\color{caseink}#2}\color{caseink}#3%
  &\footnotesize\hspace{1.80cm}\hangindent=2.48cm\hangafter=1%
    \makebox[0.68cm][l]{\bfseries\color{caseink}#4}\songfont\color{caseink}#5%
  &#1%
  \\[0.08em]%
}
% \osubsubsubrow：四級子項（a）、b））
% #1=段碼字串, #2=DE標籤, #3=DE正文, #4=ZH標籤, #5=ZH正文
\newcommand{\osubsubsubrow}[5]{%
  \footnotesize\hspace{2.10cm}\hangindent=2.68cm\hangafter=1%
    \makebox[0.58cm][l]{\bfseries\color{caseink}#2}\color{caseink}#3%
  &\footnotesize\hspace{2.10cm}\hangindent=2.68cm\hangafter=1%
    \makebox[0.58cm][l]{\bfseries\color{caseink}#4}\songfont\color{caseink}#5%
  &#1%
  \\[0.07em]%
}
% \osubsubsubsubrow：五級子項（aa）、bb））
% 標籤框 0.62cm：足以容納「aa)」在 footnotesize bold 下（約 0.50cm）並留有間距
% #1=段碼字串, #2=DE標籤, #3=DE正文, #4=ZH標籤, #5=ZH正文
\newcommand{\osubsubsubsubrow}[5]{%
  \footnotesize\hspace{2.38cm}\hangindent=3.06cm\hangafter=1%
    \makebox[0.68cm][l]{\bfseries\color{caseink}#2}\color{caseink}#3%
  &\footnotesize\hspace{2.38cm}\hangindent=3.06cm\hangafter=1%
    \makebox[0.68cm][l]{\bfseries\color{caseink}#4}\songfont\color{caseink}#5%
  &#1%
  \\[0.06em]%
}
% \outlinepage：雙語對照大綱頁包裝環境（longtable 自動跨頁）
% 三欄：德文欄 + 中文欄 + 段碼欄（最右，不折行）
% 段碼欄寬度：1.80cm，足以容納最長段碼如「309–348」
% DE/ZH 欄寬：(0.87\textwidth - 2.08cm) / 2
%   驗算：2×欄 + 0.05\textwidth + 0.28cm gap + 1.80cm = 0.87tw - 2.08cm + 0.05tw + 2.08cm = 0.92tw ✓
\newenvironment{outlinepage}{%
  \vspace*{0.40cm}%
  \setlength{\LTleft}{0.04\textwidth}%
  \setlength{\LTright}{0.04\textwidth}%
  \setlength{\tabcolsep}{0pt}%
  % 大綱列距由 longtable 的 \arraystretch 與各 row macro 結尾的 \\[...em] 共同決定。
  % 這裡控制「每一列本身的表格 strut 高度」；數值越大，A./I./1. 各項之間越像有空行。
  % A3 原本用 0.62，視覺上沒有空行；A4 也沿用同值，避免切到 A4 後項目被撐開。
  % 若日後覺得大綱太擠，只微調這個值（例如 0.68），不要改各 \omainrow/\osubrow 的縮排。
  \renewcommand{\arraystretch}{0.62}%
  \begin{longtable}{%
    >{\vspace{0pt}\raggedright\arraybackslash}p{\dimexpr(0.87\textwidth-2.08cm)/2\relax}%
    @{\hspace{0.025\textwidth}}%
    !{\color{notegray}\vrule width 0.12pt}%
    @{\hspace{0.025\textwidth}}%
    >{\vspace{0pt}\raggedright\arraybackslash}p{\dimexpr(0.87\textwidth-2.08cm)/2\relax}%
    @{\hspace{0.28cm}}%
    >{\vspace{0pt}\footnotesize\normalfont\color{pendingink}\raggedright\arraybackslash}p{1.80cm}%
    @{}%
  }%
  % 首頁欄標籤：不要橫跨置中；分別放在德文欄與中文欄的實際欄位起點。
  \footnotesize\bfseries\color{sourceink}Gliederung%
  &\footnotesize\bfseries\songfont\color{sourceink}大綱%
  &%
  \\[0.36em]%
  \endfirsthead%
  % 續頁欄標籤：同樣分欄定位，以灰色細體區分；無需標注「續」。
  \footnotesize\color{pendingink}Gliederung%
  &\footnotesize\songfont\color{pendingink}大綱%
  &%
  \\[0.24em]%
  \endhead%
}{%
  \end{longtable}%
}
 載入。
% 不含：\documentclass、\documentversion、\ggversionde/zh、\tocde/\toczh。
% 日期與首頁版本列由本檔自動產生；各文件只需手動指定 \documentversion。

\usepackage{xeCJK}
\usepackage{fontspec}
\usepackage{geometry}
\usepackage{setspace}
\usepackage{hyperref}
\usepackage{xurl}
\usepackage{bookmark}
\usepackage{enumitem}
\usepackage{xcolor}
\usepackage{titlesec}
\usepackage{array}
\usepackage{longtable}
\usepackage{tcolorbox}
\usepackage{environ}
\usepackage{pdflscape}
\usepackage{etoolbox}
\usepackage{ragged2e}
\usepackage{paracol}
\usepackage{multicol}
\usepackage{needspace}
\usepackage{fancyhdr}
\usepackage{tikz}
\usepackage{zref-abspage}
\tcbuselibrary{breakable,skins}
\usetikzlibrary{calc}
\usepackage{pgfcalendar}

\hypersetup{
  unicode=true,
  bookmarksopen=true,
  bookmarksnumbered=false,
  hidelinks
}

% ===== 自動推導，不要手動改下面 =====
% （\docmode 與 \bilingualpaper 由各文件在 % bverfge_layout.tex
% 共用排版基礎設施，供 Schwangerschaftsabbruch I 與 II 共享。
% 由各文件以 \documentclass 後 \input{../../共用LaTeX/bverfge_layout} 載入。
% 不含：\documentclass、\documentversion、\ggversionde/zh、\tocde/\toczh。
% 日期與首頁版本列由本檔自動產生；各文件只需手動指定 \documentversion。

\usepackage{xeCJK}
\usepackage{fontspec}
\usepackage{geometry}
\usepackage{setspace}
\usepackage{hyperref}
\usepackage{xurl}
\usepackage{bookmark}
\usepackage{enumitem}
\usepackage{xcolor}
\usepackage{titlesec}
\usepackage{array}
\usepackage{longtable}
\usepackage{tcolorbox}
\usepackage{environ}
\usepackage{pdflscape}
\usepackage{etoolbox}
\usepackage{ragged2e}
\usepackage{paracol}
\usepackage{multicol}
\usepackage{needspace}
\usepackage{fancyhdr}
\usepackage{tikz}
\usepackage{zref-abspage}
\tcbuselibrary{breakable,skins}
\usetikzlibrary{calc}
\usepackage{pgfcalendar}

\hypersetup{
  unicode=true,
  bookmarksopen=true,
  bookmarksnumbered=false,
  hidelinks
}

% ===== 自動推導，不要手動改下面 =====
% （\docmode 與 \bilingualpaper 由各文件在 \input{bverfge_layout} 之前設定；
%   預設 chinese / a4）
\ifdefined\docmode\else\def\docmode{chinese}\fi
\ifdefined\bilingualpaper\else\def\bilingualpaper{a4}\fi
\newif\ifshoworiginal
\newif\ifshowtranslation
\newif\ifbilinguallandscape
\newif\ifbilingualcolumns
\newif\ifintranslation
\newif\iffirstbilingualsection

\showoriginalfalse
\showtranslationfalse
\bilinguallandscapefalse
\bilingualcolumnsfalse
\intranslationfalse
\firstbilingualsectionfalse

\def\modechinese{chinese}
\def\modegerman{german}
\def\modebilingual{bilingual}
\def\bilingualpaperaThree{a3}
\def\bilingualpaperaFour{a4}

\ifx\docmode\modechinese
  \showtranslationtrue
\fi

\ifx\docmode\modegerman
  \showoriginaltrue
\fi

\ifx\docmode\modebilingual
  \showoriginaltrue
  \showtranslationtrue
  \bilinguallandscapetrue
  \bilingualcolumnstrue
\fi

% 編排模式說明：
% chinese  → \showoriginalfalse  \showtranslationtrue  （A4 直式，純中文）
% german   → \showoriginaltrue   \showtranslationfalse （A4 直式，純德文）
% bilingual→ \showoriginaltrue   \showtranslationtrue  \bilingualcolumnstrue（A3/A4 橫式對照）

\ifbilingualcolumns
  \ifx\bilingualpaper\bilingualpaperaFour
    \geometry{
      a4paper,
      landscape,
      left=1.25cm,
      right=2.25cm,
      top=1.25cm,
      bottom=1.75cm,
      footskip=8.5mm,
      marginparwidth=0.75cm,
      marginparsep=0.25cm
    }
  \else
    \geometry{
      a3paper,
      landscape,
      left=1.55cm,
      right=2.75cm,
      top=1.55cm,
      bottom=2.05cm,
      footskip=9.5mm,
      marginparwidth=0.9cm,
      marginparsep=0.35cm
    }
  \fi
\else
\geometry{
  a4paper,
  portrait,
  left=2.2cm,
  right=2.6cm,
  top=2.0cm,
  bottom=2.0cm,
  marginparwidth=0.8cm,
  marginparsep=0.3cm
}
\fi
% ===== 時間戳基礎設施 =====
\newcount\stamptimeval
\newcount\stampminuteval
\newcommand{\stamphh}{%
  \stamptimeval=\time
  \divide\stamptimeval by 60
  \ifnum\stamptimeval<10 0\fi
  \the\stamptimeval}
\newcommand{\stampmm}{%
  \stamptimeval=\time
  \stampminuteval=\time
  \divide\stampminuteval by 60
  \multiply\stampminuteval by 60
  \advance\stamptimeval by -\stampminuteval
  \ifnum\stamptimeval<10 0\fi
  \the\stamptimeval}
\newcommand{\stampwkde}[1]{%
  \ifcase#1Montag\or Dienstag\or Mittwoch\or Donnerstag\or Freitag\or Samstag\or Sonntag\fi}
\newcommand{\stampwkzh}[1]{%
  \ifcase#1 一\or 二\or 三\or 四\or 五\or 六\or 日\fi}
\newcommand{\stampmonthde}[1]{%
  \ifcase#1\or Januar\or Februar\or März\or April\or Mai\or Juni\or
  Juli\or August\or September\or Oktober\or November\or Dezember\fi}
\newcount\stampjdval
\newcount\stampwdval
\pgfcalendardatetojulian{\the\year-\the\month-\the\day}{\stampjdval}
\pgfcalendarjuliantoweekday{\the\stampjdval}{\stampwdval}
\newcommand{\timestampzh}{%
  \songfont 最後修改：\the\year 年\the\month 月\the\day 日%
  % （星期\stampwkzh{\the\stampwdval}）%
  % \space\stamphh:\stampmm
  }

\newcommand{\documentdatezh}{\the\year 年 \the\month 月 \the\day 日}
\newcommand{\documentdatede}{\the\day. \stampmonthde{\the\month} \the\year}
\newcommand{\bverfgetranslationnote}{%
  {\songfont 翻譯說明：初譯使用 Claude Code 與 GPT 協助迭代；仍請以德文原文為準。}%
}
\newcommand{\bverfgecourseinfo}{%
  {\songfont 課程：114 學年度第 2 學期；憲法專題研究（四）；授課老師：湯德宗教授；導讀日期：115 年 6 月 1 日。}%
}
\newcommand{\bverfgeeditorinfo}{%
  {\songfont 編者：王逸帆（R10A21126），國立台灣大學法律系研究所碩士班。}%
}
\newcommand{\bverfgetitleversionline}{%
  \ifbilingualcolumns
    Fassung \documentversion\enspace\textperiodcentered\enspace Stand: \documentdatede
    \quad
    {\songfont 版本 \documentversion\enspace\textperiodcentered\enspace 修訂日期：\documentdatezh\par}%
  \else
    \ifshowtranslation
      {\songfont 版本 \documentversion\enspace\textperiodcentered\enspace 修訂日期：\documentdatezh\par}%
    \else
      Fassung \documentversion\enspace\textperiodcentered\enspace Stand: \documentdatede\par
    \fi
  \fi
}
\newcommand{\bverfgetitlecornerinfo}{%
  \ifshowtranslation
    \vspace{0.72em}%
    \noindent\hfill
    \begin{minipage}{0.66\textwidth}
      \raggedleft\tiny\color{pendingink}\songfont
      \bverfgecourseinfo\par
      \bverfgeeditorinfo
    \end{minipage}%
  \fi
}
\newcommand{\bverfgetitlemeta}{%
  \vfill
  \begin{center}%
    \footnotesize\color{pendingink}%
    \bverfgetitleversionline
    \ifshowtranslation
      \vspace{0.28em}{\scriptsize\bverfgetranslationnote\par}%
    \fi
  \end{center}%
  \bverfgetitlecornerinfo
}

\pagestyle{fancy}
\fancyhf{}
\renewcommand{\headrulewidth}{0pt}
\renewcommand{\footrulewidth}{0pt}
\fancyfoot[C]{\thepage}
\ifbilingualcolumns
  \fancyfoot[R]{\scriptsize\color{pendingink}\timestampzh}%
\else
  \ifshoworiginal
  \else
    \fancyfoot[R]{\scriptsize\color{pendingink}\timestampzh\hspace{0.45cm}}%
  \fi
\fi
\setmainfont{Times New Roman}
\setsansfont{Helvetica Neue}
\setmonofont{Menlo}
\IfFileExists{/Users/iw/Library/Fonts/kaiu.ttf}{
  \setCJKmainfont[Path=/Users/iw/Library/Fonts/,AutoFakeBold=4]{kaiu.ttf}
}{
  \IfFontExistsTF{標楷體}{
    \setCJKmainfont[AutoFakeBold=4]{標楷體}
  }{
    \setCJKmainfont{Songti TC}
  }
}
\IfFontExistsTF{Songti TC}{
  \setCJKfamilyfont{song}{Songti TC}
  \setCJKmonofont{Songti TC}
}{
  \setCJKfamilyfont{song}[Path=/Users/iw/Library/Fonts/,AutoFakeBold=4]{kaiu.ttf}
  \setCJKmonofont[Path=/Users/iw/Library/Fonts/,AutoFakeBold=4]{kaiu.ttf}
}
\newcommand{\songfont}{\CJKfamily{song}}

% ===== 封面／目錄專用打字機字體（僅限封面、目錄頁使用，不影響正文）=====
\IfFontExistsTF{Radio Newsman}{%
  \newfontfamily\bverfgetitlede{Radio Newsman}%
}{%
  \newfontfamily\bverfgetitlede{Courier New}%
}
\IfFontExistsTF{Erikas Farbband}{%
  \newfontfamily\bverfgebodyde{Erikas Farbband}%
}{%
  \let\bverfgebodyde\bverfgetitlede
}
\IfFontExistsTF{Maoken Heavy Labourer}{%
  \setCJKfamilyfont{bverfgetitlecjk}{Maoken Heavy Labourer}%
}{%
  \setCJKfamilyfont{bverfgetitlecjk}{Songti TC}%
}
\IfFontExistsTF{Huiwen-mincho}{%
  \setCJKfamilyfont{bverfgebodycjk}{Huiwen-mincho}%
}{%
  \setCJKfamilyfont{bverfgebodycjk}{Songti TC}%
}
\newcommand{\bverfgetitlezh}{\bverfgetitlede\CJKfamily{bverfgetitlecjk}}
\newcommand{\bverfgebodyzh}{\bverfgebodyde\CJKfamily{bverfgebodycjk}}

\setstretch{1.35}
\setlist{itemsep=0.2em, topsep=0.4em}
\setlength{\parindent}{0pt}
\setlength{\parskip}{0.45em}
\raggedbottom
\setcounter{secnumdepth}{0}
\setcounter{tocdepth}{2}

\newcommand{\lawboxsourcefont}{\footnotesize}
\newcommand{\lawboxtransfont}{\footnotesize}
\newcommand{\lawboxsourcestretch}{1.02}
\newcommand{\lawboxtransstretch}{0.92}

\titleformat{\section}{\centering\Large\bfseries\songfont}{\thesection}{1em}{}
\titleformat{\subsection}{\large\bfseries\songfont}{\thesubsection}{1em}{}
\titleformat{\subsubsection}{\normalsize\bfseries\songfont}{\thesubsubsection}{1em}{}
\titleformat{\paragraph}{\normalsize\bfseries\songfont}{\theparagraph}{1em}{}
\titlespacing*{\section}{0pt}{1.8em}{1.0em}
\titlespacing*{\subsection}{0pt}{1.2em}{0.6em}
\titlespacing*{\subsubsection}{0pt}{1.0em}{0.45em}
\titlespacing*{\paragraph}{0pt}{0.6em}{0.4em}

\definecolor{noteblue}{HTML}{A9C9F5}
\definecolor{noteteal}{HTML}{AEE1D6}
\definecolor{notegray}{HTML}{D7DADF}
\definecolor{caseink}{HTML}{000000}
\definecolor{casepaper}{HTML}{FCFEFD}
\definecolor{sourcepaper}{HTML}{FAFBF8}
\definecolor{sourceline}{HTML}{D7DED8}
\definecolor{sourceink}{HTML}{000000}
\definecolor{pendingink}{HTML}{000000}

\tcbset{
  caseboxbase/.style={
    enhanced,
    breakable,
    width=0.985\linewidth,
    colback=casepaper,
    coltitle=caseink,
    fonttitle=\bfseries\songfont,
    boxrule=0.28pt,
    arc=1.2mm,
    left=2.4mm,
    right=2.4mm,
    top=1.5mm,
    bottom=1.5mm,
    boxsep=1.5mm,
    before skip=0.65em,
    after skip=0.65em,
    before upper={\setlength{\parindent}{0pt}\setlength{\parskip}{0.35em}},
  }
}

\newcommand{\bilingualstart}{}
\newcommand{\bilingualstop}{}
\ifbilingualcolumns
  \renewcommand{\bilingualstart}{\small\setlength{\columnsep}{1.25cm}\setlength{\columnseprule}{0.12pt}\columnratio{0.5,0.5}\multicolumnfootnotes\flushbottom\sloppy\interlinepenalty=80\clubpenalty=800\widowpenalty=800\displaywidowpenalty=800\begin{paracol}{2}\global\intranslationfalse\global\firstbilingualsectiontrue{\footnotesize\color{sourceink}Deutsch\par}\switchcolumn{\footnotesize\songfont\color{sourceink}中文譯文\par}\switchcolumn*}
  \renewcommand{\bilingualstop}{\ifintranslation\switchcolumn*\global\intranslationfalse\fi\end{paracol}\clearpage\raggedbottom}
\fi
\newif\ifrandnumbers
\randnumbersfalse
\newcounter{randnumber}
\newcounter{randnumberlocal}
\newcounter{lastsourcerandstart}
\newif\iflastsourcehasrandnumbers
\lastsourcehasrandnumbersfalse
\newif\ifinlawbox
\inlawboxfalse
\newif\iflawboxhaspendingrn
\lawboxhaspendingrnfalse
\newcommand{\lawboxpendingrn}{}
\newcommand{\startrandnumbers}{\global\randnumberstrue\setcounter{randnumber}{1}\global\lastsourcehasrandnumbersfalse}
\newcommand{\stoprandnumbers}{\global\randnumbersfalse\global\lastsourcehasrandnumbersfalse}
\newlength{\randnumberwidth}
\setlength{\randnumberwidth}{0.95cm}
\newlength{\randnumberrightoffset}
\setlength{\randnumberrightoffset}{0.85cm}
\newlength{\randnumberlawboxrightoffset}
\setlength{\randnumberlawboxrightoffset}{1.40cm}
\newcommand{\randnumberbox}[1]{%
  \leavevmode
  \makebox[0pt][l]{%
    \smash{%
      \tikz[remember picture,overlay,baseline=(rnmark.base)]{%
        \coordinate (rnmark) at (0,0);%
        \ifinlawbox
          \path let \p1=(current page.east), \p2=(rnmark.base) in
            node[
              anchor=base east,
              inner sep=0pt,
              outer sep=0pt,
              text width=\randnumberwidth,
              align=right,
              text=pendingink,
              font=\normalfont\mdseries\scriptsize\ttfamily
            ] at (\x1-\randnumberlawboxrightoffset,\y2) {{\normalfont\mdseries\scriptsize\ttfamily #1}};%
        \else
          \path let \p1=(current page.east), \p2=(rnmark.base) in
            node[
              anchor=base east,
              inner sep=0pt,
              outer sep=0pt,
              text width=\randnumberwidth,
              align=right,
              text=pendingink,
              font=\normalfont\mdseries\scriptsize\ttfamily
            ] at (\x1-\randnumberrightoffset,\y2) {{\normalfont\mdseries\scriptsize\ttfamily #1}};%
        \fi
      }%
    }%
  }%
  \ignorespaces
}
\newcommand{\lawboxstorerandnumber}[1]{%
  \gdef\lawboxpendingrn{#1}%
  \global\lawboxhaspendingrntrue
  \ignorespaces
}
\newcommand{\lawboxuserandnumber}{%
  \iflawboxhaspendingrn
    \randnumberbox{\lawboxpendingrn}%
    \global\lawboxhaspendingrnfalse
  \fi
}
\newcommand{\sourcern}[1]{%
  \ifrandnumbers
    \ifshoworiginal
      \ifshowtranslation\else
        \ifinlawbox\lawboxstorerandnumber{#1}\else\randnumberbox{#1}\fi
      \fi
    \fi
  \fi
}
\newcommand{\transrn}[1]{%
  \ifrandnumbers
    \ifshowtranslation
      \ifinlawbox\lawboxstorerandnumber{#1}\else\randnumberbox{#1}\fi
    \fi
  \fi
}
% ===== 課堂模式：德文原文縮寫註解 =====
% 日常切換只改各判決檔的一行：
%   \classroommodeon        開啟課堂版；註解掉這行即回到一般版。
% 模板預設：
%   1. 課堂模式關閉。
%   2. 課堂模式開啟時，縮寫註解包含「德文全稱＋中文譯名」。
%   3. 課堂模式開啟時，GG/條文編者註仍會自動顯示。
% 進階微調才需要在單案檔覆寫：
%   \classroomabbrzhoff     縮寫註解僅顯示德文全稱。
%   \showeditornotestrue    一般版也顯示編者註。
% 註解策略：同一實際 PDF 頁面中，同一縮寫只在第一次出現處加註。
% 這裡用 zref-abspage 的絕對頁序，不用 \thepage，避免文件內頁碼重置時誤判。
\newif\ifclassroommode
\newif\ifclassroomabbrzh
\classroommodefalse
\classroomabbrzhtrue
\newcommand{\classroommodeon}{\classroommodetrue}
\newcommand{\classroommodeoff}{\classroommodefalse}
\newcommand{\classroomabbrzhon}{\classroomabbrzhtrue}
\newcommand{\classroomabbrzhoff}{\classroomabbrzhfalse}
\newcommand{\abbrnotelabel}{\emph{Abk.}: }
\newcommand{\abbrnote}[3]{%
  \ifclassroommode
    \ifshoworiginal
      \ifintranslation\else
        \footnote{\normalfont\footnotesize\abbrnotelabel #2\ifclassroomabbrzh；{\songfont #3}\fi}%
      \fi
    \fi
  \fi
}
\newcommand{\abbrnoteonce}[4]{%
  #2%
  \ifclassroommode
    \ifshoworiginal
      \ifintranslation\else
        \begingroup
          \edef\abbrcurrentabspage{\number\numexpr\value{abspage}+1\relax}%
          \ifcsname abbrnoteseen@#1@\abbrcurrentabspage\endcsname\else
            \global\expandafter\let\csname abbrnoteseen@#1@\abbrcurrentabspage\endcsname\empty
            \abbrnote{#1}{#3}{#4}%
          \fi
        \endgroup
      \fi
    \fi
  \fi
}
\newcommand{\abbrAbs}{\abbrnoteonce{abs}{Abs.}{Absatz}{項}}
\newcommand{\abbrAAO}{\abbrnoteonce{aao}{a.a.O.}{am angegebenen Ort}{前引處}}
\newcommand{\abbrAE}{\abbrnoteonce{ae}{AE}{Alternativ-Entwurf}{替代草案}}
\newcommand{\abbrArt}{\abbrnoteonce{art}{Art.}{Artikel}{條}}
\newcommand{\abbrAufl}{\abbrnoteonce{aufl}{Aufl.}{Auflage}{版}}
\newcommand{\abbrBd}{\abbrnoteonce{bd}{Bd.}{Band}{卷}}
\newcommand{\abbrBGBl}{\abbrnoteonce{bgbl}{BGBl.}{Bundesgesetzblatt}{聯邦法律公報}}
\newcommand{\abbrBGHSt}{\abbrnoteonce{bghst}{BGHSt}{Entscheidungen des Bundesgerichtshofes in Strafsachen}{聯邦普通法院刑事判決彙編}}
\newcommand{\abbrBRDrucks}{\abbrnoteonce{brdrucks}{BRDrucks.}{Bundesrats-Drucksache}{聯邦參議院文件}}
\newcommand{\abbrBTDrucks}{\abbrnoteonce{btdrucks}{BTDrucks.}{Bundestags-Drucksache}{聯邦議院文件}}
\newcommand{\abbrBundesgesetzbl}{\abbrnoteonce{bundesgesetzbl}{Bundesgesetzbl.}{Bundesgesetzblatt}{聯邦法律公報}}
\newcommand{\abbrBvF}{\abbrnoteonce{bvf}{BvF}{Bundesverfassungsgerichtsverfahren}{聯邦憲法法院程序案號}}
\newcommand{\abbrBvL}{\abbrnoteonce{bvl}{BvL}{Bundesverfassungsgerichtsverfahren}{聯邦憲法法院程序案號}}
\newcommand{\abbrBvR}{\abbrnoteonce{bvr}{BvR}{Bundesverfassungsgerichtsverfahren}{聯邦憲法法院程序案號}}
\newcommand{\abbrBVerfG}{\abbrnoteonce{bverfg}{BVerfG}{Bundesverfassungsgericht}{聯邦憲法法院}}
\newcommand{\abbrBVerfGE}{\abbrnoteonce{bverfge}{BVerfGE}{Entscheidungen des Bundesverfassungsgerichts}{聯邦憲法法院判決集}}
\newcommand{\abbrBVerfGG}{\abbrnoteonce{bverfgg}{BVerfGG}{Bundesverfassungsgerichtsgesetz}{聯邦憲法法院法}}
\newcommand{\abbrDDR}{\abbrnoteonce{ddr}{DDR}{Deutsche Demokratische Republik}{德意志民主共和國；東德}}
\newcommand{\abbrDJT}{\abbrnoteonce{djt}{DJT}{Deutscher Juristentag}{德國法學家大會}}
\newcommand{\abbrDr}{\abbrnoteonce{dr}{Dr.}{Doktor}{Dr.}}
\newcommand{\abbrEuGRZ}{\abbrnoteonce{eugrz}{EuGRZ}{Europäische Grundrechte-Zeitschrift}{歐洲基本權期刊}}
\newcommand{\abbrEV}{\abbrnoteonce{ev}{EV}{Einigungsvertrag}{統一條約}}
\newcommand{\abbrF}{\abbrnoteonce{f}{f.}{folgende}{以下一頁；下一段}}
\newcommand{\abbrFF}{\abbrnoteonce{ff}{ff.}{fortfolgende}{以下數頁；以下各段}}
\newcommand{\abbrGBlDDR}{\abbrnoteonce{gblddr}{GBl. DDR}{Gesetzblatt der Deutschen Demokratischen Republik}{德意志民主共和國法律公報}}
\newcommand{\abbrGG}{\abbrnoteonce{gg}{GG}{Grundgesetz}{基本法}}
\newcommand{\abbrGVBl}{\abbrnoteonce{gvbl}{GVBl.}{Gesetz- und Verordnungsblatt}{法律及命令公報}}
\newcommand{\abbrJR}{\abbrnoteonce{jr}{JR}{Juristische Rundschau}{法學評論}}
\newcommand{\abbrJZ}{\abbrnoteonce{jz}{JZ}{JuristenZeitung}{法學家報}}
\newcommand{\abbrIVm}{\abbrnoteonce{ivm}{i.V.m.}{in Verbindung mit}{結合；連同}}
\newcommand{\abbrMdB}{\abbrnoteonce{mdb}{MdB}{Mitglied des Bundestages}{聯邦議院議員}}
\newcommand{\abbrMwN}{\abbrnoteonce{mwn}{m.w.N.}{mit weiteren Nachweisen}{附進一步文獻／判例出處}}
\newcommand{\abbrNF}{\abbrnoteonce{nf}{n.F.}{neue Fassung}{新版本}}
\newcommand{\abbrAF}{\abbrnoteonce{af}{a.F.}{alte Fassung}{舊版本}}
\newcommand{\abbrNr}{\abbrnoteonce{nr}{Nr.}{Nummer}{款、目或號，依語境}}
\newcommand{\abbrProf}{\abbrnoteonce{prof}{Prof.}{Professor}{教授}}
\newcommand{\abbrRdnr}{\abbrnoteonce{rdnr}{Rdnr.}{Randnummer}{邊碼；段碼}}
\newcommand{\abbrRdnrn}{\abbrnoteonce{rdnrn}{Rdnrn.}{Randnummern}{邊碼；段碼}}
\newcommand{\abbrRGBl}{\abbrnoteonce{rgbl}{RGBl.}{Reichsgesetzblatt}{帝國法律公報}}
\newcommand{\abbrRGSt}{\abbrnoteonce{rgst}{RGSt}{Entscheidungen des Reichsgerichts in Strafsachen}{帝國法院刑事判決彙編}}
\newcommand{\abbrRspr}{\abbrnoteonce{rspr}{Rspr.}{Rechtsprechung}{判例；司法實務}}
\newcommand{\abbrRVO}{\abbrnoteonce{rvo}{RVO}{Reichsversicherungsordnung}{帝國保險法}}
\newcommand{\abbrS}{\abbrnoteonce{s}{S.}{Seite}{頁}}
\newcommand{\abbrSFHG}{\abbrnoteonce{sfhg}{SFHG}{Schwangeren- und Familienhilfegesetz}{孕婦及家庭扶助法}}
\newcommand{\abbrSGBV}{\abbrnoteonce{sgbv}{SGB V}{Fünftes Buch Sozialgesetzbuch}{社會法典第五編}}
\newcommand{\abbrStAG}{\abbrnoteonce{staeg}{StÄG}{Strafrechtsänderungsgesetz}{刑法修正法}}
\newcommand{\abbrStenBer}{\abbrnoteonce{stenber}{StenBer.}{Stenographischer Bericht}{速記紀錄}}
\newcommand{\abbrStGB}{\abbrnoteonce{stgb}{StGB}{Strafgesetzbuch}{刑法}}
\newcommand{\abbrStREG}{\abbrnoteonce{streg}{StREG}{Strafrechtsreform-Ergänzungsgesetz}{刑法改革補充法}}
\newcommand{\abbrStrRG}{\abbrnoteonce{strrg}{StrRG}{Strafrechtsreformgesetz}{刑法改革法}}
\newcommand{\abbrUS}{\abbrnoteonce{us}{U.S.}{United States Reports}{美國最高法院判決彙編}}
\newcommand{\abbrVgl}{\abbrnoteonce{vgl}{vgl.}{vergleiche}{參見}}
\newcommand{\abbrVVDStRL}{\abbrnoteonce{vvdstrl}{VVDStRL}{Veröffentlichungen der Vereinigung der Deutschen Staatsrechtslehrer}{德國公法教師協會論文集}}
\newcommand{\abbrWP}{\abbrnoteonce{wp}{WP.}{Wahlperiode}{屆期}}
\newcommand{\abbrWp}{\abbrnoteonce{wp}{Wp.}{Wahlperiode}{屆期}}
\newcommand{\abbrZStrW}{\abbrnoteonce{zstrw}{ZStrW}{Zeitschrift für die gesamte Strafrechtswissenschaft}{整體刑法學期刊}}
\newcommand{\recordsourcerandstart}{%
  \ifrandnumbers
    \setcounter{lastsourcerandstart}{\value{randnumber}}%
    \global\lastsourcehasrandnumberstrue
  \else
    \global\lastsourcehasrandnumbersfalse
  \fi
}
\newlength{\biparagraphskip}
\setlength{\biparagraphskip}{0.52em}
\newlength{\lawboxblockbeforeskip}
\setlength{\lawboxblockbeforeskip}{0.72em}
\newlength{\lawboxblockafterskip}
\setlength{\lawboxblockafterskip}{0.92em}
\newif\iflawboxpartendedblock
\lawboxpartendedblockfalse
\newcommand{\sourceparstyle}{\footnotesize\color{sourceink}\setstretch{1.12}\RaggedRight\emergencystretch=3em\setlength{\parskip}{0.42em}\obeylines\selectfont}
\newcommand{\transparstyle}{\songfont\color{caseink}\setstretch{1.12}\setlength{\parindent}{0pt}\setlength{\parskip}{0.42em}}
\newcommand{\bilingualpairskip}{%
  \par\addvspace{\biparagraphskip}%
  \switchcolumn
  \par\addvspace{\biparagraphskip}%
  \switchcolumn*%
  \global\intranslationfalse
}
\newcommand{\bilingualboxpairskip}{%
  \par
  \switchcolumn
  \par
  \switchcolumn*%
  \global\intranslationfalse
}
\newcommand{\bilinguallawboxblockskip}{%
  \switchcolumn*%
  \par\addvspace{\lawboxblockafterskip}%
  \switchcolumn
  \par\addvspace{\lawboxblockafterskip}%
  \switchcolumn*%
  \global\intranslationfalse
}
\newcommand{\bikeepwithnext}[1][4]{%
  \ifbilingualcolumns
    \syncleft
    \Needspace{#1\baselineskip}%
    \switchcolumn
    \Needspace{#1\baselineskip}%
    \switchcolumn*%
    \global\intranslationfalse
  \else
    \Needspace{#1\baselineskip}%
  \fi
}
\NewEnviron{sourcepara}[1][]{
  \recordsourcerandstart
  \ifshoworiginal
    \ifbilingualcolumns
      \ifintranslation\switchcolumn*\global\intranslationfalse\fi
    \else
      \Needspace{5\baselineskip}
    \fi
    \ifbilingualcolumns\else\par\addvspace{0.35em}\fi
    {\sourceparstyle
    \noindent\BODY\par}
    \ifbilingualcolumns\else\addvspace{0.25em}\fi
    \ifbilingualcolumns
      \switchcolumn
      \global\intranslationtrue
    \fi
  \else
    \ifrandnumbers
      \begingroup
        \setbox0=\vbox{\hsize=\linewidth\sourceparstyle\noindent\BODY\par}%
      \endgroup
    \fi
  \fi
}
\NewEnviron{transpara}{
  \ifshowtranslation
    \setcounter{randnumberlocal}{\value{lastsourcerandstart}}%
    \ifbilingualcolumns
      \ifintranslation\else\switchcolumn\global\intranslationtrue\fi
      {\transparstyle\setlength{\leftskip}{0.55em}\setlength{\rightskip}{0.15em plus 1em}\addtolength{\linewidth}{-0.55em}\BODY\par}
      \switchcolumn*
      \bilingualpairskip
    \else
      {\transparstyle\BODY\par}
    \fi
  \else
    \ifbilingualcolumns
      \ifintranslation\switchcolumn*\global\intranslationfalse\fi
    \fi
  \fi
}
\NewEnviron{sourceboxpara}[1][]{
  \global\lawboxpartendedblockfalse
  \recordsourcerandstart
  \ifshoworiginal
    \ifbilingualcolumns
      \ifintranslation\switchcolumn*\global\intranslationfalse\fi
    \else
      \Needspace{3\baselineskip}
    \fi
    {\sourceparstyle
    \BODY\par}
    \ifbilingualcolumns\else
      \iflawboxpartendedblock\addvspace{\lawboxblockafterskip}\fi
    \fi
    \ifbilingualcolumns
      \switchcolumn
      \global\intranslationtrue
    \fi
  \fi
}
\NewEnviron{transboxpara}{
  \global\lawboxpartendedblockfalse
  \ifshowtranslation
    \setcounter{randnumberlocal}{\value{lastsourcerandstart}}%
    \ifbilingualcolumns
      \ifintranslation\else\switchcolumn\global\intranslationtrue\fi
      {\transparstyle\BODY\par}
      \iflawboxpartendedblock
        \bilinguallawboxblockskip
      \else
        \switchcolumn*
        \bilingualboxpairskip
      \fi
    \else
      {\transparstyle\BODY\par}
      \iflawboxpartendedblock\addvspace{\lawboxblockafterskip}\fi
    \fi
  \else
    \ifbilingualcolumns
      \ifintranslation\switchcolumn*\global\intranslationfalse\fi
    \fi
  \fi
}
\newcommand{\leitsatznum}[1]{\noindent\hangindent=3.0em\hangafter=1\makebox[2.25em][r]{\bfseries #1.}\hspace{0.55em}\ignorespaces}
\newcommand{\syncleft}{\ifbilingualcolumns\ifintranslation\switchcolumn*\global\intranslationfalse\fi\fi}
\pretocmd{\section}{\syncleft\clearpage}{}{}
\pretocmd{\subsection}{\syncleft}{}{}
\pretocmd{\subsubsection}{\syncleft}{}{}
\newcommand{\biheadingfont}{\songfont}
\newcommand{\bitocentry}[2]{%
  \ifshoworiginal
    \ifshowtranslation
      \ifstrequal{#1}{#2}{#1}{#1\space #2}%
    \else
      #1%
    \fi
  \else
    #2%
  \fi
}
\pdfstringdefDisableCommands{%
  \def\bitocentry#1#2{#1}%
}
\newcommand{\biheading}[7]{%
  \syncleft#1\bikeepwithnext[#3]\phantomsection\addcontentsline{toc}{#2}{\bitocentry{#6}{#7}}%
  \ifbilingualcolumns
    {#4 #6\par}%
    \switchcolumn
    {#4\biheadingfont #7\par}%
    \switchcolumn*%
    \global\intranslationfalse
    \par\addvspace{#5}%
    \switchcolumn
    \par\addvspace{#5}%
    \switchcolumn*%
  \else
    \ifshoworiginal{#4 #6\par}\fi
    \ifshowtranslation{#4\biheadingfont #7\par}\fi
    \addvspace{#5}%
  \fi
}
\newcommand{\termsectionheading}[1]{%
  \par\Needspace{9\baselineskip}%
  \vspace{0.65em}%
  {\songfont\bfseries #1\par}%
  \vspace{0.2em}%
}
% 章節換頁：
%   雙語 paracol 欄內不要用 \clearpage；它會在同步兩欄時吐出只有欄線/頁碼的空頁。
%   \newpage 足以讓下一個 section 起新頁，又不會強制清空所有待排材料。
%   單語模式不在 paracol 內，仍可用 \clearpage。
\newcommand{\bisectionpagebreak}{\ifbilingualcolumns\newpage\else\clearpage\fi}
\newcommand{\bisection}[2]{%
  \biheading{\iffirstbilingualsection\global\firstbilingualsectionfalse\else\bisectionpagebreak\fi}{section}{6}{\centering\Large\bfseries}{0.8em}{#1}{#2}%
}
\newcommand{\bisubsection}[2]{%
  \biheading{}{subsection}{5}{\large\bfseries}{0.55em}{#1}{#2}%
}
\newcommand{\bisubsubsection}[2]{%
  \biheading{}{subsubsection}{4}{\normalsize\bfseries}{0.45em}{#1}{#2}%
}
\newcommand{\bicompositeheading}[4]{%
  \biheading{}{subsection}{5}{\large\bfseries}{0.16em}{#1}{#2}%
  \biheading{}{subsubsection}{3}{\normalsize\bfseries}{0.45em}{#3}{#4}%
}
\newif\iflawboxfirstitem
\lawboxfirstitemfalse
\newcommand{\lawboxprespace}[1]{%
  \par
  \iflawboxfirstitem
    \global\lawboxfirstitemfalse
  \else
    \addvspace{#1}%
  \fi
}
\newtcolorbox{lawboxinner}{
  caseboxbase,
  colframe=sourceline!85!black,
  colback=white,
  left=1.05em,
  right=1.05em,
  top=0.62em,
  bottom=0.62em,
  before upper={\global\inlawboxtrue\global\lawboxfirstitemtrue\global\lawboxhaspendingrnfalse\ifintranslation\lawboxtransfont\else\lawboxsourcefont\fi\RaggedRight\setlength{\parindent}{0pt}\setlength{\parskip}{0pt}\ifintranslation\setstretch{\lawboxtransstretch}\else\setstretch{\lawboxsourcestretch}\fi},
  after upper={\global\lawboxhaspendingrnfalse\global\inlawboxfalse}
}
\tcbset{
  lawboxpartbase/.style={
    enhanced,
    breakable=false,
    width=0.985\linewidth,
    colback=white,
    frame hidden,
    boxrule=0pt,
    arc=0pt,
    left=1.05em,
    right=1.05em,
    boxsep=1.5mm,
    before skip=0pt,
    after skip=0pt,
    before upper={\global\inlawboxtrue\global\lawboxfirstitemtrue\global\lawboxhaspendingrnfalse\ifintranslation\lawboxtransfont\else\lawboxsourcefont\fi\RaggedRight\setlength{\parindent}{0pt}\setlength{\parskip}{0pt}\ifintranslation\setstretch{\lawboxtransstretch}\else\setstretch{\lawboxsourcestretch}\fi},
    after upper={\global\lawboxhaspendingrnfalse\global\inlawboxfalse},
    borderline west={0.28pt}{0pt}{sourceline!85!black},
    borderline east={0.28pt}{0pt}{sourceline!85!black},
  },
  lawboxpartfirst/.style={
    lawboxpartbase,
    top=0.62em,
    bottom=0.04em,
    before skip=\lawboxblockbeforeskip,
    borderline north={0.28pt}{0pt}{sourceline!85!black},
  },
  lawboxpartmiddle/.style={
    lawboxpartbase,
    top=0.04em,
    bottom=0.04em,
  },
  lawboxpartlast/.style={
    lawboxpartbase,
    top=0.04em,
    bottom=0.62em,
    after skip=0pt,
    borderline south={0.28pt}{0pt}{sourceline!85!black},
  }
}
\newtcolorbox{lawboxpartinner}[1]{#1}
\newtcolorbox{termboxinner}{
  caseboxbase,
  colframe=noteblue!70!black,
  colback=casepaper,
  before upper={\songfont\setlength{\parindent}{0pt}\setlength{\parskip}{0.35em}}
}
\newtcolorbox{deboxinner}{
  caseboxbase,
  colframe=notegray!72!black,
  colback=white,
  left=1.05em,
  right=1.05em,
  top=0.62em,
  bottom=0.62em,
  before upper={\global\inlawboxtrue\global\lawboxfirstitemtrue\global\lawboxhaspendingrnfalse\small\setlength{\parindent}{0pt}\setlength{\parskip}{0.35em}},
  after upper={\global\lawboxhaspendingrnfalse\global\inlawboxfalse}
}
\NewEnviron{lawbox}{
  \begin{lawboxinner}\BODY\end{lawboxinner}
}
\NewEnviron{lawboxpart}[2]{
  \begin{lawboxpartinner}{lawboxpart#1,equal height group=#2}\BODY\end{lawboxpartinner}%
  \ifstrequal{#1}{last}{\global\lawboxpartendedblocktrue}{}%
}
\NewEnviron{termbox}{
  \par\addvspace{0.55em}
  {\color{noteblue!70!black}\rule{\linewidth}{0.28pt}}\par
  \begingroup
    \songfont\small\color{caseink}
    \setlength{\parindent}{0pt}
    \setlength{\parskip}{0.24em}
    \BODY
    \par
  \endgroup
  {\color{noteblue!70!black}\rule{\linewidth}{0.28pt}}\par
  \addvspace{0.55em}
}
\NewEnviron{debox}{
  \begin{deboxinner}\BODY\end{deboxinner}
}
\providecommand{\paragraphmark}[1]{}
\newcommand{\lawboxtitleafter}{\addvspace{0.06em}}
\newcommand{\lawtitle}[1]{\lawboxprespace{0.22em}{\lawboxtransfont\bfseries\lawboxuserandnumber #1}\par\lawboxtitleafter}
\newcommand{\absno}[2]{\lawboxprespace{0.04em}\noindent\lawboxuserandnumber\hangindent=2.6em\hangafter=1\makebox[2.6em][l]{(#1)}#2\par}
\newcommand{\nrno}[2]{\lawboxprespace{0pt}\noindent\lawboxuserandnumber\hspace*{2.6em}\hangindent=4.4em\hangafter=1\makebox[1.8em][l]{#1.}#2\par}
\newcommand{\letterno}[2]{\lawboxprespace{0pt}\noindent\lawboxuserandnumber\hspace*{4.4em}\hangindent=6.0em\hangafter=1\makebox[1.6em][l]{#1)}#2\par}
\newcommand{\sourcelawtitle}[1]{\lawboxprespace{0.22em}{\lawboxsourcefont\bfseries\lawboxuserandnumber #1}\par\lawboxtitleafter}
\newcommand{\sourcelawsubtitle}[1]{\lawboxprespace{0.04em}{\lawboxsourcefont\bfseries\lawboxuserandnumber #1}\par\lawboxtitleafter}
\newcommand{\sourceabsno}[2]{\lawboxprespace{0.04em}\noindent\lawboxuserandnumber\hangindent=2.45em\hangafter=1\makebox[2.45em][l]{(#1)}#2\par}
\newcommand{\sourcenrno}[2]{\lawboxprespace{0pt}\noindent\lawboxuserandnumber\hspace*{2.45em}\hangindent=4.25em\hangafter=1\makebox[1.8em][l]{#1.}#2\par}
\newcommand{\sourceletterno}[2]{\lawboxprespace{0pt}\noindent\lawboxuserandnumber\hspace*{4.25em}\hangindent=5.85em\hangafter=1\makebox[1.6em][l]{#1)}#2\par}
\newcommand{\sourcequotepara}[1]{\lawboxprespace{0.04em}\noindent\lawboxuserandnumber\hangindent=1.2em\hangafter=1 #1\par}
\newcommand{\sourcelawline}[1]{\lawboxprespace{0.04em}\noindent\lawboxuserandnumber #1\par}
\newcommand{\lawline}[1]{\lawboxprespace{0.04em}\noindent\lawboxuserandnumber #1\par}
\newcommand{\pendingtranslation}{\par{\songfont\footnotesize\color{pendingink}（中文譯文待補。）}\par}
\newcommand{\tocpart}[1]{\par\addvspace{0.52em}{\footnotesize\bfseries\color{pendingink}#1\par}\addvspace{0.16em}}
\newcommand{\tocdashline}{\par\addvspace{0.48em}\noindent{\color{notegray}\xleaders\hbox to 0.82em{\hfil\rule{0.42em}{0.28pt}\hfil}\hskip\linewidth}\par\addvspace{0.38em}}
% \tocpanel{font-cmd}{content}：將目錄置中於所在欄寬（雙語欄適用）。
% A3 橫式使用 0.58\linewidth，A4 橫式使用 0.84\linewidth，
% 使兩種紙張下目錄欄內容實際寬度相近（均約 100–105 mm）。
\newcommand{\tocpanel}[2]{%
  \makebox[\linewidth][c]{%
    \ifx\bilingualpaper\bilingualpaperaThree
      \begin{minipage}{0.58\linewidth}%
    \else
      \begin{minipage}{0.84\linewidth}%
    \fi
    \toccolstyle
    #1#2%
    \end{minipage}%
  }%
}
% ===== 首頁簡目錄的兩層 description list 縮排 =====
% enumitem 的 description list 可把每一列拆成：
%   [label]  labelsep  body text
% 這裡的「起點」都以所在 minipage/欄位的左緣為 0。
%
% 核心規則：
%   labelindent = label 欄位從左緣開始的位置。
%   labelwidth  = label 欄位本身寬度；標籤太長（如 Abw. Meinung）就加大它。
%   labelsep    = label 欄位右緣到正文的距離。
%   leftmargin  = 正文 body text 的起點。判斷層級視覺時主要看這個。
%
% 所以：
%   想讓「Begründung / 判決理由」整列正文往右或往左：改 tocitems 的 leftmargin。
%   想讓「A. Verfahren und Gegenstand / A. 程序與審查標的」正文往右或往左：
%     改 tocsubitems 的 leftmargin。
%   想讓 A./B./C. 這些標籤本身往右或往左，但正文不動：改 tocsubitems 的 labelindent。
%   想讓 A./B./C. 與正文間距變大或變小：改 tocsubitems 的 labelsep。
%
% 層級原則：
%   tocsubitems.leftmargin 必須大於 tocitems.leftmargin，
%   否則子層級正文會看起來比「Gründe / 理由」還靠左。
\newlist{tocitems}{description}{1}
\setlist[tocitems]{%
  labelindent=0pt,    % 第一層 label 欄位起點；0pt 表示貼齊目錄欄左緣。
  leftmargin=2.54cm,  % 第一層正文起點；例如 Begründung / 判決理由 的左緣。
  labelwidth=2.16cm,  % 第一層 label 欄位寬度；需容納 Leitsätze、Abw. Meinung、不同意見等。
  labelsep=0.32cm,    % 第一層 label 與正文的水平間距；只改間距，不改正文起點。
  itemsep=0.28em,     % 第一層各項之間的垂直距離。
  topsep=0.20em,      % 第一層 list 上下與前後內容的垂直距離。
  parsep=0pt,         % 同一 item 內段落之間的距離；目錄項通常不需要。
  partopsep=0pt,      % list 若緊接段落開始時的額外距離；關閉以免目錄高度飄動。
  align=left,         % label 欄內左對齊；長短標籤不靠右貼齊。
  font=\normalfont\bfseries % label 的字型；正文不受這行影響。
}
\newlist{tocsubitems}{description}{1}
\setlist[tocsubitems]{%
  labelindent=1cm,    % 第二層 label 起點；只控制 A./B./C. 本身的位置。
  leftmargin=3.00cm,  % 第二層正文起點；必須大於 2.54cm，才會比 Begründung / 判決理由更右。
  labelwidth=0.5cm,   % 第二層 label 欄寬；A./B./C. 很短，可小於第一層。
  labelsep=1.5cm,    % 第二層 label 與正文間距；A. 與 Verfahren / 程序 的距離看這裡。
  itemsep=0.12em,     % 第二層各項較密，避免 A.–E. 佔太高。
  topsep=0.04em,      % 第二層 list 與 Gründe / 理由之間的垂直距離。
  parsep=0pt,         % 同一第二層 item 內段落距離；目錄項通常不需要。
  partopsep=0pt,      % 關閉額外段落起始距離，避免版面不穩。
  align=left,         % A./B./C. label 欄內左對齊。
  font=\normalfont\bfseries, % A./B./C. label 加粗；正文不受這行影響。
  before={}           % 不在第二層 list 前自動插入額外內容。
}
\newlist{termitems}{description}{1}
\ifbilingualcolumns
  \setlist[termitems]{%
    leftmargin=5.2cm,
    labelwidth=4.8cm,
    labelsep=0.35cm,
    itemsep=0.16em,
    topsep=0.2em,
    parsep=0pt,
    partopsep=0pt,
    font=\normalfont\bfseries,
    style=nextline
  }
\else
  \setlist[termitems]{%
    leftmargin=0pt,
    labelwidth=0pt,
    labelsep=0pt,
    itemsep=0.24em,
    topsep=0.2em,
    parsep=0pt,
    partopsep=0pt,
    font=\normalfont\bfseries,
    style=nextline
  }
\fi
\newlist{abbritems}{description}{1}
\setlist[abbritems]{%
  leftmargin=2.38cm,
  labelwidth=2.02cm,
  labelsep=0.28cm,
  itemsep=0.16em,
  topsep=0.2em,
  parsep=0pt,
  partopsep=0pt,
  align=left,
  font=\normalfont\bfseries
}
\ifbilingualcolumns
  \newenvironment{termcolumns}{\begin{multicols}{2}}{\end{multicols}}
\else
  \newenvironment{termcolumns}{}{}
\fi

% ===== 編者註腳開關 =====
% 一般版預設不顯示編者註，避免正文腳註過密。
% 若開啟 \classroommodeon，GG/條文編者註仍會自動顯示，無需另開本開關。
% 若一般版也要顯示編者註，可在單案檔 \input 後加 \showeditornotestrue。
\newif\ifshoweditornotes
\showeditornotesfalse

% ===== 目錄排版輔助 =====
\newcommand{\toccolstyle}{\raggedright\arraybackslash}
\newcommand{\tocsingleindent}{\leftskip=1.45cm\rightskip=1.2cm}

% ===== 行內引文環境（德文原文與中文譯文各一，帶左側灰色邊線）=====
\newcommand{\quoteparbreak}{\par\addvspace{0.48em}}
\newcommand{\sourcequoteinner}[1]{%
  \begin{tcolorbox}[
    enhanced, breakable, width=\dimexpr\linewidth-0.8em\relax, frame hidden,
    colback=white, coltext=caseink, boxrule=0pt,
    borderline west={0.55pt}{0.88em}{notegray},
    left=1.78em, right=0.72em, top=0.18em, bottom=0.18em,
    boxsep=0pt, before skip=0.24em, after skip=0.24em,
    before upper={\normalfont\footnotesize\color{caseink}\setlength{\parindent}{0pt}\setlength{\parskip}{0.34em}\raggedright\setlength{\rightskip}{0pt plus 1em}}
  ]%
  #1\par
  \end{tcolorbox}%
}
\newcommand{\transquoteinner}[1]{%
  \begin{tcolorbox}[
    enhanced, breakable, width=\dimexpr\linewidth-0.8em\relax, frame hidden,
    colback=white, coltext=caseink, boxrule=0pt,
    borderline west={0.55pt}{0.88em}{notegray},
    left=1.78em, right=0.72em, top=0.18em, bottom=0.18em,
    boxsep=0pt, before skip=0.24em, after skip=0.24em,
    before upper={\songfont\small\color{caseink}\setlength{\parindent}{0pt}\setlength{\parskip}{0.34em}\raggedright\setlength{\rightskip}{0pt plus 1em}}
  ]%
  #1\par
  \end{tcolorbox}%
}

% ===== 大綱（Gliederung）Rn 輔助命令 =====
\newcommand{\outrn}[2]{\enspace{\normalfont\footnotesize\color{pendingink}(Rn\,#1–#2)}}
\newcommand{\outrnsingle}[1]{\enspace{\normalfont\footnotesize\color{pendingink}(Rn\,#1)}}
\newcommand{\outrnzh}[2]{\enspace{\normalfont\footnotesize\color{pendingink}（第\,#1–#2\,段）}}
\newcommand{\outrnzhs}[1]{\enspace{\normalfont\footnotesize\color{pendingink}（第\,#1\,段）}}
% ===== 大綱雙語對照表格輔助命令（三欄：段碼 │ 德文 │ 中文）=====
% 欄 1：段碼（由欄定義自動格式化：小字、等寬、右對齊、pendingink）
% 欄 2：德文標籤＋正文
% 欄 3：中文標籤＋正文
%
% 無段碼的列（\omainrow、\osecrow），欄 1 以 & 空置。
% 有段碼的列（\osecrowrn、\osubrow、\ointrow），#1 為預格式化字串
%   例：\osubrow{2–49}{I.}{...}{I.}{...}
%       \ointrow{107}{...}{...}
%
% \opartrow：段組標題（橫跨三欄）
\newcommand{\opartrow}[2]{%
  \noalign{\vspace{0.30em}}%
  \multicolumn{3}{@{}l@{}}{%
    \footnotesize\bfseries\color{pendingink}#1\hfill\songfont#2%
  }\\[0.06em]%
}
% \odashrow：虛線分隔（橫跨三欄）
\newcommand{\odashrow}{%
  \noalign{\vspace{0.28em}}%
  \multicolumn{3}{@{}l@{}}{\color{notegray}%
    \xleaders\hbox to 0.82em{\hfil\rule{0.42em}{0.28pt}\hfil}\hskip 0.88\linewidth}%
  \\[0.24em]%
}
% \odissentpart：不同意見書區段標頭。比一般 \opartrow 更像附錄性轉場。
\newcommand{\odissentpart}[2]{%
  \noalign{\vspace{0.42em}}%
  \multicolumn{3}{@{}p{0.88\textwidth}@{}}{%
    \color{notegray}\rule{\linewidth}{0.18pt}%
  }\\[-0.18em]%
  \multicolumn{3}{@{}p{0.88\textwidth}@{}}{%
    \makebox[\linewidth][s]{%
      \footnotesize\bfseries\color{pendingink}#1%
      \hfill{\songfont #2}%
    }%
  }\\[-0.02em]%
}
% \odissentopinion：不同意見書的署名與一行摘要，右欄保留段碼範圍。
\newcommand{\odissentopinion}[5]{%
  \footnotesize\hangindent=0.52cm\hangafter=1%
    {\bfseries\color{caseink}#2}\enspace{\color{notegray}\textperiodcentered}\enspace\color{caseink}#3%
  &\footnotesize\hangindent=0.52cm\hangafter=1%
    {\songfont\bfseries\color{caseink}#4}\enspace{\color{notegray}\textperiodcentered}\enspace\songfont\color{caseink}#5%
  &#1%
  \\[0.12em]%
}
% \omainrow：頂層項目，無段碼（Leitsätze、Urteil、Gründe …）
\newcommand{\omainrow}[4]{%
  \footnotesize\hangindent=1.92cm\hangafter=1%
    \makebox[1.92cm][l]{\bfseries\color{caseink}#1}\color{caseink}#2%
  &\footnotesize\hangindent=1.92cm\hangafter=1%
    \makebox[1.92cm][l]{\bfseries\color{caseink}#3}\songfont\color{caseink}#4%
  &%
  \\[0.15em]%
}
% \omainrowrn：頂層項目，有段碼（不同意見等標籤寬於 \osecrow 框者）
\newcommand{\omainrowrn}[5]{%
  \footnotesize\hangindent=1.92cm\hangafter=1%
    \makebox[1.92cm][l]{\bfseries\color{caseink}#2}\color{caseink}#3%
  &\footnotesize\hangindent=1.92cm\hangafter=1%
    \makebox[1.92cm][l]{\bfseries\color{caseink}#4}\songfont\color{caseink}#5%
  &#1%
  \\[0.15em]%
}
% \osecrow：一級子項，無段碼（A.、C.、D. 等有子項者）
\newcommand{\osecrow}[4]{%
  \footnotesize\hspace{1.10cm}\hangindent=1.98cm\hangafter=1%
    \makebox[0.86cm][l]{\bfseries\color{caseink}#1}\color{caseink}#2%
  &\footnotesize\hspace{1.10cm}\hangindent=1.98cm\hangafter=1%
    \makebox[0.86cm][l]{\bfseries\color{caseink}#3}\songfont\color{caseink}#4%
  &%
  \\[0.11em]%
}
% \osecrowrn：一級子項，有段碼（B.、E. 等完整段落範圍者）
% #1=段碼字串（如 101–106）, #2=DE標籤, #3=DE正文, #4=ZH標籤, #5=ZH正文
\newcommand{\osecrowrn}[5]{%
  \footnotesize\hspace{1.10cm}\hangindent=1.98cm\hangafter=1%
    \makebox[0.86cm][l]{\bfseries\color{caseink}#2}\color{caseink}#3%
  &\footnotesize\hspace{1.10cm}\hangindent=1.98cm\hangafter=1%
    \makebox[0.86cm][l]{\bfseries\color{caseink}#4}\songfont\color{caseink}#5%
  &#1%
  \\[0.11em]%
}
% \osubrow：二級子項，有段碼範圍（I.、II.、A.I. …）
% #1=段碼字串, #2=DE標籤, #3=DE正文, #4=ZH標籤, #5=ZH正文
\newcommand{\osubrow}[5]{%
  \footnotesize\hspace{1.40cm}\hangindent=2.30cm\hangafter=1%
    \makebox[0.90cm][l]{\bfseries\color{caseink}#2}\color{caseink}#3%
  &\footnotesize\hspace{1.40cm}\hangindent=2.30cm\hangafter=1%
    \makebox[0.90cm][l]{\bfseries\color{caseink}#4}\songfont\color{caseink}#5%
  &#1%
  \\[0.09em]%
}
% \ointrow：引入段落，有單一段碼，無標籤
% #1=段碼字串, #2=DE正文, #3=ZH正文
\newcommand{\ointrow}[3]{%
  \footnotesize\hspace{0.52cm}\hangindent=0.52cm\hangafter=1\color{caseink}#2%
  &\footnotesize\hspace{0.52cm}\hangindent=0.52cm\hangafter=1\songfont\color{caseink}#3%
  &#1%
  \\[0.09em]%
}
% \osubsubrow：三級子項（1.、2.、3.）
% #1=段碼字串, #2=DE標籤, #3=DE正文, #4=ZH標籤, #5=ZH正文
\newcommand{\osubsubrow}[5]{%
  \footnotesize\hspace{1.80cm}\hangindent=2.48cm\hangafter=1%
    \makebox[0.68cm][l]{\bfseries\color{caseink}#2}\color{caseink}#3%
  &\footnotesize\hspace{1.80cm}\hangindent=2.48cm\hangafter=1%
    \makebox[0.68cm][l]{\bfseries\color{caseink}#4}\songfont\color{caseink}#5%
  &#1%
  \\[0.08em]%
}
% \osubsubsubrow：四級子項（a）、b））
% #1=段碼字串, #2=DE標籤, #3=DE正文, #4=ZH標籤, #5=ZH正文
\newcommand{\osubsubsubrow}[5]{%
  \footnotesize\hspace{2.10cm}\hangindent=2.68cm\hangafter=1%
    \makebox[0.58cm][l]{\bfseries\color{caseink}#2}\color{caseink}#3%
  &\footnotesize\hspace{2.10cm}\hangindent=2.68cm\hangafter=1%
    \makebox[0.58cm][l]{\bfseries\color{caseink}#4}\songfont\color{caseink}#5%
  &#1%
  \\[0.07em]%
}
% \osubsubsubsubrow：五級子項（aa）、bb））
% 標籤框 0.62cm：足以容納「aa)」在 footnotesize bold 下（約 0.50cm）並留有間距
% #1=段碼字串, #2=DE標籤, #3=DE正文, #4=ZH標籤, #5=ZH正文
\newcommand{\osubsubsubsubrow}[5]{%
  \footnotesize\hspace{2.38cm}\hangindent=3.06cm\hangafter=1%
    \makebox[0.68cm][l]{\bfseries\color{caseink}#2}\color{caseink}#3%
  &\footnotesize\hspace{2.38cm}\hangindent=3.06cm\hangafter=1%
    \makebox[0.68cm][l]{\bfseries\color{caseink}#4}\songfont\color{caseink}#5%
  &#1%
  \\[0.06em]%
}
% \outlinepage：雙語對照大綱頁包裝環境（longtable 自動跨頁）
% 三欄：德文欄 + 中文欄 + 段碼欄（最右，不折行）
% 段碼欄寬度：1.80cm，足以容納最長段碼如「309–348」
% DE/ZH 欄寬：(0.87\textwidth - 2.08cm) / 2
%   驗算：2×欄 + 0.05\textwidth + 0.28cm gap + 1.80cm = 0.87tw - 2.08cm + 0.05tw + 2.08cm = 0.92tw ✓
\newenvironment{outlinepage}{%
  \vspace*{0.40cm}%
  \setlength{\LTleft}{0.04\textwidth}%
  \setlength{\LTright}{0.04\textwidth}%
  \setlength{\tabcolsep}{0pt}%
  % 大綱列距由 longtable 的 \arraystretch 與各 row macro 結尾的 \\[...em] 共同決定。
  % 這裡控制「每一列本身的表格 strut 高度」；數值越大，A./I./1. 各項之間越像有空行。
  % A3 原本用 0.62，視覺上沒有空行；A4 也沿用同值，避免切到 A4 後項目被撐開。
  % 若日後覺得大綱太擠，只微調這個值（例如 0.68），不要改各 \omainrow/\osubrow 的縮排。
  \renewcommand{\arraystretch}{0.62}%
  \begin{longtable}{%
    >{\vspace{0pt}\raggedright\arraybackslash}p{\dimexpr(0.87\textwidth-2.08cm)/2\relax}%
    @{\hspace{0.025\textwidth}}%
    !{\color{notegray}\vrule width 0.12pt}%
    @{\hspace{0.025\textwidth}}%
    >{\vspace{0pt}\raggedright\arraybackslash}p{\dimexpr(0.87\textwidth-2.08cm)/2\relax}%
    @{\hspace{0.28cm}}%
    >{\vspace{0pt}\footnotesize\normalfont\color{pendingink}\raggedright\arraybackslash}p{1.80cm}%
    @{}%
  }%
  % 首頁欄標籤：不要橫跨置中；分別放在德文欄與中文欄的實際欄位起點。
  \footnotesize\bfseries\color{sourceink}Gliederung%
  &\footnotesize\bfseries\songfont\color{sourceink}大綱%
  &%
  \\[0.36em]%
  \endfirsthead%
  % 續頁欄標籤：同樣分欄定位，以灰色細體區分；無需標注「續」。
  \footnotesize\color{pendingink}Gliederung%
  &\footnotesize\songfont\color{pendingink}大綱%
  &%
  \\[0.24em]%
  \endhead%
}{%
  \end{longtable}%
}
 之前設定；
%   預設 chinese / a4）
\ifdefined\docmode\else\def\docmode{chinese}\fi
\ifdefined\bilingualpaper\else\def\bilingualpaper{a4}\fi
\newif\ifshoworiginal
\newif\ifshowtranslation
\newif\ifbilinguallandscape
\newif\ifbilingualcolumns
\newif\ifintranslation
\newif\iffirstbilingualsection

\showoriginalfalse
\showtranslationfalse
\bilinguallandscapefalse
\bilingualcolumnsfalse
\intranslationfalse
\firstbilingualsectionfalse

\def\modechinese{chinese}
\def\modegerman{german}
\def\modebilingual{bilingual}
\def\bilingualpaperaThree{a3}
\def\bilingualpaperaFour{a4}

\ifx\docmode\modechinese
  \showtranslationtrue
\fi

\ifx\docmode\modegerman
  \showoriginaltrue
\fi

\ifx\docmode\modebilingual
  \showoriginaltrue
  \showtranslationtrue
  \bilinguallandscapetrue
  \bilingualcolumnstrue
\fi

% 編排模式說明：
% chinese  → \showoriginalfalse  \showtranslationtrue  （A4 直式，純中文）
% german   → \showoriginaltrue   \showtranslationfalse （A4 直式，純德文）
% bilingual→ \showoriginaltrue   \showtranslationtrue  \bilingualcolumnstrue（A3/A4 橫式對照）

\ifbilingualcolumns
  \ifx\bilingualpaper\bilingualpaperaFour
    \geometry{
      a4paper,
      landscape,
      left=1.25cm,
      right=2.25cm,
      top=1.25cm,
      bottom=1.75cm,
      footskip=8.5mm,
      marginparwidth=0.75cm,
      marginparsep=0.25cm
    }
  \else
    \geometry{
      a3paper,
      landscape,
      left=1.55cm,
      right=2.75cm,
      top=1.55cm,
      bottom=2.05cm,
      footskip=9.5mm,
      marginparwidth=0.9cm,
      marginparsep=0.35cm
    }
  \fi
\else
\geometry{
  a4paper,
  portrait,
  left=2.2cm,
  right=2.6cm,
  top=2.0cm,
  bottom=2.0cm,
  marginparwidth=0.8cm,
  marginparsep=0.3cm
}
\fi
% ===== 時間戳基礎設施 =====
\newcount\stamptimeval
\newcount\stampminuteval
\newcommand{\stamphh}{%
  \stamptimeval=\time
  \divide\stamptimeval by 60
  \ifnum\stamptimeval<10 0\fi
  \the\stamptimeval}
\newcommand{\stampmm}{%
  \stamptimeval=\time
  \stampminuteval=\time
  \divide\stampminuteval by 60
  \multiply\stampminuteval by 60
  \advance\stamptimeval by -\stampminuteval
  \ifnum\stamptimeval<10 0\fi
  \the\stamptimeval}
\newcommand{\stampwkde}[1]{%
  \ifcase#1Montag\or Dienstag\or Mittwoch\or Donnerstag\or Freitag\or Samstag\or Sonntag\fi}
\newcommand{\stampwkzh}[1]{%
  \ifcase#1 一\or 二\or 三\or 四\or 五\or 六\or 日\fi}
\newcommand{\stampmonthde}[1]{%
  \ifcase#1\or Januar\or Februar\or März\or April\or Mai\or Juni\or
  Juli\or August\or September\or Oktober\or November\or Dezember\fi}
\newcount\stampjdval
\newcount\stampwdval
\pgfcalendardatetojulian{\the\year-\the\month-\the\day}{\stampjdval}
\pgfcalendarjuliantoweekday{\the\stampjdval}{\stampwdval}
\newcommand{\timestampzh}{%
  \songfont 最後修改：\the\year 年\the\month 月\the\day 日%
  % （星期\stampwkzh{\the\stampwdval}）%
  % \space\stamphh:\stampmm
  }

\newcommand{\documentdatezh}{\the\year 年 \the\month 月 \the\day 日}
\newcommand{\documentdatede}{\the\day. \stampmonthde{\the\month} \the\year}
\newcommand{\bverfgetranslationnote}{%
  {\songfont 翻譯說明：初譯使用 Claude Code 與 GPT 協助迭代；仍請以德文原文為準。}%
}
\newcommand{\bverfgecourseinfo}{%
  {\songfont 課程：114 學年度第 2 學期；憲法專題研究（四）；授課老師：湯德宗教授；導讀日期：115 年 6 月 1 日。}%
}
\newcommand{\bverfgeeditorinfo}{%
  {\songfont 編者：王逸帆（R10A21126），國立台灣大學法律系研究所碩士班。}%
}
\newcommand{\bverfgetitleversionline}{%
  \ifbilingualcolumns
    Fassung \documentversion\enspace\textperiodcentered\enspace Stand: \documentdatede
    \quad
    {\songfont 版本 \documentversion\enspace\textperiodcentered\enspace 修訂日期：\documentdatezh\par}%
  \else
    \ifshowtranslation
      {\songfont 版本 \documentversion\enspace\textperiodcentered\enspace 修訂日期：\documentdatezh\par}%
    \else
      Fassung \documentversion\enspace\textperiodcentered\enspace Stand: \documentdatede\par
    \fi
  \fi
}
\newcommand{\bverfgetitlecornerinfo}{%
  \ifshowtranslation
    \vspace{0.72em}%
    \noindent\hfill
    \begin{minipage}{0.66\textwidth}
      \raggedleft\tiny\color{pendingink}\songfont
      \bverfgecourseinfo\par
      \bverfgeeditorinfo
    \end{minipage}%
  \fi
}
\newcommand{\bverfgetitlemeta}{%
  \vfill
  \begin{center}%
    \footnotesize\color{pendingink}%
    \bverfgetitleversionline
    \ifshowtranslation
      \vspace{0.28em}{\scriptsize\bverfgetranslationnote\par}%
    \fi
  \end{center}%
  \bverfgetitlecornerinfo
}

\pagestyle{fancy}
\fancyhf{}
\renewcommand{\headrulewidth}{0pt}
\renewcommand{\footrulewidth}{0pt}
\fancyfoot[C]{\thepage}
\ifbilingualcolumns
  \fancyfoot[R]{\scriptsize\color{pendingink}\timestampzh}%
\else
  \ifshoworiginal
  \else
    \fancyfoot[R]{\scriptsize\color{pendingink}\timestampzh\hspace{0.45cm}}%
  \fi
\fi
\setmainfont{Times New Roman}
\setsansfont{Helvetica Neue}
\setmonofont{Menlo}
\IfFileExists{/Users/iw/Library/Fonts/kaiu.ttf}{
  \setCJKmainfont[Path=/Users/iw/Library/Fonts/,AutoFakeBold=4]{kaiu.ttf}
}{
  \IfFontExistsTF{標楷體}{
    \setCJKmainfont[AutoFakeBold=4]{標楷體}
  }{
    \setCJKmainfont{Songti TC}
  }
}
\IfFontExistsTF{Songti TC}{
  \setCJKfamilyfont{song}{Songti TC}
  \setCJKmonofont{Songti TC}
}{
  \setCJKfamilyfont{song}[Path=/Users/iw/Library/Fonts/,AutoFakeBold=4]{kaiu.ttf}
  \setCJKmonofont[Path=/Users/iw/Library/Fonts/,AutoFakeBold=4]{kaiu.ttf}
}
\newcommand{\songfont}{\CJKfamily{song}}

% ===== 封面／目錄專用打字機字體（僅限封面、目錄頁使用，不影響正文）=====
\IfFontExistsTF{Radio Newsman}{%
  \newfontfamily\bverfgetitlede{Radio Newsman}%
}{%
  \newfontfamily\bverfgetitlede{Courier New}%
}
\IfFontExistsTF{Erikas Farbband}{%
  \newfontfamily\bverfgebodyde{Erikas Farbband}%
}{%
  \let\bverfgebodyde\bverfgetitlede
}
\IfFontExistsTF{Maoken Heavy Labourer}{%
  \setCJKfamilyfont{bverfgetitlecjk}{Maoken Heavy Labourer}%
}{%
  \setCJKfamilyfont{bverfgetitlecjk}{Songti TC}%
}
\IfFontExistsTF{Huiwen-mincho}{%
  \setCJKfamilyfont{bverfgebodycjk}{Huiwen-mincho}%
}{%
  \setCJKfamilyfont{bverfgebodycjk}{Songti TC}%
}
\newcommand{\bverfgetitlezh}{\bverfgetitlede\CJKfamily{bverfgetitlecjk}}
\newcommand{\bverfgebodyzh}{\bverfgebodyde\CJKfamily{bverfgebodycjk}}

\setstretch{1.35}
\setlist{itemsep=0.2em, topsep=0.4em}
\setlength{\parindent}{0pt}
\setlength{\parskip}{0.45em}
\raggedbottom
\setcounter{secnumdepth}{0}
\setcounter{tocdepth}{2}

\newcommand{\lawboxsourcefont}{\footnotesize}
\newcommand{\lawboxtransfont}{\footnotesize}
\newcommand{\lawboxsourcestretch}{1.02}
\newcommand{\lawboxtransstretch}{0.92}

\titleformat{\section}{\centering\Large\bfseries\songfont}{\thesection}{1em}{}
\titleformat{\subsection}{\large\bfseries\songfont}{\thesubsection}{1em}{}
\titleformat{\subsubsection}{\normalsize\bfseries\songfont}{\thesubsubsection}{1em}{}
\titleformat{\paragraph}{\normalsize\bfseries\songfont}{\theparagraph}{1em}{}
\titlespacing*{\section}{0pt}{1.8em}{1.0em}
\titlespacing*{\subsection}{0pt}{1.2em}{0.6em}
\titlespacing*{\subsubsection}{0pt}{1.0em}{0.45em}
\titlespacing*{\paragraph}{0pt}{0.6em}{0.4em}

\definecolor{noteblue}{HTML}{A9C9F5}
\definecolor{noteteal}{HTML}{AEE1D6}
\definecolor{notegray}{HTML}{D7DADF}
\definecolor{caseink}{HTML}{000000}
\definecolor{casepaper}{HTML}{FCFEFD}
\definecolor{sourcepaper}{HTML}{FAFBF8}
\definecolor{sourceline}{HTML}{D7DED8}
\definecolor{sourceink}{HTML}{000000}
\definecolor{pendingink}{HTML}{000000}

\tcbset{
  caseboxbase/.style={
    enhanced,
    breakable,
    width=0.985\linewidth,
    colback=casepaper,
    coltitle=caseink,
    fonttitle=\bfseries\songfont,
    boxrule=0.28pt,
    arc=1.2mm,
    left=2.4mm,
    right=2.4mm,
    top=1.5mm,
    bottom=1.5mm,
    boxsep=1.5mm,
    before skip=0.65em,
    after skip=0.65em,
    before upper={\setlength{\parindent}{0pt}\setlength{\parskip}{0.35em}},
  }
}

\newcommand{\bilingualstart}{}
\newcommand{\bilingualstop}{}
\ifbilingualcolumns
  \renewcommand{\bilingualstart}{\small\setlength{\columnsep}{1.25cm}\setlength{\columnseprule}{0.12pt}\columnratio{0.5,0.5}\multicolumnfootnotes\flushbottom\sloppy\interlinepenalty=80\clubpenalty=800\widowpenalty=800\displaywidowpenalty=800\begin{paracol}{2}\global\intranslationfalse\global\firstbilingualsectiontrue{\footnotesize\color{sourceink}Deutsch\par}\switchcolumn{\footnotesize\songfont\color{sourceink}中文譯文\par}\switchcolumn*}
  \renewcommand{\bilingualstop}{\ifintranslation\switchcolumn*\global\intranslationfalse\fi\end{paracol}\clearpage\raggedbottom}
\fi
\newif\ifrandnumbers
\randnumbersfalse
\newcounter{randnumber}
\newcounter{randnumberlocal}
\newcounter{lastsourcerandstart}
\newif\iflastsourcehasrandnumbers
\lastsourcehasrandnumbersfalse
\newif\ifinlawbox
\inlawboxfalse
\newif\iflawboxhaspendingrn
\lawboxhaspendingrnfalse
\newcommand{\lawboxpendingrn}{}
\newcommand{\startrandnumbers}{\global\randnumberstrue\setcounter{randnumber}{1}\global\lastsourcehasrandnumbersfalse}
\newcommand{\stoprandnumbers}{\global\randnumbersfalse\global\lastsourcehasrandnumbersfalse}
\newlength{\randnumberwidth}
\setlength{\randnumberwidth}{0.95cm}
\newlength{\randnumberrightoffset}
\setlength{\randnumberrightoffset}{0.85cm}
\newlength{\randnumberlawboxrightoffset}
\setlength{\randnumberlawboxrightoffset}{1.40cm}
\newcommand{\randnumberbox}[1]{%
  \leavevmode
  \makebox[0pt][l]{%
    \smash{%
      \tikz[remember picture,overlay,baseline=(rnmark.base)]{%
        \coordinate (rnmark) at (0,0);%
        \ifinlawbox
          \path let \p1=(current page.east), \p2=(rnmark.base) in
            node[
              anchor=base east,
              inner sep=0pt,
              outer sep=0pt,
              text width=\randnumberwidth,
              align=right,
              text=pendingink,
              font=\normalfont\mdseries\scriptsize\ttfamily
            ] at (\x1-\randnumberlawboxrightoffset,\y2) {{\normalfont\mdseries\scriptsize\ttfamily #1}};%
        \else
          \path let \p1=(current page.east), \p2=(rnmark.base) in
            node[
              anchor=base east,
              inner sep=0pt,
              outer sep=0pt,
              text width=\randnumberwidth,
              align=right,
              text=pendingink,
              font=\normalfont\mdseries\scriptsize\ttfamily
            ] at (\x1-\randnumberrightoffset,\y2) {{\normalfont\mdseries\scriptsize\ttfamily #1}};%
        \fi
      }%
    }%
  }%
  \ignorespaces
}
\newcommand{\lawboxstorerandnumber}[1]{%
  \gdef\lawboxpendingrn{#1}%
  \global\lawboxhaspendingrntrue
  \ignorespaces
}
\newcommand{\lawboxuserandnumber}{%
  \iflawboxhaspendingrn
    \randnumberbox{\lawboxpendingrn}%
    \global\lawboxhaspendingrnfalse
  \fi
}
\newcommand{\sourcern}[1]{%
  \ifrandnumbers
    \ifshoworiginal
      \ifshowtranslation\else
        \ifinlawbox\lawboxstorerandnumber{#1}\else\randnumberbox{#1}\fi
      \fi
    \fi
  \fi
}
\newcommand{\transrn}[1]{%
  \ifrandnumbers
    \ifshowtranslation
      \ifinlawbox\lawboxstorerandnumber{#1}\else\randnumberbox{#1}\fi
    \fi
  \fi
}
% ===== 課堂模式：德文原文縮寫註解 =====
% 日常切換只改各判決檔的一行：
%   \classroommodeon        開啟課堂版；註解掉這行即回到一般版。
% 模板預設：
%   1. 課堂模式關閉。
%   2. 課堂模式開啟時，縮寫註解包含「德文全稱＋中文譯名」。
%   3. 課堂模式開啟時，GG/條文編者註仍會自動顯示。
% 進階微調才需要在單案檔覆寫：
%   \classroomabbrzhoff     縮寫註解僅顯示德文全稱。
%   \showeditornotestrue    一般版也顯示編者註。
% 註解策略：同一實際 PDF 頁面中，同一縮寫只在第一次出現處加註。
% 這裡用 zref-abspage 的絕對頁序，不用 \thepage，避免文件內頁碼重置時誤判。
\newif\ifclassroommode
\newif\ifclassroomabbrzh
\classroommodefalse
\classroomabbrzhtrue
\newcommand{\classroommodeon}{\classroommodetrue}
\newcommand{\classroommodeoff}{\classroommodefalse}
\newcommand{\classroomabbrzhon}{\classroomabbrzhtrue}
\newcommand{\classroomabbrzhoff}{\classroomabbrzhfalse}
\newcommand{\abbrnotelabel}{\emph{Abk.}: }
\newcommand{\abbrnote}[3]{%
  \ifclassroommode
    \ifshoworiginal
      \ifintranslation\else
        \footnote{\normalfont\footnotesize\abbrnotelabel #2\ifclassroomabbrzh；{\songfont #3}\fi}%
      \fi
    \fi
  \fi
}
\newcommand{\abbrnoteonce}[4]{%
  #2%
  \ifclassroommode
    \ifshoworiginal
      \ifintranslation\else
        \begingroup
          \edef\abbrcurrentabspage{\number\numexpr\value{abspage}+1\relax}%
          \ifcsname abbrnoteseen@#1@\abbrcurrentabspage\endcsname\else
            \global\expandafter\let\csname abbrnoteseen@#1@\abbrcurrentabspage\endcsname\empty
            \abbrnote{#1}{#3}{#4}%
          \fi
        \endgroup
      \fi
    \fi
  \fi
}
\newcommand{\abbrAbs}{\abbrnoteonce{abs}{Abs.}{Absatz}{項}}
\newcommand{\abbrAAO}{\abbrnoteonce{aao}{a.a.O.}{am angegebenen Ort}{前引處}}
\newcommand{\abbrAE}{\abbrnoteonce{ae}{AE}{Alternativ-Entwurf}{替代草案}}
\newcommand{\abbrArt}{\abbrnoteonce{art}{Art.}{Artikel}{條}}
\newcommand{\abbrAufl}{\abbrnoteonce{aufl}{Aufl.}{Auflage}{版}}
\newcommand{\abbrBd}{\abbrnoteonce{bd}{Bd.}{Band}{卷}}
\newcommand{\abbrBGBl}{\abbrnoteonce{bgbl}{BGBl.}{Bundesgesetzblatt}{聯邦法律公報}}
\newcommand{\abbrBGHSt}{\abbrnoteonce{bghst}{BGHSt}{Entscheidungen des Bundesgerichtshofes in Strafsachen}{聯邦普通法院刑事判決彙編}}
\newcommand{\abbrBRDrucks}{\abbrnoteonce{brdrucks}{BRDrucks.}{Bundesrats-Drucksache}{聯邦參議院文件}}
\newcommand{\abbrBTDrucks}{\abbrnoteonce{btdrucks}{BTDrucks.}{Bundestags-Drucksache}{聯邦議院文件}}
\newcommand{\abbrBundesgesetzbl}{\abbrnoteonce{bundesgesetzbl}{Bundesgesetzbl.}{Bundesgesetzblatt}{聯邦法律公報}}
\newcommand{\abbrBvF}{\abbrnoteonce{bvf}{BvF}{Bundesverfassungsgerichtsverfahren}{聯邦憲法法院程序案號}}
\newcommand{\abbrBvL}{\abbrnoteonce{bvl}{BvL}{Bundesverfassungsgerichtsverfahren}{聯邦憲法法院程序案號}}
\newcommand{\abbrBvR}{\abbrnoteonce{bvr}{BvR}{Bundesverfassungsgerichtsverfahren}{聯邦憲法法院程序案號}}
\newcommand{\abbrBVerfG}{\abbrnoteonce{bverfg}{BVerfG}{Bundesverfassungsgericht}{聯邦憲法法院}}
\newcommand{\abbrBVerfGE}{\abbrnoteonce{bverfge}{BVerfGE}{Entscheidungen des Bundesverfassungsgerichts}{聯邦憲法法院判決集}}
\newcommand{\abbrBVerfGG}{\abbrnoteonce{bverfgg}{BVerfGG}{Bundesverfassungsgerichtsgesetz}{聯邦憲法法院法}}
\newcommand{\abbrDDR}{\abbrnoteonce{ddr}{DDR}{Deutsche Demokratische Republik}{德意志民主共和國；東德}}
\newcommand{\abbrDJT}{\abbrnoteonce{djt}{DJT}{Deutscher Juristentag}{德國法學家大會}}
\newcommand{\abbrDr}{\abbrnoteonce{dr}{Dr.}{Doktor}{Dr.}}
\newcommand{\abbrEuGRZ}{\abbrnoteonce{eugrz}{EuGRZ}{Europäische Grundrechte-Zeitschrift}{歐洲基本權期刊}}
\newcommand{\abbrEV}{\abbrnoteonce{ev}{EV}{Einigungsvertrag}{統一條約}}
\newcommand{\abbrF}{\abbrnoteonce{f}{f.}{folgende}{以下一頁；下一段}}
\newcommand{\abbrFF}{\abbrnoteonce{ff}{ff.}{fortfolgende}{以下數頁；以下各段}}
\newcommand{\abbrGBlDDR}{\abbrnoteonce{gblddr}{GBl. DDR}{Gesetzblatt der Deutschen Demokratischen Republik}{德意志民主共和國法律公報}}
\newcommand{\abbrGG}{\abbrnoteonce{gg}{GG}{Grundgesetz}{基本法}}
\newcommand{\abbrGVBl}{\abbrnoteonce{gvbl}{GVBl.}{Gesetz- und Verordnungsblatt}{法律及命令公報}}
\newcommand{\abbrJR}{\abbrnoteonce{jr}{JR}{Juristische Rundschau}{法學評論}}
\newcommand{\abbrJZ}{\abbrnoteonce{jz}{JZ}{JuristenZeitung}{法學家報}}
\newcommand{\abbrIVm}{\abbrnoteonce{ivm}{i.V.m.}{in Verbindung mit}{結合；連同}}
\newcommand{\abbrMdB}{\abbrnoteonce{mdb}{MdB}{Mitglied des Bundestages}{聯邦議院議員}}
\newcommand{\abbrMwN}{\abbrnoteonce{mwn}{m.w.N.}{mit weiteren Nachweisen}{附進一步文獻／判例出處}}
\newcommand{\abbrNF}{\abbrnoteonce{nf}{n.F.}{neue Fassung}{新版本}}
\newcommand{\abbrAF}{\abbrnoteonce{af}{a.F.}{alte Fassung}{舊版本}}
\newcommand{\abbrNr}{\abbrnoteonce{nr}{Nr.}{Nummer}{款、目或號，依語境}}
\newcommand{\abbrProf}{\abbrnoteonce{prof}{Prof.}{Professor}{教授}}
\newcommand{\abbrRdnr}{\abbrnoteonce{rdnr}{Rdnr.}{Randnummer}{邊碼；段碼}}
\newcommand{\abbrRdnrn}{\abbrnoteonce{rdnrn}{Rdnrn.}{Randnummern}{邊碼；段碼}}
\newcommand{\abbrRGBl}{\abbrnoteonce{rgbl}{RGBl.}{Reichsgesetzblatt}{帝國法律公報}}
\newcommand{\abbrRGSt}{\abbrnoteonce{rgst}{RGSt}{Entscheidungen des Reichsgerichts in Strafsachen}{帝國法院刑事判決彙編}}
\newcommand{\abbrRspr}{\abbrnoteonce{rspr}{Rspr.}{Rechtsprechung}{判例；司法實務}}
\newcommand{\abbrRVO}{\abbrnoteonce{rvo}{RVO}{Reichsversicherungsordnung}{帝國保險法}}
\newcommand{\abbrS}{\abbrnoteonce{s}{S.}{Seite}{頁}}
\newcommand{\abbrSFHG}{\abbrnoteonce{sfhg}{SFHG}{Schwangeren- und Familienhilfegesetz}{孕婦及家庭扶助法}}
\newcommand{\abbrSGBV}{\abbrnoteonce{sgbv}{SGB V}{Fünftes Buch Sozialgesetzbuch}{社會法典第五編}}
\newcommand{\abbrStAG}{\abbrnoteonce{staeg}{StÄG}{Strafrechtsänderungsgesetz}{刑法修正法}}
\newcommand{\abbrStenBer}{\abbrnoteonce{stenber}{StenBer.}{Stenographischer Bericht}{速記紀錄}}
\newcommand{\abbrStGB}{\abbrnoteonce{stgb}{StGB}{Strafgesetzbuch}{刑法}}
\newcommand{\abbrStREG}{\abbrnoteonce{streg}{StREG}{Strafrechtsreform-Ergänzungsgesetz}{刑法改革補充法}}
\newcommand{\abbrStrRG}{\abbrnoteonce{strrg}{StrRG}{Strafrechtsreformgesetz}{刑法改革法}}
\newcommand{\abbrUS}{\abbrnoteonce{us}{U.S.}{United States Reports}{美國最高法院判決彙編}}
\newcommand{\abbrVgl}{\abbrnoteonce{vgl}{vgl.}{vergleiche}{參見}}
\newcommand{\abbrVVDStRL}{\abbrnoteonce{vvdstrl}{VVDStRL}{Veröffentlichungen der Vereinigung der Deutschen Staatsrechtslehrer}{德國公法教師協會論文集}}
\newcommand{\abbrWP}{\abbrnoteonce{wp}{WP.}{Wahlperiode}{屆期}}
\newcommand{\abbrWp}{\abbrnoteonce{wp}{Wp.}{Wahlperiode}{屆期}}
\newcommand{\abbrZStrW}{\abbrnoteonce{zstrw}{ZStrW}{Zeitschrift für die gesamte Strafrechtswissenschaft}{整體刑法學期刊}}
\newcommand{\recordsourcerandstart}{%
  \ifrandnumbers
    \setcounter{lastsourcerandstart}{\value{randnumber}}%
    \global\lastsourcehasrandnumberstrue
  \else
    \global\lastsourcehasrandnumbersfalse
  \fi
}
\newlength{\biparagraphskip}
\setlength{\biparagraphskip}{0.52em}
\newlength{\lawboxblockbeforeskip}
\setlength{\lawboxblockbeforeskip}{0.72em}
\newlength{\lawboxblockafterskip}
\setlength{\lawboxblockafterskip}{0.92em}
\newif\iflawboxpartendedblock
\lawboxpartendedblockfalse
\newcommand{\sourceparstyle}{\footnotesize\color{sourceink}\setstretch{1.12}\RaggedRight\emergencystretch=3em\setlength{\parskip}{0.42em}\obeylines\selectfont}
\newcommand{\transparstyle}{\songfont\color{caseink}\setstretch{1.12}\setlength{\parindent}{0pt}\setlength{\parskip}{0.42em}}
\newcommand{\bilingualpairskip}{%
  \par\addvspace{\biparagraphskip}%
  \switchcolumn
  \par\addvspace{\biparagraphskip}%
  \switchcolumn*%
  \global\intranslationfalse
}
\newcommand{\bilingualboxpairskip}{%
  \par
  \switchcolumn
  \par
  \switchcolumn*%
  \global\intranslationfalse
}
\newcommand{\bilinguallawboxblockskip}{%
  \switchcolumn*%
  \par\addvspace{\lawboxblockafterskip}%
  \switchcolumn
  \par\addvspace{\lawboxblockafterskip}%
  \switchcolumn*%
  \global\intranslationfalse
}
\newcommand{\bikeepwithnext}[1][4]{%
  \ifbilingualcolumns
    \syncleft
    \Needspace{#1\baselineskip}%
    \switchcolumn
    \Needspace{#1\baselineskip}%
    \switchcolumn*%
    \global\intranslationfalse
  \else
    \Needspace{#1\baselineskip}%
  \fi
}
\NewEnviron{sourcepara}[1][]{
  \recordsourcerandstart
  \ifshoworiginal
    \ifbilingualcolumns
      \ifintranslation\switchcolumn*\global\intranslationfalse\fi
    \else
      \Needspace{5\baselineskip}
    \fi
    \ifbilingualcolumns\else\par\addvspace{0.35em}\fi
    {\sourceparstyle
    \noindent\BODY\par}
    \ifbilingualcolumns\else\addvspace{0.25em}\fi
    \ifbilingualcolumns
      \switchcolumn
      \global\intranslationtrue
    \fi
  \else
    \ifrandnumbers
      \begingroup
        \setbox0=\vbox{\hsize=\linewidth\sourceparstyle\noindent\BODY\par}%
      \endgroup
    \fi
  \fi
}
\NewEnviron{transpara}{
  \ifshowtranslation
    \setcounter{randnumberlocal}{\value{lastsourcerandstart}}%
    \ifbilingualcolumns
      \ifintranslation\else\switchcolumn\global\intranslationtrue\fi
      {\transparstyle\setlength{\leftskip}{0.55em}\setlength{\rightskip}{0.15em plus 1em}\addtolength{\linewidth}{-0.55em}\BODY\par}
      \switchcolumn*
      \bilingualpairskip
    \else
      {\transparstyle\BODY\par}
    \fi
  \else
    \ifbilingualcolumns
      \ifintranslation\switchcolumn*\global\intranslationfalse\fi
    \fi
  \fi
}
\NewEnviron{sourceboxpara}[1][]{
  \global\lawboxpartendedblockfalse
  \recordsourcerandstart
  \ifshoworiginal
    \ifbilingualcolumns
      \ifintranslation\switchcolumn*\global\intranslationfalse\fi
    \else
      \Needspace{3\baselineskip}
    \fi
    {\sourceparstyle
    \BODY\par}
    \ifbilingualcolumns\else
      \iflawboxpartendedblock\addvspace{\lawboxblockafterskip}\fi
    \fi
    \ifbilingualcolumns
      \switchcolumn
      \global\intranslationtrue
    \fi
  \fi
}
\NewEnviron{transboxpara}{
  \global\lawboxpartendedblockfalse
  \ifshowtranslation
    \setcounter{randnumberlocal}{\value{lastsourcerandstart}}%
    \ifbilingualcolumns
      \ifintranslation\else\switchcolumn\global\intranslationtrue\fi
      {\transparstyle\BODY\par}
      \iflawboxpartendedblock
        \bilinguallawboxblockskip
      \else
        \switchcolumn*
        \bilingualboxpairskip
      \fi
    \else
      {\transparstyle\BODY\par}
      \iflawboxpartendedblock\addvspace{\lawboxblockafterskip}\fi
    \fi
  \else
    \ifbilingualcolumns
      \ifintranslation\switchcolumn*\global\intranslationfalse\fi
    \fi
  \fi
}
\newcommand{\leitsatznum}[1]{\noindent\hangindent=3.0em\hangafter=1\makebox[2.25em][r]{\bfseries #1.}\hspace{0.55em}\ignorespaces}
\newcommand{\syncleft}{\ifbilingualcolumns\ifintranslation\switchcolumn*\global\intranslationfalse\fi\fi}
\pretocmd{\section}{\syncleft\clearpage}{}{}
\pretocmd{\subsection}{\syncleft}{}{}
\pretocmd{\subsubsection}{\syncleft}{}{}
\newcommand{\biheadingfont}{\songfont}
\newcommand{\bitocentry}[2]{%
  \ifshoworiginal
    \ifshowtranslation
      \ifstrequal{#1}{#2}{#1}{#1\space #2}%
    \else
      #1%
    \fi
  \else
    #2%
  \fi
}
\pdfstringdefDisableCommands{%
  \def\bitocentry#1#2{#1}%
}
\newcommand{\biheading}[7]{%
  \syncleft#1\bikeepwithnext[#3]\phantomsection\addcontentsline{toc}{#2}{\bitocentry{#6}{#7}}%
  \ifbilingualcolumns
    {#4 #6\par}%
    \switchcolumn
    {#4\biheadingfont #7\par}%
    \switchcolumn*%
    \global\intranslationfalse
    \par\addvspace{#5}%
    \switchcolumn
    \par\addvspace{#5}%
    \switchcolumn*%
  \else
    \ifshoworiginal{#4 #6\par}\fi
    \ifshowtranslation{#4\biheadingfont #7\par}\fi
    \addvspace{#5}%
  \fi
}
\newcommand{\termsectionheading}[1]{%
  \par\Needspace{9\baselineskip}%
  \vspace{0.65em}%
  {\songfont\bfseries #1\par}%
  \vspace{0.2em}%
}
% 章節換頁：
%   雙語 paracol 欄內不要用 \clearpage；它會在同步兩欄時吐出只有欄線/頁碼的空頁。
%   \newpage 足以讓下一個 section 起新頁，又不會強制清空所有待排材料。
%   單語模式不在 paracol 內，仍可用 \clearpage。
\newcommand{\bisectionpagebreak}{\ifbilingualcolumns\newpage\else\clearpage\fi}
\newcommand{\bisection}[2]{%
  \biheading{\iffirstbilingualsection\global\firstbilingualsectionfalse\else\bisectionpagebreak\fi}{section}{6}{\centering\Large\bfseries}{0.8em}{#1}{#2}%
}
\newcommand{\bisubsection}[2]{%
  \biheading{}{subsection}{5}{\large\bfseries}{0.55em}{#1}{#2}%
}
\newcommand{\bisubsubsection}[2]{%
  \biheading{}{subsubsection}{4}{\normalsize\bfseries}{0.45em}{#1}{#2}%
}
\newcommand{\bicompositeheading}[4]{%
  \biheading{}{subsection}{5}{\large\bfseries}{0.16em}{#1}{#2}%
  \biheading{}{subsubsection}{3}{\normalsize\bfseries}{0.45em}{#3}{#4}%
}
\newif\iflawboxfirstitem
\lawboxfirstitemfalse
\newcommand{\lawboxprespace}[1]{%
  \par
  \iflawboxfirstitem
    \global\lawboxfirstitemfalse
  \else
    \addvspace{#1}%
  \fi
}
\newtcolorbox{lawboxinner}{
  caseboxbase,
  colframe=sourceline!85!black,
  colback=white,
  left=1.05em,
  right=1.05em,
  top=0.62em,
  bottom=0.62em,
  before upper={\global\inlawboxtrue\global\lawboxfirstitemtrue\global\lawboxhaspendingrnfalse\ifintranslation\lawboxtransfont\else\lawboxsourcefont\fi\RaggedRight\setlength{\parindent}{0pt}\setlength{\parskip}{0pt}\ifintranslation\setstretch{\lawboxtransstretch}\else\setstretch{\lawboxsourcestretch}\fi},
  after upper={\global\lawboxhaspendingrnfalse\global\inlawboxfalse}
}
\tcbset{
  lawboxpartbase/.style={
    enhanced,
    breakable=false,
    width=0.985\linewidth,
    colback=white,
    frame hidden,
    boxrule=0pt,
    arc=0pt,
    left=1.05em,
    right=1.05em,
    boxsep=1.5mm,
    before skip=0pt,
    after skip=0pt,
    before upper={\global\inlawboxtrue\global\lawboxfirstitemtrue\global\lawboxhaspendingrnfalse\ifintranslation\lawboxtransfont\else\lawboxsourcefont\fi\RaggedRight\setlength{\parindent}{0pt}\setlength{\parskip}{0pt}\ifintranslation\setstretch{\lawboxtransstretch}\else\setstretch{\lawboxsourcestretch}\fi},
    after upper={\global\lawboxhaspendingrnfalse\global\inlawboxfalse},
    borderline west={0.28pt}{0pt}{sourceline!85!black},
    borderline east={0.28pt}{0pt}{sourceline!85!black},
  },
  lawboxpartfirst/.style={
    lawboxpartbase,
    top=0.62em,
    bottom=0.04em,
    before skip=\lawboxblockbeforeskip,
    borderline north={0.28pt}{0pt}{sourceline!85!black},
  },
  lawboxpartmiddle/.style={
    lawboxpartbase,
    top=0.04em,
    bottom=0.04em,
  },
  lawboxpartlast/.style={
    lawboxpartbase,
    top=0.04em,
    bottom=0.62em,
    after skip=0pt,
    borderline south={0.28pt}{0pt}{sourceline!85!black},
  }
}
\newtcolorbox{lawboxpartinner}[1]{#1}
\newtcolorbox{termboxinner}{
  caseboxbase,
  colframe=noteblue!70!black,
  colback=casepaper,
  before upper={\songfont\setlength{\parindent}{0pt}\setlength{\parskip}{0.35em}}
}
\newtcolorbox{deboxinner}{
  caseboxbase,
  colframe=notegray!72!black,
  colback=white,
  left=1.05em,
  right=1.05em,
  top=0.62em,
  bottom=0.62em,
  before upper={\global\inlawboxtrue\global\lawboxfirstitemtrue\global\lawboxhaspendingrnfalse\small\setlength{\parindent}{0pt}\setlength{\parskip}{0.35em}},
  after upper={\global\lawboxhaspendingrnfalse\global\inlawboxfalse}
}
\NewEnviron{lawbox}{
  \begin{lawboxinner}\BODY\end{lawboxinner}
}
\NewEnviron{lawboxpart}[2]{
  \begin{lawboxpartinner}{lawboxpart#1,equal height group=#2}\BODY\end{lawboxpartinner}%
  \ifstrequal{#1}{last}{\global\lawboxpartendedblocktrue}{}%
}
\NewEnviron{termbox}{
  \par\addvspace{0.55em}
  {\color{noteblue!70!black}\rule{\linewidth}{0.28pt}}\par
  \begingroup
    \songfont\small\color{caseink}
    \setlength{\parindent}{0pt}
    \setlength{\parskip}{0.24em}
    \BODY
    \par
  \endgroup
  {\color{noteblue!70!black}\rule{\linewidth}{0.28pt}}\par
  \addvspace{0.55em}
}
\NewEnviron{debox}{
  \begin{deboxinner}\BODY\end{deboxinner}
}
\providecommand{\paragraphmark}[1]{}
\newcommand{\lawboxtitleafter}{\addvspace{0.06em}}
\newcommand{\lawtitle}[1]{\lawboxprespace{0.22em}{\lawboxtransfont\bfseries\lawboxuserandnumber #1}\par\lawboxtitleafter}
\newcommand{\absno}[2]{\lawboxprespace{0.04em}\noindent\lawboxuserandnumber\hangindent=2.6em\hangafter=1\makebox[2.6em][l]{(#1)}#2\par}
\newcommand{\nrno}[2]{\lawboxprespace{0pt}\noindent\lawboxuserandnumber\hspace*{2.6em}\hangindent=4.4em\hangafter=1\makebox[1.8em][l]{#1.}#2\par}
\newcommand{\letterno}[2]{\lawboxprespace{0pt}\noindent\lawboxuserandnumber\hspace*{4.4em}\hangindent=6.0em\hangafter=1\makebox[1.6em][l]{#1)}#2\par}
\newcommand{\sourcelawtitle}[1]{\lawboxprespace{0.22em}{\lawboxsourcefont\bfseries\lawboxuserandnumber #1}\par\lawboxtitleafter}
\newcommand{\sourcelawsubtitle}[1]{\lawboxprespace{0.04em}{\lawboxsourcefont\bfseries\lawboxuserandnumber #1}\par\lawboxtitleafter}
\newcommand{\sourceabsno}[2]{\lawboxprespace{0.04em}\noindent\lawboxuserandnumber\hangindent=2.45em\hangafter=1\makebox[2.45em][l]{(#1)}#2\par}
\newcommand{\sourcenrno}[2]{\lawboxprespace{0pt}\noindent\lawboxuserandnumber\hspace*{2.45em}\hangindent=4.25em\hangafter=1\makebox[1.8em][l]{#1.}#2\par}
\newcommand{\sourceletterno}[2]{\lawboxprespace{0pt}\noindent\lawboxuserandnumber\hspace*{4.25em}\hangindent=5.85em\hangafter=1\makebox[1.6em][l]{#1)}#2\par}
\newcommand{\sourcequotepara}[1]{\lawboxprespace{0.04em}\noindent\lawboxuserandnumber\hangindent=1.2em\hangafter=1 #1\par}
\newcommand{\sourcelawline}[1]{\lawboxprespace{0.04em}\noindent\lawboxuserandnumber #1\par}
\newcommand{\lawline}[1]{\lawboxprespace{0.04em}\noindent\lawboxuserandnumber #1\par}
\newcommand{\pendingtranslation}{\par{\songfont\footnotesize\color{pendingink}（中文譯文待補。）}\par}
\newcommand{\tocpart}[1]{\par\addvspace{0.52em}{\footnotesize\bfseries\color{pendingink}#1\par}\addvspace{0.16em}}
\newcommand{\tocdashline}{\par\addvspace{0.48em}\noindent{\color{notegray}\xleaders\hbox to 0.82em{\hfil\rule{0.42em}{0.28pt}\hfil}\hskip\linewidth}\par\addvspace{0.38em}}
% \tocpanel{font-cmd}{content}：將目錄置中於所在欄寬（雙語欄適用）。
% A3 橫式使用 0.58\linewidth，A4 橫式使用 0.84\linewidth，
% 使兩種紙張下目錄欄內容實際寬度相近（均約 100–105 mm）。
\newcommand{\tocpanel}[2]{%
  \makebox[\linewidth][c]{%
    \ifx\bilingualpaper\bilingualpaperaThree
      \begin{minipage}{0.58\linewidth}%
    \else
      \begin{minipage}{0.84\linewidth}%
    \fi
    \toccolstyle
    #1#2%
    \end{minipage}%
  }%
}
% ===== 首頁簡目錄的兩層 description list 縮排 =====
% enumitem 的 description list 可把每一列拆成：
%   [label]  labelsep  body text
% 這裡的「起點」都以所在 minipage/欄位的左緣為 0。
%
% 核心規則：
%   labelindent = label 欄位從左緣開始的位置。
%   labelwidth  = label 欄位本身寬度；標籤太長（如 Abw. Meinung）就加大它。
%   labelsep    = label 欄位右緣到正文的距離。
%   leftmargin  = 正文 body text 的起點。判斷層級視覺時主要看這個。
%
% 所以：
%   想讓「Begründung / 判決理由」整列正文往右或往左：改 tocitems 的 leftmargin。
%   想讓「A. Verfahren und Gegenstand / A. 程序與審查標的」正文往右或往左：
%     改 tocsubitems 的 leftmargin。
%   想讓 A./B./C. 這些標籤本身往右或往左，但正文不動：改 tocsubitems 的 labelindent。
%   想讓 A./B./C. 與正文間距變大或變小：改 tocsubitems 的 labelsep。
%
% 層級原則：
%   tocsubitems.leftmargin 必須大於 tocitems.leftmargin，
%   否則子層級正文會看起來比「Gründe / 理由」還靠左。
\newlist{tocitems}{description}{1}
\setlist[tocitems]{%
  labelindent=0pt,    % 第一層 label 欄位起點；0pt 表示貼齊目錄欄左緣。
  leftmargin=2.54cm,  % 第一層正文起點；例如 Begründung / 判決理由 的左緣。
  labelwidth=2.16cm,  % 第一層 label 欄位寬度；需容納 Leitsätze、Abw. Meinung、不同意見等。
  labelsep=0.32cm,    % 第一層 label 與正文的水平間距；只改間距，不改正文起點。
  itemsep=0.28em,     % 第一層各項之間的垂直距離。
  topsep=0.20em,      % 第一層 list 上下與前後內容的垂直距離。
  parsep=0pt,         % 同一 item 內段落之間的距離；目錄項通常不需要。
  partopsep=0pt,      % list 若緊接段落開始時的額外距離；關閉以免目錄高度飄動。
  align=left,         % label 欄內左對齊；長短標籤不靠右貼齊。
  font=\normalfont\bfseries % label 的字型；正文不受這行影響。
}
\newlist{tocsubitems}{description}{1}
\setlist[tocsubitems]{%
  labelindent=1cm,    % 第二層 label 起點；只控制 A./B./C. 本身的位置。
  leftmargin=3.00cm,  % 第二層正文起點；必須大於 2.54cm，才會比 Begründung / 判決理由更右。
  labelwidth=0.5cm,   % 第二層 label 欄寬；A./B./C. 很短，可小於第一層。
  labelsep=1.5cm,    % 第二層 label 與正文間距；A. 與 Verfahren / 程序 的距離看這裡。
  itemsep=0.12em,     % 第二層各項較密，避免 A.–E. 佔太高。
  topsep=0.04em,      % 第二層 list 與 Gründe / 理由之間的垂直距離。
  parsep=0pt,         % 同一第二層 item 內段落距離；目錄項通常不需要。
  partopsep=0pt,      % 關閉額外段落起始距離，避免版面不穩。
  align=left,         % A./B./C. label 欄內左對齊。
  font=\normalfont\bfseries, % A./B./C. label 加粗；正文不受這行影響。
  before={}           % 不在第二層 list 前自動插入額外內容。
}
\newlist{termitems}{description}{1}
\ifbilingualcolumns
  \setlist[termitems]{%
    leftmargin=5.2cm,
    labelwidth=4.8cm,
    labelsep=0.35cm,
    itemsep=0.16em,
    topsep=0.2em,
    parsep=0pt,
    partopsep=0pt,
    font=\normalfont\bfseries,
    style=nextline
  }
\else
  \setlist[termitems]{%
    leftmargin=0pt,
    labelwidth=0pt,
    labelsep=0pt,
    itemsep=0.24em,
    topsep=0.2em,
    parsep=0pt,
    partopsep=0pt,
    font=\normalfont\bfseries,
    style=nextline
  }
\fi
\newlist{abbritems}{description}{1}
\setlist[abbritems]{%
  leftmargin=2.38cm,
  labelwidth=2.02cm,
  labelsep=0.28cm,
  itemsep=0.16em,
  topsep=0.2em,
  parsep=0pt,
  partopsep=0pt,
  align=left,
  font=\normalfont\bfseries
}
\ifbilingualcolumns
  \newenvironment{termcolumns}{\begin{multicols}{2}}{\end{multicols}}
\else
  \newenvironment{termcolumns}{}{}
\fi

% ===== 編者註腳開關 =====
% 一般版預設不顯示編者註，避免正文腳註過密。
% 若開啟 \classroommodeon，GG/條文編者註仍會自動顯示，無需另開本開關。
% 若一般版也要顯示編者註，可在單案檔 \input 後加 \showeditornotestrue。
\newif\ifshoweditornotes
\showeditornotesfalse

% ===== 目錄排版輔助 =====
\newcommand{\toccolstyle}{\raggedright\arraybackslash}
\newcommand{\tocsingleindent}{\leftskip=1.45cm\rightskip=1.2cm}

% ===== 行內引文環境（德文原文與中文譯文各一，帶左側灰色邊線）=====
\newcommand{\quoteparbreak}{\par\addvspace{0.48em}}
\newcommand{\sourcequoteinner}[1]{%
  \begin{tcolorbox}[
    enhanced, breakable, width=\dimexpr\linewidth-0.8em\relax, frame hidden,
    colback=white, coltext=caseink, boxrule=0pt,
    borderline west={0.55pt}{0.88em}{notegray},
    left=1.78em, right=0.72em, top=0.18em, bottom=0.18em,
    boxsep=0pt, before skip=0.24em, after skip=0.24em,
    before upper={\normalfont\footnotesize\color{caseink}\setlength{\parindent}{0pt}\setlength{\parskip}{0.34em}\raggedright\setlength{\rightskip}{0pt plus 1em}}
  ]%
  #1\par
  \end{tcolorbox}%
}
\newcommand{\transquoteinner}[1]{%
  \begin{tcolorbox}[
    enhanced, breakable, width=\dimexpr\linewidth-0.8em\relax, frame hidden,
    colback=white, coltext=caseink, boxrule=0pt,
    borderline west={0.55pt}{0.88em}{notegray},
    left=1.78em, right=0.72em, top=0.18em, bottom=0.18em,
    boxsep=0pt, before skip=0.24em, after skip=0.24em,
    before upper={\songfont\small\color{caseink}\setlength{\parindent}{0pt}\setlength{\parskip}{0.34em}\raggedright\setlength{\rightskip}{0pt plus 1em}}
  ]%
  #1\par
  \end{tcolorbox}%
}

% ===== 大綱（Gliederung）Rn 輔助命令 =====
\newcommand{\outrn}[2]{\enspace{\normalfont\footnotesize\color{pendingink}(Rn\,#1–#2)}}
\newcommand{\outrnsingle}[1]{\enspace{\normalfont\footnotesize\color{pendingink}(Rn\,#1)}}
\newcommand{\outrnzh}[2]{\enspace{\normalfont\footnotesize\color{pendingink}（第\,#1–#2\,段）}}
\newcommand{\outrnzhs}[1]{\enspace{\normalfont\footnotesize\color{pendingink}（第\,#1\,段）}}
% ===== 大綱雙語對照表格輔助命令（三欄：段碼 │ 德文 │ 中文）=====
% 欄 1：段碼（由欄定義自動格式化：小字、等寬、右對齊、pendingink）
% 欄 2：德文標籤＋正文
% 欄 3：中文標籤＋正文
%
% 無段碼的列（\omainrow、\osecrow），欄 1 以 & 空置。
% 有段碼的列（\osecrowrn、\osubrow、\ointrow），#1 為預格式化字串
%   例：\osubrow{2–49}{I.}{...}{I.}{...}
%       \ointrow{107}{...}{...}
%
% \opartrow：段組標題（橫跨三欄）
\newcommand{\opartrow}[2]{%
  \noalign{\vspace{0.30em}}%
  \multicolumn{3}{@{}l@{}}{%
    \footnotesize\bfseries\color{pendingink}#1\hfill\songfont#2%
  }\\[0.06em]%
}
% \odashrow：虛線分隔（橫跨三欄）
\newcommand{\odashrow}{%
  \noalign{\vspace{0.28em}}%
  \multicolumn{3}{@{}l@{}}{\color{notegray}%
    \xleaders\hbox to 0.82em{\hfil\rule{0.42em}{0.28pt}\hfil}\hskip 0.88\linewidth}%
  \\[0.24em]%
}
% \odissentpart：不同意見書區段標頭。比一般 \opartrow 更像附錄性轉場。
\newcommand{\odissentpart}[2]{%
  \noalign{\vspace{0.42em}}%
  \multicolumn{3}{@{}p{0.88\textwidth}@{}}{%
    \color{notegray}\rule{\linewidth}{0.18pt}%
  }\\[-0.18em]%
  \multicolumn{3}{@{}p{0.88\textwidth}@{}}{%
    \makebox[\linewidth][s]{%
      \footnotesize\bfseries\color{pendingink}#1%
      \hfill{\songfont #2}%
    }%
  }\\[-0.02em]%
}
% \odissentopinion：不同意見書的署名與一行摘要，右欄保留段碼範圍。
\newcommand{\odissentopinion}[5]{%
  \footnotesize\hangindent=0.52cm\hangafter=1%
    {\bfseries\color{caseink}#2}\enspace{\color{notegray}\textperiodcentered}\enspace\color{caseink}#3%
  &\footnotesize\hangindent=0.52cm\hangafter=1%
    {\songfont\bfseries\color{caseink}#4}\enspace{\color{notegray}\textperiodcentered}\enspace\songfont\color{caseink}#5%
  &#1%
  \\[0.12em]%
}
% \omainrow：頂層項目，無段碼（Leitsätze、Urteil、Gründe …）
\newcommand{\omainrow}[4]{%
  \footnotesize\hangindent=1.92cm\hangafter=1%
    \makebox[1.92cm][l]{\bfseries\color{caseink}#1}\color{caseink}#2%
  &\footnotesize\hangindent=1.92cm\hangafter=1%
    \makebox[1.92cm][l]{\bfseries\color{caseink}#3}\songfont\color{caseink}#4%
  &%
  \\[0.15em]%
}
% \omainrowrn：頂層項目，有段碼（不同意見等標籤寬於 \osecrow 框者）
\newcommand{\omainrowrn}[5]{%
  \footnotesize\hangindent=1.92cm\hangafter=1%
    \makebox[1.92cm][l]{\bfseries\color{caseink}#2}\color{caseink}#3%
  &\footnotesize\hangindent=1.92cm\hangafter=1%
    \makebox[1.92cm][l]{\bfseries\color{caseink}#4}\songfont\color{caseink}#5%
  &#1%
  \\[0.15em]%
}
% \osecrow：一級子項，無段碼（A.、C.、D. 等有子項者）
\newcommand{\osecrow}[4]{%
  \footnotesize\hspace{1.10cm}\hangindent=1.98cm\hangafter=1%
    \makebox[0.86cm][l]{\bfseries\color{caseink}#1}\color{caseink}#2%
  &\footnotesize\hspace{1.10cm}\hangindent=1.98cm\hangafter=1%
    \makebox[0.86cm][l]{\bfseries\color{caseink}#3}\songfont\color{caseink}#4%
  &%
  \\[0.11em]%
}
% \osecrowrn：一級子項，有段碼（B.、E. 等完整段落範圍者）
% #1=段碼字串（如 101–106）, #2=DE標籤, #3=DE正文, #4=ZH標籤, #5=ZH正文
\newcommand{\osecrowrn}[5]{%
  \footnotesize\hspace{1.10cm}\hangindent=1.98cm\hangafter=1%
    \makebox[0.86cm][l]{\bfseries\color{caseink}#2}\color{caseink}#3%
  &\footnotesize\hspace{1.10cm}\hangindent=1.98cm\hangafter=1%
    \makebox[0.86cm][l]{\bfseries\color{caseink}#4}\songfont\color{caseink}#5%
  &#1%
  \\[0.11em]%
}
% \osubrow：二級子項，有段碼範圍（I.、II.、A.I. …）
% #1=段碼字串, #2=DE標籤, #3=DE正文, #4=ZH標籤, #5=ZH正文
\newcommand{\osubrow}[5]{%
  \footnotesize\hspace{1.40cm}\hangindent=2.30cm\hangafter=1%
    \makebox[0.90cm][l]{\bfseries\color{caseink}#2}\color{caseink}#3%
  &\footnotesize\hspace{1.40cm}\hangindent=2.30cm\hangafter=1%
    \makebox[0.90cm][l]{\bfseries\color{caseink}#4}\songfont\color{caseink}#5%
  &#1%
  \\[0.09em]%
}
% \ointrow：引入段落，有單一段碼，無標籤
% #1=段碼字串, #2=DE正文, #3=ZH正文
\newcommand{\ointrow}[3]{%
  \footnotesize\hspace{0.52cm}\hangindent=0.52cm\hangafter=1\color{caseink}#2%
  &\footnotesize\hspace{0.52cm}\hangindent=0.52cm\hangafter=1\songfont\color{caseink}#3%
  &#1%
  \\[0.09em]%
}
% \osubsubrow：三級子項（1.、2.、3.）
% #1=段碼字串, #2=DE標籤, #3=DE正文, #4=ZH標籤, #5=ZH正文
\newcommand{\osubsubrow}[5]{%
  \footnotesize\hspace{1.80cm}\hangindent=2.48cm\hangafter=1%
    \makebox[0.68cm][l]{\bfseries\color{caseink}#2}\color{caseink}#3%
  &\footnotesize\hspace{1.80cm}\hangindent=2.48cm\hangafter=1%
    \makebox[0.68cm][l]{\bfseries\color{caseink}#4}\songfont\color{caseink}#5%
  &#1%
  \\[0.08em]%
}
% \osubsubsubrow：四級子項（a）、b））
% #1=段碼字串, #2=DE標籤, #3=DE正文, #4=ZH標籤, #5=ZH正文
\newcommand{\osubsubsubrow}[5]{%
  \footnotesize\hspace{2.10cm}\hangindent=2.68cm\hangafter=1%
    \makebox[0.58cm][l]{\bfseries\color{caseink}#2}\color{caseink}#3%
  &\footnotesize\hspace{2.10cm}\hangindent=2.68cm\hangafter=1%
    \makebox[0.58cm][l]{\bfseries\color{caseink}#4}\songfont\color{caseink}#5%
  &#1%
  \\[0.07em]%
}
% \osubsubsubsubrow：五級子項（aa）、bb））
% 標籤框 0.62cm：足以容納「aa)」在 footnotesize bold 下（約 0.50cm）並留有間距
% #1=段碼字串, #2=DE標籤, #3=DE正文, #4=ZH標籤, #5=ZH正文
\newcommand{\osubsubsubsubrow}[5]{%
  \footnotesize\hspace{2.38cm}\hangindent=3.06cm\hangafter=1%
    \makebox[0.68cm][l]{\bfseries\color{caseink}#2}\color{caseink}#3%
  &\footnotesize\hspace{2.38cm}\hangindent=3.06cm\hangafter=1%
    \makebox[0.68cm][l]{\bfseries\color{caseink}#4}\songfont\color{caseink}#5%
  &#1%
  \\[0.06em]%
}
% \outlinepage：雙語對照大綱頁包裝環境（longtable 自動跨頁）
% 三欄：德文欄 + 中文欄 + 段碼欄（最右，不折行）
% 段碼欄寬度：1.80cm，足以容納最長段碼如「309–348」
% DE/ZH 欄寬：(0.87\textwidth - 2.08cm) / 2
%   驗算：2×欄 + 0.05\textwidth + 0.28cm gap + 1.80cm = 0.87tw - 2.08cm + 0.05tw + 2.08cm = 0.92tw ✓
\newenvironment{outlinepage}{%
  \vspace*{0.40cm}%
  \setlength{\LTleft}{0.04\textwidth}%
  \setlength{\LTright}{0.04\textwidth}%
  \setlength{\tabcolsep}{0pt}%
  % 大綱列距由 longtable 的 \arraystretch 與各 row macro 結尾的 \\[...em] 共同決定。
  % 這裡控制「每一列本身的表格 strut 高度」；數值越大，A./I./1. 各項之間越像有空行。
  % A3 原本用 0.62，視覺上沒有空行；A4 也沿用同值，避免切到 A4 後項目被撐開。
  % 若日後覺得大綱太擠，只微調這個值（例如 0.68），不要改各 \omainrow/\osubrow 的縮排。
  \renewcommand{\arraystretch}{0.62}%
  \begin{longtable}{%
    >{\vspace{0pt}\raggedright\arraybackslash}p{\dimexpr(0.87\textwidth-2.08cm)/2\relax}%
    @{\hspace{0.025\textwidth}}%
    !{\color{notegray}\vrule width 0.12pt}%
    @{\hspace{0.025\textwidth}}%
    >{\vspace{0pt}\raggedright\arraybackslash}p{\dimexpr(0.87\textwidth-2.08cm)/2\relax}%
    @{\hspace{0.28cm}}%
    >{\vspace{0pt}\footnotesize\normalfont\color{pendingink}\raggedright\arraybackslash}p{1.80cm}%
    @{}%
  }%
  % 首頁欄標籤：不要橫跨置中；分別放在德文欄與中文欄的實際欄位起點。
  \footnotesize\bfseries\color{sourceink}Gliederung%
  &\footnotesize\bfseries\songfont\color{sourceink}大綱%
  &%
  \\[0.36em]%
  \endfirsthead%
  % 續頁欄標籤：同樣分欄定位，以灰色細體區分；無需標注「續」。
  \footnotesize\color{pendingink}Gliederung%
  &\footnotesize\songfont\color{pendingink}大綱%
  &%
  \\[0.24em]%
  \endhead%
}{%
  \end{longtable}%
}
 載入。
% 不含：\documentclass、\documentversion、\ggversionde/zh、\tocde/\toczh。
% 日期與首頁版本列由本檔自動產生；各文件只需手動指定 \documentversion。

\usepackage{xeCJK}
\usepackage{fontspec}
\usepackage{geometry}
\usepackage{setspace}
\usepackage{hyperref}
\usepackage{xurl}
\usepackage{bookmark}
\usepackage{enumitem}
\usepackage{xcolor}
\usepackage{titlesec}
\usepackage{array}
\usepackage{longtable}
\usepackage{tcolorbox}
\usepackage{environ}
\usepackage{pdflscape}
\usepackage{etoolbox}
\usepackage{ragged2e}
\usepackage{paracol}
\usepackage{multicol}
\usepackage{needspace}
\usepackage{fancyhdr}
\usepackage{tikz}
\usepackage{zref-abspage}
\tcbuselibrary{breakable,skins}
\usetikzlibrary{calc}
\usepackage{pgfcalendar}

\hypersetup{
  unicode=true,
  bookmarksopen=true,
  bookmarksnumbered=false,
  hidelinks
}

% ===== 自動推導，不要手動改下面 =====
% （\docmode 與 \bilingualpaper 由各文件在 % bverfge_layout.tex
% 共用排版基礎設施，供 Schwangerschaftsabbruch I 與 II 共享。
% 由各文件以 \documentclass 後 % bverfge_layout.tex
% 共用排版基礎設施，供 Schwangerschaftsabbruch I 與 II 共享。
% 由各文件以 \documentclass 後 \input{../../共用LaTeX/bverfge_layout} 載入。
% 不含：\documentclass、\documentversion、\ggversionde/zh、\tocde/\toczh。
% 日期與首頁版本列由本檔自動產生；各文件只需手動指定 \documentversion。

\usepackage{xeCJK}
\usepackage{fontspec}
\usepackage{geometry}
\usepackage{setspace}
\usepackage{hyperref}
\usepackage{xurl}
\usepackage{bookmark}
\usepackage{enumitem}
\usepackage{xcolor}
\usepackage{titlesec}
\usepackage{array}
\usepackage{longtable}
\usepackage{tcolorbox}
\usepackage{environ}
\usepackage{pdflscape}
\usepackage{etoolbox}
\usepackage{ragged2e}
\usepackage{paracol}
\usepackage{multicol}
\usepackage{needspace}
\usepackage{fancyhdr}
\usepackage{tikz}
\usepackage{zref-abspage}
\tcbuselibrary{breakable,skins}
\usetikzlibrary{calc}
\usepackage{pgfcalendar}

\hypersetup{
  unicode=true,
  bookmarksopen=true,
  bookmarksnumbered=false,
  hidelinks
}

% ===== 自動推導，不要手動改下面 =====
% （\docmode 與 \bilingualpaper 由各文件在 \input{bverfge_layout} 之前設定；
%   預設 chinese / a4）
\ifdefined\docmode\else\def\docmode{chinese}\fi
\ifdefined\bilingualpaper\else\def\bilingualpaper{a4}\fi
\newif\ifshoworiginal
\newif\ifshowtranslation
\newif\ifbilinguallandscape
\newif\ifbilingualcolumns
\newif\ifintranslation
\newif\iffirstbilingualsection

\showoriginalfalse
\showtranslationfalse
\bilinguallandscapefalse
\bilingualcolumnsfalse
\intranslationfalse
\firstbilingualsectionfalse

\def\modechinese{chinese}
\def\modegerman{german}
\def\modebilingual{bilingual}
\def\bilingualpaperaThree{a3}
\def\bilingualpaperaFour{a4}

\ifx\docmode\modechinese
  \showtranslationtrue
\fi

\ifx\docmode\modegerman
  \showoriginaltrue
\fi

\ifx\docmode\modebilingual
  \showoriginaltrue
  \showtranslationtrue
  \bilinguallandscapetrue
  \bilingualcolumnstrue
\fi

% 編排模式說明：
% chinese  → \showoriginalfalse  \showtranslationtrue  （A4 直式，純中文）
% german   → \showoriginaltrue   \showtranslationfalse （A4 直式，純德文）
% bilingual→ \showoriginaltrue   \showtranslationtrue  \bilingualcolumnstrue（A3/A4 橫式對照）

\ifbilingualcolumns
  \ifx\bilingualpaper\bilingualpaperaFour
    \geometry{
      a4paper,
      landscape,
      left=1.25cm,
      right=2.25cm,
      top=1.25cm,
      bottom=1.75cm,
      footskip=8.5mm,
      marginparwidth=0.75cm,
      marginparsep=0.25cm
    }
  \else
    \geometry{
      a3paper,
      landscape,
      left=1.55cm,
      right=2.75cm,
      top=1.55cm,
      bottom=2.05cm,
      footskip=9.5mm,
      marginparwidth=0.9cm,
      marginparsep=0.35cm
    }
  \fi
\else
\geometry{
  a4paper,
  portrait,
  left=2.2cm,
  right=2.6cm,
  top=2.0cm,
  bottom=2.0cm,
  marginparwidth=0.8cm,
  marginparsep=0.3cm
}
\fi
% ===== 時間戳基礎設施 =====
\newcount\stamptimeval
\newcount\stampminuteval
\newcommand{\stamphh}{%
  \stamptimeval=\time
  \divide\stamptimeval by 60
  \ifnum\stamptimeval<10 0\fi
  \the\stamptimeval}
\newcommand{\stampmm}{%
  \stamptimeval=\time
  \stampminuteval=\time
  \divide\stampminuteval by 60
  \multiply\stampminuteval by 60
  \advance\stamptimeval by -\stampminuteval
  \ifnum\stamptimeval<10 0\fi
  \the\stamptimeval}
\newcommand{\stampwkde}[1]{%
  \ifcase#1Montag\or Dienstag\or Mittwoch\or Donnerstag\or Freitag\or Samstag\or Sonntag\fi}
\newcommand{\stampwkzh}[1]{%
  \ifcase#1 一\or 二\or 三\or 四\or 五\or 六\or 日\fi}
\newcommand{\stampmonthde}[1]{%
  \ifcase#1\or Januar\or Februar\or März\or April\or Mai\or Juni\or
  Juli\or August\or September\or Oktober\or November\or Dezember\fi}
\newcount\stampjdval
\newcount\stampwdval
\pgfcalendardatetojulian{\the\year-\the\month-\the\day}{\stampjdval}
\pgfcalendarjuliantoweekday{\the\stampjdval}{\stampwdval}
\newcommand{\timestampzh}{%
  \songfont 最後修改：\the\year 年\the\month 月\the\day 日%
  % （星期\stampwkzh{\the\stampwdval}）%
  % \space\stamphh:\stampmm
  }

\newcommand{\documentdatezh}{\the\year 年 \the\month 月 \the\day 日}
\newcommand{\documentdatede}{\the\day. \stampmonthde{\the\month} \the\year}
\newcommand{\bverfgetranslationnote}{%
  {\songfont 翻譯說明：初譯使用 Claude Code 與 GPT 協助迭代；仍請以德文原文為準。}%
}
\newcommand{\bverfgecourseinfo}{%
  {\songfont 課程：114 學年度第 2 學期；憲法專題研究（四）；授課老師：湯德宗教授；導讀日期：115 年 6 月 1 日。}%
}
\newcommand{\bverfgeeditorinfo}{%
  {\songfont 編者：王逸帆（R10A21126），國立台灣大學法律系研究所碩士班。}%
}
\newcommand{\bverfgetitleversionline}{%
  \ifbilingualcolumns
    Fassung \documentversion\enspace\textperiodcentered\enspace Stand: \documentdatede
    \quad
    {\songfont 版本 \documentversion\enspace\textperiodcentered\enspace 修訂日期：\documentdatezh\par}%
  \else
    \ifshowtranslation
      {\songfont 版本 \documentversion\enspace\textperiodcentered\enspace 修訂日期：\documentdatezh\par}%
    \else
      Fassung \documentversion\enspace\textperiodcentered\enspace Stand: \documentdatede\par
    \fi
  \fi
}
\newcommand{\bverfgetitlecornerinfo}{%
  \ifshowtranslation
    \vspace{0.72em}%
    \noindent\hfill
    \begin{minipage}{0.66\textwidth}
      \raggedleft\tiny\color{pendingink}\songfont
      \bverfgecourseinfo\par
      \bverfgeeditorinfo
    \end{minipage}%
  \fi
}
\newcommand{\bverfgetitlemeta}{%
  \vfill
  \begin{center}%
    \footnotesize\color{pendingink}%
    \bverfgetitleversionline
    \ifshowtranslation
      \vspace{0.28em}{\scriptsize\bverfgetranslationnote\par}%
    \fi
  \end{center}%
  \bverfgetitlecornerinfo
}

\pagestyle{fancy}
\fancyhf{}
\renewcommand{\headrulewidth}{0pt}
\renewcommand{\footrulewidth}{0pt}
\fancyfoot[C]{\thepage}
\ifbilingualcolumns
  \fancyfoot[R]{\scriptsize\color{pendingink}\timestampzh}%
\else
  \ifshoworiginal
  \else
    \fancyfoot[R]{\scriptsize\color{pendingink}\timestampzh\hspace{0.45cm}}%
  \fi
\fi
\setmainfont{Times New Roman}
\setsansfont{Helvetica Neue}
\setmonofont{Menlo}
\IfFileExists{/Users/iw/Library/Fonts/kaiu.ttf}{
  \setCJKmainfont[Path=/Users/iw/Library/Fonts/,AutoFakeBold=4]{kaiu.ttf}
}{
  \IfFontExistsTF{標楷體}{
    \setCJKmainfont[AutoFakeBold=4]{標楷體}
  }{
    \setCJKmainfont{Songti TC}
  }
}
\IfFontExistsTF{Songti TC}{
  \setCJKfamilyfont{song}{Songti TC}
  \setCJKmonofont{Songti TC}
}{
  \setCJKfamilyfont{song}[Path=/Users/iw/Library/Fonts/,AutoFakeBold=4]{kaiu.ttf}
  \setCJKmonofont[Path=/Users/iw/Library/Fonts/,AutoFakeBold=4]{kaiu.ttf}
}
\newcommand{\songfont}{\CJKfamily{song}}

% ===== 封面／目錄專用打字機字體（僅限封面、目錄頁使用，不影響正文）=====
\IfFontExistsTF{Radio Newsman}{%
  \newfontfamily\bverfgetitlede{Radio Newsman}%
}{%
  \newfontfamily\bverfgetitlede{Courier New}%
}
\IfFontExistsTF{Erikas Farbband}{%
  \newfontfamily\bverfgebodyde{Erikas Farbband}%
}{%
  \let\bverfgebodyde\bverfgetitlede
}
\IfFontExistsTF{Maoken Heavy Labourer}{%
  \setCJKfamilyfont{bverfgetitlecjk}{Maoken Heavy Labourer}%
}{%
  \setCJKfamilyfont{bverfgetitlecjk}{Songti TC}%
}
\IfFontExistsTF{Huiwen-mincho}{%
  \setCJKfamilyfont{bverfgebodycjk}{Huiwen-mincho}%
}{%
  \setCJKfamilyfont{bverfgebodycjk}{Songti TC}%
}
\newcommand{\bverfgetitlezh}{\bverfgetitlede\CJKfamily{bverfgetitlecjk}}
\newcommand{\bverfgebodyzh}{\bverfgebodyde\CJKfamily{bverfgebodycjk}}

\setstretch{1.35}
\setlist{itemsep=0.2em, topsep=0.4em}
\setlength{\parindent}{0pt}
\setlength{\parskip}{0.45em}
\raggedbottom
\setcounter{secnumdepth}{0}
\setcounter{tocdepth}{2}

\newcommand{\lawboxsourcefont}{\footnotesize}
\newcommand{\lawboxtransfont}{\footnotesize}
\newcommand{\lawboxsourcestretch}{1.02}
\newcommand{\lawboxtransstretch}{0.92}

\titleformat{\section}{\centering\Large\bfseries\songfont}{\thesection}{1em}{}
\titleformat{\subsection}{\large\bfseries\songfont}{\thesubsection}{1em}{}
\titleformat{\subsubsection}{\normalsize\bfseries\songfont}{\thesubsubsection}{1em}{}
\titleformat{\paragraph}{\normalsize\bfseries\songfont}{\theparagraph}{1em}{}
\titlespacing*{\section}{0pt}{1.8em}{1.0em}
\titlespacing*{\subsection}{0pt}{1.2em}{0.6em}
\titlespacing*{\subsubsection}{0pt}{1.0em}{0.45em}
\titlespacing*{\paragraph}{0pt}{0.6em}{0.4em}

\definecolor{noteblue}{HTML}{A9C9F5}
\definecolor{noteteal}{HTML}{AEE1D6}
\definecolor{notegray}{HTML}{D7DADF}
\definecolor{caseink}{HTML}{000000}
\definecolor{casepaper}{HTML}{FCFEFD}
\definecolor{sourcepaper}{HTML}{FAFBF8}
\definecolor{sourceline}{HTML}{D7DED8}
\definecolor{sourceink}{HTML}{000000}
\definecolor{pendingink}{HTML}{000000}

\tcbset{
  caseboxbase/.style={
    enhanced,
    breakable,
    width=0.985\linewidth,
    colback=casepaper,
    coltitle=caseink,
    fonttitle=\bfseries\songfont,
    boxrule=0.28pt,
    arc=1.2mm,
    left=2.4mm,
    right=2.4mm,
    top=1.5mm,
    bottom=1.5mm,
    boxsep=1.5mm,
    before skip=0.65em,
    after skip=0.65em,
    before upper={\setlength{\parindent}{0pt}\setlength{\parskip}{0.35em}},
  }
}

\newcommand{\bilingualstart}{}
\newcommand{\bilingualstop}{}
\ifbilingualcolumns
  \renewcommand{\bilingualstart}{\small\setlength{\columnsep}{1.25cm}\setlength{\columnseprule}{0.12pt}\columnratio{0.5,0.5}\multicolumnfootnotes\flushbottom\sloppy\interlinepenalty=80\clubpenalty=800\widowpenalty=800\displaywidowpenalty=800\begin{paracol}{2}\global\intranslationfalse\global\firstbilingualsectiontrue{\footnotesize\color{sourceink}Deutsch\par}\switchcolumn{\footnotesize\songfont\color{sourceink}中文譯文\par}\switchcolumn*}
  \renewcommand{\bilingualstop}{\ifintranslation\switchcolumn*\global\intranslationfalse\fi\end{paracol}\clearpage\raggedbottom}
\fi
\newif\ifrandnumbers
\randnumbersfalse
\newcounter{randnumber}
\newcounter{randnumberlocal}
\newcounter{lastsourcerandstart}
\newif\iflastsourcehasrandnumbers
\lastsourcehasrandnumbersfalse
\newif\ifinlawbox
\inlawboxfalse
\newif\iflawboxhaspendingrn
\lawboxhaspendingrnfalse
\newcommand{\lawboxpendingrn}{}
\newcommand{\startrandnumbers}{\global\randnumberstrue\setcounter{randnumber}{1}\global\lastsourcehasrandnumbersfalse}
\newcommand{\stoprandnumbers}{\global\randnumbersfalse\global\lastsourcehasrandnumbersfalse}
\newlength{\randnumberwidth}
\setlength{\randnumberwidth}{0.95cm}
\newlength{\randnumberrightoffset}
\setlength{\randnumberrightoffset}{0.85cm}
\newlength{\randnumberlawboxrightoffset}
\setlength{\randnumberlawboxrightoffset}{1.40cm}
\newcommand{\randnumberbox}[1]{%
  \leavevmode
  \makebox[0pt][l]{%
    \smash{%
      \tikz[remember picture,overlay,baseline=(rnmark.base)]{%
        \coordinate (rnmark) at (0,0);%
        \ifinlawbox
          \path let \p1=(current page.east), \p2=(rnmark.base) in
            node[
              anchor=base east,
              inner sep=0pt,
              outer sep=0pt,
              text width=\randnumberwidth,
              align=right,
              text=pendingink,
              font=\normalfont\mdseries\scriptsize\ttfamily
            ] at (\x1-\randnumberlawboxrightoffset,\y2) {{\normalfont\mdseries\scriptsize\ttfamily #1}};%
        \else
          \path let \p1=(current page.east), \p2=(rnmark.base) in
            node[
              anchor=base east,
              inner sep=0pt,
              outer sep=0pt,
              text width=\randnumberwidth,
              align=right,
              text=pendingink,
              font=\normalfont\mdseries\scriptsize\ttfamily
            ] at (\x1-\randnumberrightoffset,\y2) {{\normalfont\mdseries\scriptsize\ttfamily #1}};%
        \fi
      }%
    }%
  }%
  \ignorespaces
}
\newcommand{\lawboxstorerandnumber}[1]{%
  \gdef\lawboxpendingrn{#1}%
  \global\lawboxhaspendingrntrue
  \ignorespaces
}
\newcommand{\lawboxuserandnumber}{%
  \iflawboxhaspendingrn
    \randnumberbox{\lawboxpendingrn}%
    \global\lawboxhaspendingrnfalse
  \fi
}
\newcommand{\sourcern}[1]{%
  \ifrandnumbers
    \ifshoworiginal
      \ifshowtranslation\else
        \ifinlawbox\lawboxstorerandnumber{#1}\else\randnumberbox{#1}\fi
      \fi
    \fi
  \fi
}
\newcommand{\transrn}[1]{%
  \ifrandnumbers
    \ifshowtranslation
      \ifinlawbox\lawboxstorerandnumber{#1}\else\randnumberbox{#1}\fi
    \fi
  \fi
}
% ===== 課堂模式：德文原文縮寫註解 =====
% 日常切換只改各判決檔的一行：
%   \classroommodeon        開啟課堂版；註解掉這行即回到一般版。
% 模板預設：
%   1. 課堂模式關閉。
%   2. 課堂模式開啟時，縮寫註解包含「德文全稱＋中文譯名」。
%   3. 課堂模式開啟時，GG/條文編者註仍會自動顯示。
% 進階微調才需要在單案檔覆寫：
%   \classroomabbrzhoff     縮寫註解僅顯示德文全稱。
%   \showeditornotestrue    一般版也顯示編者註。
% 註解策略：同一實際 PDF 頁面中，同一縮寫只在第一次出現處加註。
% 這裡用 zref-abspage 的絕對頁序，不用 \thepage，避免文件內頁碼重置時誤判。
\newif\ifclassroommode
\newif\ifclassroomabbrzh
\classroommodefalse
\classroomabbrzhtrue
\newcommand{\classroommodeon}{\classroommodetrue}
\newcommand{\classroommodeoff}{\classroommodefalse}
\newcommand{\classroomabbrzhon}{\classroomabbrzhtrue}
\newcommand{\classroomabbrzhoff}{\classroomabbrzhfalse}
\newcommand{\abbrnotelabel}{\emph{Abk.}: }
\newcommand{\abbrnote}[3]{%
  \ifclassroommode
    \ifshoworiginal
      \ifintranslation\else
        \footnote{\normalfont\footnotesize\abbrnotelabel #2\ifclassroomabbrzh；{\songfont #3}\fi}%
      \fi
    \fi
  \fi
}
\newcommand{\abbrnoteonce}[4]{%
  #2%
  \ifclassroommode
    \ifshoworiginal
      \ifintranslation\else
        \begingroup
          \edef\abbrcurrentabspage{\number\numexpr\value{abspage}+1\relax}%
          \ifcsname abbrnoteseen@#1@\abbrcurrentabspage\endcsname\else
            \global\expandafter\let\csname abbrnoteseen@#1@\abbrcurrentabspage\endcsname\empty
            \abbrnote{#1}{#3}{#4}%
          \fi
        \endgroup
      \fi
    \fi
  \fi
}
\newcommand{\abbrAbs}{\abbrnoteonce{abs}{Abs.}{Absatz}{項}}
\newcommand{\abbrAAO}{\abbrnoteonce{aao}{a.a.O.}{am angegebenen Ort}{前引處}}
\newcommand{\abbrAE}{\abbrnoteonce{ae}{AE}{Alternativ-Entwurf}{替代草案}}
\newcommand{\abbrArt}{\abbrnoteonce{art}{Art.}{Artikel}{條}}
\newcommand{\abbrAufl}{\abbrnoteonce{aufl}{Aufl.}{Auflage}{版}}
\newcommand{\abbrBd}{\abbrnoteonce{bd}{Bd.}{Band}{卷}}
\newcommand{\abbrBGBl}{\abbrnoteonce{bgbl}{BGBl.}{Bundesgesetzblatt}{聯邦法律公報}}
\newcommand{\abbrBGHSt}{\abbrnoteonce{bghst}{BGHSt}{Entscheidungen des Bundesgerichtshofes in Strafsachen}{聯邦普通法院刑事判決彙編}}
\newcommand{\abbrBRDrucks}{\abbrnoteonce{brdrucks}{BRDrucks.}{Bundesrats-Drucksache}{聯邦參議院文件}}
\newcommand{\abbrBTDrucks}{\abbrnoteonce{btdrucks}{BTDrucks.}{Bundestags-Drucksache}{聯邦議院文件}}
\newcommand{\abbrBundesgesetzbl}{\abbrnoteonce{bundesgesetzbl}{Bundesgesetzbl.}{Bundesgesetzblatt}{聯邦法律公報}}
\newcommand{\abbrBvF}{\abbrnoteonce{bvf}{BvF}{Bundesverfassungsgerichtsverfahren}{聯邦憲法法院程序案號}}
\newcommand{\abbrBvL}{\abbrnoteonce{bvl}{BvL}{Bundesverfassungsgerichtsverfahren}{聯邦憲法法院程序案號}}
\newcommand{\abbrBvR}{\abbrnoteonce{bvr}{BvR}{Bundesverfassungsgerichtsverfahren}{聯邦憲法法院程序案號}}
\newcommand{\abbrBVerfG}{\abbrnoteonce{bverfg}{BVerfG}{Bundesverfassungsgericht}{聯邦憲法法院}}
\newcommand{\abbrBVerfGE}{\abbrnoteonce{bverfge}{BVerfGE}{Entscheidungen des Bundesverfassungsgerichts}{聯邦憲法法院判決集}}
\newcommand{\abbrBVerfGG}{\abbrnoteonce{bverfgg}{BVerfGG}{Bundesverfassungsgerichtsgesetz}{聯邦憲法法院法}}
\newcommand{\abbrDDR}{\abbrnoteonce{ddr}{DDR}{Deutsche Demokratische Republik}{德意志民主共和國；東德}}
\newcommand{\abbrDJT}{\abbrnoteonce{djt}{DJT}{Deutscher Juristentag}{德國法學家大會}}
\newcommand{\abbrDr}{\abbrnoteonce{dr}{Dr.}{Doktor}{Dr.}}
\newcommand{\abbrEuGRZ}{\abbrnoteonce{eugrz}{EuGRZ}{Europäische Grundrechte-Zeitschrift}{歐洲基本權期刊}}
\newcommand{\abbrEV}{\abbrnoteonce{ev}{EV}{Einigungsvertrag}{統一條約}}
\newcommand{\abbrF}{\abbrnoteonce{f}{f.}{folgende}{以下一頁；下一段}}
\newcommand{\abbrFF}{\abbrnoteonce{ff}{ff.}{fortfolgende}{以下數頁；以下各段}}
\newcommand{\abbrGBlDDR}{\abbrnoteonce{gblddr}{GBl. DDR}{Gesetzblatt der Deutschen Demokratischen Republik}{德意志民主共和國法律公報}}
\newcommand{\abbrGG}{\abbrnoteonce{gg}{GG}{Grundgesetz}{基本法}}
\newcommand{\abbrGVBl}{\abbrnoteonce{gvbl}{GVBl.}{Gesetz- und Verordnungsblatt}{法律及命令公報}}
\newcommand{\abbrJR}{\abbrnoteonce{jr}{JR}{Juristische Rundschau}{法學評論}}
\newcommand{\abbrJZ}{\abbrnoteonce{jz}{JZ}{JuristenZeitung}{法學家報}}
\newcommand{\abbrIVm}{\abbrnoteonce{ivm}{i.V.m.}{in Verbindung mit}{結合；連同}}
\newcommand{\abbrMdB}{\abbrnoteonce{mdb}{MdB}{Mitglied des Bundestages}{聯邦議院議員}}
\newcommand{\abbrMwN}{\abbrnoteonce{mwn}{m.w.N.}{mit weiteren Nachweisen}{附進一步文獻／判例出處}}
\newcommand{\abbrNF}{\abbrnoteonce{nf}{n.F.}{neue Fassung}{新版本}}
\newcommand{\abbrAF}{\abbrnoteonce{af}{a.F.}{alte Fassung}{舊版本}}
\newcommand{\abbrNr}{\abbrnoteonce{nr}{Nr.}{Nummer}{款、目或號，依語境}}
\newcommand{\abbrProf}{\abbrnoteonce{prof}{Prof.}{Professor}{教授}}
\newcommand{\abbrRdnr}{\abbrnoteonce{rdnr}{Rdnr.}{Randnummer}{邊碼；段碼}}
\newcommand{\abbrRdnrn}{\abbrnoteonce{rdnrn}{Rdnrn.}{Randnummern}{邊碼；段碼}}
\newcommand{\abbrRGBl}{\abbrnoteonce{rgbl}{RGBl.}{Reichsgesetzblatt}{帝國法律公報}}
\newcommand{\abbrRGSt}{\abbrnoteonce{rgst}{RGSt}{Entscheidungen des Reichsgerichts in Strafsachen}{帝國法院刑事判決彙編}}
\newcommand{\abbrRspr}{\abbrnoteonce{rspr}{Rspr.}{Rechtsprechung}{判例；司法實務}}
\newcommand{\abbrRVO}{\abbrnoteonce{rvo}{RVO}{Reichsversicherungsordnung}{帝國保險法}}
\newcommand{\abbrS}{\abbrnoteonce{s}{S.}{Seite}{頁}}
\newcommand{\abbrSFHG}{\abbrnoteonce{sfhg}{SFHG}{Schwangeren- und Familienhilfegesetz}{孕婦及家庭扶助法}}
\newcommand{\abbrSGBV}{\abbrnoteonce{sgbv}{SGB V}{Fünftes Buch Sozialgesetzbuch}{社會法典第五編}}
\newcommand{\abbrStAG}{\abbrnoteonce{staeg}{StÄG}{Strafrechtsänderungsgesetz}{刑法修正法}}
\newcommand{\abbrStenBer}{\abbrnoteonce{stenber}{StenBer.}{Stenographischer Bericht}{速記紀錄}}
\newcommand{\abbrStGB}{\abbrnoteonce{stgb}{StGB}{Strafgesetzbuch}{刑法}}
\newcommand{\abbrStREG}{\abbrnoteonce{streg}{StREG}{Strafrechtsreform-Ergänzungsgesetz}{刑法改革補充法}}
\newcommand{\abbrStrRG}{\abbrnoteonce{strrg}{StrRG}{Strafrechtsreformgesetz}{刑法改革法}}
\newcommand{\abbrUS}{\abbrnoteonce{us}{U.S.}{United States Reports}{美國最高法院判決彙編}}
\newcommand{\abbrVgl}{\abbrnoteonce{vgl}{vgl.}{vergleiche}{參見}}
\newcommand{\abbrVVDStRL}{\abbrnoteonce{vvdstrl}{VVDStRL}{Veröffentlichungen der Vereinigung der Deutschen Staatsrechtslehrer}{德國公法教師協會論文集}}
\newcommand{\abbrWP}{\abbrnoteonce{wp}{WP.}{Wahlperiode}{屆期}}
\newcommand{\abbrWp}{\abbrnoteonce{wp}{Wp.}{Wahlperiode}{屆期}}
\newcommand{\abbrZStrW}{\abbrnoteonce{zstrw}{ZStrW}{Zeitschrift für die gesamte Strafrechtswissenschaft}{整體刑法學期刊}}
\newcommand{\recordsourcerandstart}{%
  \ifrandnumbers
    \setcounter{lastsourcerandstart}{\value{randnumber}}%
    \global\lastsourcehasrandnumberstrue
  \else
    \global\lastsourcehasrandnumbersfalse
  \fi
}
\newlength{\biparagraphskip}
\setlength{\biparagraphskip}{0.52em}
\newlength{\lawboxblockbeforeskip}
\setlength{\lawboxblockbeforeskip}{0.72em}
\newlength{\lawboxblockafterskip}
\setlength{\lawboxblockafterskip}{0.92em}
\newif\iflawboxpartendedblock
\lawboxpartendedblockfalse
\newcommand{\sourceparstyle}{\footnotesize\color{sourceink}\setstretch{1.12}\RaggedRight\emergencystretch=3em\setlength{\parskip}{0.42em}\obeylines\selectfont}
\newcommand{\transparstyle}{\songfont\color{caseink}\setstretch{1.12}\setlength{\parindent}{0pt}\setlength{\parskip}{0.42em}}
\newcommand{\bilingualpairskip}{%
  \par\addvspace{\biparagraphskip}%
  \switchcolumn
  \par\addvspace{\biparagraphskip}%
  \switchcolumn*%
  \global\intranslationfalse
}
\newcommand{\bilingualboxpairskip}{%
  \par
  \switchcolumn
  \par
  \switchcolumn*%
  \global\intranslationfalse
}
\newcommand{\bilinguallawboxblockskip}{%
  \switchcolumn*%
  \par\addvspace{\lawboxblockafterskip}%
  \switchcolumn
  \par\addvspace{\lawboxblockafterskip}%
  \switchcolumn*%
  \global\intranslationfalse
}
\newcommand{\bikeepwithnext}[1][4]{%
  \ifbilingualcolumns
    \syncleft
    \Needspace{#1\baselineskip}%
    \switchcolumn
    \Needspace{#1\baselineskip}%
    \switchcolumn*%
    \global\intranslationfalse
  \else
    \Needspace{#1\baselineskip}%
  \fi
}
\NewEnviron{sourcepara}[1][]{
  \recordsourcerandstart
  \ifshoworiginal
    \ifbilingualcolumns
      \ifintranslation\switchcolumn*\global\intranslationfalse\fi
    \else
      \Needspace{5\baselineskip}
    \fi
    \ifbilingualcolumns\else\par\addvspace{0.35em}\fi
    {\sourceparstyle
    \noindent\BODY\par}
    \ifbilingualcolumns\else\addvspace{0.25em}\fi
    \ifbilingualcolumns
      \switchcolumn
      \global\intranslationtrue
    \fi
  \else
    \ifrandnumbers
      \begingroup
        \setbox0=\vbox{\hsize=\linewidth\sourceparstyle\noindent\BODY\par}%
      \endgroup
    \fi
  \fi
}
\NewEnviron{transpara}{
  \ifshowtranslation
    \setcounter{randnumberlocal}{\value{lastsourcerandstart}}%
    \ifbilingualcolumns
      \ifintranslation\else\switchcolumn\global\intranslationtrue\fi
      {\transparstyle\setlength{\leftskip}{0.55em}\setlength{\rightskip}{0.15em plus 1em}\addtolength{\linewidth}{-0.55em}\BODY\par}
      \switchcolumn*
      \bilingualpairskip
    \else
      {\transparstyle\BODY\par}
    \fi
  \else
    \ifbilingualcolumns
      \ifintranslation\switchcolumn*\global\intranslationfalse\fi
    \fi
  \fi
}
\NewEnviron{sourceboxpara}[1][]{
  \global\lawboxpartendedblockfalse
  \recordsourcerandstart
  \ifshoworiginal
    \ifbilingualcolumns
      \ifintranslation\switchcolumn*\global\intranslationfalse\fi
    \else
      \Needspace{3\baselineskip}
    \fi
    {\sourceparstyle
    \BODY\par}
    \ifbilingualcolumns\else
      \iflawboxpartendedblock\addvspace{\lawboxblockafterskip}\fi
    \fi
    \ifbilingualcolumns
      \switchcolumn
      \global\intranslationtrue
    \fi
  \fi
}
\NewEnviron{transboxpara}{
  \global\lawboxpartendedblockfalse
  \ifshowtranslation
    \setcounter{randnumberlocal}{\value{lastsourcerandstart}}%
    \ifbilingualcolumns
      \ifintranslation\else\switchcolumn\global\intranslationtrue\fi
      {\transparstyle\BODY\par}
      \iflawboxpartendedblock
        \bilinguallawboxblockskip
      \else
        \switchcolumn*
        \bilingualboxpairskip
      \fi
    \else
      {\transparstyle\BODY\par}
      \iflawboxpartendedblock\addvspace{\lawboxblockafterskip}\fi
    \fi
  \else
    \ifbilingualcolumns
      \ifintranslation\switchcolumn*\global\intranslationfalse\fi
    \fi
  \fi
}
\newcommand{\leitsatznum}[1]{\noindent\hangindent=3.0em\hangafter=1\makebox[2.25em][r]{\bfseries #1.}\hspace{0.55em}\ignorespaces}
\newcommand{\syncleft}{\ifbilingualcolumns\ifintranslation\switchcolumn*\global\intranslationfalse\fi\fi}
\pretocmd{\section}{\syncleft\clearpage}{}{}
\pretocmd{\subsection}{\syncleft}{}{}
\pretocmd{\subsubsection}{\syncleft}{}{}
\newcommand{\biheadingfont}{\songfont}
\newcommand{\bitocentry}[2]{%
  \ifshoworiginal
    \ifshowtranslation
      \ifstrequal{#1}{#2}{#1}{#1\space #2}%
    \else
      #1%
    \fi
  \else
    #2%
  \fi
}
\pdfstringdefDisableCommands{%
  \def\bitocentry#1#2{#1}%
}
\newcommand{\biheading}[7]{%
  \syncleft#1\bikeepwithnext[#3]\phantomsection\addcontentsline{toc}{#2}{\bitocentry{#6}{#7}}%
  \ifbilingualcolumns
    {#4 #6\par}%
    \switchcolumn
    {#4\biheadingfont #7\par}%
    \switchcolumn*%
    \global\intranslationfalse
    \par\addvspace{#5}%
    \switchcolumn
    \par\addvspace{#5}%
    \switchcolumn*%
  \else
    \ifshoworiginal{#4 #6\par}\fi
    \ifshowtranslation{#4\biheadingfont #7\par}\fi
    \addvspace{#5}%
  \fi
}
\newcommand{\termsectionheading}[1]{%
  \par\Needspace{9\baselineskip}%
  \vspace{0.65em}%
  {\songfont\bfseries #1\par}%
  \vspace{0.2em}%
}
% 章節換頁：
%   雙語 paracol 欄內不要用 \clearpage；它會在同步兩欄時吐出只有欄線/頁碼的空頁。
%   \newpage 足以讓下一個 section 起新頁，又不會強制清空所有待排材料。
%   單語模式不在 paracol 內，仍可用 \clearpage。
\newcommand{\bisectionpagebreak}{\ifbilingualcolumns\newpage\else\clearpage\fi}
\newcommand{\bisection}[2]{%
  \biheading{\iffirstbilingualsection\global\firstbilingualsectionfalse\else\bisectionpagebreak\fi}{section}{6}{\centering\Large\bfseries}{0.8em}{#1}{#2}%
}
\newcommand{\bisubsection}[2]{%
  \biheading{}{subsection}{5}{\large\bfseries}{0.55em}{#1}{#2}%
}
\newcommand{\bisubsubsection}[2]{%
  \biheading{}{subsubsection}{4}{\normalsize\bfseries}{0.45em}{#1}{#2}%
}
\newcommand{\bicompositeheading}[4]{%
  \biheading{}{subsection}{5}{\large\bfseries}{0.16em}{#1}{#2}%
  \biheading{}{subsubsection}{3}{\normalsize\bfseries}{0.45em}{#3}{#4}%
}
\newif\iflawboxfirstitem
\lawboxfirstitemfalse
\newcommand{\lawboxprespace}[1]{%
  \par
  \iflawboxfirstitem
    \global\lawboxfirstitemfalse
  \else
    \addvspace{#1}%
  \fi
}
\newtcolorbox{lawboxinner}{
  caseboxbase,
  colframe=sourceline!85!black,
  colback=white,
  left=1.05em,
  right=1.05em,
  top=0.62em,
  bottom=0.62em,
  before upper={\global\inlawboxtrue\global\lawboxfirstitemtrue\global\lawboxhaspendingrnfalse\ifintranslation\lawboxtransfont\else\lawboxsourcefont\fi\RaggedRight\setlength{\parindent}{0pt}\setlength{\parskip}{0pt}\ifintranslation\setstretch{\lawboxtransstretch}\else\setstretch{\lawboxsourcestretch}\fi},
  after upper={\global\lawboxhaspendingrnfalse\global\inlawboxfalse}
}
\tcbset{
  lawboxpartbase/.style={
    enhanced,
    breakable=false,
    width=0.985\linewidth,
    colback=white,
    frame hidden,
    boxrule=0pt,
    arc=0pt,
    left=1.05em,
    right=1.05em,
    boxsep=1.5mm,
    before skip=0pt,
    after skip=0pt,
    before upper={\global\inlawboxtrue\global\lawboxfirstitemtrue\global\lawboxhaspendingrnfalse\ifintranslation\lawboxtransfont\else\lawboxsourcefont\fi\RaggedRight\setlength{\parindent}{0pt}\setlength{\parskip}{0pt}\ifintranslation\setstretch{\lawboxtransstretch}\else\setstretch{\lawboxsourcestretch}\fi},
    after upper={\global\lawboxhaspendingrnfalse\global\inlawboxfalse},
    borderline west={0.28pt}{0pt}{sourceline!85!black},
    borderline east={0.28pt}{0pt}{sourceline!85!black},
  },
  lawboxpartfirst/.style={
    lawboxpartbase,
    top=0.62em,
    bottom=0.04em,
    before skip=\lawboxblockbeforeskip,
    borderline north={0.28pt}{0pt}{sourceline!85!black},
  },
  lawboxpartmiddle/.style={
    lawboxpartbase,
    top=0.04em,
    bottom=0.04em,
  },
  lawboxpartlast/.style={
    lawboxpartbase,
    top=0.04em,
    bottom=0.62em,
    after skip=0pt,
    borderline south={0.28pt}{0pt}{sourceline!85!black},
  }
}
\newtcolorbox{lawboxpartinner}[1]{#1}
\newtcolorbox{termboxinner}{
  caseboxbase,
  colframe=noteblue!70!black,
  colback=casepaper,
  before upper={\songfont\setlength{\parindent}{0pt}\setlength{\parskip}{0.35em}}
}
\newtcolorbox{deboxinner}{
  caseboxbase,
  colframe=notegray!72!black,
  colback=white,
  left=1.05em,
  right=1.05em,
  top=0.62em,
  bottom=0.62em,
  before upper={\global\inlawboxtrue\global\lawboxfirstitemtrue\global\lawboxhaspendingrnfalse\small\setlength{\parindent}{0pt}\setlength{\parskip}{0.35em}},
  after upper={\global\lawboxhaspendingrnfalse\global\inlawboxfalse}
}
\NewEnviron{lawbox}{
  \begin{lawboxinner}\BODY\end{lawboxinner}
}
\NewEnviron{lawboxpart}[2]{
  \begin{lawboxpartinner}{lawboxpart#1,equal height group=#2}\BODY\end{lawboxpartinner}%
  \ifstrequal{#1}{last}{\global\lawboxpartendedblocktrue}{}%
}
\NewEnviron{termbox}{
  \par\addvspace{0.55em}
  {\color{noteblue!70!black}\rule{\linewidth}{0.28pt}}\par
  \begingroup
    \songfont\small\color{caseink}
    \setlength{\parindent}{0pt}
    \setlength{\parskip}{0.24em}
    \BODY
    \par
  \endgroup
  {\color{noteblue!70!black}\rule{\linewidth}{0.28pt}}\par
  \addvspace{0.55em}
}
\NewEnviron{debox}{
  \begin{deboxinner}\BODY\end{deboxinner}
}
\providecommand{\paragraphmark}[1]{}
\newcommand{\lawboxtitleafter}{\addvspace{0.06em}}
\newcommand{\lawtitle}[1]{\lawboxprespace{0.22em}{\lawboxtransfont\bfseries\lawboxuserandnumber #1}\par\lawboxtitleafter}
\newcommand{\absno}[2]{\lawboxprespace{0.04em}\noindent\lawboxuserandnumber\hangindent=2.6em\hangafter=1\makebox[2.6em][l]{(#1)}#2\par}
\newcommand{\nrno}[2]{\lawboxprespace{0pt}\noindent\lawboxuserandnumber\hspace*{2.6em}\hangindent=4.4em\hangafter=1\makebox[1.8em][l]{#1.}#2\par}
\newcommand{\letterno}[2]{\lawboxprespace{0pt}\noindent\lawboxuserandnumber\hspace*{4.4em}\hangindent=6.0em\hangafter=1\makebox[1.6em][l]{#1)}#2\par}
\newcommand{\sourcelawtitle}[1]{\lawboxprespace{0.22em}{\lawboxsourcefont\bfseries\lawboxuserandnumber #1}\par\lawboxtitleafter}
\newcommand{\sourcelawsubtitle}[1]{\lawboxprespace{0.04em}{\lawboxsourcefont\bfseries\lawboxuserandnumber #1}\par\lawboxtitleafter}
\newcommand{\sourceabsno}[2]{\lawboxprespace{0.04em}\noindent\lawboxuserandnumber\hangindent=2.45em\hangafter=1\makebox[2.45em][l]{(#1)}#2\par}
\newcommand{\sourcenrno}[2]{\lawboxprespace{0pt}\noindent\lawboxuserandnumber\hspace*{2.45em}\hangindent=4.25em\hangafter=1\makebox[1.8em][l]{#1.}#2\par}
\newcommand{\sourceletterno}[2]{\lawboxprespace{0pt}\noindent\lawboxuserandnumber\hspace*{4.25em}\hangindent=5.85em\hangafter=1\makebox[1.6em][l]{#1)}#2\par}
\newcommand{\sourcequotepara}[1]{\lawboxprespace{0.04em}\noindent\lawboxuserandnumber\hangindent=1.2em\hangafter=1 #1\par}
\newcommand{\sourcelawline}[1]{\lawboxprespace{0.04em}\noindent\lawboxuserandnumber #1\par}
\newcommand{\lawline}[1]{\lawboxprespace{0.04em}\noindent\lawboxuserandnumber #1\par}
\newcommand{\pendingtranslation}{\par{\songfont\footnotesize\color{pendingink}（中文譯文待補。）}\par}
\newcommand{\tocpart}[1]{\par\addvspace{0.52em}{\footnotesize\bfseries\color{pendingink}#1\par}\addvspace{0.16em}}
\newcommand{\tocdashline}{\par\addvspace{0.48em}\noindent{\color{notegray}\xleaders\hbox to 0.82em{\hfil\rule{0.42em}{0.28pt}\hfil}\hskip\linewidth}\par\addvspace{0.38em}}
% \tocpanel{font-cmd}{content}：將目錄置中於所在欄寬（雙語欄適用）。
% A3 橫式使用 0.58\linewidth，A4 橫式使用 0.84\linewidth，
% 使兩種紙張下目錄欄內容實際寬度相近（均約 100–105 mm）。
\newcommand{\tocpanel}[2]{%
  \makebox[\linewidth][c]{%
    \ifx\bilingualpaper\bilingualpaperaThree
      \begin{minipage}{0.58\linewidth}%
    \else
      \begin{minipage}{0.84\linewidth}%
    \fi
    \toccolstyle
    #1#2%
    \end{minipage}%
  }%
}
% ===== 首頁簡目錄的兩層 description list 縮排 =====
% enumitem 的 description list 可把每一列拆成：
%   [label]  labelsep  body text
% 這裡的「起點」都以所在 minipage/欄位的左緣為 0。
%
% 核心規則：
%   labelindent = label 欄位從左緣開始的位置。
%   labelwidth  = label 欄位本身寬度；標籤太長（如 Abw. Meinung）就加大它。
%   labelsep    = label 欄位右緣到正文的距離。
%   leftmargin  = 正文 body text 的起點。判斷層級視覺時主要看這個。
%
% 所以：
%   想讓「Begründung / 判決理由」整列正文往右或往左：改 tocitems 的 leftmargin。
%   想讓「A. Verfahren und Gegenstand / A. 程序與審查標的」正文往右或往左：
%     改 tocsubitems 的 leftmargin。
%   想讓 A./B./C. 這些標籤本身往右或往左，但正文不動：改 tocsubitems 的 labelindent。
%   想讓 A./B./C. 與正文間距變大或變小：改 tocsubitems 的 labelsep。
%
% 層級原則：
%   tocsubitems.leftmargin 必須大於 tocitems.leftmargin，
%   否則子層級正文會看起來比「Gründe / 理由」還靠左。
\newlist{tocitems}{description}{1}
\setlist[tocitems]{%
  labelindent=0pt,    % 第一層 label 欄位起點；0pt 表示貼齊目錄欄左緣。
  leftmargin=2.54cm,  % 第一層正文起點；例如 Begründung / 判決理由 的左緣。
  labelwidth=2.16cm,  % 第一層 label 欄位寬度；需容納 Leitsätze、Abw. Meinung、不同意見等。
  labelsep=0.32cm,    % 第一層 label 與正文的水平間距；只改間距，不改正文起點。
  itemsep=0.28em,     % 第一層各項之間的垂直距離。
  topsep=0.20em,      % 第一層 list 上下與前後內容的垂直距離。
  parsep=0pt,         % 同一 item 內段落之間的距離；目錄項通常不需要。
  partopsep=0pt,      % list 若緊接段落開始時的額外距離；關閉以免目錄高度飄動。
  align=left,         % label 欄內左對齊；長短標籤不靠右貼齊。
  font=\normalfont\bfseries % label 的字型；正文不受這行影響。
}
\newlist{tocsubitems}{description}{1}
\setlist[tocsubitems]{%
  labelindent=1cm,    % 第二層 label 起點；只控制 A./B./C. 本身的位置。
  leftmargin=3.00cm,  % 第二層正文起點；必須大於 2.54cm，才會比 Begründung / 判決理由更右。
  labelwidth=0.5cm,   % 第二層 label 欄寬；A./B./C. 很短，可小於第一層。
  labelsep=1.5cm,    % 第二層 label 與正文間距；A. 與 Verfahren / 程序 的距離看這裡。
  itemsep=0.12em,     % 第二層各項較密，避免 A.–E. 佔太高。
  topsep=0.04em,      % 第二層 list 與 Gründe / 理由之間的垂直距離。
  parsep=0pt,         % 同一第二層 item 內段落距離；目錄項通常不需要。
  partopsep=0pt,      % 關閉額外段落起始距離，避免版面不穩。
  align=left,         % A./B./C. label 欄內左對齊。
  font=\normalfont\bfseries, % A./B./C. label 加粗；正文不受這行影響。
  before={}           % 不在第二層 list 前自動插入額外內容。
}
\newlist{termitems}{description}{1}
\ifbilingualcolumns
  \setlist[termitems]{%
    leftmargin=5.2cm,
    labelwidth=4.8cm,
    labelsep=0.35cm,
    itemsep=0.16em,
    topsep=0.2em,
    parsep=0pt,
    partopsep=0pt,
    font=\normalfont\bfseries,
    style=nextline
  }
\else
  \setlist[termitems]{%
    leftmargin=0pt,
    labelwidth=0pt,
    labelsep=0pt,
    itemsep=0.24em,
    topsep=0.2em,
    parsep=0pt,
    partopsep=0pt,
    font=\normalfont\bfseries,
    style=nextline
  }
\fi
\newlist{abbritems}{description}{1}
\setlist[abbritems]{%
  leftmargin=2.38cm,
  labelwidth=2.02cm,
  labelsep=0.28cm,
  itemsep=0.16em,
  topsep=0.2em,
  parsep=0pt,
  partopsep=0pt,
  align=left,
  font=\normalfont\bfseries
}
\ifbilingualcolumns
  \newenvironment{termcolumns}{\begin{multicols}{2}}{\end{multicols}}
\else
  \newenvironment{termcolumns}{}{}
\fi

% ===== 編者註腳開關 =====
% 一般版預設不顯示編者註，避免正文腳註過密。
% 若開啟 \classroommodeon，GG/條文編者註仍會自動顯示，無需另開本開關。
% 若一般版也要顯示編者註，可在單案檔 \input 後加 \showeditornotestrue。
\newif\ifshoweditornotes
\showeditornotesfalse

% ===== 目錄排版輔助 =====
\newcommand{\toccolstyle}{\raggedright\arraybackslash}
\newcommand{\tocsingleindent}{\leftskip=1.45cm\rightskip=1.2cm}

% ===== 行內引文環境（德文原文與中文譯文各一，帶左側灰色邊線）=====
\newcommand{\quoteparbreak}{\par\addvspace{0.48em}}
\newcommand{\sourcequoteinner}[1]{%
  \begin{tcolorbox}[
    enhanced, breakable, width=\dimexpr\linewidth-0.8em\relax, frame hidden,
    colback=white, coltext=caseink, boxrule=0pt,
    borderline west={0.55pt}{0.88em}{notegray},
    left=1.78em, right=0.72em, top=0.18em, bottom=0.18em,
    boxsep=0pt, before skip=0.24em, after skip=0.24em,
    before upper={\normalfont\footnotesize\color{caseink}\setlength{\parindent}{0pt}\setlength{\parskip}{0.34em}\raggedright\setlength{\rightskip}{0pt plus 1em}}
  ]%
  #1\par
  \end{tcolorbox}%
}
\newcommand{\transquoteinner}[1]{%
  \begin{tcolorbox}[
    enhanced, breakable, width=\dimexpr\linewidth-0.8em\relax, frame hidden,
    colback=white, coltext=caseink, boxrule=0pt,
    borderline west={0.55pt}{0.88em}{notegray},
    left=1.78em, right=0.72em, top=0.18em, bottom=0.18em,
    boxsep=0pt, before skip=0.24em, after skip=0.24em,
    before upper={\songfont\small\color{caseink}\setlength{\parindent}{0pt}\setlength{\parskip}{0.34em}\raggedright\setlength{\rightskip}{0pt plus 1em}}
  ]%
  #1\par
  \end{tcolorbox}%
}

% ===== 大綱（Gliederung）Rn 輔助命令 =====
\newcommand{\outrn}[2]{\enspace{\normalfont\footnotesize\color{pendingink}(Rn\,#1–#2)}}
\newcommand{\outrnsingle}[1]{\enspace{\normalfont\footnotesize\color{pendingink}(Rn\,#1)}}
\newcommand{\outrnzh}[2]{\enspace{\normalfont\footnotesize\color{pendingink}（第\,#1–#2\,段）}}
\newcommand{\outrnzhs}[1]{\enspace{\normalfont\footnotesize\color{pendingink}（第\,#1\,段）}}
% ===== 大綱雙語對照表格輔助命令（三欄：段碼 │ 德文 │ 中文）=====
% 欄 1：段碼（由欄定義自動格式化：小字、等寬、右對齊、pendingink）
% 欄 2：德文標籤＋正文
% 欄 3：中文標籤＋正文
%
% 無段碼的列（\omainrow、\osecrow），欄 1 以 & 空置。
% 有段碼的列（\osecrowrn、\osubrow、\ointrow），#1 為預格式化字串
%   例：\osubrow{2–49}{I.}{...}{I.}{...}
%       \ointrow{107}{...}{...}
%
% \opartrow：段組標題（橫跨三欄）
\newcommand{\opartrow}[2]{%
  \noalign{\vspace{0.30em}}%
  \multicolumn{3}{@{}l@{}}{%
    \footnotesize\bfseries\color{pendingink}#1\hfill\songfont#2%
  }\\[0.06em]%
}
% \odashrow：虛線分隔（橫跨三欄）
\newcommand{\odashrow}{%
  \noalign{\vspace{0.28em}}%
  \multicolumn{3}{@{}l@{}}{\color{notegray}%
    \xleaders\hbox to 0.82em{\hfil\rule{0.42em}{0.28pt}\hfil}\hskip 0.88\linewidth}%
  \\[0.24em]%
}
% \odissentpart：不同意見書區段標頭。比一般 \opartrow 更像附錄性轉場。
\newcommand{\odissentpart}[2]{%
  \noalign{\vspace{0.42em}}%
  \multicolumn{3}{@{}p{0.88\textwidth}@{}}{%
    \color{notegray}\rule{\linewidth}{0.18pt}%
  }\\[-0.18em]%
  \multicolumn{3}{@{}p{0.88\textwidth}@{}}{%
    \makebox[\linewidth][s]{%
      \footnotesize\bfseries\color{pendingink}#1%
      \hfill{\songfont #2}%
    }%
  }\\[-0.02em]%
}
% \odissentopinion：不同意見書的署名與一行摘要，右欄保留段碼範圍。
\newcommand{\odissentopinion}[5]{%
  \footnotesize\hangindent=0.52cm\hangafter=1%
    {\bfseries\color{caseink}#2}\enspace{\color{notegray}\textperiodcentered}\enspace\color{caseink}#3%
  &\footnotesize\hangindent=0.52cm\hangafter=1%
    {\songfont\bfseries\color{caseink}#4}\enspace{\color{notegray}\textperiodcentered}\enspace\songfont\color{caseink}#5%
  &#1%
  \\[0.12em]%
}
% \omainrow：頂層項目，無段碼（Leitsätze、Urteil、Gründe …）
\newcommand{\omainrow}[4]{%
  \footnotesize\hangindent=1.92cm\hangafter=1%
    \makebox[1.92cm][l]{\bfseries\color{caseink}#1}\color{caseink}#2%
  &\footnotesize\hangindent=1.92cm\hangafter=1%
    \makebox[1.92cm][l]{\bfseries\color{caseink}#3}\songfont\color{caseink}#4%
  &%
  \\[0.15em]%
}
% \omainrowrn：頂層項目，有段碼（不同意見等標籤寬於 \osecrow 框者）
\newcommand{\omainrowrn}[5]{%
  \footnotesize\hangindent=1.92cm\hangafter=1%
    \makebox[1.92cm][l]{\bfseries\color{caseink}#2}\color{caseink}#3%
  &\footnotesize\hangindent=1.92cm\hangafter=1%
    \makebox[1.92cm][l]{\bfseries\color{caseink}#4}\songfont\color{caseink}#5%
  &#1%
  \\[0.15em]%
}
% \osecrow：一級子項，無段碼（A.、C.、D. 等有子項者）
\newcommand{\osecrow}[4]{%
  \footnotesize\hspace{1.10cm}\hangindent=1.98cm\hangafter=1%
    \makebox[0.86cm][l]{\bfseries\color{caseink}#1}\color{caseink}#2%
  &\footnotesize\hspace{1.10cm}\hangindent=1.98cm\hangafter=1%
    \makebox[0.86cm][l]{\bfseries\color{caseink}#3}\songfont\color{caseink}#4%
  &%
  \\[0.11em]%
}
% \osecrowrn：一級子項，有段碼（B.、E. 等完整段落範圍者）
% #1=段碼字串（如 101–106）, #2=DE標籤, #3=DE正文, #4=ZH標籤, #5=ZH正文
\newcommand{\osecrowrn}[5]{%
  \footnotesize\hspace{1.10cm}\hangindent=1.98cm\hangafter=1%
    \makebox[0.86cm][l]{\bfseries\color{caseink}#2}\color{caseink}#3%
  &\footnotesize\hspace{1.10cm}\hangindent=1.98cm\hangafter=1%
    \makebox[0.86cm][l]{\bfseries\color{caseink}#4}\songfont\color{caseink}#5%
  &#1%
  \\[0.11em]%
}
% \osubrow：二級子項，有段碼範圍（I.、II.、A.I. …）
% #1=段碼字串, #2=DE標籤, #3=DE正文, #4=ZH標籤, #5=ZH正文
\newcommand{\osubrow}[5]{%
  \footnotesize\hspace{1.40cm}\hangindent=2.30cm\hangafter=1%
    \makebox[0.90cm][l]{\bfseries\color{caseink}#2}\color{caseink}#3%
  &\footnotesize\hspace{1.40cm}\hangindent=2.30cm\hangafter=1%
    \makebox[0.90cm][l]{\bfseries\color{caseink}#4}\songfont\color{caseink}#5%
  &#1%
  \\[0.09em]%
}
% \ointrow：引入段落，有單一段碼，無標籤
% #1=段碼字串, #2=DE正文, #3=ZH正文
\newcommand{\ointrow}[3]{%
  \footnotesize\hspace{0.52cm}\hangindent=0.52cm\hangafter=1\color{caseink}#2%
  &\footnotesize\hspace{0.52cm}\hangindent=0.52cm\hangafter=1\songfont\color{caseink}#3%
  &#1%
  \\[0.09em]%
}
% \osubsubrow：三級子項（1.、2.、3.）
% #1=段碼字串, #2=DE標籤, #3=DE正文, #4=ZH標籤, #5=ZH正文
\newcommand{\osubsubrow}[5]{%
  \footnotesize\hspace{1.80cm}\hangindent=2.48cm\hangafter=1%
    \makebox[0.68cm][l]{\bfseries\color{caseink}#2}\color{caseink}#3%
  &\footnotesize\hspace{1.80cm}\hangindent=2.48cm\hangafter=1%
    \makebox[0.68cm][l]{\bfseries\color{caseink}#4}\songfont\color{caseink}#5%
  &#1%
  \\[0.08em]%
}
% \osubsubsubrow：四級子項（a）、b））
% #1=段碼字串, #2=DE標籤, #3=DE正文, #4=ZH標籤, #5=ZH正文
\newcommand{\osubsubsubrow}[5]{%
  \footnotesize\hspace{2.10cm}\hangindent=2.68cm\hangafter=1%
    \makebox[0.58cm][l]{\bfseries\color{caseink}#2}\color{caseink}#3%
  &\footnotesize\hspace{2.10cm}\hangindent=2.68cm\hangafter=1%
    \makebox[0.58cm][l]{\bfseries\color{caseink}#4}\songfont\color{caseink}#5%
  &#1%
  \\[0.07em]%
}
% \osubsubsubsubrow：五級子項（aa）、bb））
% 標籤框 0.62cm：足以容納「aa)」在 footnotesize bold 下（約 0.50cm）並留有間距
% #1=段碼字串, #2=DE標籤, #3=DE正文, #4=ZH標籤, #5=ZH正文
\newcommand{\osubsubsubsubrow}[5]{%
  \footnotesize\hspace{2.38cm}\hangindent=3.06cm\hangafter=1%
    \makebox[0.68cm][l]{\bfseries\color{caseink}#2}\color{caseink}#3%
  &\footnotesize\hspace{2.38cm}\hangindent=3.06cm\hangafter=1%
    \makebox[0.68cm][l]{\bfseries\color{caseink}#4}\songfont\color{caseink}#5%
  &#1%
  \\[0.06em]%
}
% \outlinepage：雙語對照大綱頁包裝環境（longtable 自動跨頁）
% 三欄：德文欄 + 中文欄 + 段碼欄（最右，不折行）
% 段碼欄寬度：1.80cm，足以容納最長段碼如「309–348」
% DE/ZH 欄寬：(0.87\textwidth - 2.08cm) / 2
%   驗算：2×欄 + 0.05\textwidth + 0.28cm gap + 1.80cm = 0.87tw - 2.08cm + 0.05tw + 2.08cm = 0.92tw ✓
\newenvironment{outlinepage}{%
  \vspace*{0.40cm}%
  \setlength{\LTleft}{0.04\textwidth}%
  \setlength{\LTright}{0.04\textwidth}%
  \setlength{\tabcolsep}{0pt}%
  % 大綱列距由 longtable 的 \arraystretch 與各 row macro 結尾的 \\[...em] 共同決定。
  % 這裡控制「每一列本身的表格 strut 高度」；數值越大，A./I./1. 各項之間越像有空行。
  % A3 原本用 0.62，視覺上沒有空行；A4 也沿用同值，避免切到 A4 後項目被撐開。
  % 若日後覺得大綱太擠，只微調這個值（例如 0.68），不要改各 \omainrow/\osubrow 的縮排。
  \renewcommand{\arraystretch}{0.62}%
  \begin{longtable}{%
    >{\vspace{0pt}\raggedright\arraybackslash}p{\dimexpr(0.87\textwidth-2.08cm)/2\relax}%
    @{\hspace{0.025\textwidth}}%
    !{\color{notegray}\vrule width 0.12pt}%
    @{\hspace{0.025\textwidth}}%
    >{\vspace{0pt}\raggedright\arraybackslash}p{\dimexpr(0.87\textwidth-2.08cm)/2\relax}%
    @{\hspace{0.28cm}}%
    >{\vspace{0pt}\footnotesize\normalfont\color{pendingink}\raggedright\arraybackslash}p{1.80cm}%
    @{}%
  }%
  % 首頁欄標籤：不要橫跨置中；分別放在德文欄與中文欄的實際欄位起點。
  \footnotesize\bfseries\color{sourceink}Gliederung%
  &\footnotesize\bfseries\songfont\color{sourceink}大綱%
  &%
  \\[0.36em]%
  \endfirsthead%
  % 續頁欄標籤：同樣分欄定位，以灰色細體區分；無需標注「續」。
  \footnotesize\color{pendingink}Gliederung%
  &\footnotesize\songfont\color{pendingink}大綱%
  &%
  \\[0.24em]%
  \endhead%
}{%
  \end{longtable}%
}
 載入。
% 不含：\documentclass、\documentversion、\ggversionde/zh、\tocde/\toczh。
% 日期與首頁版本列由本檔自動產生；各文件只需手動指定 \documentversion。

\usepackage{xeCJK}
\usepackage{fontspec}
\usepackage{geometry}
\usepackage{setspace}
\usepackage{hyperref}
\usepackage{xurl}
\usepackage{bookmark}
\usepackage{enumitem}
\usepackage{xcolor}
\usepackage{titlesec}
\usepackage{array}
\usepackage{longtable}
\usepackage{tcolorbox}
\usepackage{environ}
\usepackage{pdflscape}
\usepackage{etoolbox}
\usepackage{ragged2e}
\usepackage{paracol}
\usepackage{multicol}
\usepackage{needspace}
\usepackage{fancyhdr}
\usepackage{tikz}
\usepackage{zref-abspage}
\tcbuselibrary{breakable,skins}
\usetikzlibrary{calc}
\usepackage{pgfcalendar}

\hypersetup{
  unicode=true,
  bookmarksopen=true,
  bookmarksnumbered=false,
  hidelinks
}

% ===== 自動推導，不要手動改下面 =====
% （\docmode 與 \bilingualpaper 由各文件在 % bverfge_layout.tex
% 共用排版基礎設施，供 Schwangerschaftsabbruch I 與 II 共享。
% 由各文件以 \documentclass 後 \input{../../共用LaTeX/bverfge_layout} 載入。
% 不含：\documentclass、\documentversion、\ggversionde/zh、\tocde/\toczh。
% 日期與首頁版本列由本檔自動產生；各文件只需手動指定 \documentversion。

\usepackage{xeCJK}
\usepackage{fontspec}
\usepackage{geometry}
\usepackage{setspace}
\usepackage{hyperref}
\usepackage{xurl}
\usepackage{bookmark}
\usepackage{enumitem}
\usepackage{xcolor}
\usepackage{titlesec}
\usepackage{array}
\usepackage{longtable}
\usepackage{tcolorbox}
\usepackage{environ}
\usepackage{pdflscape}
\usepackage{etoolbox}
\usepackage{ragged2e}
\usepackage{paracol}
\usepackage{multicol}
\usepackage{needspace}
\usepackage{fancyhdr}
\usepackage{tikz}
\usepackage{zref-abspage}
\tcbuselibrary{breakable,skins}
\usetikzlibrary{calc}
\usepackage{pgfcalendar}

\hypersetup{
  unicode=true,
  bookmarksopen=true,
  bookmarksnumbered=false,
  hidelinks
}

% ===== 自動推導，不要手動改下面 =====
% （\docmode 與 \bilingualpaper 由各文件在 \input{bverfge_layout} 之前設定；
%   預設 chinese / a4）
\ifdefined\docmode\else\def\docmode{chinese}\fi
\ifdefined\bilingualpaper\else\def\bilingualpaper{a4}\fi
\newif\ifshoworiginal
\newif\ifshowtranslation
\newif\ifbilinguallandscape
\newif\ifbilingualcolumns
\newif\ifintranslation
\newif\iffirstbilingualsection

\showoriginalfalse
\showtranslationfalse
\bilinguallandscapefalse
\bilingualcolumnsfalse
\intranslationfalse
\firstbilingualsectionfalse

\def\modechinese{chinese}
\def\modegerman{german}
\def\modebilingual{bilingual}
\def\bilingualpaperaThree{a3}
\def\bilingualpaperaFour{a4}

\ifx\docmode\modechinese
  \showtranslationtrue
\fi

\ifx\docmode\modegerman
  \showoriginaltrue
\fi

\ifx\docmode\modebilingual
  \showoriginaltrue
  \showtranslationtrue
  \bilinguallandscapetrue
  \bilingualcolumnstrue
\fi

% 編排模式說明：
% chinese  → \showoriginalfalse  \showtranslationtrue  （A4 直式，純中文）
% german   → \showoriginaltrue   \showtranslationfalse （A4 直式，純德文）
% bilingual→ \showoriginaltrue   \showtranslationtrue  \bilingualcolumnstrue（A3/A4 橫式對照）

\ifbilingualcolumns
  \ifx\bilingualpaper\bilingualpaperaFour
    \geometry{
      a4paper,
      landscape,
      left=1.25cm,
      right=2.25cm,
      top=1.25cm,
      bottom=1.75cm,
      footskip=8.5mm,
      marginparwidth=0.75cm,
      marginparsep=0.25cm
    }
  \else
    \geometry{
      a3paper,
      landscape,
      left=1.55cm,
      right=2.75cm,
      top=1.55cm,
      bottom=2.05cm,
      footskip=9.5mm,
      marginparwidth=0.9cm,
      marginparsep=0.35cm
    }
  \fi
\else
\geometry{
  a4paper,
  portrait,
  left=2.2cm,
  right=2.6cm,
  top=2.0cm,
  bottom=2.0cm,
  marginparwidth=0.8cm,
  marginparsep=0.3cm
}
\fi
% ===== 時間戳基礎設施 =====
\newcount\stamptimeval
\newcount\stampminuteval
\newcommand{\stamphh}{%
  \stamptimeval=\time
  \divide\stamptimeval by 60
  \ifnum\stamptimeval<10 0\fi
  \the\stamptimeval}
\newcommand{\stampmm}{%
  \stamptimeval=\time
  \stampminuteval=\time
  \divide\stampminuteval by 60
  \multiply\stampminuteval by 60
  \advance\stamptimeval by -\stampminuteval
  \ifnum\stamptimeval<10 0\fi
  \the\stamptimeval}
\newcommand{\stampwkde}[1]{%
  \ifcase#1Montag\or Dienstag\or Mittwoch\or Donnerstag\or Freitag\or Samstag\or Sonntag\fi}
\newcommand{\stampwkzh}[1]{%
  \ifcase#1 一\or 二\or 三\or 四\or 五\or 六\or 日\fi}
\newcommand{\stampmonthde}[1]{%
  \ifcase#1\or Januar\or Februar\or März\or April\or Mai\or Juni\or
  Juli\or August\or September\or Oktober\or November\or Dezember\fi}
\newcount\stampjdval
\newcount\stampwdval
\pgfcalendardatetojulian{\the\year-\the\month-\the\day}{\stampjdval}
\pgfcalendarjuliantoweekday{\the\stampjdval}{\stampwdval}
\newcommand{\timestampzh}{%
  \songfont 最後修改：\the\year 年\the\month 月\the\day 日%
  % （星期\stampwkzh{\the\stampwdval}）%
  % \space\stamphh:\stampmm
  }

\newcommand{\documentdatezh}{\the\year 年 \the\month 月 \the\day 日}
\newcommand{\documentdatede}{\the\day. \stampmonthde{\the\month} \the\year}
\newcommand{\bverfgetranslationnote}{%
  {\songfont 翻譯說明：初譯使用 Claude Code 與 GPT 協助迭代；仍請以德文原文為準。}%
}
\newcommand{\bverfgecourseinfo}{%
  {\songfont 課程：114 學年度第 2 學期；憲法專題研究（四）；授課老師：湯德宗教授；導讀日期：115 年 6 月 1 日。}%
}
\newcommand{\bverfgeeditorinfo}{%
  {\songfont 編者：王逸帆（R10A21126），國立台灣大學法律系研究所碩士班。}%
}
\newcommand{\bverfgetitleversionline}{%
  \ifbilingualcolumns
    Fassung \documentversion\enspace\textperiodcentered\enspace Stand: \documentdatede
    \quad
    {\songfont 版本 \documentversion\enspace\textperiodcentered\enspace 修訂日期：\documentdatezh\par}%
  \else
    \ifshowtranslation
      {\songfont 版本 \documentversion\enspace\textperiodcentered\enspace 修訂日期：\documentdatezh\par}%
    \else
      Fassung \documentversion\enspace\textperiodcentered\enspace Stand: \documentdatede\par
    \fi
  \fi
}
\newcommand{\bverfgetitlecornerinfo}{%
  \ifshowtranslation
    \vspace{0.72em}%
    \noindent\hfill
    \begin{minipage}{0.66\textwidth}
      \raggedleft\tiny\color{pendingink}\songfont
      \bverfgecourseinfo\par
      \bverfgeeditorinfo
    \end{minipage}%
  \fi
}
\newcommand{\bverfgetitlemeta}{%
  \vfill
  \begin{center}%
    \footnotesize\color{pendingink}%
    \bverfgetitleversionline
    \ifshowtranslation
      \vspace{0.28em}{\scriptsize\bverfgetranslationnote\par}%
    \fi
  \end{center}%
  \bverfgetitlecornerinfo
}

\pagestyle{fancy}
\fancyhf{}
\renewcommand{\headrulewidth}{0pt}
\renewcommand{\footrulewidth}{0pt}
\fancyfoot[C]{\thepage}
\ifbilingualcolumns
  \fancyfoot[R]{\scriptsize\color{pendingink}\timestampzh}%
\else
  \ifshoworiginal
  \else
    \fancyfoot[R]{\scriptsize\color{pendingink}\timestampzh\hspace{0.45cm}}%
  \fi
\fi
\setmainfont{Times New Roman}
\setsansfont{Helvetica Neue}
\setmonofont{Menlo}
\IfFileExists{/Users/iw/Library/Fonts/kaiu.ttf}{
  \setCJKmainfont[Path=/Users/iw/Library/Fonts/,AutoFakeBold=4]{kaiu.ttf}
}{
  \IfFontExistsTF{標楷體}{
    \setCJKmainfont[AutoFakeBold=4]{標楷體}
  }{
    \setCJKmainfont{Songti TC}
  }
}
\IfFontExistsTF{Songti TC}{
  \setCJKfamilyfont{song}{Songti TC}
  \setCJKmonofont{Songti TC}
}{
  \setCJKfamilyfont{song}[Path=/Users/iw/Library/Fonts/,AutoFakeBold=4]{kaiu.ttf}
  \setCJKmonofont[Path=/Users/iw/Library/Fonts/,AutoFakeBold=4]{kaiu.ttf}
}
\newcommand{\songfont}{\CJKfamily{song}}

% ===== 封面／目錄專用打字機字體（僅限封面、目錄頁使用，不影響正文）=====
\IfFontExistsTF{Radio Newsman}{%
  \newfontfamily\bverfgetitlede{Radio Newsman}%
}{%
  \newfontfamily\bverfgetitlede{Courier New}%
}
\IfFontExistsTF{Erikas Farbband}{%
  \newfontfamily\bverfgebodyde{Erikas Farbband}%
}{%
  \let\bverfgebodyde\bverfgetitlede
}
\IfFontExistsTF{Maoken Heavy Labourer}{%
  \setCJKfamilyfont{bverfgetitlecjk}{Maoken Heavy Labourer}%
}{%
  \setCJKfamilyfont{bverfgetitlecjk}{Songti TC}%
}
\IfFontExistsTF{Huiwen-mincho}{%
  \setCJKfamilyfont{bverfgebodycjk}{Huiwen-mincho}%
}{%
  \setCJKfamilyfont{bverfgebodycjk}{Songti TC}%
}
\newcommand{\bverfgetitlezh}{\bverfgetitlede\CJKfamily{bverfgetitlecjk}}
\newcommand{\bverfgebodyzh}{\bverfgebodyde\CJKfamily{bverfgebodycjk}}

\setstretch{1.35}
\setlist{itemsep=0.2em, topsep=0.4em}
\setlength{\parindent}{0pt}
\setlength{\parskip}{0.45em}
\raggedbottom
\setcounter{secnumdepth}{0}
\setcounter{tocdepth}{2}

\newcommand{\lawboxsourcefont}{\footnotesize}
\newcommand{\lawboxtransfont}{\footnotesize}
\newcommand{\lawboxsourcestretch}{1.02}
\newcommand{\lawboxtransstretch}{0.92}

\titleformat{\section}{\centering\Large\bfseries\songfont}{\thesection}{1em}{}
\titleformat{\subsection}{\large\bfseries\songfont}{\thesubsection}{1em}{}
\titleformat{\subsubsection}{\normalsize\bfseries\songfont}{\thesubsubsection}{1em}{}
\titleformat{\paragraph}{\normalsize\bfseries\songfont}{\theparagraph}{1em}{}
\titlespacing*{\section}{0pt}{1.8em}{1.0em}
\titlespacing*{\subsection}{0pt}{1.2em}{0.6em}
\titlespacing*{\subsubsection}{0pt}{1.0em}{0.45em}
\titlespacing*{\paragraph}{0pt}{0.6em}{0.4em}

\definecolor{noteblue}{HTML}{A9C9F5}
\definecolor{noteteal}{HTML}{AEE1D6}
\definecolor{notegray}{HTML}{D7DADF}
\definecolor{caseink}{HTML}{000000}
\definecolor{casepaper}{HTML}{FCFEFD}
\definecolor{sourcepaper}{HTML}{FAFBF8}
\definecolor{sourceline}{HTML}{D7DED8}
\definecolor{sourceink}{HTML}{000000}
\definecolor{pendingink}{HTML}{000000}

\tcbset{
  caseboxbase/.style={
    enhanced,
    breakable,
    width=0.985\linewidth,
    colback=casepaper,
    coltitle=caseink,
    fonttitle=\bfseries\songfont,
    boxrule=0.28pt,
    arc=1.2mm,
    left=2.4mm,
    right=2.4mm,
    top=1.5mm,
    bottom=1.5mm,
    boxsep=1.5mm,
    before skip=0.65em,
    after skip=0.65em,
    before upper={\setlength{\parindent}{0pt}\setlength{\parskip}{0.35em}},
  }
}

\newcommand{\bilingualstart}{}
\newcommand{\bilingualstop}{}
\ifbilingualcolumns
  \renewcommand{\bilingualstart}{\small\setlength{\columnsep}{1.25cm}\setlength{\columnseprule}{0.12pt}\columnratio{0.5,0.5}\multicolumnfootnotes\flushbottom\sloppy\interlinepenalty=80\clubpenalty=800\widowpenalty=800\displaywidowpenalty=800\begin{paracol}{2}\global\intranslationfalse\global\firstbilingualsectiontrue{\footnotesize\color{sourceink}Deutsch\par}\switchcolumn{\footnotesize\songfont\color{sourceink}中文譯文\par}\switchcolumn*}
  \renewcommand{\bilingualstop}{\ifintranslation\switchcolumn*\global\intranslationfalse\fi\end{paracol}\clearpage\raggedbottom}
\fi
\newif\ifrandnumbers
\randnumbersfalse
\newcounter{randnumber}
\newcounter{randnumberlocal}
\newcounter{lastsourcerandstart}
\newif\iflastsourcehasrandnumbers
\lastsourcehasrandnumbersfalse
\newif\ifinlawbox
\inlawboxfalse
\newif\iflawboxhaspendingrn
\lawboxhaspendingrnfalse
\newcommand{\lawboxpendingrn}{}
\newcommand{\startrandnumbers}{\global\randnumberstrue\setcounter{randnumber}{1}\global\lastsourcehasrandnumbersfalse}
\newcommand{\stoprandnumbers}{\global\randnumbersfalse\global\lastsourcehasrandnumbersfalse}
\newlength{\randnumberwidth}
\setlength{\randnumberwidth}{0.95cm}
\newlength{\randnumberrightoffset}
\setlength{\randnumberrightoffset}{0.85cm}
\newlength{\randnumberlawboxrightoffset}
\setlength{\randnumberlawboxrightoffset}{1.40cm}
\newcommand{\randnumberbox}[1]{%
  \leavevmode
  \makebox[0pt][l]{%
    \smash{%
      \tikz[remember picture,overlay,baseline=(rnmark.base)]{%
        \coordinate (rnmark) at (0,0);%
        \ifinlawbox
          \path let \p1=(current page.east), \p2=(rnmark.base) in
            node[
              anchor=base east,
              inner sep=0pt,
              outer sep=0pt,
              text width=\randnumberwidth,
              align=right,
              text=pendingink,
              font=\normalfont\mdseries\scriptsize\ttfamily
            ] at (\x1-\randnumberlawboxrightoffset,\y2) {{\normalfont\mdseries\scriptsize\ttfamily #1}};%
        \else
          \path let \p1=(current page.east), \p2=(rnmark.base) in
            node[
              anchor=base east,
              inner sep=0pt,
              outer sep=0pt,
              text width=\randnumberwidth,
              align=right,
              text=pendingink,
              font=\normalfont\mdseries\scriptsize\ttfamily
            ] at (\x1-\randnumberrightoffset,\y2) {{\normalfont\mdseries\scriptsize\ttfamily #1}};%
        \fi
      }%
    }%
  }%
  \ignorespaces
}
\newcommand{\lawboxstorerandnumber}[1]{%
  \gdef\lawboxpendingrn{#1}%
  \global\lawboxhaspendingrntrue
  \ignorespaces
}
\newcommand{\lawboxuserandnumber}{%
  \iflawboxhaspendingrn
    \randnumberbox{\lawboxpendingrn}%
    \global\lawboxhaspendingrnfalse
  \fi
}
\newcommand{\sourcern}[1]{%
  \ifrandnumbers
    \ifshoworiginal
      \ifshowtranslation\else
        \ifinlawbox\lawboxstorerandnumber{#1}\else\randnumberbox{#1}\fi
      \fi
    \fi
  \fi
}
\newcommand{\transrn}[1]{%
  \ifrandnumbers
    \ifshowtranslation
      \ifinlawbox\lawboxstorerandnumber{#1}\else\randnumberbox{#1}\fi
    \fi
  \fi
}
% ===== 課堂模式：德文原文縮寫註解 =====
% 日常切換只改各判決檔的一行：
%   \classroommodeon        開啟課堂版；註解掉這行即回到一般版。
% 模板預設：
%   1. 課堂模式關閉。
%   2. 課堂模式開啟時，縮寫註解包含「德文全稱＋中文譯名」。
%   3. 課堂模式開啟時，GG/條文編者註仍會自動顯示。
% 進階微調才需要在單案檔覆寫：
%   \classroomabbrzhoff     縮寫註解僅顯示德文全稱。
%   \showeditornotestrue    一般版也顯示編者註。
% 註解策略：同一實際 PDF 頁面中，同一縮寫只在第一次出現處加註。
% 這裡用 zref-abspage 的絕對頁序，不用 \thepage，避免文件內頁碼重置時誤判。
\newif\ifclassroommode
\newif\ifclassroomabbrzh
\classroommodefalse
\classroomabbrzhtrue
\newcommand{\classroommodeon}{\classroommodetrue}
\newcommand{\classroommodeoff}{\classroommodefalse}
\newcommand{\classroomabbrzhon}{\classroomabbrzhtrue}
\newcommand{\classroomabbrzhoff}{\classroomabbrzhfalse}
\newcommand{\abbrnotelabel}{\emph{Abk.}: }
\newcommand{\abbrnote}[3]{%
  \ifclassroommode
    \ifshoworiginal
      \ifintranslation\else
        \footnote{\normalfont\footnotesize\abbrnotelabel #2\ifclassroomabbrzh；{\songfont #3}\fi}%
      \fi
    \fi
  \fi
}
\newcommand{\abbrnoteonce}[4]{%
  #2%
  \ifclassroommode
    \ifshoworiginal
      \ifintranslation\else
        \begingroup
          \edef\abbrcurrentabspage{\number\numexpr\value{abspage}+1\relax}%
          \ifcsname abbrnoteseen@#1@\abbrcurrentabspage\endcsname\else
            \global\expandafter\let\csname abbrnoteseen@#1@\abbrcurrentabspage\endcsname\empty
            \abbrnote{#1}{#3}{#4}%
          \fi
        \endgroup
      \fi
    \fi
  \fi
}
\newcommand{\abbrAbs}{\abbrnoteonce{abs}{Abs.}{Absatz}{項}}
\newcommand{\abbrAAO}{\abbrnoteonce{aao}{a.a.O.}{am angegebenen Ort}{前引處}}
\newcommand{\abbrAE}{\abbrnoteonce{ae}{AE}{Alternativ-Entwurf}{替代草案}}
\newcommand{\abbrArt}{\abbrnoteonce{art}{Art.}{Artikel}{條}}
\newcommand{\abbrAufl}{\abbrnoteonce{aufl}{Aufl.}{Auflage}{版}}
\newcommand{\abbrBd}{\abbrnoteonce{bd}{Bd.}{Band}{卷}}
\newcommand{\abbrBGBl}{\abbrnoteonce{bgbl}{BGBl.}{Bundesgesetzblatt}{聯邦法律公報}}
\newcommand{\abbrBGHSt}{\abbrnoteonce{bghst}{BGHSt}{Entscheidungen des Bundesgerichtshofes in Strafsachen}{聯邦普通法院刑事判決彙編}}
\newcommand{\abbrBRDrucks}{\abbrnoteonce{brdrucks}{BRDrucks.}{Bundesrats-Drucksache}{聯邦參議院文件}}
\newcommand{\abbrBTDrucks}{\abbrnoteonce{btdrucks}{BTDrucks.}{Bundestags-Drucksache}{聯邦議院文件}}
\newcommand{\abbrBundesgesetzbl}{\abbrnoteonce{bundesgesetzbl}{Bundesgesetzbl.}{Bundesgesetzblatt}{聯邦法律公報}}
\newcommand{\abbrBvF}{\abbrnoteonce{bvf}{BvF}{Bundesverfassungsgerichtsverfahren}{聯邦憲法法院程序案號}}
\newcommand{\abbrBvL}{\abbrnoteonce{bvl}{BvL}{Bundesverfassungsgerichtsverfahren}{聯邦憲法法院程序案號}}
\newcommand{\abbrBvR}{\abbrnoteonce{bvr}{BvR}{Bundesverfassungsgerichtsverfahren}{聯邦憲法法院程序案號}}
\newcommand{\abbrBVerfG}{\abbrnoteonce{bverfg}{BVerfG}{Bundesverfassungsgericht}{聯邦憲法法院}}
\newcommand{\abbrBVerfGE}{\abbrnoteonce{bverfge}{BVerfGE}{Entscheidungen des Bundesverfassungsgerichts}{聯邦憲法法院判決集}}
\newcommand{\abbrBVerfGG}{\abbrnoteonce{bverfgg}{BVerfGG}{Bundesverfassungsgerichtsgesetz}{聯邦憲法法院法}}
\newcommand{\abbrDDR}{\abbrnoteonce{ddr}{DDR}{Deutsche Demokratische Republik}{德意志民主共和國；東德}}
\newcommand{\abbrDJT}{\abbrnoteonce{djt}{DJT}{Deutscher Juristentag}{德國法學家大會}}
\newcommand{\abbrDr}{\abbrnoteonce{dr}{Dr.}{Doktor}{Dr.}}
\newcommand{\abbrEuGRZ}{\abbrnoteonce{eugrz}{EuGRZ}{Europäische Grundrechte-Zeitschrift}{歐洲基本權期刊}}
\newcommand{\abbrEV}{\abbrnoteonce{ev}{EV}{Einigungsvertrag}{統一條約}}
\newcommand{\abbrF}{\abbrnoteonce{f}{f.}{folgende}{以下一頁；下一段}}
\newcommand{\abbrFF}{\abbrnoteonce{ff}{ff.}{fortfolgende}{以下數頁；以下各段}}
\newcommand{\abbrGBlDDR}{\abbrnoteonce{gblddr}{GBl. DDR}{Gesetzblatt der Deutschen Demokratischen Republik}{德意志民主共和國法律公報}}
\newcommand{\abbrGG}{\abbrnoteonce{gg}{GG}{Grundgesetz}{基本法}}
\newcommand{\abbrGVBl}{\abbrnoteonce{gvbl}{GVBl.}{Gesetz- und Verordnungsblatt}{法律及命令公報}}
\newcommand{\abbrJR}{\abbrnoteonce{jr}{JR}{Juristische Rundschau}{法學評論}}
\newcommand{\abbrJZ}{\abbrnoteonce{jz}{JZ}{JuristenZeitung}{法學家報}}
\newcommand{\abbrIVm}{\abbrnoteonce{ivm}{i.V.m.}{in Verbindung mit}{結合；連同}}
\newcommand{\abbrMdB}{\abbrnoteonce{mdb}{MdB}{Mitglied des Bundestages}{聯邦議院議員}}
\newcommand{\abbrMwN}{\abbrnoteonce{mwn}{m.w.N.}{mit weiteren Nachweisen}{附進一步文獻／判例出處}}
\newcommand{\abbrNF}{\abbrnoteonce{nf}{n.F.}{neue Fassung}{新版本}}
\newcommand{\abbrAF}{\abbrnoteonce{af}{a.F.}{alte Fassung}{舊版本}}
\newcommand{\abbrNr}{\abbrnoteonce{nr}{Nr.}{Nummer}{款、目或號，依語境}}
\newcommand{\abbrProf}{\abbrnoteonce{prof}{Prof.}{Professor}{教授}}
\newcommand{\abbrRdnr}{\abbrnoteonce{rdnr}{Rdnr.}{Randnummer}{邊碼；段碼}}
\newcommand{\abbrRdnrn}{\abbrnoteonce{rdnrn}{Rdnrn.}{Randnummern}{邊碼；段碼}}
\newcommand{\abbrRGBl}{\abbrnoteonce{rgbl}{RGBl.}{Reichsgesetzblatt}{帝國法律公報}}
\newcommand{\abbrRGSt}{\abbrnoteonce{rgst}{RGSt}{Entscheidungen des Reichsgerichts in Strafsachen}{帝國法院刑事判決彙編}}
\newcommand{\abbrRspr}{\abbrnoteonce{rspr}{Rspr.}{Rechtsprechung}{判例；司法實務}}
\newcommand{\abbrRVO}{\abbrnoteonce{rvo}{RVO}{Reichsversicherungsordnung}{帝國保險法}}
\newcommand{\abbrS}{\abbrnoteonce{s}{S.}{Seite}{頁}}
\newcommand{\abbrSFHG}{\abbrnoteonce{sfhg}{SFHG}{Schwangeren- und Familienhilfegesetz}{孕婦及家庭扶助法}}
\newcommand{\abbrSGBV}{\abbrnoteonce{sgbv}{SGB V}{Fünftes Buch Sozialgesetzbuch}{社會法典第五編}}
\newcommand{\abbrStAG}{\abbrnoteonce{staeg}{StÄG}{Strafrechtsänderungsgesetz}{刑法修正法}}
\newcommand{\abbrStenBer}{\abbrnoteonce{stenber}{StenBer.}{Stenographischer Bericht}{速記紀錄}}
\newcommand{\abbrStGB}{\abbrnoteonce{stgb}{StGB}{Strafgesetzbuch}{刑法}}
\newcommand{\abbrStREG}{\abbrnoteonce{streg}{StREG}{Strafrechtsreform-Ergänzungsgesetz}{刑法改革補充法}}
\newcommand{\abbrStrRG}{\abbrnoteonce{strrg}{StrRG}{Strafrechtsreformgesetz}{刑法改革法}}
\newcommand{\abbrUS}{\abbrnoteonce{us}{U.S.}{United States Reports}{美國最高法院判決彙編}}
\newcommand{\abbrVgl}{\abbrnoteonce{vgl}{vgl.}{vergleiche}{參見}}
\newcommand{\abbrVVDStRL}{\abbrnoteonce{vvdstrl}{VVDStRL}{Veröffentlichungen der Vereinigung der Deutschen Staatsrechtslehrer}{德國公法教師協會論文集}}
\newcommand{\abbrWP}{\abbrnoteonce{wp}{WP.}{Wahlperiode}{屆期}}
\newcommand{\abbrWp}{\abbrnoteonce{wp}{Wp.}{Wahlperiode}{屆期}}
\newcommand{\abbrZStrW}{\abbrnoteonce{zstrw}{ZStrW}{Zeitschrift für die gesamte Strafrechtswissenschaft}{整體刑法學期刊}}
\newcommand{\recordsourcerandstart}{%
  \ifrandnumbers
    \setcounter{lastsourcerandstart}{\value{randnumber}}%
    \global\lastsourcehasrandnumberstrue
  \else
    \global\lastsourcehasrandnumbersfalse
  \fi
}
\newlength{\biparagraphskip}
\setlength{\biparagraphskip}{0.52em}
\newlength{\lawboxblockbeforeskip}
\setlength{\lawboxblockbeforeskip}{0.72em}
\newlength{\lawboxblockafterskip}
\setlength{\lawboxblockafterskip}{0.92em}
\newif\iflawboxpartendedblock
\lawboxpartendedblockfalse
\newcommand{\sourceparstyle}{\footnotesize\color{sourceink}\setstretch{1.12}\RaggedRight\emergencystretch=3em\setlength{\parskip}{0.42em}\obeylines\selectfont}
\newcommand{\transparstyle}{\songfont\color{caseink}\setstretch{1.12}\setlength{\parindent}{0pt}\setlength{\parskip}{0.42em}}
\newcommand{\bilingualpairskip}{%
  \par\addvspace{\biparagraphskip}%
  \switchcolumn
  \par\addvspace{\biparagraphskip}%
  \switchcolumn*%
  \global\intranslationfalse
}
\newcommand{\bilingualboxpairskip}{%
  \par
  \switchcolumn
  \par
  \switchcolumn*%
  \global\intranslationfalse
}
\newcommand{\bilinguallawboxblockskip}{%
  \switchcolumn*%
  \par\addvspace{\lawboxblockafterskip}%
  \switchcolumn
  \par\addvspace{\lawboxblockafterskip}%
  \switchcolumn*%
  \global\intranslationfalse
}
\newcommand{\bikeepwithnext}[1][4]{%
  \ifbilingualcolumns
    \syncleft
    \Needspace{#1\baselineskip}%
    \switchcolumn
    \Needspace{#1\baselineskip}%
    \switchcolumn*%
    \global\intranslationfalse
  \else
    \Needspace{#1\baselineskip}%
  \fi
}
\NewEnviron{sourcepara}[1][]{
  \recordsourcerandstart
  \ifshoworiginal
    \ifbilingualcolumns
      \ifintranslation\switchcolumn*\global\intranslationfalse\fi
    \else
      \Needspace{5\baselineskip}
    \fi
    \ifbilingualcolumns\else\par\addvspace{0.35em}\fi
    {\sourceparstyle
    \noindent\BODY\par}
    \ifbilingualcolumns\else\addvspace{0.25em}\fi
    \ifbilingualcolumns
      \switchcolumn
      \global\intranslationtrue
    \fi
  \else
    \ifrandnumbers
      \begingroup
        \setbox0=\vbox{\hsize=\linewidth\sourceparstyle\noindent\BODY\par}%
      \endgroup
    \fi
  \fi
}
\NewEnviron{transpara}{
  \ifshowtranslation
    \setcounter{randnumberlocal}{\value{lastsourcerandstart}}%
    \ifbilingualcolumns
      \ifintranslation\else\switchcolumn\global\intranslationtrue\fi
      {\transparstyle\setlength{\leftskip}{0.55em}\setlength{\rightskip}{0.15em plus 1em}\addtolength{\linewidth}{-0.55em}\BODY\par}
      \switchcolumn*
      \bilingualpairskip
    \else
      {\transparstyle\BODY\par}
    \fi
  \else
    \ifbilingualcolumns
      \ifintranslation\switchcolumn*\global\intranslationfalse\fi
    \fi
  \fi
}
\NewEnviron{sourceboxpara}[1][]{
  \global\lawboxpartendedblockfalse
  \recordsourcerandstart
  \ifshoworiginal
    \ifbilingualcolumns
      \ifintranslation\switchcolumn*\global\intranslationfalse\fi
    \else
      \Needspace{3\baselineskip}
    \fi
    {\sourceparstyle
    \BODY\par}
    \ifbilingualcolumns\else
      \iflawboxpartendedblock\addvspace{\lawboxblockafterskip}\fi
    \fi
    \ifbilingualcolumns
      \switchcolumn
      \global\intranslationtrue
    \fi
  \fi
}
\NewEnviron{transboxpara}{
  \global\lawboxpartendedblockfalse
  \ifshowtranslation
    \setcounter{randnumberlocal}{\value{lastsourcerandstart}}%
    \ifbilingualcolumns
      \ifintranslation\else\switchcolumn\global\intranslationtrue\fi
      {\transparstyle\BODY\par}
      \iflawboxpartendedblock
        \bilinguallawboxblockskip
      \else
        \switchcolumn*
        \bilingualboxpairskip
      \fi
    \else
      {\transparstyle\BODY\par}
      \iflawboxpartendedblock\addvspace{\lawboxblockafterskip}\fi
    \fi
  \else
    \ifbilingualcolumns
      \ifintranslation\switchcolumn*\global\intranslationfalse\fi
    \fi
  \fi
}
\newcommand{\leitsatznum}[1]{\noindent\hangindent=3.0em\hangafter=1\makebox[2.25em][r]{\bfseries #1.}\hspace{0.55em}\ignorespaces}
\newcommand{\syncleft}{\ifbilingualcolumns\ifintranslation\switchcolumn*\global\intranslationfalse\fi\fi}
\pretocmd{\section}{\syncleft\clearpage}{}{}
\pretocmd{\subsection}{\syncleft}{}{}
\pretocmd{\subsubsection}{\syncleft}{}{}
\newcommand{\biheadingfont}{\songfont}
\newcommand{\bitocentry}[2]{%
  \ifshoworiginal
    \ifshowtranslation
      \ifstrequal{#1}{#2}{#1}{#1\space #2}%
    \else
      #1%
    \fi
  \else
    #2%
  \fi
}
\pdfstringdefDisableCommands{%
  \def\bitocentry#1#2{#1}%
}
\newcommand{\biheading}[7]{%
  \syncleft#1\bikeepwithnext[#3]\phantomsection\addcontentsline{toc}{#2}{\bitocentry{#6}{#7}}%
  \ifbilingualcolumns
    {#4 #6\par}%
    \switchcolumn
    {#4\biheadingfont #7\par}%
    \switchcolumn*%
    \global\intranslationfalse
    \par\addvspace{#5}%
    \switchcolumn
    \par\addvspace{#5}%
    \switchcolumn*%
  \else
    \ifshoworiginal{#4 #6\par}\fi
    \ifshowtranslation{#4\biheadingfont #7\par}\fi
    \addvspace{#5}%
  \fi
}
\newcommand{\termsectionheading}[1]{%
  \par\Needspace{9\baselineskip}%
  \vspace{0.65em}%
  {\songfont\bfseries #1\par}%
  \vspace{0.2em}%
}
% 章節換頁：
%   雙語 paracol 欄內不要用 \clearpage；它會在同步兩欄時吐出只有欄線/頁碼的空頁。
%   \newpage 足以讓下一個 section 起新頁，又不會強制清空所有待排材料。
%   單語模式不在 paracol 內，仍可用 \clearpage。
\newcommand{\bisectionpagebreak}{\ifbilingualcolumns\newpage\else\clearpage\fi}
\newcommand{\bisection}[2]{%
  \biheading{\iffirstbilingualsection\global\firstbilingualsectionfalse\else\bisectionpagebreak\fi}{section}{6}{\centering\Large\bfseries}{0.8em}{#1}{#2}%
}
\newcommand{\bisubsection}[2]{%
  \biheading{}{subsection}{5}{\large\bfseries}{0.55em}{#1}{#2}%
}
\newcommand{\bisubsubsection}[2]{%
  \biheading{}{subsubsection}{4}{\normalsize\bfseries}{0.45em}{#1}{#2}%
}
\newcommand{\bicompositeheading}[4]{%
  \biheading{}{subsection}{5}{\large\bfseries}{0.16em}{#1}{#2}%
  \biheading{}{subsubsection}{3}{\normalsize\bfseries}{0.45em}{#3}{#4}%
}
\newif\iflawboxfirstitem
\lawboxfirstitemfalse
\newcommand{\lawboxprespace}[1]{%
  \par
  \iflawboxfirstitem
    \global\lawboxfirstitemfalse
  \else
    \addvspace{#1}%
  \fi
}
\newtcolorbox{lawboxinner}{
  caseboxbase,
  colframe=sourceline!85!black,
  colback=white,
  left=1.05em,
  right=1.05em,
  top=0.62em,
  bottom=0.62em,
  before upper={\global\inlawboxtrue\global\lawboxfirstitemtrue\global\lawboxhaspendingrnfalse\ifintranslation\lawboxtransfont\else\lawboxsourcefont\fi\RaggedRight\setlength{\parindent}{0pt}\setlength{\parskip}{0pt}\ifintranslation\setstretch{\lawboxtransstretch}\else\setstretch{\lawboxsourcestretch}\fi},
  after upper={\global\lawboxhaspendingrnfalse\global\inlawboxfalse}
}
\tcbset{
  lawboxpartbase/.style={
    enhanced,
    breakable=false,
    width=0.985\linewidth,
    colback=white,
    frame hidden,
    boxrule=0pt,
    arc=0pt,
    left=1.05em,
    right=1.05em,
    boxsep=1.5mm,
    before skip=0pt,
    after skip=0pt,
    before upper={\global\inlawboxtrue\global\lawboxfirstitemtrue\global\lawboxhaspendingrnfalse\ifintranslation\lawboxtransfont\else\lawboxsourcefont\fi\RaggedRight\setlength{\parindent}{0pt}\setlength{\parskip}{0pt}\ifintranslation\setstretch{\lawboxtransstretch}\else\setstretch{\lawboxsourcestretch}\fi},
    after upper={\global\lawboxhaspendingrnfalse\global\inlawboxfalse},
    borderline west={0.28pt}{0pt}{sourceline!85!black},
    borderline east={0.28pt}{0pt}{sourceline!85!black},
  },
  lawboxpartfirst/.style={
    lawboxpartbase,
    top=0.62em,
    bottom=0.04em,
    before skip=\lawboxblockbeforeskip,
    borderline north={0.28pt}{0pt}{sourceline!85!black},
  },
  lawboxpartmiddle/.style={
    lawboxpartbase,
    top=0.04em,
    bottom=0.04em,
  },
  lawboxpartlast/.style={
    lawboxpartbase,
    top=0.04em,
    bottom=0.62em,
    after skip=0pt,
    borderline south={0.28pt}{0pt}{sourceline!85!black},
  }
}
\newtcolorbox{lawboxpartinner}[1]{#1}
\newtcolorbox{termboxinner}{
  caseboxbase,
  colframe=noteblue!70!black,
  colback=casepaper,
  before upper={\songfont\setlength{\parindent}{0pt}\setlength{\parskip}{0.35em}}
}
\newtcolorbox{deboxinner}{
  caseboxbase,
  colframe=notegray!72!black,
  colback=white,
  left=1.05em,
  right=1.05em,
  top=0.62em,
  bottom=0.62em,
  before upper={\global\inlawboxtrue\global\lawboxfirstitemtrue\global\lawboxhaspendingrnfalse\small\setlength{\parindent}{0pt}\setlength{\parskip}{0.35em}},
  after upper={\global\lawboxhaspendingrnfalse\global\inlawboxfalse}
}
\NewEnviron{lawbox}{
  \begin{lawboxinner}\BODY\end{lawboxinner}
}
\NewEnviron{lawboxpart}[2]{
  \begin{lawboxpartinner}{lawboxpart#1,equal height group=#2}\BODY\end{lawboxpartinner}%
  \ifstrequal{#1}{last}{\global\lawboxpartendedblocktrue}{}%
}
\NewEnviron{termbox}{
  \par\addvspace{0.55em}
  {\color{noteblue!70!black}\rule{\linewidth}{0.28pt}}\par
  \begingroup
    \songfont\small\color{caseink}
    \setlength{\parindent}{0pt}
    \setlength{\parskip}{0.24em}
    \BODY
    \par
  \endgroup
  {\color{noteblue!70!black}\rule{\linewidth}{0.28pt}}\par
  \addvspace{0.55em}
}
\NewEnviron{debox}{
  \begin{deboxinner}\BODY\end{deboxinner}
}
\providecommand{\paragraphmark}[1]{}
\newcommand{\lawboxtitleafter}{\addvspace{0.06em}}
\newcommand{\lawtitle}[1]{\lawboxprespace{0.22em}{\lawboxtransfont\bfseries\lawboxuserandnumber #1}\par\lawboxtitleafter}
\newcommand{\absno}[2]{\lawboxprespace{0.04em}\noindent\lawboxuserandnumber\hangindent=2.6em\hangafter=1\makebox[2.6em][l]{(#1)}#2\par}
\newcommand{\nrno}[2]{\lawboxprespace{0pt}\noindent\lawboxuserandnumber\hspace*{2.6em}\hangindent=4.4em\hangafter=1\makebox[1.8em][l]{#1.}#2\par}
\newcommand{\letterno}[2]{\lawboxprespace{0pt}\noindent\lawboxuserandnumber\hspace*{4.4em}\hangindent=6.0em\hangafter=1\makebox[1.6em][l]{#1)}#2\par}
\newcommand{\sourcelawtitle}[1]{\lawboxprespace{0.22em}{\lawboxsourcefont\bfseries\lawboxuserandnumber #1}\par\lawboxtitleafter}
\newcommand{\sourcelawsubtitle}[1]{\lawboxprespace{0.04em}{\lawboxsourcefont\bfseries\lawboxuserandnumber #1}\par\lawboxtitleafter}
\newcommand{\sourceabsno}[2]{\lawboxprespace{0.04em}\noindent\lawboxuserandnumber\hangindent=2.45em\hangafter=1\makebox[2.45em][l]{(#1)}#2\par}
\newcommand{\sourcenrno}[2]{\lawboxprespace{0pt}\noindent\lawboxuserandnumber\hspace*{2.45em}\hangindent=4.25em\hangafter=1\makebox[1.8em][l]{#1.}#2\par}
\newcommand{\sourceletterno}[2]{\lawboxprespace{0pt}\noindent\lawboxuserandnumber\hspace*{4.25em}\hangindent=5.85em\hangafter=1\makebox[1.6em][l]{#1)}#2\par}
\newcommand{\sourcequotepara}[1]{\lawboxprespace{0.04em}\noindent\lawboxuserandnumber\hangindent=1.2em\hangafter=1 #1\par}
\newcommand{\sourcelawline}[1]{\lawboxprespace{0.04em}\noindent\lawboxuserandnumber #1\par}
\newcommand{\lawline}[1]{\lawboxprespace{0.04em}\noindent\lawboxuserandnumber #1\par}
\newcommand{\pendingtranslation}{\par{\songfont\footnotesize\color{pendingink}（中文譯文待補。）}\par}
\newcommand{\tocpart}[1]{\par\addvspace{0.52em}{\footnotesize\bfseries\color{pendingink}#1\par}\addvspace{0.16em}}
\newcommand{\tocdashline}{\par\addvspace{0.48em}\noindent{\color{notegray}\xleaders\hbox to 0.82em{\hfil\rule{0.42em}{0.28pt}\hfil}\hskip\linewidth}\par\addvspace{0.38em}}
% \tocpanel{font-cmd}{content}：將目錄置中於所在欄寬（雙語欄適用）。
% A3 橫式使用 0.58\linewidth，A4 橫式使用 0.84\linewidth，
% 使兩種紙張下目錄欄內容實際寬度相近（均約 100–105 mm）。
\newcommand{\tocpanel}[2]{%
  \makebox[\linewidth][c]{%
    \ifx\bilingualpaper\bilingualpaperaThree
      \begin{minipage}{0.58\linewidth}%
    \else
      \begin{minipage}{0.84\linewidth}%
    \fi
    \toccolstyle
    #1#2%
    \end{minipage}%
  }%
}
% ===== 首頁簡目錄的兩層 description list 縮排 =====
% enumitem 的 description list 可把每一列拆成：
%   [label]  labelsep  body text
% 這裡的「起點」都以所在 minipage/欄位的左緣為 0。
%
% 核心規則：
%   labelindent = label 欄位從左緣開始的位置。
%   labelwidth  = label 欄位本身寬度；標籤太長（如 Abw. Meinung）就加大它。
%   labelsep    = label 欄位右緣到正文的距離。
%   leftmargin  = 正文 body text 的起點。判斷層級視覺時主要看這個。
%
% 所以：
%   想讓「Begründung / 判決理由」整列正文往右或往左：改 tocitems 的 leftmargin。
%   想讓「A. Verfahren und Gegenstand / A. 程序與審查標的」正文往右或往左：
%     改 tocsubitems 的 leftmargin。
%   想讓 A./B./C. 這些標籤本身往右或往左，但正文不動：改 tocsubitems 的 labelindent。
%   想讓 A./B./C. 與正文間距變大或變小：改 tocsubitems 的 labelsep。
%
% 層級原則：
%   tocsubitems.leftmargin 必須大於 tocitems.leftmargin，
%   否則子層級正文會看起來比「Gründe / 理由」還靠左。
\newlist{tocitems}{description}{1}
\setlist[tocitems]{%
  labelindent=0pt,    % 第一層 label 欄位起點；0pt 表示貼齊目錄欄左緣。
  leftmargin=2.54cm,  % 第一層正文起點；例如 Begründung / 判決理由 的左緣。
  labelwidth=2.16cm,  % 第一層 label 欄位寬度；需容納 Leitsätze、Abw. Meinung、不同意見等。
  labelsep=0.32cm,    % 第一層 label 與正文的水平間距；只改間距，不改正文起點。
  itemsep=0.28em,     % 第一層各項之間的垂直距離。
  topsep=0.20em,      % 第一層 list 上下與前後內容的垂直距離。
  parsep=0pt,         % 同一 item 內段落之間的距離；目錄項通常不需要。
  partopsep=0pt,      % list 若緊接段落開始時的額外距離；關閉以免目錄高度飄動。
  align=left,         % label 欄內左對齊；長短標籤不靠右貼齊。
  font=\normalfont\bfseries % label 的字型；正文不受這行影響。
}
\newlist{tocsubitems}{description}{1}
\setlist[tocsubitems]{%
  labelindent=1cm,    % 第二層 label 起點；只控制 A./B./C. 本身的位置。
  leftmargin=3.00cm,  % 第二層正文起點；必須大於 2.54cm，才會比 Begründung / 判決理由更右。
  labelwidth=0.5cm,   % 第二層 label 欄寬；A./B./C. 很短，可小於第一層。
  labelsep=1.5cm,    % 第二層 label 與正文間距；A. 與 Verfahren / 程序 的距離看這裡。
  itemsep=0.12em,     % 第二層各項較密，避免 A.–E. 佔太高。
  topsep=0.04em,      % 第二層 list 與 Gründe / 理由之間的垂直距離。
  parsep=0pt,         % 同一第二層 item 內段落距離；目錄項通常不需要。
  partopsep=0pt,      % 關閉額外段落起始距離，避免版面不穩。
  align=left,         % A./B./C. label 欄內左對齊。
  font=\normalfont\bfseries, % A./B./C. label 加粗；正文不受這行影響。
  before={}           % 不在第二層 list 前自動插入額外內容。
}
\newlist{termitems}{description}{1}
\ifbilingualcolumns
  \setlist[termitems]{%
    leftmargin=5.2cm,
    labelwidth=4.8cm,
    labelsep=0.35cm,
    itemsep=0.16em,
    topsep=0.2em,
    parsep=0pt,
    partopsep=0pt,
    font=\normalfont\bfseries,
    style=nextline
  }
\else
  \setlist[termitems]{%
    leftmargin=0pt,
    labelwidth=0pt,
    labelsep=0pt,
    itemsep=0.24em,
    topsep=0.2em,
    parsep=0pt,
    partopsep=0pt,
    font=\normalfont\bfseries,
    style=nextline
  }
\fi
\newlist{abbritems}{description}{1}
\setlist[abbritems]{%
  leftmargin=2.38cm,
  labelwidth=2.02cm,
  labelsep=0.28cm,
  itemsep=0.16em,
  topsep=0.2em,
  parsep=0pt,
  partopsep=0pt,
  align=left,
  font=\normalfont\bfseries
}
\ifbilingualcolumns
  \newenvironment{termcolumns}{\begin{multicols}{2}}{\end{multicols}}
\else
  \newenvironment{termcolumns}{}{}
\fi

% ===== 編者註腳開關 =====
% 一般版預設不顯示編者註，避免正文腳註過密。
% 若開啟 \classroommodeon，GG/條文編者註仍會自動顯示，無需另開本開關。
% 若一般版也要顯示編者註，可在單案檔 \input 後加 \showeditornotestrue。
\newif\ifshoweditornotes
\showeditornotesfalse

% ===== 目錄排版輔助 =====
\newcommand{\toccolstyle}{\raggedright\arraybackslash}
\newcommand{\tocsingleindent}{\leftskip=1.45cm\rightskip=1.2cm}

% ===== 行內引文環境（德文原文與中文譯文各一，帶左側灰色邊線）=====
\newcommand{\quoteparbreak}{\par\addvspace{0.48em}}
\newcommand{\sourcequoteinner}[1]{%
  \begin{tcolorbox}[
    enhanced, breakable, width=\dimexpr\linewidth-0.8em\relax, frame hidden,
    colback=white, coltext=caseink, boxrule=0pt,
    borderline west={0.55pt}{0.88em}{notegray},
    left=1.78em, right=0.72em, top=0.18em, bottom=0.18em,
    boxsep=0pt, before skip=0.24em, after skip=0.24em,
    before upper={\normalfont\footnotesize\color{caseink}\setlength{\parindent}{0pt}\setlength{\parskip}{0.34em}\raggedright\setlength{\rightskip}{0pt plus 1em}}
  ]%
  #1\par
  \end{tcolorbox}%
}
\newcommand{\transquoteinner}[1]{%
  \begin{tcolorbox}[
    enhanced, breakable, width=\dimexpr\linewidth-0.8em\relax, frame hidden,
    colback=white, coltext=caseink, boxrule=0pt,
    borderline west={0.55pt}{0.88em}{notegray},
    left=1.78em, right=0.72em, top=0.18em, bottom=0.18em,
    boxsep=0pt, before skip=0.24em, after skip=0.24em,
    before upper={\songfont\small\color{caseink}\setlength{\parindent}{0pt}\setlength{\parskip}{0.34em}\raggedright\setlength{\rightskip}{0pt plus 1em}}
  ]%
  #1\par
  \end{tcolorbox}%
}

% ===== 大綱（Gliederung）Rn 輔助命令 =====
\newcommand{\outrn}[2]{\enspace{\normalfont\footnotesize\color{pendingink}(Rn\,#1–#2)}}
\newcommand{\outrnsingle}[1]{\enspace{\normalfont\footnotesize\color{pendingink}(Rn\,#1)}}
\newcommand{\outrnzh}[2]{\enspace{\normalfont\footnotesize\color{pendingink}（第\,#1–#2\,段）}}
\newcommand{\outrnzhs}[1]{\enspace{\normalfont\footnotesize\color{pendingink}（第\,#1\,段）}}
% ===== 大綱雙語對照表格輔助命令（三欄：段碼 │ 德文 │ 中文）=====
% 欄 1：段碼（由欄定義自動格式化：小字、等寬、右對齊、pendingink）
% 欄 2：德文標籤＋正文
% 欄 3：中文標籤＋正文
%
% 無段碼的列（\omainrow、\osecrow），欄 1 以 & 空置。
% 有段碼的列（\osecrowrn、\osubrow、\ointrow），#1 為預格式化字串
%   例：\osubrow{2–49}{I.}{...}{I.}{...}
%       \ointrow{107}{...}{...}
%
% \opartrow：段組標題（橫跨三欄）
\newcommand{\opartrow}[2]{%
  \noalign{\vspace{0.30em}}%
  \multicolumn{3}{@{}l@{}}{%
    \footnotesize\bfseries\color{pendingink}#1\hfill\songfont#2%
  }\\[0.06em]%
}
% \odashrow：虛線分隔（橫跨三欄）
\newcommand{\odashrow}{%
  \noalign{\vspace{0.28em}}%
  \multicolumn{3}{@{}l@{}}{\color{notegray}%
    \xleaders\hbox to 0.82em{\hfil\rule{0.42em}{0.28pt}\hfil}\hskip 0.88\linewidth}%
  \\[0.24em]%
}
% \odissentpart：不同意見書區段標頭。比一般 \opartrow 更像附錄性轉場。
\newcommand{\odissentpart}[2]{%
  \noalign{\vspace{0.42em}}%
  \multicolumn{3}{@{}p{0.88\textwidth}@{}}{%
    \color{notegray}\rule{\linewidth}{0.18pt}%
  }\\[-0.18em]%
  \multicolumn{3}{@{}p{0.88\textwidth}@{}}{%
    \makebox[\linewidth][s]{%
      \footnotesize\bfseries\color{pendingink}#1%
      \hfill{\songfont #2}%
    }%
  }\\[-0.02em]%
}
% \odissentopinion：不同意見書的署名與一行摘要，右欄保留段碼範圍。
\newcommand{\odissentopinion}[5]{%
  \footnotesize\hangindent=0.52cm\hangafter=1%
    {\bfseries\color{caseink}#2}\enspace{\color{notegray}\textperiodcentered}\enspace\color{caseink}#3%
  &\footnotesize\hangindent=0.52cm\hangafter=1%
    {\songfont\bfseries\color{caseink}#4}\enspace{\color{notegray}\textperiodcentered}\enspace\songfont\color{caseink}#5%
  &#1%
  \\[0.12em]%
}
% \omainrow：頂層項目，無段碼（Leitsätze、Urteil、Gründe …）
\newcommand{\omainrow}[4]{%
  \footnotesize\hangindent=1.92cm\hangafter=1%
    \makebox[1.92cm][l]{\bfseries\color{caseink}#1}\color{caseink}#2%
  &\footnotesize\hangindent=1.92cm\hangafter=1%
    \makebox[1.92cm][l]{\bfseries\color{caseink}#3}\songfont\color{caseink}#4%
  &%
  \\[0.15em]%
}
% \omainrowrn：頂層項目，有段碼（不同意見等標籤寬於 \osecrow 框者）
\newcommand{\omainrowrn}[5]{%
  \footnotesize\hangindent=1.92cm\hangafter=1%
    \makebox[1.92cm][l]{\bfseries\color{caseink}#2}\color{caseink}#3%
  &\footnotesize\hangindent=1.92cm\hangafter=1%
    \makebox[1.92cm][l]{\bfseries\color{caseink}#4}\songfont\color{caseink}#5%
  &#1%
  \\[0.15em]%
}
% \osecrow：一級子項，無段碼（A.、C.、D. 等有子項者）
\newcommand{\osecrow}[4]{%
  \footnotesize\hspace{1.10cm}\hangindent=1.98cm\hangafter=1%
    \makebox[0.86cm][l]{\bfseries\color{caseink}#1}\color{caseink}#2%
  &\footnotesize\hspace{1.10cm}\hangindent=1.98cm\hangafter=1%
    \makebox[0.86cm][l]{\bfseries\color{caseink}#3}\songfont\color{caseink}#4%
  &%
  \\[0.11em]%
}
% \osecrowrn：一級子項，有段碼（B.、E. 等完整段落範圍者）
% #1=段碼字串（如 101–106）, #2=DE標籤, #3=DE正文, #4=ZH標籤, #5=ZH正文
\newcommand{\osecrowrn}[5]{%
  \footnotesize\hspace{1.10cm}\hangindent=1.98cm\hangafter=1%
    \makebox[0.86cm][l]{\bfseries\color{caseink}#2}\color{caseink}#3%
  &\footnotesize\hspace{1.10cm}\hangindent=1.98cm\hangafter=1%
    \makebox[0.86cm][l]{\bfseries\color{caseink}#4}\songfont\color{caseink}#5%
  &#1%
  \\[0.11em]%
}
% \osubrow：二級子項，有段碼範圍（I.、II.、A.I. …）
% #1=段碼字串, #2=DE標籤, #3=DE正文, #4=ZH標籤, #5=ZH正文
\newcommand{\osubrow}[5]{%
  \footnotesize\hspace{1.40cm}\hangindent=2.30cm\hangafter=1%
    \makebox[0.90cm][l]{\bfseries\color{caseink}#2}\color{caseink}#3%
  &\footnotesize\hspace{1.40cm}\hangindent=2.30cm\hangafter=1%
    \makebox[0.90cm][l]{\bfseries\color{caseink}#4}\songfont\color{caseink}#5%
  &#1%
  \\[0.09em]%
}
% \ointrow：引入段落，有單一段碼，無標籤
% #1=段碼字串, #2=DE正文, #3=ZH正文
\newcommand{\ointrow}[3]{%
  \footnotesize\hspace{0.52cm}\hangindent=0.52cm\hangafter=1\color{caseink}#2%
  &\footnotesize\hspace{0.52cm}\hangindent=0.52cm\hangafter=1\songfont\color{caseink}#3%
  &#1%
  \\[0.09em]%
}
% \osubsubrow：三級子項（1.、2.、3.）
% #1=段碼字串, #2=DE標籤, #3=DE正文, #4=ZH標籤, #5=ZH正文
\newcommand{\osubsubrow}[5]{%
  \footnotesize\hspace{1.80cm}\hangindent=2.48cm\hangafter=1%
    \makebox[0.68cm][l]{\bfseries\color{caseink}#2}\color{caseink}#3%
  &\footnotesize\hspace{1.80cm}\hangindent=2.48cm\hangafter=1%
    \makebox[0.68cm][l]{\bfseries\color{caseink}#4}\songfont\color{caseink}#5%
  &#1%
  \\[0.08em]%
}
% \osubsubsubrow：四級子項（a）、b））
% #1=段碼字串, #2=DE標籤, #3=DE正文, #4=ZH標籤, #5=ZH正文
\newcommand{\osubsubsubrow}[5]{%
  \footnotesize\hspace{2.10cm}\hangindent=2.68cm\hangafter=1%
    \makebox[0.58cm][l]{\bfseries\color{caseink}#2}\color{caseink}#3%
  &\footnotesize\hspace{2.10cm}\hangindent=2.68cm\hangafter=1%
    \makebox[0.58cm][l]{\bfseries\color{caseink}#4}\songfont\color{caseink}#5%
  &#1%
  \\[0.07em]%
}
% \osubsubsubsubrow：五級子項（aa）、bb））
% 標籤框 0.62cm：足以容納「aa)」在 footnotesize bold 下（約 0.50cm）並留有間距
% #1=段碼字串, #2=DE標籤, #3=DE正文, #4=ZH標籤, #5=ZH正文
\newcommand{\osubsubsubsubrow}[5]{%
  \footnotesize\hspace{2.38cm}\hangindent=3.06cm\hangafter=1%
    \makebox[0.68cm][l]{\bfseries\color{caseink}#2}\color{caseink}#3%
  &\footnotesize\hspace{2.38cm}\hangindent=3.06cm\hangafter=1%
    \makebox[0.68cm][l]{\bfseries\color{caseink}#4}\songfont\color{caseink}#5%
  &#1%
  \\[0.06em]%
}
% \outlinepage：雙語對照大綱頁包裝環境（longtable 自動跨頁）
% 三欄：德文欄 + 中文欄 + 段碼欄（最右，不折行）
% 段碼欄寬度：1.80cm，足以容納最長段碼如「309–348」
% DE/ZH 欄寬：(0.87\textwidth - 2.08cm) / 2
%   驗算：2×欄 + 0.05\textwidth + 0.28cm gap + 1.80cm = 0.87tw - 2.08cm + 0.05tw + 2.08cm = 0.92tw ✓
\newenvironment{outlinepage}{%
  \vspace*{0.40cm}%
  \setlength{\LTleft}{0.04\textwidth}%
  \setlength{\LTright}{0.04\textwidth}%
  \setlength{\tabcolsep}{0pt}%
  % 大綱列距由 longtable 的 \arraystretch 與各 row macro 結尾的 \\[...em] 共同決定。
  % 這裡控制「每一列本身的表格 strut 高度」；數值越大，A./I./1. 各項之間越像有空行。
  % A3 原本用 0.62，視覺上沒有空行；A4 也沿用同值，避免切到 A4 後項目被撐開。
  % 若日後覺得大綱太擠，只微調這個值（例如 0.68），不要改各 \omainrow/\osubrow 的縮排。
  \renewcommand{\arraystretch}{0.62}%
  \begin{longtable}{%
    >{\vspace{0pt}\raggedright\arraybackslash}p{\dimexpr(0.87\textwidth-2.08cm)/2\relax}%
    @{\hspace{0.025\textwidth}}%
    !{\color{notegray}\vrule width 0.12pt}%
    @{\hspace{0.025\textwidth}}%
    >{\vspace{0pt}\raggedright\arraybackslash}p{\dimexpr(0.87\textwidth-2.08cm)/2\relax}%
    @{\hspace{0.28cm}}%
    >{\vspace{0pt}\footnotesize\normalfont\color{pendingink}\raggedright\arraybackslash}p{1.80cm}%
    @{}%
  }%
  % 首頁欄標籤：不要橫跨置中；分別放在德文欄與中文欄的實際欄位起點。
  \footnotesize\bfseries\color{sourceink}Gliederung%
  &\footnotesize\bfseries\songfont\color{sourceink}大綱%
  &%
  \\[0.36em]%
  \endfirsthead%
  % 續頁欄標籤：同樣分欄定位，以灰色細體區分；無需標注「續」。
  \footnotesize\color{pendingink}Gliederung%
  &\footnotesize\songfont\color{pendingink}大綱%
  &%
  \\[0.24em]%
  \endhead%
}{%
  \end{longtable}%
}
 之前設定；
%   預設 chinese / a4）
\ifdefined\docmode\else\def\docmode{chinese}\fi
\ifdefined\bilingualpaper\else\def\bilingualpaper{a4}\fi
\newif\ifshoworiginal
\newif\ifshowtranslation
\newif\ifbilinguallandscape
\newif\ifbilingualcolumns
\newif\ifintranslation
\newif\iffirstbilingualsection

\showoriginalfalse
\showtranslationfalse
\bilinguallandscapefalse
\bilingualcolumnsfalse
\intranslationfalse
\firstbilingualsectionfalse

\def\modechinese{chinese}
\def\modegerman{german}
\def\modebilingual{bilingual}
\def\bilingualpaperaThree{a3}
\def\bilingualpaperaFour{a4}

\ifx\docmode\modechinese
  \showtranslationtrue
\fi

\ifx\docmode\modegerman
  \showoriginaltrue
\fi

\ifx\docmode\modebilingual
  \showoriginaltrue
  \showtranslationtrue
  \bilinguallandscapetrue
  \bilingualcolumnstrue
\fi

% 編排模式說明：
% chinese  → \showoriginalfalse  \showtranslationtrue  （A4 直式，純中文）
% german   → \showoriginaltrue   \showtranslationfalse （A4 直式，純德文）
% bilingual→ \showoriginaltrue   \showtranslationtrue  \bilingualcolumnstrue（A3/A4 橫式對照）

\ifbilingualcolumns
  \ifx\bilingualpaper\bilingualpaperaFour
    \geometry{
      a4paper,
      landscape,
      left=1.25cm,
      right=2.25cm,
      top=1.25cm,
      bottom=1.75cm,
      footskip=8.5mm,
      marginparwidth=0.75cm,
      marginparsep=0.25cm
    }
  \else
    \geometry{
      a3paper,
      landscape,
      left=1.55cm,
      right=2.75cm,
      top=1.55cm,
      bottom=2.05cm,
      footskip=9.5mm,
      marginparwidth=0.9cm,
      marginparsep=0.35cm
    }
  \fi
\else
\geometry{
  a4paper,
  portrait,
  left=2.2cm,
  right=2.6cm,
  top=2.0cm,
  bottom=2.0cm,
  marginparwidth=0.8cm,
  marginparsep=0.3cm
}
\fi
% ===== 時間戳基礎設施 =====
\newcount\stamptimeval
\newcount\stampminuteval
\newcommand{\stamphh}{%
  \stamptimeval=\time
  \divide\stamptimeval by 60
  \ifnum\stamptimeval<10 0\fi
  \the\stamptimeval}
\newcommand{\stampmm}{%
  \stamptimeval=\time
  \stampminuteval=\time
  \divide\stampminuteval by 60
  \multiply\stampminuteval by 60
  \advance\stamptimeval by -\stampminuteval
  \ifnum\stamptimeval<10 0\fi
  \the\stamptimeval}
\newcommand{\stampwkde}[1]{%
  \ifcase#1Montag\or Dienstag\or Mittwoch\or Donnerstag\or Freitag\or Samstag\or Sonntag\fi}
\newcommand{\stampwkzh}[1]{%
  \ifcase#1 一\or 二\or 三\or 四\or 五\or 六\or 日\fi}
\newcommand{\stampmonthde}[1]{%
  \ifcase#1\or Januar\or Februar\or März\or April\or Mai\or Juni\or
  Juli\or August\or September\or Oktober\or November\or Dezember\fi}
\newcount\stampjdval
\newcount\stampwdval
\pgfcalendardatetojulian{\the\year-\the\month-\the\day}{\stampjdval}
\pgfcalendarjuliantoweekday{\the\stampjdval}{\stampwdval}
\newcommand{\timestampzh}{%
  \songfont 最後修改：\the\year 年\the\month 月\the\day 日%
  % （星期\stampwkzh{\the\stampwdval}）%
  % \space\stamphh:\stampmm
  }

\newcommand{\documentdatezh}{\the\year 年 \the\month 月 \the\day 日}
\newcommand{\documentdatede}{\the\day. \stampmonthde{\the\month} \the\year}
\newcommand{\bverfgetranslationnote}{%
  {\songfont 翻譯說明：初譯使用 Claude Code 與 GPT 協助迭代；仍請以德文原文為準。}%
}
\newcommand{\bverfgecourseinfo}{%
  {\songfont 課程：114 學年度第 2 學期；憲法專題研究（四）；授課老師：湯德宗教授；導讀日期：115 年 6 月 1 日。}%
}
\newcommand{\bverfgeeditorinfo}{%
  {\songfont 編者：王逸帆（R10A21126），國立台灣大學法律系研究所碩士班。}%
}
\newcommand{\bverfgetitleversionline}{%
  \ifbilingualcolumns
    Fassung \documentversion\enspace\textperiodcentered\enspace Stand: \documentdatede
    \quad
    {\songfont 版本 \documentversion\enspace\textperiodcentered\enspace 修訂日期：\documentdatezh\par}%
  \else
    \ifshowtranslation
      {\songfont 版本 \documentversion\enspace\textperiodcentered\enspace 修訂日期：\documentdatezh\par}%
    \else
      Fassung \documentversion\enspace\textperiodcentered\enspace Stand: \documentdatede\par
    \fi
  \fi
}
\newcommand{\bverfgetitlecornerinfo}{%
  \ifshowtranslation
    \vspace{0.72em}%
    \noindent\hfill
    \begin{minipage}{0.66\textwidth}
      \raggedleft\tiny\color{pendingink}\songfont
      \bverfgecourseinfo\par
      \bverfgeeditorinfo
    \end{minipage}%
  \fi
}
\newcommand{\bverfgetitlemeta}{%
  \vfill
  \begin{center}%
    \footnotesize\color{pendingink}%
    \bverfgetitleversionline
    \ifshowtranslation
      \vspace{0.28em}{\scriptsize\bverfgetranslationnote\par}%
    \fi
  \end{center}%
  \bverfgetitlecornerinfo
}

\pagestyle{fancy}
\fancyhf{}
\renewcommand{\headrulewidth}{0pt}
\renewcommand{\footrulewidth}{0pt}
\fancyfoot[C]{\thepage}
\ifbilingualcolumns
  \fancyfoot[R]{\scriptsize\color{pendingink}\timestampzh}%
\else
  \ifshoworiginal
  \else
    \fancyfoot[R]{\scriptsize\color{pendingink}\timestampzh\hspace{0.45cm}}%
  \fi
\fi
\setmainfont{Times New Roman}
\setsansfont{Helvetica Neue}
\setmonofont{Menlo}
\IfFileExists{/Users/iw/Library/Fonts/kaiu.ttf}{
  \setCJKmainfont[Path=/Users/iw/Library/Fonts/,AutoFakeBold=4]{kaiu.ttf}
}{
  \IfFontExistsTF{標楷體}{
    \setCJKmainfont[AutoFakeBold=4]{標楷體}
  }{
    \setCJKmainfont{Songti TC}
  }
}
\IfFontExistsTF{Songti TC}{
  \setCJKfamilyfont{song}{Songti TC}
  \setCJKmonofont{Songti TC}
}{
  \setCJKfamilyfont{song}[Path=/Users/iw/Library/Fonts/,AutoFakeBold=4]{kaiu.ttf}
  \setCJKmonofont[Path=/Users/iw/Library/Fonts/,AutoFakeBold=4]{kaiu.ttf}
}
\newcommand{\songfont}{\CJKfamily{song}}

% ===== 封面／目錄專用打字機字體（僅限封面、目錄頁使用，不影響正文）=====
\IfFontExistsTF{Radio Newsman}{%
  \newfontfamily\bverfgetitlede{Radio Newsman}%
}{%
  \newfontfamily\bverfgetitlede{Courier New}%
}
\IfFontExistsTF{Erikas Farbband}{%
  \newfontfamily\bverfgebodyde{Erikas Farbband}%
}{%
  \let\bverfgebodyde\bverfgetitlede
}
\IfFontExistsTF{Maoken Heavy Labourer}{%
  \setCJKfamilyfont{bverfgetitlecjk}{Maoken Heavy Labourer}%
}{%
  \setCJKfamilyfont{bverfgetitlecjk}{Songti TC}%
}
\IfFontExistsTF{Huiwen-mincho}{%
  \setCJKfamilyfont{bverfgebodycjk}{Huiwen-mincho}%
}{%
  \setCJKfamilyfont{bverfgebodycjk}{Songti TC}%
}
\newcommand{\bverfgetitlezh}{\bverfgetitlede\CJKfamily{bverfgetitlecjk}}
\newcommand{\bverfgebodyzh}{\bverfgebodyde\CJKfamily{bverfgebodycjk}}

\setstretch{1.35}
\setlist{itemsep=0.2em, topsep=0.4em}
\setlength{\parindent}{0pt}
\setlength{\parskip}{0.45em}
\raggedbottom
\setcounter{secnumdepth}{0}
\setcounter{tocdepth}{2}

\newcommand{\lawboxsourcefont}{\footnotesize}
\newcommand{\lawboxtransfont}{\footnotesize}
\newcommand{\lawboxsourcestretch}{1.02}
\newcommand{\lawboxtransstretch}{0.92}

\titleformat{\section}{\centering\Large\bfseries\songfont}{\thesection}{1em}{}
\titleformat{\subsection}{\large\bfseries\songfont}{\thesubsection}{1em}{}
\titleformat{\subsubsection}{\normalsize\bfseries\songfont}{\thesubsubsection}{1em}{}
\titleformat{\paragraph}{\normalsize\bfseries\songfont}{\theparagraph}{1em}{}
\titlespacing*{\section}{0pt}{1.8em}{1.0em}
\titlespacing*{\subsection}{0pt}{1.2em}{0.6em}
\titlespacing*{\subsubsection}{0pt}{1.0em}{0.45em}
\titlespacing*{\paragraph}{0pt}{0.6em}{0.4em}

\definecolor{noteblue}{HTML}{A9C9F5}
\definecolor{noteteal}{HTML}{AEE1D6}
\definecolor{notegray}{HTML}{D7DADF}
\definecolor{caseink}{HTML}{000000}
\definecolor{casepaper}{HTML}{FCFEFD}
\definecolor{sourcepaper}{HTML}{FAFBF8}
\definecolor{sourceline}{HTML}{D7DED8}
\definecolor{sourceink}{HTML}{000000}
\definecolor{pendingink}{HTML}{000000}

\tcbset{
  caseboxbase/.style={
    enhanced,
    breakable,
    width=0.985\linewidth,
    colback=casepaper,
    coltitle=caseink,
    fonttitle=\bfseries\songfont,
    boxrule=0.28pt,
    arc=1.2mm,
    left=2.4mm,
    right=2.4mm,
    top=1.5mm,
    bottom=1.5mm,
    boxsep=1.5mm,
    before skip=0.65em,
    after skip=0.65em,
    before upper={\setlength{\parindent}{0pt}\setlength{\parskip}{0.35em}},
  }
}

\newcommand{\bilingualstart}{}
\newcommand{\bilingualstop}{}
\ifbilingualcolumns
  \renewcommand{\bilingualstart}{\small\setlength{\columnsep}{1.25cm}\setlength{\columnseprule}{0.12pt}\columnratio{0.5,0.5}\multicolumnfootnotes\flushbottom\sloppy\interlinepenalty=80\clubpenalty=800\widowpenalty=800\displaywidowpenalty=800\begin{paracol}{2}\global\intranslationfalse\global\firstbilingualsectiontrue{\footnotesize\color{sourceink}Deutsch\par}\switchcolumn{\footnotesize\songfont\color{sourceink}中文譯文\par}\switchcolumn*}
  \renewcommand{\bilingualstop}{\ifintranslation\switchcolumn*\global\intranslationfalse\fi\end{paracol}\clearpage\raggedbottom}
\fi
\newif\ifrandnumbers
\randnumbersfalse
\newcounter{randnumber}
\newcounter{randnumberlocal}
\newcounter{lastsourcerandstart}
\newif\iflastsourcehasrandnumbers
\lastsourcehasrandnumbersfalse
\newif\ifinlawbox
\inlawboxfalse
\newif\iflawboxhaspendingrn
\lawboxhaspendingrnfalse
\newcommand{\lawboxpendingrn}{}
\newcommand{\startrandnumbers}{\global\randnumberstrue\setcounter{randnumber}{1}\global\lastsourcehasrandnumbersfalse}
\newcommand{\stoprandnumbers}{\global\randnumbersfalse\global\lastsourcehasrandnumbersfalse}
\newlength{\randnumberwidth}
\setlength{\randnumberwidth}{0.95cm}
\newlength{\randnumberrightoffset}
\setlength{\randnumberrightoffset}{0.85cm}
\newlength{\randnumberlawboxrightoffset}
\setlength{\randnumberlawboxrightoffset}{1.40cm}
\newcommand{\randnumberbox}[1]{%
  \leavevmode
  \makebox[0pt][l]{%
    \smash{%
      \tikz[remember picture,overlay,baseline=(rnmark.base)]{%
        \coordinate (rnmark) at (0,0);%
        \ifinlawbox
          \path let \p1=(current page.east), \p2=(rnmark.base) in
            node[
              anchor=base east,
              inner sep=0pt,
              outer sep=0pt,
              text width=\randnumberwidth,
              align=right,
              text=pendingink,
              font=\normalfont\mdseries\scriptsize\ttfamily
            ] at (\x1-\randnumberlawboxrightoffset,\y2) {{\normalfont\mdseries\scriptsize\ttfamily #1}};%
        \else
          \path let \p1=(current page.east), \p2=(rnmark.base) in
            node[
              anchor=base east,
              inner sep=0pt,
              outer sep=0pt,
              text width=\randnumberwidth,
              align=right,
              text=pendingink,
              font=\normalfont\mdseries\scriptsize\ttfamily
            ] at (\x1-\randnumberrightoffset,\y2) {{\normalfont\mdseries\scriptsize\ttfamily #1}};%
        \fi
      }%
    }%
  }%
  \ignorespaces
}
\newcommand{\lawboxstorerandnumber}[1]{%
  \gdef\lawboxpendingrn{#1}%
  \global\lawboxhaspendingrntrue
  \ignorespaces
}
\newcommand{\lawboxuserandnumber}{%
  \iflawboxhaspendingrn
    \randnumberbox{\lawboxpendingrn}%
    \global\lawboxhaspendingrnfalse
  \fi
}
\newcommand{\sourcern}[1]{%
  \ifrandnumbers
    \ifshoworiginal
      \ifshowtranslation\else
        \ifinlawbox\lawboxstorerandnumber{#1}\else\randnumberbox{#1}\fi
      \fi
    \fi
  \fi
}
\newcommand{\transrn}[1]{%
  \ifrandnumbers
    \ifshowtranslation
      \ifinlawbox\lawboxstorerandnumber{#1}\else\randnumberbox{#1}\fi
    \fi
  \fi
}
% ===== 課堂模式：德文原文縮寫註解 =====
% 日常切換只改各判決檔的一行：
%   \classroommodeon        開啟課堂版；註解掉這行即回到一般版。
% 模板預設：
%   1. 課堂模式關閉。
%   2. 課堂模式開啟時，縮寫註解包含「德文全稱＋中文譯名」。
%   3. 課堂模式開啟時，GG/條文編者註仍會自動顯示。
% 進階微調才需要在單案檔覆寫：
%   \classroomabbrzhoff     縮寫註解僅顯示德文全稱。
%   \showeditornotestrue    一般版也顯示編者註。
% 註解策略：同一實際 PDF 頁面中，同一縮寫只在第一次出現處加註。
% 這裡用 zref-abspage 的絕對頁序，不用 \thepage，避免文件內頁碼重置時誤判。
\newif\ifclassroommode
\newif\ifclassroomabbrzh
\classroommodefalse
\classroomabbrzhtrue
\newcommand{\classroommodeon}{\classroommodetrue}
\newcommand{\classroommodeoff}{\classroommodefalse}
\newcommand{\classroomabbrzhon}{\classroomabbrzhtrue}
\newcommand{\classroomabbrzhoff}{\classroomabbrzhfalse}
\newcommand{\abbrnotelabel}{\emph{Abk.}: }
\newcommand{\abbrnote}[3]{%
  \ifclassroommode
    \ifshoworiginal
      \ifintranslation\else
        \footnote{\normalfont\footnotesize\abbrnotelabel #2\ifclassroomabbrzh；{\songfont #3}\fi}%
      \fi
    \fi
  \fi
}
\newcommand{\abbrnoteonce}[4]{%
  #2%
  \ifclassroommode
    \ifshoworiginal
      \ifintranslation\else
        \begingroup
          \edef\abbrcurrentabspage{\number\numexpr\value{abspage}+1\relax}%
          \ifcsname abbrnoteseen@#1@\abbrcurrentabspage\endcsname\else
            \global\expandafter\let\csname abbrnoteseen@#1@\abbrcurrentabspage\endcsname\empty
            \abbrnote{#1}{#3}{#4}%
          \fi
        \endgroup
      \fi
    \fi
  \fi
}
\newcommand{\abbrAbs}{\abbrnoteonce{abs}{Abs.}{Absatz}{項}}
\newcommand{\abbrAAO}{\abbrnoteonce{aao}{a.a.O.}{am angegebenen Ort}{前引處}}
\newcommand{\abbrAE}{\abbrnoteonce{ae}{AE}{Alternativ-Entwurf}{替代草案}}
\newcommand{\abbrArt}{\abbrnoteonce{art}{Art.}{Artikel}{條}}
\newcommand{\abbrAufl}{\abbrnoteonce{aufl}{Aufl.}{Auflage}{版}}
\newcommand{\abbrBd}{\abbrnoteonce{bd}{Bd.}{Band}{卷}}
\newcommand{\abbrBGBl}{\abbrnoteonce{bgbl}{BGBl.}{Bundesgesetzblatt}{聯邦法律公報}}
\newcommand{\abbrBGHSt}{\abbrnoteonce{bghst}{BGHSt}{Entscheidungen des Bundesgerichtshofes in Strafsachen}{聯邦普通法院刑事判決彙編}}
\newcommand{\abbrBRDrucks}{\abbrnoteonce{brdrucks}{BRDrucks.}{Bundesrats-Drucksache}{聯邦參議院文件}}
\newcommand{\abbrBTDrucks}{\abbrnoteonce{btdrucks}{BTDrucks.}{Bundestags-Drucksache}{聯邦議院文件}}
\newcommand{\abbrBundesgesetzbl}{\abbrnoteonce{bundesgesetzbl}{Bundesgesetzbl.}{Bundesgesetzblatt}{聯邦法律公報}}
\newcommand{\abbrBvF}{\abbrnoteonce{bvf}{BvF}{Bundesverfassungsgerichtsverfahren}{聯邦憲法法院程序案號}}
\newcommand{\abbrBvL}{\abbrnoteonce{bvl}{BvL}{Bundesverfassungsgerichtsverfahren}{聯邦憲法法院程序案號}}
\newcommand{\abbrBvR}{\abbrnoteonce{bvr}{BvR}{Bundesverfassungsgerichtsverfahren}{聯邦憲法法院程序案號}}
\newcommand{\abbrBVerfG}{\abbrnoteonce{bverfg}{BVerfG}{Bundesverfassungsgericht}{聯邦憲法法院}}
\newcommand{\abbrBVerfGE}{\abbrnoteonce{bverfge}{BVerfGE}{Entscheidungen des Bundesverfassungsgerichts}{聯邦憲法法院判決集}}
\newcommand{\abbrBVerfGG}{\abbrnoteonce{bverfgg}{BVerfGG}{Bundesverfassungsgerichtsgesetz}{聯邦憲法法院法}}
\newcommand{\abbrDDR}{\abbrnoteonce{ddr}{DDR}{Deutsche Demokratische Republik}{德意志民主共和國；東德}}
\newcommand{\abbrDJT}{\abbrnoteonce{djt}{DJT}{Deutscher Juristentag}{德國法學家大會}}
\newcommand{\abbrDr}{\abbrnoteonce{dr}{Dr.}{Doktor}{Dr.}}
\newcommand{\abbrEuGRZ}{\abbrnoteonce{eugrz}{EuGRZ}{Europäische Grundrechte-Zeitschrift}{歐洲基本權期刊}}
\newcommand{\abbrEV}{\abbrnoteonce{ev}{EV}{Einigungsvertrag}{統一條約}}
\newcommand{\abbrF}{\abbrnoteonce{f}{f.}{folgende}{以下一頁；下一段}}
\newcommand{\abbrFF}{\abbrnoteonce{ff}{ff.}{fortfolgende}{以下數頁；以下各段}}
\newcommand{\abbrGBlDDR}{\abbrnoteonce{gblddr}{GBl. DDR}{Gesetzblatt der Deutschen Demokratischen Republik}{德意志民主共和國法律公報}}
\newcommand{\abbrGG}{\abbrnoteonce{gg}{GG}{Grundgesetz}{基本法}}
\newcommand{\abbrGVBl}{\abbrnoteonce{gvbl}{GVBl.}{Gesetz- und Verordnungsblatt}{法律及命令公報}}
\newcommand{\abbrJR}{\abbrnoteonce{jr}{JR}{Juristische Rundschau}{法學評論}}
\newcommand{\abbrJZ}{\abbrnoteonce{jz}{JZ}{JuristenZeitung}{法學家報}}
\newcommand{\abbrIVm}{\abbrnoteonce{ivm}{i.V.m.}{in Verbindung mit}{結合；連同}}
\newcommand{\abbrMdB}{\abbrnoteonce{mdb}{MdB}{Mitglied des Bundestages}{聯邦議院議員}}
\newcommand{\abbrMwN}{\abbrnoteonce{mwn}{m.w.N.}{mit weiteren Nachweisen}{附進一步文獻／判例出處}}
\newcommand{\abbrNF}{\abbrnoteonce{nf}{n.F.}{neue Fassung}{新版本}}
\newcommand{\abbrAF}{\abbrnoteonce{af}{a.F.}{alte Fassung}{舊版本}}
\newcommand{\abbrNr}{\abbrnoteonce{nr}{Nr.}{Nummer}{款、目或號，依語境}}
\newcommand{\abbrProf}{\abbrnoteonce{prof}{Prof.}{Professor}{教授}}
\newcommand{\abbrRdnr}{\abbrnoteonce{rdnr}{Rdnr.}{Randnummer}{邊碼；段碼}}
\newcommand{\abbrRdnrn}{\abbrnoteonce{rdnrn}{Rdnrn.}{Randnummern}{邊碼；段碼}}
\newcommand{\abbrRGBl}{\abbrnoteonce{rgbl}{RGBl.}{Reichsgesetzblatt}{帝國法律公報}}
\newcommand{\abbrRGSt}{\abbrnoteonce{rgst}{RGSt}{Entscheidungen des Reichsgerichts in Strafsachen}{帝國法院刑事判決彙編}}
\newcommand{\abbrRspr}{\abbrnoteonce{rspr}{Rspr.}{Rechtsprechung}{判例；司法實務}}
\newcommand{\abbrRVO}{\abbrnoteonce{rvo}{RVO}{Reichsversicherungsordnung}{帝國保險法}}
\newcommand{\abbrS}{\abbrnoteonce{s}{S.}{Seite}{頁}}
\newcommand{\abbrSFHG}{\abbrnoteonce{sfhg}{SFHG}{Schwangeren- und Familienhilfegesetz}{孕婦及家庭扶助法}}
\newcommand{\abbrSGBV}{\abbrnoteonce{sgbv}{SGB V}{Fünftes Buch Sozialgesetzbuch}{社會法典第五編}}
\newcommand{\abbrStAG}{\abbrnoteonce{staeg}{StÄG}{Strafrechtsänderungsgesetz}{刑法修正法}}
\newcommand{\abbrStenBer}{\abbrnoteonce{stenber}{StenBer.}{Stenographischer Bericht}{速記紀錄}}
\newcommand{\abbrStGB}{\abbrnoteonce{stgb}{StGB}{Strafgesetzbuch}{刑法}}
\newcommand{\abbrStREG}{\abbrnoteonce{streg}{StREG}{Strafrechtsreform-Ergänzungsgesetz}{刑法改革補充法}}
\newcommand{\abbrStrRG}{\abbrnoteonce{strrg}{StrRG}{Strafrechtsreformgesetz}{刑法改革法}}
\newcommand{\abbrUS}{\abbrnoteonce{us}{U.S.}{United States Reports}{美國最高法院判決彙編}}
\newcommand{\abbrVgl}{\abbrnoteonce{vgl}{vgl.}{vergleiche}{參見}}
\newcommand{\abbrVVDStRL}{\abbrnoteonce{vvdstrl}{VVDStRL}{Veröffentlichungen der Vereinigung der Deutschen Staatsrechtslehrer}{德國公法教師協會論文集}}
\newcommand{\abbrWP}{\abbrnoteonce{wp}{WP.}{Wahlperiode}{屆期}}
\newcommand{\abbrWp}{\abbrnoteonce{wp}{Wp.}{Wahlperiode}{屆期}}
\newcommand{\abbrZStrW}{\abbrnoteonce{zstrw}{ZStrW}{Zeitschrift für die gesamte Strafrechtswissenschaft}{整體刑法學期刊}}
\newcommand{\recordsourcerandstart}{%
  \ifrandnumbers
    \setcounter{lastsourcerandstart}{\value{randnumber}}%
    \global\lastsourcehasrandnumberstrue
  \else
    \global\lastsourcehasrandnumbersfalse
  \fi
}
\newlength{\biparagraphskip}
\setlength{\biparagraphskip}{0.52em}
\newlength{\lawboxblockbeforeskip}
\setlength{\lawboxblockbeforeskip}{0.72em}
\newlength{\lawboxblockafterskip}
\setlength{\lawboxblockafterskip}{0.92em}
\newif\iflawboxpartendedblock
\lawboxpartendedblockfalse
\newcommand{\sourceparstyle}{\footnotesize\color{sourceink}\setstretch{1.12}\RaggedRight\emergencystretch=3em\setlength{\parskip}{0.42em}\obeylines\selectfont}
\newcommand{\transparstyle}{\songfont\color{caseink}\setstretch{1.12}\setlength{\parindent}{0pt}\setlength{\parskip}{0.42em}}
\newcommand{\bilingualpairskip}{%
  \par\addvspace{\biparagraphskip}%
  \switchcolumn
  \par\addvspace{\biparagraphskip}%
  \switchcolumn*%
  \global\intranslationfalse
}
\newcommand{\bilingualboxpairskip}{%
  \par
  \switchcolumn
  \par
  \switchcolumn*%
  \global\intranslationfalse
}
\newcommand{\bilinguallawboxblockskip}{%
  \switchcolumn*%
  \par\addvspace{\lawboxblockafterskip}%
  \switchcolumn
  \par\addvspace{\lawboxblockafterskip}%
  \switchcolumn*%
  \global\intranslationfalse
}
\newcommand{\bikeepwithnext}[1][4]{%
  \ifbilingualcolumns
    \syncleft
    \Needspace{#1\baselineskip}%
    \switchcolumn
    \Needspace{#1\baselineskip}%
    \switchcolumn*%
    \global\intranslationfalse
  \else
    \Needspace{#1\baselineskip}%
  \fi
}
\NewEnviron{sourcepara}[1][]{
  \recordsourcerandstart
  \ifshoworiginal
    \ifbilingualcolumns
      \ifintranslation\switchcolumn*\global\intranslationfalse\fi
    \else
      \Needspace{5\baselineskip}
    \fi
    \ifbilingualcolumns\else\par\addvspace{0.35em}\fi
    {\sourceparstyle
    \noindent\BODY\par}
    \ifbilingualcolumns\else\addvspace{0.25em}\fi
    \ifbilingualcolumns
      \switchcolumn
      \global\intranslationtrue
    \fi
  \else
    \ifrandnumbers
      \begingroup
        \setbox0=\vbox{\hsize=\linewidth\sourceparstyle\noindent\BODY\par}%
      \endgroup
    \fi
  \fi
}
\NewEnviron{transpara}{
  \ifshowtranslation
    \setcounter{randnumberlocal}{\value{lastsourcerandstart}}%
    \ifbilingualcolumns
      \ifintranslation\else\switchcolumn\global\intranslationtrue\fi
      {\transparstyle\setlength{\leftskip}{0.55em}\setlength{\rightskip}{0.15em plus 1em}\addtolength{\linewidth}{-0.55em}\BODY\par}
      \switchcolumn*
      \bilingualpairskip
    \else
      {\transparstyle\BODY\par}
    \fi
  \else
    \ifbilingualcolumns
      \ifintranslation\switchcolumn*\global\intranslationfalse\fi
    \fi
  \fi
}
\NewEnviron{sourceboxpara}[1][]{
  \global\lawboxpartendedblockfalse
  \recordsourcerandstart
  \ifshoworiginal
    \ifbilingualcolumns
      \ifintranslation\switchcolumn*\global\intranslationfalse\fi
    \else
      \Needspace{3\baselineskip}
    \fi
    {\sourceparstyle
    \BODY\par}
    \ifbilingualcolumns\else
      \iflawboxpartendedblock\addvspace{\lawboxblockafterskip}\fi
    \fi
    \ifbilingualcolumns
      \switchcolumn
      \global\intranslationtrue
    \fi
  \fi
}
\NewEnviron{transboxpara}{
  \global\lawboxpartendedblockfalse
  \ifshowtranslation
    \setcounter{randnumberlocal}{\value{lastsourcerandstart}}%
    \ifbilingualcolumns
      \ifintranslation\else\switchcolumn\global\intranslationtrue\fi
      {\transparstyle\BODY\par}
      \iflawboxpartendedblock
        \bilinguallawboxblockskip
      \else
        \switchcolumn*
        \bilingualboxpairskip
      \fi
    \else
      {\transparstyle\BODY\par}
      \iflawboxpartendedblock\addvspace{\lawboxblockafterskip}\fi
    \fi
  \else
    \ifbilingualcolumns
      \ifintranslation\switchcolumn*\global\intranslationfalse\fi
    \fi
  \fi
}
\newcommand{\leitsatznum}[1]{\noindent\hangindent=3.0em\hangafter=1\makebox[2.25em][r]{\bfseries #1.}\hspace{0.55em}\ignorespaces}
\newcommand{\syncleft}{\ifbilingualcolumns\ifintranslation\switchcolumn*\global\intranslationfalse\fi\fi}
\pretocmd{\section}{\syncleft\clearpage}{}{}
\pretocmd{\subsection}{\syncleft}{}{}
\pretocmd{\subsubsection}{\syncleft}{}{}
\newcommand{\biheadingfont}{\songfont}
\newcommand{\bitocentry}[2]{%
  \ifshoworiginal
    \ifshowtranslation
      \ifstrequal{#1}{#2}{#1}{#1\space #2}%
    \else
      #1%
    \fi
  \else
    #2%
  \fi
}
\pdfstringdefDisableCommands{%
  \def\bitocentry#1#2{#1}%
}
\newcommand{\biheading}[7]{%
  \syncleft#1\bikeepwithnext[#3]\phantomsection\addcontentsline{toc}{#2}{\bitocentry{#6}{#7}}%
  \ifbilingualcolumns
    {#4 #6\par}%
    \switchcolumn
    {#4\biheadingfont #7\par}%
    \switchcolumn*%
    \global\intranslationfalse
    \par\addvspace{#5}%
    \switchcolumn
    \par\addvspace{#5}%
    \switchcolumn*%
  \else
    \ifshoworiginal{#4 #6\par}\fi
    \ifshowtranslation{#4\biheadingfont #7\par}\fi
    \addvspace{#5}%
  \fi
}
\newcommand{\termsectionheading}[1]{%
  \par\Needspace{9\baselineskip}%
  \vspace{0.65em}%
  {\songfont\bfseries #1\par}%
  \vspace{0.2em}%
}
% 章節換頁：
%   雙語 paracol 欄內不要用 \clearpage；它會在同步兩欄時吐出只有欄線/頁碼的空頁。
%   \newpage 足以讓下一個 section 起新頁，又不會強制清空所有待排材料。
%   單語模式不在 paracol 內，仍可用 \clearpage。
\newcommand{\bisectionpagebreak}{\ifbilingualcolumns\newpage\else\clearpage\fi}
\newcommand{\bisection}[2]{%
  \biheading{\iffirstbilingualsection\global\firstbilingualsectionfalse\else\bisectionpagebreak\fi}{section}{6}{\centering\Large\bfseries}{0.8em}{#1}{#2}%
}
\newcommand{\bisubsection}[2]{%
  \biheading{}{subsection}{5}{\large\bfseries}{0.55em}{#1}{#2}%
}
\newcommand{\bisubsubsection}[2]{%
  \biheading{}{subsubsection}{4}{\normalsize\bfseries}{0.45em}{#1}{#2}%
}
\newcommand{\bicompositeheading}[4]{%
  \biheading{}{subsection}{5}{\large\bfseries}{0.16em}{#1}{#2}%
  \biheading{}{subsubsection}{3}{\normalsize\bfseries}{0.45em}{#3}{#4}%
}
\newif\iflawboxfirstitem
\lawboxfirstitemfalse
\newcommand{\lawboxprespace}[1]{%
  \par
  \iflawboxfirstitem
    \global\lawboxfirstitemfalse
  \else
    \addvspace{#1}%
  \fi
}
\newtcolorbox{lawboxinner}{
  caseboxbase,
  colframe=sourceline!85!black,
  colback=white,
  left=1.05em,
  right=1.05em,
  top=0.62em,
  bottom=0.62em,
  before upper={\global\inlawboxtrue\global\lawboxfirstitemtrue\global\lawboxhaspendingrnfalse\ifintranslation\lawboxtransfont\else\lawboxsourcefont\fi\RaggedRight\setlength{\parindent}{0pt}\setlength{\parskip}{0pt}\ifintranslation\setstretch{\lawboxtransstretch}\else\setstretch{\lawboxsourcestretch}\fi},
  after upper={\global\lawboxhaspendingrnfalse\global\inlawboxfalse}
}
\tcbset{
  lawboxpartbase/.style={
    enhanced,
    breakable=false,
    width=0.985\linewidth,
    colback=white,
    frame hidden,
    boxrule=0pt,
    arc=0pt,
    left=1.05em,
    right=1.05em,
    boxsep=1.5mm,
    before skip=0pt,
    after skip=0pt,
    before upper={\global\inlawboxtrue\global\lawboxfirstitemtrue\global\lawboxhaspendingrnfalse\ifintranslation\lawboxtransfont\else\lawboxsourcefont\fi\RaggedRight\setlength{\parindent}{0pt}\setlength{\parskip}{0pt}\ifintranslation\setstretch{\lawboxtransstretch}\else\setstretch{\lawboxsourcestretch}\fi},
    after upper={\global\lawboxhaspendingrnfalse\global\inlawboxfalse},
    borderline west={0.28pt}{0pt}{sourceline!85!black},
    borderline east={0.28pt}{0pt}{sourceline!85!black},
  },
  lawboxpartfirst/.style={
    lawboxpartbase,
    top=0.62em,
    bottom=0.04em,
    before skip=\lawboxblockbeforeskip,
    borderline north={0.28pt}{0pt}{sourceline!85!black},
  },
  lawboxpartmiddle/.style={
    lawboxpartbase,
    top=0.04em,
    bottom=0.04em,
  },
  lawboxpartlast/.style={
    lawboxpartbase,
    top=0.04em,
    bottom=0.62em,
    after skip=0pt,
    borderline south={0.28pt}{0pt}{sourceline!85!black},
  }
}
\newtcolorbox{lawboxpartinner}[1]{#1}
\newtcolorbox{termboxinner}{
  caseboxbase,
  colframe=noteblue!70!black,
  colback=casepaper,
  before upper={\songfont\setlength{\parindent}{0pt}\setlength{\parskip}{0.35em}}
}
\newtcolorbox{deboxinner}{
  caseboxbase,
  colframe=notegray!72!black,
  colback=white,
  left=1.05em,
  right=1.05em,
  top=0.62em,
  bottom=0.62em,
  before upper={\global\inlawboxtrue\global\lawboxfirstitemtrue\global\lawboxhaspendingrnfalse\small\setlength{\parindent}{0pt}\setlength{\parskip}{0.35em}},
  after upper={\global\lawboxhaspendingrnfalse\global\inlawboxfalse}
}
\NewEnviron{lawbox}{
  \begin{lawboxinner}\BODY\end{lawboxinner}
}
\NewEnviron{lawboxpart}[2]{
  \begin{lawboxpartinner}{lawboxpart#1,equal height group=#2}\BODY\end{lawboxpartinner}%
  \ifstrequal{#1}{last}{\global\lawboxpartendedblocktrue}{}%
}
\NewEnviron{termbox}{
  \par\addvspace{0.55em}
  {\color{noteblue!70!black}\rule{\linewidth}{0.28pt}}\par
  \begingroup
    \songfont\small\color{caseink}
    \setlength{\parindent}{0pt}
    \setlength{\parskip}{0.24em}
    \BODY
    \par
  \endgroup
  {\color{noteblue!70!black}\rule{\linewidth}{0.28pt}}\par
  \addvspace{0.55em}
}
\NewEnviron{debox}{
  \begin{deboxinner}\BODY\end{deboxinner}
}
\providecommand{\paragraphmark}[1]{}
\newcommand{\lawboxtitleafter}{\addvspace{0.06em}}
\newcommand{\lawtitle}[1]{\lawboxprespace{0.22em}{\lawboxtransfont\bfseries\lawboxuserandnumber #1}\par\lawboxtitleafter}
\newcommand{\absno}[2]{\lawboxprespace{0.04em}\noindent\lawboxuserandnumber\hangindent=2.6em\hangafter=1\makebox[2.6em][l]{(#1)}#2\par}
\newcommand{\nrno}[2]{\lawboxprespace{0pt}\noindent\lawboxuserandnumber\hspace*{2.6em}\hangindent=4.4em\hangafter=1\makebox[1.8em][l]{#1.}#2\par}
\newcommand{\letterno}[2]{\lawboxprespace{0pt}\noindent\lawboxuserandnumber\hspace*{4.4em}\hangindent=6.0em\hangafter=1\makebox[1.6em][l]{#1)}#2\par}
\newcommand{\sourcelawtitle}[1]{\lawboxprespace{0.22em}{\lawboxsourcefont\bfseries\lawboxuserandnumber #1}\par\lawboxtitleafter}
\newcommand{\sourcelawsubtitle}[1]{\lawboxprespace{0.04em}{\lawboxsourcefont\bfseries\lawboxuserandnumber #1}\par\lawboxtitleafter}
\newcommand{\sourceabsno}[2]{\lawboxprespace{0.04em}\noindent\lawboxuserandnumber\hangindent=2.45em\hangafter=1\makebox[2.45em][l]{(#1)}#2\par}
\newcommand{\sourcenrno}[2]{\lawboxprespace{0pt}\noindent\lawboxuserandnumber\hspace*{2.45em}\hangindent=4.25em\hangafter=1\makebox[1.8em][l]{#1.}#2\par}
\newcommand{\sourceletterno}[2]{\lawboxprespace{0pt}\noindent\lawboxuserandnumber\hspace*{4.25em}\hangindent=5.85em\hangafter=1\makebox[1.6em][l]{#1)}#2\par}
\newcommand{\sourcequotepara}[1]{\lawboxprespace{0.04em}\noindent\lawboxuserandnumber\hangindent=1.2em\hangafter=1 #1\par}
\newcommand{\sourcelawline}[1]{\lawboxprespace{0.04em}\noindent\lawboxuserandnumber #1\par}
\newcommand{\lawline}[1]{\lawboxprespace{0.04em}\noindent\lawboxuserandnumber #1\par}
\newcommand{\pendingtranslation}{\par{\songfont\footnotesize\color{pendingink}（中文譯文待補。）}\par}
\newcommand{\tocpart}[1]{\par\addvspace{0.52em}{\footnotesize\bfseries\color{pendingink}#1\par}\addvspace{0.16em}}
\newcommand{\tocdashline}{\par\addvspace{0.48em}\noindent{\color{notegray}\xleaders\hbox to 0.82em{\hfil\rule{0.42em}{0.28pt}\hfil}\hskip\linewidth}\par\addvspace{0.38em}}
% \tocpanel{font-cmd}{content}：將目錄置中於所在欄寬（雙語欄適用）。
% A3 橫式使用 0.58\linewidth，A4 橫式使用 0.84\linewidth，
% 使兩種紙張下目錄欄內容實際寬度相近（均約 100–105 mm）。
\newcommand{\tocpanel}[2]{%
  \makebox[\linewidth][c]{%
    \ifx\bilingualpaper\bilingualpaperaThree
      \begin{minipage}{0.58\linewidth}%
    \else
      \begin{minipage}{0.84\linewidth}%
    \fi
    \toccolstyle
    #1#2%
    \end{minipage}%
  }%
}
% ===== 首頁簡目錄的兩層 description list 縮排 =====
% enumitem 的 description list 可把每一列拆成：
%   [label]  labelsep  body text
% 這裡的「起點」都以所在 minipage/欄位的左緣為 0。
%
% 核心規則：
%   labelindent = label 欄位從左緣開始的位置。
%   labelwidth  = label 欄位本身寬度；標籤太長（如 Abw. Meinung）就加大它。
%   labelsep    = label 欄位右緣到正文的距離。
%   leftmargin  = 正文 body text 的起點。判斷層級視覺時主要看這個。
%
% 所以：
%   想讓「Begründung / 判決理由」整列正文往右或往左：改 tocitems 的 leftmargin。
%   想讓「A. Verfahren und Gegenstand / A. 程序與審查標的」正文往右或往左：
%     改 tocsubitems 的 leftmargin。
%   想讓 A./B./C. 這些標籤本身往右或往左，但正文不動：改 tocsubitems 的 labelindent。
%   想讓 A./B./C. 與正文間距變大或變小：改 tocsubitems 的 labelsep。
%
% 層級原則：
%   tocsubitems.leftmargin 必須大於 tocitems.leftmargin，
%   否則子層級正文會看起來比「Gründe / 理由」還靠左。
\newlist{tocitems}{description}{1}
\setlist[tocitems]{%
  labelindent=0pt,    % 第一層 label 欄位起點；0pt 表示貼齊目錄欄左緣。
  leftmargin=2.54cm,  % 第一層正文起點；例如 Begründung / 判決理由 的左緣。
  labelwidth=2.16cm,  % 第一層 label 欄位寬度；需容納 Leitsätze、Abw. Meinung、不同意見等。
  labelsep=0.32cm,    % 第一層 label 與正文的水平間距；只改間距，不改正文起點。
  itemsep=0.28em,     % 第一層各項之間的垂直距離。
  topsep=0.20em,      % 第一層 list 上下與前後內容的垂直距離。
  parsep=0pt,         % 同一 item 內段落之間的距離；目錄項通常不需要。
  partopsep=0pt,      % list 若緊接段落開始時的額外距離；關閉以免目錄高度飄動。
  align=left,         % label 欄內左對齊；長短標籤不靠右貼齊。
  font=\normalfont\bfseries % label 的字型；正文不受這行影響。
}
\newlist{tocsubitems}{description}{1}
\setlist[tocsubitems]{%
  labelindent=1cm,    % 第二層 label 起點；只控制 A./B./C. 本身的位置。
  leftmargin=3.00cm,  % 第二層正文起點；必須大於 2.54cm，才會比 Begründung / 判決理由更右。
  labelwidth=0.5cm,   % 第二層 label 欄寬；A./B./C. 很短，可小於第一層。
  labelsep=1.5cm,    % 第二層 label 與正文間距；A. 與 Verfahren / 程序 的距離看這裡。
  itemsep=0.12em,     % 第二層各項較密，避免 A.–E. 佔太高。
  topsep=0.04em,      % 第二層 list 與 Gründe / 理由之間的垂直距離。
  parsep=0pt,         % 同一第二層 item 內段落距離；目錄項通常不需要。
  partopsep=0pt,      % 關閉額外段落起始距離，避免版面不穩。
  align=left,         % A./B./C. label 欄內左對齊。
  font=\normalfont\bfseries, % A./B./C. label 加粗；正文不受這行影響。
  before={}           % 不在第二層 list 前自動插入額外內容。
}
\newlist{termitems}{description}{1}
\ifbilingualcolumns
  \setlist[termitems]{%
    leftmargin=5.2cm,
    labelwidth=4.8cm,
    labelsep=0.35cm,
    itemsep=0.16em,
    topsep=0.2em,
    parsep=0pt,
    partopsep=0pt,
    font=\normalfont\bfseries,
    style=nextline
  }
\else
  \setlist[termitems]{%
    leftmargin=0pt,
    labelwidth=0pt,
    labelsep=0pt,
    itemsep=0.24em,
    topsep=0.2em,
    parsep=0pt,
    partopsep=0pt,
    font=\normalfont\bfseries,
    style=nextline
  }
\fi
\newlist{abbritems}{description}{1}
\setlist[abbritems]{%
  leftmargin=2.38cm,
  labelwidth=2.02cm,
  labelsep=0.28cm,
  itemsep=0.16em,
  topsep=0.2em,
  parsep=0pt,
  partopsep=0pt,
  align=left,
  font=\normalfont\bfseries
}
\ifbilingualcolumns
  \newenvironment{termcolumns}{\begin{multicols}{2}}{\end{multicols}}
\else
  \newenvironment{termcolumns}{}{}
\fi

% ===== 編者註腳開關 =====
% 一般版預設不顯示編者註，避免正文腳註過密。
% 若開啟 \classroommodeon，GG/條文編者註仍會自動顯示，無需另開本開關。
% 若一般版也要顯示編者註，可在單案檔 \input 後加 \showeditornotestrue。
\newif\ifshoweditornotes
\showeditornotesfalse

% ===== 目錄排版輔助 =====
\newcommand{\toccolstyle}{\raggedright\arraybackslash}
\newcommand{\tocsingleindent}{\leftskip=1.45cm\rightskip=1.2cm}

% ===== 行內引文環境（德文原文與中文譯文各一，帶左側灰色邊線）=====
\newcommand{\quoteparbreak}{\par\addvspace{0.48em}}
\newcommand{\sourcequoteinner}[1]{%
  \begin{tcolorbox}[
    enhanced, breakable, width=\dimexpr\linewidth-0.8em\relax, frame hidden,
    colback=white, coltext=caseink, boxrule=0pt,
    borderline west={0.55pt}{0.88em}{notegray},
    left=1.78em, right=0.72em, top=0.18em, bottom=0.18em,
    boxsep=0pt, before skip=0.24em, after skip=0.24em,
    before upper={\normalfont\footnotesize\color{caseink}\setlength{\parindent}{0pt}\setlength{\parskip}{0.34em}\raggedright\setlength{\rightskip}{0pt plus 1em}}
  ]%
  #1\par
  \end{tcolorbox}%
}
\newcommand{\transquoteinner}[1]{%
  \begin{tcolorbox}[
    enhanced, breakable, width=\dimexpr\linewidth-0.8em\relax, frame hidden,
    colback=white, coltext=caseink, boxrule=0pt,
    borderline west={0.55pt}{0.88em}{notegray},
    left=1.78em, right=0.72em, top=0.18em, bottom=0.18em,
    boxsep=0pt, before skip=0.24em, after skip=0.24em,
    before upper={\songfont\small\color{caseink}\setlength{\parindent}{0pt}\setlength{\parskip}{0.34em}\raggedright\setlength{\rightskip}{0pt plus 1em}}
  ]%
  #1\par
  \end{tcolorbox}%
}

% ===== 大綱（Gliederung）Rn 輔助命令 =====
\newcommand{\outrn}[2]{\enspace{\normalfont\footnotesize\color{pendingink}(Rn\,#1–#2)}}
\newcommand{\outrnsingle}[1]{\enspace{\normalfont\footnotesize\color{pendingink}(Rn\,#1)}}
\newcommand{\outrnzh}[2]{\enspace{\normalfont\footnotesize\color{pendingink}（第\,#1–#2\,段）}}
\newcommand{\outrnzhs}[1]{\enspace{\normalfont\footnotesize\color{pendingink}（第\,#1\,段）}}
% ===== 大綱雙語對照表格輔助命令（三欄：段碼 │ 德文 │ 中文）=====
% 欄 1：段碼（由欄定義自動格式化：小字、等寬、右對齊、pendingink）
% 欄 2：德文標籤＋正文
% 欄 3：中文標籤＋正文
%
% 無段碼的列（\omainrow、\osecrow），欄 1 以 & 空置。
% 有段碼的列（\osecrowrn、\osubrow、\ointrow），#1 為預格式化字串
%   例：\osubrow{2–49}{I.}{...}{I.}{...}
%       \ointrow{107}{...}{...}
%
% \opartrow：段組標題（橫跨三欄）
\newcommand{\opartrow}[2]{%
  \noalign{\vspace{0.30em}}%
  \multicolumn{3}{@{}l@{}}{%
    \footnotesize\bfseries\color{pendingink}#1\hfill\songfont#2%
  }\\[0.06em]%
}
% \odashrow：虛線分隔（橫跨三欄）
\newcommand{\odashrow}{%
  \noalign{\vspace{0.28em}}%
  \multicolumn{3}{@{}l@{}}{\color{notegray}%
    \xleaders\hbox to 0.82em{\hfil\rule{0.42em}{0.28pt}\hfil}\hskip 0.88\linewidth}%
  \\[0.24em]%
}
% \odissentpart：不同意見書區段標頭。比一般 \opartrow 更像附錄性轉場。
\newcommand{\odissentpart}[2]{%
  \noalign{\vspace{0.42em}}%
  \multicolumn{3}{@{}p{0.88\textwidth}@{}}{%
    \color{notegray}\rule{\linewidth}{0.18pt}%
  }\\[-0.18em]%
  \multicolumn{3}{@{}p{0.88\textwidth}@{}}{%
    \makebox[\linewidth][s]{%
      \footnotesize\bfseries\color{pendingink}#1%
      \hfill{\songfont #2}%
    }%
  }\\[-0.02em]%
}
% \odissentopinion：不同意見書的署名與一行摘要，右欄保留段碼範圍。
\newcommand{\odissentopinion}[5]{%
  \footnotesize\hangindent=0.52cm\hangafter=1%
    {\bfseries\color{caseink}#2}\enspace{\color{notegray}\textperiodcentered}\enspace\color{caseink}#3%
  &\footnotesize\hangindent=0.52cm\hangafter=1%
    {\songfont\bfseries\color{caseink}#4}\enspace{\color{notegray}\textperiodcentered}\enspace\songfont\color{caseink}#5%
  &#1%
  \\[0.12em]%
}
% \omainrow：頂層項目，無段碼（Leitsätze、Urteil、Gründe …）
\newcommand{\omainrow}[4]{%
  \footnotesize\hangindent=1.92cm\hangafter=1%
    \makebox[1.92cm][l]{\bfseries\color{caseink}#1}\color{caseink}#2%
  &\footnotesize\hangindent=1.92cm\hangafter=1%
    \makebox[1.92cm][l]{\bfseries\color{caseink}#3}\songfont\color{caseink}#4%
  &%
  \\[0.15em]%
}
% \omainrowrn：頂層項目，有段碼（不同意見等標籤寬於 \osecrow 框者）
\newcommand{\omainrowrn}[5]{%
  \footnotesize\hangindent=1.92cm\hangafter=1%
    \makebox[1.92cm][l]{\bfseries\color{caseink}#2}\color{caseink}#3%
  &\footnotesize\hangindent=1.92cm\hangafter=1%
    \makebox[1.92cm][l]{\bfseries\color{caseink}#4}\songfont\color{caseink}#5%
  &#1%
  \\[0.15em]%
}
% \osecrow：一級子項，無段碼（A.、C.、D. 等有子項者）
\newcommand{\osecrow}[4]{%
  \footnotesize\hspace{1.10cm}\hangindent=1.98cm\hangafter=1%
    \makebox[0.86cm][l]{\bfseries\color{caseink}#1}\color{caseink}#2%
  &\footnotesize\hspace{1.10cm}\hangindent=1.98cm\hangafter=1%
    \makebox[0.86cm][l]{\bfseries\color{caseink}#3}\songfont\color{caseink}#4%
  &%
  \\[0.11em]%
}
% \osecrowrn：一級子項，有段碼（B.、E. 等完整段落範圍者）
% #1=段碼字串（如 101–106）, #2=DE標籤, #3=DE正文, #4=ZH標籤, #5=ZH正文
\newcommand{\osecrowrn}[5]{%
  \footnotesize\hspace{1.10cm}\hangindent=1.98cm\hangafter=1%
    \makebox[0.86cm][l]{\bfseries\color{caseink}#2}\color{caseink}#3%
  &\footnotesize\hspace{1.10cm}\hangindent=1.98cm\hangafter=1%
    \makebox[0.86cm][l]{\bfseries\color{caseink}#4}\songfont\color{caseink}#5%
  &#1%
  \\[0.11em]%
}
% \osubrow：二級子項，有段碼範圍（I.、II.、A.I. …）
% #1=段碼字串, #2=DE標籤, #3=DE正文, #4=ZH標籤, #5=ZH正文
\newcommand{\osubrow}[5]{%
  \footnotesize\hspace{1.40cm}\hangindent=2.30cm\hangafter=1%
    \makebox[0.90cm][l]{\bfseries\color{caseink}#2}\color{caseink}#3%
  &\footnotesize\hspace{1.40cm}\hangindent=2.30cm\hangafter=1%
    \makebox[0.90cm][l]{\bfseries\color{caseink}#4}\songfont\color{caseink}#5%
  &#1%
  \\[0.09em]%
}
% \ointrow：引入段落，有單一段碼，無標籤
% #1=段碼字串, #2=DE正文, #3=ZH正文
\newcommand{\ointrow}[3]{%
  \footnotesize\hspace{0.52cm}\hangindent=0.52cm\hangafter=1\color{caseink}#2%
  &\footnotesize\hspace{0.52cm}\hangindent=0.52cm\hangafter=1\songfont\color{caseink}#3%
  &#1%
  \\[0.09em]%
}
% \osubsubrow：三級子項（1.、2.、3.）
% #1=段碼字串, #2=DE標籤, #3=DE正文, #4=ZH標籤, #5=ZH正文
\newcommand{\osubsubrow}[5]{%
  \footnotesize\hspace{1.80cm}\hangindent=2.48cm\hangafter=1%
    \makebox[0.68cm][l]{\bfseries\color{caseink}#2}\color{caseink}#3%
  &\footnotesize\hspace{1.80cm}\hangindent=2.48cm\hangafter=1%
    \makebox[0.68cm][l]{\bfseries\color{caseink}#4}\songfont\color{caseink}#5%
  &#1%
  \\[0.08em]%
}
% \osubsubsubrow：四級子項（a）、b））
% #1=段碼字串, #2=DE標籤, #3=DE正文, #4=ZH標籤, #5=ZH正文
\newcommand{\osubsubsubrow}[5]{%
  \footnotesize\hspace{2.10cm}\hangindent=2.68cm\hangafter=1%
    \makebox[0.58cm][l]{\bfseries\color{caseink}#2}\color{caseink}#3%
  &\footnotesize\hspace{2.10cm}\hangindent=2.68cm\hangafter=1%
    \makebox[0.58cm][l]{\bfseries\color{caseink}#4}\songfont\color{caseink}#5%
  &#1%
  \\[0.07em]%
}
% \osubsubsubsubrow：五級子項（aa）、bb））
% 標籤框 0.62cm：足以容納「aa)」在 footnotesize bold 下（約 0.50cm）並留有間距
% #1=段碼字串, #2=DE標籤, #3=DE正文, #4=ZH標籤, #5=ZH正文
\newcommand{\osubsubsubsubrow}[5]{%
  \footnotesize\hspace{2.38cm}\hangindent=3.06cm\hangafter=1%
    \makebox[0.68cm][l]{\bfseries\color{caseink}#2}\color{caseink}#3%
  &\footnotesize\hspace{2.38cm}\hangindent=3.06cm\hangafter=1%
    \makebox[0.68cm][l]{\bfseries\color{caseink}#4}\songfont\color{caseink}#5%
  &#1%
  \\[0.06em]%
}
% \outlinepage：雙語對照大綱頁包裝環境（longtable 自動跨頁）
% 三欄：德文欄 + 中文欄 + 段碼欄（最右，不折行）
% 段碼欄寬度：1.80cm，足以容納最長段碼如「309–348」
% DE/ZH 欄寬：(0.87\textwidth - 2.08cm) / 2
%   驗算：2×欄 + 0.05\textwidth + 0.28cm gap + 1.80cm = 0.87tw - 2.08cm + 0.05tw + 2.08cm = 0.92tw ✓
\newenvironment{outlinepage}{%
  \vspace*{0.40cm}%
  \setlength{\LTleft}{0.04\textwidth}%
  \setlength{\LTright}{0.04\textwidth}%
  \setlength{\tabcolsep}{0pt}%
  % 大綱列距由 longtable 的 \arraystretch 與各 row macro 結尾的 \\[...em] 共同決定。
  % 這裡控制「每一列本身的表格 strut 高度」；數值越大，A./I./1. 各項之間越像有空行。
  % A3 原本用 0.62，視覺上沒有空行；A4 也沿用同值，避免切到 A4 後項目被撐開。
  % 若日後覺得大綱太擠，只微調這個值（例如 0.68），不要改各 \omainrow/\osubrow 的縮排。
  \renewcommand{\arraystretch}{0.62}%
  \begin{longtable}{%
    >{\vspace{0pt}\raggedright\arraybackslash}p{\dimexpr(0.87\textwidth-2.08cm)/2\relax}%
    @{\hspace{0.025\textwidth}}%
    !{\color{notegray}\vrule width 0.12pt}%
    @{\hspace{0.025\textwidth}}%
    >{\vspace{0pt}\raggedright\arraybackslash}p{\dimexpr(0.87\textwidth-2.08cm)/2\relax}%
    @{\hspace{0.28cm}}%
    >{\vspace{0pt}\footnotesize\normalfont\color{pendingink}\raggedright\arraybackslash}p{1.80cm}%
    @{}%
  }%
  % 首頁欄標籤：不要橫跨置中；分別放在德文欄與中文欄的實際欄位起點。
  \footnotesize\bfseries\color{sourceink}Gliederung%
  &\footnotesize\bfseries\songfont\color{sourceink}大綱%
  &%
  \\[0.36em]%
  \endfirsthead%
  % 續頁欄標籤：同樣分欄定位，以灰色細體區分；無需標注「續」。
  \footnotesize\color{pendingink}Gliederung%
  &\footnotesize\songfont\color{pendingink}大綱%
  &%
  \\[0.24em]%
  \endhead%
}{%
  \end{longtable}%
}
 之前設定；
%   預設 chinese / a4）
\ifdefined\docmode\else\def\docmode{chinese}\fi
\ifdefined\bilingualpaper\else\def\bilingualpaper{a4}\fi
\newif\ifshoworiginal
\newif\ifshowtranslation
\newif\ifbilinguallandscape
\newif\ifbilingualcolumns
\newif\ifintranslation
\newif\iffirstbilingualsection

\showoriginalfalse
\showtranslationfalse
\bilinguallandscapefalse
\bilingualcolumnsfalse
\intranslationfalse
\firstbilingualsectionfalse

\def\modechinese{chinese}
\def\modegerman{german}
\def\modebilingual{bilingual}
\def\bilingualpaperaThree{a3}
\def\bilingualpaperaFour{a4}

\ifx\docmode\modechinese
  \showtranslationtrue
\fi

\ifx\docmode\modegerman
  \showoriginaltrue
\fi

\ifx\docmode\modebilingual
  \showoriginaltrue
  \showtranslationtrue
  \bilinguallandscapetrue
  \bilingualcolumnstrue
\fi

% 編排模式說明：
% chinese  → \showoriginalfalse  \showtranslationtrue  （A4 直式，純中文）
% german   → \showoriginaltrue   \showtranslationfalse （A4 直式，純德文）
% bilingual→ \showoriginaltrue   \showtranslationtrue  \bilingualcolumnstrue（A3/A4 橫式對照）

\ifbilingualcolumns
  \ifx\bilingualpaper\bilingualpaperaFour
    \geometry{
      a4paper,
      landscape,
      left=1.25cm,
      right=2.25cm,
      top=1.25cm,
      bottom=1.75cm,
      footskip=8.5mm,
      marginparwidth=0.75cm,
      marginparsep=0.25cm
    }
  \else
    \geometry{
      a3paper,
      landscape,
      left=1.55cm,
      right=2.75cm,
      top=1.55cm,
      bottom=2.05cm,
      footskip=9.5mm,
      marginparwidth=0.9cm,
      marginparsep=0.35cm
    }
  \fi
\else
\geometry{
  a4paper,
  portrait,
  left=2.2cm,
  right=2.6cm,
  top=2.0cm,
  bottom=2.0cm,
  marginparwidth=0.8cm,
  marginparsep=0.3cm
}
\fi
% ===== 時間戳基礎設施 =====
\newcount\stamptimeval
\newcount\stampminuteval
\newcommand{\stamphh}{%
  \stamptimeval=\time
  \divide\stamptimeval by 60
  \ifnum\stamptimeval<10 0\fi
  \the\stamptimeval}
\newcommand{\stampmm}{%
  \stamptimeval=\time
  \stampminuteval=\time
  \divide\stampminuteval by 60
  \multiply\stampminuteval by 60
  \advance\stamptimeval by -\stampminuteval
  \ifnum\stamptimeval<10 0\fi
  \the\stamptimeval}
\newcommand{\stampwkde}[1]{%
  \ifcase#1Montag\or Dienstag\or Mittwoch\or Donnerstag\or Freitag\or Samstag\or Sonntag\fi}
\newcommand{\stampwkzh}[1]{%
  \ifcase#1 一\or 二\or 三\or 四\or 五\or 六\or 日\fi}
\newcommand{\stampmonthde}[1]{%
  \ifcase#1\or Januar\or Februar\or März\or April\or Mai\or Juni\or
  Juli\or August\or September\or Oktober\or November\or Dezember\fi}
\newcount\stampjdval
\newcount\stampwdval
\pgfcalendardatetojulian{\the\year-\the\month-\the\day}{\stampjdval}
\pgfcalendarjuliantoweekday{\the\stampjdval}{\stampwdval}
\newcommand{\timestampzh}{%
  \songfont 最後修改：\the\year 年\the\month 月\the\day 日%
  % （星期\stampwkzh{\the\stampwdval}）%
  % \space\stamphh:\stampmm
  }

\newcommand{\documentdatezh}{\the\year 年 \the\month 月 \the\day 日}
\newcommand{\documentdatede}{\the\day. \stampmonthde{\the\month} \the\year}
\newcommand{\bverfgetranslationnote}{%
  {\songfont 翻譯說明：初譯使用 Claude Code 與 GPT 協助迭代；仍請以德文原文為準。}%
}
\newcommand{\bverfgecourseinfo}{%
  {\songfont 課程：114 學年度第 2 學期；憲法專題研究（四）；授課老師：湯德宗教授；導讀日期：115 年 6 月 1 日。}%
}
\newcommand{\bverfgeeditorinfo}{%
  {\songfont 編者：王逸帆（R10A21126），國立台灣大學法律系研究所碩士班。}%
}
\newcommand{\bverfgetitleversionline}{%
  \ifbilingualcolumns
    Fassung \documentversion\enspace\textperiodcentered\enspace Stand: \documentdatede
    \quad
    {\songfont 版本 \documentversion\enspace\textperiodcentered\enspace 修訂日期：\documentdatezh\par}%
  \else
    \ifshowtranslation
      {\songfont 版本 \documentversion\enspace\textperiodcentered\enspace 修訂日期：\documentdatezh\par}%
    \else
      Fassung \documentversion\enspace\textperiodcentered\enspace Stand: \documentdatede\par
    \fi
  \fi
}
\newcommand{\bverfgetitlecornerinfo}{%
  \ifshowtranslation
    \vspace{0.72em}%
    \noindent\hfill
    \begin{minipage}{0.66\textwidth}
      \raggedleft\tiny\color{pendingink}\songfont
      \bverfgecourseinfo\par
      \bverfgeeditorinfo
    \end{minipage}%
  \fi
}
\newcommand{\bverfgetitlemeta}{%
  \vfill
  \begin{center}%
    \footnotesize\color{pendingink}%
    \bverfgetitleversionline
    \ifshowtranslation
      \vspace{0.28em}{\scriptsize\bverfgetranslationnote\par}%
    \fi
  \end{center}%
  \bverfgetitlecornerinfo
}

\pagestyle{fancy}
\fancyhf{}
\renewcommand{\headrulewidth}{0pt}
\renewcommand{\footrulewidth}{0pt}
\fancyfoot[C]{\thepage}
\ifbilingualcolumns
  \fancyfoot[R]{\scriptsize\color{pendingink}\timestampzh}%
\else
  \ifshoworiginal
  \else
    \fancyfoot[R]{\scriptsize\color{pendingink}\timestampzh\hspace{0.45cm}}%
  \fi
\fi
\setmainfont{Times New Roman}
\setsansfont{Helvetica Neue}
\setmonofont{Menlo}
\IfFileExists{/Users/iw/Library/Fonts/kaiu.ttf}{
  \setCJKmainfont[Path=/Users/iw/Library/Fonts/,AutoFakeBold=4]{kaiu.ttf}
}{
  \IfFontExistsTF{標楷體}{
    \setCJKmainfont[AutoFakeBold=4]{標楷體}
  }{
    \setCJKmainfont{Songti TC}
  }
}
\IfFontExistsTF{Songti TC}{
  \setCJKfamilyfont{song}{Songti TC}
  \setCJKmonofont{Songti TC}
}{
  \setCJKfamilyfont{song}[Path=/Users/iw/Library/Fonts/,AutoFakeBold=4]{kaiu.ttf}
  \setCJKmonofont[Path=/Users/iw/Library/Fonts/,AutoFakeBold=4]{kaiu.ttf}
}
\newcommand{\songfont}{\CJKfamily{song}}

% ===== 封面／目錄專用打字機字體（僅限封面、目錄頁使用，不影響正文）=====
\IfFontExistsTF{Radio Newsman}{%
  \newfontfamily\bverfgetitlede{Radio Newsman}%
}{%
  \newfontfamily\bverfgetitlede{Courier New}%
}
\IfFontExistsTF{Erikas Farbband}{%
  \newfontfamily\bverfgebodyde{Erikas Farbband}%
}{%
  \let\bverfgebodyde\bverfgetitlede
}
\IfFontExistsTF{Maoken Heavy Labourer}{%
  \setCJKfamilyfont{bverfgetitlecjk}{Maoken Heavy Labourer}%
}{%
  \setCJKfamilyfont{bverfgetitlecjk}{Songti TC}%
}
\IfFontExistsTF{Huiwen-mincho}{%
  \setCJKfamilyfont{bverfgebodycjk}{Huiwen-mincho}%
}{%
  \setCJKfamilyfont{bverfgebodycjk}{Songti TC}%
}
\newcommand{\bverfgetitlezh}{\bverfgetitlede\CJKfamily{bverfgetitlecjk}}
\newcommand{\bverfgebodyzh}{\bverfgebodyde\CJKfamily{bverfgebodycjk}}

\setstretch{1.35}
\setlist{itemsep=0.2em, topsep=0.4em}
\setlength{\parindent}{0pt}
\setlength{\parskip}{0.45em}
\raggedbottom
\setcounter{secnumdepth}{0}
\setcounter{tocdepth}{2}

\newcommand{\lawboxsourcefont}{\footnotesize}
\newcommand{\lawboxtransfont}{\footnotesize}
\newcommand{\lawboxsourcestretch}{1.02}
\newcommand{\lawboxtransstretch}{0.92}

\titleformat{\section}{\centering\Large\bfseries\songfont}{\thesection}{1em}{}
\titleformat{\subsection}{\large\bfseries\songfont}{\thesubsection}{1em}{}
\titleformat{\subsubsection}{\normalsize\bfseries\songfont}{\thesubsubsection}{1em}{}
\titleformat{\paragraph}{\normalsize\bfseries\songfont}{\theparagraph}{1em}{}
\titlespacing*{\section}{0pt}{1.8em}{1.0em}
\titlespacing*{\subsection}{0pt}{1.2em}{0.6em}
\titlespacing*{\subsubsection}{0pt}{1.0em}{0.45em}
\titlespacing*{\paragraph}{0pt}{0.6em}{0.4em}

\definecolor{noteblue}{HTML}{A9C9F5}
\definecolor{noteteal}{HTML}{AEE1D6}
\definecolor{notegray}{HTML}{D7DADF}
\definecolor{caseink}{HTML}{000000}
\definecolor{casepaper}{HTML}{FCFEFD}
\definecolor{sourcepaper}{HTML}{FAFBF8}
\definecolor{sourceline}{HTML}{D7DED8}
\definecolor{sourceink}{HTML}{000000}
\definecolor{pendingink}{HTML}{000000}

\tcbset{
  caseboxbase/.style={
    enhanced,
    breakable,
    width=0.985\linewidth,
    colback=casepaper,
    coltitle=caseink,
    fonttitle=\bfseries\songfont,
    boxrule=0.28pt,
    arc=1.2mm,
    left=2.4mm,
    right=2.4mm,
    top=1.5mm,
    bottom=1.5mm,
    boxsep=1.5mm,
    before skip=0.65em,
    after skip=0.65em,
    before upper={\setlength{\parindent}{0pt}\setlength{\parskip}{0.35em}},
  }
}

\newcommand{\bilingualstart}{}
\newcommand{\bilingualstop}{}
\ifbilingualcolumns
  \renewcommand{\bilingualstart}{\small\setlength{\columnsep}{1.25cm}\setlength{\columnseprule}{0.12pt}\columnratio{0.5,0.5}\multicolumnfootnotes\flushbottom\sloppy\interlinepenalty=80\clubpenalty=800\widowpenalty=800\displaywidowpenalty=800\begin{paracol}{2}\global\intranslationfalse\global\firstbilingualsectiontrue{\footnotesize\color{sourceink}Deutsch\par}\switchcolumn{\footnotesize\songfont\color{sourceink}中文譯文\par}\switchcolumn*}
  \renewcommand{\bilingualstop}{\ifintranslation\switchcolumn*\global\intranslationfalse\fi\end{paracol}\clearpage\raggedbottom}
\fi
\newif\ifrandnumbers
\randnumbersfalse
\newcounter{randnumber}
\newcounter{randnumberlocal}
\newcounter{lastsourcerandstart}
\newif\iflastsourcehasrandnumbers
\lastsourcehasrandnumbersfalse
\newif\ifinlawbox
\inlawboxfalse
\newif\iflawboxhaspendingrn
\lawboxhaspendingrnfalse
\newcommand{\lawboxpendingrn}{}
\newcommand{\startrandnumbers}{\global\randnumberstrue\setcounter{randnumber}{1}\global\lastsourcehasrandnumbersfalse}
\newcommand{\stoprandnumbers}{\global\randnumbersfalse\global\lastsourcehasrandnumbersfalse}
\newlength{\randnumberwidth}
\setlength{\randnumberwidth}{0.95cm}
\newlength{\randnumberrightoffset}
\setlength{\randnumberrightoffset}{0.85cm}
\newlength{\randnumberlawboxrightoffset}
\setlength{\randnumberlawboxrightoffset}{1.40cm}
\newcommand{\randnumberbox}[1]{%
  \leavevmode
  \makebox[0pt][l]{%
    \smash{%
      \tikz[remember picture,overlay,baseline=(rnmark.base)]{%
        \coordinate (rnmark) at (0,0);%
        \ifinlawbox
          \path let \p1=(current page.east), \p2=(rnmark.base) in
            node[
              anchor=base east,
              inner sep=0pt,
              outer sep=0pt,
              text width=\randnumberwidth,
              align=right,
              text=pendingink,
              font=\normalfont\mdseries\scriptsize\ttfamily
            ] at (\x1-\randnumberlawboxrightoffset,\y2) {{\normalfont\mdseries\scriptsize\ttfamily #1}};%
        \else
          \path let \p1=(current page.east), \p2=(rnmark.base) in
            node[
              anchor=base east,
              inner sep=0pt,
              outer sep=0pt,
              text width=\randnumberwidth,
              align=right,
              text=pendingink,
              font=\normalfont\mdseries\scriptsize\ttfamily
            ] at (\x1-\randnumberrightoffset,\y2) {{\normalfont\mdseries\scriptsize\ttfamily #1}};%
        \fi
      }%
    }%
  }%
  \ignorespaces
}
\newcommand{\lawboxstorerandnumber}[1]{%
  \gdef\lawboxpendingrn{#1}%
  \global\lawboxhaspendingrntrue
  \ignorespaces
}
\newcommand{\lawboxuserandnumber}{%
  \iflawboxhaspendingrn
    \randnumberbox{\lawboxpendingrn}%
    \global\lawboxhaspendingrnfalse
  \fi
}
\newcommand{\sourcern}[1]{%
  \ifrandnumbers
    \ifshoworiginal
      \ifshowtranslation\else
        \ifinlawbox\lawboxstorerandnumber{#1}\else\randnumberbox{#1}\fi
      \fi
    \fi
  \fi
}
\newcommand{\transrn}[1]{%
  \ifrandnumbers
    \ifshowtranslation
      \ifinlawbox\lawboxstorerandnumber{#1}\else\randnumberbox{#1}\fi
    \fi
  \fi
}
% ===== 課堂模式：德文原文縮寫註解 =====
% 日常切換只改各判決檔的一行：
%   \classroommodeon        開啟課堂版；註解掉這行即回到一般版。
% 模板預設：
%   1. 課堂模式關閉。
%   2. 課堂模式開啟時，縮寫註解包含「德文全稱＋中文譯名」。
%   3. 課堂模式開啟時，GG/條文編者註仍會自動顯示。
% 進階微調才需要在單案檔覆寫：
%   \classroomabbrzhoff     縮寫註解僅顯示德文全稱。
%   \showeditornotestrue    一般版也顯示編者註。
% 註解策略：同一實際 PDF 頁面中，同一縮寫只在第一次出現處加註。
% 這裡用 zref-abspage 的絕對頁序，不用 \thepage，避免文件內頁碼重置時誤判。
\newif\ifclassroommode
\newif\ifclassroomabbrzh
\classroommodefalse
\classroomabbrzhtrue
\newcommand{\classroommodeon}{\classroommodetrue}
\newcommand{\classroommodeoff}{\classroommodefalse}
\newcommand{\classroomabbrzhon}{\classroomabbrzhtrue}
\newcommand{\classroomabbrzhoff}{\classroomabbrzhfalse}
\newcommand{\abbrnotelabel}{\emph{Abk.}: }
\newcommand{\abbrnote}[3]{%
  \ifclassroommode
    \ifshoworiginal
      \ifintranslation\else
        \footnote{\normalfont\footnotesize\abbrnotelabel #2\ifclassroomabbrzh；{\songfont #3}\fi}%
      \fi
    \fi
  \fi
}
\newcommand{\abbrnoteonce}[4]{%
  #2%
  \ifclassroommode
    \ifshoworiginal
      \ifintranslation\else
        \begingroup
          \edef\abbrcurrentabspage{\number\numexpr\value{abspage}+1\relax}%
          \ifcsname abbrnoteseen@#1@\abbrcurrentabspage\endcsname\else
            \global\expandafter\let\csname abbrnoteseen@#1@\abbrcurrentabspage\endcsname\empty
            \abbrnote{#1}{#3}{#4}%
          \fi
        \endgroup
      \fi
    \fi
  \fi
}
\newcommand{\abbrAbs}{\abbrnoteonce{abs}{Abs.}{Absatz}{項}}
\newcommand{\abbrAAO}{\abbrnoteonce{aao}{a.a.O.}{am angegebenen Ort}{前引處}}
\newcommand{\abbrAE}{\abbrnoteonce{ae}{AE}{Alternativ-Entwurf}{替代草案}}
\newcommand{\abbrArt}{\abbrnoteonce{art}{Art.}{Artikel}{條}}
\newcommand{\abbrAufl}{\abbrnoteonce{aufl}{Aufl.}{Auflage}{版}}
\newcommand{\abbrBd}{\abbrnoteonce{bd}{Bd.}{Band}{卷}}
\newcommand{\abbrBGBl}{\abbrnoteonce{bgbl}{BGBl.}{Bundesgesetzblatt}{聯邦法律公報}}
\newcommand{\abbrBGHSt}{\abbrnoteonce{bghst}{BGHSt}{Entscheidungen des Bundesgerichtshofes in Strafsachen}{聯邦普通法院刑事判決彙編}}
\newcommand{\abbrBRDrucks}{\abbrnoteonce{brdrucks}{BRDrucks.}{Bundesrats-Drucksache}{聯邦參議院文件}}
\newcommand{\abbrBTDrucks}{\abbrnoteonce{btdrucks}{BTDrucks.}{Bundestags-Drucksache}{聯邦議院文件}}
\newcommand{\abbrBundesgesetzbl}{\abbrnoteonce{bundesgesetzbl}{Bundesgesetzbl.}{Bundesgesetzblatt}{聯邦法律公報}}
\newcommand{\abbrBvF}{\abbrnoteonce{bvf}{BvF}{Bundesverfassungsgerichtsverfahren}{聯邦憲法法院程序案號}}
\newcommand{\abbrBvL}{\abbrnoteonce{bvl}{BvL}{Bundesverfassungsgerichtsverfahren}{聯邦憲法法院程序案號}}
\newcommand{\abbrBvR}{\abbrnoteonce{bvr}{BvR}{Bundesverfassungsgerichtsverfahren}{聯邦憲法法院程序案號}}
\newcommand{\abbrBVerfG}{\abbrnoteonce{bverfg}{BVerfG}{Bundesverfassungsgericht}{聯邦憲法法院}}
\newcommand{\abbrBVerfGE}{\abbrnoteonce{bverfge}{BVerfGE}{Entscheidungen des Bundesverfassungsgerichts}{聯邦憲法法院判決集}}
\newcommand{\abbrBVerfGG}{\abbrnoteonce{bverfgg}{BVerfGG}{Bundesverfassungsgerichtsgesetz}{聯邦憲法法院法}}
\newcommand{\abbrDDR}{\abbrnoteonce{ddr}{DDR}{Deutsche Demokratische Republik}{德意志民主共和國；東德}}
\newcommand{\abbrDJT}{\abbrnoteonce{djt}{DJT}{Deutscher Juristentag}{德國法學家大會}}
\newcommand{\abbrDr}{\abbrnoteonce{dr}{Dr.}{Doktor}{Dr.}}
\newcommand{\abbrEuGRZ}{\abbrnoteonce{eugrz}{EuGRZ}{Europäische Grundrechte-Zeitschrift}{歐洲基本權期刊}}
\newcommand{\abbrEV}{\abbrnoteonce{ev}{EV}{Einigungsvertrag}{統一條約}}
\newcommand{\abbrF}{\abbrnoteonce{f}{f.}{folgende}{以下一頁；下一段}}
\newcommand{\abbrFF}{\abbrnoteonce{ff}{ff.}{fortfolgende}{以下數頁；以下各段}}
\newcommand{\abbrGBlDDR}{\abbrnoteonce{gblddr}{GBl. DDR}{Gesetzblatt der Deutschen Demokratischen Republik}{德意志民主共和國法律公報}}
\newcommand{\abbrGG}{\abbrnoteonce{gg}{GG}{Grundgesetz}{基本法}}
\newcommand{\abbrGVBl}{\abbrnoteonce{gvbl}{GVBl.}{Gesetz- und Verordnungsblatt}{法律及命令公報}}
\newcommand{\abbrJR}{\abbrnoteonce{jr}{JR}{Juristische Rundschau}{法學評論}}
\newcommand{\abbrJZ}{\abbrnoteonce{jz}{JZ}{JuristenZeitung}{法學家報}}
\newcommand{\abbrIVm}{\abbrnoteonce{ivm}{i.V.m.}{in Verbindung mit}{結合；連同}}
\newcommand{\abbrMdB}{\abbrnoteonce{mdb}{MdB}{Mitglied des Bundestages}{聯邦議院議員}}
\newcommand{\abbrMwN}{\abbrnoteonce{mwn}{m.w.N.}{mit weiteren Nachweisen}{附進一步文獻／判例出處}}
\newcommand{\abbrNF}{\abbrnoteonce{nf}{n.F.}{neue Fassung}{新版本}}
\newcommand{\abbrAF}{\abbrnoteonce{af}{a.F.}{alte Fassung}{舊版本}}
\newcommand{\abbrNr}{\abbrnoteonce{nr}{Nr.}{Nummer}{款、目或號，依語境}}
\newcommand{\abbrProf}{\abbrnoteonce{prof}{Prof.}{Professor}{教授}}
\newcommand{\abbrRdnr}{\abbrnoteonce{rdnr}{Rdnr.}{Randnummer}{邊碼；段碼}}
\newcommand{\abbrRdnrn}{\abbrnoteonce{rdnrn}{Rdnrn.}{Randnummern}{邊碼；段碼}}
\newcommand{\abbrRGBl}{\abbrnoteonce{rgbl}{RGBl.}{Reichsgesetzblatt}{帝國法律公報}}
\newcommand{\abbrRGSt}{\abbrnoteonce{rgst}{RGSt}{Entscheidungen des Reichsgerichts in Strafsachen}{帝國法院刑事判決彙編}}
\newcommand{\abbrRspr}{\abbrnoteonce{rspr}{Rspr.}{Rechtsprechung}{判例；司法實務}}
\newcommand{\abbrRVO}{\abbrnoteonce{rvo}{RVO}{Reichsversicherungsordnung}{帝國保險法}}
\newcommand{\abbrS}{\abbrnoteonce{s}{S.}{Seite}{頁}}
\newcommand{\abbrSFHG}{\abbrnoteonce{sfhg}{SFHG}{Schwangeren- und Familienhilfegesetz}{孕婦及家庭扶助法}}
\newcommand{\abbrSGBV}{\abbrnoteonce{sgbv}{SGB V}{Fünftes Buch Sozialgesetzbuch}{社會法典第五編}}
\newcommand{\abbrStAG}{\abbrnoteonce{staeg}{StÄG}{Strafrechtsänderungsgesetz}{刑法修正法}}
\newcommand{\abbrStenBer}{\abbrnoteonce{stenber}{StenBer.}{Stenographischer Bericht}{速記紀錄}}
\newcommand{\abbrStGB}{\abbrnoteonce{stgb}{StGB}{Strafgesetzbuch}{刑法}}
\newcommand{\abbrStREG}{\abbrnoteonce{streg}{StREG}{Strafrechtsreform-Ergänzungsgesetz}{刑法改革補充法}}
\newcommand{\abbrStrRG}{\abbrnoteonce{strrg}{StrRG}{Strafrechtsreformgesetz}{刑法改革法}}
\newcommand{\abbrUS}{\abbrnoteonce{us}{U.S.}{United States Reports}{美國最高法院判決彙編}}
\newcommand{\abbrVgl}{\abbrnoteonce{vgl}{vgl.}{vergleiche}{參見}}
\newcommand{\abbrVVDStRL}{\abbrnoteonce{vvdstrl}{VVDStRL}{Veröffentlichungen der Vereinigung der Deutschen Staatsrechtslehrer}{德國公法教師協會論文集}}
\newcommand{\abbrWP}{\abbrnoteonce{wp}{WP.}{Wahlperiode}{屆期}}
\newcommand{\abbrWp}{\abbrnoteonce{wp}{Wp.}{Wahlperiode}{屆期}}
\newcommand{\abbrZStrW}{\abbrnoteonce{zstrw}{ZStrW}{Zeitschrift für die gesamte Strafrechtswissenschaft}{整體刑法學期刊}}
\newcommand{\recordsourcerandstart}{%
  \ifrandnumbers
    \setcounter{lastsourcerandstart}{\value{randnumber}}%
    \global\lastsourcehasrandnumberstrue
  \else
    \global\lastsourcehasrandnumbersfalse
  \fi
}
\newlength{\biparagraphskip}
\setlength{\biparagraphskip}{0.52em}
\newlength{\lawboxblockbeforeskip}
\setlength{\lawboxblockbeforeskip}{0.72em}
\newlength{\lawboxblockafterskip}
\setlength{\lawboxblockafterskip}{0.92em}
\newif\iflawboxpartendedblock
\lawboxpartendedblockfalse
\newcommand{\sourceparstyle}{\footnotesize\color{sourceink}\setstretch{1.12}\RaggedRight\emergencystretch=3em\setlength{\parskip}{0.42em}\obeylines\selectfont}
\newcommand{\transparstyle}{\songfont\color{caseink}\setstretch{1.12}\setlength{\parindent}{0pt}\setlength{\parskip}{0.42em}}
\newcommand{\bilingualpairskip}{%
  \par\addvspace{\biparagraphskip}%
  \switchcolumn
  \par\addvspace{\biparagraphskip}%
  \switchcolumn*%
  \global\intranslationfalse
}
\newcommand{\bilingualboxpairskip}{%
  \par
  \switchcolumn
  \par
  \switchcolumn*%
  \global\intranslationfalse
}
\newcommand{\bilinguallawboxblockskip}{%
  \switchcolumn*%
  \par\addvspace{\lawboxblockafterskip}%
  \switchcolumn
  \par\addvspace{\lawboxblockafterskip}%
  \switchcolumn*%
  \global\intranslationfalse
}
\newcommand{\bikeepwithnext}[1][4]{%
  \ifbilingualcolumns
    \syncleft
    \Needspace{#1\baselineskip}%
    \switchcolumn
    \Needspace{#1\baselineskip}%
    \switchcolumn*%
    \global\intranslationfalse
  \else
    \Needspace{#1\baselineskip}%
  \fi
}
\NewEnviron{sourcepara}[1][]{
  \recordsourcerandstart
  \ifshoworiginal
    \ifbilingualcolumns
      \ifintranslation\switchcolumn*\global\intranslationfalse\fi
    \else
      \Needspace{5\baselineskip}
    \fi
    \ifbilingualcolumns\else\par\addvspace{0.35em}\fi
    {\sourceparstyle
    \noindent\BODY\par}
    \ifbilingualcolumns\else\addvspace{0.25em}\fi
    \ifbilingualcolumns
      \switchcolumn
      \global\intranslationtrue
    \fi
  \else
    \ifrandnumbers
      \begingroup
        \setbox0=\vbox{\hsize=\linewidth\sourceparstyle\noindent\BODY\par}%
      \endgroup
    \fi
  \fi
}
\NewEnviron{transpara}{
  \ifshowtranslation
    \setcounter{randnumberlocal}{\value{lastsourcerandstart}}%
    \ifbilingualcolumns
      \ifintranslation\else\switchcolumn\global\intranslationtrue\fi
      {\transparstyle\setlength{\leftskip}{0.55em}\setlength{\rightskip}{0.15em plus 1em}\addtolength{\linewidth}{-0.55em}\BODY\par}
      \switchcolumn*
      \bilingualpairskip
    \else
      {\transparstyle\BODY\par}
    \fi
  \else
    \ifbilingualcolumns
      \ifintranslation\switchcolumn*\global\intranslationfalse\fi
    \fi
  \fi
}
\NewEnviron{sourceboxpara}[1][]{
  \global\lawboxpartendedblockfalse
  \recordsourcerandstart
  \ifshoworiginal
    \ifbilingualcolumns
      \ifintranslation\switchcolumn*\global\intranslationfalse\fi
    \else
      \Needspace{3\baselineskip}
    \fi
    {\sourceparstyle
    \BODY\par}
    \ifbilingualcolumns\else
      \iflawboxpartendedblock\addvspace{\lawboxblockafterskip}\fi
    \fi
    \ifbilingualcolumns
      \switchcolumn
      \global\intranslationtrue
    \fi
  \fi
}
\NewEnviron{transboxpara}{
  \global\lawboxpartendedblockfalse
  \ifshowtranslation
    \setcounter{randnumberlocal}{\value{lastsourcerandstart}}%
    \ifbilingualcolumns
      \ifintranslation\else\switchcolumn\global\intranslationtrue\fi
      {\transparstyle\BODY\par}
      \iflawboxpartendedblock
        \bilinguallawboxblockskip
      \else
        \switchcolumn*
        \bilingualboxpairskip
      \fi
    \else
      {\transparstyle\BODY\par}
      \iflawboxpartendedblock\addvspace{\lawboxblockafterskip}\fi
    \fi
  \else
    \ifbilingualcolumns
      \ifintranslation\switchcolumn*\global\intranslationfalse\fi
    \fi
  \fi
}
\newcommand{\leitsatznum}[1]{\noindent\hangindent=3.0em\hangafter=1\makebox[2.25em][r]{\bfseries #1.}\hspace{0.55em}\ignorespaces}
\newcommand{\syncleft}{\ifbilingualcolumns\ifintranslation\switchcolumn*\global\intranslationfalse\fi\fi}
\pretocmd{\section}{\syncleft\clearpage}{}{}
\pretocmd{\subsection}{\syncleft}{}{}
\pretocmd{\subsubsection}{\syncleft}{}{}
\newcommand{\biheadingfont}{\songfont}
\newcommand{\bitocentry}[2]{%
  \ifshoworiginal
    \ifshowtranslation
      \ifstrequal{#1}{#2}{#1}{#1\space #2}%
    \else
      #1%
    \fi
  \else
    #2%
  \fi
}
\pdfstringdefDisableCommands{%
  \def\bitocentry#1#2{#1}%
}
\newcommand{\biheading}[7]{%
  \syncleft#1\bikeepwithnext[#3]\phantomsection\addcontentsline{toc}{#2}{\bitocentry{#6}{#7}}%
  \ifbilingualcolumns
    {#4 #6\par}%
    \switchcolumn
    {#4\biheadingfont #7\par}%
    \switchcolumn*%
    \global\intranslationfalse
    \par\addvspace{#5}%
    \switchcolumn
    \par\addvspace{#5}%
    \switchcolumn*%
  \else
    \ifshoworiginal{#4 #6\par}\fi
    \ifshowtranslation{#4\biheadingfont #7\par}\fi
    \addvspace{#5}%
  \fi
}
\newcommand{\termsectionheading}[1]{%
  \par\Needspace{9\baselineskip}%
  \vspace{0.65em}%
  {\songfont\bfseries #1\par}%
  \vspace{0.2em}%
}
% 章節換頁：
%   雙語 paracol 欄內不要用 \clearpage；它會在同步兩欄時吐出只有欄線/頁碼的空頁。
%   \newpage 足以讓下一個 section 起新頁，又不會強制清空所有待排材料。
%   單語模式不在 paracol 內，仍可用 \clearpage。
\newcommand{\bisectionpagebreak}{\ifbilingualcolumns\newpage\else\clearpage\fi}
\newcommand{\bisection}[2]{%
  \biheading{\iffirstbilingualsection\global\firstbilingualsectionfalse\else\bisectionpagebreak\fi}{section}{6}{\centering\Large\bfseries}{0.8em}{#1}{#2}%
}
\newcommand{\bisubsection}[2]{%
  \biheading{}{subsection}{5}{\large\bfseries}{0.55em}{#1}{#2}%
}
\newcommand{\bisubsubsection}[2]{%
  \biheading{}{subsubsection}{4}{\normalsize\bfseries}{0.45em}{#1}{#2}%
}
\newcommand{\bicompositeheading}[4]{%
  \biheading{}{subsection}{5}{\large\bfseries}{0.16em}{#1}{#2}%
  \biheading{}{subsubsection}{3}{\normalsize\bfseries}{0.45em}{#3}{#4}%
}
\newif\iflawboxfirstitem
\lawboxfirstitemfalse
\newcommand{\lawboxprespace}[1]{%
  \par
  \iflawboxfirstitem
    \global\lawboxfirstitemfalse
  \else
    \addvspace{#1}%
  \fi
}
\newtcolorbox{lawboxinner}{
  caseboxbase,
  colframe=sourceline!85!black,
  colback=white,
  left=1.05em,
  right=1.05em,
  top=0.62em,
  bottom=0.62em,
  before upper={\global\inlawboxtrue\global\lawboxfirstitemtrue\global\lawboxhaspendingrnfalse\ifintranslation\lawboxtransfont\else\lawboxsourcefont\fi\RaggedRight\setlength{\parindent}{0pt}\setlength{\parskip}{0pt}\ifintranslation\setstretch{\lawboxtransstretch}\else\setstretch{\lawboxsourcestretch}\fi},
  after upper={\global\lawboxhaspendingrnfalse\global\inlawboxfalse}
}
\tcbset{
  lawboxpartbase/.style={
    enhanced,
    breakable=false,
    width=0.985\linewidth,
    colback=white,
    frame hidden,
    boxrule=0pt,
    arc=0pt,
    left=1.05em,
    right=1.05em,
    boxsep=1.5mm,
    before skip=0pt,
    after skip=0pt,
    before upper={\global\inlawboxtrue\global\lawboxfirstitemtrue\global\lawboxhaspendingrnfalse\ifintranslation\lawboxtransfont\else\lawboxsourcefont\fi\RaggedRight\setlength{\parindent}{0pt}\setlength{\parskip}{0pt}\ifintranslation\setstretch{\lawboxtransstretch}\else\setstretch{\lawboxsourcestretch}\fi},
    after upper={\global\lawboxhaspendingrnfalse\global\inlawboxfalse},
    borderline west={0.28pt}{0pt}{sourceline!85!black},
    borderline east={0.28pt}{0pt}{sourceline!85!black},
  },
  lawboxpartfirst/.style={
    lawboxpartbase,
    top=0.62em,
    bottom=0.04em,
    before skip=\lawboxblockbeforeskip,
    borderline north={0.28pt}{0pt}{sourceline!85!black},
  },
  lawboxpartmiddle/.style={
    lawboxpartbase,
    top=0.04em,
    bottom=0.04em,
  },
  lawboxpartlast/.style={
    lawboxpartbase,
    top=0.04em,
    bottom=0.62em,
    after skip=0pt,
    borderline south={0.28pt}{0pt}{sourceline!85!black},
  }
}
\newtcolorbox{lawboxpartinner}[1]{#1}
\newtcolorbox{termboxinner}{
  caseboxbase,
  colframe=noteblue!70!black,
  colback=casepaper,
  before upper={\songfont\setlength{\parindent}{0pt}\setlength{\parskip}{0.35em}}
}
\newtcolorbox{deboxinner}{
  caseboxbase,
  colframe=notegray!72!black,
  colback=white,
  left=1.05em,
  right=1.05em,
  top=0.62em,
  bottom=0.62em,
  before upper={\global\inlawboxtrue\global\lawboxfirstitemtrue\global\lawboxhaspendingrnfalse\small\setlength{\parindent}{0pt}\setlength{\parskip}{0.35em}},
  after upper={\global\lawboxhaspendingrnfalse\global\inlawboxfalse}
}
\NewEnviron{lawbox}{
  \begin{lawboxinner}\BODY\end{lawboxinner}
}
\NewEnviron{lawboxpart}[2]{
  \begin{lawboxpartinner}{lawboxpart#1,equal height group=#2}\BODY\end{lawboxpartinner}%
  \ifstrequal{#1}{last}{\global\lawboxpartendedblocktrue}{}%
}
\NewEnviron{termbox}{
  \par\addvspace{0.55em}
  {\color{noteblue!70!black}\rule{\linewidth}{0.28pt}}\par
  \begingroup
    \songfont\small\color{caseink}
    \setlength{\parindent}{0pt}
    \setlength{\parskip}{0.24em}
    \BODY
    \par
  \endgroup
  {\color{noteblue!70!black}\rule{\linewidth}{0.28pt}}\par
  \addvspace{0.55em}
}
\NewEnviron{debox}{
  \begin{deboxinner}\BODY\end{deboxinner}
}
\providecommand{\paragraphmark}[1]{}
\newcommand{\lawboxtitleafter}{\addvspace{0.06em}}
\newcommand{\lawtitle}[1]{\lawboxprespace{0.22em}{\lawboxtransfont\bfseries\lawboxuserandnumber #1}\par\lawboxtitleafter}
\newcommand{\absno}[2]{\lawboxprespace{0.04em}\noindent\lawboxuserandnumber\hangindent=2.6em\hangafter=1\makebox[2.6em][l]{(#1)}#2\par}
\newcommand{\nrno}[2]{\lawboxprespace{0pt}\noindent\lawboxuserandnumber\hspace*{2.6em}\hangindent=4.4em\hangafter=1\makebox[1.8em][l]{#1.}#2\par}
\newcommand{\letterno}[2]{\lawboxprespace{0pt}\noindent\lawboxuserandnumber\hspace*{4.4em}\hangindent=6.0em\hangafter=1\makebox[1.6em][l]{#1)}#2\par}
\newcommand{\sourcelawtitle}[1]{\lawboxprespace{0.22em}{\lawboxsourcefont\bfseries\lawboxuserandnumber #1}\par\lawboxtitleafter}
\newcommand{\sourcelawsubtitle}[1]{\lawboxprespace{0.04em}{\lawboxsourcefont\bfseries\lawboxuserandnumber #1}\par\lawboxtitleafter}
\newcommand{\sourceabsno}[2]{\lawboxprespace{0.04em}\noindent\lawboxuserandnumber\hangindent=2.45em\hangafter=1\makebox[2.45em][l]{(#1)}#2\par}
\newcommand{\sourcenrno}[2]{\lawboxprespace{0pt}\noindent\lawboxuserandnumber\hspace*{2.45em}\hangindent=4.25em\hangafter=1\makebox[1.8em][l]{#1.}#2\par}
\newcommand{\sourceletterno}[2]{\lawboxprespace{0pt}\noindent\lawboxuserandnumber\hspace*{4.25em}\hangindent=5.85em\hangafter=1\makebox[1.6em][l]{#1)}#2\par}
\newcommand{\sourcequotepara}[1]{\lawboxprespace{0.04em}\noindent\lawboxuserandnumber\hangindent=1.2em\hangafter=1 #1\par}
\newcommand{\sourcelawline}[1]{\lawboxprespace{0.04em}\noindent\lawboxuserandnumber #1\par}
\newcommand{\lawline}[1]{\lawboxprespace{0.04em}\noindent\lawboxuserandnumber #1\par}
\newcommand{\pendingtranslation}{\par{\songfont\footnotesize\color{pendingink}（中文譯文待補。）}\par}
\newcommand{\tocpart}[1]{\par\addvspace{0.52em}{\footnotesize\bfseries\color{pendingink}#1\par}\addvspace{0.16em}}
\newcommand{\tocdashline}{\par\addvspace{0.48em}\noindent{\color{notegray}\xleaders\hbox to 0.82em{\hfil\rule{0.42em}{0.28pt}\hfil}\hskip\linewidth}\par\addvspace{0.38em}}
% \tocpanel{font-cmd}{content}：將目錄置中於所在欄寬（雙語欄適用）。
% A3 橫式使用 0.58\linewidth，A4 橫式使用 0.84\linewidth，
% 使兩種紙張下目錄欄內容實際寬度相近（均約 100–105 mm）。
\newcommand{\tocpanel}[2]{%
  \makebox[\linewidth][c]{%
    \ifx\bilingualpaper\bilingualpaperaThree
      \begin{minipage}{0.58\linewidth}%
    \else
      \begin{minipage}{0.84\linewidth}%
    \fi
    \toccolstyle
    #1#2%
    \end{minipage}%
  }%
}
% ===== 首頁簡目錄的兩層 description list 縮排 =====
% enumitem 的 description list 可把每一列拆成：
%   [label]  labelsep  body text
% 這裡的「起點」都以所在 minipage/欄位的左緣為 0。
%
% 核心規則：
%   labelindent = label 欄位從左緣開始的位置。
%   labelwidth  = label 欄位本身寬度；標籤太長（如 Abw. Meinung）就加大它。
%   labelsep    = label 欄位右緣到正文的距離。
%   leftmargin  = 正文 body text 的起點。判斷層級視覺時主要看這個。
%
% 所以：
%   想讓「Begründung / 判決理由」整列正文往右或往左：改 tocitems 的 leftmargin。
%   想讓「A. Verfahren und Gegenstand / A. 程序與審查標的」正文往右或往左：
%     改 tocsubitems 的 leftmargin。
%   想讓 A./B./C. 這些標籤本身往右或往左，但正文不動：改 tocsubitems 的 labelindent。
%   想讓 A./B./C. 與正文間距變大或變小：改 tocsubitems 的 labelsep。
%
% 層級原則：
%   tocsubitems.leftmargin 必須大於 tocitems.leftmargin，
%   否則子層級正文會看起來比「Gründe / 理由」還靠左。
\newlist{tocitems}{description}{1}
\setlist[tocitems]{%
  labelindent=0pt,    % 第一層 label 欄位起點；0pt 表示貼齊目錄欄左緣。
  leftmargin=2.54cm,  % 第一層正文起點；例如 Begründung / 判決理由 的左緣。
  labelwidth=2.16cm,  % 第一層 label 欄位寬度；需容納 Leitsätze、Abw. Meinung、不同意見等。
  labelsep=0.32cm,    % 第一層 label 與正文的水平間距；只改間距，不改正文起點。
  itemsep=0.28em,     % 第一層各項之間的垂直距離。
  topsep=0.20em,      % 第一層 list 上下與前後內容的垂直距離。
  parsep=0pt,         % 同一 item 內段落之間的距離；目錄項通常不需要。
  partopsep=0pt,      % list 若緊接段落開始時的額外距離；關閉以免目錄高度飄動。
  align=left,         % label 欄內左對齊；長短標籤不靠右貼齊。
  font=\normalfont\bfseries % label 的字型；正文不受這行影響。
}
\newlist{tocsubitems}{description}{1}
\setlist[tocsubitems]{%
  labelindent=1cm,    % 第二層 label 起點；只控制 A./B./C. 本身的位置。
  leftmargin=3.00cm,  % 第二層正文起點；必須大於 2.54cm，才會比 Begründung / 判決理由更右。
  labelwidth=0.5cm,   % 第二層 label 欄寬；A./B./C. 很短，可小於第一層。
  labelsep=1.5cm,    % 第二層 label 與正文間距；A. 與 Verfahren / 程序 的距離看這裡。
  itemsep=0.12em,     % 第二層各項較密，避免 A.–E. 佔太高。
  topsep=0.04em,      % 第二層 list 與 Gründe / 理由之間的垂直距離。
  parsep=0pt,         % 同一第二層 item 內段落距離；目錄項通常不需要。
  partopsep=0pt,      % 關閉額外段落起始距離，避免版面不穩。
  align=left,         % A./B./C. label 欄內左對齊。
  font=\normalfont\bfseries, % A./B./C. label 加粗；正文不受這行影響。
  before={}           % 不在第二層 list 前自動插入額外內容。
}
\newlist{termitems}{description}{1}
\ifbilingualcolumns
  \setlist[termitems]{%
    leftmargin=5.2cm,
    labelwidth=4.8cm,
    labelsep=0.35cm,
    itemsep=0.16em,
    topsep=0.2em,
    parsep=0pt,
    partopsep=0pt,
    font=\normalfont\bfseries,
    style=nextline
  }
\else
  \setlist[termitems]{%
    leftmargin=0pt,
    labelwidth=0pt,
    labelsep=0pt,
    itemsep=0.24em,
    topsep=0.2em,
    parsep=0pt,
    partopsep=0pt,
    font=\normalfont\bfseries,
    style=nextline
  }
\fi
\newlist{abbritems}{description}{1}
\setlist[abbritems]{%
  leftmargin=2.38cm,
  labelwidth=2.02cm,
  labelsep=0.28cm,
  itemsep=0.16em,
  topsep=0.2em,
  parsep=0pt,
  partopsep=0pt,
  align=left,
  font=\normalfont\bfseries
}
\ifbilingualcolumns
  \newenvironment{termcolumns}{\begin{multicols}{2}}{\end{multicols}}
\else
  \newenvironment{termcolumns}{}{}
\fi

% ===== 編者註腳開關 =====
% 一般版預設不顯示編者註，避免正文腳註過密。
% 若開啟 \classroommodeon，GG/條文編者註仍會自動顯示，無需另開本開關。
% 若一般版也要顯示編者註，可在單案檔 \input 後加 \showeditornotestrue。
\newif\ifshoweditornotes
\showeditornotesfalse

% ===== 目錄排版輔助 =====
\newcommand{\toccolstyle}{\raggedright\arraybackslash}
\newcommand{\tocsingleindent}{\leftskip=1.45cm\rightskip=1.2cm}

% ===== 行內引文環境（德文原文與中文譯文各一，帶左側灰色邊線）=====
\newcommand{\quoteparbreak}{\par\addvspace{0.48em}}
\newcommand{\sourcequoteinner}[1]{%
  \begin{tcolorbox}[
    enhanced, breakable, width=\dimexpr\linewidth-0.8em\relax, frame hidden,
    colback=white, coltext=caseink, boxrule=0pt,
    borderline west={0.55pt}{0.88em}{notegray},
    left=1.78em, right=0.72em, top=0.18em, bottom=0.18em,
    boxsep=0pt, before skip=0.24em, after skip=0.24em,
    before upper={\normalfont\footnotesize\color{caseink}\setlength{\parindent}{0pt}\setlength{\parskip}{0.34em}\raggedright\setlength{\rightskip}{0pt plus 1em}}
  ]%
  #1\par
  \end{tcolorbox}%
}
\newcommand{\transquoteinner}[1]{%
  \begin{tcolorbox}[
    enhanced, breakable, width=\dimexpr\linewidth-0.8em\relax, frame hidden,
    colback=white, coltext=caseink, boxrule=0pt,
    borderline west={0.55pt}{0.88em}{notegray},
    left=1.78em, right=0.72em, top=0.18em, bottom=0.18em,
    boxsep=0pt, before skip=0.24em, after skip=0.24em,
    before upper={\songfont\small\color{caseink}\setlength{\parindent}{0pt}\setlength{\parskip}{0.34em}\raggedright\setlength{\rightskip}{0pt plus 1em}}
  ]%
  #1\par
  \end{tcolorbox}%
}

% ===== 大綱（Gliederung）Rn 輔助命令 =====
\newcommand{\outrn}[2]{\enspace{\normalfont\footnotesize\color{pendingink}(Rn\,#1–#2)}}
\newcommand{\outrnsingle}[1]{\enspace{\normalfont\footnotesize\color{pendingink}(Rn\,#1)}}
\newcommand{\outrnzh}[2]{\enspace{\normalfont\footnotesize\color{pendingink}（第\,#1–#2\,段）}}
\newcommand{\outrnzhs}[1]{\enspace{\normalfont\footnotesize\color{pendingink}（第\,#1\,段）}}
% ===== 大綱雙語對照表格輔助命令（三欄：段碼 │ 德文 │ 中文）=====
% 欄 1：段碼（由欄定義自動格式化：小字、等寬、右對齊、pendingink）
% 欄 2：德文標籤＋正文
% 欄 3：中文標籤＋正文
%
% 無段碼的列（\omainrow、\osecrow），欄 1 以 & 空置。
% 有段碼的列（\osecrowrn、\osubrow、\ointrow），#1 為預格式化字串
%   例：\osubrow{2–49}{I.}{...}{I.}{...}
%       \ointrow{107}{...}{...}
%
% \opartrow：段組標題（橫跨三欄）
\newcommand{\opartrow}[2]{%
  \noalign{\vspace{0.30em}}%
  \multicolumn{3}{@{}l@{}}{%
    \footnotesize\bfseries\color{pendingink}#1\hfill\songfont#2%
  }\\[0.06em]%
}
% \odashrow：虛線分隔（橫跨三欄）
\newcommand{\odashrow}{%
  \noalign{\vspace{0.28em}}%
  \multicolumn{3}{@{}l@{}}{\color{notegray}%
    \xleaders\hbox to 0.82em{\hfil\rule{0.42em}{0.28pt}\hfil}\hskip 0.88\linewidth}%
  \\[0.24em]%
}
% \odissentpart：不同意見書區段標頭。比一般 \opartrow 更像附錄性轉場。
\newcommand{\odissentpart}[2]{%
  \noalign{\vspace{0.42em}}%
  \multicolumn{3}{@{}p{0.88\textwidth}@{}}{%
    \color{notegray}\rule{\linewidth}{0.18pt}%
  }\\[-0.18em]%
  \multicolumn{3}{@{}p{0.88\textwidth}@{}}{%
    \makebox[\linewidth][s]{%
      \footnotesize\bfseries\color{pendingink}#1%
      \hfill{\songfont #2}%
    }%
  }\\[-0.02em]%
}
% \odissentopinion：不同意見書的署名與一行摘要，右欄保留段碼範圍。
\newcommand{\odissentopinion}[5]{%
  \footnotesize\hangindent=0.52cm\hangafter=1%
    {\bfseries\color{caseink}#2}\enspace{\color{notegray}\textperiodcentered}\enspace\color{caseink}#3%
  &\footnotesize\hangindent=0.52cm\hangafter=1%
    {\songfont\bfseries\color{caseink}#4}\enspace{\color{notegray}\textperiodcentered}\enspace\songfont\color{caseink}#5%
  &#1%
  \\[0.12em]%
}
% \omainrow：頂層項目，無段碼（Leitsätze、Urteil、Gründe …）
\newcommand{\omainrow}[4]{%
  \footnotesize\hangindent=1.92cm\hangafter=1%
    \makebox[1.92cm][l]{\bfseries\color{caseink}#1}\color{caseink}#2%
  &\footnotesize\hangindent=1.92cm\hangafter=1%
    \makebox[1.92cm][l]{\bfseries\color{caseink}#3}\songfont\color{caseink}#4%
  &%
  \\[0.15em]%
}
% \omainrowrn：頂層項目，有段碼（不同意見等標籤寬於 \osecrow 框者）
\newcommand{\omainrowrn}[5]{%
  \footnotesize\hangindent=1.92cm\hangafter=1%
    \makebox[1.92cm][l]{\bfseries\color{caseink}#2}\color{caseink}#3%
  &\footnotesize\hangindent=1.92cm\hangafter=1%
    \makebox[1.92cm][l]{\bfseries\color{caseink}#4}\songfont\color{caseink}#5%
  &#1%
  \\[0.15em]%
}
% \osecrow：一級子項，無段碼（A.、C.、D. 等有子項者）
\newcommand{\osecrow}[4]{%
  \footnotesize\hspace{1.10cm}\hangindent=1.98cm\hangafter=1%
    \makebox[0.86cm][l]{\bfseries\color{caseink}#1}\color{caseink}#2%
  &\footnotesize\hspace{1.10cm}\hangindent=1.98cm\hangafter=1%
    \makebox[0.86cm][l]{\bfseries\color{caseink}#3}\songfont\color{caseink}#4%
  &%
  \\[0.11em]%
}
% \osecrowrn：一級子項，有段碼（B.、E. 等完整段落範圍者）
% #1=段碼字串（如 101–106）, #2=DE標籤, #3=DE正文, #4=ZH標籤, #5=ZH正文
\newcommand{\osecrowrn}[5]{%
  \footnotesize\hspace{1.10cm}\hangindent=1.98cm\hangafter=1%
    \makebox[0.86cm][l]{\bfseries\color{caseink}#2}\color{caseink}#3%
  &\footnotesize\hspace{1.10cm}\hangindent=1.98cm\hangafter=1%
    \makebox[0.86cm][l]{\bfseries\color{caseink}#4}\songfont\color{caseink}#5%
  &#1%
  \\[0.11em]%
}
% \osubrow：二級子項，有段碼範圍（I.、II.、A.I. …）
% #1=段碼字串, #2=DE標籤, #3=DE正文, #4=ZH標籤, #5=ZH正文
\newcommand{\osubrow}[5]{%
  \footnotesize\hspace{1.40cm}\hangindent=2.30cm\hangafter=1%
    \makebox[0.90cm][l]{\bfseries\color{caseink}#2}\color{caseink}#3%
  &\footnotesize\hspace{1.40cm}\hangindent=2.30cm\hangafter=1%
    \makebox[0.90cm][l]{\bfseries\color{caseink}#4}\songfont\color{caseink}#5%
  &#1%
  \\[0.09em]%
}
% \ointrow：引入段落，有單一段碼，無標籤
% #1=段碼字串, #2=DE正文, #3=ZH正文
\newcommand{\ointrow}[3]{%
  \footnotesize\hspace{0.52cm}\hangindent=0.52cm\hangafter=1\color{caseink}#2%
  &\footnotesize\hspace{0.52cm}\hangindent=0.52cm\hangafter=1\songfont\color{caseink}#3%
  &#1%
  \\[0.09em]%
}
% \osubsubrow：三級子項（1.、2.、3.）
% #1=段碼字串, #2=DE標籤, #3=DE正文, #4=ZH標籤, #5=ZH正文
\newcommand{\osubsubrow}[5]{%
  \footnotesize\hspace{1.80cm}\hangindent=2.48cm\hangafter=1%
    \makebox[0.68cm][l]{\bfseries\color{caseink}#2}\color{caseink}#3%
  &\footnotesize\hspace{1.80cm}\hangindent=2.48cm\hangafter=1%
    \makebox[0.68cm][l]{\bfseries\color{caseink}#4}\songfont\color{caseink}#5%
  &#1%
  \\[0.08em]%
}
% \osubsubsubrow：四級子項（a）、b））
% #1=段碼字串, #2=DE標籤, #3=DE正文, #4=ZH標籤, #5=ZH正文
\newcommand{\osubsubsubrow}[5]{%
  \footnotesize\hspace{2.10cm}\hangindent=2.68cm\hangafter=1%
    \makebox[0.58cm][l]{\bfseries\color{caseink}#2}\color{caseink}#3%
  &\footnotesize\hspace{2.10cm}\hangindent=2.68cm\hangafter=1%
    \makebox[0.58cm][l]{\bfseries\color{caseink}#4}\songfont\color{caseink}#5%
  &#1%
  \\[0.07em]%
}
% \osubsubsubsubrow：五級子項（aa）、bb））
% 標籤框 0.62cm：足以容納「aa)」在 footnotesize bold 下（約 0.50cm）並留有間距
% #1=段碼字串, #2=DE標籤, #3=DE正文, #4=ZH標籤, #5=ZH正文
\newcommand{\osubsubsubsubrow}[5]{%
  \footnotesize\hspace{2.38cm}\hangindent=3.06cm\hangafter=1%
    \makebox[0.68cm][l]{\bfseries\color{caseink}#2}\color{caseink}#3%
  &\footnotesize\hspace{2.38cm}\hangindent=3.06cm\hangafter=1%
    \makebox[0.68cm][l]{\bfseries\color{caseink}#4}\songfont\color{caseink}#5%
  &#1%
  \\[0.06em]%
}
% \outlinepage：雙語對照大綱頁包裝環境（longtable 自動跨頁）
% 三欄：德文欄 + 中文欄 + 段碼欄（最右，不折行）
% 段碼欄寬度：1.80cm，足以容納最長段碼如「309–348」
% DE/ZH 欄寬：(0.87\textwidth - 2.08cm) / 2
%   驗算：2×欄 + 0.05\textwidth + 0.28cm gap + 1.80cm = 0.87tw - 2.08cm + 0.05tw + 2.08cm = 0.92tw ✓
\newenvironment{outlinepage}{%
  \vspace*{0.40cm}%
  \setlength{\LTleft}{0.04\textwidth}%
  \setlength{\LTright}{0.04\textwidth}%
  \setlength{\tabcolsep}{0pt}%
  % 大綱列距由 longtable 的 \arraystretch 與各 row macro 結尾的 \\[...em] 共同決定。
  % 這裡控制「每一列本身的表格 strut 高度」；數值越大，A./I./1. 各項之間越像有空行。
  % A3 原本用 0.62，視覺上沒有空行；A4 也沿用同值，避免切到 A4 後項目被撐開。
  % 若日後覺得大綱太擠，只微調這個值（例如 0.68），不要改各 \omainrow/\osubrow 的縮排。
  \renewcommand{\arraystretch}{0.62}%
  \begin{longtable}{%
    >{\vspace{0pt}\raggedright\arraybackslash}p{\dimexpr(0.87\textwidth-2.08cm)/2\relax}%
    @{\hspace{0.025\textwidth}}%
    !{\color{notegray}\vrule width 0.12pt}%
    @{\hspace{0.025\textwidth}}%
    >{\vspace{0pt}\raggedright\arraybackslash}p{\dimexpr(0.87\textwidth-2.08cm)/2\relax}%
    @{\hspace{0.28cm}}%
    >{\vspace{0pt}\footnotesize\normalfont\color{pendingink}\raggedright\arraybackslash}p{1.80cm}%
    @{}%
  }%
  % 首頁欄標籤：不要橫跨置中；分別放在德文欄與中文欄的實際欄位起點。
  \footnotesize\bfseries\color{sourceink}Gliederung%
  &\footnotesize\bfseries\songfont\color{sourceink}大綱%
  &%
  \\[0.36em]%
  \endfirsthead%
  % 續頁欄標籤：同樣分欄定位，以灰色細體區分；無需標注「續」。
  \footnotesize\color{pendingink}Gliederung%
  &\footnotesize\songfont\color{pendingink}大綱%
  &%
  \\[0.24em]%
  \endhead%
}{%
  \end{longtable}%
}
 之前設定；
%   預設 chinese / a4）
\ifdefined\docmode\else\def\docmode{chinese}\fi
\ifdefined\bilingualpaper\else\def\bilingualpaper{a4}\fi
\newif\ifshoworiginal
\newif\ifshowtranslation
\newif\ifbilinguallandscape
\newif\ifbilingualcolumns
\newif\ifintranslation
\newif\iffirstbilingualsection

\showoriginalfalse
\showtranslationfalse
\bilinguallandscapefalse
\bilingualcolumnsfalse
\intranslationfalse
\firstbilingualsectionfalse

\def\modechinese{chinese}
\def\modegerman{german}
\def\modebilingual{bilingual}
\def\bilingualpaperaThree{a3}
\def\bilingualpaperaFour{a4}

\ifx\docmode\modechinese
  \showtranslationtrue
\fi

\ifx\docmode\modegerman
  \showoriginaltrue
\fi

\ifx\docmode\modebilingual
  \showoriginaltrue
  \showtranslationtrue
  \bilinguallandscapetrue
  \bilingualcolumnstrue
\fi

% 編排模式說明：
% chinese  → \showoriginalfalse  \showtranslationtrue  （A4 直式，純中文）
% german   → \showoriginaltrue   \showtranslationfalse （A4 直式，純德文）
% bilingual→ \showoriginaltrue   \showtranslationtrue  \bilingualcolumnstrue（A3/A4 橫式對照）

\ifbilingualcolumns
  \ifx\bilingualpaper\bilingualpaperaFour
    \geometry{
      a4paper,
      landscape,
      left=1.25cm,
      right=2.25cm,
      top=1.25cm,
      bottom=1.75cm,
      footskip=8.5mm,
      marginparwidth=0.75cm,
      marginparsep=0.25cm
    }
  \else
    \geometry{
      a3paper,
      landscape,
      left=1.55cm,
      right=2.75cm,
      top=1.55cm,
      bottom=2.05cm,
      footskip=9.5mm,
      marginparwidth=0.9cm,
      marginparsep=0.35cm
    }
  \fi
\else
\geometry{
  a4paper,
  portrait,
  left=2.2cm,
  right=2.6cm,
  top=2.0cm,
  bottom=2.0cm,
  marginparwidth=0.8cm,
  marginparsep=0.3cm
}
\fi
% ===== 時間戳基礎設施 =====
\newcount\stamptimeval
\newcount\stampminuteval
\newcommand{\stamphh}{%
  \stamptimeval=\time
  \divide\stamptimeval by 60
  \ifnum\stamptimeval<10 0\fi
  \the\stamptimeval}
\newcommand{\stampmm}{%
  \stamptimeval=\time
  \stampminuteval=\time
  \divide\stampminuteval by 60
  \multiply\stampminuteval by 60
  \advance\stamptimeval by -\stampminuteval
  \ifnum\stamptimeval<10 0\fi
  \the\stamptimeval}
\newcommand{\stampwkde}[1]{%
  \ifcase#1Montag\or Dienstag\or Mittwoch\or Donnerstag\or Freitag\or Samstag\or Sonntag\fi}
\newcommand{\stampwkzh}[1]{%
  \ifcase#1 一\or 二\or 三\or 四\or 五\or 六\or 日\fi}
\newcommand{\stampmonthde}[1]{%
  \ifcase#1\or Januar\or Februar\or März\or April\or Mai\or Juni\or
  Juli\or August\or September\or Oktober\or November\or Dezember\fi}
\newcount\stampjdval
\newcount\stampwdval
\pgfcalendardatetojulian{\the\year-\the\month-\the\day}{\stampjdval}
\pgfcalendarjuliantoweekday{\the\stampjdval}{\stampwdval}
\newcommand{\timestampzh}{%
  \songfont 最後修改：\the\year 年\the\month 月\the\day 日%
  % （星期\stampwkzh{\the\stampwdval}）%
  % \space\stamphh:\stampmm
  }

\newcommand{\documentdatezh}{\the\year 年 \the\month 月 \the\day 日}
\newcommand{\documentdatede}{\the\day. \stampmonthde{\the\month} \the\year}
\newcommand{\bverfgetranslationnote}{%
  {\songfont 翻譯說明：初譯使用 Claude Code 與 GPT 協助迭代；仍請以德文原文為準。}%
}
\newcommand{\bverfgecourseinfo}{%
  {\songfont 課程：114 學年度第 2 學期；憲法專題研究（四）；授課老師：湯德宗教授；導讀日期：115 年 6 月 1 日。}%
}
\newcommand{\bverfgeeditorinfo}{%
  {\songfont 編者：王逸帆（R10A21126），國立台灣大學法律系研究所碩士班。}%
}
\newcommand{\bverfgetitleversionline}{%
  \ifbilingualcolumns
    Fassung \documentversion\enspace\textperiodcentered\enspace Stand: \documentdatede
    \quad
    {\songfont 版本 \documentversion\enspace\textperiodcentered\enspace 修訂日期：\documentdatezh\par}%
  \else
    \ifshowtranslation
      {\songfont 版本 \documentversion\enspace\textperiodcentered\enspace 修訂日期：\documentdatezh\par}%
    \else
      Fassung \documentversion\enspace\textperiodcentered\enspace Stand: \documentdatede\par
    \fi
  \fi
}
\newcommand{\bverfgetitlecornerinfo}{%
  \ifshowtranslation
    \vspace{0.72em}%
    \noindent\hfill
    \begin{minipage}{0.66\textwidth}
      \raggedleft\tiny\color{pendingink}\songfont
      \bverfgecourseinfo\par
      \bverfgeeditorinfo
    \end{minipage}%
  \fi
}
\newcommand{\bverfgetitlemeta}{%
  \vfill
  \begin{center}%
    \footnotesize\color{pendingink}%
    \bverfgetitleversionline
    \ifshowtranslation
      \vspace{0.28em}{\scriptsize\bverfgetranslationnote\par}%
    \fi
  \end{center}%
  \bverfgetitlecornerinfo
}

\pagestyle{fancy}
\fancyhf{}
\renewcommand{\headrulewidth}{0pt}
\renewcommand{\footrulewidth}{0pt}
\fancyfoot[C]{\thepage}
\ifbilingualcolumns
  \fancyfoot[R]{\scriptsize\color{pendingink}\timestampzh}%
\else
  \ifshoworiginal
  \else
    \fancyfoot[R]{\scriptsize\color{pendingink}\timestampzh\hspace{0.45cm}}%
  \fi
\fi
\setmainfont{Times New Roman}
\setsansfont{Helvetica Neue}
\setmonofont{Menlo}
\IfFileExists{/Users/iw/Library/Fonts/kaiu.ttf}{
  \setCJKmainfont[Path=/Users/iw/Library/Fonts/,AutoFakeBold=4]{kaiu.ttf}
}{
  \IfFontExistsTF{標楷體}{
    \setCJKmainfont[AutoFakeBold=4]{標楷體}
  }{
    \setCJKmainfont{Songti TC}
  }
}
\IfFontExistsTF{Songti TC}{
  \setCJKfamilyfont{song}{Songti TC}
  \setCJKmonofont{Songti TC}
}{
  \setCJKfamilyfont{song}[Path=/Users/iw/Library/Fonts/,AutoFakeBold=4]{kaiu.ttf}
  \setCJKmonofont[Path=/Users/iw/Library/Fonts/,AutoFakeBold=4]{kaiu.ttf}
}
\newcommand{\songfont}{\CJKfamily{song}}

% ===== 封面／目錄專用打字機字體（僅限封面、目錄頁使用，不影響正文）=====
\IfFontExistsTF{Radio Newsman}{%
  \newfontfamily\bverfgetitlede{Radio Newsman}%
}{%
  \newfontfamily\bverfgetitlede{Courier New}%
}
\IfFontExistsTF{Erikas Farbband}{%
  \newfontfamily\bverfgebodyde{Erikas Farbband}%
}{%
  \let\bverfgebodyde\bverfgetitlede
}
\IfFontExistsTF{Maoken Heavy Labourer}{%
  \setCJKfamilyfont{bverfgetitlecjk}{Maoken Heavy Labourer}%
}{%
  \setCJKfamilyfont{bverfgetitlecjk}{Songti TC}%
}
\IfFontExistsTF{Huiwen-mincho}{%
  \setCJKfamilyfont{bverfgebodycjk}{Huiwen-mincho}%
}{%
  \setCJKfamilyfont{bverfgebodycjk}{Songti TC}%
}
\newcommand{\bverfgetitlezh}{\bverfgetitlede\CJKfamily{bverfgetitlecjk}}
\newcommand{\bverfgebodyzh}{\bverfgebodyde\CJKfamily{bverfgebodycjk}}

\setstretch{1.35}
\setlist{itemsep=0.2em, topsep=0.4em}
\setlength{\parindent}{0pt}
\setlength{\parskip}{0.45em}
\raggedbottom
\setcounter{secnumdepth}{0}
\setcounter{tocdepth}{2}

\newcommand{\lawboxsourcefont}{\footnotesize}
\newcommand{\lawboxtransfont}{\footnotesize}
\newcommand{\lawboxsourcestretch}{1.02}
\newcommand{\lawboxtransstretch}{0.92}

\titleformat{\section}{\centering\Large\bfseries\songfont}{\thesection}{1em}{}
\titleformat{\subsection}{\large\bfseries\songfont}{\thesubsection}{1em}{}
\titleformat{\subsubsection}{\normalsize\bfseries\songfont}{\thesubsubsection}{1em}{}
\titleformat{\paragraph}{\normalsize\bfseries\songfont}{\theparagraph}{1em}{}
\titlespacing*{\section}{0pt}{1.8em}{1.0em}
\titlespacing*{\subsection}{0pt}{1.2em}{0.6em}
\titlespacing*{\subsubsection}{0pt}{1.0em}{0.45em}
\titlespacing*{\paragraph}{0pt}{0.6em}{0.4em}

\definecolor{noteblue}{HTML}{A9C9F5}
\definecolor{noteteal}{HTML}{AEE1D6}
\definecolor{notegray}{HTML}{D7DADF}
\definecolor{caseink}{HTML}{000000}
\definecolor{casepaper}{HTML}{FCFEFD}
\definecolor{sourcepaper}{HTML}{FAFBF8}
\definecolor{sourceline}{HTML}{D7DED8}
\definecolor{sourceink}{HTML}{000000}
\definecolor{pendingink}{HTML}{000000}

\tcbset{
  caseboxbase/.style={
    enhanced,
    breakable,
    width=0.985\linewidth,
    colback=casepaper,
    coltitle=caseink,
    fonttitle=\bfseries\songfont,
    boxrule=0.28pt,
    arc=1.2mm,
    left=2.4mm,
    right=2.4mm,
    top=1.5mm,
    bottom=1.5mm,
    boxsep=1.5mm,
    before skip=0.65em,
    after skip=0.65em,
    before upper={\setlength{\parindent}{0pt}\setlength{\parskip}{0.35em}},
  }
}

\newcommand{\bilingualstart}{}
\newcommand{\bilingualstop}{}
\ifbilingualcolumns
  \renewcommand{\bilingualstart}{\small\setlength{\columnsep}{1.25cm}\setlength{\columnseprule}{0.12pt}\columnratio{0.5,0.5}\multicolumnfootnotes\flushbottom\sloppy\interlinepenalty=80\clubpenalty=800\widowpenalty=800\displaywidowpenalty=800\begin{paracol}{2}\global\intranslationfalse\global\firstbilingualsectiontrue{\footnotesize\color{sourceink}Deutsch\par}\switchcolumn{\footnotesize\songfont\color{sourceink}中文譯文\par}\switchcolumn*}
  \renewcommand{\bilingualstop}{\ifintranslation\switchcolumn*\global\intranslationfalse\fi\end{paracol}\clearpage\raggedbottom}
\fi
\newif\ifrandnumbers
\randnumbersfalse
\newcounter{randnumber}
\newcounter{randnumberlocal}
\newcounter{lastsourcerandstart}
\newif\iflastsourcehasrandnumbers
\lastsourcehasrandnumbersfalse
\newif\ifinlawbox
\inlawboxfalse
\newif\iflawboxhaspendingrn
\lawboxhaspendingrnfalse
\newcommand{\lawboxpendingrn}{}
\newcommand{\startrandnumbers}{\global\randnumberstrue\setcounter{randnumber}{1}\global\lastsourcehasrandnumbersfalse}
\newcommand{\stoprandnumbers}{\global\randnumbersfalse\global\lastsourcehasrandnumbersfalse}
\newlength{\randnumberwidth}
\setlength{\randnumberwidth}{0.95cm}
\newlength{\randnumberrightoffset}
\setlength{\randnumberrightoffset}{0.85cm}
\newlength{\randnumberlawboxrightoffset}
\setlength{\randnumberlawboxrightoffset}{1.40cm}
\newcommand{\randnumberbox}[1]{%
  \leavevmode
  \makebox[0pt][l]{%
    \smash{%
      \tikz[remember picture,overlay,baseline=(rnmark.base)]{%
        \coordinate (rnmark) at (0,0);%
        \ifinlawbox
          \path let \p1=(current page.east), \p2=(rnmark.base) in
            node[
              anchor=base east,
              inner sep=0pt,
              outer sep=0pt,
              text width=\randnumberwidth,
              align=right,
              text=pendingink,
              font=\normalfont\mdseries\scriptsize\ttfamily
            ] at (\x1-\randnumberlawboxrightoffset,\y2) {{\normalfont\mdseries\scriptsize\ttfamily #1}};%
        \else
          \path let \p1=(current page.east), \p2=(rnmark.base) in
            node[
              anchor=base east,
              inner sep=0pt,
              outer sep=0pt,
              text width=\randnumberwidth,
              align=right,
              text=pendingink,
              font=\normalfont\mdseries\scriptsize\ttfamily
            ] at (\x1-\randnumberrightoffset,\y2) {{\normalfont\mdseries\scriptsize\ttfamily #1}};%
        \fi
      }%
    }%
  }%
  \ignorespaces
}
\newcommand{\lawboxstorerandnumber}[1]{%
  \gdef\lawboxpendingrn{#1}%
  \global\lawboxhaspendingrntrue
  \ignorespaces
}
\newcommand{\lawboxuserandnumber}{%
  \iflawboxhaspendingrn
    \randnumberbox{\lawboxpendingrn}%
    \global\lawboxhaspendingrnfalse
  \fi
}
\newcommand{\sourcern}[1]{%
  \ifrandnumbers
    \ifshoworiginal
      \ifshowtranslation\else
        \ifinlawbox\lawboxstorerandnumber{#1}\else\randnumberbox{#1}\fi
      \fi
    \fi
  \fi
}
\newcommand{\transrn}[1]{%
  \ifrandnumbers
    \ifshowtranslation
      \ifinlawbox\lawboxstorerandnumber{#1}\else\randnumberbox{#1}\fi
    \fi
  \fi
}
% ===== 課堂模式：德文原文縮寫註解 =====
% 日常切換只改各判決檔的一行：
%   \classroommodeon        開啟課堂版；註解掉這行即回到一般版。
% 模板預設：
%   1. 課堂模式關閉。
%   2. 課堂模式開啟時，縮寫註解包含「德文全稱＋中文譯名」。
%   3. 課堂模式開啟時，GG/條文編者註仍會自動顯示。
% 進階微調才需要在單案檔覆寫：
%   \classroomabbrzhoff     縮寫註解僅顯示德文全稱。
%   \showeditornotestrue    一般版也顯示編者註。
% 註解策略：同一實際 PDF 頁面中，同一縮寫只在第一次出現處加註。
% 這裡用 zref-abspage 的絕對頁序，不用 \thepage，避免文件內頁碼重置時誤判。
\newif\ifclassroommode
\newif\ifclassroomabbrzh
\classroommodefalse
\classroomabbrzhtrue
\newcommand{\classroommodeon}{\classroommodetrue}
\newcommand{\classroommodeoff}{\classroommodefalse}
\newcommand{\classroomabbrzhon}{\classroomabbrzhtrue}
\newcommand{\classroomabbrzhoff}{\classroomabbrzhfalse}
\newcommand{\abbrnotelabel}{\emph{Abk.}: }
\newcommand{\abbrnote}[3]{%
  \ifclassroommode
    \ifshoworiginal
      \ifintranslation\else
        \footnote{\normalfont\footnotesize\abbrnotelabel #2\ifclassroomabbrzh；{\songfont #3}\fi}%
      \fi
    \fi
  \fi
}
\newcommand{\abbrnoteonce}[4]{%
  #2%
  \ifclassroommode
    \ifshoworiginal
      \ifintranslation\else
        \begingroup
          \edef\abbrcurrentabspage{\number\numexpr\value{abspage}+1\relax}%
          \ifcsname abbrnoteseen@#1@\abbrcurrentabspage\endcsname\else
            \global\expandafter\let\csname abbrnoteseen@#1@\abbrcurrentabspage\endcsname\empty
            \abbrnote{#1}{#3}{#4}%
          \fi
        \endgroup
      \fi
    \fi
  \fi
}
\newcommand{\abbrAbs}{\abbrnoteonce{abs}{Abs.}{Absatz}{項}}
\newcommand{\abbrAAO}{\abbrnoteonce{aao}{a.a.O.}{am angegebenen Ort}{前引處}}
\newcommand{\abbrAE}{\abbrnoteonce{ae}{AE}{Alternativ-Entwurf}{替代草案}}
\newcommand{\abbrArt}{\abbrnoteonce{art}{Art.}{Artikel}{條}}
\newcommand{\abbrAufl}{\abbrnoteonce{aufl}{Aufl.}{Auflage}{版}}
\newcommand{\abbrBd}{\abbrnoteonce{bd}{Bd.}{Band}{卷}}
\newcommand{\abbrBGBl}{\abbrnoteonce{bgbl}{BGBl.}{Bundesgesetzblatt}{聯邦法律公報}}
\newcommand{\abbrBGHSt}{\abbrnoteonce{bghst}{BGHSt}{Entscheidungen des Bundesgerichtshofes in Strafsachen}{聯邦普通法院刑事判決彙編}}
\newcommand{\abbrBRDrucks}{\abbrnoteonce{brdrucks}{BRDrucks.}{Bundesrats-Drucksache}{聯邦參議院文件}}
\newcommand{\abbrBTDrucks}{\abbrnoteonce{btdrucks}{BTDrucks.}{Bundestags-Drucksache}{聯邦議院文件}}
\newcommand{\abbrBundesgesetzbl}{\abbrnoteonce{bundesgesetzbl}{Bundesgesetzbl.}{Bundesgesetzblatt}{聯邦法律公報}}
\newcommand{\abbrBvF}{\abbrnoteonce{bvf}{BvF}{Bundesverfassungsgerichtsverfahren}{聯邦憲法法院程序案號}}
\newcommand{\abbrBvL}{\abbrnoteonce{bvl}{BvL}{Bundesverfassungsgerichtsverfahren}{聯邦憲法法院程序案號}}
\newcommand{\abbrBvR}{\abbrnoteonce{bvr}{BvR}{Bundesverfassungsgerichtsverfahren}{聯邦憲法法院程序案號}}
\newcommand{\abbrBVerfG}{\abbrnoteonce{bverfg}{BVerfG}{Bundesverfassungsgericht}{聯邦憲法法院}}
\newcommand{\abbrBVerfGE}{\abbrnoteonce{bverfge}{BVerfGE}{Entscheidungen des Bundesverfassungsgerichts}{聯邦憲法法院判決集}}
\newcommand{\abbrBVerfGG}{\abbrnoteonce{bverfgg}{BVerfGG}{Bundesverfassungsgerichtsgesetz}{聯邦憲法法院法}}
\newcommand{\abbrDDR}{\abbrnoteonce{ddr}{DDR}{Deutsche Demokratische Republik}{德意志民主共和國；東德}}
\newcommand{\abbrDJT}{\abbrnoteonce{djt}{DJT}{Deutscher Juristentag}{德國法學家大會}}
\newcommand{\abbrDr}{\abbrnoteonce{dr}{Dr.}{Doktor}{Dr.}}
\newcommand{\abbrEuGRZ}{\abbrnoteonce{eugrz}{EuGRZ}{Europäische Grundrechte-Zeitschrift}{歐洲基本權期刊}}
\newcommand{\abbrEV}{\abbrnoteonce{ev}{EV}{Einigungsvertrag}{統一條約}}
\newcommand{\abbrF}{\abbrnoteonce{f}{f.}{folgende}{以下一頁；下一段}}
\newcommand{\abbrFF}{\abbrnoteonce{ff}{ff.}{fortfolgende}{以下數頁；以下各段}}
\newcommand{\abbrGBlDDR}{\abbrnoteonce{gblddr}{GBl. DDR}{Gesetzblatt der Deutschen Demokratischen Republik}{德意志民主共和國法律公報}}
\newcommand{\abbrGG}{\abbrnoteonce{gg}{GG}{Grundgesetz}{基本法}}
\newcommand{\abbrGVBl}{\abbrnoteonce{gvbl}{GVBl.}{Gesetz- und Verordnungsblatt}{法律及命令公報}}
\newcommand{\abbrJR}{\abbrnoteonce{jr}{JR}{Juristische Rundschau}{法學評論}}
\newcommand{\abbrJZ}{\abbrnoteonce{jz}{JZ}{JuristenZeitung}{法學家報}}
\newcommand{\abbrIVm}{\abbrnoteonce{ivm}{i.V.m.}{in Verbindung mit}{結合；連同}}
\newcommand{\abbrMdB}{\abbrnoteonce{mdb}{MdB}{Mitglied des Bundestages}{聯邦議院議員}}
\newcommand{\abbrMwN}{\abbrnoteonce{mwn}{m.w.N.}{mit weiteren Nachweisen}{附進一步文獻／判例出處}}
\newcommand{\abbrNF}{\abbrnoteonce{nf}{n.F.}{neue Fassung}{新版本}}
\newcommand{\abbrAF}{\abbrnoteonce{af}{a.F.}{alte Fassung}{舊版本}}
\newcommand{\abbrNr}{\abbrnoteonce{nr}{Nr.}{Nummer}{款、目或號，依語境}}
\newcommand{\abbrProf}{\abbrnoteonce{prof}{Prof.}{Professor}{教授}}
\newcommand{\abbrRdnr}{\abbrnoteonce{rdnr}{Rdnr.}{Randnummer}{邊碼；段碼}}
\newcommand{\abbrRdnrn}{\abbrnoteonce{rdnrn}{Rdnrn.}{Randnummern}{邊碼；段碼}}
\newcommand{\abbrRGBl}{\abbrnoteonce{rgbl}{RGBl.}{Reichsgesetzblatt}{帝國法律公報}}
\newcommand{\abbrRGSt}{\abbrnoteonce{rgst}{RGSt}{Entscheidungen des Reichsgerichts in Strafsachen}{帝國法院刑事判決彙編}}
\newcommand{\abbrRspr}{\abbrnoteonce{rspr}{Rspr.}{Rechtsprechung}{判例；司法實務}}
\newcommand{\abbrRVO}{\abbrnoteonce{rvo}{RVO}{Reichsversicherungsordnung}{帝國保險法}}
\newcommand{\abbrS}{\abbrnoteonce{s}{S.}{Seite}{頁}}
\newcommand{\abbrSFHG}{\abbrnoteonce{sfhg}{SFHG}{Schwangeren- und Familienhilfegesetz}{孕婦及家庭扶助法}}
\newcommand{\abbrSGBV}{\abbrnoteonce{sgbv}{SGB V}{Fünftes Buch Sozialgesetzbuch}{社會法典第五編}}
\newcommand{\abbrStAG}{\abbrnoteonce{staeg}{StÄG}{Strafrechtsänderungsgesetz}{刑法修正法}}
\newcommand{\abbrStenBer}{\abbrnoteonce{stenber}{StenBer.}{Stenographischer Bericht}{速記紀錄}}
\newcommand{\abbrStGB}{\abbrnoteonce{stgb}{StGB}{Strafgesetzbuch}{刑法}}
\newcommand{\abbrStREG}{\abbrnoteonce{streg}{StREG}{Strafrechtsreform-Ergänzungsgesetz}{刑法改革補充法}}
\newcommand{\abbrStrRG}{\abbrnoteonce{strrg}{StrRG}{Strafrechtsreformgesetz}{刑法改革法}}
\newcommand{\abbrUS}{\abbrnoteonce{us}{U.S.}{United States Reports}{美國最高法院判決彙編}}
\newcommand{\abbrVgl}{\abbrnoteonce{vgl}{vgl.}{vergleiche}{參見}}
\newcommand{\abbrVVDStRL}{\abbrnoteonce{vvdstrl}{VVDStRL}{Veröffentlichungen der Vereinigung der Deutschen Staatsrechtslehrer}{德國公法教師協會論文集}}
\newcommand{\abbrWP}{\abbrnoteonce{wp}{WP.}{Wahlperiode}{屆期}}
\newcommand{\abbrWp}{\abbrnoteonce{wp}{Wp.}{Wahlperiode}{屆期}}
\newcommand{\abbrZStrW}{\abbrnoteonce{zstrw}{ZStrW}{Zeitschrift für die gesamte Strafrechtswissenschaft}{整體刑法學期刊}}
\newcommand{\recordsourcerandstart}{%
  \ifrandnumbers
    \setcounter{lastsourcerandstart}{\value{randnumber}}%
    \global\lastsourcehasrandnumberstrue
  \else
    \global\lastsourcehasrandnumbersfalse
  \fi
}
\newlength{\biparagraphskip}
\setlength{\biparagraphskip}{0.52em}
\newlength{\lawboxblockbeforeskip}
\setlength{\lawboxblockbeforeskip}{0.72em}
\newlength{\lawboxblockafterskip}
\setlength{\lawboxblockafterskip}{0.92em}
\newif\iflawboxpartendedblock
\lawboxpartendedblockfalse
\newcommand{\sourceparstyle}{\footnotesize\color{sourceink}\setstretch{1.12}\RaggedRight\emergencystretch=3em\setlength{\parskip}{0.42em}\obeylines\selectfont}
\newcommand{\transparstyle}{\songfont\color{caseink}\setstretch{1.12}\setlength{\parindent}{0pt}\setlength{\parskip}{0.42em}}
\newcommand{\bilingualpairskip}{%
  \par\addvspace{\biparagraphskip}%
  \switchcolumn
  \par\addvspace{\biparagraphskip}%
  \switchcolumn*%
  \global\intranslationfalse
}
\newcommand{\bilingualboxpairskip}{%
  \par
  \switchcolumn
  \par
  \switchcolumn*%
  \global\intranslationfalse
}
\newcommand{\bilinguallawboxblockskip}{%
  \switchcolumn*%
  \par\addvspace{\lawboxblockafterskip}%
  \switchcolumn
  \par\addvspace{\lawboxblockafterskip}%
  \switchcolumn*%
  \global\intranslationfalse
}
\newcommand{\bikeepwithnext}[1][4]{%
  \ifbilingualcolumns
    \syncleft
    \Needspace{#1\baselineskip}%
    \switchcolumn
    \Needspace{#1\baselineskip}%
    \switchcolumn*%
    \global\intranslationfalse
  \else
    \Needspace{#1\baselineskip}%
  \fi
}
\NewEnviron{sourcepara}[1][]{
  \recordsourcerandstart
  \ifshoworiginal
    \ifbilingualcolumns
      \ifintranslation\switchcolumn*\global\intranslationfalse\fi
    \else
      \Needspace{5\baselineskip}
    \fi
    \ifbilingualcolumns\else\par\addvspace{0.35em}\fi
    {\sourceparstyle
    \noindent\BODY\par}
    \ifbilingualcolumns\else\addvspace{0.25em}\fi
    \ifbilingualcolumns
      \switchcolumn
      \global\intranslationtrue
    \fi
  \else
    \ifrandnumbers
      \begingroup
        \setbox0=\vbox{\hsize=\linewidth\sourceparstyle\noindent\BODY\par}%
      \endgroup
    \fi
  \fi
}
\NewEnviron{transpara}{
  \ifshowtranslation
    \setcounter{randnumberlocal}{\value{lastsourcerandstart}}%
    \ifbilingualcolumns
      \ifintranslation\else\switchcolumn\global\intranslationtrue\fi
      {\transparstyle\setlength{\leftskip}{0.55em}\setlength{\rightskip}{0.15em plus 1em}\addtolength{\linewidth}{-0.55em}\BODY\par}
      \switchcolumn*
      \bilingualpairskip
    \else
      {\transparstyle\BODY\par}
    \fi
  \else
    \ifbilingualcolumns
      \ifintranslation\switchcolumn*\global\intranslationfalse\fi
    \fi
  \fi
}
\NewEnviron{sourceboxpara}[1][]{
  \global\lawboxpartendedblockfalse
  \recordsourcerandstart
  \ifshoworiginal
    \ifbilingualcolumns
      \ifintranslation\switchcolumn*\global\intranslationfalse\fi
    \else
      \Needspace{3\baselineskip}
    \fi
    {\sourceparstyle
    \BODY\par}
    \ifbilingualcolumns\else
      \iflawboxpartendedblock\addvspace{\lawboxblockafterskip}\fi
    \fi
    \ifbilingualcolumns
      \switchcolumn
      \global\intranslationtrue
    \fi
  \fi
}
\NewEnviron{transboxpara}{
  \global\lawboxpartendedblockfalse
  \ifshowtranslation
    \setcounter{randnumberlocal}{\value{lastsourcerandstart}}%
    \ifbilingualcolumns
      \ifintranslation\else\switchcolumn\global\intranslationtrue\fi
      {\transparstyle\BODY\par}
      \iflawboxpartendedblock
        \bilinguallawboxblockskip
      \else
        \switchcolumn*
        \bilingualboxpairskip
      \fi
    \else
      {\transparstyle\BODY\par}
      \iflawboxpartendedblock\addvspace{\lawboxblockafterskip}\fi
    \fi
  \else
    \ifbilingualcolumns
      \ifintranslation\switchcolumn*\global\intranslationfalse\fi
    \fi
  \fi
}
\newcommand{\leitsatznum}[1]{\noindent\hangindent=3.0em\hangafter=1\makebox[2.25em][r]{\bfseries #1.}\hspace{0.55em}\ignorespaces}
\newcommand{\syncleft}{\ifbilingualcolumns\ifintranslation\switchcolumn*\global\intranslationfalse\fi\fi}
\pretocmd{\section}{\syncleft\clearpage}{}{}
\pretocmd{\subsection}{\syncleft}{}{}
\pretocmd{\subsubsection}{\syncleft}{}{}
\newcommand{\biheadingfont}{\songfont}
\newcommand{\bitocentry}[2]{%
  \ifshoworiginal
    \ifshowtranslation
      \ifstrequal{#1}{#2}{#1}{#1\space #2}%
    \else
      #1%
    \fi
  \else
    #2%
  \fi
}
\pdfstringdefDisableCommands{%
  \def\bitocentry#1#2{#1}%
}
\newcommand{\biheading}[7]{%
  \syncleft#1\bikeepwithnext[#3]\phantomsection\addcontentsline{toc}{#2}{\bitocentry{#6}{#7}}%
  \ifbilingualcolumns
    {#4 #6\par}%
    \switchcolumn
    {#4\biheadingfont #7\par}%
    \switchcolumn*%
    \global\intranslationfalse
    \par\addvspace{#5}%
    \switchcolumn
    \par\addvspace{#5}%
    \switchcolumn*%
  \else
    \ifshoworiginal{#4 #6\par}\fi
    \ifshowtranslation{#4\biheadingfont #7\par}\fi
    \addvspace{#5}%
  \fi
}
\newcommand{\termsectionheading}[1]{%
  \par\Needspace{9\baselineskip}%
  \vspace{0.65em}%
  {\songfont\bfseries #1\par}%
  \vspace{0.2em}%
}
% 章節換頁：
%   雙語 paracol 欄內不要用 \clearpage；它會在同步兩欄時吐出只有欄線/頁碼的空頁。
%   \newpage 足以讓下一個 section 起新頁，又不會強制清空所有待排材料。
%   單語模式不在 paracol 內，仍可用 \clearpage。
\newcommand{\bisectionpagebreak}{\ifbilingualcolumns\newpage\else\clearpage\fi}
\newcommand{\bisection}[2]{%
  \biheading{\iffirstbilingualsection\global\firstbilingualsectionfalse\else\bisectionpagebreak\fi}{section}{6}{\centering\Large\bfseries}{0.8em}{#1}{#2}%
}
\newcommand{\bisubsection}[2]{%
  \biheading{}{subsection}{5}{\large\bfseries}{0.55em}{#1}{#2}%
}
\newcommand{\bisubsubsection}[2]{%
  \biheading{}{subsubsection}{4}{\normalsize\bfseries}{0.45em}{#1}{#2}%
}
\newcommand{\bicompositeheading}[4]{%
  \biheading{}{subsection}{5}{\large\bfseries}{0.16em}{#1}{#2}%
  \biheading{}{subsubsection}{3}{\normalsize\bfseries}{0.45em}{#3}{#4}%
}
\newif\iflawboxfirstitem
\lawboxfirstitemfalse
\newcommand{\lawboxprespace}[1]{%
  \par
  \iflawboxfirstitem
    \global\lawboxfirstitemfalse
  \else
    \addvspace{#1}%
  \fi
}
\newtcolorbox{lawboxinner}{
  caseboxbase,
  colframe=sourceline!85!black,
  colback=white,
  left=1.05em,
  right=1.05em,
  top=0.62em,
  bottom=0.62em,
  before upper={\global\inlawboxtrue\global\lawboxfirstitemtrue\global\lawboxhaspendingrnfalse\ifintranslation\lawboxtransfont\else\lawboxsourcefont\fi\RaggedRight\setlength{\parindent}{0pt}\setlength{\parskip}{0pt}\ifintranslation\setstretch{\lawboxtransstretch}\else\setstretch{\lawboxsourcestretch}\fi},
  after upper={\global\lawboxhaspendingrnfalse\global\inlawboxfalse}
}
\tcbset{
  lawboxpartbase/.style={
    enhanced,
    breakable=false,
    width=0.985\linewidth,
    colback=white,
    frame hidden,
    boxrule=0pt,
    arc=0pt,
    left=1.05em,
    right=1.05em,
    boxsep=1.5mm,
    before skip=0pt,
    after skip=0pt,
    before upper={\global\inlawboxtrue\global\lawboxfirstitemtrue\global\lawboxhaspendingrnfalse\ifintranslation\lawboxtransfont\else\lawboxsourcefont\fi\RaggedRight\setlength{\parindent}{0pt}\setlength{\parskip}{0pt}\ifintranslation\setstretch{\lawboxtransstretch}\else\setstretch{\lawboxsourcestretch}\fi},
    after upper={\global\lawboxhaspendingrnfalse\global\inlawboxfalse},
    borderline west={0.28pt}{0pt}{sourceline!85!black},
    borderline east={0.28pt}{0pt}{sourceline!85!black},
  },
  lawboxpartfirst/.style={
    lawboxpartbase,
    top=0.62em,
    bottom=0.04em,
    before skip=\lawboxblockbeforeskip,
    borderline north={0.28pt}{0pt}{sourceline!85!black},
  },
  lawboxpartmiddle/.style={
    lawboxpartbase,
    top=0.04em,
    bottom=0.04em,
  },
  lawboxpartlast/.style={
    lawboxpartbase,
    top=0.04em,
    bottom=0.62em,
    after skip=0pt,
    borderline south={0.28pt}{0pt}{sourceline!85!black},
  }
}
\newtcolorbox{lawboxpartinner}[1]{#1}
\newtcolorbox{termboxinner}{
  caseboxbase,
  colframe=noteblue!70!black,
  colback=casepaper,
  before upper={\songfont\setlength{\parindent}{0pt}\setlength{\parskip}{0.35em}}
}
\newtcolorbox{deboxinner}{
  caseboxbase,
  colframe=notegray!72!black,
  colback=white,
  left=1.05em,
  right=1.05em,
  top=0.62em,
  bottom=0.62em,
  before upper={\global\inlawboxtrue\global\lawboxfirstitemtrue\global\lawboxhaspendingrnfalse\small\setlength{\parindent}{0pt}\setlength{\parskip}{0.35em}},
  after upper={\global\lawboxhaspendingrnfalse\global\inlawboxfalse}
}
\NewEnviron{lawbox}{
  \begin{lawboxinner}\BODY\end{lawboxinner}
}
\NewEnviron{lawboxpart}[2]{
  \begin{lawboxpartinner}{lawboxpart#1,equal height group=#2}\BODY\end{lawboxpartinner}%
  \ifstrequal{#1}{last}{\global\lawboxpartendedblocktrue}{}%
}
\NewEnviron{termbox}{
  \par\addvspace{0.55em}
  {\color{noteblue!70!black}\rule{\linewidth}{0.28pt}}\par
  \begingroup
    \songfont\small\color{caseink}
    \setlength{\parindent}{0pt}
    \setlength{\parskip}{0.24em}
    \BODY
    \par
  \endgroup
  {\color{noteblue!70!black}\rule{\linewidth}{0.28pt}}\par
  \addvspace{0.55em}
}
\NewEnviron{debox}{
  \begin{deboxinner}\BODY\end{deboxinner}
}
\providecommand{\paragraphmark}[1]{}
\newcommand{\lawboxtitleafter}{\addvspace{0.06em}}
\newcommand{\lawtitle}[1]{\lawboxprespace{0.22em}{\lawboxtransfont\bfseries\lawboxuserandnumber #1}\par\lawboxtitleafter}
\newcommand{\absno}[2]{\lawboxprespace{0.04em}\noindent\lawboxuserandnumber\hangindent=2.6em\hangafter=1\makebox[2.6em][l]{(#1)}#2\par}
\newcommand{\nrno}[2]{\lawboxprespace{0pt}\noindent\lawboxuserandnumber\hspace*{2.6em}\hangindent=4.4em\hangafter=1\makebox[1.8em][l]{#1.}#2\par}
\newcommand{\letterno}[2]{\lawboxprespace{0pt}\noindent\lawboxuserandnumber\hspace*{4.4em}\hangindent=6.0em\hangafter=1\makebox[1.6em][l]{#1)}#2\par}
\newcommand{\sourcelawtitle}[1]{\lawboxprespace{0.22em}{\lawboxsourcefont\bfseries\lawboxuserandnumber #1}\par\lawboxtitleafter}
\newcommand{\sourcelawsubtitle}[1]{\lawboxprespace{0.04em}{\lawboxsourcefont\bfseries\lawboxuserandnumber #1}\par\lawboxtitleafter}
\newcommand{\sourceabsno}[2]{\lawboxprespace{0.04em}\noindent\lawboxuserandnumber\hangindent=2.45em\hangafter=1\makebox[2.45em][l]{(#1)}#2\par}
\newcommand{\sourcenrno}[2]{\lawboxprespace{0pt}\noindent\lawboxuserandnumber\hspace*{2.45em}\hangindent=4.25em\hangafter=1\makebox[1.8em][l]{#1.}#2\par}
\newcommand{\sourceletterno}[2]{\lawboxprespace{0pt}\noindent\lawboxuserandnumber\hspace*{4.25em}\hangindent=5.85em\hangafter=1\makebox[1.6em][l]{#1)}#2\par}
\newcommand{\sourcequotepara}[1]{\lawboxprespace{0.04em}\noindent\lawboxuserandnumber\hangindent=1.2em\hangafter=1 #1\par}
\newcommand{\sourcelawline}[1]{\lawboxprespace{0.04em}\noindent\lawboxuserandnumber #1\par}
\newcommand{\lawline}[1]{\lawboxprespace{0.04em}\noindent\lawboxuserandnumber #1\par}
\newcommand{\pendingtranslation}{\par{\songfont\footnotesize\color{pendingink}（中文譯文待補。）}\par}
\newcommand{\tocpart}[1]{\par\addvspace{0.52em}{\footnotesize\bfseries\color{pendingink}#1\par}\addvspace{0.16em}}
\newcommand{\tocdashline}{\par\addvspace{0.48em}\noindent{\color{notegray}\xleaders\hbox to 0.82em{\hfil\rule{0.42em}{0.28pt}\hfil}\hskip\linewidth}\par\addvspace{0.38em}}
% \tocpanel{font-cmd}{content}：將目錄置中於所在欄寬（雙語欄適用）。
% A3 橫式使用 0.58\linewidth，A4 橫式使用 0.84\linewidth，
% 使兩種紙張下目錄欄內容實際寬度相近（均約 100–105 mm）。
\newcommand{\tocpanel}[2]{%
  \makebox[\linewidth][c]{%
    \ifx\bilingualpaper\bilingualpaperaThree
      \begin{minipage}{0.58\linewidth}%
    \else
      \begin{minipage}{0.84\linewidth}%
    \fi
    \toccolstyle
    #1#2%
    \end{minipage}%
  }%
}
% ===== 首頁簡目錄的兩層 description list 縮排 =====
% enumitem 的 description list 可把每一列拆成：
%   [label]  labelsep  body text
% 這裡的「起點」都以所在 minipage/欄位的左緣為 0。
%
% 核心規則：
%   labelindent = label 欄位從左緣開始的位置。
%   labelwidth  = label 欄位本身寬度；標籤太長（如 Abw. Meinung）就加大它。
%   labelsep    = label 欄位右緣到正文的距離。
%   leftmargin  = 正文 body text 的起點。判斷層級視覺時主要看這個。
%
% 所以：
%   想讓「Begründung / 判決理由」整列正文往右或往左：改 tocitems 的 leftmargin。
%   想讓「A. Verfahren und Gegenstand / A. 程序與審查標的」正文往右或往左：
%     改 tocsubitems 的 leftmargin。
%   想讓 A./B./C. 這些標籤本身往右或往左，但正文不動：改 tocsubitems 的 labelindent。
%   想讓 A./B./C. 與正文間距變大或變小：改 tocsubitems 的 labelsep。
%
% 層級原則：
%   tocsubitems.leftmargin 必須大於 tocitems.leftmargin，
%   否則子層級正文會看起來比「Gründe / 理由」還靠左。
\newlist{tocitems}{description}{1}
\setlist[tocitems]{%
  labelindent=0pt,    % 第一層 label 欄位起點；0pt 表示貼齊目錄欄左緣。
  leftmargin=2.54cm,  % 第一層正文起點；例如 Begründung / 判決理由 的左緣。
  labelwidth=2.16cm,  % 第一層 label 欄位寬度；需容納 Leitsätze、Abw. Meinung、不同意見等。
  labelsep=0.32cm,    % 第一層 label 與正文的水平間距；只改間距，不改正文起點。
  itemsep=0.28em,     % 第一層各項之間的垂直距離。
  topsep=0.20em,      % 第一層 list 上下與前後內容的垂直距離。
  parsep=0pt,         % 同一 item 內段落之間的距離；目錄項通常不需要。
  partopsep=0pt,      % list 若緊接段落開始時的額外距離；關閉以免目錄高度飄動。
  align=left,         % label 欄內左對齊；長短標籤不靠右貼齊。
  font=\normalfont\bfseries % label 的字型；正文不受這行影響。
}
\newlist{tocsubitems}{description}{1}
\setlist[tocsubitems]{%
  labelindent=1cm,    % 第二層 label 起點；只控制 A./B./C. 本身的位置。
  leftmargin=3.00cm,  % 第二層正文起點；必須大於 2.54cm，才會比 Begründung / 判決理由更右。
  labelwidth=0.5cm,   % 第二層 label 欄寬；A./B./C. 很短，可小於第一層。
  labelsep=1.5cm,    % 第二層 label 與正文間距；A. 與 Verfahren / 程序 的距離看這裡。
  itemsep=0.12em,     % 第二層各項較密，避免 A.–E. 佔太高。
  topsep=0.04em,      % 第二層 list 與 Gründe / 理由之間的垂直距離。
  parsep=0pt,         % 同一第二層 item 內段落距離；目錄項通常不需要。
  partopsep=0pt,      % 關閉額外段落起始距離，避免版面不穩。
  align=left,         % A./B./C. label 欄內左對齊。
  font=\normalfont\bfseries, % A./B./C. label 加粗；正文不受這行影響。
  before={}           % 不在第二層 list 前自動插入額外內容。
}
\newlist{termitems}{description}{1}
\ifbilingualcolumns
  \setlist[termitems]{%
    leftmargin=5.2cm,
    labelwidth=4.8cm,
    labelsep=0.35cm,
    itemsep=0.16em,
    topsep=0.2em,
    parsep=0pt,
    partopsep=0pt,
    font=\normalfont\bfseries,
    style=nextline
  }
\else
  \setlist[termitems]{%
    leftmargin=0pt,
    labelwidth=0pt,
    labelsep=0pt,
    itemsep=0.24em,
    topsep=0.2em,
    parsep=0pt,
    partopsep=0pt,
    font=\normalfont\bfseries,
    style=nextline
  }
\fi
\newlist{abbritems}{description}{1}
\setlist[abbritems]{%
  leftmargin=2.38cm,
  labelwidth=2.02cm,
  labelsep=0.28cm,
  itemsep=0.16em,
  topsep=0.2em,
  parsep=0pt,
  partopsep=0pt,
  align=left,
  font=\normalfont\bfseries
}
\ifbilingualcolumns
  \newenvironment{termcolumns}{\begin{multicols}{2}}{\end{multicols}}
\else
  \newenvironment{termcolumns}{}{}
\fi

% ===== 編者註腳開關 =====
% 一般版預設不顯示編者註，避免正文腳註過密。
% 若開啟 \classroommodeon，GG/條文編者註仍會自動顯示，無需另開本開關。
% 若一般版也要顯示編者註，可在單案檔 \input 後加 \showeditornotestrue。
\newif\ifshoweditornotes
\showeditornotesfalse

% ===== 目錄排版輔助 =====
\newcommand{\toccolstyle}{\raggedright\arraybackslash}
\newcommand{\tocsingleindent}{\leftskip=1.45cm\rightskip=1.2cm}

% ===== 行內引文環境（德文原文與中文譯文各一，帶左側灰色邊線）=====
\newcommand{\quoteparbreak}{\par\addvspace{0.48em}}
\newcommand{\sourcequoteinner}[1]{%
  \begin{tcolorbox}[
    enhanced, breakable, width=\dimexpr\linewidth-0.8em\relax, frame hidden,
    colback=white, coltext=caseink, boxrule=0pt,
    borderline west={0.55pt}{0.88em}{notegray},
    left=1.78em, right=0.72em, top=0.18em, bottom=0.18em,
    boxsep=0pt, before skip=0.24em, after skip=0.24em,
    before upper={\normalfont\footnotesize\color{caseink}\setlength{\parindent}{0pt}\setlength{\parskip}{0.34em}\raggedright\setlength{\rightskip}{0pt plus 1em}}
  ]%
  #1\par
  \end{tcolorbox}%
}
\newcommand{\transquoteinner}[1]{%
  \begin{tcolorbox}[
    enhanced, breakable, width=\dimexpr\linewidth-0.8em\relax, frame hidden,
    colback=white, coltext=caseink, boxrule=0pt,
    borderline west={0.55pt}{0.88em}{notegray},
    left=1.78em, right=0.72em, top=0.18em, bottom=0.18em,
    boxsep=0pt, before skip=0.24em, after skip=0.24em,
    before upper={\songfont\small\color{caseink}\setlength{\parindent}{0pt}\setlength{\parskip}{0.34em}\raggedright\setlength{\rightskip}{0pt plus 1em}}
  ]%
  #1\par
  \end{tcolorbox}%
}

% ===== 大綱（Gliederung）Rn 輔助命令 =====
\newcommand{\outrn}[2]{\enspace{\normalfont\footnotesize\color{pendingink}(Rn\,#1–#2)}}
\newcommand{\outrnsingle}[1]{\enspace{\normalfont\footnotesize\color{pendingink}(Rn\,#1)}}
\newcommand{\outrnzh}[2]{\enspace{\normalfont\footnotesize\color{pendingink}（第\,#1–#2\,段）}}
\newcommand{\outrnzhs}[1]{\enspace{\normalfont\footnotesize\color{pendingink}（第\,#1\,段）}}
% ===== 大綱雙語對照表格輔助命令（三欄：段碼 │ 德文 │ 中文）=====
% 欄 1：段碼（由欄定義自動格式化：小字、等寬、右對齊、pendingink）
% 欄 2：德文標籤＋正文
% 欄 3：中文標籤＋正文
%
% 無段碼的列（\omainrow、\osecrow），欄 1 以 & 空置。
% 有段碼的列（\osecrowrn、\osubrow、\ointrow），#1 為預格式化字串
%   例：\osubrow{2–49}{I.}{...}{I.}{...}
%       \ointrow{107}{...}{...}
%
% \opartrow：段組標題（橫跨三欄）
\newcommand{\opartrow}[2]{%
  \noalign{\vspace{0.30em}}%
  \multicolumn{3}{@{}l@{}}{%
    \footnotesize\bfseries\color{pendingink}#1\hfill\songfont#2%
  }\\[0.06em]%
}
% \odashrow：虛線分隔（橫跨三欄）
\newcommand{\odashrow}{%
  \noalign{\vspace{0.28em}}%
  \multicolumn{3}{@{}l@{}}{\color{notegray}%
    \xleaders\hbox to 0.82em{\hfil\rule{0.42em}{0.28pt}\hfil}\hskip 0.88\linewidth}%
  \\[0.24em]%
}
% \odissentpart：不同意見書區段標頭。比一般 \opartrow 更像附錄性轉場。
\newcommand{\odissentpart}[2]{%
  \noalign{\vspace{0.42em}}%
  \multicolumn{3}{@{}p{0.88\textwidth}@{}}{%
    \color{notegray}\rule{\linewidth}{0.18pt}%
  }\\[-0.18em]%
  \multicolumn{3}{@{}p{0.88\textwidth}@{}}{%
    \makebox[\linewidth][s]{%
      \footnotesize\bfseries\color{pendingink}#1%
      \hfill{\songfont #2}%
    }%
  }\\[-0.02em]%
}
% \odissentopinion：不同意見書的署名與一行摘要，右欄保留段碼範圍。
\newcommand{\odissentopinion}[5]{%
  \footnotesize\hangindent=0.52cm\hangafter=1%
    {\bfseries\color{caseink}#2}\enspace{\color{notegray}\textperiodcentered}\enspace\color{caseink}#3%
  &\footnotesize\hangindent=0.52cm\hangafter=1%
    {\songfont\bfseries\color{caseink}#4}\enspace{\color{notegray}\textperiodcentered}\enspace\songfont\color{caseink}#5%
  &#1%
  \\[0.12em]%
}
% \omainrow：頂層項目，無段碼（Leitsätze、Urteil、Gründe …）
\newcommand{\omainrow}[4]{%
  \footnotesize\hangindent=1.92cm\hangafter=1%
    \makebox[1.92cm][l]{\bfseries\color{caseink}#1}\color{caseink}#2%
  &\footnotesize\hangindent=1.92cm\hangafter=1%
    \makebox[1.92cm][l]{\bfseries\color{caseink}#3}\songfont\color{caseink}#4%
  &%
  \\[0.15em]%
}
% \omainrowrn：頂層項目，有段碼（不同意見等標籤寬於 \osecrow 框者）
\newcommand{\omainrowrn}[5]{%
  \footnotesize\hangindent=1.92cm\hangafter=1%
    \makebox[1.92cm][l]{\bfseries\color{caseink}#2}\color{caseink}#3%
  &\footnotesize\hangindent=1.92cm\hangafter=1%
    \makebox[1.92cm][l]{\bfseries\color{caseink}#4}\songfont\color{caseink}#5%
  &#1%
  \\[0.15em]%
}
% \osecrow：一級子項，無段碼（A.、C.、D. 等有子項者）
\newcommand{\osecrow}[4]{%
  \footnotesize\hspace{1.10cm}\hangindent=1.98cm\hangafter=1%
    \makebox[0.86cm][l]{\bfseries\color{caseink}#1}\color{caseink}#2%
  &\footnotesize\hspace{1.10cm}\hangindent=1.98cm\hangafter=1%
    \makebox[0.86cm][l]{\bfseries\color{caseink}#3}\songfont\color{caseink}#4%
  &%
  \\[0.11em]%
}
% \osecrowrn：一級子項，有段碼（B.、E. 等完整段落範圍者）
% #1=段碼字串（如 101–106）, #2=DE標籤, #3=DE正文, #4=ZH標籤, #5=ZH正文
\newcommand{\osecrowrn}[5]{%
  \footnotesize\hspace{1.10cm}\hangindent=1.98cm\hangafter=1%
    \makebox[0.86cm][l]{\bfseries\color{caseink}#2}\color{caseink}#3%
  &\footnotesize\hspace{1.10cm}\hangindent=1.98cm\hangafter=1%
    \makebox[0.86cm][l]{\bfseries\color{caseink}#4}\songfont\color{caseink}#5%
  &#1%
  \\[0.11em]%
}
% \osubrow：二級子項，有段碼範圍（I.、II.、A.I. …）
% #1=段碼字串, #2=DE標籤, #3=DE正文, #4=ZH標籤, #5=ZH正文
\newcommand{\osubrow}[5]{%
  \footnotesize\hspace{1.40cm}\hangindent=2.30cm\hangafter=1%
    \makebox[0.90cm][l]{\bfseries\color{caseink}#2}\color{caseink}#3%
  &\footnotesize\hspace{1.40cm}\hangindent=2.30cm\hangafter=1%
    \makebox[0.90cm][l]{\bfseries\color{caseink}#4}\songfont\color{caseink}#5%
  &#1%
  \\[0.09em]%
}
% \ointrow：引入段落，有單一段碼，無標籤
% #1=段碼字串, #2=DE正文, #3=ZH正文
\newcommand{\ointrow}[3]{%
  \footnotesize\hspace{0.52cm}\hangindent=0.52cm\hangafter=1\color{caseink}#2%
  &\footnotesize\hspace{0.52cm}\hangindent=0.52cm\hangafter=1\songfont\color{caseink}#3%
  &#1%
  \\[0.09em]%
}
% \osubsubrow：三級子項（1.、2.、3.）
% #1=段碼字串, #2=DE標籤, #3=DE正文, #4=ZH標籤, #5=ZH正文
\newcommand{\osubsubrow}[5]{%
  \footnotesize\hspace{1.80cm}\hangindent=2.48cm\hangafter=1%
    \makebox[0.68cm][l]{\bfseries\color{caseink}#2}\color{caseink}#3%
  &\footnotesize\hspace{1.80cm}\hangindent=2.48cm\hangafter=1%
    \makebox[0.68cm][l]{\bfseries\color{caseink}#4}\songfont\color{caseink}#5%
  &#1%
  \\[0.08em]%
}
% \osubsubsubrow：四級子項（a）、b））
% #1=段碼字串, #2=DE標籤, #3=DE正文, #4=ZH標籤, #5=ZH正文
\newcommand{\osubsubsubrow}[5]{%
  \footnotesize\hspace{2.10cm}\hangindent=2.68cm\hangafter=1%
    \makebox[0.58cm][l]{\bfseries\color{caseink}#2}\color{caseink}#3%
  &\footnotesize\hspace{2.10cm}\hangindent=2.68cm\hangafter=1%
    \makebox[0.58cm][l]{\bfseries\color{caseink}#4}\songfont\color{caseink}#5%
  &#1%
  \\[0.07em]%
}
% \osubsubsubsubrow：五級子項（aa）、bb））
% 標籤框 0.62cm：足以容納「aa)」在 footnotesize bold 下（約 0.50cm）並留有間距
% #1=段碼字串, #2=DE標籤, #3=DE正文, #4=ZH標籤, #5=ZH正文
\newcommand{\osubsubsubsubrow}[5]{%
  \footnotesize\hspace{2.38cm}\hangindent=3.06cm\hangafter=1%
    \makebox[0.68cm][l]{\bfseries\color{caseink}#2}\color{caseink}#3%
  &\footnotesize\hspace{2.38cm}\hangindent=3.06cm\hangafter=1%
    \makebox[0.68cm][l]{\bfseries\color{caseink}#4}\songfont\color{caseink}#5%
  &#1%
  \\[0.06em]%
}
% \outlinepage：雙語對照大綱頁包裝環境（longtable 自動跨頁）
% 三欄：德文欄 + 中文欄 + 段碼欄（最右，不折行）
% 段碼欄寬度：1.80cm，足以容納最長段碼如「309–348」
% DE/ZH 欄寬：(0.87\textwidth - 2.08cm) / 2
%   驗算：2×欄 + 0.05\textwidth + 0.28cm gap + 1.80cm = 0.87tw - 2.08cm + 0.05tw + 2.08cm = 0.92tw ✓
\newenvironment{outlinepage}{%
  \vspace*{0.40cm}%
  \setlength{\LTleft}{0.04\textwidth}%
  \setlength{\LTright}{0.04\textwidth}%
  \setlength{\tabcolsep}{0pt}%
  % 大綱列距由 longtable 的 \arraystretch 與各 row macro 結尾的 \\[...em] 共同決定。
  % 這裡控制「每一列本身的表格 strut 高度」；數值越大，A./I./1. 各項之間越像有空行。
  % A3 原本用 0.62，視覺上沒有空行；A4 也沿用同值，避免切到 A4 後項目被撐開。
  % 若日後覺得大綱太擠，只微調這個值（例如 0.68），不要改各 \omainrow/\osubrow 的縮排。
  \renewcommand{\arraystretch}{0.62}%
  \begin{longtable}{%
    >{\vspace{0pt}\raggedright\arraybackslash}p{\dimexpr(0.87\textwidth-2.08cm)/2\relax}%
    @{\hspace{0.025\textwidth}}%
    !{\color{notegray}\vrule width 0.12pt}%
    @{\hspace{0.025\textwidth}}%
    >{\vspace{0pt}\raggedright\arraybackslash}p{\dimexpr(0.87\textwidth-2.08cm)/2\relax}%
    @{\hspace{0.28cm}}%
    >{\vspace{0pt}\footnotesize\normalfont\color{pendingink}\raggedright\arraybackslash}p{1.80cm}%
    @{}%
  }%
  % 首頁欄標籤：不要橫跨置中；分別放在德文欄與中文欄的實際欄位起點。
  \footnotesize\bfseries\color{sourceink}Gliederung%
  &\footnotesize\bfseries\songfont\color{sourceink}大綱%
  &%
  \\[0.36em]%
  \endfirsthead%
  % 續頁欄標籤：同樣分欄定位，以灰色細體區分；無需標注「續」。
  \footnotesize\color{pendingink}Gliederung%
  &\footnotesize\songfont\color{pendingink}大綱%
  &%
  \\[0.24em]%
  \endhead%
}{%
  \end{longtable}%
}
